% !TeX program = pdflatex
% Master DT. Kompilować z books/pl, po functional_analysis.tex.
\documentclass[11pt,a4paper]{article}
\usepackage[T1]{fontenc}
\usepackage[utf8]{inputenc}
\usepackage[polish]{babel}
\usepackage{lmodern}
\usepackage{amsmath,amssymb,amsthm,mathtools}
\usepackage{microtype}
\usepackage{geometry}
\usepackage{enumitem}
\usepackage{tocloft}
\usepackage{xcolor}
\usepackage{tikz}
\usetikzlibrary{arrows.meta}
\usepackage{xr-hyper}
\usepackage[unicode]{hyperref}
\geometry{margin=2.5cm}
\hypersetup{colorlinks=true,linkcolor=blue!45!black,urlcolor=blue!45!black,pdftitle={Teoria dystrybucji i analiza harmoniczna}}
\setlength{\parindent}{0pt}
\setlength{\parskip}{4pt}
\setlength{\emergencystretch}{2em}
\setlength{\cftsecnumwidth}{3em}
\setlength{\cftsubsecnumwidth}{4.5em}
% Moduły 15A/15B mają numery podpodsekcji w rodzaju 15B.12.8.
\setlength{\cftsubsubsecnumwidth}{6em}
\setlist[enumerate]{topsep=4pt,itemsep=2pt}
\setlist[itemize]{topsep=4pt,itemsep=2pt}
\numberwithin{equation}{section}
\theoremstyle{plain}
\newtheorem{twierdzenie}{Twierdzenie}[section]
\newtheorem{stwierdzenie}[twierdzenie]{Stwierdzenie}
\newtheorem{lemat}[twierdzenie]{Lemat}
\newtheorem{wniosek}[twierdzenie]{Wniosek}
\theoremstyle{definition}
\newtheorem{definicja}[twierdzenie]{Definicja}
\newtheorem{przyklad}[twierdzenie]{Przykład}
\theoremstyle{remark}
\newtheorem{uwaga}[twierdzenie]{Uwaga}
\newcommand{\R}{\mathbb R}
\newcommand{\C}{\mathbb C}
\newcommand{\N}{\mathbb N}
\newcommand{\K}{\mathbb K}
\newcommand{\cl}[1]{\overline{#1}}
\newcommand{\TVS}{\mathrm{TVS}}
\DeclareMathOperator{\lin}{lin}
\DeclareMathOperator{\supp}{supp}
\DeclareMathOperator{\id}{id}
\title{Teoria dystrybucji i analiza harmoniczna}
\author{}
\date{Wersja robocza po restrukturyzacji -- 28 listopada 2026}

% BEGIN GLOBAL CODE LISTINGS (generated from latex/pl/code-listings.tex)
\usepackage{xcolor}
\usepackage{listings}
\definecolor{codebg}{RGB}{248,249,251}
\definecolor{codeborder}{RGB}{180,191,203}
\definecolor{codecomment}{RGB}{61,113,79}
\definecolor{codekeyword}{RGB}{41,78,132}
\lstdefinelanguage{Julia}{
  morekeywords={abstract,baremodule,begin,break,catch,const,continue,do,else,elseif,
    end,export,false,finally,for,function,global,if,import,in,isa,let,local,macro,
    module,mutable,primitive,quote,return,struct,true,try,type,using,where,while,
    missing,nothing,println,print},
  sensitive=true,morecomment=[l]{\#},morecomment=[n]{\#=}{=\#},
  morestring=[b]",
}
\lstdefinestyle{skrypt}{
  backgroundcolor=\color{codebg},frame=single,rulecolor=\color{codeborder},
  framesep=5pt,frameround=tttt,framerule=.5pt,
  basicstyle=\ttfamily\footnotesize,keywordstyle=\color{codekeyword}\bfseries,
  commentstyle=\color{codecomment}\itshape,stringstyle=\color{purple!65!black},
  numbers=left,numberstyle=\tiny\color{gray!75!black},numbersep=8pt,
  stepnumber=1,breaklines=true,breakatwhitespace=false,
  columns=fullflexible,keepspaces=true,showstringspaces=false,tabsize=4,
  linewidth=.90\linewidth,captionpos=b,aboveskip=11pt,belowskip=12pt,
  xleftmargin=12pt,xrightmargin=2pt,
}
\lstdefinestyle{pythonpde}{style=skrypt,language=Python}
\lstdefinestyle{python}{style=skrypt,language=Python}
\lstdefinestyle{julia}{style=skrypt,language=Julia}
\lstset{style=skrypt}
\renewcommand{\lstlistingname}{Kod}
\renewcommand{\lstlistlistingname}{Spis kodów}
% END GLOBAL CODE LISTINGS

% AF pozostaje odrębnym źródłem trzech wcześniejszych wyników.
% Kompilacja: najpierw functional_analysis.tex, następnie ten skrypt.
% Prefiks zapobiega kolizji z 211 etykietami własnych rozdziałów.
\externaldocument[AF-ext-]{functional_analysis}
\makeatletter
\newcommand{\importFAreference}[1]{%
  \@ifundefined{r@AF-ext-#1}{%
    \PackageWarningNoLine{distribution-theory}{Missing AF reference #1; compile functional_analysis first}%
  }{\expandafter\let\csname r@#1\expandafter\endcsname\csname r@AF-ext-#1\endcsname}}
\importFAreference{cor:frechet-barrelled}
\importFAreference{thm:banach-steinhaus-lcs}
\importFAreference{thm:point-closed-subspace}
\makeatother
\providecommand{\atlasref}[2]{#2}
\begin{document}
\maketitle
\noindent Numeracja docelowa: pominięto rozdziały, których materiał jest jeszcze w przygotowaniu.\par
\tableofcontents
\clearpage
\setcounter{part}{0}
\part{Podstawy teorii dystrybucji}
\setcounter{section}{0}
% Wersja źródłowa: audyt II, jednostka 17

\section{Przestrzenie funkcji gładkich i przestrzeń Schwartza}
\label{sec:gladkie-schwartz17}

\providecommand{\atlasref}[2]{#2}


\subsection{Topologia przestrzeni \texorpdfstring{$C^\infty(\Omega)$}{Cinfty(Omega)}}

Niech $d\ge1$, niech $\Omega\subseteq\R^d$ będzie niepustym zbiorem
otwartym, a $\K\in\{\R,\C\}$. Dla wielowskaźnika
$\alpha=(\alpha_1,\dots,\alpha_d)\in\N_0^d$ oznaczamy
$|\alpha|=\sum_j\alpha_j$, $\alpha!=\prod_j\alpha_j!$,
$x^\alpha=\prod_jx_j^{\alpha_j}$ oraz
$\partial^\alpha=\partial_1^{\alpha_1}\cdots\partial_d^{\alpha_d}$.
Dla $\gamma,\beta\in\N_0^d$ piszemy $\gamma\le\beta$,
gdy $\gamma_j\le\beta_j$ dla każdego $j$;
wtedy $\binom{\beta}{\gamma}=\prod_{j=1}^d\binom{\beta_j}{\gamma_j}$.
Symbole $e_i$ oznaczają wektory standardowej bazy $\R^d$,
traktowane także jako wielowskaźniki. Piszemy $K\Subset\Omega$,
gdy $K$ jest zwartym podzbiorem $\Omega$.
Przez $C^\infty(\Omega;\K)$ rozumiemy przestrzeń funkcji
$\Omega\to\K$ posiadających ciągłe pochodne cząstkowe wszystkich rzędów.
Przez $|x|$ dla $x\in\R^d$ rozumiemy normę euklidesową,
a przez $\lambda_d$ miarę Lebesgue'a na $\R^d$.
Nie narzucamy ograniczoności funkcji ani ich pochodnych na całym $\Omega$.

\begin{definicja}[Seminormy zbieżności lokalnej wszystkich pochodnych]
\label{def:seminormy-cinf17}
Niech $d\ge1$, $\Omega\subseteq\R^d$ będzie niepustym zbiorem
otwartym, a $\K=\R$ lub $\K=\C$.
Dla zwartego $K\subseteq\Omega$ i liczby $m\in\N_0$ określamy
\[
 p_{K,m}(f)=\max_{|\alpha|\le m}\sup_{x\in K}|\partial^\alpha f(x)|,
 \qquad f\in C^\infty(\Omega;\K).
\]
Dla $K=\varnothing$ przyjmujemy $p_{K,m}=0$.
Nierówność trójkąta i jednorodność wynikają z odpowiednich
własności wartości bezwzględnej, supremum i maksimum.
Topologią standardową na $C^\infty(\Omega;\K)$ nazywamy
topologię generowaną przez rodzinę tych seminorm.
Sieć $(f_\iota)$ zbiega do $f$ w tej topologii wtedy i tylko wtedy,
gdy dla każdego zwartego $K\subseteq\Omega$ i każdego
$\alpha\in\N_0^d$ pochodne $\partial^\alpha f_\iota$
zbiegają jednostajnie do $\partial^\alpha f$ na $K$.
\end{definicja}

\begin{lemat}[Wyczerpanie zbiorami zwartymi]
\label{lem:exhaustion17}
Dla każdego niepustego otwartego $\Omega\subseteq\R^d$ istnieje
ciąg zbiorów zwartych $(K_n)_{n\ge1}$ taki, że
$K_n\subseteq\operatorname{int}K_{n+1}$, $\bigcup_nK_n=\Omega$
i każdy zwarty $K\subseteq\Omega$ jest zawarty w pewnym $K_n$.
\end{lemat}
\begin{proof}
Dla $\Omega\ne\R^d$ weźmy
$K_n=\{x\in\Omega:|x|\le n,\
\operatorname{dist}(x,\R^d\setminus\Omega)\ge1/n\}$;
dla $\Omega=\R^d$ przyjmijmy $K_n=\{x:|x|\le n\}$.
Zbiory te są zwarte: warunek odległości jest domknięty,
a dodatnia odległość od dopełnienia wymusza przynależność do $\Omega$.
Nierówności definiujące $K_n$ są ściślejsze niż analogiczne dla
$K_{n+1}$, stąd $K_n\subseteq\operatorname{int}K_{n+1}$.
Każdy $x\in\Omega$ ma dodatnią odległość od dopełnienia
i skończoną normę, więc należy do pewnego $K_n$.
Jeżeli zwarty $K\subseteq\Omega$ nie byłby zawarty w żadnym $K_n$,
można byłoby wybrać $x_n\in K\setminus K_n$.
Ze zwartości $K$ istnieje podciąg $(x_{n_j})$ zbieżny
do $x\in K$. Pewne $K_N$ zawiera $x$ w swoim wnętrzu,
bo $x\in K_{N-1}\subseteq\operatorname{int}K_N$
dla dostatecznie dużego $N$.
Dla dużych $j$ mamy $n_j\ge N$ i $x_{n_j}\in K_N\subseteq K_{n_j}$,
sprzeczność. Gdy dopełnienie jest puste, argument jest identyczny.
\end{proof}

\begin{twierdzenie}[Przeliczalna rodzina seminorm]
\label{thm:countable-cinf17}
Niech $\Omega\subseteq\R^d$ będzie niepusty i otwarty,
a $(K_n)$ ciągiem z Lematu~\ref{lem:exhaustion17}.
Wówczas topologia z Definicji~\ref{def:seminormy-cinf17}
jest generowana przez rosnący ciąg seminorm
$q_n=p_{K_n,n}$, $n\ge1$ (w szczególności $q_n=0$,
jeżeli $K_n=\varnothing$). Jest ona metryzowalna i Hausdorffa.
Ciąg $(f_j)$ zbiega w niej do $f$ wtedy i tylko wtedy, gdy
$q_n(f_j-f)\to0$ dla każdego $n\ge1$.
\end{twierdzenie}
\begin{proof}
Każde $q_n$ jest jedną z seminorm $p_{K,m}$.
Z drugiej strony dla ustalonych $K\Subset\Omega$ i $m$
wybierzmy $n\ge m$ z $K\subseteq K_n$.
Wtedy $p_{K,m}\le q_n$, więc obie rodziny wyznaczają tę samą
topologię. Ponieważ $K_n\subseteq K_{n+1}$, seminormy $q_n$
są niemalejące. Jeśli $f\ne0$, to $f(x)\ne0$ dla pewnego
$x\in\Omega$ i $x\in K_n$ dla pewnego $n$, zatem $q_n(f)>0$.
Topologia jest więc Hausdorffa. Jest metryzowalna na mocy
kryterium dla przeliczalnej rodziny seminorm; można użyć metryki
$\sum_{n\ge1}2^{-n}\min\{1,q_n(f-g)\}$.
Z definicji topologii wyznaczonej przez seminormy wynika
podane kryterium zbieżności.
\end{proof}

\begin{twierdzenie}[Zupełność $C^\infty(\Omega)$]
\label{thm:frechet-cinf17}
Dla każdego niepustego otwartego $\Omega\subseteq\R^d$ przestrzeń
$C^\infty(\Omega;\K)$ ze standardową topologią jest przestrzenią
Frécheta.
\end{twierdzenie}
\begin{proof}
Metryzowalność, lokalna wypukłość i warunek Hausdorffa wynikają
z Twierdzenia~\ref{thm:countable-cinf17}; pozostaje zupełność.
Niech $(f_j)$ będzie ciągiem Cauchy'ego dla wszystkich $q_n$.
Dla każdego wielowskaźnika $\alpha$ pochodne
$\partial^\alpha f_j$ są jednostajnie Cauchy'ego na każdym
$K_n$, gdy $n\ge|\alpha|$. Z zupełności przestrzeni
$C(K_n;\K)$ mają granice jednostajne $g_{\alpha,n}$.
Na przecięciach niepustych $K_n$ granice się zgadzają
na mocy jednoznaczności granicy punktowej. Określają więc
funkcje $g_\alpha:\Omega\to\K$, ciągłe, ponieważ każdy punkt
$\Omega$ leży we wnętrzu pewnego $K_n$ z $n\ge|\alpha|$; w szczególności
$f_j\to g_0$ lokalnie jednostajnie.

Wykażemy, że $g_\alpha=\partial^\alpha g_0$.
Ustalmy $x\in\Omega$ i zamkniętą kulę $\overline B(x,r)\Subset\Omega$.
Dla $i\in\{1,\dots,d\}$ i $|t|<r$, z twierdzenia
Newtona--Leibniza na odcinku od $x$ do $x+te_i$ zachodzi
\[
 \partial^\alpha f_j(x+te_i)-\partial^\alpha f_j(x)
 =\int_0^t\partial^{\alpha+e_i}f_j(x+se_i)\,\mathrm ds.
\]
Wybierając $K_n$ zawierający tę kulę i $n\ge|\alpha|+1$,
możemy przejść do granicy po obu stronach: błąd całki
nie przekracza $|t|\sup_{K_n}|\partial^{\alpha+e_i}f_j-
g_{\alpha+e_i}|$, który dąży do zera. Dla $t\ne0$ mamy zatem
\[
\frac{g_\alpha(x+te_i)-g_\alpha(x)}{t}
=\frac1t\int_0^t g_{\alpha+e_i}(x+se_i)\,\mathrm ds.
\]
Z ciągłości $g_{\alpha+e_i}$ prawa strona dąży do
$g_{\alpha+e_i}(x)$ przy $t\to0$, również dla $t<0$.
Stąd $\partial_i g_\alpha=g_{\alpha+e_i}$ na $\Omega$.
Indukcja względem $|\alpha|$ daje $g_0\in C^\infty(\Omega)$
i wszystkie wymagane tożsamości. Jednostajna zbieżność
pochodnych na $K_n$ oznacza $q_n(f_j-g_0)\to0$ dla każdego $n$.
Ciąg jest Cauchy'ego dla metryki z poprzedniego dowodu
wtedy i tylko wtedy, gdy jest Cauchy'ego dla każdej $q_n$:
w jedną stronę wystarczy rozważyć $n$-ty składnik sumy,
w drugą najpierw oszacować ogon szeregu przez $\sum_{n>N}2^{-n}$,
a następnie skończoną sumę dla $n\le N$.
Wykazaliśmy zatem zupełność tej metryki zgodnej z topologią
przestrzeni i niezmienniczej względem przesunięć.
\end{proof}

\begin{stwierdzenie}[Operacje i ograniczoność]
\label{prop:operations-cinf17}
Niech $\Omega\subseteq\R^d$ będzie niepusty i otwarty.
W przestrzeni $C^\infty(\Omega;\K)$ różniczkowanie
$\partial^\alpha$ jest operatorem liniowym ciągłym,
a mnożenie $(f,g)\mapsto fg$ jest odwzorowaniem dwuliniowym
łącznie ciągłym. Zbiór $B\subseteq C^\infty(\Omega;\K)$
jest ograniczony w sensie topologii przestrzeni liniowej
topologicznej wtedy i tylko wtedy, gdy
$\sup_{f\in B}p_{K,m}(f)<\infty$ dla wszystkich zwartych
$K\Subset\Omega$ i wszystkich $m\in\N_0$.
\end{stwierdzenie}
\begin{proof}
Z nierówności $p_{K,m}(\partial^\alpha f)\le p_{K,m+|\alpha|}(f)$
wynika ciągłość różniczkowania. Ze wzoru Leibniza
$\partial^\beta(fg)=\sum_{\gamma\le\beta}\binom{\beta}{\gamma}
\partial^\gamma f\,\partial^{\beta-\gamma}g$
oraz $\sum_{\gamma\le\beta}\binom{\beta}{\gamma}=2^{|\beta|}$
wynika $p_{K,m}(fg)\le 2^m p_{K,m}(f)p_{K,m}(g)$,
co dowodzi łącznej ciągłości mnożenia w zerze;
dwuliniowość i ciągłość przy dowolnej parze wynikają przez
rozwinięcie $(f+h)(g+k)-fg$. Kryterium ograniczoności
wynika z ogólnego kryterium seminormowego dla przestrzeni
lokalnie wypukłych: rodzina obrazów zbioru ograniczonego przez
każdą ciągłą seminormę jest ograniczona w $\R$ i odwrotnie.
\end{proof}

\subsection{Funkcje odcięcia i przestrzeń dualna topologiczna}

\noindent\emph{Powiązania tematyczne w Bibliotece matematycznej:} \atlasref{ME-10::thm:cc-in-lp-10}{ME-10: gęstość gładkich funkcji o zwartym nośniku w przestrzeniach funkcji całkowalnych}.


\begin{lemat}[Gładka funkcja odcięcia]
\label{lem:cutoffs17}
Jeżeli $K\subseteq\Omega\subseteq\R^d$, gdzie $K$ jest zwarty,
a $\Omega$ otwarty, to istnieje $\chi\in C_c^\infty(\Omega;\R)$
taka, że $0\le\chi\le1$ i $\chi=1$ w pewnym otoczeniu $K$.
Tu $C_c^\infty$ oznacza funkcje gładkie o zwartym nośniku.
\end{lemat}
\begin{proof}
Jeżeli $K=\varnothing$, wystarczy funkcja zerowa.
Dla $K\ne\varnothing$ wybierzmy otwarte zbiory $V,W$ takie, że
$K\subseteq V$, $\overline V\subseteq W$, $\overline W\Subset\Omega$.
Dla każdego $a\in K$ wybieramy $r_a>0$ tak,
aby $\overline B(a,3r_a)\subseteq\Omega$. Ze zwartości $K$
wybieramy skończone pokrycie kulami $B(a_j,r_{a_j})$.
Jako $V$ bierzemy sumę tych kul, a jako $W$ sumę kul
$B(a_j,2r_{a_j})$. Domknięcia skończonych sum są sumami
domknięć, więc otrzymujemy żądane zawierania.
Odległości $\operatorname{dist}(\overline V,\R^d\setminus W)$
i $\operatorname{dist}(K,\R^d\setminus V)$ są dodatnie,
ponieważ $\overline V$ i $K$ są zwarte
i rozłączne z odpowiednimi zbiorami domkniętymi.
Wybieramy $\delta>0$ mniejszą od obu odległości, tak że zbiór
$\{x:\operatorname{dist}(x,\overline V)\le\delta\}$
leży w $W$ oraz $\operatorname{dist}(K,\R^d\setminus V)>\delta$.
Niech $\rho\in C_c^\infty(\R^d)$ będzie nieujemna,
niech $\operatorname{supp}\rho\subseteq B(0,\delta)$ i
$\int_{\R^d}\rho(x)\,\mathrm d\lambda_d(x)=1$.
Taką funkcję otrzymujemy przez unormowanie
$x\mapsto\exp[-1/(1-4|x|^2/\delta^2)]$ dla $|x|<\delta/2$,
przedłużonej zerem dla $|x|\ge\delta/2$.
Gładkość przedłużenia wynika z faktu, że każda pochodna
funkcji $e^{-1/t}$ dla $t>0$ jest wielomianem w $1/t$
pomnożonym przez $e^{-1/t}$, a więc dąży do zera dla $t\downarrow0$.
Nośnik leży w $\overline B(0,\delta/2)\subset B(0,\delta)$.
Definiujemy $\chi(x)=\int_{\R^d}\chi_V(y)\rho(x-y)\,\mathrm d\lambda_d(y)$,
gdzie $\chi_V$ jest funkcją charakterystyczną $V$.
Przy każdym różniczkowaniu ilorazy różnicowe pochodnych
$\rho(x-y)$ są ograniczone przez supremum odpowiedniej
pochodnej wyższego rzędu, na mocy twierdzenia o wartości średniej.
Po pomnożeniu przez $\chi_V$ otrzymujemy całkowalną majorantę,
gdyż $V$ jest ograniczony. Z twierdzenia o zbieżności
zmajoryzowanej, indukcyjnie,
$\partial^\alpha\chi(x)=\int_V\partial^\alpha\rho(x-y)\,\mathrm d\lambda_d(y)$;
ciągłość tych pochodnych wynika z tego samego twierdzenia.
Stąd $\chi\in C^\infty(\R^d;\R)$.
Po podstawieniu $z=x-y$ mamy
$\chi(x)=\int_{\R^d}\chi_V(x-z)\rho(z)\,\mathrm d\lambda_d(z)$.
Ponieważ $0\le\chi_V\le1$, $\rho\ge0$ i całka $\rho$ jest równa $1$,
otrzymujemy $0\le\chi(x)\le1$.
Jeżeli $\operatorname{dist}(x,K)<\delta/2$ i
$z\in\operatorname{supp}\rho$, to
$\operatorname{dist}(x-z,K)\le\operatorname{dist}(x,K)+|z|<\delta$.
Z wyboru $\delta$ wynika $x-z\in V$, więc w tym otoczeniu
$K$ całka jest równa $1$, czyli $\chi(x)=1$. Jeżeli $\operatorname{dist}(x,\overline V)>\delta/2$,
to $\rho(x-y)=0$ dla każdego $y\in V$.
Stąd $\operatorname{supp}\chi\subseteq
\{x:\operatorname{dist}(x,\overline V)\le\delta/2\}$;
jest to zbiór zwarty zawarty w $W\Subset\Omega$,
co kończy dowód.
\end{proof}

\begin{wniosek}[Gęstość funkcji o zwartym nośniku]
\label{cor:density-cinf17}
Dla każdego niepustego otwartego $\Omega\subseteq\R^d$
przestrzeń $C_c^\infty(\Omega;\K)$ jest gęsta w
$C^\infty(\Omega;\K)$ ze standardową topologią.
\end{wniosek}
\begin{proof}
Dla wyczerpania $(K_n)$ wybierzmy z Lematu~\ref{lem:cutoffs17}
$\chi_n\in C_c^\infty(\Omega)$ równą $1$ na otoczeniu $K_n$.
Dla każdego $f\in C^\infty(\Omega)$ funkcje $\chi_n f$
mają zwarty nośnik. Dla ustalonego $K_m$ i dowolnego
rzędu pochodnych mamy $\chi_n f=f$ na otoczeniu $K_m$
dla wszystkich $n\ge m$. Zatem $\chi_n f\to f$
w topologii $C^\infty(\Omega)$.
\end{proof}

\begin{twierdzenie}[Kryterium ciągłości funkcjonału]
\label{thm:dual-cinf17}
Niech $\Omega\subseteq\R^d$ będzie niepusty i otwarty,
a $T:C^\infty(\Omega;\K)\to\K$ liniowy.
Następujące warunki są równoważne:
\begin{enumerate}
\item $T$ jest ciągły;
\item istnieją zwarty $K\Subset\Omega$, $m\in\N_0$ i $C>0$ takie, że
$|T(f)|\le C\max_{|\alpha|\le m}\sup_{x\in K}|\partial^\alpha f(x)|$
dla każdego $f\in C^\infty(\Omega;\K)$.
\end{enumerate}
W szczególności ciągły funkcjonał $T$ zależy tylko od wartości
pochodnych rzędów co najwyżej $m$ na pewnym zwartym zbiorze $K$:
jeżeli $\partial^\alpha f|_K=\partial^\alpha g|_K$ dla każdego
$|\alpha|\le m$, to $T(f)=T(g)$.
W szczególności jeśli $f$ znika w otoczeniu $K$, to $T(f)=0$.
\end{twierdzenie}
\begin{proof}
Wybierzmy rosnącą rodzinę seminorm $q_n$ z
Twierdzenia~\ref{thm:countable-cinf17}.
Ciągłość $T$ w zerze oznacza, że pewne otoczenie
$\{f:q_n(f)<\varepsilon\}$ jest odwzorowywane w
$\{z:|z|<1\}$. Jeśli $q_n(f)>0$, to dla każdego $0<a<\varepsilon/q_n(f)$
mamy $q_n(af)<\varepsilon$, zatem $a|T(f)|<1$.
Przechodząc z $a$ do $\varepsilon/q_n(f)$ od dołu,
otrzymujemy $|T(f)|\le\varepsilon^{-1}q_n(f)$.
Jeśli $q_n(f)=0$, to $q_n(af)=0$ dla każdego $a>0$,
więc $a|T(f)|<1$ dla wszystkich $a>0$, co daje $T(f)=0$.
To dowodzi warunku 2 z $K=K_n$ i $m=n$.
W razie potrzeby zwiększamy $n$, aby $K_n$ był niepusty;
nierówność pozostaje prawdziwa, gdyż $q_n$ są niemalejące.
Odwrotna implikacja wynika z ciągłości seminormy $p_{K,m}$.
Jeżeli wskazane pochodne $f$ i $g$ zgadzają się na $K$,
to $p_{K,m}(f-g)=0$, więc z liniowości
$|T(f)-T(g)|=|T(f-g)|\le Cp_{K,m}(f-g)=0$.
Jeśli $f=0$ na otoczeniu $K$, wszystkie jej pochodne
znikają na $K$ i można przyjąć $g=0$.
\end{proof}

\begin{przyklad}[Ewaluacja i pochodne ewaluacji]
\label{ex:evaluation-cinf17}
Dla $a\in\Omega$ i $\alpha\in\N_0^d$ funkcjonały
$T_{a,\alpha}(f)=(-1)^{|\alpha|}\partial^\alpha f(a)$
są ciągłe na $C^\infty(\Omega;\K)$, ponieważ
$|T_{a,\alpha}(f)|\le p_{\{a\},|\alpha|}(f)$.
Dla $\alpha=0$
przyjmują postać ewaluacji $f\mapsto f(a)$.
W teorii dystrybucji pojawią się jako pochodne
dystrybucji Diraca. Obecnie są po prostu przykładami
elementów przestrzeni dualnej topologicznej.
\end{przyklad}

\subsection{Przestrzeń Schwartza \texorpdfstring{$\mathcal S(\R^d)$}{S(Rd)}}

\noindent\emph{Powiązania tematyczne w Bibliotece matematycznej:} \atlasref{PDE-08::def:pde8-schwartz}{PDE-08: przestrzeń Schwartza w analizie równań różniczkowych}.


\begin{definicja}[Przestrzeń Schwartza]
\label{def:schwartz17}
Dla $d\ge1$ i $\K\in\{\R,\C\}$ definiujemy
\[
\mathcal S(\R^d;\K)=\left\{f\in C^\infty(\R^d;\K):
\sup_{x\in\R^d}|x^\alpha\partial^\beta f(x)|<\infty
\ \text{dla wszystkich }\alpha,\beta\in\N_0^d\right\}.
\]
Jej standardową topologię określa rodzina seminorm
$s_{\alpha,\beta}(f)=\sup_x|x^\alpha\partial^\beta f(x)|$.
Równoważną rodzinę tworzą
\[
 Q_{N,m}(f)=\max_{|\beta|\le m}\sup_{x\in\R^d}
 (1+|x|)^N|\partial^\beta f(x)|,
 \qquad N,m\in\N_0.
\]
Przyjmujemy też rosnący ciąg $Q_n=Q_{n,n}$, $n\ge1$.
\end{definicja}

\begin{lemat}[Równoważność seminorm Schwartza]
\label{lem:schwartz-seminorms17}
Seminormy $s_{\alpha,\beta}$ oraz $Q_{N,m}$ generują tę samą
topologię. Ponadto $Q_n$ są niemalejące, oddzielają punkty
i wyznaczają topologię przestrzeni Schwartza.
\end{lemat}
\begin{proof}
Ponieważ $|x^\alpha|\le |x|^{|\alpha|}\le(1+|x|)^{|\alpha|}$,
mamy $s_{\alpha,\beta}\le Q_{|\alpha|,|\beta|}$.
Z drugiej strony nierówność
$(1+|x|)^N\le C_{d,N}\bigl(1+\sum_{|\alpha|=N}|x^\alpha|\bigr)$
dla $N\ge1$ wynika z oszacowania
$(1+r)^N\le 2^N\max\{1,r^N\}$ dla $r\ge0$ oraz
$|x|^N\le d^{N/2}\sum_{j=1}^d|x_j|^N$.
Dla $N=0$ wystarczy stała $1$. W szczególności dla $N\ge1$
\[
Q_{N,m}(f)\le 2^N\sum_{|\beta|\le m}
\left(s_{0,\beta}(f)+d^{N/2}\sum_{j=1}^d s_{Ne_j,\beta}(f)\right).
\]
Dla $N=0$ mamy $Q_{0,m}=\max_{|\beta|\le m}s_{0,\beta}$.
Obie rodziny seminorm wyznaczają więc tę samą topologię.
Rodzina $Q_n$ jest współkońcowa względem $Q_{N,m}$,
gdy $n\ge\max\{N,m\}$. Oddzielanie punktów wynika już z $Q_1$.
\end{proof}

\begin{twierdzenie}[Zupełność przestrzeni Schwartza]
\label{thm:schwartz-frechet17}
Dla $\K=\R$ lub $\K=\C$ przestrzeń $\mathcal S(\R^d;\K)$
ze standardową topologią jest przestrzenią Frécheta.
Naturalne włożenie $\mathcal S(\R^d;\K)\hookrightarrow
C^\infty(\R^d;\K)$ jest ciągłe.
\end{twierdzenie}
\begin{proof}
Metryzowalność, lokalna wypukłość i warunek Hausdorffa wynikają z
Lematu~\ref{lem:schwartz-seminorms17}.
Niech $(f_j)$ będzie ciągiem Cauchy'ego dla wszystkich $Q_n$.
Ponieważ na każdym zwartym zbiorze i dla każdego rzędu pochodnych
seminormy $C^\infty$ są kontrolowane przez odpowiednie $Q_n$,
ciąg jest Cauchy'ego w $C^\infty(\R^d)$.
Na mocy Twierdzenia~\ref{thm:frechet-cinf17} istnieje
$f\in C^\infty(\R^d)$, do której $f_j$ zbiega lokalnie
wraz ze wszystkimi pochodnymi.
Ustalmy $N,m$ i $\varepsilon>0$. Istnieje $j_0$ takie, że
$Q_{N,m}(f_j-f_k)\le\varepsilon$ dla $j,k\ge j_0$.
Dla ustalonego $j\ge j_0$, $x\in\R^d$ i $|\beta|\le m$
mamy $(1+|x|)^N|\partial^\beta f_j(x)-\partial^\beta f_k(x)|
\le\varepsilon$ dla każdego $k\ge j_0$.
Przechodząc z $k$ do nieskończoności przy ustalonym $x$,
a następnie biorąc supremum po $x$ i maksimum po $\beta$,
otrzymujemy $Q_{N,m}(f_j-f)\le\varepsilon$.
Ponieważ $Q_{N,m}(f_j)<\infty$, wynika stąd
$Q_{N,m}(f)<\infty$. Wszystkie takie seminormy są skończone,
więc $f\in\mathcal S$, a $f_j\to f$ w topologii Schwartza.
Ciągłość włożenia wynika z nierówności
$p_{K,m}(f)\le Q_{0,m}(f)$ dla każdego $K\Subset\R^d$.
\end{proof}

\begin{przyklad}[Funkcje w przestrzeni Schwartza]
\label{ex:schwartz-functions17}
Niech $d\ge1$ i $\K\in\{\R,\C\}$.
Każda funkcja z $C_c^\infty(\R^d;\K)$ należy do
$\mathcal S(\R^d;\K)$, ponieważ wszystkie jej pochodne
znikają poza jej zwartym nośnikiem, a na nim
ważone pochodne są ciągłe, więc ograniczone.
Dla każdego wielomianu $P:\R^d\to\K$ funkcja
$f(x)=P(x)e^{-|x|^2}$ także należy do $\mathcal S(\R^d;\K)$.
Istotnie, $\partial_j(Pe^{-|x|^2})=
(\partial_jP-2x_jP)e^{-|x|^2}$, więc przez indukcję
każda pochodna ma postać wielomianu razy $e^{-|x|^2}$.
Po przemnożeniu przez dowolny jednomian otrzymujemy
funkcję ograniczoną, gdyż $(1+r)^k e^{-r^2}\to0$
dla każdego $k\in\N_0$ przy $r\to\infty$.
Natomiast $f(x)=(1+|x|^2)^{-1}$ nie należy do $\mathcal S$:
dla $t\ge1$ mamy $t^3|f(te_1)|=t^3/(1+t^2)\ge t/2$,
zatem $s_{3e_1,0}(f)=+\infty$.
\end{przyklad}

\begin{stwierdzenie}[Różniczkowanie, wielomiany i normy $L^p$]
\label{prop:schwartz-operators17}
Dla każdego $\alpha\in\N_0^d$ operatory
$f\mapsto\partial^\alpha f$ oraz $f\mapsto x^\alpha f$
są ciągłymi operatorami liniowymi na $\mathcal S(\R^d;\K)$.
Mnożenie $\mathcal S\times\mathcal S\to\mathcal S$ jest
łącznie ciągłe. Dla każdego $1\le p\le\infty$
kanoniczne włożenie $\mathcal S(\R^d;\K)\to L^p(\R^d;\K)$
jest ciągłe.
\end{stwierdzenie}
\begin{proof}
Różniczkowanie spełnia
$Q_{N,m}(\partial^\alpha f)\le Q_{N,m+|\alpha|}(f)$.
Dla mnożenia przez $x^\alpha$ wzór Leibniza daje
\[
\partial^\beta(x^\alpha f)=
\sum_{\substack{\gamma\le\beta\\\gamma\le\alpha}}
\binom{\beta}{\gamma}\frac{\alpha!}{(\alpha-\gamma)!}
x^{\alpha-\gamma}\partial^{\beta-\gamma}f.
\]
Ponieważ $|x^{\alpha-\gamma}|\le(1+|x|)^{|\alpha|}$,
po pomnożeniu przez wagę i wzięciu supremum otrzymujemy
$Q_{N,m}(x^\alpha f)\le C_{\alpha,m}Q_{N+|\alpha|,m}(f)$,
gdzie stała jest maksimum sum współczynników powyższego wzoru
dla $|\beta|\le m$. Analogiczny wzór daje
$Q_{N,m}(fg)\le C_{d,m}Q_{N,m}(f)Q_{0,m}(g)$,
a więc łączną ciągłość iloczynu.
Dla $p=\infty$ mamy $\|f\|_\infty\le Q_{0,0}(f)$.
Dla $1\le p<\infty$ wybieramy liczbę całkowitą $N>d/p$.
Wtedy
$\|f\|_{L^p}^p\le Q_{N,0}(f)^p
\int_{\R^d}(1+|x|)^{-Np}\,\mathrm d\lambda_d(x)$,
a całka jest skończona: na kuli jednostkowej całkujemy
funkcję ograniczoną, zaś w pierścieniu
$2^k\le|x|<2^{k+1}$, $k\ge0$, całka jest nie większa niż
$\lambda_d(B(0,1))2^d\,2^{k(d-Np)}$.
Szereg tych oszacowań jest zbieżny, bo $Np>d$.
Po wyciągnięciu pierwiastka stopnia $p$ dostajemy
żądane oszacowanie normy, a więc ciągłość włożenia.
\end{proof}

\begin{twierdzenie}[Gęstość funkcji o zwartym nośniku w $\mathcal S$]
\label{thm:density-schwartz17}
Przestrzeń $C_c^\infty(\R^d;\K)$ jest gęsta w
$\mathcal S(\R^d;\K)$ dla standardowej topologii Schwartza.
\end{twierdzenie}
\begin{proof}
Ustalmy $f\in\mathcal S(\R^d;\K)$.
Z Lematu~\ref{lem:cutoffs17} dla $K=\overline B(0,1)$
i $\Omega=B(0,2)$ wybieramy funkcję odcięcia, którą
przedłużamy zerem na $\R^d$ (jej nośnik jest zwarto zawarty
w $B(0,2)$, więc przedłużenie jest gładkie).
Otrzymujemy $\chi\in C_c^\infty(\R^d;\R)$ taką, że
$0\le\chi\le1$, $\chi=1$ na $\overline B(0,1)$
i $\operatorname{supp}\chi\subseteq B(0,2)$.
Niech $\chi_R(x)=\chi(x/R)$ dla $R\ge1$.
Funkcja $\chi_R f$ ma zwarty nośnik. Dla wielowskaźnika $\beta$
ze wzoru Leibniza
\[
\partial^\beta((1-\chi_R)f)
=(1-\chi_R)\partial^\beta f+
\sum_{0<\gamma\le\beta}\binom{\beta}{\gamma}
 (-1)R^{-|\gamma|}(\partial^\gamma\chi)(x/R)
 \partial^{\beta-\gamma}f.
\]
Nośnik pierwszego składnika jest zawarty w $\{|x|\ge R\}$,
a nośniki pozostałych w $\{R\le|x|\le2R\}$.
Dla $|\beta|\le m$ na nośniku każdego składnika mamy
$(1+|x|)^{-1}\le(1+R)^{-1}$. Ponadto
$|1-\chi_R|\le1$ i $R^{-|\gamma|}\le1$.
Po pomnożeniu przez $(1+|x|)^N$ pierwszy składnik jest
ograniczony przez $(1+R)^{-1}Q_{N+1,m}(f)$, a składnik
sumy o indeksie $\gamma$ przez tę samą wielkość
pomnożoną przez $\binom{\beta}{\gamma}\|\partial^\gamma\chi\|_\infty$.
Sumując, otrzymujemy
$Q_{N,m}((1-\chi_R)f)\le C_m(1+R)^{-1}Q_{N+1,m}(f)$,
gdzie $C_m$ jest maksimum odpowiednich skończonych sum
dla $|\beta|\le m$ i nie zależy od $R$ ani $f$.
Stąd $Q_{N,m}(f-\chi_Rf)\to0$ przy $R\to\infty$.
\end{proof}

\begin{stwierdzenie}[Brak normowalności]
\label{prop:non-normable17}
Jeśli $d\ge1$ i $\Omega\subseteq\R^d$ jest niepusty i otwarty,
to ani $C^\infty(\Omega;\K)$, ani $\mathcal S(\R^d;\K)$
ze swoimi standardowymi topologiami nie są normowalne.
\end{stwierdzenie}
\begin{proof}
Wystarczy pokazać, że żadne otoczenie zera nie jest ograniczone,
a następnie zastosować kryterium normowalności Kolmogorowa.
Najpierw rozważmy funkcje zespolone. Każde otoczenie zera
w $C^\infty(\Omega)$ zawiera zbiór $U=\{f:q_N(f)<\varepsilon\}$
dla pewnych $N\ge1$ i $\varepsilon>0$.
Wybierzmy kulę domkniętą o dodatnim promieniu zawartą w $\Omega$
i, z Lematu~\ref{lem:cutoffs17}, funkcję
$\psi\in C_c^\infty(\Omega;\R)$ równą $1$ w jej otoczeniu.
Jej nośnik zawiera się w pewnym $K_M$ z $M>N+1$.
Niech $f_j(x)=c j^{-N}\psi(x)e^{ijx_1}$.
Wzór Leibniza zapewnia $q_N(f_j)\le Cc$,
więc małym $c>0$ można zagwarantować $f_j\in U$.
Na wnętrzu wybranej kuli wszystkie pochodne $\psi$
dodatniego rzędu znikają, zatem
$\partial_1^{N+1}f_j(x)=c j i^{N+1}e^{ijx_1}$,
której moduł jest równy $cj$. Stąd
$q_M(f_j)\to\infty$, więc $U$ nie jest ograniczony.
Dla przestrzeni Schwartza wybieramy analogiczną funkcję
$\psi\in C_c^\infty(\R^d;\R)$, równą $1$ na otoczeniu kuli.
Na jej zwartym nośniku waga $(1+|x|)^N$ jest ograniczona,
więc z oszacowania Leibniza otrzymujemy $Q_N(f_j)\le C_Nc$.
Na wnętrzu kuli mamy natomiast
$Q_{N+1}(f_j)\ge cj$, ponieważ waga jest nie mniejsza od $1$.
Dla dowolnego otoczenia zera zawierającego
$\{f:Q_N(f)<\varepsilon\}$ wybieramy zatem
$0<c<\varepsilon/C_N$, uzyskując nieograniczony podzbiór
tego otoczenia.
Przypadek rzeczywisty otrzymujemy, stosując
$\cos(jx_1)$ zamiast $e^{ijx_1}$ i supremum na małej kuli,
gdzie $\psi$ jest stała i niezerowa; przy dużym $j$
cosinus lub sinus odpowiedniej pochodnej przyjmuje tam
wartość bezwzględną $1$.
\end{proof}

\begin{uwaga}[Przygotowanie teorii dystrybucji]
\label{rem:distribution-next17}
Ciągłe funkcjonały na $C^\infty(\Omega)$ mają oszacowania
na pochodne na jednym zwartym zbiorze (Twierdzenie~\ref{thm:dual-cinf17}).
Dla liniowego funkcjonału $T:\mathcal S(\R^d;\K)\to\K$
analogiczny warunek ciągłości ma postać
$|T(f)|\le C Q_{N,m}(f)$ dla pewnych $C>0$, $N,m\in\N_0$
i wszystkich $f\in\mathcal S(\R^d;\K)$.
Istotnie, przy ciągłości pewne otoczenie
$\{f:Q_n(f)<\varepsilon\}$ jest odwzorowywane w kulę jednostkową
ciała skalarów. Skalowanie, jak w dowodzie
Twierdzenia~\ref{thm:dual-cinf17}, daje
$|T(f)|\le\varepsilon^{-1}Q_{n,n}(f)$.
Odwrotnie, takie oszacowanie zapewnia ciągłość w zerze,
a z liniowości również w każdym punkcie.
Przestrzenie ciągłych funkcjonałów liniowych na tych
przestrzeniach funkcji nie są tożsame.
W następnym rozdziale skonstruujemy $\mathcal D(\Omega)$
z topologią granicy induktywnej przestrzeni funkcji o nośniku
zawartym w ustalonym zbiorze zwartym. Dopiero wtedy zdefiniujemy
ogólną przestrzeń dystrybucji $\mathcal D'(\Omega)$.
\end{uwaga}

\setcounter{section}{1}
% Wersja źródłowa: audyt II, jednostka 18
\setcounter{section}{1}
\section{Funkcje testowe i ścisłe granice induktywne}
\label{sec:funkcje-testowe18}

\providecommand{\atlasref}[2]{#2}


\subsection{Funkcje o ustalonym zwartym nośniku}

\noindent\emph{Powiązania tematyczne w Bibliotece matematycznej:} \atlasref{PDE-07::def:pde7-test}{PDE-07: funkcje testowe}.


W całym rozdziale zakładamy, że $d\ge1$ jest liczbą całkowitą,
$\Omega\subseteq\R^d$ jest niepustym zbiorem otwartym,
a $\K\in\{\R,\C\}$. Oznaczenia wielowskaźników i seminorm
$p_{K,m}$ przyjmujemy jak w rozdziale~\ref{sec:gladkie-schwartz17}.
Dla $f\in C^\infty(\Omega;\K)$ definiujemy jego \emph{nośnik}
przez $\supp f=\overline{\{x\in\Omega:f(x)\ne0\}}^{\,\Omega}$,
tj. domknięcie względne w $\Omega$. Zapis $K\Subset\Omega$
oznacza, że $K$ jest zwartym podzbiorem $\Omega$.

\begin{definicja}[Przestrzeń $\mathcal D_K(\Omega)$]
\label{def:DK18}
Dla każdego zwartego $K\Subset\Omega$ kładziemy
\[
 \mathcal D_K(\Omega;\K)=\{\varphi\in C^\infty(\Omega;\K):
 \supp\varphi\subseteq K\}.
\]
Wyposażamy tę przestrzeń w topologię wyznaczoną przez seminormy
\[p_m(\varphi)=\max_{|\alpha|\le m}\sup_{x\in\Omega}
|\partial^\alpha\varphi(x)|,\qquad m\in\N_0.\]
Supremum jest skończone: na otwartym zbiorze $\Omega\setminus K$
funkcja i wszystkie jej pochodne znikają, a na $K$ są ciągłe
i ograniczone. Jeśli $K=\varnothing$, to $\mathcal D_K=\{0\}$.
Również gdy $\operatorname{int}K=\varnothing$, mamy
$\mathcal D_K=\{0\}$: dla niezerowej funkcji $\varphi$
zbiór $\{x\in\Omega:\varphi(x)\ne0\}$ jest niepusty,
otwarty i zawarty w $K$, co przeczy pustemu wnętrzu.
\end{definicja}

\begin{twierdzenie}[Struktura etapów]
\label{thm:DK-frechet18}
Dla każdego zwartego $K\Subset\Omega$ przestrzeń
$\mathcal D_K(\Omega;\K)$ jest przestrzenią Fr\'echeta.
Jeżeli $K\subseteq L\Subset\Omega$ są zwarte, to
$\mathcal D_K$ jest domkniętą podprzestrzenią $\mathcal D_L$,
a włożenie $\mathcal D_K\hookrightarrow\mathcal D_L$
jest izomorfizmem na swój obraz z topologią względną.
\end{twierdzenie}
\begin{proof}
Na mocy Twierdzenia~\ref{thm:frechet-cinf17}
$C^\infty(\Omega;\K)$ jest przestrzenią Frécheta.
Dla $x\in\Omega$ ewaluacja $\varphi\mapsto\varphi(x)$
jest ciągła, gdyż $|\varphi(x)|\le p_{\{x\},0}(\varphi)$.
Ponieważ $K$ jest domknięty w $\Omega$, mamy
\[
\mathcal D_K=\bigcap_{x\in\Omega\setminus K}
\ker(\varphi\mapsto\varphi(x)).
\]
Jest to przecięcie domkniętych podprzestrzeni liniowych,
a więc domknięta podprzestrzeń $C^\infty(\Omega;\K)$.
Dla każdego zwartego $L\Subset\Omega$ oraz $m\in\N_0$
na $\mathcal D_K$ mamy $p_{L,m}\le p_m$.
Z drugiej strony $p_m=p_{K,m}$ na $\mathcal D_K$,
gdyż wszystkie pochodne znikają poza $K$
(dla pustego $K$ obie seminormy są zerowe).
Obie rodziny wyznaczają zatem tę samą topologię.
Domknięta podprzestrzeń przestrzeni Fr\'echeta jest przestrzenią
Fr\'echeta. Ten sam argument, zastosowany do $K\subseteq L$,
daje domkniętość $\mathcal D_K$ w $\mathcal D_L$;
seminormy $p_m$ są na obu etapach określone tym samym wzorem.
\end{proof}

\begin{definicja}[Algebraiczna przestrzeń funkcji testowych]
\label{def:D-algebraic18}
Oznaczamy przez $\mathcal D(\Omega;\K)=C_c^\infty(\Omega;\K)$
przestrzeń wszystkich funkcji gładkich o zwartym nośniku.
Jeśli $(K_n)_{n\ge1}$ jest wyczerpaniem zwartymi zbiorami,
$K_n\subseteq\operatorname{int}K_{n+1}$ i każdy zwarty
$K\Subset\Omega$ jest zawarty w pewnym $K_n$, to
$\mathcal D(\Omega)=\bigcup_{n\ge1}\mathcal D_{K_n}(\Omega)$.
Takie wyczerpanie skonstruowano w Lemacie~\ref{lem:exhaustion17}.
Równość wynika stąd, że nośnik każdej funkcji testowej
jest zawarty w pewnym $K_n$, a każda funkcja z $\mathcal D_{K_n}$
ma zwarty nośnik.
\end{definicja}

\subsection{Topologia granicy induktywnej}

\noindent\emph{Powiązania tematyczne w Bibliotece matematycznej:} \atlasref{KAT-08::def:colimit}{KAT-08: kogranice i ich własność uniwersalna (właściwa kategoria ma tu znaczenie)}; \atlasref{GT-04::thm:topog-koncowa-uniwersalnosc}{GT-04: topologia końcowa (porównanie z granicą induktywną lokalnie wypukłą)}.


\begin{definicja}[Lokalnie wypukła granica induktywna]
\label{def:LF18}
Niech $(E_n,\tau_n)_{n\ge1}$ będzie ciągiem przestrzeni
lokalnie wypukłych nad tym samym ciałem $\K\in\{\R,\C\}$,
przy czym $E_n\subseteq E_{n+1}$ jako przestrzenie liniowe
i każde włożenie
$i_n:E_n\hookrightarrow E_{n+1}$ jest liniowe i ciągłe.
Na $E=\bigcup_nE_n$ rozważamy \emph{najsilniejszą lokalnie
wypukłą topologię}, dla której wszystkie odwzorowania
$E_n\hookrightarrow E$ są ciągłe. Oznaczamy ją
$\operatorname{ind}_n E_n$ i nazywamy topologią granicy induktywnej.
Na etapie ogólnej konstrukcji dopuszczamy topologię,
która nie musi spełniać warunku Hausdorffa.
Warunek ten zachodzi wtedy i tylko wtedy, gdy ciągłe
seminormy granicy oddzielają punkty. Jeśli $E_n$ są
przestrzeniami Frécheta i granica jest Hausdorffa,
nazywamy ją przestrzenią typu $(LF)$.
Granica jest \emph{ścisła}, jeśli
każde $E_n$ ma topologię względną z $E_{n+1}$ i jest w nim
podprzestrzenią domkniętą.
\end{definicja}

\begin{twierdzenie}[Konstrukcja przez seminormy i własność uniwersalna]
\label{thm:LF-universal18}
Niech $(E_n,\tau_n)_{n\ge1}$ będzie rosnącym ciągiem
przestrzeni lokalnie wypukłych nad tym samym ciałem $\K$,
z ciągłymi włożeniami liniowymi, a $E=\bigcup_{n\ge1}E_n$
niech będzie wyposażona w topologię granicy induktywnej
z Definicji~\ref{def:LF18}. Rozważmy rodzinę
$\mathcal P$ wszystkich seminorm $p:E\to[0,\infty)$ takich,
że $p|_{E_n}$ jest ciągła dla każdego $n$.
Topologia generowana przez $\mathcal P$ jest topologią
granicy induktywnej. Dla dowolnej przestrzeni lokalnie
wypukłej $F$ i odwzorowania liniowego $T:E\to F$ mamy
\[
 T\text{ jest ciągłe}\quad\Longleftrightarrow\quad
 T|_{E_n}:E_n\to F\text{ jest ciągłe dla każdego }n.
\]
\end{twierdzenie}
\begin{proof}
Oznaczmy przez $\tau_{\mathcal P}$ topologię wyznaczoną przez
$\mathcal P$. Przeciwobraz podstawowego otoczenia zera
$\bigcap_{j=1}^r\{x\in E:p_j(x)<\varepsilon_j\}$
przez włożenie $E_n\hookrightarrow E$ jest przecięciem
otoczeń zera określonych przez ciągłe seminormy $p_j|_{E_n}$.
Włożenie jest zatem ciągłe w topologii $\tau_{\mathcal P}$. Niech $\rho$ będzie
inną lokalnie wypukłą topologią na $E$, dla której włożenia
są ciągłe. Każda seminorma ciągła dla $\rho$ należy do
$\mathcal P$, gdyż jej ograniczenie do $E_n$ jest ciągłe.
Stąd $\rho$ jest nie silniejsza od naszej topologii.
Jeśli $T$ jest ciągłe, to każde $T|_{E_n}$ jest ciągłe
jako złożenie z ciągłym włożeniem.
Odwrotnie, niech wszystkie $T|_{E_n}$ będą ciągłe.
Dla każdej ciągłej seminormy $q$ na $F$ funkcja $q\circ T$
jest seminormą na $E$ dzięki liniowości $T$.
Jej ograniczenie do $E_n$ jest ciągłe, ponieważ jest równe
$q\circ(T|_{E_n})$. Zatem $q\circ T\in\mathcal P$,
a więc jest ciągła na $E$. Ponieważ ciągłe seminormy $q$
wyznaczają topologię $F$, ciągłość wszystkich $q\circ T$
daje ciągłość $T$.
\end{proof}

\begin{definicja}[Standardowa topologia $\mathcal D(\Omega)$]
\label{def:D-LF18}
Na $\mathcal D(\Omega;\K)$ przyjmujemy topologię
\[
 \mathcal D(\Omega;\K)
   =\operatorname{ind}_{n\to\infty}\mathcal D_{K_n}(\Omega;\K),
\]
gdzie $(K_n)$ jest wyczerpaniem z
Definicji~\ref{def:D-algebraic18}. Jest to \emph{standardowa
topologia funkcji testowych}. Nie należy jej utożsamiać
ani z topologią podprzestrzeni dziedziczoną po
$C^\infty(\Omega)$, ani automatycznie z końcową topologią
w kategorii wszystkich przestrzeni topologicznych.
\end{definicja}

\begin{stwierdzenie}[Niezależność od wyczerpania]
\label{prop:D-independence18}
Dla niepustego otwartego $\Omega\subseteq\R^d$ dwa rosnące
wyczerpania $\Omega$ zbiorami zwartymi, współkońcowe
wśród zwartych podzbiorów $\Omega$ prowadzą do tej samej
topologii na $\mathcal D(\Omega)$, nawet jeśli nie mają
identycznych wyrazów.
\end{stwierdzenie}
\begin{proof}
Niech $(K_n)$, $(L_j)$ będą takimi wyczerpaniami.
Seminorma $p$ na $\mathcal D(\Omega)$ jest ciągła na
każdym $\mathcal D_{K_n}$ wtedy i tylko wtedy, gdy jest
ciągła na każdym $\mathcal D_{L_j}$: dla ustalonego $K_n$
istnieje $L_j$ zawierający $K_n$, a ograniczenie seminormy
ciągłej na $\mathcal D_{L_j}$ do $\mathcal D_{K_n}$ jest
ciągłe. W drugą stronę argument jest symetryczny.
Korzystamy z Twierdzenia~\ref{thm:LF-universal18}.
\end{proof}

\begin{twierdzenie}[Topologia na każdym etapie i oddzielanie punktów]
\label{thm:D-stage18}
Dla dowolnego zwartego $K\Subset\Omega$ topologia względna
indukowana na $\mathcal D_K$ przez $\mathcal D(\Omega)$
jest topologią jego przestrzeni Fr\'echeta.
Ponadto $\mathcal D(\Omega)$ jest lokalnie wypukłą
przestrzenią liniowo-topologiczną Hausdorffa.
\end{twierdzenie}
\begin{proof}
Dla dowolnego zwartego $K\Subset\Omega$ wybieramy $n$
takie, że $K\subseteq K_n$. Włożenie $\mathcal D_K\to\mathcal D$
jest złożeniem dwóch ciągłych włożeń
$\mathcal D_K\to\mathcal D_{K_n}$ oraz
$\mathcal D_{K_n}\to\mathcal D$. Jego ciągłość oznacza, że
topologia względna nie jest silniejsza od topologii
$\mathcal D_K$. Dla każdego $m\in\N_0$ określmy na
całej $\mathcal D$ seminormę
$p_m(\varphi)=\max_{|\alpha|\le m}\sup_{x\in\Omega}
|\partial^\alpha\varphi(x)|$. Jej ograniczenie do
każdego $\mathcal D_{K_n}$ jest ciągłe, a zatem $p_m$
jest ciągła na $\mathcal D$ z
Twierdzenia~\ref{thm:LF-universal18}. Seminormy te
po ograniczeniu do $\mathcal D_K$ wyznaczają jego wyjściową
topologię Frécheta. Ponieważ są ciągłe w topologii względnej
indukowanej przez $\mathcal D$, topologia wyjściowa nie
jest od niej silniejsza. Obie topologie są więc równe.
Ponieważ $p_0(\varphi)>0$ dla $\varphi\ne0$, topologia
jest Hausdorffa; lokalna wypukłość wynika z konstrukcji.
Etapy $\mathcal D_{K_n}$ są przestrzeniami Frécheta
i są domkniętymi podprzestrzeniami kolejnych etapów
na mocy Twierdzenia~\ref{thm:DK-frechet18}.
Wobec tego $\mathcal D(\Omega;\K)$ jest ścisłą przestrzenią
typu $(LF)$ w sensie Definicji~\ref{def:LF18}.
\end{proof}

\begin{uwaga}[Przestrzeń dualna topologiczna]
\label{rem:D-dual18}
Własność uniwersalna jest szczególnie użyteczna dla
funkcjonałów: odwzorowanie liniowe $T:\mathcal D(\Omega)\to\K$
jest ciągłe wtedy i tylko wtedy, gdy każde ograniczenie
$T|_{\mathcal D_K}$ jest ciągłe. Ponieważ $\mathcal D_K$
ma topologię seminorm $p_m$, równoważnie: dla każdego
$K\Subset\Omega$ istnieją $C_K>0$ i $m_K\in\N_0$ takie, że
\[
 |T(\varphi)|\le C_K p_{m_K}(\varphi)
 \qquad(\varphi\in\mathcal D_K(\Omega)).
\]
Istotnie, jeśli $T|_{\mathcal D_K}$ jest ciągłe, to dla pewnych
$m_K$ i $\varepsilon>0$ warunek $p_{m_K}(\varphi)<\varepsilon$
pociąga $|T(\varphi)|<1$, gdyż seminormy $p_m$ są niemalejące.
Skalując przez $0<a<\varepsilon/p_{m_K}(\varphi)$
i przechodząc do kresu, dostajemy oszacowanie z
$C_K=\varepsilon^{-1}$; gdy $p_{m_K}(\varphi)=0$,
dowolnie duże $a>0$ dają $T(\varphi)=0$.
Odwrotnie, oszacowania zapewniają ciągłość każdego ograniczenia,
a własność uniwersalna daje ciągłość $T$.
Wielkości $C_K$ i $m_K$ mogą zależeć od $K$.
To będzie podstawowe kryterium w teorii dystrybucji.
\end{uwaga}


\begin{przyklad}[Dystrybucja o nieograniczonym rzędzie globalnym]
\label{ex:variable-order18}
Niech $\Omega=\R$. Dla $\varphi\in\mathcal D(\R)$ określamy
$T(\varphi)=\sum_{n=1}^\infty \varphi^{(n)}(n)$.
Suma jest skończona dla każdej $\varphi$, gdyż tylko skończenie
wiele liczb naturalnych może należeć do zwartego nośnika.
Na każdym $\mathcal D_K$ funkcjonał jest sumą skończonej
liczby ciągłych ewaluacji pochodnych, więc $T$ jest ciągły
na $\mathcal D(\R)$.

Nie istnieje jednak jeden rząd $m\in\N_0$ właściwy dla wszystkich
zbiorów zwartych, nawet z osobno dobranymi stałymi $C_K$.
Ustalmy $m$ i wybierzmy liczbę całkowitą $n>m$.
Niech $\chi\in C_c^\infty((-1/4,1/4))$ będzie równa $1$
w pobliżu zera, a $\psi(t)=t^n\chi(t)$. Dla $0<\varepsilon<1$
kładziemy $\varphi_\varepsilon(x)=\varepsilon^n
\psi((x-n)/\varepsilon)$. Wszystkie te funkcje mają
nośniki w $K=[n-1/4,n+1/4]$, który nie zawiera żadnej
innej liczby całkowitej niż $n$. Dla $k\in\N_0$ reguła
łańcuchowa daje
$\varphi_\varepsilon^{(k)}(x)=
\varepsilon^{n-k}\psi^{(k)}((x-n)/\varepsilon)$.
Zatem jedyny potencjalnie niezerowy składnik sumy to
$T(\varphi_\varepsilon)=\varphi_\varepsilon^{(n)}(n)
=\psi^{(n)}(0)=n!$, gdyż $\psi(t)=t^n$ w otoczeniu zera.
Ponieważ $0<\varepsilon<1$ i $n>m$,
$p_m(\varphi_\varepsilon)\le
\varepsilon^{n-m}\max_{0\le k\le m}\|\psi^{(k)}\|_\infty\to0$.
Oszacowanie $|T(\varphi_\varepsilon)|\le C_Kp_m(\varphi_\varepsilon)$
dla tego ustalonego $K$ prowadziłoby do $n!\le0$.
To wyklucza oszacowanie z jednym globalnie ustalonym $m$.
\end{przyklad}

\subsection{Zbiory ograniczone, zbieżność ciągów i zupełność}

\begin{twierdzenie}[Dokładne kryterium ograniczoności]
\label{thm:D-bounded18}
Niech $B\subseteq\mathcal D(\Omega;\K)$. Następujące
warunki są równoważne:
\begin{enumerate}
\item $B$ jest ograniczony w topologii $\mathcal D(\Omega)$;
\item istnieje zwarty $K\Subset\Omega$, który zawiera
nośniki wszystkich $\varphi\in B$, a ponadto dla każdego
$m\in\N_0$ zachodzi $\sup_{\varphi\in B}p_m(\varphi)<\infty$.
\end{enumerate}
\end{twierdzenie}
\begin{proof}
Załóżmy warunek 2. Z kryterium ograniczoności w przestrzeni
lokalnie wypukłej, której topologię generują seminormy $p_m$,
wynika, że $B$ jest ograniczony w $\mathcal D_K$.
Jego obraz przez ciągłe liniowe włożenie do $\mathcal D$
jest więc ograniczony.

Na odwrót, niech $B$ będzie ograniczony w $\mathcal D$.
Ponieważ każda seminorma $p_m$ z dowodu
Twierdzenia~\ref{thm:D-stage18} jest ciągła, mamy
$\sup_{\varphi\in B}p_m(\varphi)<\infty$.
Pozostaje wykazać wspólność nośnika. Załóżmy przeciwnie,
że $B$ nie zawiera się w żadnym $\mathcal D_{K_n}$.
Dla każdego $n$ wybierzmy $\varphi_n\in B\setminus\mathcal D_{K_n}$.
Ponieważ $K_n$ jest domknięty, istnieje
$x_n\in\Omega\setminus K_n$ z $\varphi_n(x_n)\ne0$
(w przeciwnym razie także nośnik byłby zawarty w $K_n$).
Z Lematu~\ref{lem:cutoffs17}, zastosowanego do $\{x_n\}$
w otwartym zbiorze $\Omega\setminus K_n$, wybieramy
$\theta_n\in C_c^\infty(\Omega\setminus K_n;\R)$,
$0\le\theta_n\le1$, $\theta_n(x_n)=1$.
Przedłużamy je zerem do $\Omega$; przedłużenia są gładkie
dzięki zwartości nośników w $\Omega\setminus K_n$.
Rodzina nośników jest lokalnie skończona: każdy punkt
$\Omega$ należy do wnętrza pewnego $K_M$, a dla $n\ge M$
nośnik $\theta_n$ nie przecina $K_M$. Ustalmy $a_n>n/|\varphi_n(x_n)|$ i określmy
$w(x)=1+\sum_{n\ge1}a_n\theta_n(x)$. Suma jest lokalnie
skończona, zatem $w$ jest dodatnią funkcją gładką.
Seminorma $p_w(\varphi)=\sup_{x\in\Omega}w(x)|\varphi(x)|$
jest skończona dla każdej funkcji testowej. Dla niepustego zwartego
$K$ mamy $p_w|_{\mathcal D_K}\le(\sup_Kw)p_0$,
a na $\mathcal D_\varnothing=\{0\}$ seminorma jest zerowa, więc
z Twierdzenia~\ref{thm:LF-universal18} wynika jej ciągłość
na $\mathcal D$. Jednak $p_w(\varphi_n)>n$,
co przeczy ograniczoności $B$. Wobec tego istnieje
$K_n$ zawierający wszystkie nośniki, a oszacowania
$p_m$ kończą dowód.
\end{proof}

\begin{twierdzenie}[Zbieżność ciągów funkcji testowych]
\label{thm:D-sequences18}
Dla ciągu $(\varphi_j)\subseteq\mathcal D(\Omega;\K)$
i $\varphi\in\mathcal D(\Omega;\K)$ następujące warunki
są równoważne:
\begin{enumerate}
\item $\varphi_j\to\varphi$ w topologii $\mathcal D(\Omega)$;
\item istnieje zwarty $K\Subset\Omega$ zawierający
$\supp\varphi$ i wszystkie $\supp\varphi_j$ oraz
$p_m(\varphi_j-\varphi)\to0$ dla każdego $m\in\N_0$.
\end{enumerate}
\end{twierdzenie}
\begin{proof}
Dla każdej ciągłej seminormy $q$ na $\mathcal D$
z założenia zbieżności mamy $q(\varphi_j-\varphi)\to0$.
Nierówność $q(\varphi_j)\le q(\varphi)+q(\varphi_j-\varphi)$
daje ograniczoność ciągu wartości $q(\varphi_j)$.
Zatem $\{\varphi\}\cup\{\varphi_j:j\ge1\}$ jest ograniczony
w $\mathcal D$ na mocy kryterium seminormowego.
Twierdzenie~\ref{thm:D-bounded18} daje wspólny zwarty $K$.
Ponieważ topologia na $\mathcal D_K$ jest względna,
zbieżność w $\mathcal D$ oznacza zbieżność w $\mathcal D_K$,
a tę charakteryzują seminormy $p_m$.
Implikacja odwrotna wynika z ciągłości
$\mathcal D_K\hookrightarrow\mathcal D$.
\end{proof}

\begin{wniosek}[Zupełność ciągowa]
\label{cor:D-seqcomplete18}
Każdy ciąg Cauchy'ego w $\mathcal D(\Omega;\K)$ jest
zbieżny w tej przestrzeni.
\end{wniosek}
\begin{proof}
Każdy ciąg Cauchy'ego jest ograniczony w przestrzeni lokalnie
wypukłej: dla każdej ciągłej seminormy $q$ nierówność
$|q(\varphi_j)-q(\varphi_k)|\le q(\varphi_j-\varphi_k)$
pokazuje, że $(q(\varphi_j))$ jest rzeczywistym ciągiem
Cauchy'ego, a więc jest ograniczony.
Wybieramy wspólny $K$ z
Twierdzenia~\ref{thm:D-bounded18}. Ograniczenia topologii
na $\mathcal D_K$ i $\mathcal D$ są jednakowe, więc ciąg
jest Cauchy'ego w przestrzeni Fr\'echeta $\mathcal D_K$;
tam ma granicę, do której zbiega również w $\mathcal D$.
\end{proof}

\begin{twierdzenie}[Zupełność $\mathcal D(\Omega)$]
\label{thm:D-complete18}
Przestrzeń $\mathcal D(\Omega;\K)$ jest zupełna:
każdy filtr Cauchy'ego względem jej kanonicznej struktury
jednostajnej jest zbieżny. W szczególności jest zupełną
ścisłą przestrzenią typu $(LF)$.
\end{twierdzenie}
\begin{proof}
Przeprowadzimy dowód dla ogólnej rosnącej sumy
$E=\bigcup_nE_n$ przestrzeni zupełnych lokalnie wypukłych,
która ma topologię granicy induktywnej i indukuje na
każdym $E_n$ topologię wyjściową. W naszym przypadku
$E_n=\mathcal D_{K_n}$, a wszystkie te własności już
udowodniono.

Niech $\mathcal F$ będzie filtrem Cauchy'ego w $E$.
Przez $\mathcal N$ oznaczmy filtr otoczeń zera, a przez
$\mathcal G$ filtr generowany przez zbiory $F+V$,
$F\in\mathcal F$, $V\in\mathcal N$. Zbiory te są niepuste
i dla dwóch takich zbiorów ich przecięcie zawiera
$(F_1\cap F_2)+(V_1\cap V_2)$, więc tworzą bazę filtru.
Ponieważ $F\subseteq F+V$, mamy $\mathcal G\subseteq\mathcal F$.
Dla dowolnego otoczenia $U$ zera wybierzmy symetryczne
otoczenie $V$ takie, że $V+V+V\subseteq U$.
Istnieje $F\in\mathcal F$ z $F-F\subseteq V$,
a wtedy $(F+V)-(F+V)\subseteq V+V+V\subseteq U$.
To dowodzi, że $\mathcal G$ jest filtrem Cauchy'ego.
Pokażemy, że dla pewnego $n$ każdy element $\mathcal G$
przecina $E_n$.

Załóżmy przeciwnie. Dla każdego $n$ istnieją
$F_n\in\mathcal F$ i absolutnie wypukłe otoczenie
$V_n$ zera w $E$ takie, że
$(F_n+V_n)\cap E_n=\varnothing$. Zastępując $V_n$
odpowiednimi skończonymi przecięciami, możemy zakładać
$V_{n+1}\subseteq V_n$. Definiujemy
$W=\operatorname{aco}(\bigcup_n(V_n\cap E_n))$,
gdzie $\operatorname{aco}$ oznacza absolutnie wypukłą
otoczkę zbioru. Zbiór $W$ jest absolutnie wypukły,
pochłaniający w $E$, a dla każdego $n$ zbiór
$W\cap E_n$ zawiera otoczenie $V_n\cap E_n$ zera w $E_n$.
Z konstrukcji topologii granicy induktywnej wynika, że
$W$ jest otoczeniem zera: jego funkcjonał Minkowskiego
$p_W(z)=\inf\{t>0:z\in tW\}$ jest seminormą,
ponieważ $W$ jest absolutnie wypukły i pochłaniający.
Na każdym $E_n$ jest ciągły: dla dowolnego $\eta>0$
otoczenie $(\eta/2)(V_n\cap E_n)$ jest zawarte
w $\{z\in E_n:p_W(z)<\eta\}$, dzięki jednorodności
i nierówności $p_W(v)\le1$ dla $v\in V_n\cap E_n$.
Z Twierdzenia~\ref{thm:LF-universal18} (zastosowanego
do rodziny seminorm) wynika ciągłość $p_W$ na $E$.
Ponieważ $\{z\in E:p_W(z)<1\}\subseteq W$,
zbiór $W$ jest otoczeniem zera.

Każdy $w\in W$ jest skończoną sumą $\sum_j c_jv_j$,
gdzie $v_j\in V_{k_j}\cap E_{k_j}$ i $\sum_j|c_j|\le1$.
Dla ustalonego $n$ suma składników z $k_j\le n$
należy do $E_n$. Pozostałe $v_j$ należą do $V_n$,
więc ich ważona suma należy do $V_n$ z absolutnej wypukłości.
Dlatego $W\subseteq E_n+V_n$.
Gdyby $f+w\in E_n$ dla $f\in F_n$ i $w=u+v$,
$u\in E_n$, $v\in V_n$, to $f+v=(f+w)-u\in E_n$,
wbrew $(F_n+V_n)\cap E_n=\varnothing$.
Stąd $(F_n+W)\cap E_n=\varnothing$.
Ponieważ $\mathcal F$ jest filtrem Cauchy'ego, istnieje
$F_0\in\mathcal F$ takie, że $F_0-F_0\subseteq W$.
Wybierzmy $x\in F_0$; należy on do pewnego $E_n$.
Istnieje $y\in F_0\cap F_n$, bo oba zbiory należą do
filtru $\mathcal F$. Ale
$x=y+(x-y)\in(F_n+W)\cap E_n$, sprzeczność.

Zatem istnieje $n$ takie, że każdy element $\mathcal G$
przecina $E_n$. Przecięcia tych elementów z $E_n$
tworzą bazę właściwego filtru na $E_n$,
bo każde skończone przecięcie również jest niepuste.
Dla dowolnego otoczenia zera w $E_n$ wybieramy otoczenie
$U$ zera w $E$, którego przecięcie z $E_n$ jest w nim
zawarte. Istnieje $G\in\mathcal G$ z $G-G\subseteq U$,
zatem $(G\cap E_n)-(G\cap E_n)\subseteq U\cap E_n$.
Filtr na $E_n$ jest więc Cauchy'ego. Zupełność $E_n$
daje jego granicę $x\in E_n$. Każde otoczenie $x$ w $E$
przecina każdy $G\in\mathcal G$, ponieważ przecięcie
tego otoczenia z $E_n$ należy do filtru na $E_n$,
podobnie jak $G\cap E_n$. Punkt $x$ jest zatem punktem
skupienia filtru $\mathcal G$ w $E$. Filtr Cauchy'ego mający punkt skupienia
jest zbieżny do tego punktu: dla dowolnego otoczenia
$U$ zera bierzemy $V$ z $V+V\subseteq U$ oraz
$G\in\mathcal G$ z $G-G\subseteq V$; skoro
$(x+V)\cap G\ne\varnothing$, to $G\subseteq x+U$.
Stąd $\mathcal G\to x$, a ponieważ $\mathcal F$ jest
silniejszy, również $\mathcal F\to x$.
\end{proof}

\subsection{Dlaczego topologia \texorpdfstring{$\mathcal D$}{D} jest istotna?}

\begin{przyklad}[Ucieczka nośników]
\label{ex:escaping18}
Wybierzmy ciąg niezerowych funkcji $\psi_n\in\mathcal D(\Omega)$
mających nośniki poza $K_n$. Każda taka funkcja istnieje,
ponieważ $\Omega\setminus K_n$ jest niepustym zbiorem otwartym:
gdyby $\Omega=K_n$, to $\Omega$ byłby jednocześnie zwarty
i otwarty w $\R^d$, co jest niemożliwe dla niepustego
zbioru przy $d\ge1$. Wystarczy zastosować
Lemat~\ref{lem:cutoffs17} do jednego punktu tego dopełnienia
i przedłużyć otrzymaną funkcję zerem do $\Omega$.
Dla dowolnego zwartego $K\Subset\Omega$ i ustalonego $m$
mamy $p_{K,m}(\psi_n)=0$ dla dostatecznie dużych $n$.
Zatem $\psi_n\to0$ w $C^\infty(\Omega)$.
Nie zachodzi jednak $\psi_n\to0$ w $\mathcal D(\Omega)$,
ponieważ nośniki nie mieszczą się w żadnym wspólnym
zbiorze zwartym. Ponadto włożenie
$\mathcal D(\Omega)\to C^\infty(\Omega)$ jest ciągłe,
ponieważ $p_{K,m}\le p_m$ i $p_m$ jest ciągła na $\mathcal D$.
Topologia funkcji testowych jest zatem ściśle silniejsza
od topologii względnej z $C^\infty$.
\end{przyklad}

\begin{twierdzenie}[Niemetryzowalność]
\label{thm:D-nonmetric18}
Dla każdego niepustego otwartego $\Omega\subseteq\R^d$,
$d\ge1$, przestrzeń $\mathcal D(\Omega;\K)$ z jej
standardową topologią nie jest metryzowalna.
\end{twierdzenie}
\begin{proof}
Każda podprzestrzeń $E_n=\mathcal D_{K_n}$ jest domknięta
w $\mathcal D$: istotnie,
$E_n=\bigcap_{x\in\Omega\setminus K_n}\{\varphi:\varphi(x)=0\}$,
a ewaluacje $\varphi\mapsto\varphi(x)$ są ciągłe
na każdym etapie, więc i na $\mathcal D$.
Każde $E_n$ jest właściwe na mocy konstrukcji funkcji
$\psi_n$ z Przykładu~\ref{ex:escaping18}.
Właściwa podprzestrzeń liniowa ma puste wnętrze:
gdyby zawierała niepusty zbiór otwarty, po przesunięciu
zawierałaby otoczenie zera, które jest pochłaniające,
a przez liniowość zawierałaby całą przestrzeń.
Gdyby $\mathcal D$ miała przeliczalną malejącą bazę
otoczeń zera $(U_n)$, moglibyśmy wybrać
$\varphi_n\in U_n\setminus E_n$. Wtedy
$\varphi_n\to0$, lecz nośniki $\varphi_n$ nie mieściłyby
się w żadnym wspólnym zbiorze zwartym, wbrew
Twierdzeniu~\ref{thm:D-sequences18}.
\end{proof}

\begin{twierdzenie}[Względna zwartość zbiorów ograniczonych]
\label{thm:D-Montel18}
Każdy zbiór ograniczony $B\subseteq\mathcal D(\Omega;\K)$
jest względnie zwarty w standardowej topologii.
\end{twierdzenie}
\begin{proof}
Z Twierdzenia~\ref{thm:D-bounded18} istnieje zwarty $K$
taki, że $B\subseteq\mathcal D_K$ i wszystkie seminormy
$p_m$ są jednostajnie ograniczone na $B$.
Jeśli $K=\varnothing$, teza jest oczywista.
Przedłużmy funkcje z $\mathcal D_K$ zerem poza $\Omega$
do funkcji gładkich na $\R^d$: każdy punkt spoza $\Omega$
ma otoczenie rozłączne z domkniętym zbiorem $K$,
na którym przedłużenie jest identycznie zerowe. Dla każdej pochodnej
$\partial^\alpha$ pochodne rzędu $|\alpha|+1$ mają
jednostajną normę supremum, toteż rodzina
$\{\partial^\alpha\varphi:\varphi\in B\}$ jest
jednostajnie ograniczona i równociągła na $\R^d$
(z twierdzenia o wartości średniej). Na zwartym $K$
możemy zastosować twierdzenie Arzel\`a--Ascoliego.
Wybierając kolejno podciągi dla przeliczalnej rodziny
wielowskaźników i stosując argument diagonalny,
uzyskujemy podciąg dowolnego ciągu elementów $B$,
którego wszystkie pochodne są jednostajnie zbieżne na $K$.
Poza $K$ wszystkie pochodne są zerowe. Dla każdego $m$
istnieje tylko skończenie wiele wielowskaźników $|\alpha|\le m$,
więc otrzymany podciąg jest Cauchy'ego dla $p_m$.
Z zupełności $\mathcal D_K$ wynika istnienie granicy w $\mathcal D_K$, a więc także
w $\mathcal D$. Skoro $\mathcal D_K$ jest metryzowalna,
każdy ciąg z domknięcia $B$ w $\mathcal D_K$ można
przybliżyć wyraz po wyrazie ciągiem z $B$, z błędami
metrycznymi mniejszymi niż $1/j$. Z tego ostatniego
wybieramy zbieżny podciąg, więc odpowiedni podciąg
pierwszego ciągu zbiega do tej samej granicy.
Domknięcie $B$ w $\mathcal D_K$ jest zatem zwarte ciągowo,
a przez metryzowalność jest zwarte. Jego obraz w $\mathcal D$
jest zwarty i domknięty, ponieważ $\mathcal D$ jest Hausdorffa.
Zawiera $B$, zatem domknięcie $B$ w $\mathcal D$ jest
domkniętym podzbiorem zbioru zwartego, czyli jest zwarte.
\end{proof}

\begin{stwierdzenie}[Naturalne operatory na funkcjach testowych]
\label{prop:D-operators18}
Niech $a\in C^\infty(\Omega;\K)$ i
$\alpha\in\N_0^d$. Odwzorowania
$\varphi\mapsto\partial^\alpha\varphi$ oraz
$\varphi\mapsto a\varphi$ są ciągłymi operatorami
liniowymi $\mathcal D(\Omega)\to\mathcal D(\Omega)$.
Dla $f\in L^1_{\mathrm{loc}}(\Omega;\K)$ odwzorowanie
$T_f(\varphi)=\int_\Omega f(x)\varphi(x)\,\mathrm d\lambda_d(x)$
jest ciągłym funkcjonałem liniowym na $\mathcal D(\Omega)$.
\end{stwierdzenie}
\begin{proof}
Każdy z dwóch pierwszych operatorów zachowuje
$\mathcal D_K$, ponieważ funkcja i jej pochodne znikają
poza $K$. Dla $m\in\N_0$ mamy
$p_m(\partial^\alpha\varphi)\le p_{m+|\alpha|}(\varphi)$,
a ze wzoru Leibniza
$p_m(a\varphi)\le 2^m p_{K,m}(a)p_m(\varphi)$
(dla pustego $K$ obie strony są zerowe).
Każdy operator jest zatem ciągły jako odwzorowanie
$\mathcal D_K\to\mathcal D_K$. Złożenie z ciągłym
włożeniem $\mathcal D_K\to\mathcal D$ jest ciągłe.
Stosując Twierdzenie~\ref{thm:LF-universal18} do tych
ograniczeń, otrzymujemy ciągłość na $\mathcal D$.
Całka definiująca $T_f$ jest bezwzględnie zbieżna
i niezależna od reprezentanta klasy $f$, ponieważ
$\varphi$ jest ograniczona i ma zwarty nośnik.
Dla $\varphi\in\mathcal D_K$ mamy
$|T_f(\varphi)|\le
\bigl(\int_K|f(x)|\,\mathrm d\lambda_d(x)\bigr)p_0(\varphi)$,
a więc każde ograniczenie $T_f|_{\mathcal D_K}$
jest ciągłe. Ponownie stosujemy własność uniwersalną.
\end{proof}

\begin{center}
\begin{tikzpicture}[scale=0.82]
 \draw[thick,rounded corners=14pt] (-4,-1.7) rectangle (4.6,1.8);
 \draw[thick,rounded corners=12pt,blue!55!black] (-3.0,-1.25) rectangle (2.6,1.3);
 \draw[thick,rounded corners=10pt,blue!55!black] (-2.1,-0.85) rectangle (1.15,0.85);
 \draw[thick,rounded corners=7pt,blue!55!black] (-1.2,-0.43) rectangle (0.25,0.43);
 \node at (-0.48,0) {$K_1$};
 \node at (0.75,0.64) {$K_2$};
 \node at (-2.55,0.94) {$K_3$};
 \node at (-3.55,1.49) {$\Omega$};
 \draw[orange!75!black,thick,fill=orange!15] (3.60,-0.72) ellipse (0.31 and 0.25);
 \node[orange!65!black,font=\footnotesize] at (3.60,0.05) {$\supp\psi_3$};
 \draw[orange!75!black,->] (3.60,-0.18) -- (3.60,-0.44);
\end{tikzpicture}

\small\emph{Schemat:} funkcje testowe mogą mieć nośniki poza
kolejnymi zbiorami zwartymi, a wtedy nie tworzą ciągu zbieżnego
w topologii $\mathcal D(\Omega)$, nawet jeśli zbiegają do zera
w $C^\infty(\Omega)$.
\end{center}

\begin{uwaga}[Przejście do dystrybucji]
\label{rem:D-distributions-next18}
W kolejnym rozdziale zdefiniujemy przestrzeń dystrybucji
$\mathcal D'(\Omega)$ jako przestrzeń dualną topologiczną
$\mathcal D(\Omega)$. Podstawowe oszacowania funkcjonałów
uzyskaliśmy już w Uwadze~\ref{rem:D-dual18}; będziemy też
odróżniać topologię słabą-* od topologii zbieżności
jednostajnej na zbiorach ograniczonych funkcji testowych.
\end{uwaga}

% Materiał z AF 25; fragment do rozdziału o funkcjach testowych.
\subsection{Porównanie topologii przestrzeni funkcji testowych}
\label{sec:synteza25}

\begin{definicja}[Notacja porównawcza przestrzeni funkcji]
\label{def:trzy-przestrzenie25}
Zachowujemy przestrzenie $\mathcal E(\Omega)=C^\infty(\Omega)$,
$\mathcal D_K(\Omega)$, $\mathcal D(\Omega)$ i $\mathcal S(\R^d)$
z definicji \ref{def:DK18}, \ref{def:D-LF18}
oraz \ref{def:schwartz17}. Na $\mathcal E$ i $\mathcal D_K$
używamy seminorm
$p_{K,m}(f)=\max_{|\alpha|\le m}\sup_{x\in K}|\partial^\alpha f(x)|$.
Na $\mathcal S$ zapis $p_{N,m}(f)$ oznacza normę
$Q_{N,m}(f)$ z definicji \ref{def:schwartz17}.
Przestrzenie sprzężone topologicznie oznaczamy przez
$\mathcal E'$, $\mathcal D'$ i $\mathcal S'$; indeksy
$\sigma(X',X)$ i $\beta(X',X)$ oznaczają odpowiednio
topologię słabą-* (zbieżność punktowa na $X$)
i silną (zbieżność jednostajna na zbiorach ograniczonych).
Wszystkie przestrzenie rozpatrujemy nad wspólnym ciałem $\K$.
\end{definicja}

\begin{twierdzenie}[Ciągłe i gęste zanurzenia]
\label{thm:zanurzenia25}
Niech $d\ge1$. Odwzorowania identycznościowe
\[
 \mathcal D(\R^d)\longrightarrow\mathcal S(\R^d)
 \longrightarrow\mathcal E(\R^d)
\]
są ciągłymi, różnowartościowymi odwzorowaniami liniowymi o gęstych
obrazach. Dla dowolnego niepustego zbioru otwartego $\Omega\subseteq\R^d$
takie samo stwierdzenie zachodzi dla zanurzenia
$\mathcal D(\Omega)\to\mathcal E(\Omega)$.
Żadne z trzech wskazanych odwzorowań nie jest
izomorfizmem liniowo-topologicznym na swój obraz z topologią
podprzestrzeni; ostatnie stwierdzenie zachodzi dla każdego
niepustego zbioru otwartego $\Omega$.
\end{twierdzenie}
\begin{proof}
Dla ustalonego zwartego $K\Subset\R^d$ oraz $N,m\in\N_0$
mamy $p_{N,m}(f)\le (1+\sup_{x\in K}|x|)^N p_{K,m}(f)$
dla $f\in\mathcal D_K(\R^d)$. Stąd ciągłość
$\mathcal D_K\to\mathcal S$; własność uniwersalna granicy
induktywnej daje ciągłość $\mathcal D\to\mathcal S$.
Ciągłość $\mathcal S\to\mathcal E$ wynika z oszacowania
$p_{K,m}(f)\le p_{0,m}(f)$. Analogicznie sprawdzamy
ciągłość $\mathcal D(\Omega)\to\mathcal E(\Omega)$.

Aby wykazać gęstość, wybierzmy $\chi\in\mathcal D(\R^d)$
równe $1$ na kuli jednostkowej i zdefiniujmy
$\chi_R(x)=\chi(x/R)$ dla $R\ge1$. Dla $f\in\mathcal S(\R^d)$ i $|\alpha|\le m$
wzór Leibniza daje
\[
 \partial^\alpha((1-\chi_R)f)
 =\sum_{\gamma\le\alpha}\binom{\alpha}{\gamma}
 \partial^\gamma(1-\chi_R)\partial^{\alpha-\gamma}f.
\]
Wszystkie czynniki $\partial^\gamma(1-\chi_R)$
znikają dla $|x|<R$. Dla $\gamma\ne0$ są ograniczone
przez $R^{-|\gamma|}\sup|\partial^\gamma\chi|$,
a dla $\gamma=0$ przez $1+\sup|\chi|$.
Ponieważ $R\ge1$ i $(1+|x|)^{-1}\le(1+R)^{-1}$
na pozostałej części przestrzeni, otrzymujemy
\[
 p_{N,m}((1-\chi_R)f)
 \le \frac{C_{\chi,m}}{1+R}p_{N+1,m}(f).
\]
Zatem $\chi_Rf\to f$ w $\mathcal S$ i $\chi_Rf\in\mathcal D$.
Dla $f\in\mathcal E(\R^d)$ mamy $\chi_Rf=f$ w otoczeniu
każdego ustalonego zbioru zwartego przy dostatecznie dużym $R$,
a więc $\chi_Rf\to f$ w $\mathcal E$. W ogólnej
$\Omega$ używamy funkcji odcięcia $\chi_j\in\mathcal D(\Omega)$,
równej $1$ w otoczeniu $K_j$, gdzie $(K_j)$ jest
wyczerpaniem $\Omega$ zbiorami zwartymi. Każdy zwarty $L\Subset\Omega$ mieści się w pewnym $K_J$,
więc dla $j\ge J$ wszystkie pochodne różnicy znikają na $L$.
Stąd $\chi_j f\to f$ w $\mathcal E(\Omega)$.

Na $\R^d$ niech $\psi\in\mathcal D(\R^d)$ będzie niezerowa
i $f_n(x)=\psi(x-ne_1)$. Wówczas $f_n\to0$ w
$\mathcal E(\R^d)$, bo nośniki ostatecznie omijają każdy
ustalony zbiór zwarty, ale nie w $\mathcal S(\R^d)$,
gdyż $p_{0,0}(f_n)=p_{0,0}(\psi)>0$.
Z kolei $g_n(x)=e^{-n^2}\psi(x-ne_1)$ zbiega do zera w
$\mathcal S(\R^d)$: dla $\supp\psi\subseteq\overline B(0,A)$
mamy
\[
 p_{N,m}(g_n)\le e^{-n^2}(1+n+A)^N p_{0,m}(\psi)
 \longrightarrow0.
\]
Nie zbiega do zera w $\mathcal D(\R^d)$, gdyż nośniki
nie mieszczą się w jednym zbiorze zwartym
(Twierdzenie~\ref{thm:D-sequences18}).

Dla dowolnego niepustego otwartego $\Omega\subseteq\R^d$
wybieramy wyczerpanie $(K_j)$ jak w definicji.
Zbiór $\Omega\setminus K_j$ jest niepusty i otwarty,
bo żaden niepusty otwarty podzbiór $\R^d$ nie jest zwarty.
Wybieramy niezerowe $h_j\in\mathcal D(\Omega)$ takie,
że $\supp h_j\subseteq\Omega\setminus K_j$.
Wówczas $h_j\to0$ w $\mathcal E(\Omega)$, gdyż dla
każdego zwartego $L\Subset\Omega$ nośniki ostatecznie
omijają jego otoczenie (wybieramy $L\subseteq\operatorname{int}K_J$).
Gdyby $h_j\to0$ w $\mathcal D(\Omega)$, wszystkie nośniki
byłyby zawarte w pewnym zwartym $L\Subset\Omega$,
a więc w pewnym $K_J$. Dla $j\ge J$ oznaczałoby to
$h_j=0$, sprzecznie z wyborem. Topologia $\mathcal D(\Omega)$
jest zatem ściśle silniejsza od odziedziczonej z $\mathcal E(\Omega)$.
\end{proof}

\setcounter{section}{2}
% Wersja źródłowa: audyt II, jednostka 19

\section{Przestrzeń dystrybucji i jej topologie}
\label{sec:dystrybucje19}

\providecommand{\atlasref}[2]{#2}


\subsection{Definicja i pierwsze przykłady}

\noindent\emph{Powiązanie tematyczne w Bibliotece matematycznej:} \atlasref{PDE-07::def:pde7-distribution}{PDE-07: dystrybucje w teorii równań różniczkowych}.


W całym rozdziale zakładamy, że $d\ge1$ jest liczbą całkowitą,
$\Omega\subseteq\R^d$ jest niepustym zbiorem otwartym,
a $\K\in\{\R,\C\}$.
Przestrzeń $\mathcal D(\Omega;\K)$ ma topologię skonstruowaną
w rozdziale~\ref{sec:funkcje-testowe18}. W szczególności
jest lokalnie wypukła, Hausdorffa i zupełna, ale nie jest
metryzowalna. Dla $K\Subset\Omega$ zapisujemy
$\mathcal D_K=\{\varphi\in\mathcal D(\Omega):\supp\varphi\subseteq K\}$
oraz $p_m(\varphi)=\max_{|\alpha|\le m}\sup_{x\in\Omega}
|\partial^\alpha\varphi(x)|$. Symbolem $\mathcal D'(\Omega;\K)$
oznaczamy zbiór \emph{ciągłych $\K$-liniowych} funkcjonałów
na $\mathcal D(\Omega;\K)$.

\begin{definicja}[Dystrybucja]
\label{def:distribution19}
\emph{Dystrybucją} na $\Omega$ nazywamy element
$T\in\mathcal D'(\Omega;\K)$; wartość $T$ na funkcji
testowej $\varphi$ zapisujemy $\langle T,\varphi\rangle$.
Przyjmujemy konwencję $\K$-liniowości względem $\varphi$:
$\langle T,a\varphi+b\psi\rangle=a\langle T,\varphi\rangle+
b\langle T,\psi\rangle$ dla dowolnych $a,b\in\K$
i $\varphi,\psi\in\mathcal D(\Omega;\K)$.
Nie wprowadzamy tutaj sprzężenia
zespolonego w drugim argumencie.
\end{definicja}

\begin{twierdzenie}[Lokalne oszacowanie charakteryzujące dystrybucje]
\label{thm:distribution-estimate19}
Niech $T:\mathcal D(\Omega;\K)\to\K$ będzie odwzorowaniem
liniowym. Następujące warunki są równoważne:
\begin{enumerate}
\item $T$ jest ciągły w standardowej topologii $\mathcal D(\Omega)$;
\item dla każdego $K\Subset\Omega$ istnieją $C_K\ge0$
i $m_K\in\N_0$ takie, że
\[|\langle T,\varphi\rangle|\le C_K p_{m_K}(\varphi)
\qquad(\varphi\in\mathcal D_K(\Omega)).\]
\end{enumerate}
\end{twierdzenie}
\begin{proof}
Z własności uniwersalnej granicy induktywnej
(Twierdzenie~\ref{thm:LF-universal18}) ciągłość $T$
jest równoważna ciągłości każdego
$T|_{\mathcal D_K}$. Topologia $\mathcal D_K$ jest wyznaczona
przez rosnącą rodzinę seminorm $(p_m)_{m\in\N_0}$.
Jeśli $T|_{\mathcal D_K}$ jest ciągły, istnieje $m$ oraz
$\varepsilon>0$ takie, że $p_m(\varphi)<\varepsilon$
implikuje $|T(\varphi)|<1$. Dla $p_m(\varphi)>0$ stosujemy
tę implikację do $a\varphi$ dla
$a=\varepsilon/(2p_m(\varphi))>0$.
Wówczas $p_m(a\varphi)=\varepsilon/2$ oraz
$a|T(\varphi)|<1$, czyli
$|T(\varphi)|\le(2/\varepsilon)p_m(\varphi)$.
Jeśli $p_m(\varphi)=0$, to dla każdego $a>0$ mamy
$a|T(\varphi)|<1$, skąd $T(\varphi)=0$. Odwrotnie, podane oszacowanie
zapewnia ciągłość na każdym etapie, a więc na całej
$\mathcal D(\Omega)$.
\end{proof}

\begin{definicja}[Rząd dystrybucji na zwartym zbiorze]
\label{def:order19}
Dla $T\in\mathcal D'(\Omega;\K)$, zwartego $K\Subset\Omega$
i $m\in\N_0$ mówimy,
że $T$ ma \emph{rząd nie większy niż $m$ na $K$}, jeżeli
istnieje $C_K\ge0$ takie, że
$|\langle T,\varphi\rangle|\le C_K p_m(\varphi)$
dla $\varphi\in\mathcal D_K$. Mówimy, że $T$ ma
\emph{rząd skończony na $\Omega$}, jeżeli można dobrać
jedno $m$ właściwe dla wszystkich $K\Subset\Omega$
(przy czym $C_K$ może zależeć od $K$).
Dokładnym rzędem na $K$ nazywamy najmniejszą liczbę
$m\in\N_0$, dla której istnieje $C_K\ge0$ spełniająca
powyższe oszacowanie dla wszystkich $\varphi\in\mathcal D_K$.
Istnienie takiego minimum wynika z
Twierdzenia~\ref{thm:distribution-estimate19} i dobrego
uporządkowania $\N_0$.
\end{definicja}

\begin{uwaga}[Zależność od zbioru zwartego]
\label{rem:order-local19}
Twierdzenie~\ref{thm:distribution-estimate19} gwarantuje
\emph{lokalnie} skończony rząd każdej dystrybucji,
nie zaś istnienie jednego rzędu na całej $\Omega$.
Przykład~\ref{ex:variable-order18} konstruuje dystrybucję
$\langle T,\varphi\rangle=\sum_{n\ge1}\varphi^{(n)}(n)$
na $\R$, która nie ma skończonego rzędu globalnego.
\end{uwaga}

\begin{przyklad}[Dystrybucje regularne]
\label{ex:regular19}
Jeżeli $f\in L^1_{\mathrm{loc}}(\Omega;\K)$, to
\[\langle T_f,\varphi\rangle=\int_\Omega
f(x)\varphi(x)\,\mathrm d\lambda_d(x)\]
określa funkcjonał liniowy, niezależny od wyboru reprezentanta
klasy $f$. Dla $\varphi\in\mathcal D_K$ całka jest bezwzględnie
zbieżna i zachodzi
\[|\langle T_f,\varphi\rangle|\le
\left(\int_K|f(x)|\,\mathrm d\lambda_d(x)\right)p_0(\varphi).\]
Na mocy Twierdzenia~\ref{thm:distribution-estimate19}
jest to dystrybucja. Tak więc dystrybucje regularne
mają rząd nie większy niż $0$.
Zwróćmy uwagę, że nie jest potrzebne założenie
$f\in L^1(\Omega)$; np. $f(x)=e^{|x|^2}$ na $\R^d$
definiuje dystrybucję.
\end{przyklad}

\begin{stwierdzenie}[Wierność utożsamienia funkcji z dystrybucją]
\label{prop:regular-injective19}
Jeśli $f,g\in L^1_{\mathrm{loc}}(\Omega;\K)$ i $T_f=T_g$,
to $f=g$ prawie wszędzie w $\Omega$ względem miary Lebesgue'a.
\end{stwierdzenie}
\begin{proof}
Z liniowości całki $T_f-T_g=T_{f-g}$, więc po zastąpieniu
$f$ przez $f-g$ wystarczy rozważyć przypadek $T_f=0$.
Wybierzmy niezerową, nieujemną funkcję
$\rho_0\in C_c^\infty(B(0,1);\R)$, korzystając
z Lematu~\ref{lem:cutoffs17}, i przedłużmy ją zerem.
Jej całka jest dodatnia, gdyż $\rho_0$ jest ciągła
i dodatnia na pewnym niepustym zbiorze otwartym.
Niech $\rho=\rho_0/\int_{\R^d}\rho_0\,\mathrm d\lambda_d$.
Wówczas $\supp\rho\subseteq\overline{B(0,1)}$ i
$\int_{\R^d}\rho(x)\,\mathrm d\lambda_d(x)=1$. Określmy
$\rho_\varepsilon(x)=\varepsilon^{-d}\rho(x/\varepsilon)$.
Dla $x\in\Omega$ oraz $\varepsilon>0$ takich, że
$\overline B(x,\varepsilon)\subseteq\Omega$, definiujemy
$(f*\rho_\varepsilon)(x)=\int_\Omega f(y)\rho_\varepsilon(x-y)
\,\mathrm d\lambda_d(y)$.
Funkcja $y\mapsto\rho_\varepsilon(x-y)$ należy do
$\mathcal D(\Omega;\K)$, zatem
$(f*\rho_\varepsilon)(x)=0$.

Na każdym $K\Subset\Omega$ funkcje
$f*\rho_\varepsilon$ zbiegają do $f$ w $L^1(K)$.
Dla niepustego $K$ wybierzmy otwarty $V$ taki, że
$K\subseteq V$ i $\overline V\Subset\Omega$.
Można go skonstruować jako skończoną sumę kul o domknięciach
zawartych w $\Omega$, pokrywających $K$.
Przez $g$ oznaczmy przedłużenie $f\chi_V$ zerem na $\R^d$.
Wtedy $g\in L^1(\R^d)$, gdyż $f$ jest całkowalna na
$\overline V$. Odległość $K$ od $\R^d\setminus V$ jest dodatnia.
Jeśli $\varepsilon$ jest od niej mniejsza, to
$\overline B(x,\varepsilon)\subseteq V$ dla wszystkich $x\in K$,
więc sploty $f*\rho_\varepsilon$ i $g*\rho_\varepsilon$
są równe na $K$. Ponadto, po podstawieniu $y=x-\varepsilon z$,
z normalizacji $\rho$ i twierdzenia Tonellego dostajemy
\[\|g*\rho_\varepsilon-g\|_{L^1(\R^d)}\le
\int_{\R^d}\rho(z)\,
\|g(\cdot-\varepsilon z)-g\|_{L^1(\R^d)}\,
\mathrm d\lambda_d(z)\longrightarrow0,\]
co wynika z ciągłości translacji w $L^1$ i twierdzenia
o zbieżności zmajoryzowanej: dla każdego $z$
norma różnicy translacji dąży do zera, a cała funkcja
podcałkowa jest ograniczona przez
$2\|g\|_{L^1}\rho(z)$, która jest całkowalna.
Ponieważ $g=f$ prawie wszędzie na $K$ i
$f*\rho_\varepsilon=0$ na $K$, otrzymujemy
$\int_K|f|\,\mathrm d\lambda_d=0$.
Stosujemy to do przeliczalnego wyczerpania $(K_n)$
z Lematu~\ref{lem:exhaustion17}. Każdy zbiór
$A_n=\{x\in K_n:f(x)\ne0\}$ ma miarę zero,
a $\{x\in\Omega:f(x)\ne0\}=\bigcup_n A_n$ dla ustalonego
reprezentanta $f$. Z przeliczalnej podaddytywności miary
wynika $f=0$ prawie wszędzie w $\Omega$.
\end{proof}

\begin{przyklad}[Dystrybucje punktowe]
\label{ex:delta19}
Dla $a\in\Omega$ i wielowskaźnika $\alpha\in\N_0^d$
funkcjonały
\[\langle\delta_a,\varphi\rangle=\varphi(a),\qquad
\langle T_{a,\alpha},\varphi\rangle
=(-1)^{|\alpha|}\partial^\alpha\varphi(a)\]
są liniowe, a z oszacowania
$|T_{a,\alpha}(\varphi)|\le p_{|\alpha|}(\varphi)$
wynika ich ciągłość (dla $\alpha=0$ otrzymujemy $\delta_a$).
Mają rząd nie większy niż $|\alpha|$.
Żaden z tych funkcjonałów nie jest postaci $T_f$
dla $f\in L^1_{\mathrm{loc}}(\Omega;\K)$.
Istotnie, wybierzmy $\eta\in C_c^\infty(B(0,1);\R)$
równą $1$ w otoczeniu zera (Lemat~\ref{lem:cutoffs17}).
Przedłużmy $\eta$ zerem na $\R^d$ i połóżmy
$\psi(x)=x^\alpha\eta(x)$ dla $x\in\R^d$.
Wtedy $\partial^\alpha\psi(0)=\alpha!\ne0$.
Dla $0<\varepsilon<1$ tak małego, że
$\overline B(a,\varepsilon)\subseteq\Omega$, połóżmy
$\varphi_\varepsilon(x)=\varepsilon^{|\alpha|}
\psi((x-a)/\varepsilon)$. Wtedy
$T_{a,\alpha}(\varphi_\varepsilon)=(-1)^{|\alpha|}\alpha!$
nie zależy od $\varepsilon$, natomiast dla każdego
$f\in L^1_{\mathrm{loc}}(\Omega)$ mamy
$|T_f(\varphi_\varepsilon)|\le
\varepsilon^{|\alpha|}\|\psi\|_\infty
\int_{B(a,\varepsilon)}|f(x)|\,\mathrm d\lambda_d(x)
\to0$.
Ostatnia granica wynika z bezwzględnej ciągłości całki
funkcji $|f|$ na ustalonej zwartej kuli wokół $a$,
ponieważ $\lambda_d(B(a,\varepsilon))\to0$.
Argument obejmuje również $\alpha=0$, czyli $\delta_a$.
Znak wybrano zgodnie z definicją
pochodnej dystrybucyjnej w kolejnym rozdziale:
$T_{a,\alpha}=\partial^\alpha\delta_a$.
\end{przyklad}

\begin{stwierdzenie}[Dokładny rząd pochodnej delty]
\label{prop:delta-order19}
Niech $a\in\Omega$, $\alpha\in\N_0^d$ i niech
$T_{a,\alpha}$ będzie dystrybucją z Przykładu~\ref{ex:delta19}.
Jeśli zwarty $K\Subset\Omega$ zawiera otoczenie punktu $a$,
to $T_{a,\alpha}$ ma na $K$ dokładnie rząd $|\alpha|$.
\end{stwierdzenie}
\begin{proof}
Oszacowanie rzędu $|\alpha|$ wynika z definicji.
Dla $\alpha=0$ rząd nie może być mniejszy niż $0$.
Niech zatem $m\in\N_0$, $m<|\alpha|$.
Wybierzmy $\psi$ jak w Przykładzie~\ref{ex:delta19}.
Dla małego $\varepsilon>0$ funkcja
$\varphi_\varepsilon(x)=\varepsilon^{|\alpha|}
\psi((x-a)/\varepsilon)$ ma nośnik w $K$. Ponadto
$|T_{a,\alpha}(\varphi_\varepsilon)|
=|\partial^\alpha\psi(0)|>0$, podczas gdy
dla każdego $|\beta|\le m$ zachodzi
$\partial^\beta\varphi_\varepsilon(x)=
\varepsilon^{|\alpha|-|\beta|}(\partial^\beta\psi)((x-a)/\varepsilon)$.
Biorąc supremum i maksimum otrzymujemy
$p_m(\varphi_\varepsilon)\le C_m\varepsilon^{|\alpha|-m}
\to0$, gdzie $C_m=\max_{|\beta|\le m}\|\partial^\beta\psi\|_\infty$.
Nie może więc zachodzić oszacowanie rzędu $m$.
\end{proof}

\begin{przyklad}[Wartość główna Cauchy'ego]
\label{ex:principal-value19}
Dla $\varphi\in\mathcal D(\R;\K)$ definiujemy
\[\left\langle\operatorname{pv}\frac1x,\varphi\right\rangle
=\lim_{\varepsilon\downarrow0}
\int_{\{x\in\R:|x|>\varepsilon\}}
\frac{\varphi(x)}{x}\,\mathrm d\lambda_1(x)
=\int_0^\infty\frac{\varphi(t)-\varphi(-t)}{t}\,
\mathrm d\lambda_1(t).\]
Równość wynika z podstawienia $x=-t$ w całce
po $(-\infty,-\varepsilon)$ i przejścia do granicy,
uzasadnionego poniższym oszacowaniem.
Dla $\supp\varphi\subseteq[-R,R]$ oraz $R\ge1$
z tożsamości $\varphi(t)-\varphi(-t)=\int_{-t}^t\varphi'(s)\,\mathrm ds$
(w tym dla funkcji zespolonych) wynika
$|\varphi(t)-\varphi(-t)|\le2t\,p_1(\varphi)$ dla
$0<t\le1$, natomiast dla $1\le t\le R$ licznik
ma moduł nie większy niż $2p_0(\varphi)$. Dlatego
\[\left|\left\langle\operatorname{pv}\frac1x,
\varphi\right\rangle\right|\le
2p_1(\varphi)+2(\log R)p_0(\varphi).\]
W szczególności całka różnicy jest bezwzględnie zbieżna,
a granica istnieje. Dla każdego zwartego $K\subseteq\R$
wybieramy $R\ge1$ z $K\subseteq[-R,R]$ i stosujemy
$p_0\le p_1$. Liniowość wynika z liniowości całki,
a ciągłość z Twierdzenia~\ref{thm:distribution-estimate19}.
Jest to dystrybucja, chociaż $1/x\notin L^1_{\mathrm{loc}}(\R)$.
Symbol $\operatorname{pv}$ oznacza \emph{wartość główną},
a nie zwykłą, bezwzględnie zbieżną całkę.
\end{przyklad}

\subsection{Słaba-* i silna topologia na przestrzeni dystrybucji}

\begin{definicja}[Topologia słaba-*]
\label{def:weakstar19}
Na $\mathcal D'(\Omega)$ definiujemy
$\sigma(\mathcal D'(\Omega),\mathcal D(\Omega))$ jako
najsłabszą lokalnie wypukłą topologię, dla której
wszystkie funkcjonały $T\mapsto\langle T,\varphi\rangle$,
$\varphi\in\mathcal D(\Omega)$, są ciągłe.
Jest ona wyznaczona przez seminormy
$q_\varphi(T)=|\langle T,\varphi\rangle|$ dla $\varphi\in\mathcal D(\Omega)$.
Rodzina ta oddziela punkty, gdyż dla $T\ne S$ istnieje
$\varphi$ z $\langle T-S,\varphi\rangle\ne0$.
Topologia słaba-* jest zatem Hausdorffa.
Sieć $(T_i)$ zbiega w tej topologii do $T$ dokładnie wtedy,
gdy $\langle T_i,\varphi\rangle\to\langle T,\varphi\rangle$
dla każdej $\varphi\in\mathcal D(\Omega)$.
\end{definicja}

\begin{definicja}[Topologia silna]
\label{def:strong19}
Na $\mathcal D'(\Omega)$ definiujemy
\emph{silną topologię przestrzeni dualnej}
$\beta(\mathcal D'(\Omega),\mathcal D(\Omega))$
jako topologię zbieżności jednostajnej na zbiorach
ograniczonych $B\subseteq\mathcal D(\Omega)$.
Wyznaczają ją seminormy
\[q_B(T)=\sup_{\varphi\in B}|\langle T,\varphi\rangle|.\]
Dla $B=\varnothing$ przyjmujemy $q_B=0$.
Wartości $q_B$ są skończone, ponieważ ciągły funkcjonał
liniowy przenosi zbiory ograniczone na ograniczone.
Dla $B=\{\varphi\}$ otrzymujemy $q_B=q_\varphi$,
zatem topologia silna jest nie słabsza od słabej-*
i także jest Hausdorffa. Sieć $(T_i)$ zbiega silnie do $T$
wtedy i tylko wtedy, gdy $q_B(T_i-T)\to0$
dla każdego ograniczonego $B\subseteq\mathcal D(\Omega)$.
\end{definicja}

\begin{twierdzenie}[Beczkowość przestrzeni testowej]
\label{thm:D-barrelled19}
Przestrzeń $\mathcal D(\Omega;\K)$ jest beczkowa.
Każda rodzina $\mathcal A\subseteq\mathcal D'(\Omega;\K)$
ograniczona punktowo, tj. spełniająca
\[\sup_{T\in\mathcal A}|\langle T,\varphi\rangle|<\infty
\quad\text{dla każdego }\varphi\in\mathcal D(\Omega),\]
jest równociągła.
\end{twierdzenie}
\begin{proof}
Niech $U\subseteq\mathcal D(\Omega)$ będzie
beczką, tj. zbiorem domkniętym, absolutnie wypukłym
i pochłaniającym. Ustalmy wyczerpanie $(K_n)$ z Lematu~\ref{lem:exhaustion17}.
Zbiór $U\cap\mathcal D_{K_n}$ jest domknięty dzięki ciągłości
włożenia etapu, a absolutna wypukłość jest dziedziczona
przez przecięcie z podprzestrzenią liniową. Dla każdego
$\varphi\in\mathcal D_{K_n}$ istnieje $t>0$ z $\varphi\in tU$,
więc $\varphi/t\in U\cap\mathcal D_{K_n}$.
Jest to zatem beczka w przestrzeni Frécheta $\mathcal D_{K_n}$,
która jest beczkowa na mocy Wniosku~\ref{cor:frechet-barrelled}. Zatem
$U\cap\mathcal D_{K_n}$ jest otoczeniem zera
w $\mathcal D_{K_n}$.
Funkcjonał Minkowskiego $p_U$ jest seminormą, gdyż $U$
jest absolutnie wypukły i pochłaniający.
Na $U\cap\mathcal D_{K_n}$ mamy $p_U\le1$, zatem dla
każdego $\varepsilon>0$ seminorma $p_U$ jest mniejsza niż
$\varepsilon$ na otoczeniu $(\varepsilon/2)(U\cap\mathcal D_{K_n})$.
Jej ograniczenie do każdego etapu jest więc ciągłe. Własność uniwersalna
granicy induktywnej pokazuje, że $p_U$ jest ciągłą
seminormą na $\mathcal D(\Omega)$; zbiór
$\{\varphi:p_U(\varphi)<1\}\subseteq U$ jest otoczeniem
zera. Przestrzeń $\mathcal D(\Omega)$ jest więc beczkowa.
Druga część wynika z Twierdzenia~\ref{thm:banach-steinhaus-lcs}
dla rodziny ciągłych operatorów liniowych o wartościach w $\K$.
\end{proof}

\begin{wniosek}[Ograniczoność i równociągłość rodzin dystrybucji]
\label{cor:bounded-distributions19}
Dla rodziny $\mathcal A\subseteq\mathcal D'(\Omega)$
następujące warunki są równoważne:
\begin{enumerate}
\item $\mathcal A$ jest punktowo ograniczona na $\mathcal D(\Omega)$;
\item $\mathcal A$ jest równociągła na $\mathcal D(\Omega)$;
\item $\mathcal A$ jest ograniczona w silnej topologii
$\beta(\mathcal D',\mathcal D)$.
\end{enumerate}
\end{wniosek}
\begin{proof}
Pierwszy warunek implikuje drugi przez
Twierdzenie~\ref{thm:D-barrelled19}. Jeżeli
$\mathcal A$ jest równociągła, to dla każdego
$\varepsilon>0$ istnieje otoczenie $V$ zera takie,
że $|\langle T,\varphi\rangle|\le\varepsilon$ dla
$T\in\mathcal A$ i $\varphi\in V$.
Każdy zbiór ograniczony $B$ leży w $\lambda V$
dla pewnego $\lambda>0$, więc
$\sup_{T\in\mathcal A}q_B(T)\le\lambda\varepsilon$.
To daje ograniczoność silną. Wreszcie singletony
$\{\varphi\}$ są ograniczone, więc trzeci warunek
implikuje pierwszy.
\end{proof}

\begin{wniosek}[Jednostajne lokalne oszacowania rodziny dystrybucji]
\label{cor:uniform-order19}
Jeżeli rodzina $\mathcal A\subseteq\mathcal D'(\Omega)$ jest
punktowo ograniczona, to dla każdego $K\Subset\Omega$
istnieją $m_K\in\N_0$ i $C_K\ge0$ takie, że
$|\langle T,\varphi\rangle|\le C_Kp_{m_K}(\varphi)$
jednocześnie dla wszystkich $T\in\mathcal A$
i $\varphi\in\mathcal D_K$.
\end{wniosek}
\begin{proof}
Z Twierdzenia~\ref{thm:D-barrelled19} rodzina $\mathcal A$
jest równociągła na $\mathcal D(\Omega)$, a więc jej
ograniczenia do $\mathcal D_K$ są równociągłe.
Istnieje zatem otoczenie zera $U$ w $\mathcal D_K$,
na którym $|\langle T,\varphi\rangle|\le1$ dla wszystkich
$T\in\mathcal A$. Ponieważ seminormy $p_m$ są rosnące,
istnieją $m_K$ i $\varepsilon_K>0$ takie, że
$\{\varphi:p_{m_K}(\varphi)<\varepsilon_K\}\subseteq U$.
Jeśli $p_{m_K}(\varphi)>0$, podstawiamy
$a\varphi$ dla $a=\varepsilon_K/(2p_{m_K}(\varphi))$.
Wówczas $a|\langle T,\varphi\rangle|\le1$
jednocześnie dla wszystkich $T\in\mathcal A$,
co daje tezę z $C_K=2/\varepsilon_K$.
Jeśli $p_{m_K}(\varphi)=0$, to $a\varphi\in U$
dla każdego $a>0$, więc $\langle T,\varphi\rangle=0$
dla wszystkich $T\in\mathcal A$.
\end{proof}

\begin{twierdzenie}[Zbieżność ciągów: słaba-* = silna]
\label{thm:weak-strong-sequences19}
Niech $(T_n)_{n\ge1}\subseteq\mathcal D'(\Omega;\K)$
i $T\in\mathcal D'(\Omega;\K)$. Następujące warunki
są równoważne:
\begin{enumerate}
\item $T_n\to T$ w $\sigma(\mathcal D',\mathcal D)$;
\item $T_n\to T$ w $\beta(\mathcal D',\mathcal D)$.
\end{enumerate}
\end{twierdzenie}
\begin{proof}
Implikacja (2)$\Rightarrow$(1) wynika z definicji.
Zakładamy (1). Rodzina $\{T_n-T:n\ge1\}$ jest
punktowo ograniczona, gdyż dla ustalonej funkcji
testowej ciąg liczb $\langle T_n-T,\varphi\rangle$
jest zbieżny. Z Wniosku~\ref{cor:bounded-distributions19}
rodzina ta jest równociągła. Ustalmy zbiór ograniczony
$B\subseteq\mathcal D(\Omega)$ i $\varepsilon>0$.
Jego domknięcie $K=\overline B^{\,\mathcal D}$ jest zwarte
na mocy Twierdzenia~\ref{thm:D-Montel18}.
Jeśli $B=\varnothing$, to $q_B=0$, więc załóżmy $B\ne\varnothing$.
Z równociągłości wybierzmy otwarte, zrównoważone otoczenie
zera $V$ takie, że
$|\langle T_n-T,v\rangle|<\varepsilon/2$
za każdym razem, gdy $n\ge1$ i $v\in V$.
Zwartość $K$ pozwala wybrać
$\varphi_1,\ldots,\varphi_N\in K$ z
$K\subseteq\bigcup_{j=1}^N(\varphi_j+V)$.
Zbieżność punktowa zapewnia indeks $n_0$ taki, że
$|\langle T_n-T,\varphi_j\rangle|<\varepsilon/2$
dla wszystkich $n\ge n_0$ i $j\le N$.
Dla dowolnego $\varphi\in B$ piszemy
$\varphi=\varphi_j+v$ z $v\in V$, a więc dla $n\ge n_0$ z liniowości i nierówności trójkąta
$|\langle T_n-T,\varphi\rangle|
\le|\langle T_n-T,\varphi_j\rangle|+|\langle T_n-T,v\rangle|
<\varepsilon$.
Biorąc supremum, otrzymujemy $q_B(T_n-T)\le\varepsilon$
dla $n\ge n_0$. Ponieważ $\varepsilon>0$ jest dowolne,
$q_B(T_n-T)\to0$, co kończy dowód.
\end{proof}

\begin{uwaga}[Dlaczego nie wystarczy zbieżność sieci?]
\label{rem:net-distinction19}
Powyższy wynik dotyczy ciągów. Dla ogólnej sieci
zbieżność punktowa nie musi zapewniać punktowej
ograniczoności \emph{całej} rodziny jej wyrazów.
Nie wolno więc mechanicznie zastępować w tym dowodzie
ciągu przez sieć. Poniżej wykażemy, że obie topologie
są istotnie różne.
\end{uwaga}

\begin{twierdzenie}[Ścisła różnica topologii]
\label{thm:weak-strong-different19}
Jeżeli $\Omega\subseteq\R^d$ jest niepustym zbiorem
otwartym, to na $\mathcal D'(\Omega)$ topologia
$\beta(\mathcal D',\mathcal D)$ jest ściśle silniejsza
od $\sigma(\mathcal D',\mathcal D)$. Istnieje sieć
dystrybucji zbieżna słabo-* do zera, która nie
zbiega silnie do zera.
\end{twierdzenie}
\begin{proof}
Wybierzmy $a\in\Omega$ i $r>0$ tak, że
$\overline B(a,r)\Subset\Omega$. Definiujemy
$\chi(x)=\exp(-1/(r^2-|x-a|^2))$ dla $|x-a|<r$
i $\chi(x)=0$ poza kulą.
Gładkość na brzegu wynika z zanikania wszystkich pochodnych
funkcji $t\mapsto e^{-1/t}$ przy $t\downarrow0$,
jak w dowodzie Lematu~\ref{lem:cutoffs17}.
Mamy $\supp\chi=\overline B(a,r)$ i $\chi>0$ w $B(a,r)$.
Dla $j\ge0$ połóżmy
$\psi_j(x)=2^{-j}((x_1-a_1)/r)^j\chi(x)$.
Niech $u(x)=(x_1-a_1)/r$. Dla $0\le k\le j$ mamy
$\partial_1^k u^j=r^{-k}j!/(j-k)!\,u^{j-k}$,
a dla $k>j$ pochodna jest zerowa.
Na nośniku $\chi$ zachodzi $|u|\le1$.
Każda pochodna $u^j$ zawierająca różniczkowanie względem
którejkolwiek z pozostałych współrzędnych jest zerowa.
Z reguły Leibniza i ograniczoności pochodnych $\chi$
wynika więc $p_m(\psi_j)\le C_m 2^{-j}(1+j)^m$,
gdzie $C_m$ nie zależy od $j$. Ciąg po prawej stronie
jest ograniczony dla każdego $m$, gdyż wykładniczy
spadek dominuje dowolną ustaloną potęgę.
Ponadto $\supp\psi_j\subseteq\overline B(a,r)$,
więc z Twierdzenia~\ref{thm:D-bounded18}
zbiór $B_0=\{\psi_j:j\ge0\}$ jest ograniczony w $\mathcal D$.
Jeżeli $\sum_{j=0}^N c_j\psi_j=0$, to na kuli dzielimy
przez dodatnią $\chi$ i ograniczamy się do $x=a+rt e_1$,
$-1<t<1$. Otrzymujemy wielomian
$\sum_{j=0}^N c_j2^{-j}t^j=0$ dla każdego $t\in(-1,1)$.
Wszystkie $c_j$ są zerowe, więc funkcje są liniowo niezależne.

Niech $I$ będzie zbiorem par $(F,n)$, gdzie $F$ jest
skończonym podzbiorem $\mathcal D(\Omega)$, a $n\in\N$,
z porządkiem $(F,n)\preceq(G,m)$ wtedy i tylko wtedy,
gdy $F\subseteq G$ i $n\le m$. Jest to zbiór skierowany:
wspólnym ograniczeniem górnym par $(F,n)$ i $(G,m)$
jest $(F\cup G,\max\{n,m\})$.
Dla każdej pary wybieramy $\psi\in B_0\setminus\operatorname{lin}(F)$,
co jest możliwe, gdyż $B_0$ zawiera nieskończenie wiele
liniowo niezależnych funkcji, a $\operatorname{lin}(F)$
jest skończenie wymiarowa.
Skończenie wymiarowa podprzestrzeń
$\operatorname{lin}(F)$ jest domknięta, zatem
z Twierdzenia~\ref{thm:point-closed-subspace}
istnieje ciągły funkcjonał liniowy $L$ z
$L|_{\operatorname{lin}(F)}=0$ i $L(\psi)\ne0$.
Definiujemy $T_{F,n}=(n/L(\psi))L$. Wtedy
$T_{F,n}|_{\operatorname{lin}(F)}=0$ i
$\langle T_{F,n},\psi\rangle=n$.

Dla ustalonej $\varphi\in\mathcal D(\Omega)$ i
$(F,n)\succeq(\{\varphi\},1)$ mamy
$\langle T_{F,n},\varphi\rangle=0$. Sieć
$(T_{F,n})_{(F,n)\in I}$ zbiega więc słabo-* do zera.
Jednak $q_{B_0}(T_{F,n})\ge n$ dla każdej pary.
Dla każdego $M>0$, wybierając liczbę całkowitą $n_0>M$,
otrzymujemy $q_{B_0}(T_{F,n})>M$ dla
$(F,n)\succeq(\varnothing,n_0)$. Seminorma dąży więc
do nieskończoności, a sieć nie zbiega silnie do zera. Ponieważ topologia silna
jest nie słabsza od słabej-*, a obie mają różne sieci zbieżne,
jest od niej ściśle silniejsza.
\end{proof}

\subsection{Granice regularizacji i zastosowania}

\begin{przyklad}[Delty uciekające do nieskończoności]
\label{ex:delta-escape19}
Niech $\Omega=\R^d$, a $(a_n)\subseteq\R^d$
spełnia $|a_n|\to\infty$. Wtedy $\delta_{a_n}\to0$
\emph{silnie} w $\mathcal D'(\R^d)$.
Istotnie, każdy ograniczony $B\subseteq\mathcal D(\R^d)$
ma wspólny zwarty nośnik $K$, zatem dla dużych $n$
$\sup_{\varphi\in B}|\varphi(a_n)|=0$.
\end{przyklad}

\begin{przyklad}[Ciągłość delty względem położenia]
\label{ex:delta-continuity19}
Jeżeli $a_n\to a$ w $\Omega$, to
$\delta_{a_n}\to\delta_a$ silnie w $\mathcal D'(\Omega)$.
Dla ograniczonego $B\subseteq\mathcal D(\Omega)$
i $n$ dostatecznie dużych odcinek $[a,a_n]$ leży
w $\Omega$, ponieważ $a_n$ od pewnego miejsca
leży w kuli o środku $a$ zawartej w $\Omega$.
Dla każdej $\varphi\in B$ stosujemy tożsamość
$\varphi(a_n)-\varphi(a)=\int_0^1
\sum_{j=1}^d(a_{n,j}-a_j)\partial_j\varphi(a+t(a_n-a))\,
\mathrm dt$.
Nierówność Cauchy'ego--Schwarza obowiązuje również
dla zespolonych wartości pochodnych. Biorąc supremum po $B$,
otrzymujemy
\[q_B(\delta_{a_n}-\delta_a)\le
|a_n-a|\sqrt d\sup_{\varphi\in B}p_1(\varphi),
\]
gdyż $|\nabla\varphi(x)|\le\sqrt d\,p_1(\varphi)$.
Prawa strona dąży do zera na mocy kryterium
ograniczoności z Twierdzenia~\ref{thm:D-bounded18}.
\end{przyklad}

\begin{twierdzenie}[Aproksymacja jednostki i zbieżność do delty]
\label{thm:mollifier-delta19}
Niech $a\in\Omega$ i niech
$\rho\in C_c^\infty(\R^d)$ ma nośnik w
$\overline{B(0,1)}$, spełnia $\rho\ge0$ i
$\int_{\R^d}\rho(x)\,\mathrm d\lambda_d(x)=1$.
Dla $n\ge1$ określmy
$f_n(x)=n^d\rho(n(x-a))$ dla $x\in\Omega$.
Wtedy $T_{f_n}\to\delta_a$ silnie w $\mathcal D'(\Omega)$.
Co więcej, istnieje stała $C_\rho\ge0$ zależna tylko od
$\rho$ i $d$, taka że dla każdego ograniczonego zbioru
$B\subseteq\mathcal D(\Omega)$ i wszystkich dostatecznie
dużych $n$ (niezależnie od $B$) zachodzi
\[q_B(T_{f_n}-\delta_a)\le
\frac{C_\rho}{n}\sup_{\varphi\in B}p_1(\varphi).\]
\end{twierdzenie}
\begin{proof}
Wybieramy $r>0$ z $\overline B(a,r)\subseteq\Omega$,
a następnie $n_0$ tak, że $1/n_0<r$.
Dla $n\ge n_0$ podstawienie $y=n(x-a)$ daje
\[\langle T_{f_n},\varphi\rangle
=\int_{\overline B(0,1)}\rho(y)\varphi(a+y/n)\,
\mathrm d\lambda_d(y).\]
Całkujemy po kuli, gdyż poza nią $\rho=0$,
a dla $y$ w tej kuli argument $\varphi$ należy do $\Omega$.
Dla $y\in\supp\rho$ oraz dostatecznie dużego $n$
odcinek od $a$ do $a+y/n$ leży w $\Omega$. Zatem
z twierdzenia Newtona--Leibniza i nierówności Cauchy'ego--Schwarza
otrzymujemy
$|\varphi(a+y/n)-\varphi(a)|\le
\sqrt d\,|y|\,p_1(\varphi)/n$.
Ponieważ $\int\rho\,\mathrm d\lambda_d=1$,
po odjęciu $\varphi(a)$ całkujemy tę różnicę
z nieujemną wagą $\rho$.
Stąd teza z
$C_\rho=\sqrt d\int_{\R^d}|y|\rho(y)\,
\mathrm d\lambda_d(y)$. Ograniczoność $B$ zapewnia
$\sup_{\varphi\in B}p_1(\varphi)<\infty$, więc prawa
strona dąży do zera.
\end{proof}



\begin{twierdzenie}[Zbieżność w $L^1_{\mathrm{loc}}$ a zbieżność dystrybucji]
\label{thm:L1-to-strong19}
Niech $f_n,f\in L^1_{\mathrm{loc}}(\Omega;\K)$ i załóżmy,
że $f_n\to f$ w $L^1(K)$ dla każdego $K\Subset\Omega$.
Wówczas $T_{f_n}\to T_f$ w silnej topologii
$\beta(\mathcal D'(\Omega),\mathcal D(\Omega))$.
\end{twierdzenie}
\begin{proof}
Niech $B\subseteq\mathcal D(\Omega)$ będzie ograniczony.
Jeśli $B=\varnothing$, teza jest natychmiastowa.
Z Twierdzenia~\ref{thm:D-bounded18} istnieje
$K\Subset\Omega$ zawierający wszystkie nośniki funkcji
z $B$, przy czym $M_B=\sup_{\varphi\in B}p_0(\varphi)<\infty$.
Otrzymujemy
\[q_B(T_{f_n}-T_f)\le
M_B\int_K|f_n(x)-f(x)|\,\mathrm d\lambda_d(x)
\longrightarrow0.\]
Jest to dokładnie definicja zbieżności silnej.
\end{proof}

\begin{przyklad}[Silna zbieżność bez zbieżności lokalnej w $L^1$]
\label{ex:oscillations19}
Na $\Omega=\R$ weźmy $f_n(x)=\sin(nx)$.
Dla każdej $\varphi\in\mathcal D(\R)$ wybieramy przedział
$[-R,R]$ zawierający jej nośnik we wnętrzu.
Całkowanie przez części, z zerowymi wyrazami brzegowymi,
daje
\[\langle T_{f_n},\varphi\rangle=
\frac1n\int_\R\cos(nx)\varphi'(x)\,
\mathrm d\lambda_1(x).\]
Niech $B\subseteq\mathcal D(\R)$ będzie ograniczony.
Wspólny nośnik funkcji z $B$ mieści się w $[-R,R]$
dla pewnego $R>0$, a $\sup_{\varphi\in B}p_1(\varphi)<\infty$.
Dlatego
\[q_B(T_{f_n})\le\frac{2R}{n}
\sup_{\varphi\in B}p_1(\varphi)\longrightarrow0.\]
Zatem $T_{f_n}\to0$ \emph{silnie} w $\mathcal D'(\R)$,
ale $\int_{-1}^1|\sin(nx)|\,\mathrm d\lambda_1(x)$
nie dąży do zera. Istotnie, przez parzystość $|\sin|$
i podstawienie $t=nx$ otrzymujemy
$\int_{-1}^1|\sin(nx)|\,\mathrm d\lambda_1(x)
=(2/n)\int_0^n|\sin t|\,\mathrm d\lambda_1(t)$.
Zapisując $n=k\pi+r$ z $k=\lfloor n/\pi\rfloor$ i $0\le r<\pi$,
mamy $\int_0^n|\sin t|\,\mathrm d\lambda_1(t)=2k+\int_0^r\sin t\,\mathrm d\lambda_1(t)$.
Ostatnia całka leży w $[0,2]$, zaś $k/n\to1/\pi$,
zatem granica wynosi $4/\pi$.
Zbieżność dystrybucji jest więc znacznie słabsza niż
zbieżność reprezentujących je funkcji w $L^1_{\mathrm{loc}}$.
\end{przyklad}

% Materiał z AF 25; fragment do rozdziału o dystrybucjach.

\setcounter{section}{3}
% Kandydat restrukturyzacji; nie jest aktywnym rozdziałem Atlasu.
\section{Lokalność, nośniki i przedłużanie dystrybucji}
\providecommand{\atlasref}[2]{#2}

\subsection{Nośnik i lokalność}

\begin{definicja}[Zerowanie się dystrybucji na zbiorze otwartym]
\label{def:distribution-zero19}
Niech $T\in\mathcal D'(\Omega)$, a $U\subseteq\Omega$
będzie otwarty. Dla $U=\varnothing$ przyjmujemy
$\mathcal D(U;\K)=\{0\}$ i $T|_U=0$.
Dla niepustego $U$ i $\varphi\in\mathcal D(U;\K)$ oznaczmy przez
$e_U\varphi$ jej przedłużenie zerem na $\Omega$.
Jest ono gładkie, gdyż nośnik $\varphi$ jest zwarty
i zawarty w $U$. Na każdym etapie $\mathcal D_K(U)$
odwzorowanie $e_U$ zachowuje seminormy $p_m$ i ma obraz
w $\mathcal D_K(\Omega)$, więc jest ciągłe.
Z Twierdzenia~\ref{thm:LF-universal18} wynika ciągłość
$e_U:\mathcal D(U;\K)\to\mathcal D(\Omega;\K)$.
Definiujemy ograniczenie $T|_U$ przez
$\langle T|_U,\varphi\rangle=\langle T,e_U\varphi\rangle$.
Jest ono dystrybucją jako złożenie ciągłych odwzorowań liniowych.
Piszemy $T|_U=0$, jeśli powyższa wartość jest zerowa
dla każdej $\varphi\in\mathcal D(U;\K)$.
Nośnikiem $T$ jest domknięty względnie zbiór
\[\supp T=\Omega\setminus\bigcup\{U\subseteq\Omega:
U\text{ jest otwarty i }T|_U=0\}.\]
\end{definicja}

\begin{lemat}[Lokalność dystrybucji]
\label{lem:distribution-local19}
Niech $T\in\mathcal D'(\Omega)$ i
$\varphi\in\mathcal D(\Omega)$ spełnia
$\supp\varphi\cap\supp T=\varnothing$. Wtedy
$\langle T,\varphi\rangle=0$.
\end{lemat}
\begin{proof}
Jeśli $\varphi=0$, teza jest oczywista.
Niech $L=\supp\varphi\ne\varnothing$.
Każdy $x\in L$ ma otoczenie otwarte $U_x\subseteq\Omega$,
na którym $T$ się zeruje. Wybierzmy kulę $B_x$
o środku $x$ i dodatnim promieniu, z $\overline B_x\subseteq U_x$.
Ze zwartości $L$ wybieramy skończone pokrycie
$B_{x_1},\ldots,B_{x_N}$.
Z Lematu~\ref{lem:cutoffs17} wybieramy nieujemne
$\eta_j\in C_c^\infty(U_{x_j};\R)$ równe $1$
w otoczeniu $\overline B_{x_j}$ i przedłużamy je zerem.
Niech $h=\sum_{j=1}^N\eta_j$ i $O=\{x\in\Omega:h(x)>0\}$.
Zbiór $O$ jest otwarty i zawiera $L$.
Na $O$ definiujemy $\varphi_j=\eta_j\varphi/h$,
a na $\Omega\setminus O$ przyjmujemy $\varphi_j=0$.
Przedłużenie jest gładkie, ponieważ każdy punkt spoza $O$
ma otoczenie rozłączne z domkniętym zbiorem $L$,
na którym $\varphi_j$ jest identycznie zerowa.
Mamy $\supp\varphi_j\subseteq L\cap\supp\eta_j\Subset U_{x_j}$
oraz $\sum_j\varphi_j=\varphi$ na $O$ i poza $O$.
Zatem $\langle T,\varphi_j\rangle=0$ dla każdego $j$.
Sumując, otrzymujemy tezę.

W szczególności $T$ zeruje się na otwartym zbiorze
$\Omega\setminus\supp T$, gdyż każda funkcja testowa
na tym zbiorze, przedłużona zerem, ma nośnik rozłączny
z $\supp T$. Jest to więc największy otwarty zbiór
zerowania się $T$.
\end{proof}

\begin{stwierdzenie}[Nośniki przykładów]
\label{prop:support-examples19}
Dla $a\in\Omega$ i $\alpha\in\N_0^d$ mamy
$\supp\delta_a=\supp T_{a,\alpha}=\{a\}$.
Dla $f\in L^1_{\mathrm{loc}}(\Omega)$ nośnik $T_f$
jest dopełnieniem największego otwartego podzbioru
$U\subseteq\Omega$, na którym $f=0$ prawie wszędzie.
\end{stwierdzenie}
\begin{proof}
Gdy otwarty $U\subseteq\Omega$ nie zawiera $a$,
przedłużenie każdej funkcji z $\mathcal D(U)$ znika
w otoczeniu $a$, zatem wszystkie jej pochodne w $a$ są zerowe.
Jeśli $a\in U$, wybieramy $\eta\in C_c^\infty(U;\R)$
równą $1$ w otoczeniu $a$. Funkcja
$\varphi(x)=(x-a)^\alpha\eta(x)$ spełnia
$\partial^\alpha\varphi(a)=\alpha!\ne0$.
Stąd żadne otoczenie $a$ nie jest zbiorem zerowania
$T_{a,\alpha}$, także dla $\alpha=0$.
Dla $T_f$ mamy $(T_f)|_U=T_{f|_U}$ z definicji ograniczenia.
Jeśli $f=0$ prawie wszędzie w $U$, to $(T_f)|_U=0$.
Odwrotnie, ze Stwierdzenia~\ref{prop:regular-injective19}
na $U$ wynika, że $(T_f)|_U=0$ pociąga $f=0$ prawie wszędzie w $U$.
Największy taki otwarty zbiór istnieje na mocy Lematu
o lokalności, a jego dopełnienie jest nośnikiem $T_f$.
\end{proof}

\begin{twierdzenie}[Zwarty nośnik implikuje skończony rząd]
\label{thm:compact-support-order19}
Jeżeli $T\in\mathcal D'(\Omega;\K)$ i
$S=\supp T\Subset\Omega$, to istnieją
$m\in\N_0$, $C\ge0$ i $K\Subset\Omega$ takie, że
\[|\langle T,\varphi\rangle|\le C p_{K,m}(\varphi)
\quad(\varphi\in\mathcal D(\Omega)),\qquad
p_{K,m}(\varphi)=\max_{|\alpha|\le m}
\sup_{x\in K}|\partial^\alpha\varphi(x)|.\]
W szczególności $T$ ma skończony rząd na całej $\Omega$.
Ponadto $T$ przedłuża się w sposób jednoznaczny do
ciągłego funkcjonału liniowego na $C^\infty(\Omega;\K)$.
\end{twierdzenie}
\begin{proof}
Jeśli $S=\varnothing$, Lemat~\ref{lem:distribution-local19}
daje $T=0$. Możemy wtedy wybrać $m=0$, $C=0$
i dowolny niepusty zwarty $K\Subset\Omega$.
Przedłużeniem jest funkcjonał zerowy; jego jednoznaczność
wynika z gęstości $\mathcal D(\Omega)$ w $C^\infty(\Omega)$
na mocy Wniosku~\ref{cor:density-cinf17}.
W przeciwnym przypadku z Lematu~\ref{lem:cutoffs17}
wybieramy $\chi\in C_c^\infty(\Omega;\R)$ równą $1$
na pewnym otoczeniu $S$. Wtedy
$\supp((1-\chi)\varphi)\cap S=\varnothing$, więc z
Lematu~\ref{lem:distribution-local19} wynika
$\langle T,\varphi\rangle=\langle T,\chi\varphi\rangle$.
Niech $K=\supp\chi$. Z Twierdzenia~\ref{thm:distribution-estimate19}
istnieją $m,C_0$ takie, że
$|\langle T,\psi\rangle|\le C_0p_m(\psi)$ dla
$\psi\in\mathcal D_K$. Reguła Leibniza daje
$p_m(\chi\varphi)\le 2^m p_m(\chi)p_{K,m}(\varphi)$,
ponieważ wszystkie pochodne $\chi$ znikają poza $K$.
Otrzymujemy żądane oszacowanie z $C=C_0 2^m p_m(\chi)$.
Dla każdego zwartego $L\Subset\Omega$ i $\varphi\in\mathcal D_L$
mamy $p_{K,m}(\varphi)\le p_m(\varphi)$,
więc jeden rząd $m$ wystarcza dla wszystkich $L$.
Dla $u\in C^\infty(\Omega)$ definiujemy
$\widetilde T(u)=\langle T,\chi u\rangle$.
Odwzorowanie to jest liniowe, a powyższe oszacowanie
iloczynu obowiązuje również dla $u\in C^\infty(\Omega)$.
Zatem $|\widetilde T(u)|\le C p_{K,m}(u)$,
co zapewnia ciągłość w topologii $C^\infty(\Omega)$, a wzór z poprzedniego akapitu
pokazuje zgodność z $T$ na funkcjach testowych.
Przestrzeń $\mathcal D(\Omega)$ jest gęsta w
$C^\infty(\Omega)$ w tej topologii (Wniosek~\ref{cor:density-cinf17}),
więc przedłużenie jest jednoznaczne.
\end{proof}

\subsection{Lokalność i problemy przedłużania dystrybucji}

\begin{twierdzenie}[Ograniczenie i brak automatycznego przedłużenia]
\label{thm:restriction25}
Jeżeli $U\subseteq\Omega\subseteq\R^d$ są niepustymi
zbiorami otwartymi, to odwzorowanie ograniczenia
$\mathcal D'(\Omega)\to\mathcal D'(U)$ jest liniowe
i ciągłe w topologiach słabych-* oraz silnych.
Nie dla każdej pary $U\subseteq\Omega$ odwzorowanie
ograniczenia jest surjektywne. Nawet dla przedziałów
$U=(0,1)$ i $\Omega=(-1,1)$ istnieją dystrybucje
regularne na $U$, które nie mają przedłużenia na $\Omega$.
\end{twierdzenie}
\begin{proof}
Funkcję testową $\varphi\in\mathcal D(U)$ przedłużamy
przez zero do $\Omega$ i oznaczamy przedłużenie przez
$j\varphi=\widetilde\varphi$. Jeżeli
$\supp\varphi\subseteq L\Subset U$, to $L$ jest również
zwarty w $\Omega$, a każdy punkt $\Omega\setminus U$
ma otoczenie rozłączne z $L$. Przedłużenie jest więc
gładkie także na brzegu $U$ w $\Omega$ i
$\supp(j\varphi)=\supp\varphi$, gdzie oba nośniki
utożsamiamy jako podzbiory $\Omega$.
Na każdym niezerowym etapie $\mathcal D_L(U)$
odwzorowanie $j$ ma wartości w $\mathcal D_L(\Omega)$
i zachowuje seminormy $p_{L,m}$, ponieważ pochodne
na $U$ nie zmieniają się. Na etapie zerowym ciągłość
jest oczywista. Własność uniwersalna granicy induktywnej
(Twierdzenie~\ref{thm:LF-universal18}) zapewnia ciągłość
liniowego odwzorowania $j:\mathcal D(U)\to\mathcal D(\Omega)$.

Ograniczenie $R_U T=T|_U$ definiujemy przez
$\langle R_U T,\varphi\rangle=\langle T,j\varphi\rangle$.
Jest to ciągły funkcjonał na $\mathcal D(U)$ jako
złożenie ciągłych odwzorowań liniowych.
Dla każdego ograniczonego $B\subseteq\mathcal D(U)$
mamy dokładną równość
\[
 \sup_{\varphi\in B}|\langle R_U T,\varphi\rangle|
 =\sup_{\eta\in j(B)}|\langle T,\eta\rangle|.
\]
Ponieważ $j(B)$ jest ograniczony, otrzymujemy ciągłość
silną, a dla $B$ jednoelementowego także słabą-*.
Nie zakładamy tu gęstości obrazu $j$.

Dla otwartego $V\subseteq U$ kolejne przedłużenia przez
zero dają to samo odwzorowanie co bezpośrednie
przedłużenie z $V$ do $\Omega$. Zatem
$(T|_U)|_V=T|_V$. Z definicji nośnika i
Lematu~\ref{lem:distribution-local19} wynika
$\supp(T|_U)=U\cap\supp T$, z nośnikiem po lewej
stronie rozumianym względnie w $U$: punkt $x\in U$
nie należy do tego nośnika wtedy i tylko wtedy,
gdy istnieje otwarte otoczenie $V\subseteq U$ punktu
$x$, na którym $T|_V=0$. W szczególności
$\ker R_U=\{T\in\mathcal D'(\Omega):
\supp T\subseteq\Omega\setminus U\}$.
Jeżeli przedłużenie istnieje, jego jednoznaczność
nie jest na ogół zapewniona: do dowolnego przedłużenia
można dodać element tego jądra.

Aby wykazać drugie stwierdzenie, weźmy
$U=(0,1)$, $\Omega=(-1,1)$ i dystrybucję regularną
$T_f$ wyznaczoną przez $f(x)=\exp(1/x^2)$ na $U$.
Funkcja $f$ jest gładka na $U$, a więc lokalnie
całkowalna. Ustalmy $K=[-1/2,1/2]\Subset\Omega$.
Gdyby istniało przedłużenie
$\widetilde T\in\mathcal D'(\Omega)$ z
$\widetilde T|_U=T_f$, to z
Twierdzenia~\ref{thm:distribution-estimate19}
dla tego ustalonego $K$ istniałyby $m\in\N_0$
i $C>0$ takie, że
$|\langle\widetilde T,\varphi\rangle|\le C p_{K,m}(\varphi)$
dla wszystkich $\varphi\in\mathcal D_K(\Omega)$.
Niech $\psi\in\mathcal D((-1,1))$ będzie nieujemna
i niezerowa. Przedłużamy ją przez zero na $\R$.
Dla całkowitych $n\ge3$ przyjmujemy
$r_n=1/(10n^2)$ i
$\varphi_n(x)=\psi((x-1/n)/r_n)$.
Nośnik zawiera się w przedziale
$[1/n-r_n,1/n+r_n]\subset(0,1/2)$, zatem
$\varphi_n$ jest testem zarówno na $U$, jak i na
ustalonym etapie $\mathcal D_K(\Omega)$.
Dla każdego $0\le k\le m$ mamy
$\varphi_n^{(k)}(x)=r_n^{-k}\psi^{(k)}((x-1/n)/r_n)$.
Ponieważ $r_n\le1$, otrzymujemy
\[
 p_{K,m}(\varphi_n)
 \le 10^m n^{2m}\max_{0\le k\le m}
 \sup_{\R}|\psi^{(k)}|.
\]
Z drugiej strony dla $x=1/n+r_nt$, $|t|\le1$,
mamy $0<x\le(11/10)n^{-1}$, a więc
$1/x^2\ge(100/121)n^2\ge n^2/2$.
Podstawienie w całce i nieujemność $\psi$ dają
\[
 \begin{aligned}
 \langle T_f,\varphi_n\rangle
 &=r_n\int_{-1}^1 e^{1/(1/n+r_nt)^2}\psi(t)\,
 \mathrm d\lambda_1(t)\\
 &\ge \frac{e^{n^2/2}}{10n^2}
 \int_{-1}^1\psi(t)\,\mathrm d\lambda_1(t).
 \end{aligned}
\]
Stała $c=\frac1{10}\int_{-1}^1\psi(t)\,\mathrm d\lambda_1(t)$
jest ściśle dodatnia. Ponieważ
$\langle\widetilde T,\varphi_n\rangle=
\langle T_f,\varphi_n\rangle$, z oszacowania na
$\mathcal D_K(\Omega)$ wynikałoby
$c e^{n^2/2}\le C' n^{2m+2}$ dla wszystkich $n\ge3$,
gdzie $C'$ nie zależy od $n$. Jednak
$e^{n^2/2}\ge (n^2/2)^{m+2}/(m+2)!$,
więc po podzieleniu przez $n^{2m+2}$ lewa strona
rośnie co najmniej jak dodatnia stała razy $n^2$.
Otrzymujemy sprzeczność. Brak przedłużenia dotyczy
wszystkich dystrybucji, nie tylko dystrybucji regularnych.
\end{proof}

\setcounter{section}{4}
% Rozdział aktywny po restrukturyzacji DT.
% Wersja źródłowa: audyt II, jednostka 20

\section{Pochodne dystrybucyjne i transpozycja operatorów}
\label{sec:pochodne20}

\providecommand{\atlasref}[2]{#2}


\subsection{Definicja pochodnej i ciągłość w topologiach dualnych}

\noindent\emph{Powiązanie tematyczne w Bibliotece matematycznej:} \atlasref{PDE-07::def:pde7-derivative}{PDE-07: definicja pochodnej dystrybucyjnej}.


W tym rozdziale $d\ge1$ jest liczbą całkowitą,
$\Omega\subseteq\R^d$ jest niepustym zbiorem otwartym,
a $\K\in\{\R,\C\}$. Dystrybucje są
ciągłymi $\K$-liniowymi funkcjonałami na $\mathcal D(\Omega;\K)$, zgodnie
z konwencją przyjętą w rozdziale~\ref{sec:dystrybucje19}.
Dla wielowskaźnika $\alpha=(\alpha_1,\ldots,\alpha_d)\in\N_0^d$
piszemy $|\alpha|=\sum_j\alpha_j$,
$\partial^\alpha=\partial_1^{\alpha_1}\cdots\partial_d^{\alpha_d}$
oraz $\binom\alpha\beta=\prod_j\binom{\alpha_j}{\beta_j}$
dla $\beta\le\alpha$ współrzędnie. Przez $e_j$ oznaczamy
$j$-ty wektor bazy standardowej $\R^d$, a przez $\mathcal D_K$
i $p_m$ przestrzenie i seminormy z rozdziału~\ref{sec:funkcje-testowe18}.
Pominięcie $\K$ w oznaczeniu przestrzeni funkcji lub dystrybucji
oznacza, że nadal pracujemy nad ustalonym ciałem $\K$.

\begin{lemat}[Różniczkowanie funkcji testowych]
\label{lem:test-derivative20}
Dla każdego $\alpha\in\N_0^d$ odwzorowanie
$D_\alpha:\mathcal D(\Omega;\K)\to\mathcal D(\Omega;\K)$,
$D_\alpha\varphi=\partial^\alpha\varphi$, jest ciągłe i liniowe.
Ponadto dla każdego zbioru ograniczonego $B\subseteq\mathcal D(\Omega)$
zbiór $D_\alpha(B)$ jest ograniczony.
\end{lemat}
\begin{proof}
Dla każdego $K\Subset\Omega$ odwzorowanie $D_\alpha$ zachowuje
$\mathcal D_K(\Omega)$, ponieważ w otoczeniu każdego punktu
spoza $\supp\varphi$ funkcja $\varphi$ jest identycznie zerowa,
a więc wszystkie jej pochodne również są tam zerowe.
Przy seminormach z rozdziału~\ref{sec:funkcje-testowe18} mamy
$p_m(D_\alpha\varphi)\le p_{m+|\alpha|}(\varphi)$.
Ograniczenie $D_\alpha|_{\mathcal D_K}$ jest więc ciągłe.
Po złożeniu z ciągłym włożeniem $\mathcal D_K\to\mathcal D(\Omega)$
otrzymujemy ciągłość na każdym etapie, a z
Twierdzenia~\ref{thm:LF-universal18} ciągłość na całej $\mathcal D(\Omega)$.
Liniowość wynika z liniowości różniczkowania.
Jeśli $B$ jest ograniczony i $U$ jest otoczeniem zera
w $\mathcal D(\Omega)$, to $D_\alpha^{-1}(U)$ jest
otoczeniem zera. Istnieje $t>0$ takie, że
$B\subseteq tD_\alpha^{-1}(U)$, a z liniowości
$D_\alpha(B)\subseteq tU$. Zatem $D_\alpha(B)$ jest ograniczony.
\end{proof}

\begin{definicja}[Pochodna dystrybucyjna]
\label{def:dist-deriv20}
Niech $T\in\mathcal D'(\Omega;\K)$ i $\alpha\in\N_0^d$.
\emph{Pochodną dystrybucyjną} $\partial^\alpha T$ określamy wzorem
\[
 \langle\partial^\alpha T,\varphi\rangle
 =(-1)^{|\alpha|}\langle T,\partial^\alpha\varphi\rangle,
 \qquad\varphi\in\mathcal D(\Omega;\K).
\]
Prawa strona jest ciągłym funkcjonałem liniowym względem $\varphi$,
na mocy Lematu~\ref{lem:test-derivative20}.
Dla $j\in\{1,\ldots,d\}$ i $\alpha=e_j$ piszemy również $\partial_jT$.
\end{definicja}

\begin{twierdzenie}[Ciągłość i algebra pochodnych]
\label{thm:dist-deriv-cont20}
Niech $T\in\mathcal D'(\Omega;\K)$ oraz
$\alpha,\beta\in\N_0^d$. Wówczas:
\begin{enumerate}
\item $\partial^\alpha T$ jest dystrybucją dla każdego
$T\in\mathcal D'(\Omega)$;
\item $\partial^\alpha\partial^\beta T
=\partial^{\alpha+\beta}T$, w szczególności
$\partial_j\partial_kT=\partial_k\partial_jT$
dla $j,k\in\{1,\ldots,d\}$;
\item operator $\partial^\alpha:\mathcal D'(\Omega)\to
\mathcal D'(\Omega)$ jest ciągły zarówno w topologii słabej-*
$\sigma(\mathcal D',\mathcal D)$, jak i silnej
$\beta(\mathcal D',\mathcal D)$;
\item $\supp(\partial^\alpha T)\subseteq\supp T$.
Jeśli dla zwartego $K\Subset\Omega$ i $m\in\N_0$
dystrybucja $T$ ma na $K$ rząd nie większy niż $m$,
to $\partial^\alpha T$ ma na $K$ rząd nie większy niż
$m+|\alpha|$;
\item $\operatorname{sing\,supp}(\partial^\alpha T)
\subseteq\operatorname{sing\,supp}T$.
\end{enumerate}
\end{twierdzenie}
\begin{proof}
Punkt pierwszy wynika z ciągłości $D_\alpha$ na
$\mathcal D(\Omega)$ i złożenia funkcjonałów.
Dla każdej $\varphi\in\mathcal D(\Omega)$ dwukrotne
zastosowanie definicji daje
\[
\begin{aligned}
\langle\partial^\alpha\partial^\beta T,\varphi\rangle
&=(-1)^{|\alpha|+|\beta|}
  \langle T,\partial^\beta\partial^\alpha\varphi\rangle\\
&=(-1)^{|\alpha+\beta|}
  \langle T,\partial^{\alpha+\beta}\varphi\rangle
 =\langle\partial^{\alpha+\beta}T,\varphi\rangle.
\end{aligned}
\]
W drugiej równości korzystamy z przemienności pochodnych
funkcji gładkich i z $|\alpha+\beta|=|\alpha|+|\beta|$.
W szczególności można zamienić kolejność dwóch pochodnych
pierwszego rzędu. Operator $\partial^\alpha$ jest liniowy
z definicji.
Dla ustalonej $\varphi$ zachodzi
$q_\varphi(\partial^\alpha T)=q_{\partial^\alpha\varphi}(T)$,
więc otrzymujemy ciągłość słabą-*.
Dla zbioru ograniczonego $B\subseteq\mathcal D(\Omega)$
mamy natomiast
\[
 q_B(\partial^\alpha T)=q_{D_\alpha(B)}(T).
\]
Z lematu~\ref{lem:test-derivative20} zbiór $D_\alpha(B)$ jest
ograniczony, co daje ciągłość silną. Gdy $T$ zeruje się na
otwartym $U\subseteq\Omega$, zeruje się tam także $\partial^\alpha T$.
Istotnie, jeśli $e_U$ oznacza przedłużenie zerem z definicji
\ref{def:distribution-zero19}, to dla $\varphi\in\mathcal D(U)$
mamy $\partial^\alpha(e_U\varphi)=e_U(\partial^\alpha\varphi)$.
Równość zachodzi na $U$ i w otoczeniu każdego punktu
spoza $U$, gdzie obie strony znikają dzięki zwartemu nośnikowi
$\varphi$ w $U$. Stąd
$\langle(\partial^\alpha T)|_U,\varphi\rangle
=(-1)^{|\alpha|}\langle T|_U,\partial^\alpha\varphi\rangle=0$.
Stosując to do $U=\Omega\setminus\supp T$, otrzymujemy
inkluzję nośników.
Wreszcie dla $\varphi\in\mathcal D_K$ mamy
$|\langle\partial^\alpha T,\varphi\rangle|
\le C_Kp_m(\partial^\alpha\varphi)
\le C_Kp_{m+|\alpha|}(\varphi)$.

Aby uzasadnić ostatni punkt, ustalmy
$x\notin\operatorname{sing\,supp}T$. Z
Definicji~\ref{def:dt-singular-support} istnieje
otwarte otoczenie $U$ punktu $x$ oraz
$f\in C^\infty(U;\K)$ takie, że $T|_U=T_f$.
Udowodniona zgodność przedłużania zerem z różniczkowaniem
daje $(\partial^\alpha T)|_U=\partial^\alpha(T|_U)$.
Dla gładkiej $f$ całkowanie przez części
(szczegółowo przeprowadzone w następnym stwierdzeniu)
daje $\partial^\alpha(T|_U)=T_{\partial^\alpha f}$.
Ponieważ $\partial^\alpha f$ jest gładka na $U$,
punkt $x$ nie należy do nośnika osobliwego
$\partial^\alpha T$.
\end{proof}

\begin{stwierdzenie}[Zgodność z pochodną klasyczną]
\label{prop:classical-weak20}
Niech $m\in\N$, $f\in C^m(\Omega;\K)$, a $\alpha\in\N_0^d$
niech spełnia $|\alpha|\le m$. Wtedy
$\partial^\alpha T_f=T_{\partial^\alpha f}$.
Dla $f,g\in L^1_{\mathrm{loc}}(\Omega;\K)$
oraz $j\in\{1,\ldots,d\}$ tożsamość
$\partial_jT_f=T_g$ jest równoważna temu, że $g$ jest
\emph{słabą pochodną} $f$ w kierunku $x_j$, tj.
\[
 \int_\Omega f(x)\,\partial_j\varphi(x)\,\mathrm d\lambda_d(x)
 =-\int_\Omega g(x)\varphi(x)\,\mathrm d\lambda_d(x)
 \quad(\varphi\in\mathcal D(\Omega)).
\]
\end{stwierdzenie}
\begin{proof}
Ustalmy $\varphi\in\mathcal D(\Omega)$ i wybierzmy
$\chi\in C_c^\infty(\Omega;\R)$ równą $1$ w otoczeniu
$\supp\varphi$ (Lemat~\ref{lem:cutoffs17}).
Przedłużenie $F$ funkcji $\chi f$ zerem na $\R^d$
jest klasy $C^m$ i ma zwarty nośnik. Niech $\Phi$ będzie
przedłużeniem $\varphi$ zerem na $\R^d$.
Na otoczeniu $\supp\varphi$ funkcje $F$ i $f$ oraz ich pochodne
do rzędu $m$ są zgodne. Wszystkie pochodne $\Phi$
znikają poza $\supp\varphi$.
Dla $j\in\{1,\ldots,d\}$ i funkcji $F$ klasy co najmniej $C^1$
twierdzenie Fubiniego i całkowanie przez części względem $x_j$
dają
\[
\int_{\R^d}F\,\partial_j\Phi\,\mathrm d\lambda_d
=-\int_{\R^d}(\partial_jF)\Phi\,\mathrm d\lambda_d.
\]
Możemy całkować po dostatecznie dużym prostokącie zawierającym
oba nośniki we wnętrzu, więc wyrazy brzegowe znikają.
Całki są bezwzględnie zbieżne dzięki ciągłości i zwartym nośnikom.
Powtarzając argument dla kolejnych pochodnych zgodnie z
wielowskaźnikiem $\alpha$, dostajemy
\[
\begin{aligned}
\langle\partial^\alpha T_f,\varphi\rangle
&=(-1)^{|\alpha|}\int_{\R^d}F\partial^\alpha\Phi\,\mathrm d\lambda_d\\
&=\int_{\R^d}(\partial^\alpha F)\Phi\,\mathrm d\lambda_d
=\int_\Omega(\partial^\alpha f)\varphi\,\mathrm d\lambda_d
=\langle T_{\partial^\alpha f},\varphi\rangle.
\end{aligned}
\]
Dla $\alpha=0$ jest to tożsamość bez całkowania przez części.
Nie potrzeba żadnych warunków brzegowych na $f$.

W drugiej części z definicji
$\langle\partial_jT_f,\varphi\rangle
=-\int_\Omega f\partial_j\varphi\,\mathrm d\lambda_d$,
natomiast $\langle T_g,\varphi\rangle
=\int_\Omega g\varphi\,\mathrm d\lambda_d$.
Równość dystrybucji zachodzi dokładnie wtedy, gdy te wartości
są równe dla wszystkich funkcji testowych.
Jeżeli słaba pochodna istnieje, jest jednoznaczna jako element
$L^1_{\mathrm{loc}}(\Omega;\K)$ na mocy
Stwierdzenia~\ref{prop:regular-injective19}.
\end{proof}

\subsection{Skoki, funkcja Heaviside'a i osobliwości punktowe}

\begin{przyklad}[Heaviside i delta Diraca]
\label{ex:heaviside20}
Na $\R$ niech $H(x)=0$ dla $x<0$ oraz $H(x)=1$ dla $x>0$;
przyjmijmy $H(0)=0$. Wartość w $0$ nie ma znaczenia
dla wyznaczonej dystrybucji, ponieważ $\lambda_1(\{0\})=0$.
Dla dowolnej $\varphi\in\mathcal D(\R)$ otrzymujemy
\[
 \langle (T_H)',\varphi\rangle
 =-\int_0^\infty\varphi'(x)\,\mathrm d\lambda_1(x)
 =\varphi(0)=\langle\delta_0,\varphi\rangle.
\]
Całka ze znakiem minus jest równa $\varphi(0)$, gdyż $\varphi$ znika
dla dostatecznie dużych $x$.
Zatem $(T_H)'=\delta_0$; w dalszym tekście piszemy krótko
$H'=\delta_0$, utożsamiając funkcje lokalnie całkowalne
z odpowiadającymi im dystrybucjami.
Ponieważ $\operatorname{sgn}x=2H(x)-1$ prawie wszędzie,
a pochodna stałej dystrybucji jest zerowa
(Stwierdzenie~\ref{prop:classical-weak20}), mamy
$(\operatorname{sgn}x)'=2\delta_0$.
Ponadto
\[
\begin{aligned}
\langle(T_{|\cdot|})',\varphi\rangle
&=\int_{-\infty}^0 x\varphi'(x)\,\mathrm d\lambda_1(x)
 -\int_0^\infty x\varphi'(x)\,\mathrm d\lambda_1(x)\\
&=-\int_{-\infty}^0\varphi(x)\,\mathrm d\lambda_1(x)
 +\int_0^\infty\varphi(x)\,\mathrm d\lambda_1(x)
 =\langle T_{\operatorname{sgn}},\varphi\rangle .
\end{aligned}
\]
Wyrazy brzegowe w zerze znikają przez czynnik $x$,
a w nieskończoności przez zwarty nośnik $\varphi$.
Stąd $|x|'=\operatorname{sgn}x$ i $|x|''=2\delta_0$
w sensie dystrybucji.
\end{przyklad}

\begin{center}
\begin{tikzpicture}[scale=0.82,>=Stealth]
\draw[->] (-3,0) -- (3.3,0) node[right] {$x$};
\draw[->] (0,-0.5) -- (0,1.9) node[above] {$H(x)$};
\draw[very thick,blue!60!black] (-2.9,0)--(0,0);
\draw[very thick,blue!60!black] (0,1.3)--(3.0,1.3);
\draw[fill=blue!60!black,draw=blue!60!black,thick] (0,0) circle (2.3pt);
\draw[fill=white,draw=blue!60!black,thick] (0,1.3) circle (2.3pt);
\node[left] at (-0.12,1.3) {$1$};
\node[below left] at (-0.08,-0.05) {$0$};
\draw[gray,dashed] (0,0.12)--(0,1.2);
\end{tikzpicture}

\small\emph{Skok funkcji $H$ z przyjętą wartością $H(0)=0$.
Zmiana wartości w zerze nie wpływa na tożsamość dystrybucyjną
$H'=\delta_0$.}
\end{center}

\begin{twierdzenie}[Pochodna funkcji ze skokiem]
\label{thm:jump20}
Niech $a\in\R$, a $f\in L^1_{\mathrm{loc}}(\R;\K)$ będzie
klasy $C^1$ na $(-\infty,a)$ i $(a,\infty)$. Załóżmy ponadto,
że funkcja ma skończone granice jednostronne
$f(a-),f(a+)$, a klasyczna pochodna określona poza $a$
należy do $L^1_{\mathrm{loc}}(\R;\K)$.
Niech $g(x)=f'(x)$ dla $x\ne a$ (wartość w $a$ jest dowolna).
Wówczas
\[
 (T_f)'=T_g+\bigl(f(a+)-f(a-)\bigr)\delta_a.
\]
\end{twierdzenie}
\begin{proof}
Ustalmy $\varphi\in\mathcal D(\R)$ i $R>0$ takie, że
$\supp\varphi\cup\{a\}\subset(-R,R)$.
Dla $0<\varepsilon<R-|a|$ całkujemy przez części na
$[-R,a-\varepsilon]$ i $[a+\varepsilon,R]$:
\[
\begin{aligned}
-\int_{-R}^{a-\varepsilon}f\varphi'\,\mathrm d\lambda_1
&=-f(a-\varepsilon)\varphi(a-\varepsilon)
 +\int_{-R}^{a-\varepsilon}g\varphi\,\mathrm d\lambda_1,\\
-\int_{a+\varepsilon}^Rf\varphi'\,\mathrm d\lambda_1
&= f(a+\varepsilon)\varphi(a+\varepsilon)
 +\int_{a+\varepsilon}^Rg\varphi\,\mathrm d\lambda_1.
\end{aligned}
\]
Wyrazy brzegowe w $\pm R$ znikają.
Skończone granice jednostronne $f$ zapewniają ograniczoność
$|f|$ przez pewne $M$ w przebitym otoczeniu $a$.
Pominięty wkład do całki z $f\varphi'$ ma moduł
nie większy niż $2\varepsilon M p_1(\varphi)$,
więc dąży do zera.
Ponieważ $g\varphi$ jest całkowalna na $[-R,R]$,
z bezwzględnej ciągłości całki wkład z
$(a-\varepsilon,a+\varepsilon)$ również dąży do zera.
Ciągłość $\varphi$ i granice jednostronne $f$ dają
granicę wyrazów brzegowych
$(f(a+)-f(a-))\varphi(a)$.
Sumując i przechodząc do granicy, otrzymujemy
\[
 -\int_\R f(x)\varphi'(x)\,\mathrm d\lambda_1(x)
 =\int_\R g(x)\varphi(x)\,\mathrm d\lambda_1(x)
 +(f(a+)-f(a-))\varphi(a).
\]
Jest to żądana tożsamość dystrybucyjna.
\end{proof}

\begin{przyklad}[Funkcja częściowo gładka]
\label{ex:piecewise20}
Niech $f(x)=e^x$ dla $x<0$, $f(x)=2e^{-x}$ dla $x>0$
oraz $f(0)=0$. Funkcja jest lokalnie całkowalna,
jest klasy $C^1$ poza zerem, ma granice $f(0-)=1$ i $f(0+)=2$,
a jej pochodna poza zerem jest lokalnie całkowalna.
Z Twierdzenia~\ref{thm:jump20} otrzymujemy, ponieważ
$f(0+)-f(0-)=1$,
\[
 (T_f)'=T_{e^x\chi_{(-\infty,0)}
 -2e^{-x}\chi_{(0,\infty)}}+\delta_0.
\]
Jest to równość dystrybucji na $\R$; jej prawa strona
zawiera składnik osobliwy, którego nie wykrywa klasyczne
różniczkowanie poza punktem skoku.
\end{przyklad}

\begin{przyklad}[Pochodne delty]
\label{ex:delta-derivative20}
Dla $a\in\Omega$ i $\alpha\in\N_0^d$ zachodzi
$\langle\partial^\alpha\delta_a,\varphi\rangle
=(-1)^{|\alpha|}\partial^\alpha\varphi(a)$.
Dla $\Omega=\R$ funkcjonał $\delta_0'$ jest więc różny
od zwykłej miary Diraca: dla funkcji testowej
$\varphi(x)=x\eta(x)$, gdzie $\eta\in C_c^\infty(\R;\R)$
jest równa $1$ w otoczeniu zera, mamy
$\langle\delta_0,\varphi\rangle=0$ oraz
$\langle\delta_0',\varphi\rangle=-1$.
Ze Stwierdzenia~\ref{prop:delta-order19} wynika,
że $\delta_0'$ ma rząd dokładnie $1$ na każdym zwartym
przedziale zawierającym otoczenie zera.
\end{przyklad}

\subsection{Mnożenie przez funkcję gładką i reguła Leibniza}

\begin{definicja}[Mnożenie dystrybucji]
\label{def:multiply20}
Dla $g\in C^\infty(\Omega;\K)$ oraz $T\in\mathcal D'(\Omega;\K)$
definiujemy funkcjonał $gT$ na $\mathcal D(\Omega;\K)$ wzorem
$\langle gT,\varphi\rangle=\langle T,g\varphi\rangle$
dla każdej $\varphi\in\mathcal D(\Omega;\K)$.
Iloczyn $g\varphi$ jest gładki i ma nośnik zawarty w
$\supp\varphi$, więc należy do $\mathcal D(\Omega;\K)$.
Ciągłość funkcjonału $gT$ wykażemy w następnym twierdzeniu.
\end{definicja}

\begin{twierdzenie}[Własności mnożenia]
\label{thm:multiply20}
Niech $g,h\in C^\infty(\Omega;\K)$, $T\in\mathcal D'(\Omega;\K)$
i $\alpha\in\N_0^d$. Wtedy:
\begin{enumerate}
\item $gT$ jest dobrze określoną dystrybucją,
$g(hT)=(gh)T$ oraz $1T=T$;
\item $\supp(gT)\subseteq\supp g\cap\supp T$;
\item operator $M_g:T\mapsto gT$ jest ciągły w topologiach
słabej-* i silnej; jeśli $T$ ma rząd nie większy niż $m$
na zwartym $K\Subset\Omega$, to $gT$ ma na $K$ rząd
nie większy niż $m$;
\item zachodzi \emph{wzór Leibniza}
\[
 \partial^\alpha(gT)=
 \sum_{\beta\le\alpha}\binom\alpha\beta
 (\partial^\beta g)\,\partial^{\alpha-\beta}T.
\]
\end{enumerate}
\end{twierdzenie}
\begin{proof}
Oznaczmy przez $N_g$ operator mnożenia przez $g$ na
przestrzeni funkcji testowych, aby odróżnić go od operatora
$M_g$ działającego na dystrybucjach.
Dla każdego niepustego zwartego $K\Subset\Omega$
odwzorowanie $N_g$ przekształca $\mathcal D_K$ w siebie.
Ustalmy $m\in\N_0$. Jeśli $\varphi\in\mathcal D_K$
i $\gamma\in\N_0^d$, $|\gamma|\le m$, to
klasyczna reguła Leibniza daje
$\partial^\gamma(g\varphi)=\sum_{\beta\le\gamma}
\binom\gamma\beta(\partial^\beta g)
(\partial^{\gamma-\beta}\varphi)$.
Wszystkie składniki znikają poza $K$, a na $K$ moduł
$\partial^{\gamma-\beta}\varphi$ nie przekracza $p_m(\varphi)$.
Biorąc supremum i maksimum, otrzymujemy oszacowanie
\[
 p_m(g\varphi)\le C_{K,m,g}\,p_m(\varphi),
 \quad C_{K,m,g}=\max_{|\gamma|\le m}
 \sum_{\beta\le\gamma}\binom\gamma\beta
 \sup_{x\in K}|\partial^\beta g(x)|.
\]
Stała jest skończona ze zwartości $K$ i ciągłości pochodnych $g$.
Dla pustego $K$ etap składa się wyłącznie z funkcji zerowej.
Po złożeniu z włożeniem etapu w $\mathcal D(\Omega)$
otrzymujemy ciągłość $N_g$ na każdym etapie, a z
Twierdzenia~\ref{thm:LF-universal18} ciągłość na $\mathcal D(\Omega)$.
Wobec tego $gT=T\circ N_g$ jest ciągłym funkcjonałem liniowym.
Dla każdej funkcji testowej $\varphi$ mamy
$\langle g(hT),\varphi\rangle=\langle hT,g\varphi\rangle
=\langle T,hg\varphi\rangle=\langle(gh)T,\varphi\rangle$
oraz $\langle1T,\varphi\rangle=\langle T,\varphi\rangle$.
Liniowość $M_g$ wynika z liniowości parowania względem $T$.

Jeśli $g$ zeruje się na otwartym $U\subseteq\Omega$, to
$g(e_U\varphi)=0$ dla $\varphi\in\mathcal D(U)$.
Jeśli natomiast $T|_U=0$, to
$g(e_U\varphi)=e_U((g|_U)\varphi)$, gdzie
$(g|_U)\varphi\in\mathcal D(U)$, zatem
$\langle gT,e_U\varphi\rangle=0$.
Stosując te obserwacje odpowiednio do
$U=\Omega\setminus\supp g$ i $U=\Omega\setminus\supp T$,
otrzymujemy $\supp(gT)\subseteq\supp g\cap\supp T$.
Ciągłość słaba-* wynika z tożsamości
$q_\varphi(gT)=q_{g\varphi}(T)$, a silna z
$q_B(gT)=q_{N_g(B)}(T)$, gdzie $B$ jest dowolnym zbiorem
ograniczonym w $\mathcal D(\Omega)$.
Zbiór $N_g(B)$ jest ograniczony, ponieważ $N_g$ jest
ciągły i liniowy (argument jak w Lemacie~\ref{lem:test-derivative20}).
Jeśli $|\langle T,\psi\rangle|\le C_Kp_m(\psi)$
dla $\psi\in\mathcal D_K$, to $g\varphi\in\mathcal D_K$
dla $\varphi\in\mathcal D_K$, a wcześniejsze oszacowanie
mnożenia daje
\[
 |\langle gT,\varphi\rangle|
 \le C_Kp_m(g\varphi)
 \le C_K C_{K,m,g}p_m(\varphi).
\]
To jest oszacowanie rzędu nie większego niż $m$ na $K$.
Dla $j\in\{1,\ldots,d\}$ i $\alpha=e_j$ obliczamy
\[
 \begin{aligned}
 \langle\partial_j(gT),\varphi\rangle
 &=-\langle T,g\partial_j\varphi\rangle\\
 &=\langle\partial_jT,g\varphi\rangle
   +\langle T,(\partial_jg)\varphi\rangle,
 \end{aligned}
\]
czyli $\partial_j(gT)=g\partial_jT+(\partial_jg)T$.
Dla $\alpha=0$ wzór jest tożsamością $gT=gT$.
Załóżmy, że wzór zachodzi dla wszystkich wielowskaźników
rzędu $n$, i ustalmy $\gamma$ rzędu $n+1$.
Wybierzmy $j$ z $\gamma_j\ge1$ i połóżmy $\alpha=\gamma-e_j$.
Korzystając z przemienności pochodnych
(Twierdzenie~\ref{thm:dist-deriv-cont20}) oraz udowodnionej
reguły pierwszego rzędu dla każdego składnika, otrzymujemy
\[
\begin{aligned}
\partial^\gamma(gT)
={}&\sum_{\beta\le\alpha}\binom\alpha\beta
(\partial^{\beta+e_j}g)\partial^{\alpha-\beta}T\\
&+\sum_{\beta\le\alpha}\binom\alpha\beta
(\partial^\beta g)\partial^{\alpha-\beta+e_j}T.
\end{aligned}
\]
W pierwszej sumie podstawiamy $\nu=\beta+e_j$, a w drugiej
$\nu=\beta$. Przyjmując $\binom\alpha\mu=0$, gdy którakolwiek
współrzędna $\mu$ jest ujemna lub większa od odpowiedniej
współrzędnej $\alpha$, otrzymujemy
\[
\partial^\gamma(gT)=\sum_{\nu\le\gamma}
\left(\binom\alpha{\nu-e_j}+\binom\alpha\nu\right)
(\partial^\nu g)\partial^{\gamma-\nu}T.
\]
Ze zwykłej tożsamości Pascala dla $j$-tej współrzędnej,
po pomnożeniu przez $\prod_{k\ne j}\binom{\alpha_k}{\nu_k}$,
wynika $\binom\alpha{\nu-e_j}+\binom\alpha\nu=\binom\gamma\nu$.
Obejmuje to także skrajne wartości $\nu_j=0$ i $\nu_j=\gamma_j$.
Dowód indukcyjny jest zakończony.
\end{proof}

\begin{przyklad}[Mnożenie dystrybucji punktowych]
\label{ex:multiply-delta20}
Dla $a\in\R$ i $g\in C^\infty(\R;\K)$ mamy
\[
 g\delta_a=g(a)\delta_a,\qquad
 g\delta_a'=g(a)\delta_a'-g'(a)\delta_a.
\]
W szczególności $x\delta_0=0$ oraz $x\delta_0'=-\delta_0$.
Istotnie, dla $\varphi\in\mathcal D(\R)$ zachodzi
$\langle g\delta_a,\varphi\rangle=g(a)\varphi(a)$ oraz
\[
\begin{aligned}
\langle g\delta_a',\varphi\rangle
&=-(g\varphi)'(a)=-g'(a)\varphi(a)-g(a)\varphi'(a)\\
&=\langle g(a)\delta_a'-g'(a)\delta_a,\varphi\rangle.
\end{aligned}
\]
Podstawiając $g(x)=x$ i $a=0$, otrzymujemy obie
podane tożsamości szczególne.

Ogólniej, dla $a\in\Omega$, $\alpha\in\N_0^d$
i $g\in C^\infty(\Omega;\K)$ mamy
\[
 g\partial^\alpha\delta_a
 =\sum_{\beta\le\alpha}(-1)^{|\beta|}
 \binom{\alpha}{\beta}(\partial^\beta g)(a)
 \partial^{\alpha-\beta}\delta_a.
\]
Rzeczywiście, z definicji mnożenia, definicji pochodnej
i klasycznej reguły Leibniza dla dowolnego testu
$\varphi\in\mathcal D(\Omega)$ wynika
\[
\begin{aligned}
 \langle g\partial^\alpha\delta_a,\varphi\rangle
 &=(-1)^{|\alpha|}\partial^\alpha(g\varphi)(a)\\
 &=\sum_{\beta\le\alpha}(-1)^{|\alpha|}
  \binom{\alpha}{\beta}
  (\partial^\beta g)(a)(\partial^{\alpha-\beta}\varphi)(a)\\
 &=\sum_{\beta\le\alpha}(-1)^{|\beta|}
  \binom{\alpha}{\beta}(\partial^\beta g)(a)
  \langle\partial^{\alpha-\beta}\delta_a,\varphi\rangle.
\end{aligned}
\]
W ostatnim przejściu wykorzystaliśmy
$|\alpha|=|\beta|+|\alpha-\beta|$.
W połączeniu z Twierdzeniem~\ref{thm:dt-point-support-structure}
wzór opisuje mnożenie przez funkcję gładką każdej
dystrybucji skupionej w punkcie. W szczególności
mnożnik, którego wszystkie pochodne do rzędu $m$
znikają w $a$, anihiluje każdą dystrybucję
o nośniku zawartym w $\{a\}$ i rzędzie nie większym niż $m$.
\end{przyklad}

\begin{uwaga}[Dlaczego nie mnożymy dowolnych dystrybucji?]
\label{rem:no-product20}
Powyższa definicja wymaga \emph{gładkiego} mnożnika.
Nie ustanawia iloczynu dwóch dowolnych dystrybucji.
Przykładowo wyrażenie $H\delta_0$ nie ma kanonicznego
znaczenia na podstawie samej struktury $\mathcal D'(\R)$:
po zastąpieniu $\delta_0$ przez
$\rho_\varepsilon(x)=\varepsilon^{-1}\rho(x/\varepsilon)$,
gdzie $\rho\in C_c^\infty(\R;\R)$, $\rho\ge0$ oraz
$\int_\R\rho(x)\,\mathrm d\lambda_1(x)=1$,
otrzymujemy w sensie dystrybucji
\[
 H\rho_\varepsilon\longrightarrow c_\rho\delta_0,
 \qquad c_\rho=\int_0^\infty\rho(t)\,\mathrm d\lambda_1(t).
\]
Dla $\varepsilon>0$ iloczyn funkcji $H\rho_\varepsilon$
jest lokalnie całkowalny. Podstawienie $x=\varepsilon t$
daje dla każdej $\varphi\in\mathcal D(\R)$
\[
\langle T_{H\rho_\varepsilon},\varphi\rangle
=\int_0^\infty\rho(t)\varphi(\varepsilon t)\,\mathrm d\lambda_1(t).
\]
Całka dąży do $c_\rho\varphi(0)$ na mocy twierdzenia
o zbieżności zmajoryzowanej z majorantą
$\rho(t)p_0(\varphi)$.
Nawet silna zbieżność zachodzi: z
$|\varphi(\varepsilon t)-\varphi(0)|\le\varepsilon t p_1(\varphi)$
dla $t\ge0$ otrzymujemy
$q_B(T_{H\rho_\varepsilon}-c_\rho\delta_0)
\le\varepsilon\sup_{\varphi\in B}p_1(\varphi)
\int_0^\infty t\rho(t)\,\mathrm d\lambda_1(t)$
dla każdego niepustego ograniczonego $B\subseteq\mathcal D(\R)$.
Wszystkie takie $\rho_\varepsilon$ zbiegają do $\delta_0$
silnie, co wynika z analogicznego oszacowania całki po $\R$
z czynnikiem $|t|$. Wybierając znormalizowaną nieujemną
funkcję testową o nośniku w $(-1,0)$, otrzymujemy $c_\rho=0$,
a o nośniku w $(0,1)$ otrzymujemy $c_\rho=1$.
Istnienie obu funkcji wynika z Lematu~\ref{lem:cutoffs17}
i normalizacji przez dodatnią całkę.
Zatem takie regularizacje nie wyznaczają jednoznacznego
iloczynu $H\delta_0$; nie jest to zakaz konstruowania
iloczynów w dodatkowo określonych strukturach.
\end{uwaga}

\subsection{Operatory różniczkowe i ich transpozycja}

\begin{definicja}[Formalny operator transponowany]
\label{def:formal-transpose20}
Niech $m\in\N_0$ i $a_\alpha\in C^\infty(\Omega;\K)$
dla $|\alpha|\le m$. Dla funkcji gładkich i dystrybucji
rozważamy operator $P=\sum_{|\alpha|\le m}a_\alpha\partial^\alpha$.
Jego \emph{formalny operator transponowany} na $\mathcal D(\Omega)$
definiujemy wzorem
\[
 P^{\mathrm t}\varphi=
 \sum_{|\alpha|\le m}(-1)^{|\alpha|}
 \partial^\alpha(a_\alpha\varphi).
\]
Przy naszej konwencji parowania nie występuje sprzężenie
współczynników zespolonych. Nie należy zatem utożsamiać
formalnej transpozycji z formalnym sprzężeniem względem
zespolonego iloczynu skalarnego w $L^2$, ani z operatorem
sprzężonym określonym dla konkretnej dziedziny operatora.
W szczególnych przypadkach wzory mogą się pokrywać.
\end{definicja}

\begin{twierdzenie}[Transpozycja jako konstrukcja operatora na dystrybucjach]
\label{thm:transpose20}
Dla operatora $P$ z definicji~\ref{def:formal-transpose20}
oraz każdej dystrybucji $T\in\mathcal D'(\Omega;\K)$
zachodzą:
\begin{enumerate}
\item $P^{\mathrm t}:\mathcal D(\Omega)\to\mathcal D(\Omega)$
jest ciągły i liniowy;
\item wzór $\langle PT,\varphi\rangle=
\langle T,P^{\mathrm t}\varphi\rangle$ dla każdej
$\varphi\in\mathcal D(\Omega)$ wyznacza dystrybucję
zgodną z $\sum_{|\alpha|\le m}a_\alpha\partial^\alpha T$;
\item $P$ jest ciągły w topologiach słabej-* i silnej;
\item $\supp(PT)\subseteq\supp T$;
\item jeżeli $f\in C^m(\Omega)$, to $PT_f=T_{Pf}$;
\item $\operatorname{sing\,supp}(PT)\subseteq
\operatorname{sing\,supp}T$, a jeżeli $T$ ma rząd
nie większy niż $r$ na zwartym $K\Subset\Omega$,
to $PT$ ma na $K$ rząd nie większy niż $r+m$.
\end{enumerate}
\end{twierdzenie}
\begin{proof}
Każdy składnik $P^{\mathrm t}$ jest równy
$(-1)^{|\alpha|}D_\alpha N_{a_\alpha}$, gdzie
$N_{a_\alpha}$ oznacza mnożenie funkcji testowych przez
$a_\alpha$. Obie operacje są ciągłe i liniowe na mocy
Lematu~\ref{lem:test-derivative20} i dowodu
Twierdzenia~\ref{thm:multiply20}. Suma jest skończona,
więc $P^{\mathrm t}$ jest ciągły i liniowy.
Ponieważ każda z operacji nie powiększa nośnika,
$\supp(P^{\mathrm t}\varphi)\subseteq\supp\varphi$.
Złożenie z $T$ jest ciągłym funkcjonałem liniowym.
Ponadto dla każdej funkcji testowej zachodzi
\[
\begin{aligned}
\left\langle\sum_{|\alpha|\le m}a_\alpha\partial^\alpha T,
\varphi\right\rangle
&=\sum_{|\alpha|\le m}\langle\partial^\alpha T,
a_\alpha\varphi\rangle\\
&=\sum_{|\alpha|\le m}(-1)^{|\alpha|}
\langle T,\partial^\alpha(a_\alpha\varphi)\rangle
=\langle T,P^{\mathrm t}\varphi\rangle.
\end{aligned}
\]
To dowodzi zgodności obu konstrukcji $PT$.
Zależności
$q_\varphi(PT)=q_{P^{\mathrm t}\varphi}(T)$ oraz
$q_B(PT)=q_{P^{\mathrm t}(B)}(T)$
dają obie ciągłości, ponieważ $P^{\mathrm t}(B)$ jest ograniczony.
Niech $U=\Omega\setminus\supp T$. Dla
$\varphi\in\mathcal D(U)$ funkcja $P^{\mathrm t}(e_U\varphi)$
ma nośnik zawarty w $\supp(e_U\varphi)\Subset U$,
jest więc przedłużeniem zerem funkcji testowej na $U$.
Wobec $T|_U=0$ mamy $\langle PT,e_U\varphi\rangle=0$,
a zatem $\supp(PT)\subseteq\supp T$.
Dla $f\in C^m(\Omega;\K)$, korzystając ze
Stwierdzenia~\ref{prop:classical-weak20} (dla $m=0$
pochodna zerowego rzędu jest identycznością), dostajemy
$\partial^\alpha T_f=T_{\partial^\alpha f}$.
Ponadto $a_\alpha T_{\partial^\alpha f}
=T_{a_\alpha\partial^\alpha f}$, gdyż obie strony mają
wartość $\int_\Omega a_\alpha(\partial^\alpha f)\varphi\,
\mathrm d\lambda_d$ na dowolnej funkcji testowej $\varphi$.
Sumowanie daje $PT_f=T_{Pf}$.

Jeśli $x\notin\operatorname{sing\,supp}T$, to istnieje
otwarte otoczenie $U$ punktu $x$, na którym $T|_U=T_f$
dla $f\in C^\infty(U;\K)$.
Przedłużanie testów zerem jest zgodne zarówno
z różniczkowaniem, jak i mnożeniem przez współczynniki
ograniczone do $U$, jak wykazano w dowodach
Twierdzeń~\ref{thm:dist-deriv-cont20}
i~\ref{thm:multiply20}.
Zatem $(PT)|_U=P_U(T|_U)$, gdzie $P_U$ oznacza
operator o współczynnikach $a_\alpha|_U$.
Z poprzedniego punktu $(PT)|_U=T_{P_U f}$,
a $P_U f$ jest gładka na $U$.
Stąd żądana inkluzja nośników osobliwych.
Wreszcie, gdy $T$ ma rząd nie większy niż $r$ na $K$,
Twierdzenie~\ref{thm:dist-deriv-cont20} daje rząd
nie większy niż $r+|\alpha|$ dla $\partial^\alpha T$,
a Twierdzenie~\ref{thm:multiply20} zachowuje to
ograniczenie po pomnożeniu przez $a_\alpha$.
Każdy składnik ma więc na $K$ rząd nie większy niż
$r+m$; sumując skończenie wiele odpowiednich stałych
otrzymujemy oszacowanie tego samego rzędu dla $PT$.
\end{proof}

\begin{przyklad}[Operator pierwszego rzędu]
\label{ex:first-order-P20}
W jednym wymiarze niech $a,b\in C^\infty(\R;\K)$
i $P=a(x)\frac{\mathrm d}{\mathrm dx}+b(x)$.
Z Definicji~\ref{def:formal-transpose20} otrzymujemy
$P^{\mathrm t}\varphi=-(a\varphi)'+b\varphi
=-a\varphi'+(b-a')\varphi$.
W $\R^d$ dla $b=(b_1,\ldots,b_d)\in C^\infty(\R^d;\K^d)$
i $c\in C^\infty(\R^d;\K)$ operator
$P=b\cdot\nabla+c=\sum_{j=1}^d b_j\partial_j+c$
ma formalną transpozycję
\[
P^{\mathrm t}\varphi=-\sum_{j=1}^d\partial_j(b_j\varphi)
+c\varphi=-b\cdot\nabla\varphi+(c-\operatorname{div}b)\varphi.
\]
Tutaj $\operatorname{div}b=\sum_{j=1}^d\partial_j b_j$,
a iloczyn $b\cdot\nabla$ nie zawiera sprzężenia zespolonego.
Druga równość wynika z klasycznej reguły Leibniza.
\end{przyklad}

\begin{wniosek}[Komutator pochodnej i mnożenia]
\label{cor:commutator20}
Dla $g\in C^\infty(\Omega)$ i $j\in\{1,\ldots,d\}$ mamy
$[\partial_j,M_g]=M_{\partial_jg}$ na $\mathcal D'(\Omega)$,
gdzie $[A,B]=AB-BA$ jest komutatorem operatorów.
\end{wniosek}
\begin{proof}
Jest to wzór Leibniza pierwszego rzędu:
$\partial_j(gT)-g\partial_jT=(\partial_jg)T$.
\end{proof}

\subsection{Równania różniczkowe w sensie dystrybucji}

\begin{twierdzenie}[Jądro pochodnej na przedziale]
\label{thm:zero-derivative20}
Niech $I\subseteq\R$ będzie niepustym przedziałem otwartym,
a $T\in\mathcal D'(I;\K)$. Wówczas $T'=0$
wtedy i tylko wtedy, gdy istnieje $c\in\K$ takie, że
$T=T_c$, czyli
$\langle T,\varphi\rangle=c\int_I\varphi(x)\,\mathrm d\lambda_1(x)$
dla każdej $\varphi\in\mathcal D(I)$.
\end{twierdzenie}
\begin{proof}
Wybierzmy niezerową nieujemną funkcję
$\eta_0\in C_c^\infty(I;\R)$ na mocy
Lematu~\ref{lem:cutoffs17}. Jej całka jest skończona
dzięki zwartemu nośnikowi i dodatnia, ponieważ $\eta_0$
jest dodatnia na pewnym niepustym otwartym podprzedziale.
Zatem
funkcja $\eta=\eta_0/\int_I\eta_0\,\mathrm d\lambda_1$
należy do $\mathcal D(I;\K)$ i ma całkę $1$.
Dla $\varphi\in\mathcal D(I)$ oznaczmy
$m_\varphi=\int_I\varphi(x)\,\mathrm d\lambda_1(x)$
i $h=\varphi-m_\varphi\eta$. Wtedy
$\int_I h\,\mathrm d\lambda_1=m_\varphi-m_\varphi=0$.
Przedłużamy $h$ zerem na $\R$ i określamy
$\psi(x)=\int_{-\infty}^xh(t)\,\mathrm d\lambda_1(t)$.
Przedłużenie jest gładkie, gdyż $\supp h\Subset I$.
Wybierzmy $u<v$ z $[u,v]\Subset I$ i
$\supp h\subseteq[u,v]$. Dla $x<u$ całka jest zerowa,
a dla $x>v$ jest równa całce po całej prostej, czyli również
zero. Z twierdzenia Newtona--Leibniza $\psi'=h$,
więc $\psi$ jest gładka i $\supp\psi\subseteq[u,v]$.
Jej ograniczenie do $I$ należy zatem do $\mathcal D(I)$. Jeśli $T'=0$, to
$0=\langle T',\psi\rangle=-\langle T,\psi'\rangle$,
więc $\langle T,\varphi\rangle=m_\varphi\langle T,\eta\rangle$.
Wybieramy $c=\langle T,\eta\rangle$; testowanie na $\eta$
pokazuje także jednoznaczność tej stałej.
Odwrotnie, dla $c\in\K$ i $\varphi\in\mathcal D(I)$
wybieramy przedział $[u,v]\Subset I$ zawierający
$\supp\varphi$ we wnętrzu. Wtedy
$\langle T_c',\varphi\rangle=-c\int_u^v\varphi'(x)\,
\mathrm d\lambda_1(x)=-c(\varphi(v)-\varphi(u))=0$.
\end{proof}

\begin{twierdzenie}[Istnienie pierwotnej dystrybucji]
\label{thm:primitive20}
Dla każdej dystrybucji $S\in\mathcal D'(\R;\K)$
istnieje $T\in\mathcal D'(\R;\K)$ takie, że $T'=S$.
Wszystkie takie pierwotne różnią się o stałą dystrybucję.
\end{twierdzenie}
\begin{proof}
Ustalmy $\eta\in\mathcal D(\R;\K)$ o całce $1$,
skonstruowaną jak w dowodzie
Twierdzenia~\ref{thm:zero-derivative20}, i
zdefiniujmy $m_\varphi=\int_\R\varphi(x)\,\mathrm d\lambda_1(x)$
oraz
\[
 R\varphi(x)=\int_{-\infty}^x
 \bigl(\varphi(t)-m_\varphi\eta(t)\bigr)\,\mathrm d\lambda_1(t).
\]
Całka funkcji $\varphi-m_\varphi\eta$ jest zerowa,
więc argument o nośniku z dowodu
Twierdzenia~\ref{thm:zero-derivative20} daje
$R\varphi\in\mathcal D(\R)$.
Liniowość $R$ wynika z liniowości całek.
Dla każdego zwartego $K\Subset\R$ wybieramy zwarty przedział
$L=[u,v]$ o dodatniej długości $\ell=v-u$, zawierający
$K\cup\supp\eta$. Wówczas $R(\mathcal D_K)\subseteq\mathcal D_L$.
Dla $\varphi\in\mathcal D_K$ mamy
$|m_\varphi|\le\lambda_1(K)p_0(\varphi)$.
Ponieważ $R\varphi$ znika poza $L$,
\[
p_0(R\varphi)\le\ell\bigl(p_0(\varphi)
+\lambda_1(K)p_0(\varphi)p_0(\eta)\bigr).
\]
Ustalmy $m\in\N_0$. Dla $1\le j\le m$ z twierdzenia Newtona--Leibniza
otrzymujemy $(R\varphi)^{(j)}=\varphi^{(j-1)}
-m_\varphi\eta^{(j-1)}$, zatem
$p_0((R\varphi)^{(j)})\le
(1+\lambda_1(K)p_m(\eta))p_m(\varphi)$.
Oba oszacowania dają dla każdego $m\in\N_0$
\[
p_m(R\varphi)\le
\max\{1,\ell\}(1+\lambda_1(K)p_m(\eta))p_m(\varphi).
\]
Ograniczenie $R|_{\mathcal D_K}:\mathcal D_K\to\mathcal D_L$
jest więc ciągłe, także gdy $m=0$ lub $K=\varnothing$.
Po złożeniu z ciągłym włożeniem $\mathcal D_L\to\mathcal D(\R)$
i zastosowaniu Twierdzenia~\ref{thm:LF-universal18}
otrzymujemy ciągłość $R:\mathcal D(\R)\to\mathcal D(\R)$.
Definiujemy
$\langle T,\varphi\rangle=-\langle S,R\varphi\rangle$.
Wtedy $T=-S\circ R$ jest ciągłym funkcjonałem liniowym,
czyli $T\in\mathcal D'(\R;\K)$.
Dla $\varphi\in\mathcal D(\R)$ zachodzi
$m_{\varphi'}=\int_\R\varphi'\,\mathrm d\lambda_1=0$,
a z twierdzenia Newtona--Leibniza
$R(\varphi')(x)=\int_{-\infty}^x\varphi'(t)\,
\mathrm d\lambda_1(t)=\varphi(x)$, gdyż $\varphi$ znika
dla dostatecznie małych argumentów.
Stąd
$\langle T',\varphi\rangle=-\langle T,\varphi'\rangle
=\langle S,R\varphi'\rangle=\langle S,\varphi\rangle$.
Jeżeli $T_1'=T_2'=S$, to $(T_1-T_2)'=0$, a z
Twierdzenia~\ref{thm:zero-derivative20} wynika $T_1-T_2=T_c$
dla pewnego $c\in\K$. Odwrotnie, dodanie $T_c$ do
dowolnej pierwotnej nie zmienia jej pochodnej.
\end{proof}

\begin{wniosek}[Równanie pierwszego rzędu]
\label{cor:ode20}
Niech $a\in C^\infty(\R;\K)$ i
$S\in\mathcal D'(\R;\K)$. Wszystkie dystrybucyjne
rozwiązania równania $T'+aT=S$ mają postać
$T=e^{-A}V$, gdzie $A(x)=\int_0^xa(t)\,\mathrm d\lambda_1(t)$,
a $V$ jest dowolną pierwotną dystrybucją $e^AS$.
Dwa rozwiązania różnią się o dystrybucję regularną
$T_{c e^{-A}}$ dla pewnego $c\in\K$.
\end{wniosek}
\begin{proof}
Z twierdzenia Newtona--Leibniza $A'=a$, a funkcje $e^A$
i $e^{-A}$ są gładkie, także dla zespolonego $a$.
Z Twierdzenia~\ref{thm:multiply20} operatory mnożenia
przez te funkcje są wzajemnie odwrotne.
Reguła Leibniza daje dla każdej dystrybucji $T$
$(e^AT)'=e^AT'+(e^A)'T=e^A(T'+aT)$.
Zatem $T'+aT=S$ wtedy i tylko wtedy, gdy
$(e^AT)'=e^AS$. Z Twierdzenia~\ref{thm:primitive20}
istnieje pierwotna $V$ dystrybucji $e^AS$ i wszystkie
pozostałe mają postać $V+T_c$, $c\in\K$.
Mnożąc przez $e^{-A}$, otrzymujemy wszystkie rozwiązania
$e^{-A}V+T_{ce^{-A}}$.
\end{proof}

\begin{przyklad}[Rozwiązania dystrybucyjne i skoki]
\label{ex:forced-ode20}
Dla $\lambda\in\K$ rozwiązujemy równanie
$T'+\lambda T=\delta_0$ na $\R$.
Funkcja $f(x)=e^{-\lambda x}H(x)$ jest lokalnie całkowalna,
jest klasy $C^1$ poza zerem, ma granice $f(0-)=0$ i
$f(0+)=1$, a jej klasyczna pochodna poza zerem jest równa
$-\lambda f$ i jest lokalnie całkowalna.
Z Twierdzenia~\ref{thm:jump20} dla $E_\lambda=T_f$ mamy
$E_\lambda'=-\lambda E_\lambda+\delta_0$, czyli
$E_\lambda'+\lambda E_\lambda=\delta_0$.
Z Wniosku~\ref{cor:ode20}, przy $A(x)=\lambda x$,
wszystkie rozwiązania mają więc postać
$T=E_\lambda+cT_{e^{-\lambda x}}$, $c\in\K$.
Dla $\lambda=0$ otrzymujemy równanie $T'=\delta_0$
i pierwotną $T_H$.
Ponadto dla $x_+(x)=xH(x)$ z reguły Leibniza mamy
$(T_{x_+})'=T_H+x\delta_0=T_H$, ponieważ
$x\delta_0=0$ (Przykład~\ref{ex:multiply-delta20}).
Różniczkując ponownie, dostajemy
$\frac{\mathrm d^2}{\mathrm dx^2}T_{x_+}=\delta_0$.
\end{przyklad}

\begin{twierdzenie}[Znikanie gradientu na zbiorze spójnym]
\label{thm:zero-gradient20}
Niech $\Omega\subseteq\R^d$ będzie niepustym spójnym
zbiorem otwartym i $T\in\mathcal D'(\Omega;\K)$.
Jeżeli $\partial_jT=0$ dla każdego $j=1,\ldots,d$,
to istnieje $c\in\K$ takie, że
$\langle T,\varphi\rangle=c\int_\Omega\varphi(x)\,
\mathrm d\lambda_d(x)$ dla każdej $\varphi\in\mathcal D(\Omega)$.
\end{twierdzenie}
\begin{proof}
Najpierw rozważmy otwarty prostokąt
$Q=I_1\times\cdots\times I_d$ o domknięciu zwartym w $\Omega$.
Przez $T_Q=T|_Q$ oznaczmy ograniczenie dystrybucji do $Q$.
Zgodność różniczkowania z przedłużaniem zerem, wykazana
w dowodzie Twierdzenia~\ref{thm:dist-deriv-cont20},
daje $\partial_jT_Q=0$ dla każdego $j$.
Dla każdego $j$ wybieramy $\eta_j\in\mathcal D(I_j)$
o całce $1$ jak w dowodzie Twierdzenia~\ref{thm:zero-derivative20}
i dla $\varphi\in\mathcal D(Q)$ określamy operator
\[
 (P_j\varphi)(x)=\eta_j(x_j)
 \int_{I_j}\varphi(x_1,\ldots,x_{j-1},t,x_{j+1},\ldots,x_d)
 \mathrm d\lambda_1(t).
\]
Dla $\varphi\in\mathcal D(Q)$ wybieramy zwarte przedziały
$L_k\Subset I_k$ zawierające rzuty $\supp\varphi$
oraz $\supp\eta_k$ dla $k=1,\ldots,d$.
Całkowanie względem $t$ w definicji $P_j$ można ograniczyć
do $L_j$. Różniczkowanie pod całką względem pozostałych
zmiennych jest dozwolone, gdyż całkujemy funkcję gładką
po ustalonym zwartym przedziale; lokalnie wszystkie
pochodne są ograniczone na odpowiednich iloczynach zwartych.
Zatem $P_j\varphi$ jest gładka.
Jeżeli $x_j\notin L_j$, znika czynnik $\eta_j(x_j)$,
a jeżeli $x_k\notin L_k$ dla pewnego $k\ne j$, znika cała
funkcja podcałkowa. Stąd $\supp(P_j\varphi)\subseteq
L_1\times\cdots\times L_d$, czyli $P_j\varphi\in\mathcal D(Q)$.
Liniowość $P_j$ wynika z liniowości całki.
Dla $j\ne k$ oba złożenia $P_jP_k$ i $P_kP_j$ dają
czynnik $\eta_j(x_j)\eta_k(x_k)$ pomnożony przez całkę
podwójną funkcji $\varphi$ po $I_j\times I_k$ przy pozostałych
zmiennych ustalonych. Twierdzenie Fubiniego (całka jest
bezwzględnie zbieżna) daje ich równość.

Dla $\psi\in\mathcal D(Q)$ oznaczmy
$h_j=(\operatorname{Id}-P_j)\psi$.
Z definicji $P_j$ i $\int_{I_j}\eta_j\,\mathrm d\lambda_1=1$
wynika, że całka $h_j$ po $x_j$ jest zerowa dla każdego
ustalenia pozostałych zmiennych.
Przedłużamy $h_j$ zerem na $\R^d$ i określamy
\[
\theta_j(x)=\int_{-\infty}^{x_j}
h_j(x_1,\ldots,x_{j-1},t,x_{j+1},\ldots,x_d)\,
\mathrm d\lambda_1(t).
\]
Wybieramy zwarte przedziały $L_k=[u_k,v_k]\Subset I_k$
tak, aby ich iloczyn zawierał $\supp h_j$.
Całka w definicji $\theta_j$ jest zerowa, gdy
$x_j<u_j$ albo $x_j>v_j$ (w drugim przypadku dzięki
zerowej całce po całej zmiennej $x_j$), a także gdy
$x_k\notin L_k$ dla pewnego $k\ne j$.
Zatem $\supp\theta_j\subseteq L_1\times\cdots\times L_d$.
Z twierdzenia Newtona--Leibniza i różniczkowania pod całką
względem pozostałych zmiennych $\theta_j$ jest gładka i
$\partial_j\theta_j=h_j$. Ograniczenie $\theta_j$ do $Q$
jest funkcją testową. Ponieważ $\partial_jT_Q=0$,
\[
\langle T_Q,(\operatorname{Id}-P_j)\psi\rangle
=\langle T_Q,\partial_j\theta_j\rangle
=-\langle\partial_jT_Q,\theta_j\rangle=0.
\]
Ostatnia równość mówi, że
$\langle T_Q,\psi\rangle=\langle T_Q,P_j\psi\rangle$
dla każdej $\psi\in\mathcal D(Q)$ i każdego $j$.
Stosując ją kolejno do $\varphi,P_1\varphi,\ldots,
P_{d-1}\cdots P_1\varphi$ i korzystając z przemienności
operatorów, otrzymujemy
$\langle T_Q,\varphi\rangle=\langle T_Q,P_1\cdots P_d\varphi\rangle$.
Wielokrotne zastosowanie definicji $P_j$ i twierdzenia
Fubiniego daje
$P_1\cdots P_d\varphi=(\int_Q\varphi\,\mathrm d\lambda_d)
\eta_1\cdots\eta_d$, gdzie iloczyn $\eta_1\cdots\eta_d$
oznacza funkcję $x\mapsto\prod_{j=1}^d\eta_j(x_j)$.
Ma ona zwarty nośnik w $Q$ i całkę $1$.
Zatem $T_Q=T_{c_Q}$, gdzie
$c_Q=\langle T_Q,\eta_1\cdots\eta_d\rangle$.
Jeśli $Q\cap Q'\ne\varnothing$, wybieramy funkcję testową
$\zeta$ o całce $1$ i nośniku w $Q\cap Q'$ (lemat o
funkcjach odcięcia i normalizacja). Testując oba ograniczenia
na przedłużeniach $\zeta$ zerem, otrzymujemy
$c_Q=\langle T,\zeta\rangle=c_{Q'}$, gdzie w środkowym
wyrażeniu $\zeta$ jest przedłużona zerem na $\Omega$.

Każdy punkt $x\in\Omega$ należy do otwartego prostokąta
$Q$ o domknięciu zwartym w $\Omega$. Przypisujemy mu
$c(x)=c_Q$; poprzedni argument pokazuje, że wartość nie
zależy od wyboru prostokąta. Funkcja $c(x)$ jest lokalnie
stała, gdyż na każdym takim $Q$ jest równa $c_Q$.
Ustalmy $x_0\in\Omega$. Zbiór
$\{x\in\Omega:c(x)=c(x_0)\}$ i jego dopełnienie są otwarte.
Pierwszy zbiór jest niepusty, więc ze spójności $\Omega$
jest równy $\Omega$. Oznaczmy wspólną stałą przez $c$.
Dystrybucja $T-T_c$ zeruje się na każdym takim prostokącie,
a prostokąty pokrywają $\Omega$. Z definicji nośnika
$\supp(T-T_c)=\varnothing$, a z Lematu~\ref{lem:distribution-local19}
wynika $T-T_c=0$. To daje żądaną równość na całym $\Omega$.
Stała $c$ jest jednoznaczna, gdyż można wybrać funkcję
testową o całce $1$ w dowolnym prostokącie $Q\Subset\Omega$
i obliczyć na niej wartość $T$.
\end{proof}

\begin{uwaga}[Granice twierdzeń]
\label{rem:limits20}
Niech $d\ge2$, $I\subseteq\R$ będzie niepustym przedziałem
otwartym, a $V\subseteq\R^{d-1}$ niepustym zbiorem otwartym.
Niech $T\in\mathcal D'(I\times V;\K)$.
Warunek $\partial_1T=0$ oznacza, że $T$ zależy od funkcji
testowej wyłącznie przez jej całkę po pierwszej współrzędnej.
Istotnie, wybierzmy $\eta\in\mathcal D(I)$ o całce $1$.
Dla $\varphi\in\mathcal D(I\times V)$ połóżmy
$A\varphi(y)=\int_I\varphi(t,y)\,\mathrm d\lambda_1(t)$.
Nośnik $A\varphi$ mieści się w zwartym rzucie
$\supp\varphi$ na $V$, a gładkość wynika z różniczkowania
pod całką po ustalonym zwartym podprzedziale $I$,
zawierającym rzut nośnika na $I$.
Funkcja $h(t,y)=\varphi(t,y)-\eta(t)A\varphi(y)$ ma zerową
całkę po $I$ dla każdego $y$.
Po przedłużeniu zerem jej pierwotna względem $t$,
$\theta(t,y)=\int_{-\infty}^t h(s,y)\,\mathrm d\lambda_1(s)$,
jest gładka. Jej nośnik jest zawarty w $[u,v]\times K$,
gdzie $[u,v]\Subset I$ zawiera rzut $\supp h$ na $I$,
a $K\Subset V$ jest rzutem $\supp h$ na $V$.
Zatem $\theta|_{I\times V}$ jest funkcją testową i
$\langle T,h\rangle=-\langle\partial_1T,\theta\rangle=0$.
Otrzymujemy
\[
\langle T,\varphi\rangle
=\left\langle S,\int_I\varphi(t,\cdot)\,\mathrm d\lambda_1(t)\right\rangle,
\qquad
\langle S,\psi\rangle=\langle T,\eta\otimes\psi\rangle,
\]
gdzie $(\eta\otimes\psi)(t,y)=\eta(t)\psi(y)$.
Odwzorowanie $\psi\mapsto\eta\otimes\psi$ jest ciągłe
na każdym etapie $\mathcal D_K(V)$: nośnik obrazu mieści
się w $\supp\eta\times K$, a jego seminorma $p_m$ nie
przekracza $p_m(\eta)p_m(\psi)$. Z własności uniwersalnej
granicy induktywnej jest ciągłe na $\mathcal D(V)$.
Stąd $S\in\mathcal D'(V;\K)$. Dystrybucja $S$ jest jednoznacznie wyznaczona,
ponieważ $A(\eta\otimes\psi)=\psi$ dla każdej $\psi\in\mathcal D(V)$.
Odwrotnie, gdy $T$ ma podaną reprezentację, to
$A(\partial_1\varphi)=0$ z twierdzenia Newtona--Leibniza
i zwartego nośnika $\varphi$, więc $\partial_1T=0$.
Dla przykładu regularna dystrybucja odpowiadająca funkcji
$(t,y)\mapsto y_1$ ma zerową pierwszą pochodną, ale
nie jest stała na niepustym zbiorze otwartym $I\times V$.
Dla układu $\partial_jT=S_j$, gdzie
$T,S_1,\ldots,S_d\in\mathcal D'(\Omega;\K)$, konieczne są warunki zgodności
$\partial_kS_j=\partial_jS_k$, wynikające z przemienności
pochodnych dystrybucyjnych. Ich wystarczalność globalna
może zależeć od topologii zbioru otwartego; nie wynika z tego
rozdziału.
\end{uwaga}

\begin{uwaga}[Związek z operatorami na TVS]
\label{rem:tvs-derivative20}
Konstrukcja $T\mapsto PT$ nie wymaga normy na
$\mathcal D(\Omega)$ ani na $\mathcal D'(\Omega)$.
Opiera się na ciągłości operatora $P^{\mathrm t}$ na
lokalnie wypukłej przestrzeni funkcji testowych.
W ten sposób pochodna dystrybucyjna jest szczególnym
przypadkiem transpozycji ciągłego operatora liniowego
między przestrzeniami liniowo-topologicznymi.
\end{uwaga}

\setcounter{section}{5}
% Kandydat restrukturyzacji; nie jest aktywnym rozdziałem Atlasu.
\section{Dystrybucje o zwartym nośniku}
\providecommand{\atlasref}[2]{#2}

Pracujemy nad $\K\in\{\R,\C\}$. Przez $\mathcal E'(\Omega)$ rozumiemy przestrzeń sprzężoną do $C^\infty(\Omega)$; jej identyfikację z dystrybucjami o zwartym nośniku udowodnimy poniżej.

\subsection{Dystrybucje o zwartym nośniku i funkcje gładkie}

\begin{lemat}[Przedłużenie dystrybucji o zwartym nośniku]
\label{lem:compact-action21}
Niech $T\in\mathcal D'(\R^d;\K)$ i
$K=\supp T\Subset\R^d$. Wówczas dla każdej funkcji
$f\in C^\infty(\R^d;\K)$ istnieje dobrze określona wartość
$\langle T,f\rangle:=\langle T,\chi f\rangle$, gdzie
$\chi\in\mathcal D(\R^d)$ jest równa $1$ na otoczeniu $K$.
Wartość ta nie zależy od wyboru $\chi$. Istnieje także
niepusty zbiór zwarty $L$ i liczba $m\in\N_0$ oraz $C>0$ takie, że
\[
 |\langle T,f\rangle|\le
 C\max_{|\alpha|\le m}\sup_{x\in L}|\partial^\alpha f(x)|
 \quad(f\in C^\infty(\R^d)).
\]
W szczególności $T$ wyznacza ciągły funkcjonał na
przestrzeni Frécheta $C^\infty(\R^d)$.
\end{lemat}
\begin{proof}
Funkcja $\chi$ istnieje na mocy Lematu~\ref{lem:cutoffs17}.
Jeżeli $\chi_1,\chi_2$ są równe $1$ w otoczeniach $K$,
to $(\chi_1-\chi_2)f$ zeruje się na otoczeniu $K$;
jej nośnik jest rozłączny z $K$, więc z Lematu~\ref{lem:distribution-local19}
wynika $\langle T,(\chi_1-\chi_2)f\rangle=0$.
Gdy $f\in\mathcal D(\R^d)$, identyczny argument dla
$(1-\chi)f$ pokazuje, że przedłużenie zgadza się z pierwotnym
działaniem $T$. Liniowość wynika z użycia jednej ustalonej $\chi$.
Dla ustalonej $\chi$ wybieramy niepusty zwarty zbiór $L$
zawierający $\supp\chi$.
Ciągłość $T$ na $\mathcal D_L$ daje nierówność
$|\langle T,\psi\rangle|\le C_0p_m(\psi)$ dla funkcji
$\psi\in\mathcal D_L$. Podstawiamy $\psi=\chi f$
i stosujemy wzór Leibniza. Dokładniej, dla $|\alpha|\le m$ mamy
$\partial^\alpha(\chi f)=\sum_{\beta\le\alpha}
\binom\alpha\beta(\partial^\beta\chi)(\partial^{\alpha-\beta}f)$.
Poza $L$ wszystkie składniki znikają. Na $L$ ich moduły
ograniczamy przez iloczyny odpowiednich supremów.
Możemy zatem przyjąć
\[
 C=C_0\max\left\{1,\max_{|\alpha|\le m}
 \sum_{\beta\le\alpha}\binom\alpha\beta
 \sup_{x\in\R^d}|\partial^\beta\chi(x)|\right\}.
\]
Otrzymujemy żądane oszacowanie, a zatem ciągłość na
$C^\infty(\R^d)$ także w przypadku $T=0$.
\end{proof}

\begin{uwaga}[Przestrzeń dualna przestrzeni $C^\infty$]
\label{rem:Eprime21}
Odwrotnie, niech $F$ będzie ciągłym funkcjonałem liniowym na
$C^\infty(\R^d;\K)$. Ciągłość daje oszacowanie
$|F(f)|\le C\max_{|\alpha|\le m}\sup_{x\in L}|\partial^\alpha f(x)|$
dla pewnego zwartego $L\subseteq\R^d$, $m\in\N_0$ i $C>0$:
wystarczy zastąpić skończoną rodzinę seminorm jedną seminormą,
biorąc unię odpowiednich zbiorów zwartych i maksimum rzędów
(oraz w razie potrzeby powiększając $L$ do zbioru niepustego).
Na każdym etapie $\mathcal D_M$ prawa strona jest ograniczona
przez $Cp_m(f)$. Własność uniwersalna granicy induktywnej
(Twierdzenie~\ref{thm:LF-universal18}) daje ciągłość ograniczenia
$T=F|_{\mathcal D}$, a znikanie $F$ na funkcjach testowych
o nośniku rozłącznym z $L$ daje $\supp T\subseteq L$.
Wybierając $\chi\in\mathcal D(\R^d)$ równą $1$ w otoczeniu $L$, mamy
$F(f-\chi f)=0$ z oszacowania, zatem $F(f)=\langle T,\chi f\rangle$.
To pokazuje również jednoznaczność ciągłego przedłużenia.
Otrzymujemy naturalną identyfikację liniową
$\mathcal E'(\R^d)\cong C^\infty(\R^d)'$ przez ograniczenie
i przedłużenie, przy czym poniżej
nie będziemy utożsamiać bez wyjaśnienia różnych możliwych
topologii na tej przestrzeni dualnej.
\end{uwaga}

\subsection{Przestrzeń sprzężona do funkcji gładkich na zbiorze otwartym}
\begin{twierdzenie}[Dystrybucje o zwartym nośniku]
\label{thm:compact-support-E25}
Niech $\Omega\subseteq\R^d$ będzie niepustym zbiorem otwartym.
Obraz kanonicznego odwzorowania
$\mathcal E'(\Omega)\to\mathcal D'(\Omega)$
składa się dokładnie z dystrybucji o zwartym nośniku.
Dokładniej, dla każdego $F\in\mathcal E'(\Omega)$ istnieją
niepusty zwarty $K\Subset\Omega$, $m\in\N_0$ i $C>0$ takie, że
\[
 |F(f)|\le C\max_{|\alpha|\le m}\sup_{x\in K}
 |\partial^\alpha f(x)|\qquad(f\in\mathcal E(\Omega)).
\]
Odwrotnie, jeżeli $T\in\mathcal D'(\Omega)$ ma zwarty nośnik,
to istnieje dokładnie jeden funkcjonał ciągły
$F\in\mathcal E'(\Omega)$ równy $T$ na $\mathcal D(\Omega)$.
\end{twierdzenie}
\begin{proof}
Z kryterium ciągłości funkcjonału na przestrzeni
$\mathcal E(\Omega)$ (Twierdzenie~\ref{thm:dual-cinf17})
otrzymujemy oszacowanie przez skończenie wiele
seminorm $p_{K_j,m_j}$. Biorąc $K=\bigcup_jK_j$
i $m=\max_jm_j$, otrzymujemy oszacowanie z tezy
(dla funkcjonału zerowego można wybrać dowolny
niepusty zwarty $K$). Ograniczenie
$T=F|_{\mathcal D(\Omega)}$ jest dystrybucją, ponieważ
odwzorowanie $\mathcal D(\Omega)\to\mathcal E(\Omega)$
jest ciągłe na mocy twierdzenia~\ref{thm:zanurzenia25}. Gdy
$\varphi\in\mathcal D(\Omega\setminus K)$, jej
przedłużenie przez zero wraz ze wszystkimi pochodnymi
znika w otoczeniu $K$. Zatem $T(\varphi)=0$,
czyli $\supp T\subseteq K$.
Nośnik jest tu nośnikiem dystrybucji $T$,
a nie zbiorem punktów przestrzeni argumentów funkcjonału $F$.

Niech teraz $K=\supp T\Subset\Omega$. Jeśli $K$ jest
pusty, to $T=0$ z lokalności dystrybucji
(Lemat~\ref{lem:distribution-local19}) i wystarcza $F=0$.
W przeciwnym przypadku wybierzmy
$\chi\in\mathcal D(\Omega)$ równą $1$ w otoczeniu $K$
(Lemat~\ref{lem:cutoffs17}) i połóżmy $L=\supp\chi$.
Kładziemy $F(f)=\langle T,\chi f\rangle$ dla
$f\in\mathcal E(\Omega)$. Ponieważ $\supp(\chi f)\subseteq L$,
z lokalnego oszacowania dystrybucji
(Twierdzenie~\ref{thm:distribution-estimate19})
istnieją $C_L>0$ i $m\in\N_0$ takie, że
\[
 |F(f)|\le C_L p_{L,m}(\chi f)
 \le C_L 2^m p_{L,m}(\chi)p_{L,m}(f).
\]
Ostatnia nierówność wynika z reguły Leibniza i tożsamości
$\sum_{\gamma\le\alpha}\binom{\alpha}{\gamma}=2^{|\alpha|}$.
Stąd $F$ jest liniowym funkcjonałem ciągłym.
Wartość nie zależy od wyboru $\chi$: dla dwóch odcięć
$\chi_1,\chi_2$ funkcja $(\chi_1-\chi_2)f$ jest testowa
i znika w otoczeniu $K$. Jej nośnik jest więc zwartym
podzbiorem $\Omega\setminus K$. Z definicji nośnika
i lokalności dystrybucji wynika, że $T$ ją anihiluje. Jeżeli $f\in\mathcal D(\Omega)$, to
$T(f)=T(\chi f)$, ponieważ $(1-\chi)f$ znika w
otoczeniu $K$. Jednoznaczność przedłużenia wynika
z gęstości $\mathcal D(\Omega)$ w $\mathcal E(\Omega)$.
Otrzymana identyfikacja jest zgodna z przedłużeniem
z Lematu~\ref{lem:compact-action21} oraz Uwagą~
\ref{rem:Eprime21}; nie nadajemy obrazowi automatycznie
topologii podprzestrzeni $\mathcal D'$.
\end{proof}

\setcounter{section}{6}
% Wersja źródłowa: audyt II, jednostka 21

\section{Splot dystrybucji i regularyzacja}
\label{sec:splot21}

\providecommand{\atlasref}[2]{#2}


W całym rozdziale $d\ge1$ jest liczbą całkowitą,
$\K\in\{\R,\C\}$, a dystrybucje
są ciągłymi $\K$-liniowymi funkcjonałami na
$\mathcal D(\R^d;\K)$. Rozpatrujemy splot na całej
$\R^d$; dla ogólnego zbioru otwartego $\Omega$ nie ma
jednej operacji splotu dwóch dystrybucji z
$\mathcal D'(\Omega)$ bez dodatkowych ustaleń.
Oznaczymy przez $\mathcal E'(\R^d)$ przestrzeń dystrybucji
o zwartym nośniku. Nie należy mylić jej ze zbiorem funkcji
testowych $\mathcal D(\R^d)$. Pominięcie $\K$ w oznaczeniu
przestrzeni oznacza pracę nad ustalonym ciałem $\K$.
Dla $f\in L^1_{\mathrm{loc}}(\R^d;\K)$ piszemy $T_f$ na
odpowiadającą mu dystrybucję regularną. W zapisie splotu
$T*\rho$ dla $\rho\in\mathcal D(\R^d)$ rozumiemy
$T*T_\rho$; jeżeli splot jest regularny, jego funkcję
reprezentującą będziemy wskazywać jawnie.
Dla zbiorów $K,F\subseteq\R^d$ oznaczamy
$K+F=\{k+f:k\in K,\ f\in F\}$ i
$K-F=\{k-f:k\in K,\ f\in F\}$.

\subsection{Iterowane działania dystrybucji i definicja splotu}

\begin{lemat}[Przestawianie dwóch działań dystrybucji]
\label{lem:iterated-distributions21}
Niech $k,l\ge1$ będą liczbami całkowitymi,
$T\in\mathcal D'(\R^k)$,
$S\in\mathcal D'(\R^l)$ i
$\Phi\in\mathcal D(\R^{k+l})$. Funkcje
\[
 x\longmapsto\langle S_y,\Phi(x,y)\rangle,
 \qquad
 y\longmapsto\langle T_x,\Phi(x,y)\rangle
\]
należą odpowiednio do $\mathcal D(\R^k)$ i
$\mathcal D(\R^l)$. Ponadto
\[
 \langle T_x,\langle S_y,\Phi(x,y)\rangle\rangle
 =\langle S_y,\langle T_x,\Phi(x,y)\rangle\rangle.
\]
\end{lemat}
\begin{proof}
Ustalmy zwarte prostokąty $K_x\subseteq\R^k$ i
$K_y\subseteq\R^l$, których wnętrza zawierają odpowiednie
rzuty $\supp\Phi$. Dla każdego $x$ funkcja $\Phi(x,\cdot)$
należy do $\mathcal D_{K_y}$. Ilorazy różnicowe względem
$x_j$ zbiegają w każdej seminormie $p_m$ tej przestrzeni
do $\partial_{x_j}\Phi(x,\cdot)$: po różniczkowaniu
względem $y$ błąd jest ograniczony przez moduł ciągłości
odpowiedniej pochodnej na iloczynie zwartego otoczenia $x$
i $K_y$, na mocy twierdzenia Newtona--Leibniza.
Powtarzając ten argument dla kolejnych pochodnych i
stosując ciągłość $S$ na $\mathcal D_{K_y}$, otrzymujemy
\[
 \partial_x^\alpha\langle S_y,\Phi(x,y)\rangle
 =\langle S_y,\partial_x^\alpha\Phi(x,y)\rangle.
\]
Ciągłość tych pochodnych wynika z jednostajnej ciągłości
odpowiednich pochodnych $\Phi$ na tych samych zbiorach.
Analogicznie postępujemy z $T$. Nośniki obu otrzymanych
funkcji mieszczą się w odpowiednich rzutach $\supp\Phi$,
więc są to funkcje testowe.

Istnieją $C_T,C_S>0$ i $m_T,m_S\in\N_0$ takie, że
$|\langle T,u\rangle|\le C_Tp_{m_T}(u)$ dla
$u\in\mathcal D_{K_x}$ i analogicznie dla $S$ na
$\mathcal D_{K_y}$. Z powyższego wzoru na pochodne mamy
\[
 |\langle T_x,\langle S_y,\Phi(x,y)\rangle\rangle|
 \le C_TC_S p_{m_T+m_S}(\Phi).
\]
To samo oszacowanie zachodzi dla odwrotnej kolejności.
W szczególności oba działania są ciągłe na każdym etapie
o ustalonym zwartym nośniku.

Pozostaje uzasadnić równość. Dla $\Phi(x,y)=u(x)v(y)$,
gdzie $u,v$ są funkcjami testowymi, obie strony są równe
$\langle T,u\rangle\langle S,v\rangle$.
Pokażemy, że skończone sumy takich funkcji aproksymują
dowolną $\Phi$ w jednym etapie przestrzeni $\mathcal D$.
Wybierzmy nieujemne rzeczywiste funkcje testowe
$\rho$ na $\R^k$ i $\sigma$ na $\R^l$, o całkach $1$
i nośnikach w kulach jednostkowych, przez normalizację
funkcji odcięcia z Lematu~\ref{lem:cutoffs17}.
Niech $\rho_\varepsilon(x)=\varepsilon^{-k}\rho(x/\varepsilon)$
i $\sigma_\varepsilon(y)=\varepsilon^{-l}\sigma(y/\varepsilon)$.
Dla $0<\varepsilon\le1$ określamy
\[
 \Phi_\varepsilon(x,y)=
 \int_{\R^k\times\R^l}\Phi(u,v)
 \rho_\varepsilon(x-u)\sigma_\varepsilon(y-v)
 \,\mathrm d\lambda_{k+l}(u,v).
\]
Wszystkie te funkcje mają nośniki w ustalonym zwartym
prostokącie $L_x\times L_y$, gdzie $L_x,L_y$ zawierają
odpowiednio $K_x+\overline{B(0,1)}$ i
$K_y+\overline{B(0,1)}$. Po zmianie zmiennych całka
jest średnią przesuniętych funkcji $\Phi$ z wagą
$\rho(u)\sigma(v)$. Dla każdego $m$ jednostajna ciągłość
pochodnych $\Phi$ rzędu co najwyżej $m$ daje
$p_m(\Phi_\varepsilon-\Phi)\to0$.

Dla ustalonego $\varepsilon$ całkę definiującą
$\Phi_\varepsilon$ można ograniczyć do $K_x\times K_y$.
Jej sumy Riemanna są skończonymi sumami funkcji postaci
$c\,\rho_\varepsilon(x-u)\sigma_\varepsilon(y-v)$.
Dla każdego skończonego $m$ sumy te zbiegają w $p_m$:
pochodne względem $(x,y)$ można przenieść pod całkę,
a funkcje podcałkowe po różniczkowaniu są jednostajnie
ciągłe na zwartym iloczynie prostokątów zmiennych
$(x,y)$ i $(u,v)$. Poza $L_x\times L_y$ zarówno sumy,
jak i całka znikają. Dla $\varepsilon=1/n$ wybieramy
sumę $F_n$ z błędem $p_n(F_n-\Phi_{1/n})<1/n$.
Wtedy dla każdego ustalonego $m$ i $n\ge m$ mamy
$p_m(F_n-\Phi)\le 1/n+p_m(\Phi_{1/n}-\Phi)\to0$.
Nośniki wszystkich $F_n$ zawierają się w jednym zwartym
zbiorze $L_x\times L_y$, a funkcje te są sumami
iloczynów funkcji testowych nad ustalonym ciałem $\K$.
Stosując uzyskane oszacowania, tym razem na etapie
o nośniku zawartym w $L_x\times L_y$, do różnicy $F_n-\Phi$,
przechodzimy do granicy w równości obu działań na $F_n$.
\end{proof}

\begin{lemat}[Operator testowy indukowany dystrybucją]
\label{lem:convolution-transpose21}
Niech $T\in\mathcal E'(\R^d)$, $K=\supp T$, i niech
$\varphi\in\mathcal D(\R^d)$. Funkcja
\[
 (A_T\varphi)(y)=\langle T_x,\varphi(x+y)\rangle
\]
należy do $\mathcal D(\R^d)$ i spełnia
$\supp(A_T\varphi)\subseteq\supp\varphi-K$.
Operator $A_T:\mathcal D(\R^d)\to\mathcal D(\R^d)$
jest liniowy i ciągły, a obrazy zbiorów ograniczonych
są ograniczone.
\end{lemat}
\begin{proof}
W działaniu $T_x$ korzystamy z przedłużenia z
Lematu~\ref{lem:compact-action21}. Dla $y$ z ustalonego
zbioru zwartego funkcje $x\mapsto\varphi(x+y)$ mają
wspólny zwarty nośnik; można więc różniczkować po $y$:
$\partial_y^\beta(A_T\varphi)(y)=
\langle T_x,\partial^\beta\varphi(x+y)\rangle$.
Jeżeli $y\notin\supp\varphi-K$, funkcja
$\varphi(\cdot+y)$ zeruje się na otoczeniu $K$,
stąd wartość działania $T$ jest zerowa. Zatem obraz
jest funkcją testową.

Dla ustalonego $L\Subset\R^d$, $m\in\N_0$ i
$\varphi\in\mathcal D_L$ oszacowanie z
Lematu~\ref{lem:compact-action21}, zastosowane do
przesunięć funkcji $\varphi$, daje
\[
 p_m(A_T\varphi)\le C_{T,L,m}\,p_{m+r}(\varphi),
 \qquad \supp(A_T\varphi)\subseteq L-K,
\]
gdzie $r$ jest rzędem oszacowania z Lematu~\ref{lem:compact-action21}.
Istotnie, dla $|\beta|\le m$ stosujemy to oszacowanie do
$f(x)=\partial^\beta\varphi(x+y)$; każda pochodna $f$
rzędu co najwyżej $r$ ma moduł ograniczony przez
$p_{m+r}(\varphi)$ niezależnie od $y$.
Można więc wziąć stałą z tego lematu niezależną od $L,m$.
Zbiór $L-K$ jest zwarty jako obraz $L\times K$ przez
ciągłe odwzorowanie $(u,v)\mapsto u-v$; dla $K=\varnothing$
operator jest zerowy.
To dowodzi ciągłości na każdym etapie $\mathcal D_L$,
a z własności uniwersalnej granicy induktywnej
wynika ciągłość na $\mathcal D$. Ciągłe odwzorowanie liniowe
przenosi zbiory ograniczone na ograniczone: dla otoczenia
zera $U$ w przestrzeni docelowej $A_T^{-1}(U)$ jest
otoczeniem zera, a z $B\subseteq t A_T^{-1}(U)$
wynika $A_T(B)\subseteq tU$.
\end{proof}

\begin{definicja}[Splot dystrybucji]
\label{def:convolution21}
Niech $T\in\mathcal E'(\R^d)$ i
$S\in\mathcal D'(\R^d)$. Ich \emph{splotem} nazywamy
dystrybucję $T*S\in\mathcal D'(\R^d)$ określoną wzorem
\[
 \langle T*S,\varphi\rangle
 =\langle S_y,\langle T_x,\varphi(x+y)\rangle\rangle
 =\langle S,A_T\varphi\rangle,
 \qquad\varphi\in\mathcal D(\R^d).
\]
Dla $S\in\mathcal D'(\R^d)$ i $T\in\mathcal E'(\R^d)$
definiujemy również $S*T:=T*S$, jeżeli $S$ nie ma
zwartego nośnika. Gdy oba czynniki mają zwarty nośnik,
zgodność obu możliwych kolejności zostanie dowiedziona
poniżej. W szczególności nie definiujemy tutaj splotu
dowolnych dwóch dystrybucji bez warunków na nośniki.
\end{definicja}

\begin{twierdzenie}[Istnienie, ciągłość i nośnik splotu]
\label{thm:convolution-properties21}
Niech $T\in\mathcal E'(\R^d)$ i
$S\in\mathcal D'(\R^d)$. Wówczas:
\begin{enumerate}
\item splot $T*S$ jest dobrze określoną dystrybucją;
\item $\supp(T*S)\subseteq\supp T+\supp S$;
\item dla ustalonego $T$ odwzorowanie
$S\mapsto T*S$ jest ciągłym operatorem liniowym na
$\mathcal D'(\R^d)$ zarówno dla topologii słabej-*, jak
i dla topologii silnej;
\item jeżeli także $S$ ma zwarty nośnik, to $T*S$ ma
zwarty nośnik.
\end{enumerate}
\end{twierdzenie}
\begin{proof}
Pierwszy punkt jest konsekwencją ciągłości $A_T$:
złożenie ciągłego funkcjonału $S$ z $A_T$ jest ciągłym
funkcjonałem na $\mathcal D$. Oznaczmy $K=\supp T$, $F=\supp S$. Suma $K+F$
jest domknięta: jeżeli $k_n+f_n\to z$, $k_n\in K$,
$f_n\in F$, to z ciągu $(k_n)$ wybieramy podciąg zbieżny
do $k\in K$, a odpowiedni podciąg $(f_n)$ zbiega do
$z-k\in F$. Dla pustego składnika suma jest pusta.
Jeżeli $\supp\varphi$ jest rozłączny z $K+F$, to
$(\supp\varphi-K)\cap F=\varnothing$, gdyż
$y=u-k\in F$ oznaczałoby $u=k+y\in K+F$.
Z Lematu~\ref{lem:convolution-transpose21} wynika więc
$\supp(A_T\varphi)\cap F=\varnothing$, a z
Lematu~\ref{lem:distribution-local19} otrzymujemy
$\langle S,A_T\varphi\rangle=0$.
Stosując to do wszystkich funkcji testowych o nośniku
w otwartym dopełnieniu $K+F$, dostajemy
$\supp(T*S)\subseteq K+F$.

Dla ustalonej funkcji testowej $\varphi$ zachodzi
$|\langle T*S,\varphi\rangle|=
|\langle S,A_T\varphi\rangle|$, co zapewnia
ciągłość słabą-*. Dla zbioru ograniczonego
$B\subseteq\mathcal D$ mamy
\[
 \sup_{\varphi\in B}|\langle T*S,\varphi\rangle|
 =\sup_{\psi\in A_T(B)}|\langle S,\psi\rangle|.
\]
Lemat~\ref{lem:convolution-transpose21} daje
ograniczoność $A_T(B)$, a więc ciągłość silną.
Jeśli oba nośniki są zwarte, ich suma jest zwarta,
co daje ostatni punkt.
\end{proof}

\subsection{Przemienność, łączność i różniczkowanie}

\begin{twierdzenie}[Przemienność i łączność splotu]
\label{thm:convolution-algebra21}
Niech $T,S,U\in\mathcal D'(\R^d)$.
Jeśli co najmniej jeden z czynników $T,S$ ma zwarty nośnik,
to $T*S=S*T$, a wynik nie zależy od sposobu użycia
Definicji~\ref{def:convolution21}. Jeśli co najmniej dwa
spośród $T,S,U$ mają zwarty nośnik, to
\[
 (T*S)*U=T*(S*U).
\]
W szczególności $\mathcal E'(\R^d)$ jest przemienną
algebrą z jedynką $\delta_0$ względem splotu,
a $\mathcal D'(\R^d)$ jest modułem nad tą algebrą.
\end{twierdzenie}
\begin{proof}
Jeżeli $T$ i $S$ mają zwarty nośnik, wybieramy funkcje
odcięcia $\chi_T,\chi_S$, równe $1$ na otoczeniach
odpowiednich nośników. Dla $\varphi\in\mathcal D$ funkcja
$\Phi(x,y)=\chi_T(x)\chi_S(y)\varphi(x+y)$
należy do $\mathcal D(\R^{2d})$. Dla każdego $y$
mnożenie przez $\chi_T$ nie zmienia działania $T_x$
na $\varphi(x+y)$, na mocy Lematu~\ref{lem:compact-action21}.
Zatem $\langle T_x,\Phi(x,y)\rangle=\chi_S(y)A_T\varphi(y)$.
Ponieważ $A_T\varphi$ jest funkcją testową, mnożenie jej
przez $\chi_S$ nie zmienia działania $S$. Otrzymujemy
$\langle S_y,\langle T_x,\Phi(x,y)\rangle\rangle
=\langle T*S,\varphi\rangle$. Analogicznie w odwrotnej
kolejności otrzymujemy $\langle S*T,\varphi\rangle$.
Równość wynika z Lematu~\ref{lem:iterated-distributions21}.
Jeśli tylko jeden czynnik ma zwarty nośnik,
przemienność jest częścią definicji, a powyższy
argument zapewnia jej niesprzeczność.

Załóżmy najpierw, że $T,S\in\mathcal E'(\R^d)$,
a $U\in\mathcal D'(\R^d)$. Z Twierdzenia~
\ref{thm:convolution-properties21} splot $T*S$ ma zwarty
nośnik, więc wszystkie potrzebne sploty są określone.
Dla $\varphi\in\mathcal D(\R^d)$ i ustalonego $z\in\R^d$
funkcja $\varphi_z(w)=\varphi(w+z)$ jest testowa.
Z definicji splotu mamy
\[
\begin{aligned}
 (A_{T*S}\varphi)(z)
 &=\langle T*S,\varphi_z\rangle\\
 &=\langle S_y,\langle T_x,\varphi(x+y+z)\rangle\rangle\\
 &=\langle S_y,(A_T\varphi)(y+z)\rangle
 =(A_S A_T\varphi)(z).
\end{aligned}
\]
Wszystkie działania są dozwolone na mocy
Lematu~\ref{lem:convolution-transpose21}: wewnętrzne
działanie daje przesunięcie funkcji testowej $A_T\varphi$.
Zatem $A_{T*S}=A_S A_T$. Z przemienności dla zwartych
nośników mamy również $A_S A_T=A_T A_S$.
W konsekwencji
\[
\begin{aligned}
 \langle (T*S)*U,\varphi\rangle
 &=\langle U,A_{T*S}\varphi\rangle
 =\langle U,A_S A_T\varphi\rangle\\
 &=\langle S*U,A_T\varphi\rangle
 =\langle T*(S*U),\varphi\rangle .
\end{aligned}
\]
Dla przypadku, w którym $T,U$ mają zwarte nośniki,
przemienność i wykazana już tożsamość dają
\[
\begin{aligned}
 (T*S)*U&=U*(T*S)=(U*T)*S\\
 &=(T*U)*S=T*(U*S)=T*(S*U).
\end{aligned}
\]
Gdy $S,U$ mają zwarte nośniki, analogicznie
\[
 (T*S)*U=U*(S*T)=(U*S)*T=T*(S*U).
\]
W każdym z tych wzorów przynajmniej jeden czynnik
każdego splotu ma zwarty nośnik, gdyż splot dwóch
dystrybucji o zwartych nośnikach ma zwarty nośnik.
W ten sposób wykazaliśmy wszystkie przypadki łączności.

Dla każdej $\varphi\in\mathcal D$ zachodzi
$A_{\delta_0}\varphi=\varphi$, więc $\delta_0*T=T$.
Liniowość w drugim czynniku wynika z liniowości
parowania, a w pierwszym z liniowości
$T\mapsto A_T$ na $\mathcal E'(\R^d)$: dla kombinacji
liniowej dwóch dystrybucji o zwartych nośnikach można
użyć jednego odcięcia równego $1$ w otoczeniu zbioru
$\supp T_1\cup\supp T_2$ dla dwóch dystrybucji
$T_1,T_2\in\mathcal E'(\R^d)$.
Nośnik każdej ich kombinacji liniowej zawiera się
w tym zbiorze, więc z Lematu~\ref{lem:compact-action21}
możemy użyć wspólnego odcięcia we wszystkich parowaniach.
Stąd $A_{bT_1+cT_2}=bA_{T_1}+cA_{T_2}$ dla $b,c\in\K$.
Stąd przemienna algebra z jedynką oraz struktura
modułu, zgodnie z udowodnioną łącznością.
\end{proof}

\begin{definicja}[Translacja dystrybucji]
\label{def:translation21}
Dla $a\in\R^d$ i $T\in\mathcal D'(\R^d)$ definiujemy
$\tau_aT$ wzorem
$\langle\tau_aT,\varphi\rangle=
\langle T,\varphi(\cdot+a)\rangle$.
Odwzorowanie $\varphi\mapsto\varphi(\cdot+a)$
przekształca $\mathcal D_K$ w $\mathcal D_{K-a}$ i
zachowuje wszystkie seminormy $p_m$. Jest więc ciągłe
na $\mathcal D$ z własności uniwersalnej granicy
induktywnej; $\tau_aT$ jest dobrze określoną dystrybucją.
Dla $f\in L^1_{\mathrm{loc}}(\R^d;\K)$ podstawienie
$z=x+a$ w całce daje
$\langle\tau_aT_f,\varphi\rangle
=\int_{\R^d}f(z-a)\varphi(z)\,\mathrm d\lambda_d(z)$.
Zatem $\tau_aT_f=T_{f(\cdot-a)}$.
\end{definicja}

\begin{przyklad}[Dystrybucje Diraca i przesunięcia]
\label{ex:delta-convolution21}
Dla $a,b\in\R^d$ i $T\in\mathcal D'$ mamy
$\delta_a*T=\tau_aT$, $\delta_a*\delta_b=\delta_{a+b}$.
Dla wielowskaźników $\alpha,\beta\in\N_0^d$ otrzymujemy
$\partial^\alpha\delta_0*\partial^\beta\delta_0
=\partial^{\alpha+\beta}\delta_0$.
Istotnie, $A_{\delta_a}\varphi(y)=\varphi(a+y)$,
co daje pierwszą tożsamość, a dla $T=\delta_b$
także drugą. Dla pochodnych delty mamy
$A_{\partial^\alpha\delta_0}\varphi(y)
=(-1)^{|\alpha|}\partial^\alpha\varphi(y)$.
Zastosowanie $\partial^\beta\delta_0$ do tej funkcji
daje $(-1)^{|\alpha|+|\beta|}
\partial^{\alpha+\beta}\varphi(0)$, czyli działanie
$\partial^{\alpha+\beta}\delta_0$.
\end{przyklad}

\begin{twierdzenie}[Pochodne splotu]
\label{thm:convolution-derivative21}
Niech $T,S\in\mathcal D'(\R^d)$ i niech co najmniej
jeden z nich ma zwarty nośnik. Dla każdego
$\alpha\in\N_0^d$ zachodzi
\[
 \partial^\alpha(T*S)
 =(\partial^\alpha T)*S
 =T*(\partial^\alpha S).
\]
W szczególności dla wielomianu $P$ w $d$ zmiennych
i operatora stałowspółczynnikowego $P(\partial)$
operator splotu z ustaloną dystrybucją o zwartym
nośniku komutuje z $P(\partial)$.
\end{twierdzenie}
\begin{proof}
Załóżmy, że $T$ ma zwarty nośnik $K$. Ustalmy
$j\in\{1,\ldots,d\}$.
Z Twierdzenia~\ref{thm:dist-deriv-cont20} wszystkie
pochodne $T$ mają nośniki zawarte w $K$, więc ich sploty
z $S$ są określone. Dla $f\in C^\infty(\R^d)$
wzór na pochodną dystrybucji pozostaje prawdziwy
dla przedłużonego działania:
\[
 \langle\partial_jT,f\rangle
 =-\langle T,\partial_jf\rangle.
\]
Istotnie, bierzemy $\chi\in\mathcal D$ równą $1$
w otoczeniu $K$. Wówczas
$\langle\partial_jT,\chi f\rangle
=-\langle T,\chi\partial_jf\rangle
-\langle T,(\partial_j\chi)f\rangle$.
Ostatni składnik znika, gdyż nośnik funkcji
$(\partial_j\chi)f$ jest rozłączny z $K$.
Pierwszy składnik jest przedłużonym działaniem
$-T$ na $\partial_jf$. Dla
$\varphi\in\mathcal D$ i $j\in\{1,\ldots,d\}$ mamy
\[
 \begin{aligned}
 \langle\partial_j(T*S),\varphi\rangle
 &=-\langle S_y,\langle T_x,\partial_j\varphi(x+y)\rangle\rangle\\
 &=\langle S_y,\langle(\partial_jT)_x,\varphi(x+y)\rangle\rangle
 =\langle(\partial_jT)*S,\varphi\rangle.
 \end{aligned}
\]
W drugim kroku różniczkowanie dotyczy zmiennej $x$.
Aby uzyskać $T*(\partial_jS)$, zauważamy, że
$\partial_{y_j}A_T\varphi=A_T(\partial_j\varphi)$,
więc definicja pochodnej $S$ daje tę samą wartość.
Dla $\alpha=0$ wzór jest oczywisty. Indukcja względem
$|\alpha|$, wykorzystująca zachowanie zwartości nośnika
pochodnych $T$, daje obie tożsamości dla każdego
wielowskaźnika. Gdy zwarty nośnik ma tylko $S$,
stosujemy przemienność. Dla wielomianu
$P(\xi)=\sum_\alpha c_\alpha\xi^\alpha$ sumujemy
równości $T*(\partial^\alpha S)=\partial^\alpha(T*S)$
z wagami $c_\alpha$, co daje ostatnie stwierdzenie.
\end{proof}

\begin{przyklad}[Inkluzja nośników może być ostra]
\label{ex:support-strict21}
Niech $T=\delta_0-\delta_1\in\mathcal E'(\R)$,
a $S=T_1$ będzie dystrybucją odpowiadającą funkcji
stałej $1$. Wtedy $T*S=T_1-T_1=0$,
mimo że $\supp T+\supp S=\R$. Dlatego
inkluzji nośników nie można zastąpić równością.
\end{przyklad}

\subsection{Splot z funkcją testową i aproksymacja jedności}

\noindent\emph{Powiązanie tematyczne w Bibliotece matematycznej:} \atlasref{PDE-07::thm:pde7-mollify}{PDE-07: wygładzanie przez splot}.


\begin{twierdzenie}[Wygładzanie dystrybucji]
\label{thm:smoothing21}
Niech $T\in\mathcal D'(\R^d)$ i
$\rho\in\mathcal D(\R^d)$. Wówczas splot
$T*\rho$ jest dystrybucją regularną odpowiadającą
funkcji $u\in C^\infty(\R^d)$ danej przez
\[
 u(x)=\langle T_y,\rho(x-y)\rangle.
\]
Dla każdego wielowskaźnika $\alpha$ zachodzi
$\partial^\alpha u(x)
=\langle T_y,(\partial^\alpha\rho)(x-y)\rangle$.
Ponadto $\supp u\subseteq\supp T+\supp\rho$.
\end{twierdzenie}
\begin{proof}
Ustalmy zwarty zbiór $M\subseteq\R^d$ i oznaczmy
$K=\supp\rho$. Dla $x\in M$ funkcje
$y\mapsto\rho(x-y)$ mają nośniki w $M-K$.
Przy badaniu ilorazów różnicowych powiększamy $M$
do zwartego otoczenia badanego punktu. Względem każdej
seminormy $p_m$ iloraz różnicowy w kierunku $e_j$
zbiega do $y\mapsto(\partial_j\rho)(x-y)$:
z twierdzenia Newtona--Leibniza błąd jest ograniczony
przez moduł ciągłości pochodnych $\rho$ rzędu co najwyżej
$m+1$. Wszystkie te funkcje mają wspólny zwarty nośnik.
Ciągłość $T$ na odpowiednim etapie $\mathcal D_L$
pozwala przenieść różniczkowanie pod znak parowania.
Powtarzając argument dla kolejnych pochodnych,
otrzymujemy podany wzór; ciągłość tych pochodnych
wynika z jednostajnej ciągłości pochodnych $\rho$.
Jeśli $x\notin\supp T+K$, to
$(x-K)\cap\supp T=\varnothing$. Lokalność dystrybucji
(Lemat~\ref{lem:distribution-local19}) daje $u(x)=0$.
Ponieważ $\supp T+K$ jest domknięty, otrzymujemy
żądaną inkluzję nośników.

Ustalmy $\varphi\in\mathcal D(\R^d)$ i połóżmy
$G(x,y)=\rho(x-y)\varphi(x)$. Jego nośnik zawiera się
w zwartym zbiorze
\[
 \supp\varphi\times(\supp\varphi-K),
\]
więc $G\in\mathcal D(\R^{2d})$.
Niech $T_1$ oznacza dystrybucję regularną odpowiadającą
funkcji stałej $1$. Stosujemy
Lemat~\ref{lem:iterated-distributions21} do $T_1$ w zmiennej
$x$ i $T$ w zmiennej $y$. Daje to
\[
 \begin{aligned}
 \int_{\R^d}u(x)\varphi(x)\,\mathrm d\lambda_d(x)
 &=\left\langle T_y,
      \int_{\R^d}\rho(x-y)\varphi(x)\,\mathrm d\lambda_d(x)
   \right\rangle\\
 &=\langle T,\check\rho*\varphi\rangle
 =\langle T* T_\rho,\varphi\rangle.
 \end{aligned}
\]
Tutaj $\check\rho(z)=\rho(-z)$, a klasyczny splot
funkcji definiujemy przez
$(\check\rho*\varphi)(y)=\int_{\R^d}\check\rho(y-x)
\varphi(x)\,\mathrm d\lambda_d(x)$.
Jest to dokładnie $A_{T_\rho}\varphi(y)$: podstawienie
$x=y+z$ daje
\[
 A_{T_\rho}\varphi(y)=
 \int_{\R^d}\rho(z)\varphi(y+z)\,\mathrm d\lambda_d(z).
\]
Ostatnia równość wynika więc z definicji splotu
dystrybucji i jego przemienności. Nie wymagamy
zwartego nośnika $T$. W rezultacie $T*T_\rho=T_u$.
\end{proof}

\begin{definicja}[Standardowa aproksymacja jedności]
\label{def:mollifier21}
Wybierzmy $\rho\in\mathcal D(\R^d)$ taką, że
$\rho\ge0$, $\supp\rho\subseteq\overline{B(0,1)}$
oraz $\int_{\R^d}\rho(x)\,\mathrm d\lambda_d(x)=1$.
Istnienie $\rho$ wynika z Lematu~\ref{lem:cutoffs17}:
wybieramy nieujemną funkcję testową w kuli jednostkowej,
równą $1$ na pewnej mniejszej kuli, i dzielimy ją przez
jej dodatnią całkę. Dla $\varepsilon>0$ piszemy
$\rho_\varepsilon(x)=\varepsilon^{-d}\rho(x/\varepsilon)$.
Zmiana zmiennych $x=\varepsilon z$ pokazuje, że
$\int_{\R^d}\rho_\varepsilon(x)\,\mathrm d\lambda_d(x)=1$,
a $\supp\rho_\varepsilon\subseteq\overline{B(0,\varepsilon)}$.
\end{definicja}

\begin{twierdzenie}[Regularyzacja jest silnie zbieżna]
\label{thm:strong-mollification21}
Dla każdego $T\in\mathcal D'(\R^d)$ i każdej
standardowej aproksymacji jedności z
Definicji~\ref{def:mollifier21} dystrybucje
$T*\rho_\varepsilon$ są reprezentowane przez funkcje
$C^\infty(\R^d)$ oraz
\[
 T*\rho_\varepsilon\longrightarrow T
 \quad\text{w topologii silnej }\beta(\mathcal D',\mathcal D)
 \quad(\varepsilon\downarrow0).
\]
Jeśli $T$ ma zwarty nośnik $K$, to
$\supp(T*\rho_\varepsilon)\subseteq K+
\overline{B(0,\varepsilon)}$, więc dla małych
$\varepsilon$ funkcje wygładzające mają nośniki
w jednym zbiorze zwartym.
\end{twierdzenie}
\begin{proof}
Z Twierdzenia~\ref{thm:smoothing21} mamy wygładzanie.
Dla $\varphi\in\mathcal D$ zachodzi
\[
 \langle T*\rho_\varepsilon-T,\varphi\rangle
 =\langle T,\check\rho_\varepsilon*\varphi-\varphi\rangle.
\]
Niech $B\subseteq\mathcal D$ będzie niepustym zbiorem
ograniczonym (dla pustego zbioru warunek jest trywialny).
Z Twierdzenia~\ref{thm:D-bounded18}
wynika istnienie $K\Subset\R^d$ zawierającego nośniki
wszystkich $\varphi\in B$, przy czym
$M_m(B)=\sup_{\varphi\in B}p_m(\varphi)<\infty$
dla wszystkich $m$. Dla $0<\varepsilon\le1$
różnice $\check\rho_\varepsilon*\varphi-\varphi$
mają nośniki w jednym zwartym zbiorze
$L=K+\overline{B(0,1)}$. Dla $m\in\N_0$ i $|\alpha|\le m$ z normalizacji jądra
i podstawienia $z=x/\varepsilon$ mamy
\[
\begin{aligned}
 \partial^\alpha(\check\rho_\varepsilon*\varphi-\varphi)(y)
 =\int_{\R^d}\rho(z)
 \bigl(\partial^\alpha\varphi(y+\varepsilon z)
 -\partial^\alpha\varphi(y)\bigr)\,\mathrm d\lambda_d(z).
\end{aligned}
\]
Różniczkowanie pod całką jest uzasadnione zwartym
nośnikiem $\rho$ i ograniczonością pochodnych $\varphi$.
Z twierdzenia Newtona--Leibniza wzdłuż odcinka
$t\mapsto y+t\varepsilon z$ otrzymujemy
\[
 |\partial^\alpha\varphi(y+\varepsilon z)
 -\partial^\alpha\varphi(y)|
 \le\varepsilon\|z\|_1 p_{m+1}(\varphi),
 \qquad \|z\|_1=\sum_{j=1}^d|z_j|.
\]
Nierówność obowiązuje również dla funkcji zespolonych.
Biorąc supremum po $y$ i $\varphi\in B$ oraz maksimum
po $|\alpha|\le m$, dostajemy
\[
 \sup_{\varphi\in B}p_m(\check\rho_\varepsilon*\varphi
 -\varphi)
 \le C_{\rho,d}\varepsilon\,M_{m+1}(B),
 \quad
 C_{\rho,d}=\int_{\R^d}\|z\|_1\rho(z)\,\mathrm d\lambda_d(z).
\]
Ciągłość $T|_{\mathcal D_L}$ oznacza, że dla pewnych
$m,C$ mamy $|\langle T,\psi\rangle|\le Cp_m(\psi)$
gdy $\psi\in\mathcal D_L$. Dlatego
\[
 \sup_{\varphi\in B}|\langle T*\rho_\varepsilon-T,
 \varphi\rangle|
 \le CC_{\rho,d}\varepsilon M_{m+1}(B)\longrightarrow0.
\]
Ponieważ $B$ jest dowolnym zbiorem ograniczonym,
jest to zbieżność silna. Oszacowanie nośnika wynika
z Twierdzenia~\ref{thm:smoothing21}.
\end{proof}

\begin{wniosek}[Gęstość funkcji testowych dla zwartych nośników]
\label{cor:dense-test-compact21}
Niech $T\in\mathcal E'(\R^d)$. Istnieje ciąg
$(f_n)\subseteq\mathcal D(\R^d)$ taki, że
$T_{f_n}\to T$ silnie zarówno w $\mathcal D'(\R^d)$,
jak i w $\mathcal E'(\R^d)$.
Można wziąć $f_n(x)=\langle T_y,\rho_{1/n}(x-y)\rangle$,
czyli funkcję reprezentującą dystrybucję $T*T_{\rho_{1/n}}$.
\end{wniosek}
\begin{proof}
Funkcje $f_n$ są gładkie i mają nośniki zwarte
z poprzedniego twierdzenia. Zbieżność silna w $\mathcal D'(\R^d)$ także
wynika z tego twierdzenia.
Jeśli $K=\supp T$, wszystkie nośniki $f_n$ i nośnik $T$
zawierają się w jednym zwartym zbiorze
$K+\overline{B(0,1)}$. Z Twierdzenia~\ref{thm:dt-Eprime-fixed-support}
obie topologie silne na tej podprzestrzeni pokrywają się.
Stąd zbieżność silna w $\mathcal E'(\R^d)$.
\end{proof}

\begin{przyklad}[Regularyzacja delty Diraca]
\label{ex:delta-smooth21}
Dla $a\in\R^d$ i standardowej aproksymacji jedności
$\rho_\varepsilon$ zachodzi
$\delta_a*\rho_\varepsilon=T_{\rho_\varepsilon(\cdot-a)}$.
Funkcje te stają się coraz bardziej skupione
w otoczeniu $a$, a jako dystrybucje zbiegają
silnie do $\delta_a$. Nie ma tu zbieżności w topologii $L^1_{\mathrm{loc}}$
do żadnej funkcji $f$: gdyby taka zbieżność zachodziła,
to dla każdej $\varphi\in\mathcal D(\R^d)$
\[
 |\langle T_{\rho_\varepsilon(\cdot-a)}-T_f,\varphi\rangle|
 \le p_0(\varphi)
 \int_{\supp\varphi}|\rho_\varepsilon(x-a)-f(x)|
 \,\mathrm d\lambda_d(x)\longrightarrow0.
\]
Jednoznaczność granicy słabej-* dałaby $T_f=\delta_a$,
co jest niemożliwe. Istotnie, ograniczając obie dystrybucje do
$\R^d\setminus\{a\}$ i stosując
Stwierdzenie~\ref{prop:regular-injective19}, otrzymalibyśmy, że
$f=0$ prawie wszędzie poza $\{a\}$, a zatem prawie wszędzie
w $\R^d$, podczas gdy $\delta_a\ne0$.
\end{przyklad}

\subsection{Ograniczenia operacji splotu i powiązania z dalszą teorią}

\begin{uwaga}[Kiedy splot nie jest zdefiniowany]
\label{rem:convolution-limits21}
Warunek zwartego nośnika przynajmniej jednego
czynnika jest wystarczający, ale niekonieczny.
Dla funkcji odpowiednio całkowalnych splot może istnieć
bez tego warunku. Natomiast nie można rozszerzyć splotu na wszystkie pary
dystrybucji tak, aby zachować jego dotychczasowe wartości
i obie równości dotyczące różniczkowania.
Już dla $d=1$ niech $H$ będzie funkcją Heaviside'a,
a $T_1$ dystrybucją funkcji stałej $1$.
Z rozdziału~\ref{sec:pochodne20} mamy $\partial T_H=\delta_0$ i
$\partial T_1=0$. Gdyby istniało takie rozszerzenie,
to jednocześnie
\[
 \partial(T_1*T_H)=(\partial T_1)*T_H=0
 \quad\text{oraz}\quad
 \partial(T_1*T_H)=T_1*\delta_0=T_1,
\]
co jest sprzeczne z $T_1\ne0$.
Dla $d>1$ zamiast $T_H$ bierzemy funkcjonał $V$
określony dla $\varphi\in\mathcal D(\R^d)$ przez
$\langle V,\varphi\rangle=\int_0^\infty
\varphi(t,0,\ldots,0)\,\mathrm d\lambda_1(t)$.
Na każdym etapie $\mathcal D_K$ działanie to jest
ograniczone przez $p_0(\varphi)$ razy długość
przedziału zawierającego rzut $K$ na pierwszą oś,
więc jest ciągłe na każdym etapie. Liniowość wynika z
liniowości całki, a własność uniwersalna granicy induktywnej
daje $V\in\mathcal D'(\R^d)$. Dla $\varphi\in\mathcal D(\R^d)$
z twierdzenia Newtona--Leibniza mamy
\[
 \langle\partial_1V,\varphi\rangle
 =-\int_0^\infty(\partial_1\varphi)(t,0,\ldots,0)
 \,\mathrm d\lambda_1(t)=\varphi(0).
\]
Zatem $\partial_1 V=\delta_0$, podczas gdy
$\partial_1 T_1=0$. Uzyskujemy tę samą sprzeczność.

Na przykład formalny splot dwóch funkcji stałych
na $\R^d$ wymagałby całkowania stałej $1$
po całej $\R^d$, co nie daje skończonej wartości.
Dalsze rozszerzenia opierają się między innymi
na warunkach dotyczących nośników, zapewniających
właściwość odwzorowania dodawania na ich iloczynie
(tzn. zwartość przeciwobrazu każdego zbioru zwartego),
lub na dodatkowych własnościach analitycznych czynników.
\end{uwaga}

\subsection{Splot przy właściwym dodawaniu nośników}

\begin{definicja}[Warunek właściwości dodawania]
\label{def:dt-convolution-proper}
Niech $K,F\subseteq\R^d$ będą domknięte.
Mówimy, że dodawanie jest właściwe na $K\times F$,
jeśli dla każdego zwartego $L\subseteq\R^d$ zbiór
\[
 P_L=\{(x,y)\in K\times F:x+y\in L\}
\]
jest zwarty w $\R^{2d}$. Jest to właściwość odwzorowania
$(x,y)\mapsto x+y$ ograniczonego do $K\times F$.
Warunek zachodzi w szczególności, gdy $K$ jest zwarty:
wtedy $P_L$ jest domknięty jako przecięcie $K\times F$
z przeciwobrazem domkniętego $L$ przez ciągłe dodawanie.
Ponadto $P_L\subseteq K\times(L-K)$, ponieważ
$x+y\in L$ i $x\in K$ implikują $y\in L-K$.
Zbiór $L-K$ jest zwarty jako obraz zwartego $L\times K$
przez ciągłe odejmowanie. Zatem $P_L$ jest domkniętym
podzbiorem zbioru zwartego i jest zwarty.
Analogicznie, gdy $F$ jest zwarty.
\end{definicja}

\begin{twierdzenie}[Splot przy właściwym dodawaniu]
\label{thm:dt-convolution-proper}
Niech $T,S\in\mathcal D'(\R^d)$ i niech
$K=\supp T$, $F=\supp S$ spełniają warunek
Definicji~\ref{def:dt-convolution-proper}.
Wtedy istnieje dobrze określona dystrybucja
$T*S$ dana dla $\varphi\in\mathcal D(\R^d)$ wzorem
\[
 \langle T*S,\varphi\rangle
 =\langle T_x,\langle S_y,\chi(x,y)\varphi(x+y)\rangle\rangle,
\]
gdzie $\chi\in\mathcal D(\R^{2d};\R)$ jest równa $1$
w otoczeniu $P_{\supp\varphi}$.
Wartość nie zależy od wyboru $\chi$.
Konstrukcja jest zgodna z Definicją~\ref{def:convolution21},
jest przemienna i spełnia
\[
 \supp(T*S)\subseteq K+F,\qquad
 \partial^\alpha(T*S)=(\partial^\alpha T)*S
 =T*(\partial^\alpha S)
 \quad(\alpha\in\N_0^d).
\]
\end{twierdzenie}
\begin{proof}
Dla $\Phi\in\mathcal D(\R^{2d})$ oznaczmy
\[
 \langle W,\Phi\rangle
 =\langle T_x,\langle S_y,\Phi(x,y)\rangle\rangle.
\]
Lemat~\ref{lem:iterated-distributions21} zapewnia
liniowość i oszacowanie na każdym ustalonym zwartym
nośniku, a Twierdzenie~\ref{thm:distribution-estimate19}
daje $W\in\mathcal D'(\R^{2d})$.
Wykażemy, że $\supp W\subseteq K\times F$.
Jeżeli $(a,b)\notin K\times F$, to $a\notin K$ lub
$b\notin F$. W pierwszym przypadku wybieramy otwarte
otoczenie $U$ punktu $a$ rozłączne z $K$ oraz otwarte
otoczenie $V$ punktu $b$. Dla każdej funkcji testowej
$\Phi$ o nośniku zwartym w $U\times V$ funkcja
$x\mapsto\langle S_y,\Phi(x,y)\rangle$ ma nośnik
w zwartym rzucie $\supp\Phi$ na $U$, rozłącznym z $K$.
Z Lematu~\ref{lem:distribution-local19} otrzymujemy
$\langle W,\Phi\rangle=0$. W drugim przypadku
zamieniamy kolejność działań na mocy
Lematu~\ref{lem:iterated-distributions21} i stosujemy
ten sam argument z $S$ na zewnątrz.
Z definicji nośnika wynika żądana inkluzja.

Dla ustalonej $\varphi$ zbiór $P_{\supp\varphi}$
jest zwarty. Lemat~\ref{lem:cutoffs17} daje odpowiednią
funkcję $\chi$; jeśli zbiór ten jest pusty, można wziąć
$\chi=0$. Iloczyn $\chi(x,y)\varphi(x+y)$ jest funkcją
testową, gdyż jest gładki i ma nośnik w $\supp\chi$.
Dla dwóch dopuszczalnych odcięć $\chi_1,\chi_2$
funkcja
\[
 (\chi_1(x,y)-\chi_2(x,y))\varphi(x+y)
\]
znika w otoczeniu każdego punktu $K\times F$:
na $P_{\supp\varphi}$ wynika to z równości obu odcięć
z $1$ w otoczeniu tego zbioru, a poza nim
z $x+y\notin\supp\varphi$.
Suma tych otoczeń jest otwarta i zawiera $K\times F$,
a więc zawiera $\supp W$. Nośnik różnicy obu testów
jest rozłączny z $\supp W$. Lemat~\ref{lem:distribution-local19}
zapewnia niezależność wyniku od odcięcia.

Dla $\varphi,\psi\in\mathcal D(\R^d)$ i $b,c\in\K$
wybieramy $L=\supp\varphi\cup\supp\psi$.
Ponieważ $\supp(b\varphi+c\psi)\subseteq L$,
jedno odcięcie $\chi_L$ równe $1$ w otoczeniu $P_L$
jest dopuszczalne dla wszystkich trzech testów.
Z liniowości $W$ otrzymujemy
\[
 \langle T*S,b\varphi+c\psi\rangle
 =b\langle T*S,\varphi\rangle+c\langle T*S,\psi\rangle.
\]
Aby wykazać ciągłość, ustalmy zwarty $L\subseteq\R^d$
i wybierzmy $\chi_L$ równą $1$ w otoczeniu $P_L$.
Można jej używać dla wszystkich $\varphi\in\mathcal D_L$.
Jeśli $\chi_L=0$, otrzymujemy na tym etapie funkcjonał
zerowy. W przeciwnym przypadku z ciągłości $W$ na etapie
$\mathcal D_{\supp\chi_L}$ istnieją $C>0$ i $m\in\N_0$,
dla których
\[
 |\langle W,\chi_L(\varphi\circ a)\rangle|
 \le C p_m(\chi_L(\varphi\circ a))
 \le C2^m p_m(\chi_L)p_m(\varphi),
 \qquad a(x,y)=x+y.
\]
Istotnie, pochodna $\partial_x^\mu\partial_y^\nu
(\varphi\circ a)$ jest równa
$(\partial^{\mu+\nu}\varphi)\circ a$, a reguła Leibniza
dla iloczynu z $\chi_L$ daje drugą nierówność.
Twierdzenie~\ref{thm:distribution-estimate19} zapewnia
więc, że $T*S$ jest dystrybucją.

Jeżeli $T=0$ lub $S=0$, obie konstrukcje dają
dystrybucję zerową; możemy więc w dowodzie zgodności
założyć, że oba nośniki są niepuste.
Dla testu zerowego zgodność jest natychmiastowa.
Jeżeli $K$ jest zwarty, z Lematu~\ref{lem:cutoffs17}
wybieramy $\chi_T$ równą $1$ w otoczeniu $K$. Dla $L=\supp\varphi$ i
$M=\supp\chi_T$ zbiór $L-M$ jest zwarty.
Z Lematu~\ref{lem:cutoffs17} wybieramy
$\chi_S$ równą $1$ w otoczeniu $L-M$.
Iloczyn $\chi(x,y)=\chi_T(x)\chi_S(y)$ jest dopuszczalnym
odcięciem, ponieważ dla $(x,y)\in P_L$ mamy
$x\in K\subseteq M$ oraz $y\in L-M$.
Iloczyn otwartych otoczeń $K$ i $L-M$, na których
odpowiednie odcięcia są równe $1$, jest otoczeniem
$P_L$, na którym $\chi=1$.
Ponadto $\chi_S(y)=1$ wszędzie tam, gdzie
$\chi_T(x)\varphi(x+y)$ może być niezerowa.
Z Lematu~\ref{lem:iterated-distributions21} i
Lematu~\ref{lem:compact-action21} otrzymujemy
\[
 \langle W,\chi(\varphi\circ a)\rangle
 =\langle S_y,\langle T_x,\chi_T(x)\varphi(x+y)\rangle\rangle
 =\langle S,A_T\varphi\rangle.
\]
Jest to dokładnie poprzednia definicja.
Przemienność w ogólnym przypadku wynika z zamiany
zmiennych $(x,y)\mapsto(y,x)$, zamiany odcięcia
$\chi(x,y)$ na $\chi(y,x)$ oraz równości iterowanych
działań z Lematu~\ref{lem:iterated-distributions21}.
Zapewnia ona również zgodność z poprzednią definicją,
gdy tylko $F$ jest zwarty.

Zbiór $K+F$ jest domknięty. Istotnie, jeśli
$z_n=x_n+y_n\to z$ dla $x_n\in K$, $y_n\in F$,
to $L=\{z\}\cup\{z_n:n\ge1\}$ jest zwarty.
Ciąg $(x_n,y_n)$ należy do zwartego $P_L$, ma więc
podciąg zbieżny do $(x,y)\in P_L$. Z ciągłości
dodawania wynika $z=x+y\in K+F$.
Jeżeli jeden z nośników jest pusty, suma jest pusta
i także domknięta. Dla testu o nośniku rozłącznym
z $K+F$ zbiór $P_{\supp\varphi}$ jest pusty,
więc można użyć $\chi=0$. Stąd inkluzja nośników.

Pozostaje różniczkowanie. Dla $j=1,\ldots,d$
dystrybucja na $\R^{2d}$ odpowiadająca parze
$(\partial_jT,S)$ jest równa $\partial_{x_j}W$:
dla funkcji testowej $\Phi$ definicja pochodnej
i wzór na różniczkowanie iterowanego działania
z dowodu Lematu~\ref{lem:iterated-distributions21}
dają
\[
 \langle\partial_{x_j}W,\Phi\rangle
 =-\langle T_x,\langle S_y,\partial_{x_j}\Phi(x,y)\rangle\rangle
 =\langle(\partial_jT)_x,\langle S_y,\Phi(x,y)\rangle\rangle.
\]
Z Twierdzenia~\ref{thm:dist-deriv-cont20}
$\supp(\partial_jT)\subseteq K$, więc właściwość
dodawania zachodzi również dla pary $(\partial_jT,S)$
(jej przeciwobrazy zbiorów zwartych są domkniętymi
podzbiorami odpowiednich $P_L$).
Dla odcięcia $\chi$ wybranego dla $\supp\varphi$
możemy użyć tego samego odcięcia dla $\partial_j\varphi$,
której nośnik nie powiększa się, i dla pary
$(\partial_jT,S)$. Otrzymujemy
\[
\begin{aligned}
 \langle(\partial_jT)*S,\varphi\rangle
 &=-\langle W,\partial_{x_j}(\chi(\varphi\circ a))\rangle\\
 &=-\langle W,\chi((\partial_j\varphi)\circ a)\rangle
   -\langle W,(\partial_{x_j}\chi)(\varphi\circ a)\rangle\\
 &=-\langle T*S,\partial_j\varphi\rangle.
\end{aligned}
\]
Ostatni składnik w drugiej równości znika, ponieważ
$(\partial_{x_j}\chi)(\varphi\circ a)$ znika
na otoczeniu $K\times F$: na $P_{\supp\varphi}$
pochodna $\chi$ jest zerowa, a poza tym zbiorem
$\varphi\circ a$ jest lokalnie zerowa.
Zatem nośnik tego testu jest rozłączny z $\supp W$.
Lemat~\ref{lem:distribution-local19} daje więc ostatnią
równość.
Przemienność daje analogiczny wzór dla $\partial_j S$,
a iterowanie równości pierwszego rzędu dowodzi tezy
dla każdego wielowskaźnika (dla $\alpha=0$ jest to
tożsamość).
\end{proof}

\begin{przyklad}[Splot na półprostej]
\label{ex:dt-convolution-halfline}
Jeśli $K,F\subseteq[0,\infty)$ są domknięte,
to dodawanie jest właściwe na $K\times F$:
dla zwartego $L\subseteq\R$ wybieramy $R>0$ takie,
że $L\subseteq[-R,R]$. Z $x,y\ge0$ i $x+y\in L$
wynika $0\le x,y\le R$, więc $P_L$ jest domkniętym
podzbiorem $[0,R]^2$.

W szczególności dla funkcji Heaviside'a $H$
z Przykładu~\ref{ex:heaviside20} mamy
$\supp T_H=[0,\infty)$: na $(-\infty,0)$ funkcja
$H$ jest zerowa, natomiast każde otoczenie punktu
$[0,\infty)$ zawiera niepusty otwarty przedział
w $(0,\infty)$, na którym można wybrać nieujemny,
niezerowy test o dodatniej całce
(Lemat~\ref{lem:cutoffs17}).
Działanie $T_H$ na tym teście jest dodatnie.
Z powyższego kryterium istnieje $T_H*T_H$,
mimo że żaden z czynników nie ma zwartego nośnika.
Dla $\varphi\in\mathcal D(\R)$ o nośniku w $[-R,R]$
i dopuszczalnego odcięcia z poprzedniego twierdzenia
mamy, z definicji dystrybucji regularnej,
\[
 \langle T_H*T_H,\varphi\rangle
 =\int_0^\infty\int_0^\infty
 \chi(x,y)\varphi(x+y)\,\mathrm d\lambda_1(y)\,
 \mathrm d\lambda_1(x).
\]
Na pierwszej ćwiartce $\chi=1$ w każdym punkcie,
w którym $\varphi(x+y)\ne0$, więc można usunąć $\chi$.
Całka jest bezwzględnie zbieżna: funkcja podcałkowa
znika poza trójkątem $0\le x,y$, $x+y\le R$
i ma moduł nie większy niż $p_0(\varphi)$.
Podstawiając $u=x+y$, $v=x$ (moduł wyznacznika
jakobianu wynosi $1$) i stosując twierdzenie Fubiniego,
otrzymujemy
\[
 \langle T_H*T_H,\varphi\rangle
 =\int_0^\infty\int_0^u\varphi(u)\,
 \mathrm d\lambda_1(v)\,\mathrm d\lambda_1(u)
 =\int_0^\infty u\varphi(u)\,\mathrm d\lambda_1(u).
\]
Zatem $T_H*T_H=T_{x_+}$, gdzie $x_+(x)=xH(x)$.
Wynik jest zgodny z rachunkiem pochodnych:
$(T_H*T_H)'=\delta_0*T_H=T_H$, zgodnie z
Twierdzeniem~\ref{thm:dt-convolution-proper}
i Przykładami~\ref{ex:heaviside20}
oraz~\ref{ex:delta-convolution21}.
Niezależnie, pochodną $T_{x_+}$ obliczono
w Przykładzie~\ref{ex:forced-ode20}.
\end{przyklad}

\begin{uwaga}[Dlaczego jest to zagadnienie analizy funkcjonalnej]
\label{rem:convolution-tvs21}
Kluczowe kroki nie polegały na rozwiązywaniu równań
różniczkowych. Użyliśmy ciągłości operatora
$A_T$ na przestrzeni funkcji testowych, transpozycji
na jej przestrzeń dualną, przenoszenia zbiorów ograniczonych
oraz zbieżności silnej w tej przestrzeni dualnej.
Są to konstrukcje z teorii przestrzeni liniowo-topologicznych.
Dalsze zastosowania splotu do rozwiązywania równań
różniczkowych pozostawiamy odrębnemu skryptowi.
\end{uwaga}

% Materiał z AF 25; fragment do rozdziału o splocie.

\subsection{Aproksymacja w przestrzeniach Lebesgue'a}

\begin{twierdzenie}[Gęstość funkcji testowych]
\label{thm:dt-test-density}
Dla $1\le p<\infty$ funkcje $C_c^\infty(\R^d)$
są gęste w $L^p(\R^d)$.
\end{twierdzenie}
\begin{proof}
Dla $f\in L^p(\R^d)$ i $n\ge1$ połóżmy
$f_n=f\chi_{\{x:|x|\le n,\ |f(x)|\le n\}}$.
Mamy $|f_n-f|^p\le |f|^p$ i $f_n\to f$ prawie wszędzie.
Twierdzenie o zbieżności zmajoryzowanej daje
$\|f_n-f\|_p\to0$.
Każdą ograniczoną funkcję $f_n$ można aproksymować
jednostajnie (z dokładnością do zbioru miary zero)
na $\overline{B(0,n)}$ funkcjami prostymi, dzieląc
ograniczony zbiór wartości na przedziały (lub kwadraty
w przypadku zespolonym) o średnicy dążącej do zera.
Przyjmujemy wartości zerowe poza kulą. Dla takiej
funkcji prostej $s$ zachodzi
\[
 \|f_n-s\|_p\le
 \lambda_d(\overline{B(0,n)})^{1/p}\|f_n-s\|_\infty.
\]
Wystarczy zatem aproksymować funkcje charakterystyczne
ograniczonych zbiorów mierzalnych.
Dla mierzalnego ograniczonego zbioru $E$ regularność
miary Lebesgue'a daje zwarty $F\subset E$ oraz
ograniczony otwarty $O\supset E$ z
$\lambda_d(O\setminus F)<\eta$ dla dowolnego $\eta>0$.
Gdy $F\ne\varnothing$, wybieramy $r>0$ z
$2r<\operatorname{dist}(F,\R^d\setminus O)$.
Niech $\rho$ będzie nieujemną funkcją testową o nośniku
w $B(0,r/2)$ i całce $1$ (istnienie wynika z
Definicji~\ref{def:mollifier21} po odpowiednim przeskalowaniu).
Oznaczmy $F_r=\{x:\operatorname{dist}(x,F)<r\}$.
Z Twierdzenia~\ref{thm:smoothing21} funkcja
\[
 \psi(x)=\int_{\R^d}\chi_{F_r}(x-y)\rho(y)
 \,\mathrm d\lambda_d(y)
\]
jest gładką reprezentantką splotu $T_{\chi_{F_r}}*T_\rho$.
Ponieważ $0\le\chi_{F_r}\le1$ i całka $\rho$ wynosi $1$,
mamy $0\le\psi\le1$.
Dla $x\in F$ i $y\in\supp\rho$ zachodzi
$\operatorname{dist}(x-y,F)\le |y|<r/2<r$,
a więc $\psi(x)=1$.
Jeżeli $\operatorname{dist}(x,F)>3r/2$, to dla
$y\in\supp\rho$ mamy
$\operatorname{dist}(x-y,F)\ge
\operatorname{dist}(x,F)-|y|>r$.
Zatem $\psi(x)=0$ poza zbiorem
$\{x:\operatorname{dist}(x,F)\le3r/2\}$,
który jest zwarty i zawarty w $O$ dzięki wyborowi $r$.
W konsekwencji $\psi$ jest funkcją testową. Gdy $F=\varnothing$,
bierzemy $\psi=0$. W obu przypadkach
$\|\chi_E-\psi\|_p^p\le\lambda_d(O\setminus F)<\eta$.
Dla funkcji prostej $s=\sum_{j=1}^N c_j\chi_{E_j}$
i zadanej dokładności $\epsilon>0$ wybieramy testy
$\psi_j$ z $\|\chi_{E_j}-\psi_j\|_p<
\epsilon/(1+\sum_{j=1}^N|c_j|)$.
Wtedy $\sum_{j=1}^Nc_j\psi_j$ jest funkcją testową
i z nierówności trójkąta jej odległość od $s$ w $L^p$
jest mniejsza niż $\epsilon$.
Wraz z pierwszą częścią dowodu daje to tezę.
\end{proof}

\begin{twierdzenie}[Aproksymacja jedności w $L^p$]
\label{thm:dt-approx-lp}
Niech $1\le p<\infty$, $f\in L^p(\R^d)$ i niech
$\rho_\varepsilon$ będzie standardową aproksymacją
jedności z definicji \ref{def:mollifier21}.
Wtedy $\|f*\rho_\varepsilon-f\|_p\to0$.
Jeżeli $f$ ma zwarty nośnik, to $f*\rho_\varepsilon$
jest funkcją testową.
\end{twierdzenie}
\begin{proof}
Najpierw dla $v\in C_c^\infty$ mamy
$\|v(\cdot-y)-v\|_p\to0$ przy $y\to0$:
dla $|y|\le1$ różnica ma nośnik w jednym zbiorze
zwartym, a jednostajna ciągłość daje zbieżność
jednostajną. Dla dowolnego $f\in L^p$
i takiej funkcji $v$ nierówność trójkąta
oraz niezmienniczość normy względem translacji dają
\[
 \|f(\cdot-y)-f\|_p\le2\|f-v\|_p+\|v(\cdot-y)-v\|_p.
\]
Gęstość z twierdzenia \ref{thm:dt-test-density}
dowodzi ciągłości translacji w normie $L^p$.
Ponieważ $p\ge1$, nierówność Höldera na każdym
zbiorze zwartym daje $f\in L^1_{\mathrm{loc}}$.
Wybieramy mierzalną borelowską reprezentantkę $f$.
Dla każdego $x$ całka
\[
 (f*\rho_\varepsilon)(x)=
 \int_{\R^d}f(x-y)\rho_\varepsilon(y)\,\mathrm d\lambda_d(y)
\]
jest bezwzględnie zbieżna, gdyż jądro jest ograniczone
i ma zwarty nośnik. Nie zależy od reprezentantki $f$,
ponieważ translacja i odbicie zachowują zbiory miary zero.
Z Twierdzenia~\ref{thm:smoothing21}, po podstawieniu
$z=x-y$, jest to gładka reprezentantka splotu dystrybucji.
Niech $\nu_\varepsilon$ będzie miarą probabilistyczną
$\rho_\varepsilon(y)\,\mathrm d\lambda_d(y)$.
Z twierdzenia Tonellego i niezmienniczości miary
Lebesgue'a względem translacji otrzymujemy
\[
 \int_{\R^d}\int_{\R^d}|f(x-y)|^p
 \,\mathrm d\nu_\varepsilon(y)\,\mathrm d\lambda_d(x)
 =\int_{\R^d}\|f\|_p^p\,\mathrm d\nu_\varepsilon(y)
 =\|f\|_p^p.
\]
Stąd dla prawie każdego $x$ wewnętrzna całka jest
skończona i nierówność Jensena daje
$\|f*\rho_\varepsilon\|_p^p\le\|f\|_p^p$.
Stosując tę samą nierówność do $f(x-y)-f(x)$
i ponownie twierdzenie Tonellego, otrzymujemy
\[
 \|f*\rho_\varepsilon-f\|_p^p
 \le \int_{\R^d}\rho_\varepsilon(y)
            \|f(\cdot-y)-f\|_p^p\,\mathrm d\lambda_d(y)
 \le\sup_{|y|\le\varepsilon}\|f(\cdot-y)-f\|_p^p.
\]
W ostatnim oszacowaniu użyliśmy
$\supp\rho_\varepsilon\subseteq\overline{B(0,\varepsilon)}$
oraz $\nu_\varepsilon(\R^d)=1$.
Ostatnie supremum dąży do zera przez ciągłość
translacji w zerze.
W przypadku zwartego nośnika gładkość i zwartość
nośnika splotu wynikają z twierdzenia
\ref{thm:smoothing21}.
\end{proof}


\setcounter{section}{7}
% Rozdział aktywny po restrukturyzacji DT.

\section{Skoki powierzchniowe i ograniczenia rachunku dystrybucji}
\label{sec:dt-jumps-products}
\providecommand{\atlasref}[2]{#2}
Rozwijamy wzór skoku z rozdziału \ref{sec:pochodne20}
na przypadek hiperpowierzchni. Por.
\atlasref{PDE-07::thm:pde7-jump}{RRC-07: wzór skoku}.
\subsection{Pochodne przez hiperpowierzchnię}
\begin{twierdzenie}[Wzór skoku przez hiperpowierzchnię]
\label{thm:dt-hypersurface-jump}
Niech $d\ge1$, $\K\in\{\R,\C\}$, a $\Omega\subseteq\R^d$
będzie niepustym zbiorem otwartym. Niech
$S\subseteq\Omega$ będzie względnie domkniętą,
zorientowaną, zanurzoną hiperpowierzchnią klasy $C^1$
bez brzegu w $\Omega$. Zakładamy, że w każdym
punkcie $S$ istnieje otoczenie, w którym $S$ jest
wykresem funkcji klasy $C^1$ (po obrocie współrzędnych)
i rozdziela to otoczenie na dwie strony $\Omega_-$
i $\Omega_+$. Nazwy stron są zgodne z orientacją:
jednostkowa normalna $\nu$ jest skierowana
od strony minus do strony plus.
Niech $u\in C^1(\Omega\setminus S;\K)$ i niech lokalnie
przy każdym punkcie $S$ ograniczenia $u$ do obu stron
mają przedłużenia klasy $C^1$ na ich domknięcia.
Oznaczmy ślady na $S$ przez $u_-$ i $u_+$.
Dla $j=1,\ldots,d$ obie strony poniższego wzoru
należą do $\mathcal D'(\Omega;\K)$.
Wtedy, w sensie dystrybucji,
\begin{equation}\label{eq:dt-hypersurface-jump}
 \partial_jT_u=T_{(\partial_ju)_{\mathrm{reg}}}
 +(u_+-u_-)\nu_j\delta_S,
\end{equation}
gdzie $(\partial_ju)_{\mathrm{reg}}$ oznacza klasyczną
pochodną poza $S$, a dla ciągłej funkcji $b:S\to\K$ definiujemy
dystrybucję ważonej miary powierzchniowej przez
\[
 \langle b\delta_S,\varphi\rangle
 =\int_S b(s)\varphi(s)\,\mathrm d\sigma(s),
 \qquad \sigma=\mathcal H^{d-1}\!\restriction_S.
\]
Wartości funkcji $u$ i $(\partial_ju)_{\mathrm{reg}}$ na $S$
nie wpływają na ich dystrybucje regularne.
Zapis $b\delta_S$ oznacza ważoną miarę, nie mnożenie
dystrybucji przez dowolną funkcję ciągłą na $\Omega$.
\end{twierdzenie}
\begin{proof}
W lokalnych współrzędnych wykresowych hiperpowierzchnia
ma miarę Lebesgue'a zero z twierdzenia Fubiniego:
każdy przekrój w kierunku osi wykresu jest co najwyżej
jednopunktowy. Przeliczalne pokrycie takimi otoczeniami
daje $\lambda_d(S)=0$. Na każdym zwartym $K\Subset\Omega$
zbiór $K\cap S$ jest zwarty dzięki względnemu domknięciu
$S$. Można go pokryć skończenie wieloma mniejszymi
otoczeniami wykresowymi o ograniczonej podstawie
i ograniczonej pochodnej funkcji wykresu.
Wzór na miarę powierzchni wykresu
$\sqrt{1+|\nabla h|^2}\,\mathrm d\lambda_{d-1}$
daje $\sigma(K\cap S)<\infty$.
W wymiarze $d=1$ stosujemy zamiast tego miarę
liczącą: $K\cap S$ jest skończonym zbiorem punktów.
Z lokalnych przedłużeń klasy $C^1$ wynika,
że $u$ i jej pochodne klasyczne poza $S$ są
ograniczone na dostatecznie małych otoczeniach
punktów $S$ po obu stronach wykresu. Poza $S$
jest tak z ciągłości. Skończone pokrycie $K$
takimi otoczeniami daje ograniczoność na $K\setminus S$,
a stąd lokalną całkowalność $u$
i $(\partial_ju)_{\mathrm{reg}}$ po nadaniu im
wartości zerowych na $S$.
Jeśli $K\cap S\ne\varnothing$, to
\[
 |\langle b\delta_S,\varphi\rangle|
 \le \sup_{K\cap S}|b|\,\sigma(K\cap S)p_{K,0}(\varphi)
 \qquad(\varphi\in\mathcal D_K(\Omega)).
\]
Supremum jest skończone z ciągłości $b$ na
zwartym $K\cap S$. Gdy przecięcie jest puste,
funkcjonał znika na $\mathcal D_K(\Omega)$. Zatem wyraz powierzchniowy
jest liniowy i spełnia kryterium ciągłości
z Twierdzenia~\ref{thm:distribution-estimate19},
ponieważ $p_{K,0}(\varphi)=p_0(\varphi)$ dla
$\varphi\in\mathcal D_K(\Omega)$.
Ma rząd globalny zero w sensie Definicji~\ref{def:order19}.

Ustalmy $\varphi\in\mathcal D(\Omega)$.
Dla każdego punktu $S$ wybieramy najpierw
dostatecznie mały cylinder wykresowy $C$
o podstawie prostokątnej, ze zwartym domknięciem
w otoczeniu, w którym obowiązują założenia.
W lokalnych współrzędnych zapisujemy
$C=D\times(c_-,c_+)$, gdzie $D$ jest otwartym
prostokątem w $\R^{d-1}$, a wykres $x_d=h(x')$
spełnia $c_-<h(x')<c_+$ na $\overline D$.
Taki wybór jest możliwy dzięki ciągłości $h$
i zmniejszeniu podstawy wokół danego punktu.
Oznaczamy części cylindra po obu stronach
wykresu przez $C_-$ i $C_+$ zgodnie z orientacją.
Najpierw rozważmy test $\varphi$ o nośniku
zawartym w jednym z tak wybranych cylindrów.
Wzór Gaussa na obu częściach cylindra stosujemy
do pola $u\varphi\,e_j$, gdzie $e_j$ jest $j$-tym
wektorem standardowej bazy. Obie części mają
brzeg kawałkami klasy $C^1$, a pole wraz z pierwszymi
pochodnymi przedłuża się ciągle do odpowiedniego brzegu.
Dla funkcji zespolonych stosujemy wzór Gaussa osobno
do części rzeczywistej i urojonej pola.
Obrót współrzędnych nie zmienia wzoru Gaussa
dla pola zapisanego w pierwotnych współrzędnych.
W wymiarze $d=1$ cylinder zastępujemy przedziałem
zawierającym dokładnie jeden punkt $S$ i stosujemy
zwykłe całkowanie przez części na jego dwóch
podprzedziałach, z tymi samymi znakami normalnych.
Na sztucznych ścianach cylindra pole znika,
ponieważ $\varphi$ ma nośnik we wnętrzu.
Normalna zewnętrzna części minus na $S$ to $\nu$,
a części plus to $-\nu$. W konsekwencji
\[
 \begin{aligned}
 -\int_{C_-}u(x)\partial_j\varphi(x)\,\mathrm d\lambda_d(x)
 &=\int_{C_-}(\partial_ju)(x)\varphi(x)\,\mathrm d\lambda_d(x)
   -\int_{S\cap C} u_-\varphi\nu_j\,\mathrm d\sigma,\\
 -\int_{C_+}u(x)\partial_j\varphi(x)\,\mathrm d\lambda_d(x)
 &=\int_{C_+}(\partial_ju)(x)\varphi(x)\,\mathrm d\lambda_d(x)
   +\int_{S\cap C} u_+\varphi\nu_j\,\mathrm d\sigma.
 \end{aligned}
\]
Tutaj całki dotyczą części lokalnego cylindra.
Sumując, otrzymujemy~\eqref{eq:dt-hypersurface-jump} dla takiego
testu. Dla testu o nośniku rozłącznym z $S$ wzór
wynika z klasycznego całkowania przez części.

W ogólnym przypadku pokrywamy zwarty nośnik
$\varphi$ skończoną rodziną uprzednio wybranych
cylindrów wykresowych oraz otoczeń rozłącznych
z $S$. Korzystamy z gładkiego
podziału jedności podporządkowanego temu pokryciu
w otoczeniu nośnika. Dokładniej, wybieramy skończenie
wiele nieujemnych funkcji testowych $\eta_i$
o nośnikach w tych otoczeniach, których suma $A$
jest dodatnia w otoczeniu nośnika $\varphi$
(istnienie wynika z lematu~\ref{lem:cutoffs17}:
dla każdego punktu nośnika wybieramy kulę o domknięciu
zawartym w jednym otoczeniu pokrycia i nieujemną funkcję
odcięcia równą $1$ na mniejszej kuli; zwartość pozwala
wybrać skończenie wiele takich kul).
Funkcje $\varphi_i=\varphi\eta_i/A$ tam, gdzie $A>0$,
przedłużamy przez zero poza ten zbiór.
Są gładkie, gdyż nośnik $\varphi$ jest zawarty
w otwartym zbiorze $\{A>0\}$, oraz mają zwarte
nośniki w odpowiednich otoczeniach.
Mamy $\varphi=\sum_i\varphi_i$. Stosując wzór
lokalny do każdej $\varphi_i$ i sumując skończenie
wiele równości, otrzymujemy tezę dla $\varphi$.
\end{proof}

\begin{wniosek}[Skok strumienia i równanie ciągłości]
Niech $F=(F_1,\ldots,F_d)$ będzie polem rzeczywistym,
którego każda składowa spełnia założenia
twierdzenia~\ref{thm:dt-hypersurface-jump}. Wtedy
\[
 \operatorname{div}T_F
 =T_{(\operatorname{div}F)_{\mathrm{reg}}}
 +[F]\cdot\nu\,\delta_S,\qquad [F]=F_+-F_-.
\]
Równanie $\operatorname{div}T_F=0$ jest równoważne
warunkom $\operatorname{div}F=0$ poza $S$ oraz
$[F]\cdot\nu=0$ na $S$.
\end{wniosek}
\begin{proof}
Sumujemy wzór~\eqref{eq:dt-hypersurface-jump} dla składowych
$F_j$ i pochodnych $\partial_j$.
Jeśli $\operatorname{div}T_F=0$, testy o nośnikach
w $\Omega\setminus S$ dają zerowanie klasycznej
dywergencji prawie wszędzie na tym zbiorze na mocy
stwierdzenia~\ref{prop:regular-injective19}; z ciągłości
jest ona zerowa punktowo. Ponieważ $\lambda_d(S)=0$,
pozostaje $b\delta_S=0$, gdzie $b=[F]\cdot\nu$ jest
ciągłą funkcją rzeczywistą na $S$.
Gdyby $b(s_0)\ne0$, w dostatecznie małym otoczeniu
$s_0$ funkcja $b$ miałaby stały znak i moduł
ograniczony od dołu przez dodatnią stałą.
Wybierając nieujemną funkcję testową o nośniku
w tym otoczeniu, dodatnią w otoczeniu $s_0$,
otrzymalibyśmy niezerową całkę powierzchniową:
każdy niepusty otwarty fragment wykresu ma dodatnią
miarę powierzchniową (dla $d=1$ miarę liczącą).
Sprzeczność dowodzi $b=0$. Odwrotnie, gdy oba
warunki zachodzą, oba składniki wzoru znikają.
Jest to mechanizm stojący za warunkiem
Rankine'a--Hugoniota w teorii praw zachowania.
\end{proof}

\subsection{Funkcje charakterystyczne i pochodne drugiego rzędu}

\begin{wniosek}[Pochodna funkcji charakterystycznej]
\label{cor:dt-characteristic-boundary}
Niech $G\subseteq\R^d$ będzie zbiorem otwartym
o brzegu klasy $C^1$, to znaczy lokalnie brzeg jest
wykresem klasy $C^1$, a $G$ w tym otoczeniu jest
dokładnie jedną z dwóch otwartych stron wykresu.
Oznaczmy przez $n$ jednostkową normalną zewnętrzną do $G$.
Wtedy
\[
 \partial_j T_{\chi_G}=-n_j\delta_{\partial G},
 \qquad j=1,\ldots,d.
\]
W szczególności dla pola testowego
$X\in C_c^\infty(\R^d;\R^d)$
\[
 \int_G \operatorname{div}X(x)\,\mathrm d\lambda_d(x)
 =\int_{\partial G} X(s)\cdot n(s)\,
 \mathrm d\mathcal H^{d-1}(s).
\]
\end{wniosek}
\begin{proof}
Stosujemy Twierdzenie~\ref{thm:dt-hypersurface-jump}
do $S=\partial G$, wybierając stronę minus wewnątrz $G$,
stronę plus na zewnątrz i normalną $\nu=n$.
Ślady $\chi_G$ są odpowiednio $1$ i $0$,
a klasyczne pochodne poza $S$ są zerowe.
Zatem wzór skoku daje pierwszą równość.
Z definicji pochodnej dystrybucyjnej
\[
 -\int_G\partial_j X_j(x)\,\mathrm d\lambda_d(x)
 =-\int_{\partial G}n_j(s)X_j(s)\,
 \mathrm d\mathcal H^{d-1}(s).
\]
Sumując po $j$ i mnożąc przez $-1$, otrzymujemy
drugą równość. Nie wymagamy ograniczoności $G$:
wszystkie całki dotyczą zwartego nośnika pola testowego.
Jest to dystrybucyjna postać wzoru Gaussa,
a nie niezależny dowód twierdzenia Gaussa,
którego użyliśmy w dowodzie wzoru skoku.
\end{proof}

\begin{przyklad}[Kula i miara powierzchniowa]
\label{ex:dt-ball-boundary}
Dla $G=B(a,R)\subseteq\R^d$, $R>0$, normalna
zewnętrzna w punkcie $s\in\partial G$ wynosi
$n(s)=(s-a)/R$, ponieważ jest to jednostkowy
wektor radialny skierowany na zewnątrz kuli.
Z Wniosku~\ref{cor:dt-characteristic-boundary}
\[
 \langle\partial_j T_{\chi_{B(a,R)}},\varphi\rangle
 =-\frac1R\int_{\partial B(a,R)}
 (s_j-a_j)\varphi(s)\,\mathrm d\mathcal H^{d-1}(s).
\]
Dla każdego $j=1,\ldots,d$ ta dystrybucja
nie jest regularna. Istotnie, gdyby
$\partial_jT_{\chi_G}=T_f$ dla
$f\in L^1_{\mathrm{loc}}$, to z zerowania dystrybucji
poza $\partial G$ i Stwierdzenia~\ref{prop:regular-injective19}
mielibyśmy $f=0$ prawie wszędzie poza $\partial G$.
Brzeg ma miarę Lebesgue'a zero, więc $T_f=0$.
Tymczasem na małym otoczeniu punktu $a+Re_j$,
gdzie $n_j>0$, nieujemna funkcja testowa dodatnia
w tym punkcie daje ściśle ujemną wartość powyższego
parowania. Istnienie testu wynika z Lematu~\ref{lem:cutoffs17},
a dodatniość miary fragmentu brzegu z lokalnego wzoru
na miarę wykresu użytego w dowodzie wzoru skoku.
\end{przyklad}

\begin{twierdzenie}[Drugie pochodne i laplasjan przy skoku]
\label{thm:dt-second-interface}
Niech $\Omega,S,u,\nu$ spełniają założenia
Twierdzenia~\ref{thm:dt-hypersurface-jump}.
Załóżmy dodatkowo, że $u$ jest klasy $C^2$ poza $S$,
a lokalne przedłużenia do domknięć obu stron są klasy $C^2$.
Piszemy $[u]=u_+-u_-$ oraz
$[\partial_j u]=(\partial_j u)_+-(\partial_j u)_-$.
Wtedy dla $i,j=1,\ldots,d$ zachodzi
\[
 \partial_i\partial_j T_u
 =T_{(\partial_i\partial_j u)_{\mathrm{reg}}}
 +[\partial_j u]\nu_i\delta_S
 +\partial_i([u]\nu_j\delta_S).
\]
W szczególności
\[
 \Delta T_u
 =T_{(\Delta u)_{\mathrm{reg}}}
 +[\nabla u]\cdot\nu\,\delta_S
 +\operatorname{div}([u]\nu\delta_S).
\]
Ostatni składnik jest określony przez
\[
 \langle\operatorname{div}([u]\nu\delta_S),\varphi\rangle
 =-\int_S [u](s)\nu(s)\cdot\nabla\varphi(s)\,
 \mathrm d\mathcal H^{d-1}(s).
\]
Jeżeli $[u]=0$ na $S$, to
\[
 \Delta T_u=T_{(\Delta u)_{\mathrm{reg}}}
 +[\partial_\nu u]\delta_S,\qquad
 [\partial_\nu u]=[\nabla u]\cdot\nu.
\]
\end{twierdzenie}
\begin{proof}
Wszystkie klasyczne pochodne $u$ rzędów nie większych
niż dwa są lokalnie całkowalne, przez ten sam argument
z lokalnymi przedłużeniami i zwartymi zbiorami co w
Twierdzeniu~\ref{thm:dt-hypersurface-jump}.
Ślady $[u]$ i $[\partial_j u]$ są ciągłe na $S$,
podobnie jak $\nu$. Zatem ich iloczyny z odpowiednimi
składowymi normalnej definiują dystrybucje powierzchniowe
rzędu zerowego przez wzór z tego twierdzenia.
Stosujemy $\partial_i$ do równości
\[
 \partial_j T_u=T_{(\partial_j u)_{\mathrm{reg}}}
 +[u]\nu_j\delta_S.
\]
Do pierwszego składnika stosujemy ponownie wzór skoku,
tym razem dla funkcji $(\partial_j u)_{\mathrm{reg}}$,
która spełnia założenia klasy $C^1$ dzięki założeniu
klasy $C^2$ dla $u$. Otrzymujemy wzór na drugą pochodną.
Sumując go dla $i=j$, dostajemy wzór na laplasjan,
gdzie dywergencja dystrybucji wektorowej oznacza sumę
pochodnych jej składowych. Definicja pochodnej
dystrybucyjnej daje podany wzór na parowanie z testem.
Jeśli $[u]=0$, dystrybucja wektorowa $[u]\nu\delta_S$
jest zerowa, podobnie jak jej dywergencja.
\end{proof}

\begin{przyklad}[Ciągły profil z załamaniem i skok funkcji]
\label{ex:dt-flat-interface}
Niech $S=\{x\in\R^d:x_d=0\}$, $\nu=e_d$,
a $d\ge2$. Dla $u(x)=|x_d|$ ślady $u$ są równe zero,
klasyczny laplasjan poza $S$ jest zerowy, a
$[\nabla u]=2e_d$. Z Twierdzenia~\ref{thm:dt-second-interface}
wynika
\[
 \Delta T_{|x_d|}=2\delta_S,\qquad
 \langle\delta_S,\varphi\rangle
 =\int_{\R^{d-1}}\varphi(x',0)\,\mathrm d\lambda_{d-1}(x').
\]
Dla $v(x)=H(x_d)$ mamy natomiast $[v]=1$,
$[\nabla v]=0$ i klasyczny laplasjan równy zero poza $S$.
Zatem
\[
 \Delta T_v=\partial_d\delta_S,\qquad
 \langle\Delta T_v,\varphi\rangle
 =-\int_{\R^{d-1}}(\partial_d\varphi)(x',0)\,
 \mathrm d\lambda_{d-1}(x').
\]
Pierwszy przykład daje miarę powierzchniową, drugi
jej pochodną w kierunku normalnym. Ta druga dystrybucja
nie jest rzędu zerowego: wybierzmy
$\eta\in\mathcal D(\R^{d-1})$ o całce $1$ i
$\theta\in\mathcal D(\R)$ z $\theta'(0)=1$
(np. $\theta(t)=t\zeta(t)$, gdzie
$\zeta\in\mathcal D(\R;\R)$ jest równa $1$ w otoczeniu $0$).
Funkcję $\eta$ otrzymujemy przez normalizację nieujemnego
niezerowego testu, a $\zeta$ z Lematu~\ref{lem:cutoffs17}.
Dla $\varphi_\varepsilon(x',x_d)
=\eta(x')\varepsilon\theta(x_d/\varepsilon)$,
$0<\varepsilon\le1$, nośniki są zawarte w jednym
zbiorze zwartym, a $p_0(\varphi_\varepsilon)\to0$.
Jednocześnie
$\langle\partial_d\delta_S,\varphi_\varepsilon\rangle=-1$.
To przeczy oszacowaniu rzędu zerowego na tym zbiorze.
Z drugiej strony dla zwartego $K\subseteq\R^d$
i $\varphi\in\mathcal D_K(\R^d)$ zachodzi
\[
 |\langle\partial_d\delta_S,\varphi\rangle|
 \le\lambda_{d-1}(\pi(K))p_1(\varphi),
 \qquad \pi(x',x_d)=x'.
\]
Istotnie, całkowana pochodna znika poza zwartym
rzutem $\pi(K)$ i jej moduł nie przekracza $p_1(\varphi)$.
Miara tego rzutu jest skończona.
Z Definicji~\ref{def:order19} rząd globalny wynosi
więc dokładnie $1$.
\end{przyklad}

\subsection{Przeszkoda Schwartza i kwadrat delty}

Mnożenie dystrybucji przez funkcję gładką jest
określone w definicji \ref{def:multiply20}.
Nie istnieje jednak rozszerzenie tej operacji
na wszystkie dystrybucje z pełnym zestawem
klasycznych własności iloczynu. Poniższe argumenty odróżniają dwa rodzaje przeszkód
algebraicznych od niepowodzenia aproksymacji.
\begin{stwierdzenie}[Przeszkoda w rozszerzeniu iloczynu]
\label{prop:dt-schwartz-obstruction}
Na $\mathcal D'(\R)$ nie istnieje przemienny, łączny
i dwuliniowy iloczyn $\star$ spełniający jednocześnie
następujące warunki:
\begin{enumerate}
\item $T_f\star T_g=T_{fg}$ dla wszystkich
$f,g\in L^\infty_{\mathrm{loc}}(\R)$;
\item $(S\star T)'=S'\star T+S\star T'$ dla wszystkich
$S,T\in\mathcal D'(\R)$.
\end{enumerate}
\end{stwierdzenie}
\begin{proof}
Niech $A=T_H$, gdzie $H$ jest funkcją Heaviside'a.
Ponieważ $H^2=H$ prawie wszędzie, pierwszy warunek
daje $A\star A=A$. Z przykładu~\ref{ex:heaviside20}
mamy $A'=\delta_0$. Różniczkując równość
$A\star A=A$ za pomocą drugiego warunku i korzystając
z przemienności, otrzymujemy
\[
 2(A\star\delta_0)=\delta_0.
\]
Mnożymy tę równość przez $A$. Dwuliniowość i łączność
dają
\[
 A\star\delta_0
 =2A\star(A\star\delta_0)
 =2(A\star A)\star\delta_0
 =2A\star\delta_0.
\]
Stąd $A\star\delta_0=0$, a pierwsza równość
wymusza $\delta_0=0$. Jest to sprzeczne
z działaniem delty na test o wartości $1$ w zerze.
\end{proof}

\begin{stwierdzenie}[Przeszkoda wynikająca z mnożenia przez funkcje gładkie]
\label{prop:dt-smooth-product-obstruction}
Nie istnieje przemienny, łączny, dwuliniowy iloczyn
$\star$ na $\mathcal D'(\R)$ taki, że
\[
 T_g\star T=gT
 \qquad(g\in C^\infty(\R;\K),\ T\in\mathcal D'(\R;\K)),
\]
gdzie prawa strona oznacza mnożenie z
Definicji~\ref{def:multiply20}.
Nie jest tu potrzebna reguła Leibniza.
\end{stwierdzenie}
\begin{proof}
Niech $P=\operatorname{pv}(1/x)$ będzie dystrybucją
wartości głównej z Przykładu~\ref{ex:principal-value19}.
Dla $\varphi\in\mathcal D(\R)$ z definicji mnożenia
i wartości głównej mamy
\[
 \langle xP,\varphi\rangle
 =\lim_{\varepsilon\downarrow0}
 \int_{|x|>\varepsilon}\varphi(x)\,\mathrm d\lambda_1(x)
 =\int_\R\varphi(x)\,\mathrm d\lambda_1(x).
\]
Ostatnia równość wynika z oszacowania modułu całki
po $[-\varepsilon,\varepsilon]$ przez
$2\varepsilon p_0(\varphi)$. Zatem $xP=T_1$.
Jednocześnie $x\delta_0=0$, ponieważ
$\langle\delta_0,x\varphi\rangle=0$.
Z żądanej zgodności z mnożeniem przez funkcje gładkie
$T_1$ jest elementem neutralnym, $T_x\star P=T_1$
oraz $T_x\star\delta_0=0$. Stąd
\[
 \delta_0=T_1\star\delta_0
 =(T_x\star P)\star\delta_0
 =P\star(T_x\star\delta_0)=0.
\]
Ostatni krok wynika z dwuliniowości, a przegrupowanie
z przemienności i łączności. Sprzeczność z
$\delta_0\ne0$ dowodzi tezy.
\end{proof}

Powyższe dwa stwierdzenia dotyczą różnych założeń
o zgodności iloczynu z dotychczasowym rachunkiem:
pierwsze wymaga zachowania iloczynu funkcji lokalnie
ograniczonych i reguły Leibniza, drugie zachowania
mnożenia każdej dystrybucji przez funkcję gładką.
Nie wykluczają one częściowo określonych iloczynów
ani osadzenia dystrybucji w większej algebrze,
w której rezygnuje się z odpowiednich wymagań zgodności.

\begin{przyklad}[Skąd bierze się problem z $\delta^2$?]
\label{ex:dt-delta-square}
Niech $\rho\in\mathcal D(\R;\R)$ będzie nieujemna,
$\supp\rho\subseteq[-1,1]$ i
$\int_\R\rho\,\mathrm d\lambda_1=1$.
Dla $\varepsilon>0$ określamy
$\rho_\varepsilon(x)=\varepsilon^{-1}\rho(x/\varepsilon)$.
Z twierdzenia~\ref{thm:strong-mollification21}
(dla standardowej aproksymacji jedności) wynika
$T_{\rho_\varepsilon}=\delta_0*T_{\rho_\varepsilon}
\to\delta_0$ silnie w $\mathcal D'(\R)$;
równość wynika z Przykładu~\ref{ex:delta-convolution21}.
Bezpośrednio zbieżność słabą-* widać także stąd, że
dla testu $\psi$ mamy
$\langle T_{\rho_\varepsilon},\psi\rangle=
\int\rho(y)\psi(\varepsilon y)\,\mathrm d\lambda_1(y)
\to\psi(0)$ przez jednostajną zbieżność na nośniku $\rho$.
Wybieramy rzeczywisty, nieujemny test $\varphi$
taki, że $\varphi(0)>0$.
Podstawienie $x=\varepsilon y$ daje
\[
 \varepsilon\langle T_{\rho_\varepsilon^2},\varphi\rangle
 =\int_\R \rho(y)^2\varphi(\varepsilon y)\,
 \mathrm d\lambda_1(y).
\]
Ponieważ nośnik $\rho$ jest zwarty,
$\varphi(\varepsilon y)\to\varphi(0)$ jednostajnie
na nim. Po odjęciu liczby
\[
 c=\varphi(0)\int_\R\rho(y)^2\,\mathrm d\lambda_1(y)
\]
moduł różnicy jest ograniczony przez
\[
 \left(\int_\R\rho(y)^2\,\mathrm d\lambda_1(y)\right)
 \sup_{y\in\supp\rho}|\varphi(\varepsilon y)-\varphi(0)|.
\]
Pierwszy czynnik jest skończony, a drugi dąży
do zera. Zatem prawa strona dąży do $c>0$.
Dodatniość wynika z $\varphi(0)>0$ oraz z faktu,
że ciągła funkcja $\rho$ o całce $1$ nie jest zerowa.
Dla dostatecznie małego $\varepsilon$ otrzymujemy
zatem
$\langle T_{\rho_\varepsilon^2},\varphi\rangle
\ge c/(2\varepsilon)\to+\infty$.
Kwadraty tej aproksymacji delty nie mają więc
granicy w $\mathcal D'(\R)$, ponieważ zbieżność
dystrybucyjna wymaga skończonej granicy działania
na każdym ustalonym teście.
Jest to przeszkoda dla definiowania iloczynu
przez zwykłe przejście do granicy gładkich
aproksymacji, odrębna od powyższego argumentu
algebraicznego.
\end{przyklad}

\setcounter{part}{1}
\part{Transformacja Fouriera}
\setcounter{section}{8}
% Rozdział aktywny po restrukturyzacji DT.
\section{Transformacja Fouriera w przestrzeniach Schwartza i funkcji całkowalnych}
\providecommand{\atlasref}[2]{#2}

\subsection{Transformacja Fouriera w przestrzeni Schwartza}

\noindent\emph{Powiązania tematyczne w Bibliotece matematycznej:} \atlasref{PDE-08::thm:pde8-inversion}{PDE-08: twierdzenie o odwracaniu transformacji Fouriera}.


W całym rozdziale $d\ge1$ jest liczbą całkowitą,
pracujemy nad $\C$ i przyjmujemy stałą
$c_d=(2\pi)^{-d/2}$. Piszemy
$\mathcal S=\mathcal S(\R^d;\C)$ i
$L^p=L^p(\R^d;\C)$, o ile nie wskazano inaczej.
Korzystamy z seminorm $s_{\alpha,\beta}$ i $Q_{N,m}$
przestrzeni Schwartza z rozdziału~\ref{sec:gladkie-schwartz17}.
Wzór z całką jest zwykłą całką Lebesgue'a,
bez potrzeby odwoływania się do dystrybucji.

\begin{definicja}[Transformacja Fouriera]
\label{def:fourier-schwartz17}
Dla $f\in\mathcal S(\R^d;\C)$ definiujemy
\[
\widehat f(\xi)=\mathcal Ff(\xi)=
(2\pi)^{-d/2}\int_{\R^d}e^{-ix\cdot\xi}f(x)\,
\mathrm d\lambda_d(x),\qquad\xi\in\R^d.
\]
Całka jest bezwzględnie zbieżna, gdyż $\mathcal S\subseteq L^1$
na mocy Stwierdzenia~\ref{prop:schwartz-operators17}.
\end{definicja}

\begin{twierdzenie}[Transformacja Fouriera zachowuje $\mathcal S$]
\label{thm:fourier-schwartz17}
Dla każdego $d\ge1$ transformacja Fouriera jest ciągłym
izomorfizmem liniowo-topologicznym
$\mathcal F:\mathcal S(\R^d;\C)\to\mathcal S(\R^d;\C)$.
Jej odwrotność wyraża wzór
\[
\mathcal F^{-1}g(x)=(2\pi)^{-d/2}
\int_{\R^d}e^{ix\cdot\xi}g(\xi)\,
\mathrm d\lambda_d(\xi).
\]
Dla każdego $f\in\mathcal S$ oraz $j=1,\dots,d$ zachodzą
$\mathcal F(\partial_jf)(\xi)=i\xi_j\widehat f(\xi)$
i $\mathcal F(x_jf)(\xi)=i\partial_{\xi_j}\widehat f(\xi)$.
\end{twierdzenie}
\begin{proof}
Przy różniczkowaniu względem $\xi_j$ iloraz różnicowy
czynnika wykładniczego jest ograniczony przez $|x_j|$.
Indukcyjnie majorantą przy kolejnych różniczkowaniach
jest $|x^{\beta+e_j}f(x)|$, całkowalna na mocy
Stwierdzenia~\ref{prop:schwartz-operators17}.
Twierdzenie o zbieżności zmajoryzowanej daje więc
$\partial_\xi^\beta\widehat f=(-i)^{|\beta|}
\mathcal F(x^\beta f)$, a także ciągłość tych pochodnych.
Dla pochodnej $\partial_j f$ stosujemy twierdzenie Fubiniego
i całkowanie przez części względem $x_j$ na $[-R,R]$.
Przy ustalonych pozostałych współrzędnych wyrazy brzegowe
dążą do zera, gdyż $f$ szybko maleje; całki jednowymiarowe
są bezwzględnie zbieżne. Następnie całkujemy po pozostałych
współrzędnych, korzystając z całkowalności $f$ i $\partial_j f$.
Otrzymujemy $\mathcal F(\partial_j f)=i\xi_j\widehat f$.
Iteracja daje $\mathcal F(\partial^\alpha f)=(i\xi)^\alpha\widehat f$.
Połączenie tych tożsamości prowadzi do
\[
\xi^\alpha\partial_\xi^\beta\widehat f(\xi)
=(-i)^{|\alpha|+|\beta|}
\mathcal F\bigl(\partial^\alpha(x^\beta f)\bigr)(\xi).
\]
Wartość bezwzględną prawej strony ograniczamy przez
$(2\pi)^{-d/2}\|\partial^\alpha(x^\beta f)\|_{L^1}$.
Z oszacowań w Stwierdzeniu~\ref{prop:schwartz-operators17}
otrzymujemy, dla pewnej stałej $C_{\alpha,\beta,d}$,
$s_{\alpha,\beta}(\widehat f)\le
C_{\alpha,\beta,d}Q_{d+1+|\beta|,|\alpha|}(f)$,
stosując kolejno oszacowanie normy $L^1$, różniczkowania
i mnożenia przez wielomian. Dokładniej,
\[
 \|\partial^\alpha(x^\beta f)\|_1
 \le C_d Q_{d+1,0}(\partial^\alpha(x^\beta f))
 \le C_d Q_{d+1,|\alpha|}(x^\beta f)
 \le C_{\alpha,\beta,d}'Q_{d+1+|\beta|,|\alpha|}(f).
\]
Tutaj $C_d=\int_{\R^d}(1+|x|)^{-d-1}\,\mathrm d\lambda_d(x)$
jest skończona na mocy tego samego stwierdzenia.
Wszystkie seminormy $s_{\alpha,\beta}(\widehat f)$ są więc
skończone i zależą w sposób ciągły od $f$.

Pozostaje udowodnić wzór odwrotny. Dla $\varepsilon>0$ definiujemy
\[
I_\varepsilon(x)=(2\pi)^{-d/2}\int_{\R^d}
 e^{-\varepsilon|\xi|^2/2}e^{ix\cdot\xi}
 \widehat f(\xi)\,\mathrm d\lambda_d(\xi).
\]
Ponieważ $f\in L^1$ i czynnik Gaussa należy do $L^1$,
możemy zastosować twierdzenie Fubiniego. Jednowymiarową całkę Gaussa obliczamy następująco:
$I(u)=\int_\R e^{-\varepsilon t^2/2}e^{iut}\,\mathrm d\lambda_1(t)$
można różniczkować pod całką z majorantą
$|t|e^{-\varepsilon t^2/2}$. Ponieważ
$(e^{-\varepsilon t^2/2})'=-\varepsilon t e^{-\varepsilon t^2/2}$,
całkowanie przez części, z zanikającymi wyrazami brzegowymi, daje
\[
I'(u)=-\frac{i}{\varepsilon}
\int_\R (e^{-\varepsilon t^2/2})'e^{iut}\,\mathrm d\lambda_1(t)
=-\frac u\varepsilon I(u).
\]
Ponadto $I(0)=\sqrt{2\pi/\varepsilon}$ z całki Gaussa
(po podstawieniu $s=\sqrt\varepsilon t$).
Dla przypomnienia, jeżeli
$J=\int_\R e^{-s^2/2}\,\mathrm d\lambda_1(s)$,
to $0<J<\infty$ (dla $|s|\ge1$ porównujemy funkcję
podcałkową z $e^{-|s|/2}$).
Z twierdzenia Tonellego i podstawienia biegunowego
\[
 J^2=\int_{\R^2} e^{-(s^2+t^2)/2}\,\mathrm d\lambda_2(s,t)
 =2\pi\int_0^\infty re^{-r^2/2}\,\mathrm d\lambda_1(r)
 =2\pi.
\]
Ostatnią całkę obliczamy podstawieniem $v=r^2/2$.
Z równania różniczkowego dla $I$ wynika
$(e^{u^2/(2\varepsilon)}I(u))'=0$, zatem
$I(u)=\sqrt{2\pi/\varepsilon}e^{-u^2/(2\varepsilon)}$.
Zastosowanie tego wzoru w każdej współrzędnej daje
\[
I_\varepsilon(x)=\int_{\R^d}f(y)
(2\pi\varepsilon)^{-d/2}e^{-|x-y|^2/(2\varepsilon)}
\,\mathrm d\lambda_d(y).
\]
Oznaczmy $G_\varepsilon(z)=(2\pi\varepsilon)^{-d/2}
e^{-|z|^2/(2\varepsilon)}$. Podstawienie $z=\sqrt\varepsilon w$
daje $\int G_\varepsilon\,\mathrm d\lambda_d=1$ oraz
$\int_{|z|\ge\delta}G_\varepsilon(z)\,\mathrm d\lambda_d(z)
=\int_{|w|\ge\delta/\sqrt\varepsilon}G_1(w)\,\mathrm d\lambda_d(w)
\to0$ dla każdego $\delta>0$, z całkowalności $G_1$.
Dla ustalonego $x$ i $\eta>0$ wybieramy $\delta>0$ tak,
aby $|f(x-z)-f(x)|<\eta$ dla $|z|<\delta$.
Wtedy
\[
|I_\varepsilon(x)-f(x)|\le
\eta+2\|f\|_\infty
\int_{|z|\ge\delta}G_\varepsilon(z)\,\mathrm d\lambda_d(z).
\]
Stąd granica górna błędu nie przekracza $\eta$, a ponieważ
$\eta$ jest dowolne, $I_\varepsilon(x)\to f(x)$.
Z drugiej strony $\widehat f\in L^1$, zatem
z twierdzenia o zbieżności zmajoryzowanej
$I_\varepsilon(x)$ dąży do całki ze wzoru odwrotnego.
Otrzymujemy $\mathcal F^{-1}\mathcal Ff=f$.
Ten sam argument, ze zmienionym znakiem wykładnika,
daje $\mathcal F\mathcal F^{-1}g=g$ dla $g\in\mathcal S$.
Ciągłość operatora odwrotnego wynika z identycznych
oszacowań seminorm po zmianie znaku wykładnika.
\end{proof}

\begin{wniosek}[Odbicie i iteracje transformacji Fouriera]
\label{cor:dt-fourier-iteration}
Niech $R:\mathcal S(\R^d;\C)\to\mathcal S(\R^d;\C)$
będzie operatorem odbicia $(Rf)(x)=f(-x)$.
Wtedy
\[
 \mathcal F^{-1}=R\mathcal F=\mathcal F R,\qquad
 \mathcal F^2=R,\qquad \mathcal F^4=I.
\]
Ponadto dla $f\in\mathcal S$ zachodzi
\[
 \widehat{\overline f}(\xi)=\overline{\widehat f(-\xi)}.
\]
\end{wniosek}
\begin{proof}
Odbicie jest ciągłym automorfizmem $\mathcal S$,
ponieważ zachowuje wszystkie seminormy $Q_{N,m}$
i spełnia $R^2=I$. W istocie
$\partial^\alpha(Rf)(x)=(-1)^{|\alpha|}(\partial^\alpha f)(-x)$
i $|-x|=|x|$.
Z wzoru odwrotnego w Twierdzeniu~\ref{thm:fourier-schwartz17}
otrzymujemy $\mathcal F^{-1}=R\mathcal F$.
Podstawienie $y=-x$ w całce Fourierowskiej daje
$\mathcal F R=R\mathcal F$, ponieważ miara Lebesgue'a
jest niezmiennicza względem odbicia. Zatem
$I=\mathcal F^{-1}\mathcal F=R\mathcal F^2$,
a mnożenie przez $R$ daje $\mathcal F^2=R$.
Stąd $\mathcal F^4=R^2=I$.
Sprzężenie zespolone przenosimy pod znak bezwzględnie
zbieżnej całki, co daje ostatnią tożsamość.
\end{proof}

\begin{przyklad}[Funkcja Gaussa i wizualizacja w Plotly]
\label{ex:fourier-gauss17}
Dla $a>0$ oraz $f_a(x)=e^{-ax^2/2}$, $x\in\R$,
stosując obliczoną w poprzednim dowodzie całkę $I(u)$
z $\varepsilon=a$ i $u=-\xi$, otrzymujemy
$\widehat f_a(\xi)=a^{-1/2}e^{-\xi^2/(2a)}$.
W szczególności zwiększenie skupienia funkcji $f_a$
odpowiada rozszerzeniu jej transformaty Fouriera.
Poniższy kod Python rysuje obie funkcje z użyciem biblioteki Plotly.
\begin{lstlisting}[language=Python,caption={Para funkcji Gaussa: funkcja i transformata Fouriera}]
import numpy as np
import plotly.graph_objects as go

a = 2.0
x = np.linspace(-5.0, 5.0, 1000)
f = np.exp(-a * x**2 / 2)
f_hat = a**(-0.5) * np.exp(-x**2 / (2 * a))

fig = go.Figure()
fig.add_scatter(x=x, y=f, mode="lines", name="f(x)")
fig.add_scatter(x=x, y=f_hat, mode="lines", name="Fourier(f)(x)")
fig.update_layout(title="Para Gaussa", xaxis_title="x / xi",
                  yaxis_title="wartosc")
fig.show()
\end{lstlisting}
\end{przyklad}

\begin{stwierdzenie}[Splot funkcji szybko malejących]
\label{prop:convolution-schwartz17}
Dla $f,g\in\mathcal S(\R^d;\C)$ definiujemy
$(f*g)(x)=\int_{\R^d}f(y)g(x-y)\,\mathrm d\lambda_d(y)$.
Wówczas $f*g\in\mathcal S$, odwzorowanie
$(f,g)\mapsto f*g$ jest dwuliniowe i łącznie ciągłe
na $\mathcal S\times\mathcal S$, a ponadto
$\mathcal F(f*g)=(2\pi)^{d/2}\widehat f\,\widehat g$.
\end{stwierdzenie}
\begin{proof}
Z $\mathcal S\subset L^1\cap L^\infty$ wynika
bezwzględna zbieżność splotu, gdyż całka modułu
nie przekracza $\|f\|_1\|g\|_\infty$.
Ilorazy różnicowe pochodnych $g$ względem $x_j$
są ograniczone przez $\|\partial^{\beta+e_j}g\|_\infty$,
więc po pomnożeniu przez $|f(y)|$ otrzymujemy majorantę
całkowalną względem $y$. Indukcja i twierdzenie o zbieżności
zmajoryzowanej dają $\partial^\beta(f*g)=f*(\partial^\beta g)$
i ciągłość wszystkich tych pochodnych.
Nierówność $1+|x|\le(1+|y|)(1+|x-y|)$ prowadzi do
\[
 Q_{N,m}(f*g)\le Q_{N,m}(g)
 \int_{\R^d}(1+|y|)^N|f(y)|\,\mathrm d\lambda_d(y)
 \le C_{d,N}Q_{N+d+1,0}(f)Q_{N,m}(g).
\]
To dowodzi jednocześnie szybkiego zanikania i ciągłości
splotu. Dla jasności, przy ustalonych $f_0,g_0\in\mathcal S$
rozkładamy różnicę
\[
 f*g-f_0*g_0=(f-f_0)*(g-g_0)
 +(f-f_0)*g_0+f_0*(g-g_0).
\]
Powyższe oszacowanie zastosowane do trzech składników
pokazuje, że dla każdego $N,m$ seminorma różnicy
dąży do zera przy $(f,g)\to(f_0,g_0)$ w topologii
produktowej $\mathcal S\times\mathcal S$.
Twierdzenie Tonellego i niezmienniczość miary Lebesgue'a
względem przesunięć dają
$\int_{\R^d}\int_{\R^d}|f(y)g(x-y)|\,\mathrm d\lambda_d(x)
\,\mathrm d\lambda_d(y)=\|f\|_1\|g\|_1<\infty$.
Możemy więc zastosować twierdzenie Fubiniego i podstawić $z=x-y$:
\[
\begin{aligned}
\mathcal F(f*g)(\xi)
&=(2\pi)^{-d/2}
\int_{\R^d} f(y)e^{-iy\cdot\xi}\,\mathrm d\lambda_d(y)
\int_{\R^d}g(z)e^{-iz\cdot\xi}\,\mathrm d\lambda_d(z)\\
&=(2\pi)^{d/2}\widehat f(\xi)\widehat g(\xi).
\end{aligned}
\]
\end{proof}

\begin{stwierdzenie}[Transformata iloczynu i splotu]
\label{prop:dt-fourier-product}
Dla $f,g\in\mathcal S(\R^d;\C)$ zachodzą
\[
 \mathcal F(fg)=c_d(\widehat f*\widehat g),\qquad
 \mathcal F^{-1}(fg)=c_d(\mathcal F^{-1}f*\mathcal F^{-1}g),
 \qquad c_d=(2\pi)^{-d/2}.
\]
W szczególności transformacja Fouriera zamienia
splot na mnożenie punktowe z czynnikiem $c_d^{-1}$,
a mnożenie punktowe na splot z czynnikiem $c_d$.
\end{stwierdzenie}
\begin{proof}
Reguła Leibniza daje
$Q_{N,m}(fg)\le 2^m Q_{N,m}(f)Q_{0,m}(g)$,
więc $fg\in\mathcal S$. Rzeczywiście, w każdym składniku
pochodnej iloczynu rzędu co najwyżej $m$ przypisujemy
całą wagę $(1+|x|)^N$ pochodnej $f$, a suma
współczynników dwumianowych dla wielowskaźnika
$\alpha$ wynosi $2^{|\alpha|}\le2^m$.
Wstawiamy wzór odwrotny dla $g$ z
Twierdzenia~\ref{thm:fourier-schwartz17}:
\[
\begin{aligned}
 \widehat{fg}(\xi)
 &=c_d^2\int_{\R^d}\int_{\R^d}
 f(x)\widehat g(\eta)e^{-ix\cdot(\xi-\eta)}
 \,\mathrm d\lambda_d(\eta)\,\mathrm d\lambda_d(x)\\
 &=c_d\int_{\R^d}\widehat f(\xi-\eta)\widehat g(\eta)
 \,\mathrm d\lambda_d(\eta).
\end{aligned}
\]
Zamiana kolejności całkowania jest dozwolona na mocy
twierdzenia Fubiniego, ponieważ całka modułu bez czynnika
$c_d^2$ wynosi $\|f\|_1\|\widehat g\|_1<\infty$.
Druga równość wynika z pierwszej i Wniosku~\ref{cor:dt-fourier-iteration}:
odbicie spełnia $R(u*v)=(Ru)*(Rv)$ przez podstawienie
$y=-z$ w całce definiującej splot.
Regułę transformowania splotu udowodniono w
Stwierdzeniu~\ref{prop:convolution-schwartz17}.
\end{proof}

\begin{przyklad}[Gauss wielowymiarowy i zmiana skali]
\label{ex:dt-fourier-gaussian-matrix}
Niech $A$ będzie rzeczywistą symetryczną macierzą
dodatnio określoną rozmiaru $d\times d$.
Dla $f_A(x)=\exp(-x^TAx/2)$ mamy
\[
 \widehat f_A(\xi)=(\det A)^{-1/2}
 \exp(-\xi^TA^{-1}\xi/2).
\]
Istotnie, z twierdzenia spektralnego dla macierzy
symetrycznych istnieje macierz ortogonalna $U$
oraz $a_1,\ldots,a_d>0$ takie, że
$A=U^T\operatorname{diag}(a_1,\ldots,a_d)U$.
W szczególności $x^TAx\ge(\min_j a_j)|x|^2$,
więc $f_A\in\mathcal S$: każda jej pochodna jest
iloczynem wielomianu i $f_A$, co wynika indukcyjnie
z reguły iloczynu.
Podstawienie $y=Ux$ zachowuje miarę Lebesgue'a,
a $x\cdot\xi=y\cdot(U\xi)$.
Twierdzenie Fubiniego pozwala rozdzielić całkę na
$d$ jednowymiarowych całek (ich moduły są całkowalne).
Z Przykładu~\ref{ex:fourier-gauss17} otrzymujemy
\[
 \widehat f_A(\xi)
 =\prod_{j=1}^d a_j^{-1/2}
 \exp\left(-\frac{(U\xi)_j^2}{2a_j}\right).
\]
Ponieważ $\det A=\prod_j a_j$ i
$A^{-1}=U^T\operatorname{diag}(a_1^{-1},\ldots,a_d^{-1})U$,
dostajemy tezę. W szczególności kierunki silniejszego
skupienia funkcji odpowiadają kierunkom słabszego
zanikania transformaty.
\end{przyklad}

\subsection{Transformacja całkowa w przestrzeni \texorpdfstring{$L^1$}{L1}}

Przyjmujemy w całym skrypcie normalizację unitarną z
Definicji~\ref{def:fourier-schwartz17}. Rozróżniamy
transformację (operator) i transformatę (jego wartość).
W $L^1$ całka określa funkcję w każdym punkcie, natomiast
w $L^2$ twierdzenie Plancherela określa w pierwszej
kolejności klasę funkcji równych prawie wszędzie.
Dla $f\in L^1(\R^d)$ definiujemy
\[
 \widehat f(\xi)=c_d\int_{\R^d}f(x)e^{-ix\cdot\xi}
 \,\mathrm d\lambda_d(x),\qquad c_d=(2\pi)^{-d/2}.
\]
Całka jest bezwzględnie zbieżna dla każdego $\xi$.
Jeżeli dwie reprezentantki klasy $f$ są równe prawie
wszędzie, to ich iloczyny z wykładnikiem również
są równe prawie wszędzie dla każdego ustalonego $\xi$.
Transformata jest więc dobrze określona na klasach
funkcji w $L^1$, a dla $f\in\mathcal S$ definicje
pokrywają się.
Intuicyjnie jest to pomiar amplitudy i fazy poszczególnych
częstotliwości; przy dużej częstotliwości wkłady z różnych
punktów mogą się znosić, nawet gdy moduł jądra jest równy
jeden. To zjawisko opisuje następne twierdzenie.

\begin{twierdzenie}[Lemat Riemanna--Lebesgue'a]
\label{thm:dt-fourier-l1-rl}
Transformacja całkowa jest ograniczonym operatorem
liniowym z $L^1(\R^d)$ do $C_0(\R^d)$, gdzie $C_0$
oznacza przestrzeń ciągłych funkcji zanikających w
nieskończoności, z normą supremum. Zachodzi
$\|\widehat f\|_\infty\le c_d\|f\|_1$.
\end{twierdzenie}
\begin{proof}
Liniowość i oszacowanie wynikają bezpośrednio z całki.
Ciągłość jest konsekwencją twierdzenia o zbieżności
zmajoryzowanej z majorantą $|f|$.
Aby wykazać zanikanie, dla $h\in\R^d$ oznaczmy
$\tau_h f(x)=f(x-h)$. Z niezmienniczości miary
Lebesgue'a względem translacji
$\widehat{\tau_h f}(\xi)=e^{-ih\cdot\xi}\widehat f(\xi)$.
Dla $\xi\ne0$ wybieramy $h=\pi\xi/|\xi|^2$.
Wtedy $e^{-ih\cdot\xi}=-1$ oraz
\[
 2|\widehat f(\xi)|
 =|\widehat{f-\tau_h f}(\xi)|
 \le c_d\|f-\tau_h f\|_1.
\]
Gdy $|\xi|\to\infty$, mamy $h\to0$, a prawa strona
dąży do zera na mocy ciągłości translacji w $L^1$,
wykazanej w dowodzie Twierdzenia~\ref{thm:dt-approx-lp}.
Por. także
\atlasref{ME-14::thm:translation-lp-14}{ciągłość translacji w przestrzeniach Lebesgue'a}.
\end{proof}

\begin{wniosek}[Jednostajna ciągłość transformaty]
\label{cor:dt-fourier-uniform-continuity}
Dla $f\in L^1(\R^d)$ i $h\in\R^d$ zachodzi
\[
 \sup_{\xi\in\R^d}|\widehat f(\xi+h)-\widehat f(\xi)|
 \le c_d\int_{\R^d}|f(x)|\,|e^{-ix\cdot h}-1|
 \,\mathrm d\lambda_d(x).
\]
W szczególności $\widehat f$ jest jednostajnie ciągła.
\end{wniosek}
\begin{proof}
Odejmujemy od siebie całki definiujące transformatę
w punktach $\xi+h$ i $\xi$, wyłączając wspólny
czynnik $e^{-ix\cdot\xi}$ o module $1$.
Nierówność trójkąta pod całką daje oszacowanie
niezależne od $\xi$. Przy $h\to0$ funkcja podcałkowa
dąży punktowo do zera, a jej moduł nie przekracza
$2|f(x)|$. Twierdzenie o zbieżności zmajoryzowanej
dowodzi jednostajnej ciągłości.
\end{proof}

\begin{stwierdzenie}[Reguły transformacji całkowej]
\label{prop:dt-fourier-l1-rules}
Dla $f,g\in L^1(\R^d)$, $h,\eta\in\R^d$ i odwracalnej
rzeczywistej macierzy $A$ zachodzą tożsamości
\[
\begin{aligned}
 \widehat{\tau_h f}(\xi)&=e^{-ih\cdot\xi}\widehat f(\xi),\\
 \mathcal F(e^{i\eta\cdot x}f(x))(\xi)&=\widehat f(\xi-\eta),\\
 \mathcal F(f(A\,\cdot))(\xi)&=
 |\det A|^{-1}\widehat f(A^{-T}\xi),\\
 \widehat{f*g}(\xi)&=c_d^{-1}\widehat f(\xi)\widehat g(\xi).
\end{aligned}
\]
Jeżeli ponadto $x^\beta f\in L^1$ dla wszystkich
wielowskaźników $|\beta|\le m$, to $\widehat f\in C^m$
oraz $\partial_\xi^\beta\widehat f=(-i)^{|\beta|}
\mathcal F(x^\beta f)$.
\end{stwierdzenie}
\begin{proof}
Pierwsza tożsamość wystąpiła w dowodzie
Twierdzenia~\ref{thm:dt-fourier-l1-rl}; druga wynika z
połączenia wykładników. W trzeciej stosujemy podstawienie
$y=Ax$, dla którego $x\cdot\xi=y\cdot A^{-T}\xi$.
W ostatniej całkowalność modułu funkcji podcałkowej
wynika z
\[
 \int_{\R^d}\int_{\R^d}|f(y)g(x-y)|
 \,\mathrm d\lambda_d(y)\,\mathrm d\lambda_d(x)
 =\|f\|_1\|g\|_1.
\]
Twierdzenie Tonellego daje również bezwzględną zbieżność
całki definiującej $(f*g)(x)$ dla prawie każdego $x$
i oszacowanie $\|f*g\|_1\le\|f\|_1\|g\|_1$.
Możemy używać borelowskich reprezentantek $f,g$;
wtedy funkcja $(x,y)\mapsto f(y)g(x-y)$ jest mierzalna.
Zmiana reprezentantki nie zmienia splotu prawie
wszędzie, ponieważ zmiana o funkcję zerową prawie
wszędzie ma, z tego samego oszacowania, splot o normie
$L^1$ równej zeru.
Twierdzenie Fubiniego i podstawienie $z=x-y$ pozwalają
rozłożyć całkę na iloczyn dwóch całek. Czynnik wynosi
$c_d/c_d^2=c_d^{-1}$.
Dla pochodnych różniczkujemy pod całką, korzystając z
$|(e^{-itx_j}-1)/t|\le |x_j|$. Na każdym kolejnym etapie
majorantą jest $|x^{\beta+e_j}f(x)|$; ciągłość
pochodnych wynika ponownie ze zbieżności zmajoryzowanej.
\end{proof}

\begin{wniosek}[Transformata pochodnej słabej w $L^1$]
\label{cor:dt-fourier-weak-derivative-l1}
Niech $f,g\in L^1(\R^d;\C)$ i niech
$\partial_jT_f=T_g$ w $\mathcal D'(\R^d)$
dla ustalonego $j\in\{1,\ldots,d\}$.
Wtedy dla każdego $\xi\in\R^d$
\[
 \widehat g(\xi)=i\xi_j\widehat f(\xi).
\]
W szczególności jeśli wszystkie pochodne słabe
$\partial^\alpha f$ dla $|\alpha|\le m$ należą do $L^1$,
to $\mathcal F(\partial^\alpha f)=(i\xi)^\alpha\widehat f$
dla każdego takiego wielowskaźnika.
\end{wniosek}
\begin{proof}
Nie można bezpośrednio użyć funkcji wykładniczej
jako testu w definicji pochodnej dystrybucyjnej,
gdyż nie ma zwartego nośnika.
Z Lematu~\ref{lem:cutoffs17} wybieramy
$\chi\in\mathcal D(\R^d;\R)$, $0\le\chi\le1$,
równą $1$ w otoczeniu zera. Dla $R>0$ i
$\xi\in\R^d$ bierzemy test
$\varphi_R(x)=\chi(x/R)e^{-ix\cdot\xi}$.
Z definicji pochodnej dystrybucyjnej
\[
\begin{aligned}
 \int_{\R^d}g(x)\varphi_R(x)\,\mathrm d\lambda_d(x)
 &=i\xi_j\int_{\R^d} f(x)\varphi_R(x)\,\mathrm d\lambda_d(x)\\
 &\quad-\frac1R\int_{\R^d}f(x)(\partial_j\chi)(x/R)
 e^{-ix\cdot\xi}\,\mathrm d\lambda_d(x).
\end{aligned}
\]
Ostatni składnik ma moduł nie większy niż
$R^{-1}\|f\|_1\|\partial_j\chi\|_\infty$, więc zanika.
W dwóch pozostałych całkach stosujemy twierdzenie
o zbieżności zmajoryzowanej, ponieważ
$\chi(x/R)\to1$ punktowo, a majoranty to odpowiednio
$|g|$ i $|f|$.
Mnożąc otrzymaną równość przez $c_d$, dostajemy tezę.
Iterowanie tego wyniku po kolejnych pochodnych
pierwszego rzędu daje wzór dla wielowskaźników:
jeśli $|\alpha|>0$ i $\alpha_j>0$, stosujemy wzór
do pochodnej słabej $\partial^{\alpha-e_j}f$,
której pochodna w kierunku $j$ jest równa
$\partial^\alpha f$ jako dystrybucja regularna.
Obie funkcje należą do $L^1$ z założenia.
Indukcja rozpoczyna się od $\alpha=0$.
W szczególności z Twierdzenia~\ref{thm:dt-fourier-l1-rl}
mamy $\xi^\alpha\widehat f(\xi)\to0$ przy $|\xi|\to\infty$.
\end{proof}

\begin{lemat}[Odwracanie z regularyzacją Gaussa]
\label{lem:dt-fourier-gaussian-l1}
Niech $f\in L^1(\R^d)$ i $\varepsilon>0$.
Definiujemy jądro Gaussa
\[
 G_\varepsilon(x)=(2\pi\varepsilon)^{-d/2}
 e^{-|x|^2/(2\varepsilon)}.
\]
Wtedy dla każdego $x$ zachodzi
\[
 c_d\int_{\R^d}e^{ix\cdot\xi}
 e^{-\varepsilon|\xi|^2/2}\widehat f(\xi)
 \,\mathrm d\lambda_d(\xi)=(G_\varepsilon*f)(x).
\]
Ponadto $G_\varepsilon*f\to f$ w $L^1$ dla
$\varepsilon\downarrow0$.
\end{lemat}
\begin{proof}
Lewa strona jest bezwzględnie zbieżna na mocy
Twierdzenia~\ref{thm:dt-fourier-l1-rl}.
Podstawiamy definicję $\widehat f$ i stosujemy Fubiniego,
gdyż $|f(y)|e^{-\varepsilon|\xi|^2/2}$ jest całkowalne
na przestrzeni produktowej. Całkę po $\xi$ obliczamy
wzorem Gaussa z dowodu Twierdzenia~\ref{thm:fourier-schwartz17};
łącznie z czynnikiem $c_d^2$ daje to $G_\varepsilon(x-y)$.
Splot po prawej stronie jest określony dla każdego $x$,
ponieważ $G_\varepsilon$ jest ograniczona i $f\in L^1$.
Z tego samego dowodu wiemy, że $\int G_\varepsilon=1$
oraz całka $G_\varepsilon$ poza dowolną kulą o dodatnim
promieniu dąży do zera. Po podstawieniu $z=x-y$
możemy więc zapisać
\[
 (G_\varepsilon*f)(x)-f(x)=
 \int_{\R^d}G_\varepsilon(y)(f(x-y)-f(x))\,\mathrm d\lambda_d(y)
\]
dla prawie każdego $x$.
Nierówność trójkąta pod całką i twierdzenie Tonellego
dla nieujemnej funkcji
$G_\varepsilon(y)|f(x-y)-f(x)|$ dają
\[
 \|G_\varepsilon*f-f\|_1
 \le \int_{\R^d}G_\varepsilon(y)
 \|\tau_yf-f\|_1\,\mathrm d\lambda_d(y).
\]
Dla dowolnego $\delta>0$ część po $|y|<\delta$ nie
przekracza $\sup_{|y|<\delta}\|\tau_yf-f\|_1$,
a pozostała nie przekracza
$2\|f\|_1\int_{|y|\ge\delta}G_\varepsilon\,\mathrm d\lambda_d$.
Ciągłość translacji z dowodu Twierdzenia~\ref{thm:dt-approx-lp}
pozwala najpierw wybrać $\delta$,
a następnie zmniejszyć $\varepsilon$, tak aby obie
części były dowolnie małe.
\end{proof}

\begin{wniosek}[Punktowa zbieżność regularyzacji Gaussa]
\label{cor:dt-gaussian-pointwise}
Niech $f\in L^1(\R^d)$ i niech wybrana reprezentantka
$f$ będzie ciągła w punkcie $x_0\in\R^d$.
Wtedy
\[
 \lim_{\varepsilon\downarrow0}(G_\varepsilon*f)(x_0)
 =f(x_0),
\]
gdzie $G_\varepsilon$ jest jądrem z
Lematu~\ref{lem:dt-fourier-gaussian-l1}.
Nie zakładamy, że $\widehat f\in L^1$ ani że $f$
jest ograniczona na całej przestrzeni.
\end{wniosek}
\begin{proof}
Dla $\eta>0$ wybieramy $\delta>0$ tak, aby
$|f(x_0-y)-f(x_0)|<\eta$ dla $|y|<\delta$.
Korzystamy z $\int G_\varepsilon\,\mathrm d\lambda_d=1$.
Nierówność trójkąta i rozbicie całki dają
\[
 \begin{aligned}
 |(G_\varepsilon*f)(x_0)-f(x_0)|
 &\le \eta+\|f\|_1
 \sup_{|y|\ge\delta}G_\varepsilon(y)\\
 &\quad+|f(x_0)|\int_{|y|\ge\delta}G_\varepsilon(y)
 \,\mathrm d\lambda_d(y).
 \end{aligned}
\]
W drugim składniku użyliśmy niezmienniczości miary
Lebesgue'a względem translacji i odbicia, aby ograniczyć
całkę z $|f(x_0-y)|$ przez $\|f\|_1$.
Dla ustalonego $\delta>0$
\[
 \sup_{|y|\ge\delta}G_\varepsilon(y)
 =(2\pi\varepsilon)^{-d/2}e^{-\delta^2/(2\varepsilon)}
 \longrightarrow0,
\]
gdyż funkcja wykładnicza zanika szybciej niż dowolna
ujemna potęga $\varepsilon$.
Trzeci składnik dąży do zera na mocy własności
jądra Gaussa wykazanej w dowodzie
Lematu~\ref{lem:dt-fourier-gaussian-l1}.
Granica górna błędu nie przekracza więc $\eta$,
a $\eta>0$ jest dowolna.
\end{proof}

\begin{twierdzenie}[Odwrócenie i jednoznaczność w $L^1$]
\label{thm:dt-fourier-l1-inversion}
Jeżeli $f,\widehat f\in L^1(\R^d)$, to
\[
 f(x)=c_d\int_{\R^d}e^{ix\cdot\xi}\widehat f(\xi)
 \,\mathrm d\lambda_d(\xi)
\]
dla prawie każdego $x$, a prawa strona jest ciągłym,
zanikającym w nieskończoności reprezentantem $f$.
Jeżeli $f$ jest ciągła, równość zachodzi wszędzie.
Niezależnie od całkowalności $\widehat f$ transformacja
na $L^1$ jest różnowartościowa.
\end{twierdzenie}
\begin{proof}
Niech $g$ oznacza prawą stronę wzoru odwrotnego.
Twierdzenie~\ref{thm:dt-fourier-l1-rl}, ze zmienionym
znakiem argumentu, daje $g\in C_0$.
Z Lematu~\ref{lem:dt-fourier-gaussian-l1} oraz
zbieżności zmajoryzowanej mamy
\[
 \|G_\varepsilon*f-g\|_\infty
 \le c_d\int_{\R^d}
 (1-e^{-\varepsilon|\xi|^2/2})|\widehat f(\xi)|
 \,\mathrm d\lambda_d(\xi)\longrightarrow0.
\]
Z drugiej strony $G_\varepsilon*f\to f$ w $L^1$.
Na każdej kuli skończonej miary zbieżność jednostajna
implikuje zbieżność w $L^1$ na tej kuli. Jednoznaczność
granicy daje $g=f$ prawie wszędzie na każdej takiej kuli,
a ich przeliczalne wyczerpanie daje równość na $\R^d$.
Jeżeli $f$ jest ciągła, różnica dwóch funkcji ciągłych
równa zeru prawie wszędzie jest identycznie zerowa.
Jeżeli $\widehat f=0$, nawet bez założenia o
całkowalności transformaty, lewa strona tożsamości
z Lematu~\ref{lem:dt-fourier-gaussian-l1} jest zerowa
dla każdego $\varepsilon$, więc $G_\varepsilon*f=0$.
Przejście do granicy w $L^1$ daje $f=0$.
Liniowość implikuje różnowartościowość.
\end{proof}

\begin{twierdzenie}[Obraz transformacji całkowej w $C_0$]
\label{thm:dt-fourier-l1-range}
Dla każdego $d\ge1$ operator
$\mathcal F:L^1(\R^d)\to C_0(\R^d)$ ma normę $c_d$.
Jego obraz jest gęstą, właściwą podprzestrzenią liniową
$C_0(\R^d)$, a operator odwrotny określony na tym obrazie
nie jest ciągły względem normy supremum.
\end{twierdzenie}
\begin{proof}
Oszacowanie $\|\mathcal F\|\le c_d$ pochodzi z
Twierdzenia~\ref{thm:dt-fourier-l1-rl}. Dla nieujemnej
funkcji $f\in L^1$ o normie $1$ mamy
$\widehat f(0)=c_d$, więc zachodzi równość.

Z Twierdzenia~\ref{thm:fourier-schwartz17}
obraz zawiera $\mathcal S$, w szczególności
$C_c^\infty(\R^d)$.
Aby wykazać gęstość, niech $v\in C_0(\R^d)$ i $\eta>0$.
Wybieramy $R>0$ tak, aby $|v(x)|<\eta$ poza $B(0,R)$,
oraz, z Lematu~\ref{lem:cutoffs17}, funkcję
$\chi\in\mathcal D(\R^d;\R)$ równą $1$ na
$\overline{B(0,R)}$ i spełniającą $0\le\chi\le1$.
Funkcja $w=\chi v$ jest ciągła o zwartym nośniku
i $\|w-v\|_\infty\le\eta$.
Jest ona jednostajnie ciągła: na dostatecznie dużej
domkniętej kuli wynika to ze zwartości, a dla dwóch
punktów o odległości mniejszej niż $1$, z których
jeden leży poza kulą zawierającą jednostkowe otoczenie
nośnika, obie wartości są zerowe.
Niech $\rho_\varepsilon$ będzie aproksymacją jedności
z Definicji~\ref{def:mollifier21}.
Z Twierdzenia~\ref{thm:smoothing21} splot
$w*\rho_\varepsilon$ jest funkcją testową, a
\[
 \|w*\rho_\varepsilon-w\|_\infty
 \le\sup_{|y|\le\varepsilon}\sup_x|w(x-y)-w(x)|
 \longrightarrow0.
\]
Dowodzi to gęstości.

Dla wykazania braku ciągłości odwrotności skonstruujemy
ciąg funkcji o normie $L^1$ równej $1$, których transformaty
dążą jednostajnie do zera.
Wybieramy nieujemne $\psi\in C_c^\infty(\R^d)$ o normie
$L^1$ równej $1$ oraz o nośniku w $B(0,1)$.
Dla $N=2^k$, $k\ge0$, określamy wielomiany rekurencyjnie:
\[
 P_0(z)=Q_0(z)=1,\qquad
 \begin{cases}
 P_{k+1}(z)=P_k(z)+z^{2^k}Q_k(z),\\
 Q_{k+1}(z)=P_k(z)-z^{2^k}Q_k(z).
 \end{cases}
\]
Indukcyjnie oba wielomiany mają stopień $2^k-1$
i wszystkie współczynniki równe $1$ lub $-1$:
przesunięcie o $2^k$ rozdziela zakresy potęg obu składników.
Dla $|z|=1$ tożsamość
$|a+b|^2+|a-b|^2=2(|a|^2+|b|^2)$ daje
\[
 |P_k(z)|^2+|Q_k(z)|^2=2^{k+1}=2N.
\]
Jeżeli $P_k(z)=\sum_{j=0}^{N-1}b_jz^j$,
połóżmy
\[
 f_k(x)=N^{-1}\sum_{j=0}^{N-1}b_j\psi(x-3je_1).
\]
Nośniki składników są parami rozłączne, stąd
$\|f_k\|_1=N^{-1}\sum_j|b_j|\|\psi\|_1=1$.
Z reguły translacji w
Stwierdzeniu~\ref{prop:dt-fourier-l1-rules}
\[
 \widehat f_k(\xi)=N^{-1}\widehat\psi(\xi)
 P_k(e^{-3i\xi_1}),
 \qquad
 \|\widehat f_k\|_\infty\le c_d\sqrt{2/N}\longrightarrow0.
\]
Różnowartościowość z Twierdzenia~\ref{thm:dt-fourier-l1-inversion}
pozwala mówić o odwrotności na obrazie. Powyższy ciąg
dowodzi, że nie jest ona ciągła w zerze.
Gdyby obraz był całym $C_0(\R^d)$, twierdzenie
o odwzorowaniu otwartym dla przestrzeni Banacha
$L^1$ i $C_0$ wymuszałoby ciągłość odwrotności.
Sprzeczność dowodzi, że obraz jest właściwy.
\end{proof}

\begin{przyklad}[Prostokąt i funkcja wykładnicza]
\label{ex:dt-fourier-l1-pairs}
Dla $a>0$ w wymiarze jeden mamy
\[
 \mathcal F\chi_{[-a,a]}(\xi)
 =\sqrt{\frac{2}{\pi}}\frac{\sin(a\xi)}{\xi},
 \qquad
 \mathcal F(e^{-a|x|})(\xi)
 =\sqrt{\frac{2}{\pi}}\frac{a}{a^2+\xi^2}.
\]
Pierwszą wartość przy $\xi=0$ rozumiemy jako
$\sqrt{2/\pi}\,a$. Pierwszy wzór wynika z całkowania
$e^{-ix\xi}$ po $[-a,a]$. W drugim rozdzielamy całkę
na półproste, otrzymując
$c_1((a+i\xi)^{-1}+(a-i\xi)^{-1})$.
Transformata prostokąta nie należy do $L^1$: na
przedziałach $a\xi\in[n\pi+\pi/6,n\pi+5\pi/6]$
dla $n\ge1$ licznik ma moduł co najmniej $1/2$,
więc suma całek jest ograniczona z dołu przez
dodatnią stałą razy $\sum_{n\ge1}(n+1)^{-1}$.
W drugim przypadku obie funkcje należą do $L^1$,
a wzór odwrotny obowiązuje w każdym punkcie.
\end{przyklad}

\begin{uwaga}[Zgodność z normalizacją RRC]
\label{rem:dt-fourier-normalization-rrc}
W skrypcie RRC-08 stosowana jest transformata
$\mathcal F_0f(\xi)=\int e^{-ix\cdot\xi}f(x)
\,\mathrm d\lambda_d(x)$. Mamy $\mathcal F=c_d\mathcal F_0$;
we wzorze odwrotnym dla $\mathcal F_0$ występuje
$(2\pi)^{-d}$, a we wzorze na splot nie ma dodatkowego
czynnika. To nie są odmienne twierdzenia, lecz różne
normalizacje tego samego operatora. Por.
\atlasref{PDE-08::prop:pde8-RL}{RRC-08: lemat Riemanna--Lebesgue'a}
oraz
\atlasref{PDE-08::thm:pde8-inversion}{RRC-08: odwrócenie transformaty}.
\end{uwaga}

\setcounter{section}{9}
% Wersja źródłowa: audyt II, jednostka 22

\section{Dystrybucje temperowane i transformacja Fouriera}
\label{sec:temperowane22}

\providecommand{\atlasref}[2]{#2}


\noindent\textbf{Fakty z wcześniejszych rozdziałów.}
Przyjmujemy, że przestrzeń Schwartza jest zupełną przestrzenią Frécheta,
$\mathcal D(\R^d)$ jest w niej gęsta, a transformacja Fouriera
jest jej ciągłym automorfizmem z odwrotnością i wzorem
$\mathcal F^2\varphi=\varphi(-\cdot)$.
Dystrybucja o zwartym nośniku $K$ działa również na funkcjach
$C^\infty$ poprzez $\langle T,f\rangle=\langle T,\chi f\rangle$,
gdzie $\chi=1$ w otoczeniu $K$, i spełnia oszacowanie
skończonego rzędu na ustalonym zwartym otoczeniu $K$.
Korzystamy z Twierdzeń~\ref{thm:schwartz-frechet17},
\ref{thm:density-schwartz17}, \ref{thm:fourier-schwartz17}
i Lematu~\ref{lem:compact-action21}. Wzór na $\mathcal F^2$
udowodniono we Wniosku~\ref{cor:dt-fourier-iteration}.

W całym rozdziale $d\ge1$ jest liczbą całkowitą.
Rozważamy funkcje oraz dystrybucje zespolone na $\R^d$,
chyba że jawnie dopuszczamy także ciało $\R$.
Przestrzeń Schwartza $\mathcal S=\mathcal S(\R^d;\C)$
wyposażamy w topologię Frécheta z rozdziału~\ref{sec:gladkie-schwartz17}.
Dla $N,m\in\N_0$ przyjmujemy seminormy (również oznaczane
w dalszej części przez $p_{N,m}=Q_{N,m}$)
$Q_{N,m}(\varphi)=\max_{|\alpha|\le m}\sup_{x\in\R^d}
(1+|x|)^N|\partial^\alpha\varphi(x)|$.
Przez $\mathcal S'$ rozumiemy przestrzeń \emph{ciągłych,
$\C$-liniowych} funkcjonałów na $\mathcal S$.
W parowaniu $\langle T,\varphi\rangle$ nie sprzęgamy zespolenie
funkcji testowej. Przypominamy, że
$\mathcal D(\R^d)=C_c^\infty(\R^d)$ jest gęsta w $\mathcal S$.

\subsection{Dystrybucje temperowane jako przestrzeń sprzężona}

\noindent\emph{Powiązania tematyczne w Bibliotece matematycznej:} \atlasref{PDE-08::def:pde8-tempered}{PDE-08: definicja dystrybucji temperowanej}.


\begin{definicja}[Dystrybucja temperowana]
\label{def:tempered22}
Dystrybucją temperowaną na $\R^d$ nazywamy element
$T\in\mathcal S'$, czyli ciągły liniowy funkcjonał
$T:\mathcal S\to\C$. Rozpatrujemy na $\mathcal S'$ topologię
słabą-* $\sigma(\mathcal S',\mathcal S)$, wyznaczoną przez
seminormy $T\mapsto|\langle T,\varphi\rangle|$
($\varphi\in\mathcal S$), oraz topologię silną
$\beta(\mathcal S',\mathcal S)$, wyznaczoną przez
$P_B(T)=\sup_{\varphi\in B}|\langle T,\varphi\rangle|$,
gdzie $B$ przebiega wszystkie niepuste zbiory ograniczone w $\mathcal S$.
Dla każdego ciągłego funkcjonału $T$ liczba $P_B(T)$
jest skończona, ponieważ ciągłe odwzorowania liniowe
przenoszą zbiory ograniczone na zbiory ograniczone.
\end{definicja}

\begin{twierdzenie}[Kryterium wzrostu seminormowego]
\label{thm:tempered-estimate22}
Niech $T:\mathcal S(\R^d;\C)\to\C$ będzie liniowy.
Następujące warunki są równoważne:
\begin{enumerate}
\item $T\in\mathcal S'(\R^d)$;
\item istnieją $C>0$ oraz $N,m\in\N_0$ takie, że
$|\langle T,\varphi\rangle|\le C Q_{N,m}(\varphi)$
dla wszystkich $\varphi\in\mathcal S$.
\end{enumerate}
W szczególności ograniczenie $T|_{\mathcal D}$ jest dystrybucją,
a odwzorowanie $\mathcal S'\to\mathcal D'$ dane przez
ograniczenie jest liniowe i różnowartościowe.
\end{twierdzenie}
\begin{proof}
Rodzina $Q_{j,j}$, $j\in\N_0$, jest rosnąca
i wyznacza topologię $\mathcal S$. Jeśli $T$ jest ciągły,
to dla pewnych $j$ i $\varepsilon>0$ warunek
$Q_{j,j}(\varphi)<\varepsilon$ implikuje
$|T(\varphi)|<1$. Liniowość i skalowanie dają
$|T(\varphi)|\le (2/\varepsilon)Q_{j,j}(\varphi)$;
gdy seminorma znika, wniosek wynika ze skalowania przez
dowolną liczbę. Odwrotna implikacja jest natychmiastowa.
Włożenie $j:\mathcal D\hookrightarrow\mathcal S$ jest ciągłe:
dla niepustego zwartego $K\subseteq\R^d$ i
$\varphi\in\mathcal D_K$ mamy
$Q_{N,m}(\varphi)\le (1+\max_{x\in K}|x|)^N p_m(\varphi)$.
Dla $K=\varnothing$ etap jest przestrzenią zerową.
Ciągłość $j$ wynika z własności uniwersalnej granicy
induktywnej (Twierdzenie~\ref{thm:LF-universal18}).
Stąd ograniczenie $T|_{\mathcal D}=T\circ j$ jest ciągłe.
Odwzorowanie $j^{\mathrm t}:T\mapsto T\circ j$ jest
liniowe i różnowartościowe: jeśli dwa ciągłe funkcjonały
zgadzają się na gęstej podprzestrzeni $\mathcal D$,
to ich różnica znika na całej $\mathcal S$.
\end{proof}

\begin{twierdzenie}[Rodziny ograniczone dystrybucji temperowanych]
\label{thm:dt-tempered-bounded-families}
Niech $\mathcal A\subseteq\mathcal S'$. Następujące warunki
są równoważne:
\begin{enumerate}
\item $\mathcal A$ jest ograniczona w topologii słabej-*;
\item $\mathcal A$ jest ograniczona w topologii silnej;
\item istnieją $C>0$ i $N,m\in\N_0$ takie, że
\[
 |\langle T,\varphi\rangle|\le C Q_{N,m}(\varphi)
 \qquad(T\in\mathcal A,\ \varphi\in\mathcal S).
\]
\end{enumerate}
\end{twierdzenie}
\begin{proof}
Jeżeli $\mathcal A=\varnothing$, wszystkie warunki są
spełnione. W pozostałych przypadkach ograniczoność silna
implikuje słabą-* (singleton w $\mathcal S$ jest ograniczony).
Ograniczoność słaba-* oznacza, że dla każdej $\varphi\in\mathcal S$
zbiór $\{\langle T,\varphi\rangle:T\in\mathcal A\}$ jest
ograniczony w $\C$. Przestrzeń $\mathcal S$ jest przestrzenią
Frécheta, a zatem jest beczkowa
(Wniosek~\ref{cor:frechet-barrelled}).
Z twierdzenia Banacha--Steinhausa
(Twierdzenie~\ref{thm:banach-steinhaus-lcs}, z przestrzenią
wartości $\C$) wynika równociągłość rodziny $\mathcal A$. Istnieją więc
$j\in\N_0$ oraz $\varepsilon>0$ takie, że
$Q_{j,j}(\varphi)<\varepsilon$ implikuje
$|\langle T,\varphi\rangle|<1$ dla wszystkich $T\in\mathcal A$.
Dla $\varphi\ne0$ stosujemy to do
$\varepsilon\varphi/(2Q_{j,j}(\varphi))$.
Dostajemy wspólne oszacowanie z $C=2/\varepsilon$
oraz $N=m=j$; dla $\varphi=0$ nierówność jest oczywista.
Wreszcie wspólne oszacowanie implikuje
\[
 \sup_{T\in\mathcal A}P_B(T)
 \le C\sup_{\varphi\in B}Q_{N,m}(\varphi)<\infty
\]
dla każdego niepustego ograniczonego $B\subseteq\mathcal S$.
Jest to dokładnie ograniczoność w topologii silnej.
\end{proof}

\begin{uwaga}
Wspólne oszacowanie oznacza kontrolę zarówno wzrostu
w nieskończoności (parametr $N$), jak i liczby pochodnych
testu (parametr $m$). Nie można na ogół wybierać tych
parametrów osobno dla kolejnych elementów rodziny, jeżeli
chcemy wnioskować o jej ograniczoności. Przykładowo
rodzina $(\delta_{ne_1})_{n\ge1}$ spełnia wspólne
oszacowanie z $N=m=0$.
\end{uwaga}

\begin{przyklad}[Dystrybucje regularne i miary]
\label{ex:tempered-regular22}
Jeżeli $f\in L^1_{\mathrm{loc}}(\R^d)$ i
dla pewnych $C>0$, $M\ge0$ zachodzi
$|f(x)|\le C(1+|x|)^M$ prawie wszędzie,
to $T_f(\varphi)=\int_{\R^d}f(x)\varphi(x)\,\mathrm d\lambda_d(x)$
określa element $\mathcal S'$. Istotnie, dla całkowitego $N>M+d$ mamy
\[
 \int_{\R^d}|f(x)\varphi(x)|\,\mathrm d\lambda_d(x)
 \le C Q_{N,0}(\varphi)
 \int_{\R^d}(1+|x|)^{M-N}\,\mathrm d\lambda_d(x).
\]
Ostatnia całka jest skończona: na kuli jednostkowej
funkcja jest ograniczona, a na pierścieniu
$2^k\le |x|<2^{k+1}$ całkę ogranicza stała razy
$2^{k(d+M-N)}$. Szereg tych oszacowań jest zbieżny.
Z Twierdzenia~\ref{thm:tempered-estimate22} wynika
temperowalność $T_f$.
Bardziej ogólnie, jeżeli $\mu$ jest zespoloną lokalnie skończoną
miarą borelowską i dla pewnego $N\ge0$ zachodzi
$\int_{\R^d}(1+|x|)^{-N}\,\mathrm d|\mu|(x)<\infty$,
to $\langle T_\mu,\varphi\rangle=
\int_{\R^d}\varphi(x)\,\mathrm d\mu(x)$
definiuje dystrybucję temperowaną. Istotnie, biorąc
liczbę całkowitą $n\ge N$, mamy
\[
 \left|\int_{\R^d}\varphi\,\mathrm d\mu\right|
 \le Q_{n,0}(\varphi)
 \int_{\R^d}(1+|x|)^{-N}\,\mathrm d|\mu|(x).
\]
Całka jest bezwzględnie zbieżna względem $|\mu|$,
a oszacowanie dowodzi ciągłości funkcjonału.
Przykładami są funkcja stała $1$, wielomiany i miary Diraca.
\end{przyklad}

\begin{przyklad}[Dystrybucja nietemperowana]
\label{ex:non-tempered22}
Dystrybucja regularna wyznaczona przez
$f(x)=e^{|x|^2}$ należy do $\mathcal D'(\R^d)$,
ale nie jest ograniczeniem żadnego elementu $\mathcal S'$.
Aby to wykazać, weźmy $\chi\in\mathcal D(\R^d)$,
$0\le\chi\le1$, równą $1$ na $\overline B(0,1)$,
i funkcje $\varphi_R(x)=\chi(x/R)e^{-|x|^2/2}$.
Dla $R\to\infty$ zachodzi
$\varphi_R\to e^{-|x|^2/2}$ w $\mathcal S$
na mocy dowodu Twierdzenia~\ref{thm:density-schwartz17}:
dla $g(x)=e^{-|x|^2/2}\in\mathcal S$ i $R\ge1$ zachodzi
$Q_{N,m}(g-\chi(\cdot/R)g)\le
C_m(1+R)^{-1}Q_{N+1,m}(g)$.
Przynależność $g$ do $\mathcal S$ wynika stąd, że
każda jego pochodna jest wielomianem pomnożonym przez $g$.
Jednak
$\langle T_f,\varphi_R\rangle
=\int_{\R^d}\chi(x/R)e^{|x|^2/2}\,\mathrm d\lambda_d(x)
\to+\infty$,
gdyż dla $R\ge1$ całka jest nie mniejsza niż
$\int_{B(0,R)}1\,\mathrm d\lambda_d=
R^d\lambda_d(B(0,1))$.
Gdyby istniało ciągłe przedłużenie na $\mathcal S$,
wartości te musiałyby zbiegać do skończonej liczby.
\end{przyklad}

\begin{stwierdzenie}[Podprzestrzenie klasycznych dystrybucji]
\label{prop:schwartz-embeddings22}
Istnieją naturalne, różnowartościowe i ciągłe odwzorowania
\[
 \mathcal D\hookrightarrow\mathcal S\hookrightarrow
 L^2(\R^d)\hookrightarrow\mathcal S'\hookrightarrow\mathcal D',
\]
przy czym obie przestrzenie dualne wyposażamy albo
w topologie silne, albo w topologie słabe-*,
gdzie funkcję $f\in L^2$ identyfikujemy z funkcjonałem
$\langle T_f,\varphi\rangle=\int_{\R^d}f(x)\varphi(x)
\,\mathrm d\lambda_d(x)$.
Ponadto każda dystrybucja o zwartym nośniku jest temperowana:
$\mathcal E'(\R^d)\subseteq\mathcal S'(\R^d)$.
\end{stwierdzenie}
\begin{proof}
Pierwsze włożenie wykazano w Twierdzeniu~
\ref{thm:tempered-estimate22}. Drugie jest ciągłe na mocy
Stwierdzenia~\ref{prop:schwartz-operators17} i różnowartościowe,
ponieważ ciągła funkcja równa zeru prawie wszędzie znika
wszędzie. Dla $f\in L^2$ nierówność Cauchy'ego--Schwarza daje
$|\langle T_f,\varphi\rangle|\le\|f\|_2\|\varphi\|_2$;
także $\|\varphi\|_2\le C_dQ_{d+1,0}(\varphi)$.
Dla ograniczonego $B\subseteq\mathcal S$ otrzymujemy
\[
 P_B(T_f)\le C_d\|f\|_2\sup_{\varphi\in B}Q_{d+1,0}(\varphi).
\]
Prawa strona jest skończona i dowodzi ciągłości silnej,
a dla singletonów także słabej-*.
Jeśli $T_f=0$, to $\int f\varphi\,\mathrm d\lambda_d=0$
dla wszystkich $\varphi\in\mathcal D$.
Z gęstości $\mathcal D$ w $L^2$
(Twierdzenie~\ref{thm:dt-test-density}) wybieramy $\varphi_n\to
\overline f$ w $L^2$; nierówność Cauchy'ego--Schwarza daje
$\int |f|^2\,\mathrm d\lambda_d=0$, więc $f=0$ w $L^2$.
Ostatnie odwzorowanie jest różnowartościowe z
Twierdzenia~\ref{thm:tempered-estimate22}. Dla
$B\subseteq\mathcal D$ ograniczonego mamy
$P_B(T|_{\mathcal D})=P_{j(B)}(T)$, gdzie $j(B)$ jest
ograniczony w $\mathcal S$ dzięki ciągłości $j$.
Daje to ciągłość silną, a oceny na pojedynczych
funkcjach testowych dają ciągłość słabą-*.
Jeżeli $T\in\mathcal E'$ ma zwarty nośnik,
to jego przedłużenie z Lematu~\ref{lem:compact-action21}
spełnia $|T(\varphi)|\le C Q_{0,m}(\varphi)$.
Jest więc ciągłe na $\mathcal S$ i zgadza się z $T$
na $\mathcal D$; identyfikujemy $T$ z tym jednoznacznym
przedłużeniem.
\end{proof}

\subsection{Operatory na przestrzeni dystrybucji temperowanych}

\begin{twierdzenie}[Różniczkowanie i mnożenie przez wielomian]
\label{thm:tempered-operators22}
Dla każdego $\alpha,\beta\in\N_0^d$ operatory
$T\mapsto\partial^\alpha T$ i $T\mapsto x^\beta T$
są liniowe i ciągłe na $\mathcal S'$ zarówno dla topologii
$\sigma(\mathcal S',\mathcal S)$, jak i
$\beta(\mathcal S',\mathcal S)$, gdzie
\[
\langle\partial^\alpha T,\varphi\rangle
=(-1)^{|\alpha|}\langle T,\partial^\alpha\varphi\rangle,
\qquad\langle x^\beta T,\varphi\rangle
=\langle T,x^\beta\varphi\rangle.
\]
To samo dotyczy mnożenia przez każdą funkcję
$a\in C^\infty(\R^d;\C)$, dla której dla każdego
$\gamma\in\N_0^d$ istnieją $C_\gamma>0$ oraz
$M_\gamma\in\N_0$ takie, że
$|\partial^\gamma a(x)|\le C_\gamma(1+|x|)^{M_\gamma}$
dla wszystkich $x\in\R^d$.
\end{twierdzenie}
\begin{proof}
Różniczkowanie i mnożenie przez $x^\beta$ są ciągłymi
operatorami liniowymi $\mathcal S\to\mathcal S$,
co wykazano w Stwierdzeniu~\ref{prop:schwartz-operators17}.
Dla ogólnego mnożnika $a$ ustalmy $N,m\in\N_0$,
$M_m=\max_{|\gamma|\le m}M_\gamma$ i
$C_m=\max_{|\beta|\le m}\sum_{\gamma\le\beta}
\binom\beta\gamma C_\gamma$.
Ze wzoru Leibniza, dla $|\beta|\le m$, dostajemy
\[
 (1+|x|)^N|\partial^\beta(a\varphi)(x)|
 \le C_m Q_{N+M_m,m}(\varphi).
\]
Po wzięciu supremum oraz maksimum otrzymujemy
$Q_{N,m}(a\varphi)\le C_m Q_{N+M_m,m}(\varphi)$.
Zatem odwzorowanie $\varphi\mapsto a\varphi$ jest
ciągłym operatorem na $\mathcal S$.
Dla dowolnego ciągłego operatora liniowego
$A:\mathcal S\to\mathcal S$ jego transpozycja
$A^{\mathrm t}T=T\circ A$ spełnia
$P_B(A^{\mathrm t}T)=P_{A(B)}(T)$.
Ciągłość $A$ zapewnia ograniczoność $A(B)$, skąd
ciągłość silna; dla singletonów otrzymujemy ciągłość
słabą-*. Stosujemy to do $A\varphi=(-1)^{|\alpha|}
\partial^\alpha\varphi$ i $A\varphi=a\varphi$,
zgodnie z definicjami pochodnej i mnożenia dystrybucji.
Ograniczenie tych działań do $\mathcal D$ zgadza się
z działaniami na zwykłych dystrybucjach.
\end{proof}

\begin{twierdzenie}[Charakterystyka mnożników przestrzeni Schwartza]
\label{thm:dt-schwartz-multiplier-characterization}
Niech $a\in C^\infty(\R^d;\C)$. Następujące warunki
są równoważne:
\begin{enumerate}
\item $af\in\mathcal S$ dla każdego $f\in\mathcal S$;
\item operator $M_a:f\mapsto af$ jest ciągły na $\mathcal S$;
\item dla każdego $\alpha\in\N_0^d$ istnieją
$C_\alpha>0$ i $N_\alpha\in\N_0$ takie, że
\[
 |\partial^\alpha a(x)|\le C_\alpha(1+|x|)^{N_\alpha}
 \qquad(x\in\R^d).
\]
\end{enumerate}
\end{twierdzenie}
\begin{proof}
Implikację z warunku 3 do 2 wykazano w dowodzie
Twierdzenia~\ref{thm:tempered-operators22}, a warunek 2
implikuje 1 z definicji operatora na $\mathcal S$.
Pokażemy, że warunek 1 implikuje 3.
Gdyby warunek 3 nie zachodził, dla pewnego ustalonego
wielowskaźnika $\alpha$ i każdego $n\ge1$ funkcja
$x\mapsto (1+|x|)^{-n}|\partial^\alpha a(x)|$
byłaby nieograniczona. Ponieważ jest ciągła, jest
ograniczona na każdym zbiorze zwartym.
Dla $n=1$ wybieramy punkt poza kulą domkniętą
o promieniu $1$, gdzie iloraz przekracza $1$.
Dla $n\ge2$ wybieramy punkt poza kulą domkniętą
o promieniu $\max(n,|x_{n-1}|+2)$, gdzie odpowiedni
iloraz przekracza $1$. Jest to możliwe, ponieważ
ograniczoność ilorazu poza taką kulą wraz z jego
ciągłością na kuli dawałaby ograniczoność globalną.
Otrzymujemy zatem punkty $x_n$ takie, że
$|x_n|>n$, $|x_n|>|x_{n-1}|+2$ dla $n\ge2$ oraz
\[
 |\partial^\alpha a(x_n)|>(1+|x_n|)^n.
\]
Z Lematu~\ref{lem:cutoffs17} wybieramy
$\chi\in\mathcal D(B(0,1/2);\R)$ równą $1$ w otoczeniu zera
i przedłużamy ją zerem na $\R^d$.
Niech
\[
 f(x)=\sum_{n=1}^\infty
 (1+|x_n|)^{-n/2}\chi(x-x_n).
\]
Nośniki składników są parami rozłączne, ponieważ
$|x_n-x_k|\ge\bigl||x_n|-|x_k|\bigr|>2$ dla $n\ne k$.
Rodzina nośników jest lokalnie skończona dzięki
$|x_n|\to\infty$, zatem $f$ jest gładka.
Dla $|x-x_n|\le1/2$ mamy
$1+|x|\le(3/2)(1+|x_n|)$. W konsekwencji dla $N,m\in\N_0$
\[
 Q_{N,m}(f)\le (3/2)^N Q_{0,m}(\chi)
 \sup_{n\ge1}(1+|x_n|)^{N-n/2}<\infty.
\]
Istotnie, dla $n\ge2N$ wyrażenie pod supremum jest
nie większe niż $1$, a pozostałych indeksów jest
skończenie wiele. Poza sumą nośników wszystkie
pochodne $f$ znikają.
Stąd $f\in\mathcal S$. W otoczeniu $x_n$ funkcja $f$
jest stała i równa $(1+|x_n|)^{-n/2}$, więc
\[
 |\partial^\alpha(af)(x_n)|
 =(1+|x_n|)^{-n/2}|\partial^\alpha a(x_n)|
 >(1+|x_n|)^{n/2}\longrightarrow\infty.
\]
To przeczy $af\in\mathcal S$, gdyż $Q_{0,|\alpha|}(af)$
musiałaby być skończona. Otrzymujemy warunek 3.
\end{proof}

\begin{lemat}[Względna zwartość zbiorów ograniczonych w $\mathcal S$]
\label{lem:montel-schwartz22}
Każdy zbiór ograniczony w $\mathcal S(\R^d)$ ma zwarte
domknięcie w jej topologii Frécheta.
\end{lemat}
\begin{proof}
Niech $B\subseteq\mathcal S$ będzie ograniczony i
$(\varphi_n)\subseteq B$ dowolnym ciągiem.
Dla każdej liczby całkowitej $R\ge1$ i każdego $m\ge0$
rodziny pochodnych $\partial^\alpha\varphi_n$ dla
$|\alpha|\le m$ są jednostajnie ograniczone i
jednostajnie równociągłe na $\overline B(0,R)$:
równociągłość wynika z ograniczenia pochodnych
rzędu $m+1$ i całkowania wzdłuż odcinka w kuli:
dla $x,y$ w tej kuli oraz $|\alpha|\le m$ mamy
$|\partial^\alpha\varphi_n(x)-\partial^\alpha\varphi_n(y)|
\le\|x-y\|_1\sup_{\psi\in B}Q_{0,m+1}(\psi)$. Z twierdzenia Arzelà--Ascoliego
i argumentu diagonalnego uzyskujemy podciąg
$(\varphi_{n_k})$, którego wszystkie pochodne zbiegają
jednostajnie na wszystkich kulach domkniętych.
Nie musimy osobno identyfikować granic pochodnych:
poniższe oszacowanie da zbieżność w zupełnej
przestrzeni Schwartza.
Dla ustalonych $N,m$ i $R\ge1$ poza $B(0,R)$ zachodzi
\[
(1+|x|)^N|\partial^\alpha(\varphi_{n_k}-\varphi_{n_l})(x)|
\le R^{-1}\bigl(Q_{N+1,m}(\varphi_{n_k})
+Q_{N+1,m}(\varphi_{n_l})\bigr)\quad(|\alpha|\le m).
\]
Prawa strona jest jednostajnie mała przy dużym $R$;
wewnątrz kuli $B(0,R)$ zbieżność wszystkich pochodnych
jest jednostajna. Dlatego $(\varphi_{n_k})$ jest ciągiem
Cauchy'ego we wszystkich seminormach $Q_{N,m}$.
Zupełność $\mathcal S$ daje zbieżność do elementu
$\mathcal S$. Ponieważ przestrzeń jest metryzowalna,
ciągowa względna zwartość zbioru jest równoważna
zwartości jego domknięcia. Istotnie, dla ciągu $(g_n)$
w domknięciu $B$ wybieramy $f_n\in B$ w odległości
metrycznej mniejszej niż $1/n$ od $g_n$.
Zbieżny podciąg $(f_{n_k})$ daje zbieżny podciąg
$(g_{n_k})$ o tej samej granicy, należącej do domknięcia.
Zwartość wynika z kryterium ciągowego dla przestrzeni
metrycznych.
\end{proof}

\begin{twierdzenie}[Zbieżność słaba-* i silna ciągów]
\label{thm:seq-topologies-schwartz22}
Niech $(T_n)$ będzie ciągiem dystrybucji temperowanych,
a $T\in\mathcal S'$. Wtedy
\[
T_n\longrightarrow T\text{ w }\sigma(\mathcal S',\mathcal S)
\quad\Longleftrightarrow\quad
T_n\longrightarrow T\text{ w }\beta(\mathcal S',\mathcal S).
\]
Mimo to topologie $\sigma(\mathcal S',\mathcal S)$ i
$\beta(\mathcal S',\mathcal S)$ nie są równe.
\end{twierdzenie}
\begin{proof}
Implikacja ze zbieżności silnej jest natychmiastowa.
Załóżmy zbieżność słabą-*. Rodzina
$\{T_n-T:n\ge1\}$ jest punktowo ograniczona.
Z Twierdzenia~\ref{thm:dt-tempered-bounded-families}
otrzymujemy wspólne oszacowanie seminormowe tej rodziny,
a więc jej równociągłość.
Dla dowolnego niepustego ograniczonego $B\subseteq\mathcal S$
Lemat~\ref{lem:montel-schwartz22} daje zwartość
$K=\overline B$. Równociągłość funkcjonałów i
punktowa zbieżność do zera zapewniają zbieżność
jednostajną na $K$: dla danego $\varepsilon>0$
wybieramy otoczenie $V$ z
$\sup_n|\langle T_n-T,v\rangle|<\varepsilon/3$
dla $v\in V$. Możemy wybrać $V$ otwarte dzięki
wspólnemu oszacowaniu seminormowemu. Z pokrycia
$K$ zbiorami $\varphi+V$, $\varphi\in K$, wybieramy
skończone podpokrycie $\varphi_j+V$, $j=1,\ldots,r$.
Następnie dla dostatecznie dużego $n$ wszystkie
$|\langle T_n-T,\varphi_j\rangle|<\varepsilon/3$.
Dla każdego $\varphi\in B$ istnieje $j$ taki, że
$\varphi-\varphi_j\in V$; liniowość i nierówność
trójkąta dają
$|\langle T_n-T,\varphi\rangle|<2\varepsilon/3$.
Po wzięciu supremum otrzymujemy
$P_B(T_n-T)\le2\varepsilon/3<\varepsilon$.

Aby wykazać różność topologii, wybierzmy niezerową funkcję
$\psi\in\mathcal D$ o nośniku w $B(0,1)$ i zdefiniujmy
$\psi_n(x)=e^{-n^2}\psi(x-3ne_1)$ dla $n\ge1$.
Zbiór $B=\{\psi_n:n\ge1\}$ jest ograniczony w
$\mathcal S$: dla ustalonych $N,m$ mamy
$Q_{N,m}(\psi_n)\le e^{-n^2}(2+3n)^N Q_{0,m}(\psi)$,
co jest ograniczone niezależnie od $n$.
Jego liniowa otoczka jest nieskończenie wymiarowa,
ponieważ nośniki są parami rozłączne i każda $\psi_n$
jest niezerowa.
Otoczenie silne $\{T:P_B(T)<1/2\}$ nie zawiera
żadnego słabo-* otoczenia zera: dla dowolnie wielu
$\varphi_1,\ldots,\varphi_r\in\mathcal S$
można wybrać $\psi_n$ poza ich liniową otoczką.
Podprzestrzeń $M=\operatorname{span}\{\varphi_1,\ldots,
\varphi_r\}$ jest skończenie wymiarowa, a więc domknięta
w przestrzeni lokalnie wypukłej Hausdorffa $\mathcal S$.
Ponieważ $\psi_n\notin M$, twierdzenie o oddzielaniu
punktu od domkniętej podprzestrzeni
(Twierdzenie~\ref{thm:point-closed-subspace})
daje ciągły funkcjonał liniowy zerujący się na $M$
i niezerowy w $\psi_n$. Dzieląc go przez wartość
w $\psi_n$, otrzymujemy $T\in\mathcal S'$ taki, że
$T(\varphi_j)=0$ dla wszystkich $j$ oraz $T(\psi_n)=1$.
Należy on do dowolnego podstawowego słabo-* otoczenia
wyznaczonego przez te $\varphi_j$, ale nie do
$\{T:P_B(T)<1/2\}$. Ponieważ topologia silna jest
co najmniej tak mocna jak słaba-*, dowodzi to ich różności.
\end{proof}

\subsection{Rozszerzenie transformacji Fouriera}

\noindent\emph{Powiązania tematyczne w Bibliotece matematycznej:} \atlasref{PDE-08::def:pde8-distFT}{PDE-08: transformacja Fouriera dystrybucji}.


\begin{definicja}[Transformacja Fouriera dystrybucji temperowanej]
\label{def:fourier-tempered22}
Dla $T\in\mathcal S'(\R^d)$ definiujemy
$\widehat T=\mathcal F T\in\mathcal S'$ przez
\[
\langle\mathcal F T,\varphi\rangle
=\langle T,\mathcal F\varphi\rangle
\qquad(\varphi\in\mathcal S).
\]
Występuje tu transformacja Fouriera na $\mathcal S$
z Definicji~\ref{def:fourier-schwartz17}, z normalizacją
$(2\pi)^{-d/2}$ i jądrem $e^{-ix\cdot\xi}$.
Ponieważ $\mathcal F:\mathcal S\to\mathcal S$ jest ciągła
(Twierdzenie~\ref{thm:fourier-schwartz17}), złożenie
$T\circ\mathcal F$ jest ciągłym $\C$-liniowym funkcjonałem.
Definicja nie zawiera sprzężenia zespolonego ani
transformacji odwrotnej.
\end{definicja}

\begin{twierdzenie}[Transformacja Fouriera jest automorfizmem]
\label{thm:fourier-auto-tempered22}
Transformacja $\mathcal F:\mathcal S'\to\mathcal S'$
jest ciągłym izomorfizmem liniowo-topologicznym
zarówno dla topologii $\sigma(\mathcal S',\mathcal S)$,
jak i $\beta(\mathcal S',\mathcal S)$.
Jej odwrotność dana jest transpozycją
$\mathcal F^{-1}$ na $\mathcal S$:
$\langle\mathcal F^{-1}T,\varphi\rangle
=\langle T,\mathcal F^{-1}\varphi\rangle$.
Dla każdego $T\in\mathcal S'$ i $j\in\{1,\ldots,d\}$
zachodzą wzory
\[
\mathcal F(\partial_jT)=i\xi_j\,\mathcal F T,
\qquad
\mathcal F(x_jT)=i\partial_{\xi_j}(\mathcal F T),
\qquad\mathcal F^2T=T(-\cdot),
\]
gdzie $\langle T(-\cdot),\varphi\rangle=
\langle T,\varphi(-\cdot)\rangle$.
\end{twierdzenie}
\begin{proof}
Z Twierdzenia~\ref{thm:fourier-schwartz17} oba operatory
$\mathcal F$ i $\mathcal F^{-1}$ na $\mathcal S$ są
ciągłe. Dla $\varphi\in\mathcal S$ mamy
$\langle\mathcal F^{-1}\mathcal FT,\varphi\rangle
=\langle T,\mathcal F\mathcal F^{-1}\varphi\rangle
=\langle T,\varphi\rangle$.
Analogicznie
$\langle\mathcal F\mathcal F^{-1}T,\varphi\rangle
=\langle T,\mathcal F^{-1}\mathcal F\varphi\rangle
=\langle T,\varphi\rangle$.
Są to zatem wzajemnie odwrotne operatory na $\mathcal S'$.
Ciągłość słaba-* wynika z działania na pojedynczej
funkcji testowej. Dla silnej mamy
$P_B(\mathcal FT)=P_{\mathcal F(B)}(T)$;
$\mathcal F(B)$ jest ograniczony, gdy $B$ jest ograniczony,
co daje ciągłość. Odwrotność rozpatrujemy analogicznie.
Ustalmy $\varphi\in\mathcal S$, oznaczając jej argument
przez $\xi$, a argument $\mathcal F\varphi$ przez $x$.
Z Twierdzenia~\ref{thm:fourier-schwartz17} wynika
$\partial_{x_j}\mathcal F\varphi
=\mathcal F(-i\xi_j\varphi)$ oraz
$x_j\mathcal F\varphi=\mathcal F(-i\partial_{\xi_j}\varphi)$.
Zatem
\[
\begin{aligned}
\langle\mathcal F(\partial_jT),\varphi\rangle
&=-\langle T,\partial_{x_j}\mathcal F\varphi\rangle\\
&=\langle T,\mathcal F(i\xi_j\varphi)\rangle
=\langle i\xi_j\mathcal FT,\varphi\rangle,\\
\langle\mathcal F(x_jT),\varphi\rangle
&=\langle T,x_j\mathcal F\varphi\rangle\\
&=\langle T,\mathcal F(-i\partial_{\xi_j}\varphi)\rangle\\
&=-i\langle\mathcal FT,\partial_{\xi_j}\varphi\rangle
=\langle i\partial_{\xi_j}\mathcal FT,\varphi\rangle.
\end{aligned}
\]
Wzór odwrotny na $\mathcal S$ daje
$\mathcal F^{-1}\varphi(x)=\mathcal F\varphi(-x)$,
więc $\mathcal F^2\varphi(x)=\varphi(-x)$.
W konsekwencji
$\langle\mathcal F^2T,\varphi\rangle
=\langle T,\mathcal F^2\varphi\rangle
=\langle T,\varphi(-\cdot)\rangle$.
W szczególności $\mathcal F^4T=T$.
\end{proof}

\begin{przyklad}[Funkcja stała, delta Diraca i pochodne]
\label{ex:fourier-delta22}
Dla $a\in\R^d$ i $\alpha\in\N_0^d$ mamy w $\mathcal S'$
wzory
\[
\begin{aligned}
\mathcal F\delta_a&=(2\pi)^{-d/2}T_{e^{-ia\cdot\xi}},\\
\mathcal F T_1&=(2\pi)^{d/2}\delta_0,\\
\mathcal F(\partial^\alpha\delta_0)
&=(2\pi)^{-d/2}T_{(i\xi)^\alpha}.
\end{aligned}
\]
Przez $T_f$ oznaczamy dystrybucję regularną
reprezentowaną przez $f$, w szczególności $T_1$
odpowiada funkcji stałej $1$.
Pierwszy wynika z
$\langle\mathcal F\delta_a,\varphi\rangle
=(2\pi)^{-d/2}\int e^{-ia\cdot\xi}\varphi(\xi)
\,\mathrm d\lambda_d(\xi)$.
Dla $a=0$ mamy $\mathcal F\delta_0=(2\pi)^{-d/2}T_1$.
Ponieważ odbicie $\delta_0$ jest równe $\delta_0$,
zastosowanie $\mathcal F$ daje
$\delta_0=(2\pi)^{-d/2}\mathcal F T_1$.
Iterując wzór na transformatę pochodnej z
Twierdzenia~\ref{thm:fourier-auto-tempered22}, dostajemy
$\mathcal F(\partial^\alpha\delta_0)
=(i\xi)^\alpha\mathcal F\delta_0$, co daje trzeci wzór.
W szczególności delta Diraca nie jest funkcją,
ale jej transformata Fouriera jest funkcją gładką.
\end{przyklad}

\begin{stwierdzenie}[Zgodność z transformacją całkową]
\label{prop:dt-tempered-fourier-l1}
Jeżeli $f\in L^1(\R^d;\C)$, to zarówno $T_f$, jak i
$T_{\widehat f}$ są dystrybucjami temperowanymi, przy czym
\[
 \mathcal F T_f=T_{\widehat f}.
\]
Po prawej stronie $\widehat f$ jest transformatą całkową
z rozdziału o transformacji Fouriera funkcji całkowalnych.
\end{stwierdzenie}
\begin{proof}
Mamy $|\langle T_f,\varphi\rangle|
\le\|f\|_1 Q_{0,0}(\varphi)$, co daje temperowalność $T_f$.
Z Twierdzenia~\ref{thm:dt-fourier-l1-rl} funkcja $\widehat f$
jest ciągła i ograniczona, więc $T_{\widehat f}\in\mathcal S'$
na mocy Przykładu~\ref{ex:tempered-regular22}.
Dla $\varphi\in\mathcal S$ mamy, przy $c_d=(2\pi)^{-d/2}$,
\[
\begin{aligned}
 \langle\mathcal FT_f,\varphi\rangle
 &=\langle T_f,\mathcal F\varphi\rangle\\
 &=c_d\int_{\R^d}\int_{\R^d}f(x)e^{-ix\cdot\xi}
 \varphi(\xi)\,\mathrm d\lambda_d(\xi)\,\mathrm d\lambda_d(x)\\
 &=\int_{\R^d}\widehat f(\xi)\varphi(\xi)
 \,\mathrm d\lambda_d(\xi).
\end{aligned}
\]
Zamiana kolejności całkowania wynika z twierdzenia Fubiniego:
całka modułu z czynnikiem $c_d$ wynosi
$c_d\|f\|_1\|\varphi\|_1<\infty$.
Wobec tego obie definicje dają tę samą dystrybucję.
\end{proof}

\begin{stwierdzenie}[Przesunięcia i modulacje]
\label{prop:dt-tempered-fourier-shift-modulation}
Dla $h,\eta\in\R^d$ i $T\in\mathcal S'$ definiujemy
\[
 \langle\tau_h T,\varphi\rangle
 =\langle T,\varphi(\cdot+h)\rangle,\qquad
 M_\eta T=e^{i\eta\cdot x}T.
\]
Są to ciągłe automorfizmy liniowe w obu topologiach
przestrzeni $\mathcal S'$. Zachodzą tożsamości
\[
 \mathcal F(\tau_h T)=e^{-ih\cdot\xi}\mathcal FT,
 \qquad \mathcal F(M_\eta T)=\tau_\eta(\mathcal FT).
\]
\end{stwierdzenie}
\begin{proof}
Dla przesunięcia testu mamy
$Q_{N,m}(\varphi(\cdot+h))\le
(1+|h|)^NQ_{N,m}(\varphi)$, ponieważ
$1+|x|\le(1+|h|)(1+|x+h|)$.
Ciągłość mnożenia przez $e^{i\eta\cdot x}$ wynika z
Twierdzenia~\ref{thm:tempered-operators22}, gdyż wszystkie
pochodne tego mnożnika są ograniczone. Ciągłość transpozycji
wynika z tego samego twierdzenia. Operatory odwrotne to
odpowiednio $\tau_{-h}$ oraz $M_{-\eta}$.
Dla dystrybucji regularnej przesunięcie odpowiada
funkcji $x\mapsto f(x-h)$, przez podstawienie w całce.
Z definicji transformacji testów bezpośrednio otrzymujemy
\[
 (\mathcal F\varphi)(x+h)
 =\mathcal F(e^{-ih\cdot\xi}\varphi)(x),
 \qquad
 e^{i\eta\cdot x}(\mathcal F\varphi)(x)
 =\mathcal F(\varphi(\cdot+\eta))(x).
\]
Obie całki są bezwzględnie zbieżne, a druga równość
wynika z podstawienia $\zeta=\xi+\eta$.
Stosując te równości w definicji transpozycji, dostajemy
\[
\begin{aligned}
 \langle\mathcal F(\tau_hT),\varphi\rangle
 &=\langle T,(\mathcal F\varphi)(\cdot+h)\rangle
 =\langle\mathcal FT,e^{-ih\cdot\xi}\varphi\rangle,\\
 \langle\mathcal F(M_\eta T),\varphi\rangle
 &=\langle T,e^{i\eta\cdot x}\mathcal F\varphi\rangle
 =\langle\mathcal FT,\varphi(\cdot+\eta)\rangle .
\end{aligned}
\]
To dowodzi obu wzorów.
\end{proof}

\begin{przyklad}[Fale płaskie]
Dla $a,\eta\in\R^d$ mamy
\[
 \mathcal F\delta_a=c_d T_{e^{-ia\cdot\xi}},
 \qquad
 \mathcal F T_{e^{i\eta\cdot x}}=c_d^{-1}\delta_\eta,
 \qquad c_d=(2\pi)^{-d/2}.
\]
Pierwszy wzór wynika z obliczenia
$\langle\mathcal F\delta_a,\varphi\rangle
=(\mathcal F\varphi)(a)$.
W drugim używamy $\mathcal F T_1=c_d^{-1}\delta_0$
z Przykładu~\ref{ex:fourier-delta22} i poprzedniego
stwierdzenia. Fala o jednej częstotliwości ma więc
transformatę skupioną w jednym punkcie, mimo że nie
jest funkcją całkowalną na $\R^d$.
\end{przyklad}

\begin{stwierdzenie}[Odwracalna liniowa zmiana współrzędnych]
\label{prop:dt-tempered-linear-change}
Niech $A\in GL_d(\R)$ i $T\in\mathcal S'(\R^d)$.
Definiujemy $L_A T$ przez
\[
 \langle L_A T,\varphi\rangle
 =|\det A|^{-1}\langle T,\varphi\circ A^{-1}\rangle,
 \qquad \varphi\in\mathcal S.
\]
Operator $L_A$ jest automorfizmem liniowo-topologicznym
w obu topologiach $\mathcal S'$, a jego odwrotnością
jest $L_{A^{-1}}$. Jeżeli $f\in L^1_{\mathrm{loc}}(\R^d;\C)$ i dystrybucja
regularna $T_f$ jest temperowana, to $T_{f\circ A}$
jest temperowana oraz $L_A T_f=T_{f\circ A}$.
Obie dystrybucje regularne rozumiemy tu jako
jednoznaczne ciągłe przedłużenia z $\mathcal D$ na $\mathcal S$.
Zachodzi
\[
 \mathcal F(L_A T)=|\det A|^{-1}L_{A^{-T}}(\mathcal F T).
\]
\end{stwierdzenie}
\begin{proof}
Dla każdej macierzy odwracalnej $B$ operator
$C_B\varphi=\varphi\circ B$ zachowuje $\mathcal S$
i jest ciągły. Istotnie, z reguły łańcuchowej każda
pochodna rzędu $k$ funkcji $\varphi\circ B$ jest
skończoną kombinacją pochodnych rzędu $k$ funkcji
$\varphi$, obliczanych w punkcie $Bx$, ze stałymi
współczynnikami zależnymi od $B$.
Ponadto $1+|x|\le\max(1,\|B^{-1}\|)(1+|Bx|)$.
Stąd dla $N,m\in\N_0$ istnieje stała $C_{B,N,m}$ taka,
że $Q_{N,m}(C_B\varphi)\le C_{B,N,m}Q_{N,m}(\varphi)$.
Zastosowanie tego argumentu do $B^{-1}$ daje ciągłość
odwrotności. Definicja $L_A$ jest więc transpozycją
ciągłego operatora z czynnikiem skalarnym; ciągłość
w obu topologiach wynika z argumentu w dowodzie
Twierdzenia~\ref{thm:tempered-operators22}.
Bezpośrednie złożenie definicji dla $A$ i $A^{-1}$
daje operator identycznościowy w obu kolejnościach.
Dla $f\in L^1_{\mathrm{loc}}$ podstawienie $y=Ax$
na zbiorach zwartych pokazuje lokalną całkowalność
$f\circ A$, a dla $\varphi\in\mathcal D$ daje
\[
 \langle L_A T_f,\varphi\rangle
 =|\det A|^{-1}\int_{\R^d}f(y)\varphi(A^{-1}y)
 \,\mathrm d\lambda_d(y)
 =\int_{\R^d}f(Ax)\varphi(x)\,\mathrm d\lambda_d(x).
\]
Zatem $L_A T_f$ jest ciągłym przedłużeniem dystrybucji
regularnej $T_{f\circ A}$ na $\mathcal S$.
Jednoznaczność wynika z Twierdzenia~\ref{thm:tempered-estimate22}.
Nie twierdzimy, że całka z $f\varphi$ jest bezwzględnie
zbieżna dla każdego $\varphi\in\mathcal S$, gdy jedynym
założeniem o wzroście $f$ jest temperowalność $T_f$.

Z reguły zmiany zmiennych w transformacji całkowej
(Stwierdzenie~\ref{prop:dt-fourier-l1-rules}) wynika
\[
 \mathcal F(\varphi\circ A^T)(x)
 =|\det A|^{-1}(\mathcal F\varphi)(A^{-1}x).
\]
Wobec tego
\[
\begin{aligned}
 \langle\mathcal F(L_AT),\varphi\rangle
 &=|\det A|^{-1}
 \langle T,(\mathcal F\varphi)\circ A^{-1}\rangle\\
 &=\langle T,\mathcal F(\varphi\circ A^T)\rangle\\
 &=\langle\mathcal F T,\varphi\circ A^T\rangle\\
 &=|\det A|^{-1}\langle L_{A^{-T}}(\mathcal FT),\varphi\rangle .
\end{aligned}
\]
Ostatni krok wykorzystuje
$|\det A^{-T}|^{-1}=|\det A|$. Dowodzi to tezy.
\end{proof}

\begin{przyklad}[Transformata wielomianu]
\label{ex:dt-tempered-polynomial}
Dla $\alpha\in\N_0^d$ zachodzi
\[
 \mathcal F T_{x^\alpha}
 =(2\pi)^{d/2}i^{|\alpha|}\partial^\alpha\delta_0.
\]
Wielomian wyznacza dystrybucję temperowaną na mocy
Przykładu~\ref{ex:tempered-regular22}.
Z Przykładu~\ref{ex:fourier-delta22} mamy
$\mathcal FT_1=(2\pi)^{d/2}\delta_0$.
Ponieważ $x^\alpha T_1=T_{x^\alpha}$, iteracja
wzoru na transformatę mnożenia przez współrzędną
z Twierdzenia~\ref{thm:fourier-auto-tempered22}
daje podaną równość.
\end{przyklad}

\subsection{Splot z funkcją szybko malejącą}

\begin{twierdzenie}[Splot z elementem przestrzeni Schwartza]
\label{thm:schwartz-convolution22}
Niech $T\in\mathcal S'(\R^d)$ i
$\rho\in\mathcal S(\R^d)$. Funkcja
$u(x)=\langle T_y,\rho(x-y)\rangle$ należy do
$C^\infty(\R^d)$, a każda jej pochodna ma
co najwyżej wzrost wielomianowy. Dystrybucja
regularna $T_u$ jest temperowana i oznaczamy ją $T*\rho$.
Dla każdego $\alpha\in\N_0^d$ zachodzi
\[
\mathcal F(T*\rho)=(2\pi)^{d/2}(\mathcal FT)\,(\mathcal F\rho),
\qquad \partial^\alpha(T*\rho)=T*(\partial^\alpha\rho).
\]
\end{twierdzenie}
\begin{proof}
Dla $h\in\R^d$ oraz $N,m\in\N_0$ zachodzi
$Q_{N,m}(\rho(h-\cdot))\le(1+|h|)^N Q_{N,m}(\rho)$,
gdyż $1+|y|\le(1+|h|)(1+|h-y|)$.
Ta sama nierówność obowiązuje dla każdej pochodnej
$\rho$. Z twierdzenia Newtona--Leibniza zastosowanego
dwukrotnie, dla $0<|t|\le1$ i $j\in\{1,\ldots,d\}$,
otrzymujemy
\[
\begin{aligned}
 Q_{N,m}\left(
 \frac{\rho(x+te_j-\cdot)-\rho(x-\cdot)}{t}
 -(\partial_j\rho)(x-\cdot)\right)
 \le |t|(2+|x|)^N Q_{N,m+2}(\rho).
\end{aligned}
\]
Istotnie, po każdej pochodnej względem $y$ resztę
ograniczamy przez $|t|$ razy supremum odpowiedniej
pochodnej drugiego rzędu względem parametru na
odcinku od $x$ do $x+te_j$ i stosujemy poprzednią
nierówność wagową. Jednokrotne zastosowanie twierdzenia
Newtona--Leibniza daje również ciągłość przesunięć
w każdej seminormie. Powtarzając argument dla
pochodnych $\rho$, otrzymujemy gładkość funkcji
$x\mapsto\rho(x-\cdot)$ o wartościach w $\mathcal S$.
Ciągłość $T$ pozwala więc różniczkować pod parowaniem:
\[
 \partial^\alpha u(x)=\langle T_y,
 (\partial^\alpha\rho)(x-y)\rangle.
\]
Z Twierdzenia~\ref{thm:tempered-estimate22} wybieramy
$C>0$, $N,m\in\N_0$ z $|T(\varphi)|\le C Q_{N,m}(\varphi)$.
Poprzednie oszacowanie przesunięcia daje
\[
 |\partial^\alpha u(x)|\le
 C(1+|x|)^N Q_{N,m+|\alpha|}(\rho).
\]
W szczególności $T_u\in\mathcal S'$ na mocy
Przykładu~\ref{ex:tempered-regular22}.

Ustalmy $\psi\in\mathcal S$ i niech $G_x(y)=
\psi(x)\rho(x-y)$. Odwzorowanie $x\mapsto G_x$
jest ciągłe w $\mathcal S$ i
\[
 Q_{N,m}(G_x)\le |\psi(x)|(1+|x|)^N Q_{N,m}(\rho).
\]
Prawa strona jest całkowalna po $\R^d$ dla każdych $N,m$.
Dla $R\ge1$ określmy $I_R(y)=\int_{[-R,R]^d}G_x(y)
\,\mathrm d\lambda_d(x)$. Różniczkowanie pod całką na
tym zwartym prostokącie jest dozwolone, a powyższe
oszacowanie dowodzi, że $I_R\in\mathcal S$.
Sumy Riemanna tej całki zbiegają do $I_R$ w każdej
seminormie $Q_{N,m}$: błąd jest nie większy niż
objętość prostokąta razy supremum
$Q_{N,m}(G_x-G_{x'})$ dla $x,x'$ z tej samej części
podziału. Supremum dąży do zera z jednostajnej
ciągłości w danej seminormie na prostokącie.
Stosując $T$ do sum Riemanna i przechodząc do granicy
otrzymujemy
\[
 \langle T,I_R\rangle=
 \int_{[-R,R]^d}\psi(x)u(x)\,\mathrm d\lambda_d(x).
\]
Dla $R'<R''$ oszacowanie ogona ma postać
\[
 Q_{N,m}(I_{R''}-I_{R'})
 \le Q_{N,m}(\rho)
 \int_{\R^d\setminus[-R',R']^d}
 (1+|x|)^N|\psi(x)|\,\mathrm d\lambda_d(x).
\]
W szczególności $I_n$ jest ciągiem Cauchy'ego w $\mathcal S$.
Z jej zupełności granica istnieje w $\mathcal S$, a
zbieżność punktowa identyfikuje ją z funkcją
$\check\rho*\psi$, gdzie $\check\rho(y)=\rho(-y)$:
\[
 (\check\rho*\psi)(y)=
 \int_{\R^d}\rho(x-y)\psi(x)\,\mathrm d\lambda_d(x).
\]
Z ciągłości $T$ oraz bezwzględnej całkowalności
$u\psi$ otrzymujemy
$\langle T_u,\psi\rangle=\langle T,\check\rho*\psi\rangle$.
Stwierdzenie~\ref{prop:convolution-schwartz17} zapewnia
ciągłość operatora $\psi\mapsto\check\rho*\psi$ na
$\mathcal S$. Jego transpozycją jest więc ciągły
operator $T\mapsto T*\rho$ dla obu topologii na
$\mathcal S'$, zgodnie z dowodem
Twierdzenia~\ref{thm:tempered-operators22}.
Jeśli $\rho\in\mathcal D$, konstrukcja zgadza się na
$\mathcal D$ z Twierdzeniem~\ref{thm:smoothing21}.

Pozostaje obliczyć stałą i znak w formule Fouriera.
Dla $\psi\in\mathcal S$ mamy, przy $c_d=(2\pi)^{-d/2}$,
\[
\begin{aligned}
(\check\rho*\mathcal F\psi)(y)
 &=\int_{\R^d}\rho(z)(\mathcal F\psi)(y+z)
 \,\mathrm d\lambda_d(z)\\
 &=c_d\int_{\R^d}e^{-iy\cdot\xi}\psi(\xi)
 \int_{\R^d}\rho(z)e^{-iz\cdot\xi}\,\mathrm d\lambda_d(z)
 \,\mathrm d\lambda_d(\xi)\\
 &=(2\pi)^{d/2}\mathcal F((\mathcal F\rho)\psi)(y).
\end{aligned}
\]
Twierdzenie Fubiniego stosujemy, ponieważ całka
modułu podwójnego jest ograniczona przez
$c_d\|\rho\|_1\|\psi\|_1$. Iloczyn
$(\mathcal F\rho)\psi$ należy do $\mathcal S$.
Zatem
\[
\begin{aligned}
\langle\mathcal F(T*\rho),\psi\rangle
&=\langle T,\check\rho*\mathcal F\psi\rangle\\
&=(2\pi)^{d/2}\langle\mathcal FT,
 (\mathcal F\rho)\psi\rangle.
\end{aligned}
\]
To jest pierwszy wzór, a drugi wynika ze wzoru na
$\partial^\alpha u$ i zgodności pochodnej dystrybucji
regularnej z klasyczną pochodną gładkiej funkcji
(wystarczy sprawdzić na $\mathcal D$, gęstej w $\mathcal S$).
\end{proof}

\begin{twierdzenie}[Aproksymacja dystrybucji temperowanych funkcjami testowymi]
\label{thm:dt-tempered-test-density}
Dla każdej $T\in\mathcal S'(\R^d)$ istnieje ciąg
$(f_n)\subseteq\mathcal D(\R^d)$ taki, że
$T_{f_n}\to T$ w silnej topologii $\beta(\mathcal S',\mathcal S)$.
Co więcej, ciąg można wybrać tak, by dla każdego
niepustego ograniczonego $B\subseteq\mathcal S$
zachodziło $P_B(T_{f_n}-T)=O(n^{-1})$.
\end{twierdzenie}
\begin{proof}
Niech $\chi\in\mathcal D(\R^d;\R)$ będzie równa $1$
na $\overline B(0,1)$, $0\le\chi\le1$,
$\supp\chi\subseteq B(0,2)$, jak w dowodzie
Twierdzenia~\ref{thm:density-schwartz17}.
Piszemy $\chi_n(x)=\chi(x/n)$. Wybieramy standardową
aproksymację jedności $\rho_{1/n}$ z
Definicji~\ref{def:mollifier21}.
Dystrybucja $\chi_nT$ ma zwarty nośnik, ponieważ jej
ograniczenie do $\mathcal D$ znika poza $\supp\chi_n$
na mocy definicji mnożenia przez funkcję gładką.
Z Twierdzenia~\ref{thm:smoothing21} splot
\[
 T_{f_n}=(\chi_nT)*\rho_{1/n}
\]
jest reprezentowany przez funkcję testową $f_n$,
gdyż jego nośnik zawiera się w zwartym zbiorze
$\supp\chi_n+\supp\rho_{1/n}$.
Z dowodu Twierdzenia~\ref{thm:schwartz-convolution22},
dla $\varphi\in\mathcal S$, mamy
\[
 \langle T_{f_n}-T,\varphi\rangle
 =\langle T,\chi_n(\check\rho_{1/n}*\varphi)-\varphi\rangle.
\]
Niech $N,m\in\N_0$. Dla $0<\varepsilon\le1$,
z normalizacji $\rho$ i podstawienia zmiennej otrzymujemy
\[
 (\check\rho_\varepsilon*\varphi)(x)-\varphi(x)
 =\int_{\R^d}\rho(z)
 [\varphi(x+\varepsilon z)-\varphi(x)]\,\mathrm d\lambda_d(z).
\]
Wzór pozostaje prawdziwy po zastosowaniu każdej pochodnej
$\partial^\alpha$ do obu stron, ponieważ całkujemy po
zwartym nośniku $\rho$, a pochodne $\varphi$ są ograniczone.
Dla $|z|\le1$, $0\le t\le1$ zachodzi
$1+|x|\le2(1+|x+t\varepsilon z|)$.
Stosując twierdzenie Newtona--Leibniza do różnicy
$\partial^\alpha\varphi(x+\varepsilon z)
-\partial^\alpha\varphi(x)$, dla $|\alpha|\le m$ dostajemy
\[
 Q_{N,m}(\check\rho_\varepsilon*\varphi-\varphi)
 \le 2^N\varepsilon Q_{N,m+1}(\varphi)
 \int_{\R^d}\|z\|_1\rho(z)\,\mathrm d\lambda_d(z),
\]
gdzie $\|z\|_1=\sum_{j=1}^d|z_j|$.
Z reguły Leibniza dla $|\alpha|\le m$ mamy
\[
 (1+|x|)^N|\partial^\alpha(\chi_n g)(x)|
 \le Q_{N,m}(g)\sum_{\gamma\le\alpha}
 \binom{\alpha}{\gamma}n^{-|\gamma|}
 \|\partial^\gamma\chi\|_\infty.
\]
Ponieważ $n^{-|\gamma|}\le1$ dla $n\ge1$, po wzięciu
supremum i maksimum otrzymujemy
$Q_{N,m}(\chi_n g)\le C_m Q_{N,m}(g)$
dla $g\in\mathcal S$, gdzie $C_m$ nie zależy od $n$. Z dowodu
Twierdzenia~\ref{thm:density-schwartz17} mamy też
$Q_{N,m}(\chi_n\varphi-\varphi)
\le C_m'(1+n)^{-1}Q_{N+1,m}(\varphi)$.
Rozkładając różnicę w parowaniu na
$\chi_n(\check\rho_{1/n}*\varphi-\varphi)
+(\chi_n\varphi-\varphi)$, otrzymujemy zatem
\[
\begin{aligned}
 &Q_{N,m}(\chi_n(\check\rho_{1/n}*\varphi)-\varphi)\\
 &\qquad\le \frac{C_{N,m,\rho,\chi}}{n}
 \bigl(Q_{N,m+1}(\varphi)+Q_{N+1,m}(\varphi)\bigr).
\end{aligned}
\]
Z Twierdzenia~\ref{thm:tempered-estimate22} wybieramy
$N,m,C$ takie, że $|T(g)|\le C Q_{N,m}(g)$.
Dla dowolnego niepustego zbioru ograniczonego
$B\subseteq\mathcal S$ obie seminormy po prawej stronie
ostatniego oszacowania są jednostajnie ograniczone na $B$.
Dokładniej, dla każdego takiego $B$ istnieje stała
$C_{T,B,\rho,\chi}\ge0$ niezależna od $n$, taka że
\[
 P_B(T_{f_n}-T)\le \frac{C_{T,B,\rho,\chi}}{n}
 \qquad(n\ge1).
\]
W szczególności $P_B(T_{f_n}-T)\to0$.
Jest to dokładnie zbieżność silna.
\end{proof}

\subsection{Łańcuch przestrzeni sprzężonych i kryterium temperowania}

\begin{lemat}[Transpozycja gęstego zanurzenia]
\label{lem:transpozycja25}
Niech $X,Y$ będą przestrzeniami lokalnie wypukłymi
Hausdorffa, a $j:X\to Y$ ciągłym odwzorowaniem liniowym
o gęstym obrazie. Wówczas transpozycja
$j':Y'\to X'$, $j'(T)=T\circ j$, jest liniowa,
różnowartościowa i ciągła zarówno dla topologii słabych-*, jak i dla
silnych topologii przestrzeni dualnych.
\end{lemat}
\begin{proof}
Dla każdego $T\in Y'$ złożenie $T\circ j$ jest ciągłym
funkcjonałem liniowym na $X$, więc $j'$ jest dobrze
określone; liniowość wynika z liniowości składania
względem pierwszego czynnika. Jeżeli $T\circ j=0$,
to $T=0$ przez gęstość $j(X)$ i ciągłość $T$.
Dla $x\in X$ mamy
$|\langle j'T,x\rangle|=|\langle T,jx\rangle|$,
co daje ciągłość w topologiach słabych-*.
Jeśli $B\subseteq X$ jest ograniczony, to $j(B)$ jest
ograniczony w $Y$ (dla każdej ciągłej seminormy $p$
na $Y$ seminorma $p\circ j$ jest ciągła na $X$).
Dla seminorm silnych zachodzi dokładna równość
$\sup_{x\in B}|\langle j'T,x\rangle|
= \sup_{y\in j(B)}|\langle T,y\rangle|$.
To dowodzi ciągłości w topologiach silnych.
\end{proof}

\begin{wniosek}[Łańcuch przestrzeni sprzężonych]
\label{cor:dual-lancuch25}
Dla każdego $d\ge1$ zachodzą kanoniczne, ciągłe i
różnowartościowe odwzorowania liniowe
\[
 \mathcal E'(\R^d)\longrightarrow\mathcal S'(\R^d)
 \longrightarrow\mathcal D'(\R^d),
\]
zarówno w topologiach słabych-*, jak i silnych.
Dla każdego niepustego zbioru otwartego $\Omega\subseteq\R^d$
kanoniczne odwzorowanie
$\mathcal E'(\Omega)\to\mathcal D'(\Omega)$
ma te same własności. Samo stwierdzenie ciągłości i
różnowartościowości \emph{nie oznacza}, że są to zanurzenia
topologiczne na obrazy z topologią podprzestrzeni.
\end{wniosek}
\begin{proof}
Stosujemy Lemat~\ref{lem:transpozycja25} do odwzorowań
z Twierdzenia~\ref{thm:zanurzenia25}. Jeżeli $i:X\to Y$
oraz $j:Y\to Z$ są tymi kolejnymi odwzorowaniami,
to $(j\circ i)'=i'\circ j'$, gdyż dla $T\in Z'$
mamy $(T\circ j)\circ i=T\circ(j\circ i)$.
Dlatego ograniczenie z $\mathcal E'$ do $\mathcal D'$
przez $\mathcal S'$ jest tym samym odwzorowaniem co
bezpośrednie ograniczenie. W przypadku zespolonym
transpozycja jest liniowa, bez sprzężenia zespolonego.
Ciągłość i różnowartościowość nie zapewniają ciągłości
odwrotności na obrazie; poniżej podamy kontrprzykłady.
\end{proof}
\subsection{Dystrybucje temperowane jako podprzestrzeń \texorpdfstring{$\mathcal D'$}{D'}}

\begin{twierdzenie}[Kryterium temperowania]
\label{thm:temp-criterion25}
Niech $d\ge1$ będzie liczbą całkowitą, $\K\in\{\R,\C\}$
i $T\in\mathcal D'(\R^d)$ nad $\K$.
Następujące warunki są równoważne:
\begin{enumerate}[label=\textup{(\roman*)}]
\item $T$ należy do obrazu kanonicznego odwzorowania
$\mathcal S'(\R^d)\to\mathcal D'(\R^d)$;
\item istnieją $C>0$, $N,m\in\N_0$ takie, że
$|\langle T,\varphi\rangle|\le C p_{N,m}(\varphi)$
dla każdej $\varphi\in\mathcal D(\R^d)$.
\end{enumerate}
Jeżeli te warunki zachodzą, przedłużenie
$\widetilde T\in\mathcal S'(\R^d)$ jest jednoznaczne
i spełnia $|\langle\widetilde T,f\rangle|\le C p_{N,m}(f)$
dla wszystkich $f\in\mathcal S(\R^d)$, z tymi samymi
stałymi co w~\textup{(ii)}. W szczególności każda
dystrybucja o zwartym nośniku jest temperowana.
\end{twierdzenie}
\begin{proof}
Niech najpierw $\widetilde T\in\mathcal S'$ będzie
przedłużeniem $T$. Z ciągłości w zerze i współkońcowości
rosnącego ciągu norm $p_{j,j}$ istnieją $j\in\N_0$ oraz
$\delta>0$ takie, że $p_{j,j}(f)<\delta$ implikuje
$|\widetilde T(f)|<1$. Dla $f\ne0$ stosujemy ten warunek
do $\delta f/(2p_{j,j}(f))$, otrzymując
$|\widetilde T(f)|\le (2/\delta)p_{j,j}(f)$.
Dla $f=0$ nierówność jest oczywista. Ograniczenie
do $\mathcal D$ daje warunek~\textup{(ii)}.
Jest to argument z Twierdzenia~\ref{thm:tempered-estimate22},
działający nad obydwoma ciałami.

Załóżmy teraz~\textup{(ii)}. Dla każdego $f\in\mathcal S$
wybieramy ciąg $(\varphi_j)$ w $\mathcal D$ zbieżny do $f$
w $\mathcal S$ (Twierdzenie~\ref{thm:zanurzenia25}).
Ponieważ
\[
 |T(\varphi_j)-T(\varphi_k)|
 \le C p_{N,m}(\varphi_j-\varphi_k)\longrightarrow0
 \quad(j,k\to\infty),
\]
ciąg wartości ma granicę w zupełnym ciele $\K$.
Kładziemy $\widetilde T(f)=\lim_jT(\varphi_j)$.
Dla drugiego ciągu $(\psi_j)$ zbieżnego do $f$ mamy
$|T(\varphi_j)-T(\psi_j)|\le
C p_{N,m}(\varphi_j-\psi_j)\to0$, więc definicja nie zależy
od wyboru ciągu. Dla $f,g\in\mathcal S$ i $a,b\in\K$
można aproksymować $af+bg$ przez $a\varphi_j+b\psi_j$,
gdzie teraz $\varphi_j\to f$ i $\psi_j\to g$.
Przechodząc do granicy, otrzymujemy liniowość.
Ponadto $|p_{N,m}(\varphi_j)-p_{N,m}(f)|
\le p_{N,m}(\varphi_j-f)$, więc
\[
 |\widetilde T(f)|\le C p_{N,m}(f).
\]
Daje to ciągłość; na $\mathcal D$ wystarczy wziąć
ciąg stały, aby uzyskać $\widetilde T|_{\mathcal D}=T$.
Jednoznaczność wynika z gęstości $\mathcal D$.

Jeżeli $T$ ma zwarty nośnik, to z
Twierdzenia~\ref{thm:compact-support-E25} otrzymujemy
$|T(\varphi)|\le C p_{K,m}(\varphi)\le C p_{0,m}(\varphi)$
na funkcjach testowych, a zatem~\textup{(ii)} zachodzi z $N=0$.
Przedłużenie jest ograniczeniem funkcjonału na $\mathcal E$
do $\mathcal S$, przez jednoznaczność; odpowiada to
Stwierdzeniu~\ref{prop:schwartz-embeddings22}.
\end{proof}

\begin{przyklad}[Dystrybucja nietemperowana]
\label{ex:non-tempered25}
Na $\R$ funkcja $f(x)=e^x$ jest lokalnie całkowalna,
więc wyznacza dystrybucję regularną $T_f\in\mathcal D'(\R)$.
Nie jest ona temperowana. Wybierzmy nieujemną, niezerową
$\psi\in\mathcal D((-1,1))$, przedłużoną zerem na $\R$,
i niech $\varphi_n(x)=e^{-n/2}\psi(x-n)$ dla $n\ge1$.
Dla $x\in\supp\varphi_n$ mamy $|x|\le n+1$; pochodne
$\varphi_n$ znikają poza tym nośnikiem. Stąd dla wszystkich
$N,m\in\N_0$
\[
 p_{N,m}(\varphi_n)\le e^{-n/2}(n+2)^N p_{0,m}(\psi)
 \longrightarrow0,
\]
gdyż funkcja wykładnicza rośnie szybciej niż każdy
wielomian. Jednak po podstawieniu $y=x-n$ otrzymujemy
\[
 \langle T_f,\varphi_n\rangle=
 e^{n/2}\int_\R e^y\psi(y)\,\mathrm d\lambda_1(y)
 \longrightarrow+\infty.
\]
Całka jest skończona i ściśle dodatnia: $\psi$ jest
nieujemna i dodatnia na pewnym niepustym przedziale.
Ciągłość przedłużenia na $\mathcal S$ wymagałaby
$\langle T_f,\varphi_n\rangle\to0$, co daje sprzeczność.
Przykład działa nad $\R$ i $\C$; razem z
Przykładem~\ref{ex:non-tempered22} pokazuje on dwie
metody wykrywania braku temperowalności (przesuwanie
nośnika oraz powiększanie odcięcia).
\end{przyklad}
\begin{przyklad}[Zbieżność i ucieczka nośnika]
\label{ex:escaping-deltas25}
Na $\R$ niech $\delta_n$ oznacza deltę Diraca w punkcie
$n$ dla $n\ge1$. Każda z tych dystrybucji ma zwarty
nośnik, więc należy także do $\mathcal E'$ i $\mathcal S'$
na mocy Twierdzeń~\ref{thm:compact-support-E25}
i~\ref{thm:temp-criterion25}.
Wtedy $\delta_n\to0$ silnie w $\mathcal D'(\R)$:
dla ograniczonego $B\subseteq\mathcal D(\R)$ nośniki
wszystkich jego elementów mieszczą się w jednym zwartym
przedziale (Twierdzenie~\ref{thm:D-bounded18}), więc
$\sup_{\varphi\in B}|\varphi(n)|=0$ dla dużych $n$.
Również $\delta_n\to0$ silnie w $\mathcal S'(\R)$, bo dla
ograniczonego zbioru $B\subseteq\mathcal S$ mamy
$C_B=\sup_{\varphi\in B}p_{2,0}(\varphi)<\infty$ oraz
$\sup_{\varphi\in B}|\varphi(n)|\le C_B(1+n)^{-2}$.
Jednak $\delta_n\not\to0$ nawet słabo-* w
$\mathcal E'(\R)$, ponieważ $\langle\delta_n,1\rangle=1$.

Z kolei $T_n=e^n\delta_n$ zbiega silnie do zera w
$\mathcal D'(\R)$ (na każdym ograniczonym zbiorze
seminorma silna jest od pewnego indeksu zerowa).
Nie zbiega do zera nawet słabo-* w $\mathcal S'(\R)$.
Niech $r(x)=\sqrt{1+x^2}$ i $\varphi(x)=e^{-r(x)/2}$.
Wykażemy, że $\varphi\in\mathcal S(\R)$.
Dla każdego $k\in\N_0$ istnieje wielomian $P_k(x,z)$
taki, że
\[
 \varphi^{(k)}(x)=e^{-r(x)/2}r(x)^{-2k}P_k(x,r(x)).
\]
Dla $k=0$ przyjmujemy $P_0=1$. Ponieważ $r'=x/r$,
różniczkowanie wzoru daje rekurencję
\[
 P_{k+1}=z^2\partial_x P_k+xz\partial_z P_k
          -2kxP_k-\tfrac12 xzP_k.
\]
Zatem wszystkie $P_k$ są wielomianami. Ponieważ
$1\le r(x)\le1+|x|$ i $r(x)\ge|x|$, każda pochodna
jest ograniczona przez stałą razy wielomianowy wzrost
$(1+|x|)^{M_k}$ pomnożony przez $e^{-|x|/2}$.
Wszystkie seminormy Schwartza są więc skończone.
Jednak dla $n\ge1$ mamy $r(n)\le n+1$, toteż
$\langle T_n,\varphi\rangle=e^{n-r(n)/2}
\ge e^{(n-1)/2}\to+\infty$.

Dla dowolnego niepustego zbioru otwartego $\Omega$
wybierzmy wyczerpanie $(K_j)$ i punkty
$a_j\in\Omega\setminus K_j$.
Wówczas $\delta_{a_j}\to0$ silnie w $\mathcal D'(\Omega)$,
gdyż każdy ograniczony zbiór testów ma wspólny zwarty
nośnik zawarty w pewnym $K_J$.
Natomiast $\langle\delta_{a_j},1\rangle=1$ na
$\mathcal E(\Omega)$ dla wszystkich $j$.
Zatem wszystkie odwzorowania z
Wniosku~\ref{cor:dual-lancuch25} nie są izomorfizmami
topologicznymi na swoje obrazy z topologią podprzestrzeni,
ani dla topologii słabych-*, ani dla silnych.
\end{przyklad}

\subsection{Zgodność operacji i granice ich działania}

\begin{stwierdzenie}[Różniczkowanie i mnożenie]
\label{prop:ops25}
Niech $\Omega\subseteq\R^d$ będzie niepustym zbiorem otwartym.
Dla każdego wielowskaźnika $\alpha\in\N_0^d$
operator $\partial^\alpha$ działa liniowo i ciągle na
$\mathcal E(\Omega)$ i $\mathcal D(\Omega)$, a przez
transpozycję także na $\mathcal E'(\Omega)$ i
$\mathcal D'(\Omega)$ w topologiach słabych-* i silnych.
Dla $a\in C^\infty(\Omega)$ mnożenie $f\mapsto af$
i odpowiadające mu mnożenie dystrybucji działają ciągle
na tych przestrzeniach.
Różniczkowanie działa liniowo i ciągle również na
$\mathcal S(\R^d)$ i na $\mathcal S'(\R^d)$ w obu
topologiach dualnych. Na tych przestrzeniach
te same własności mnożenia zachodzą dla
$a\in C^\infty(\R^d)$, jeżeli dla każdego
wielowskaźnika $\alpha$ istnieją $C_\alpha>0$ i
$N_\alpha\in\N_0$ takie, że
$|\partial^\alpha a(x)|\le C_\alpha(1+|x|)^{N_\alpha}$
dla wszystkich $x\in\R^d$.
\end{stwierdzenie}
\begin{proof}
Dla $\mathcal E$ i $\mathcal D$ oszacowania seminorm
dla różniczkowania i mnożenia wynikają z reguły Leibniza,
jak w stwierdzeniu~\ref{prop:operations-cinf17}
i twierdzeniach~\ref{thm:dist-deriv-cont20}, \ref{thm:multiply20}.
Dla $\mathcal S$ i $\mathcal S'$ stosujemy stwierdzenie
\ref{prop:schwartz-operators17} oraz twierdzenie
\ref{thm:tempered-operators22}, w tym oszacowanie
mnożenia przez funkcję o pochodnych wzrostu wielomianowego.
Dla ciągłego operatora $A:X\to X$ transpozycja spełnia
$P_B(A'T)=P_{A(B)}(T)$, gdzie
$P_B(T)=\sup_{f\in B}|T(f)|$ i $A(B)$ jest ograniczony.
Daje to ciągłość silną, a dla singletonów także słabą-*.
Stosujemy ten argument w szczególności do
$X=\mathcal E$ i operatorów $(-1)^{|\alpha|}\partial^\alpha$
oraz $f\mapsto af$. Operatory na funkcjach komutują
z kanonicznymi odwzorowaniami przestrzeni funkcji,
więc przez transpozycję odpowiadające operacje
na przestrzeniach sprzężonych komutują z ograniczeniami
z wniosku~\ref{cor:dual-lancuch25} (w przypadku $\mathcal S'$
przy podanym warunku wzrostu mnożnika).
\end{proof}

\begin{przyklad}[Nie każdy mnożnik gładki zachowuje przestrzeń Schwartza]
\label{ex:not-multiplier25}
Funkcja $a(x)=e^{|x|^2}$ jest gładka, ale mnożenie przez nią
nie zachowuje przestrzeni Schwartza: dla
$f(x)=e^{-|x|^2}\in\mathcal S(\R^d)$ iloczyn $af=1$
nie należy do $\mathcal S(\R^d)$. Przynależność funkcji Gaussa do $\mathcal S$ wynika
z tego, że każda jej pochodna jest wielomianem
pomnożonym przez tę funkcję (indukcja z reguły
iloczynu), a spadek Gaussa dominuje każdą potęgę.
Funkcja $1$ nie spełnia warunku $p_{1,0}(1)<\infty$.
W przestrzeniach $\mathcal D$ i $\mathcal E$
mnożenie przez $a$ jest natomiast ciągłe.
Co więcej, zwykła dystrybucja $aT_1=T_a$ nie jest
temperowana, choć $T_1$ jest temperowana
(Przykłady~\ref{ex:tempered-regular22}
i~\ref{ex:non-tempered22}).
\end{przyklad}

\begin{uwaga}[Transformacja Fouriera i splot]
\label{rem:fourier-conv25}
Nad $\C$ transformacja Fouriera jest automorfizmem
liniowo-topologicznym $\mathcal S(\R^d)$ i
$\mathcal S'(\R^d)$ w obu topologiach dualnych
(Twierdzenia~\ref{thm:fourier-schwartz17}
i~\ref{thm:fourier-auto-tempered22}).
Nie jest zdefiniowana przez tę samą transpozycję
na całej $\mathcal D'$: niezerowa funkcja testowa
nie może mieć transformaty Fouriera o zwartym nośniku.
Istotnie, dla $f\in\mathcal D$ jej transformata
jest rzeczywiście analityczna: rozwinięcie wykładnika
w szereg potęgowy w otoczeniu dowolnego punktu
można całkować wyraz po wyrazie, ponieważ $f$ ma
zwarty nośnik, a szereg jest jednostajnie bezwzględnie
zbieżny na nim dla parametrów z każdego ograniczonego
otoczenia. Gdyby transformata miała zwarty nośnik,
znikałaby na niepustym otwartym podzbiorze $\R^d$,
a z jednoznaczności funkcji rzeczywiście analitycznej
znikałaby na całym spójnym $\R^d$.
Różnowartościowość transformaty na $\mathcal S$
dawałaby $f=0$. Zatem wyrażenie
$\langle T,\mathcal F f\rangle$ nie musi mieć sensu
dla $T\in\mathcal D'$ i $f\in\mathcal D$.

Dla ustalonego $T\in\mathcal E'(\R^d)$ splot
$S\mapsto T*S$ jest ciągłym operatorem liniowym
na $\mathcal D'(\R^d)$ w obu topologiach
(Twierdzenie~\ref{thm:convolution-properties21}).
Aby sprawdzić, że zachowuje również $\mathcal S'$,
niech $L$ będzie zwartym zbiorem i $C>0$, $m\in\N_0$
stałymi oszacowania przedłużenia $T$ na $\mathcal E$
z Twierdzenia~\ref{thm:compact-support-E25}.
Dla $f\in\mathcal S$ określamy
$(A_Tf)(y)=\langle T_x,f(x+y)\rangle$,
korzystając z przedłużenia na $\mathcal E$.
Z ciągłości $T$ na $\mathcal E$ i różniczkowania
przesunięć w jej topologii otrzymujemy
$\partial_y^\beta A_T f(y)=
\langle T_x,\partial^\beta f(x+y)\rangle$:
ilorazy różnicowe wraz z pochodnymi względem $x$
zbiegają jednostajnie na każdym zwartym zbiorze
przez wzór całkowy na różnicę funkcji gładkiej.
Ponieważ $1+|y|\le(1+|x|)(1+|x+y|)$, dla
$R=\max_{x\in L}|x|$ otrzymujemy
\[
 p_{N,k}(A_Tf)\le C(1+R)^N p_{N,k+m}(f).
\]
Zatem $A_T:\mathcal S\to\mathcal S$ jest liniowe
i ciągłe. Jego transpozycja
$\langle T*S,f\rangle=\langle S,A_Tf\rangle$
jest ciągła na $\mathcal S'$ w obu topologiach
i zgadza się na $\mathcal D$ z definicją splotu
z rozdziału~\ref{sec:splot21}.
Nie wynika stąd istnienie splotu dwóch dowolnych
dystrybucji (Uwaga~\ref{rem:convolution-limits21}).
\end{uwaga}

\begin{twierdzenie}[Transformata splotu z dystrybucją o zwartym nośniku]
\label{thm:dt-tempered-compact-convolution}
Niech $U\in\mathcal E'(\R^d)$ oraz $T\in\mathcal S'(\R^d)$.
Definiujemy funkcję
\[
 a(\xi)=c_d\langle U_x,e^{-ix\cdot\xi}\rangle,
 \qquad c_d=(2\pi)^{-d/2},
\]
korzystając z przedłużenia działania $U$ na $\mathcal E$.
Wtedy $a$ jest gładkim mnożnikiem przestrzeni Schwartza,
$\mathcal F U=T_a$, a
\[
 \mathcal F(U*T)=c_d^{-1}a\,\mathcal F T
 \quad\text{w }\mathcal S'.
\]
Splot po lewej stronie jest zgodny ze splotem
dystrybucji z Rozdziału~\ref{sec:splot21}.
\end{twierdzenie}
\begin{proof}
W dalszym rozdziale, w Twierdzeniu~\ref{thm:compact-fourier22},
udowodnimy niezależnie od wzoru na splot, że
$\mathcal FU=T_a$ i że istnieje $m\in\N_0$ takie,
że dla każdego $\alpha\in\N_0^d$ zachodzi
$|\partial^\alpha a(\xi)|\le C_\alpha(1+|\xi|)^m$.
Stosujemy ten wynik tutaj, aby nie powtarzać dowodu
zamiany parowania z całką. Z Twierdzenia
\ref{thm:dt-schwartz-multiplier-characterization}
wynika, że $a$ jest mnożnikiem $\mathcal S$.

Dla ustalonego $y\in\R^d$ i $\varphi\in\mathcal S$
funkcja $\psi_y(\xi)=e^{-iy\cdot\xi}\varphi(\xi)$
należy do $\mathcal S$ na mocy
Twierdzenia~\ref{thm:tempered-operators22}
(jego dowodu dotyczącego mnożników).
Bezpośrednio z definicji całkowej mamy
$(\mathcal F\psi_y)(x)=(\mathcal F\varphi)(x+y)$.
Korzystając z $\mathcal FU=T_a$, dostajemy
\[
 \begin{aligned}
 A_U(\mathcal F\varphi)(y)
 &=\langle U_x,(\mathcal F\varphi)(x+y)\rangle\\
 &=\langle U,\mathcal F\psi_y\rangle
 =\langle\mathcal FU,\psi_y\rangle\\
 &=\int_{\R^d}a(\xi)e^{-iy\cdot\xi}\varphi(\xi)
 \,\mathrm d\lambda_d(\xi)
 =c_d^{-1}\mathcal F(a\varphi)(y),
 \end{aligned}
\]
gdzie $A_Ug(y)=\langle U_x,g(x+y)\rangle$.
Całka jest bezwzględnie zbieżna, ponieważ
$a\varphi\in\mathcal S\subseteq L^1$.
Operator $A_U$ zachowuje $\mathcal S$ i jest ciągły,
jak wykazano w poprzedniej uwadze
\ref{rem:fourier-conv25}. Stąd $U*T$ jest dystrybucją temperowaną.
Korzystamy z definicji splotu przez transpozycję $A_U$:
\[
\begin{aligned}
 \langle\mathcal F(U*T),\varphi\rangle
 &=\langle T,A_U(\mathcal F\varphi)\rangle\\
 &=c_d^{-1}\langle T,\mathcal F(a\varphi)\rangle\\
 &=c_d^{-1}\langle\mathcal F T,a\varphi\rangle .
\end{aligned}
\]
Jest to żądana równość, a zgodność splotów na
$\mathcal D$ wynika z poprzedniej uwagi.
\end{proof}

\begin{uwaga}[Zakres rozdziału]
\label{rem:scope22}
Transformacja Fouriera jest tu automorfizmem przestrzeni
dystrybucji temperowanych i podstawowym narzędziem teorii
operatorów. Jej zastosowania do rozwiązań fundamentalnych rozwiniemy
w dalszych rozdziałach tego skryptu; zastosowania do
konkretnych zagadnień RRC omawia także skrypt o równaniach
różniczkowych cząstkowych.
Dalsze rozdziały rozwiną teorię przestrzeni Montela,
przestrzeni nuklearnych i produktów tensorowych topologicznych.
\end{uwaga}

% Materiał z AF 25; fragment do rozdziału o dystrybucjach temperowanych.

\setcounter{section}{10}
% Kandydat restrukturyzacji; nie jest aktywnym rozdziałem Atlasu.
\section{Transformacja Fouriera w przestrzeni $L^2$ i twierdzenie Plancherela}
\providecommand{\atlasref}[2]{#2}

\begin{twierdzenie}[Twierdzenie Plancherela]
\label{thm:plancherel22}
Transformacja Fouriera z normalizacją $(2\pi)^{-d/2}$
przedłuża się jednoznacznie do operatora unitarnego
$\mathcal F:L^2(\R^d;\C)\to L^2(\R^d;\C)$,
czyli dla $f,g\in L^2$ zachodzi
$\int_{\R^d}\widehat f(\xi)\overline{\widehat g(\xi)}
\,\mathrm d\lambda_d(\xi)
=\int_{\R^d}f(x)\overline{g(x)}\,\mathrm d\lambda_d(x)$.
Dla $f\in L^2$ zachodzi ponadto
$\mathcal F T_f=T_{\widehat f}$, gdzie $\widehat f$
oznacza transformatę w $L^2$.
\end{twierdzenie}
\begin{proof}
Najpierw ustalmy $f,\varphi\in\mathcal S$.
Ponieważ $f,\varphi\in L^1$, moduł funkcji
podcałkowej w całce podwójnej jest całkowalny:
\[
 \int_{\R^d}\int_{\R^d}|f(x)\varphi(\xi)|
 \,\mathrm d\lambda_d(x)\,\mathrm d\lambda_d(\xi)
 =\|f\|_1\|\varphi\|_1.
\]
Twierdzenie Fubiniego i symetria jądra
$e^{-ix\cdot\xi}$ dają
\[
 \int_{\R^d}\widehat f(\xi)\varphi(\xi)
 \,\mathrm d\lambda_d(\xi)
 =\int_{\R^d}f(x)(\mathcal F\varphi)(x)
 \,\mathrm d\lambda_d(x).
\]
Jest to zgodność $\mathcal F T_f=T_{\widehat f}$ na
$\mathcal S$; parowanie jest biliniowe, a nie
półtoraliniowe.

Niech teraz $f,g\in\mathcal S$. Stosujemy powyższą
tożsamość do $\varphi=\overline{\widehat g}\in\mathcal S$.
Ze wzoru odwrotnego w Twierdzeniu~\ref{thm:fourier-schwartz17}
wynika
$\mathcal F(\overline{\widehat g})(x)
=\overline{(\mathcal F^{-1}\widehat g)(x)}
=\overline{g(x)}$.
Otrzymujemy
\[
 \int_{\R^d}\widehat f(\xi)\overline{\widehat g(\xi)}
 \,\mathrm d\lambda_d(\xi)
 =\int_{\R^d}f(x)\overline{g(x)}\,\mathrm d\lambda_d(x).
\]
W szczególności $\|\widehat f\|_2=\|f\|_2$.
Przestrzeń $\mathcal S$ jest gęsta w $L^2$, bo zawiera
gęstą podprzestrzeń $\mathcal D$. Dla $f\in L^2$
wybieramy $f_n\in\mathcal S$ z $f_n\to f$ w $L^2$.
Ciąg $(\widehat f_n)$ jest Cauchy'ego w $L^2$,
ponieważ $\|\widehat f_n-\widehat f_m\|_2
=\|f_n-f_m\|_2$. Z zupełności istnieje granica;
dla innego ciągu przybliżającego $(g_n)$ tę samą granicę
zapewnia $\|\widehat f_n-\widehat g_n\|_2
=\|f_n-g_n\|_2\to0$. Definiujemy $\widehat f$ jako
tak uzyskaną granicę. Liniowość i zachowanie normy
przechodzą na granicę; jednoznaczność przedłużenia
wynika z gęstości $\mathcal S$.
Obraz izometrii jest domknięty: jeśli ciąg obrazów
jest zbieżny, ciąg przeciwobrazów jest Cauchy'ego i
jego granica ma żądany obraz. Obraz zawiera
$\mathcal F(\mathcal S)=\mathcal S$, więc jest całym
$L^2$. Zachowanie iloczynu skalarnego wynika z ciągłości
iloczynu skalarnego na $L^2$ i tożsamości na $\mathcal S$.
Otrzymujemy zatem operator unitarny.

Na koniec, dla $f_n\to f$ w $L^2$ oraz ograniczonego
$B\subseteq\mathcal S$ mamy
\[
 P_B(T_{f_n}-T_f)
 \le C_d\|f_n-f\|_2
 \sup_{\varphi\in B}Q_{d+1,0}(\varphi)\longrightarrow0.
\]
Analogicznie $T_{\widehat f_n}\to T_{\widehat f}$ silnie
w $\mathcal S'$, bo $\widehat f_n\to\widehat f$ w $L^2$.
Z ciągłości $\mathcal F$ w topologii silnej i
wykazanej zgodności na $\mathcal S$ otrzymujemy
$\mathcal FT_f=\lim_n\mathcal FT_{f_n}
=\lim_nT_{\widehat f_n}=T_{\widehat f}$.
Jednoznaczność granicy wynika z warunku Hausdorffa
dla topologii silnej na $\mathcal S'$.
\end{proof}

\begin{stwierdzenie}[Zgodność transformat całkowej, dystrybucyjnej i unitarnej]
\label{prop:dt-fourier-l1-l2-compat}
Dla $f\in L^1$ transformacja dystrybucyjna spełnia
$\mathcal FT_f=T_{\widehat f}$, gdzie $\widehat f$
jest transformatą całkową. Gdy $f\in L^1\cap L^2$,
transformata całkowa jest reprezentantem transformaty
unitarnej z Twierdzenia~\ref{thm:plancherel22}.
\end{stwierdzenie}
\begin{proof}
Dla $\varphi\in\mathcal S$ całkowalność modułu
$|f(x)\varphi(\xi)|$ wynika z $f,\varphi\in L^1$.
Z twierdzenia Fubiniego otrzymujemy
\[
 \int_{\R^d}\widehat f(\xi)\varphi(\xi)\,\mathrm d\lambda_d(\xi)
 =\int_{\R^d}f(x)(\mathcal F\varphi)(x)\,\mathrm d\lambda_d(x).
\]
Jest to definicyjna tożsamość transformacji dystrybucyjnej.
Dla $f\in L^1\cap L^2$ twierdzenie Plancherela daje
tę samą dystrybucję regularną dla transformaty w $L^2$.
Oba reprezentanty są lokalnie całkowalne; ich równość
prawie wszędzie wynika ze
Stwierdzenia~\ref{prop:regular-injective19}.
\end{proof}

\setcounter{section}{11}
% Kandydat restrukturyzacji; nie jest aktywnym rozdziałem Atlasu.
\section{Splot funkcji i nierówność Younga}
\providecommand{\atlasref}[2]{#2}

\subsection{Nierówność Younga dla splotu}

\begin{twierdzenie}[Nierówność Younga]
\label{thm:dt-young-general}
Niech $1\le p,q,r\le\infty$ i
$1+1/r=1/p+1/q$, z konwencją $1/\infty=0$.
Dla $f\in L^p(\R^d)$ i $g\in L^q(\R^d)$ całka
\[
 (f*g)(x)=\int_{\R^d}f(y)g(x-y)\,\mathrm d\lambda_d(y)
\]
jest bezwzględnie zbieżna dla prawie każdego $x$,
określa element $L^r$ niezależny od reprezentantów
i spełnia $\|f*g\|_r\le\|f\|_p\|g\|_q$.
\end{twierdzenie}
\begin{proof}
Dla $r=\infty$ mamy $1/p+1/q=1$ i bezpośrednio
stosujemy nierówność Höldera do całki po $y$.
Niech zatem $r<\infty$. Wtedy $p,q<\infty$ oraz
$p,q\le r$. Rozkładamy moduł iloczynu na trzy czynniki:
\[
 |f(y)g(x-y)|
 =\bigl(|f(y)|^p|g(x-y)|^q\bigr)^{1/r}
 |f(y)|^{1-p/r}|g(x-y)|^{1-q/r}.
\]
Jeżeli któryś z dwóch ostatnich wykładników jest zerowy,
odpowiedni czynnik zastępujemy jednością.
Stosujemy nierówność Höldera dla trzech czynników,
z odwrotnościami wykładników
$1/r$, $1/p-1/r$ oraz $1/q-1/r$.
Ich suma jest równa jedności z założenia.
Otrzymujemy
\[
 \int_{\R^d}|f(y)g(x-y)|\,\mathrm d\lambda_d(y)
 \le A(x)^{1/r}\|f\|_p^{1-p/r}\|g\|_q^{1-q/r},
\]
gdzie
$A(x)=\int_{\R^d}|f(y)|^p|g(x-y)|^q\,\mathrm d\lambda_d(y)$.
Z twierdzenia Tonellego
$\int_{\R^d} A(x)\,\mathrm d\lambda_d(x)=\|f\|_p^p\|g\|_q^q$,
więc $A(x)<\infty$ prawie wszędzie.
Podnosimy oszacowanie do potęgi $r$, całkujemy po $x$
i otrzymujemy $\|f*g\|_r^r\le\|f\|_p^r\|g\|_q^r$.
Przypadek zerowej normy daje splot zerowy prawie wszędzie.
Mierzalność i niezależność od reprezentantów wynikają
z twierdzenia Tonellego zastosowanego do modułów i
do zbiorów, na których zmieniono reprezentanty.
\end{proof}

\begin{uwaga}[Dwie nierówności Younga i ostrość]
\label{rem:dt-young-scope}
Nie należy utożsamiać nierówności Younga dla splotu
z nierównością dla liczb dodatnich
$ab\le a^p/p+b^{p'}/p'$.
Stała jeden w Twierdzeniu~\ref{thm:dt-young-general}
wystarcza do konstrukcji ciągłych operatorów splotowych;
w ogólności nie jest ona optymalna. Przypadek $p=1$
obejmuje oszacowanie stosowane przy przybliżaniu jedynki.
\end{uwaga}

\setcounter{section}{12}
% Kandydat restrukturyzacji; nie jest aktywnym rozdziałem Atlasu.
\section{Interpolacja zespolona i nierówność Hausdorffa--Younga}
\providecommand{\atlasref}[2]{#2}

\subsection{Interpolacja zespolona i nierówność Hausdorffa--Younga}

Wynik dla $L^1$ daje kontrolę normy supremum transformaty,
a twierdzenie Plancherela zachowanie normy $L^2$.
Pomiędzy tymi przypadkami można uzyskać oszacowania w
$L^{p'}$, gdzie $p'$ jest wykładnikiem sprzężonym do $p$.
Potrzebujemy jednak interpolacji norm operatorów, nie tylko
interpolacji norm jednej funkcji.

\begin{lemat}[Lemat o trzech prostych]
\label{lem:dt-three-lines}
Niech $F$ będzie ograniczoną funkcją ciągłą na domkniętym
pasie $0\le\operatorname{Re}z\le1$, holomorficzną wewnątrz
pasa. Załóżmy, że dla każdego $t\in\R$ zachodzą
$|F(it)|\le M_0$ oraz $|F(1+it)|\le M_1$.
Dla $0<\theta<1$ mamy
\[
 |F(\theta)|\le M_0^{1-\theta}M_1^\theta.
\]
\end{lemat}
\begin{proof}
Najpierw zakładamy $M_0,M_1>0$. Funkcja
$G(z)=F(z)\exp(-(1-z)\log M_0-z\log M_1)$
jest ograniczona na pasie i ma moduł co najwyżej jeden
na jego pionowych brzegach. Dla $\varepsilon>0$ połóżmy
$H_\varepsilon(z)=G(z)\exp(\varepsilon(z^2-1))$.
Jeżeli $z=x+iy$, $0\le x\le1$, to moduł dodatkowego
czynnika wynosi $\exp(\varepsilon(x^2-y^2-1))$.
Na pionowych brzegach jest więc nie większy od jedności.
Na poziomych odcinkach $y=\pm R$ moduł $H_\varepsilon$
nie przekracza $\|G\|_\infty e^{-\varepsilon R^2}$,
co dla dużych $R$ jest mniejsze od jedności.
Z zasady maksimum modułu na prostokącie otrzymujemy
$|H_\varepsilon(\theta)|\le1$.
Przejście do granicy $\varepsilon\downarrow0$ daje
$|G(\theta)|\le1$. Dla $M_j=0$ stosujemy udowodniony
wynik ze stałymi $M_j+\delta$ i przechodzimy do
granicy $\delta\downarrow0$.
\end{proof}

\begin{twierdzenie}[Twierdzenie Riesza--Thorina]
\label{thm:dt-riesz-thorin}
Niech $(X,\mu)$ i $(Y,\nu)$ będą przestrzeniami
z miarami $\sigma$-skończonymi. Niech $T$ będzie
operatorem liniowym określonym na zespolonych funkcjach
prostych o nośniku skończonej miary w $X$, o wartościach
w klasach funkcji mierzalnych w $Y$.
Załóżmy, że dla $j=0,1$ oraz wszystkich takich $f$
\[
 \|Tf\|_{q_j}\le M_j\|f\|_{p_j},
 \qquad 1\le p_j<\infty,\quad 1\le q_j\le\infty.
\]
Dla $0<\theta<1$ definiujemy $p,q$ przez
\[
 \frac1p=\frac{1-\theta}{p_0}+\frac{\theta}{p_1},
 \qquad
 \frac1q=\frac{1-\theta}{q_0}+\frac{\theta}{q_1},
\]
przyjmując $1/\infty=0$. Wówczas $T$ przedłuża się
jednoznacznie do ograniczonego operatora
$L^p(X)\to L^q(Y)$, którego norma nie przekracza
$M_0^{1-\theta}M_1^\theta$.
\end{twierdzenie}
\begin{proof}
Przyjmijmy najpierw dodatnie $M_j$. Wystarczy rozważyć
niezerowe funkcje proste $f,g$ o nośnikach skończonej
miary, unormowane przez $\|f\|_p=\|g\|_{q'}=1$.
Na zbiorach, na których funkcje są niezerowe, piszemy
$f=\omega_f|f|$ oraz $g=\omega_g|g|$, gdzie
$|\omega_f|=|\omega_g|=1$. Dla $z$ w pasie definiujemy
\[
\begin{aligned}
 a(z)&=p\bigl((1-z)/p_0+z/p_1\bigr),\\
 f_z&=\omega_f\exp(a(z)\log|f|).
\end{aligned}
\]
Poza zbiorem $\{f\ne0\}$ przyjmujemy $f_z=0$.
Jeśli $q'<\infty$, definiujemy analogicznie
\[
 b(z)=q'\bigl((1-z)/q_0'+z/q_1'\bigr),
 \qquad
 g_z=\omega_g\exp(b(z)\log|g|),
\]
z wartością zero poza $\{g\ne0\}$.
Jeżeli $q'=\infty$, to $q=q_0=q_1=1$ i bierzemy
po prostu $g_z=g$.
Mamy $f_\theta=f$, $g_\theta=g$. Na brzegach pasa
$\|f_{j+it}\|_{p_j}=1$. Istotnie, podniesienie
modułu do potęgi $p_j$ daje $|f|^p$.
Podobnie $\|g_{j+it}\|_{q_j'}\le1$; jeśli
$q_j'=\infty$ w przypadku $q'<\infty$, moduł
na niezerowym nośniku jest równy jeden.

Rozważmy biliniowe parowanie
\[
 F(z)=\int_Y (Tf_z)(y)g_z(y)\,\mathrm d\nu(y).
\]
Jest ono skończone: rozpisując $f_z,g_z$ na skończone
sumy funkcji charakterystycznych, otrzymujemy całki
$T\chi_A$ po zbiorach skończonej miary, które istnieją
na mocy oszacowań brzegowych i nierówności Höldera.
Liniowość $T$ pokazuje, że $F$ jest skończoną sumą
funkcji wykładniczych o wykładnikach liniowych w $z$.
Ponieważ współczynniki przy $z$ w tych wykładnikach
są rzeczywiste, $F$ jest holomorficzna i ograniczona
na pasie. Na jego brzegach, z nierówności Höldera,
\[
 |F(j+it)|\le
 \|Tf_{j+it}\|_{q_j}\|g_{j+it}\|_{q_j'}\le M_j.
\]
Lemat~\ref{lem:dt-three-lines} daje
$|\int_Y (Tf)g\,\mathrm d\nu|\le M_0^{1-\theta}M_1^\theta$.
Po uwzględnieniu dowolnych norm $f,g$ otrzymujemy
\[
 \left|\int_Y(Tf)g\,\mathrm d\nu\right|
 \le M_0^{1-\theta}M_1^\theta\|f\|_p\|g\|_{q'}.
\]

Wyjaśnijmy, dlaczego wystarcza testowanie funkcjami
prostymi. Niech $u$ będzie funkcją mierzalną skończoną
prawie wszędzie, lokalnie całkowalną na zbiorach
skończonej miary, oraz niech
$|\int ug\,\mathrm d\nu|\le C\|g\|_{q'}$ dla takich
funkcji prostych $g$. Jeżeli $1<q<\infty$, wybieramy
rosnące wyczerpanie $Y$ zbiorami $B_n$ skończonej miary
i kładziemy $E_n=B_n\cap\{|u|\le n\}$. Testujemy
ograniczoną funkcją
$\chi_{E_n}\overline{\omega_u}|u|^{q-1}$, z fazą
$\omega_u=u/|u|$ na $\{u\ne0\}$ i zerem w pozostałych
punktach. Jest ona jednostajną granicą funkcji prostych
o nośnikach w $E_n$, więc nierówność przechodzi do
granicy zarówno w całce, jak i w normie $L^{q'}$.
Dostajemy
\[
 \int_{E_n}|u|^q\,\mathrm d\nu
 \le C\left(\int_{E_n}|u|^q\,\mathrm d\nu\right)^{1/q'}.
\]
Wynika stąd $\|\chi_{E_n}u\|_q\le C$, a twierdzenie
o zbieżności monotonicznej daje $\|u\|_q\le C$.
Dla $q=1$ stosujemy test $\chi_{E_n}\overline{\omega_u}$
o normie supremum nie większej od jedności.
Dla $q=\infty$ i $\eta>0$, gdyby zbiór
$\{|u|>C+\eta\}$ miał dodatnią miarę, miałby podzbiór
$B$ o dodatniej skończonej mierze. Test
$\chi_B\overline{\omega_u}$, znów przybliżony jednostajnie
funkcjami prostymi, prowadziłby do sprzeczności
$(C+\eta)\nu(B)\le C\nu(B)$.
Stosując ten argument do $u=Tf$, otrzymujemy żądane
oszacowanie normy. Przypadek zerowego $M_j$ wynika
z zastąpienia go przez $M_j+\delta$.
Ponieważ $p<\infty$, funkcje proste o nośniku skończonej
miary są gęste w $L^p$, a $L^q$ jest zupełna.
Przedłużenie przez granice jest zatem dobrze określone,
ograniczone i jednoznaczne.
\end{proof}

\begin{twierdzenie}[Nierówność Hausdorffa--Younga]
\label{thm:dt-hausdorff-young}
Niech $1\le p\le2$, a $p'$ będzie wykładnikiem sprzężonym.
Dla unitarnej normalizacji transformacji Fouriera
z Definicji~\ref{def:fourier-schwartz17} zachodzi
\[
 \|\widehat f\|_{p'}
 \le (2\pi)^{-d(1/p-1/2)}\|f\|_p,
 \qquad f\in L^p(\R^d).
\]
Transformata jest jednoznacznym ciągłym przedłużeniem
operatora całkowego z funkcji prostych o nośniku
skończonej miary, zgodnym z transformacją na $\mathcal S'$.
\end{twierdzenie}
\begin{proof}
Dla $p=1$ stosujemy Twierdzenie~\ref{thm:dt-fourier-l1-rl},
dla $p=2$ Twierdzenie~\ref{thm:plancherel22}.
Dla $1<p<2$ korzystamy z Twierdzenia~\ref{thm:dt-riesz-thorin}
z parametrami
\[
 (p_0,q_0,M_0)=(1,\infty,c_d),\qquad
 (p_1,q_1,M_1)=(2,2,1).
\]
Oba oszacowania dotyczą tego samego operatora dzięki
Stwierdzeniu~\ref{prop:dt-fourier-l1-l2-compat}.
Mamy $1/p=1-\theta/2$, $1/q=\theta/2$, zatem
$q=p'$ i $1-\theta=2/p-1$.
Norma interpolowana wynosi
$c_d^{\,2/p-1}=(2\pi)^{-d(1/p-1/2)}$.

Aby sprawdzić zgodność z dystrybucjami, wybieramy
funkcje proste $f_n$ o nośnikach skończonej miary
zbiegające do $f$ w $L^p$. Wówczas
$\widehat f_n\to\widehat f$ w $L^{p'}$.
Dla $\varphi\in\mathcal S$ nierówność Höldera daje
\[
 \left|\int_{\R^d}(f_n-f)\mathcal F\varphi\,
 \mathrm d\lambda_d\right|
 \le\|f_n-f\|_p\|\mathcal F\varphi\|_{p'}\to0,
\]
a analogicznie parowanie $\widehat f_n-\widehat f$
z $\varphi$ dąży do zera. Każde $f_n$ należy do $L^1$,
więc stosujemy
Stwierdzenie~\ref{prop:dt-fourier-l1-l2-compat}
i przechodzimy do granicy w tożsamości biliniowej.
Dostajemy $\mathcal FT_f=T_{\widehat f}$.
\end{proof}

\begin{uwaga}[Zakres i stała]
\label{rem:dt-hy-constant}
Powyższa stała pochodzi z interpolacji dwóch norm
brzegowych. Nie twierdzimy, że jest optymalna dla
$1<p<2$. Ostre oszacowanie Babenki--Becknera
wymaga odrębnego dowodu.
Nie można natomiast rozszerzyć samego oszacowania
$\|\widehat f\|_{p'}\le C\|f\|_p$ na $p>2$
przez formalną zamianę wykładników. Kontrprzykład
do takiego rozszerzenia zostanie omówiony poniżej.
\end{uwaga}

\begin{stwierdzenie}[Brak analogicznego oszacowania dla $p>2$]
\label{prop:dt-hy-fails}
Dla $2<p\le\infty$ nie istnieje stała $C$ taka, że
$\|\widehat f\|_{p'}\le C\|f\|_p$ dla wszystkich
$f\in\mathcal S(\R^d)$.
\end{stwierdzenie}
\begin{proof}
Niech $\psi$ będzie niezerową funkcją gładką
o nośniku w kuli o promieniu $1/4$ i niech
$\phi=\mathcal F^{-1}\psi\in\mathcal S$.
Dla $N\ge1$ i znaków $\varepsilon_j\in\{-1,1\}$
kładziemy
\[
 f_N(x)=\phi(x)\sum_{j=1}^N\varepsilon_j e^{ijx_1}.
\]
Transformaty składników mają parami rozłączne nośniki,
więc dla $q=p'<2$ (także $q=1$)
\[
 \|\widehat f_N\|_q=N^{1/q}\|\psi\|_q.
\]
Dla skończonego $p>2$ wybierzmy znaki niezależnie,
z równymi prawdopodobieństwami obu wartości.
Dla $m$ całkowitego z $2m\ge p$ rozwijamy
$\left|\sum_j\varepsilon_j e^{ijx_1}\right|^{2m}$.
Wartość oczekiwana składnika znika, jeśli którykolwiek
indeks występuje nieparzystą liczbę razy. Pozostałe
składniki mają moduł jeden i najwyżej $m$ różnych
indeksów. Ich liczba nie przekracza
$(2m)^{2m}N^m$ (wybór wzorca indeksów i samych indeksów).
Nierówność Höldera na przestrzeni probabilistycznej daje
zatem
\[
 \mathbb E\left|\sum_j\varepsilon_j e^{ijx_1}\right|^p
 \le (2m)^p N^{p/2}.
\]
Całkując i stosując Tonellego, znajdujemy wybór znaków
taki, że $\|f_N\|_p\le 2m\sqrt N\|\phi\|_p$.
Gdyby rozważane oszacowanie zachodziło, to
$N^{1/q-1/2}$ byłoby ograniczone niezależnie od $N$,
co jest niemożliwe, gdy $q<2$.

Dla $p=\infty$ korzystamy z jawnej konstrukcji
wielomianów Rudina--Shapira:
$P_0=Q_0=1$ oraz
\[
 P_{n+1}(z)=P_n(z)+z^{2^n}Q_n(z),\qquad
 Q_{n+1}(z)=P_n(z)-z^{2^n}Q_n(z).
\]
Indukcyjnie oba mają po $2^n$ współczynników
równych $1$ lub $-1$, ponieważ nośniki potęg
w obu składnikach są rozłączne. Dla $|z|=1$ mamy
$|P_{n+1}|^2+|Q_{n+1}|^2=2(|P_n|^2+|Q_n|^2)$,
więc $|P_n(z)|\le\sqrt{2^{n+1}}$.
Biorąc $N=2^n$ i $f_N(x)=\phi(x)e^{ix_1}P_n(e^{ix_1})$,
otrzymujemy $\|f_N\|_\infty\le\sqrt{2N}\|\phi\|_\infty$,
podczas gdy $\|\widehat f_N\|_1=N\|\psi\|_1$.
To również wyklucza istnienie stałej $C$.
\end{proof}

\setcounter{section}{13}
% Kandydat restrukturyzacji; nie jest aktywnym rozdziałem Atlasu.
\section{Transformacja Fouriera miar}
\providecommand{\atlasref}[2]{#2}

\subsection{Transformacja Fouriera miar}

\begin{definicja}[Transformata Fouriera skończonej miary]
\label{def:dt-fourier-measure}
Dla skończonej zespolonej miary Radona $\mu$ na $\R^d$
definiujemy
\[
 \widehat\mu(\xi)=c_d\int_{\R^d}e^{-ix\cdot\xi}
 \,\mathrm d\mu(x),\qquad c_d=(2\pi)^{-d/2}.
\]
Normą miary jest jej całkowita wariacja
$\|\mu\|=|\mu|(\R^d)$.
Por. \atlasref{ME-16::thm:complex-variation-16}{wariacja miary zespolonej}.
\end{definicja}

\begin{twierdzenie}[Ciągłość i zgodność z dystrybucjami]
\label{thm:dt-fourier-measure-cont}
Transformata $\widehat\mu$ jest ograniczona i jednostajnie
ciągła, $\|\widehat\mu\|_\infty\le c_d\|\mu\|$.
Dystrybucja temperowana
$T_\mu(\varphi)=\int\varphi\,\mathrm d\mu$
spełnia $\mathcal FT_\mu=T_{\widehat\mu}$.
Jeśli $\mu=f\lambda_d$, $f\in L^1$, to $\widehat\mu=\widehat f$.
\end{twierdzenie}
\begin{proof}
Oszacowanie wynika z nierówności
$|\int u\,\mathrm d\mu|\le\int|u|\,\mathrm d|\mu|$.
Dla dowolnego $h\in\R^d$ zachodzi
\[
 \sup_{\xi\in\R^d}|\widehat\mu(\xi+h)-\widehat\mu(\xi)|
 \le c_d\int_{\R^d}|e^{-ix\cdot h}-1|\,\mathrm d|\mu|(x).
\]
Prawa strona dąży do zera przez zbieżność zmajoryzowaną
z majorantą $2$, ponieważ $|\mu|$ jest skończona.
Ponieważ $|T_\mu(\varphi)|\le\|\mu\|\|\varphi\|_\infty$,
funkcjonał jest ciągły na $\mathcal S$.
Dla $\varphi\in\mathcal S$ moduł jądra podwójnego
parowania jest całkowalny względem
$|\mu|\otimes\lambda_d$, z całką ograniczoną przez
$c_d\|\mu\|\|\varphi\|_1$. Twierdzenie Fubiniego daje
\[
 \int_{\R^d}\widehat\mu(\xi)\varphi(\xi)\,\mathrm d\lambda_d(\xi)
 =\int_{\R^d}(\mathcal F\varphi)(x)\,\mathrm d\mu(x).
\]
To jest definicja transformaty dystrybucyjnej.
Ostatnia równość wynika z podstawienia gęstości $f$.
\end{proof}

\begin{twierdzenie}[Jednoznaczność transformaty miary]
\label{thm:dt-fourier-measure-unique}
Jeśli $\mu,\nu$ są skończonymi zespolonymi miarami Radona
i $\widehat\mu=\widehat\nu$, to $\mu=\nu$.
\end{twierdzenie}
\begin{proof}
Z Twierdzenia~\ref{thm:dt-fourier-measure-cont} mamy
$\mathcal FT_{\mu-\nu}=0$. Automorfizm $\mathcal F$ na
$\mathcal S'$ z Twierdzenia~\ref{thm:fourier-auto-tempered22}
daje $T_{\mu-\nu}=0$.
Dla $u\in C_c(\R^d)$ wybieramy nieujemną funkcję testową
$\rho$ o całce jeden i kładziemy
$\rho_\varepsilon(x)=\varepsilon^{-d}\rho(x/\varepsilon)$.
Funkcja $u*\rho_\varepsilon$ jest gładka o zwartym nośniku;
jej zbieżność jednostajna do $u$ wynika z jednostajnej
ciągłości $u$ i zwężającego się nośnika jądra:
\[
 \|u*\rho_\varepsilon-u\|_\infty
 \le\sup_{y\in\operatorname{supp}\rho_\varepsilon}
        \|\tau_yu-u\|_\infty\to0.
\]
Z ograniczoności miary otrzymujemy
$\int u\,\mathrm d(\mu-\nu)=0$ dla każdego $u\in C_c$.
Jednoznaczność w twierdzeniu Riesza--Markowa--Kakutaniego
(por. \atlasref{ME-21::cor:measures-separate-21}{rozdzielanie miar przez funkcje ciągłe})
daje $\mu=\nu$.
\end{proof}

\begin{stwierdzenie}[Splot miar i splot miary z funkcją]
\label{prop:dt-fourier-measure-conv}
Dla skończonych zespolonych miar Radona definiujemy
$\mu*\nu$ jako obraz miary produktowej $\mu\otimes\nu$
przez odwzorowanie $(x,y)\mapsto x+y$. Wówczas
$\|\mu*\nu\|\le\|\mu\|\|\nu\|$ oraz
\[
 \widehat{\mu*\nu}=c_d^{-1}\widehat\mu\,\widehat\nu.
\]
Dla $f\in L^p(\R^d)$, $1\le p\le\infty$, splot
$\mu*f(x)=\int f(x-y)\,\mathrm d\mu(y)$ jest określony
prawie wszędzie i $\|\mu*f\|_p\le\|\mu\|\|f\|_p$.
Jeżeli $f\in L^1$, to
$\widehat{\mu*f}=c_d^{-1}\widehat\mu\,\widehat f$.
\end{stwierdzenie}
\begin{proof}
Wariacja miary produktowej jest ograniczona przez
$|\mu|\otimes|\nu|$ (można to uzyskać, przedstawiając
obie miary przez ich gęstości względem wariacji).
Wariacja obrazu nie przekracza obrazu wariacji, co
wynika z definicji przez rozłączne podziały.
Stąd oszacowanie normy splotu.
Rozkład $e^{-i(x+y)\cdot\xi}=e^{-ix\cdot\xi}e^{-iy\cdot\xi}$
i twierdzenie Fubiniego dają wzór na transformatę.

Dla $\mu\ne0$ i $1\le p<\infty$ stosujemy nierówność
Jensena względem miary probabilistycznej $|\mu|/\|\mu\|$:
\[
 \left(\int|f(x-y)|\,\mathrm d|\mu|(y)\right)^p
 \le \|\mu\|^{p-1}\int|f(x-y)|^p\,\mathrm d|\mu|(y).
\]
Po całkowaniu po $x$ i zastosowaniu Tonellego prawa strona
ma całkę $\|\mu\|^p\|f\|_p^p$. Dowodzi to również istnienia
całki prawie wszędzie. Dla $p=\infty$ wybieramy
reprezentanta ograniczonego przez $\|f\|_\infty$;
zmiana reprezentanta nie zmienia splotu prawie wszędzie
na mocy Tonellego. Dla $\mu=0$ wszystko jest oczywiste.
Jeśli $f\in L^1$, twierdzenie Fubiniego i podstawienie
$z=x-y$ dają ostatni wzór; moduł jądra podwójnej całki
ma całkę $c_d\|\mu\|\|f\|_1$.
\end{proof}

\begin{przyklad}[Miara Diraca a lemat Riemanna--Lebesgue'a]
\label{ex:dt-fourier-measure-atoms}
Mamy $\widehat{\delta_a}(\xi)=c_de^{-ia\cdot\xi}$,
więc transformata skończonej miary nie musi zanikać
w nieskończoności. Dla $\sum_{j=1}^N b_j\delta_{a_j}$
transformata jest sumą $c_d\sum_j b_je^{-ia_j\cdot\xi}$.
W przypadku gęstości całkowalnej zanikanie zapewnia
Twierdzenie~\ref{thm:dt-fourier-l1-rl}. Zatem hipoteza
o bezwzględnej ciągłości miary z całkowalną gęstością
ma tu istotne znaczenie.
\end{przyklad}

\setcounter{section}{14}
% Kandydat restrukturyzacji; nie jest aktywnym rozdziałem Atlasu.
\section{Transformacja Fouriera funkcji radialnych}
\providecommand{\atlasref}[2]{#2}

\subsection{Transformacja Fouriera funkcji radialnych}

Niech $f(x)=F(|x|)$ będzie funkcją radialną w $L^1(\R^d)$.
Niezmienniczość względem ortogonalnych zmian współrzędnych
w Stwierdzeniu~\ref{prop:dt-fourier-l1-rules} daje
$\widehat f(O\xi)=\widehat f(\xi)$ dla $O\in O(d)$.
Można zatem zastąpić całkę wielowymiarową całką po
promieniu. W wymiarze jeden otrzymujemy bezpośrednio
\[
 \widehat f(\xi)=\sqrt{2/\pi}\int_0^\infty
 F(r)\cos(r\xi)\,\mathrm d\lambda_1(r).
\]

\begin{lemat}[Całka po sferze]
\label{lem:dt-fourier-sphere-kernel}
Niech $d\ge2$ i $\sigma_{d-1}$ będzie miarą powierzchniową
na sferze jednostkowej $S^{d-1}$, czyli miarą Hausdorffa
o wymiarze $d-1$. Dla $t\in\R$ mamy
\[
 \int_{S^{d-1}}e^{-it\omega_1}\,\mathrm d\sigma_{d-1}(\omega)
 =|S^{d-2}|\int_{-1}^1 e^{-itu}(1-u^2)^{(d-3)/2}
 \,\mathrm d\lambda_1(u).
\]
Dla $d=2$ przyjmujemy $|S^0|=2$.
\end{lemat}
\begin{proof}
Poza dwoma biegunami stosujemy parametryzację
$\omega=(u,\sqrt{1-u^2}\,\eta)$, $\eta\in S^{d-2}$,
$-1<u<1$. Wektor pochodnej względem $u$ ma normę
$(1-u^2)^{-1/2}$ i jest prostopadły do kierunków
stycznych w $\eta$, ponieważ te ostatnie są prostopadłe
do $\eta$. Długości w $d-2$ kierunkach stycznych
skaluje czynnik $\sqrt{1-u^2}$. Z wzoru na pole
powierzchni iloczyn czynników wynosi
$(1-u^2)^{(d-3)/2}$. Bieguny mają miarę zero.
Dla $d=2$ parametryzacja składa się z dwóch łuków
odpowiadających dwóm punktom $S^0$, a ten sam
rachunek daje element długości $(1-u^2)^{-1/2}$.
Całkowana funkcja zależy wyłącznie od $u$, co daje tezę.
\end{proof}

\begin{definicja}[Funkcja Bessela pierwszego rodzaju]
\label{def:dt-bessel-j}
Dla $\alpha\ge0$ i $t>0$ definiujemy
\[
 J_\alpha(t)=\sum_{m=0}^\infty
 \frac{(-1)^m(t/2)^{2m+\alpha}}
      {m!\,\Gamma(m+\alpha+1)}.
\]
Tutaj $\Gamma(a)=\int_0^\infty
s^{a-1}e^{-s}\,\mathrm d\lambda_1(s)$ dla $a>0$.
Stosunek modułów kolejnych składników szeregu wynosi
$t^2/(4(m+1)(m+\alpha+1))$, ponieważ całkowanie
przez części daje $\Gamma(a+1)=a\Gamma(a)$.
Szereg zatem zbiega bezwzględnie i jednostajnie na
zwartych przedziałach dodatnich $t$.
\end{definicja}

\begin{twierdzenie}[Wzór radialny z funkcją Bessela]
\label{thm:dt-fourier-radial}
Dla $d\ge2$, $\alpha=d/2-1$ oraz $f(x)=F(|x|)\in L^1$,
przy $\rho=|\xi|>0$ zachodzi
\[
 \widehat f(\xi)=\rho^{-\alpha}\int_0^\infty
 F(r)J_\alpha(r\rho)r^{\alpha+1}\,\mathrm d\lambda_1(r).
\]
Przy $\rho=0$ wartość wynosi
$c_d|S^{d-1}|\int_0^\infty F(r)r^{d-1}\,\mathrm d\lambda_1(r)$.
\end{twierdzenie}
\begin{proof}
Najpierw obliczymy jądro sferyczne. Dla $a,b>0$
podstawienie $s=vw$, $t=v(1-w)$ w iloczynie całek
definiujących $\Gamma(a)\Gamma(b)$ ma jakobian $v$.
Twierdzenie Tonellego daje wzór beta:
\[
 \int_0^1 w^{a-1}(1-w)^{b-1}\,\mathrm d\lambda_1(w)
 =\frac{\Gamma(a)\Gamma(b)}{\Gamma(a+b)}.
\]
Ponadto całka Gaussa daje $\Gamma(1/2)=\sqrt{\pi}$,
a całkowanie $e^{-|x|^2}$ we współrzędnych biegunowych
w $\R^n$ daje
$\pi^{n/2}=|S^{n-1}|\Gamma(n/2)/2$.
W szczególności $|S^{d-2}|=2\pi^{(d-1)/2}/\Gamma((d-1)/2)$.
Rozwijamy wykładnik w Lemacie~\ref{lem:dt-fourier-sphere-kernel}
w szereg, który jest jednostajnie bezwzględnie zbieżny
dla $-1\le u\le1$ przy ustalonym $t$.
Składniki nieparzyste całkują się do zera, a podstawienie
$w=u^2$ i wzór beta dają
\[
 \int_{-1}^1 u^{2m}(1-u^2)^{(d-3)/2}\,\mathrm d\lambda_1(u)
 =\frac{\Gamma(m+1/2)\Gamma((d-1)/2)}{\Gamma(m+d/2)}.
\]
Z rekurencji funkcji gamma
$\Gamma(m+1/2)=\sqrt{\pi}(2m)!/(4^m m!)$.
Podstawiając do szeregu i mnożąc przez $|S^{d-2}|$,
otrzymujemy dla $t>0$
\[
 \int_{S^{d-1}}e^{-it\omega_1}\,\mathrm d\sigma_{d-1}
 =(2\pi)^{d/2}t^{-\alpha}J_\alpha(t).
\]
Istotnie, po podstawieniu definicji $J_\alpha$ po obu
stronach współczynnik przy $(-1)^mt^{2m}$ jest równy
$2\pi^{d/2}/(4^m m!\Gamma(m+d/2))$.

Wracamy do transformacji $f$. Obracamy $\xi$ na
$\rho e_1$ i stosujemy współrzędne biegunowe:
\[
 \widehat f(\xi)=c_d\int_0^\infty F(r)r^{d-1}
 \left(\int_{S^{d-1}}e^{-ir\rho\omega_1}
 \,\mathrm d\sigma_{d-1}(\omega)\right)\mathrm d\lambda_1(r).
\]
Całkowalność wynika z $f\in L^1$ i modułu jądra równego
jedności. Po wstawieniu obliczonej całki sferycznej
czynniki $(2\pi)^{d/2}$ i $c_d$ skracają się, a
$d-1-\alpha=\alpha+1$. Daje to wzór.
Wartość w zerze wynika wprost z całki definiującej
transformatę; nie trzeba dzielić przez $\rho=0$.
\end{proof}

\begin{przyklad}[Transformata kuli]
\label{ex:dt-fourier-ball}
Dla $R>0$ i $d\ge2$ niech $f=\chi_{\{|x|\le R\}}$.
Z Twierdzenia~\ref{thm:dt-fourier-radial} wynika dla
$\rho=|\xi|>0$
\[
 \widehat f(\xi)=R^{d/2}\rho^{-d/2}J_{d/2}(R\rho).
\]
Aby obliczyć całkę radialną, różniczkujemy szereg
definiujący $t^{\alpha+1}J_{\alpha+1}(t)$ na zwartych
przedziałach dodatnich $t$. Zbieżność jednostajna
szeregu pochodnych wynika z kryterium ilorazowego.
Korzystając z $\Gamma(m+\alpha+2)=(m+\alpha+1)\Gamma(m+\alpha+1)$,
dostajemy
$(t^{\alpha+1}J_{\alpha+1}(t))'=t^{\alpha+1}J_\alpha(t)$.
Podstawienie $t=r\rho$ i granica w zerze (równa zero
z szeregu) dają wskazany wzór. W szczególności dla $d=2$
jest to $R J_1(R\rho)/\rho$. W zerze transformata jest
równa $c_d\lambda_d(B(0,R))$, zgodnie z definicją.
\end{przyklad}

\setcounter{section}{15}
% Kandydat restrukturyzacji; nie jest aktywnym rozdziałem Atlasu.
\section{Transformata dystrybucji o zwartym nośniku}
\providecommand{\atlasref}[2]{#2}

\begin{twierdzenie}[Transformata dystrybucji o zwartym nośniku]
\label{thm:compact-fourier22}
Jeżeli $T\in\mathcal E'(\R^d)$ ma zwarty nośnik $K$,
to jego transformata Fouriera jest dystrybucją regularną
reprezentowaną przez gładką funkcję
\[
F_T(\xi)=(2\pi)^{-d/2}\langle T_x,e^{-ix\cdot\xi}\rangle.
\]
Dla pewnego $m\in\N_0$ i każdego wielowskaźnika $\beta$
istnieje $C_\beta>0$ takie, że
$|\partial_\xi^\beta F_T(\xi)|
\le C_\beta(1+|\xi|)^m$
dla wszystkich $\xi\in\R^d$.
W działaniu na funkcji wykładniczej używamy
przedłużenia z Lematu~\ref{lem:compact-action21}.
\end{twierdzenie}
\begin{proof}
Wybierzmy $\chi\in\mathcal D$ równą $1$ na
otoczeniu $K$. Dla $\beta\in\N_0^d$ funkcje
$x\mapsto\chi(x)(-ix)^\beta e^{-ix\cdot\xi}$
mają nośniki w ustalonym zwartym zbiorze $\supp\chi$.
Na zwartym otoczeniu ustalonego parametru $\xi$
wszystkie pochodne względem $x$ i parametru są ciągłe
i jednostajnie ciągłe na iloczynie tych dwóch zbiorów.
Stosując twierdzenie Newtona--Leibniza do ilorazów
różnicowych względem parametru w seminormach
$\mathcal D_{\supp\chi}$, dostajemy
\[
 \partial_\xi^\beta F_T(\xi)
 =(2\pi)^{-d/2}
 \langle T_x,(-ix)^\beta e^{-ix\cdot\xi}\rangle.
\]
Ciągłość tych pochodnych wynika z tego samego argumentu.
Z Lematu~\ref{lem:compact-action21} wybieramy $C,m,L$
niezależne od $\beta,\xi$, gdzie $L$ jest zwarty i
$|T(f)|\le C\max_{|\alpha|\le m}\sup_{x\in L}
|\partial^\alpha f(x)|$.
Dla $|\alpha|\le m$ wzór Leibniza daje
\[
 \partial_x^\alpha\bigl((-ix)^\beta e^{-ix\cdot\xi}\bigr)
 =\sum_{\gamma\le\alpha}\binom\alpha\gamma
 \partial_x^\gamma((-ix)^\beta)
 (-i\xi)^{\alpha-\gamma}e^{-ix\cdot\xi}.
\]
Ponieważ $x\in L$, pochodne wielomianu są ograniczone,
a moduł czynnika wykładniczego dla rzeczywistych $\xi$
wynosi $1$. Ponadto $|\xi^{\alpha-\gamma}|
\le(1+|\xi|)^m$. Po zsumowaniu skończonej liczby
składników otrzymujemy $C_\beta(1+|\xi|)^m$ jako
oszacowanie modułu $\partial_\xi^\beta F_T(\xi)$.
W szczególności $T_{F_T}$ jest temperowana.

Aby sprawdzić równość dystrybucji, ustalmy
$\varphi\in\mathcal S$ i $c_d=(2\pi)^{-d/2}$.
Dla $R\ge1$ połóżmy
\[
 J_R(x)=c_d\chi(x)\int_{[-R,R]^d}
 e^{-ix\cdot\xi}\varphi(\xi)\,\mathrm d\lambda_d(\xi).
\]
Sumy Riemanna zbiegają do $J_R$ w
$\mathcal D_{\supp\chi}$: dla każdego rzędu pochodnych
względem $x$ wynika to z jednostajnej ciągłości funkcji
podcałkowej i jej pochodnych na iloczynie zwartych
zbiorów $\supp\chi$ i $[-R,R]^d$.
Stąd
\[
 \langle T,J_R\rangle=
 \int_{[-R,R]^d}F_T(\xi)\varphi(\xi)
 \,\mathrm d\lambda_d(\xi).
\]
Dla każdego $q\in\N_0$ wzór Leibniza daje
\[
 p_q(J_R-\chi\mathcal F\varphi)
 \le C_{\chi,q,d}\int_{\R^d\setminus[-R,R]^d}
 (1+|\xi|)^q|\varphi(\xi)|\,\mathrm d\lambda_d(\xi)
 \longrightarrow0.
\]
Pochodne pod całką istnieją na mocy całkowalności
każdej wagi wielomianowej razy $|\varphi|$.
Wszystkie funkcje po lewej mają nośniki w $\supp\chi$,
więc zbieżność zachodzi w jednym etapie $\mathcal D$.
Z ciągłości $T$ i bezwzględnej całkowalności
$F_T\varphi$ dostajemy
\[
 \langle\mathcal FT,\varphi\rangle
 =\langle T,\chi\mathcal F\varphi\rangle
 =\int_{\R^d}F_T(\xi)\varphi(\xi)
 \,\mathrm d\lambda_d(\xi).
\]
Pierwsza równość wykorzystuje jednoznaczne przedłużenie
$T$ na $\mathcal S$ i regułę odcięcia.
Zatem $\mathcal FT=T_{F_T}$.
\end{proof}

\setcounter{part}{2}
\part{Struktura topologiczna i dalsza teoria dystrybucji}
\setcounter{section}{16}
% Rozdział aktywny po restrukturyzacji DT.
% Wersja źródłowa: audyt II, jednostka 23

\section{Przestrzenie Montela i nuklearne}
\label{sec:montel-nuklearne23}

\providecommand{\atlasref}[2]{#2}


\noindent\textbf{Wymagane wyniki.} Korzystamy z topologii
przestrzeni funkcji testowych i przestrzeni Schwartza,
twierdzeń Banacha--Steinhausa i Arzelà--Ascoliego,
teorii przestrzeni Hilberta oraz analizy Fouriera.
\medskip

Niech $d\ge1$ będzie liczbą całkowitą, a $\Omega\subseteq\R^d$
niepustym zbiorem otwartym. W poprzednich rozdziałach
zbudowaliśmy przestrzenie
$\mathcal E(\Omega)=C^\infty(\Omega)$,
$\mathcal D(\Omega)$ i $\mathcal S(\R^d)$ wraz z ich
przestrzeniami sprzężonymi. Obecnie wyjaśnimy, dlaczego w tych
przestrzeniach ograniczoność prowadzi do zwartości oraz w jakim
sensie operator liniowy może być reprezentowany przez jedno jądro
dystrybucyjne. Wszystkie przestrzenie liniowo-topologiczne w tym
rozdziale są nad $\K\in\{\R,\C\}$ i są Hausdorffa. W argumentach
wykorzystujących rozwinięcia Fouriera pracujemy nad $\C$;
wersje rzeczywiste wynikają przez kompleksyfikację.

\subsection{Własność Montela i jej znaczenie}

\noindent\emph{Powiązania tematyczne w Bibliotece matematycznej:} \atlasref{GT-07::thm:topog-zwartosc-sieci}{GT-07: zwartość i podsieci zbieżne}.


\begin{definicja}[Przestrzeń Montela]
\label{def:montel23}
Przestrzenią Montela nazywamy beczkową przestrzeń lokalnie wypukłą
$E$, w której każdy zbiór ograniczony ma zwarte domknięcie.
Jeżeli dodatkowo $E$ jest przestrzenią Frécheta, mówimy o
\emph{przestrzeni Frécheta--Montela}. W tej klasie beczkowość
wynika już z twierdzenia Baire'a, więc wystarczy wymagać
względnej zwartości zbiorów ograniczonych.
\end{definicja}

\begin{twierdzenie}[Klasyczne przestrzenie Montela]
\label{thm:montel-test23}
Dla każdej liczby całkowitej $d\ge1$ i każdego
niepustego zbioru otwartego $\Omega\subseteq\R^d$
przestrzenie $\mathcal E(\Omega)$,
$\mathcal S(\R^d)$ oraz $\mathcal D(\Omega)$, wyposażone
w swoje standardowe topologie lokalnie wypukłe, są przestrzeniami
Montela.
\end{twierdzenie}
\begin{proof}
Dla $\mathcal S$ względną zwartość zbiorów ograniczonych
wykazano w Lemacie~\ref{lem:montel-schwartz22} (dowód
jest taki sam nad $\R$ i nad $\C$). Z Twierdzenia~
\ref{thm:schwartz-frechet17} przestrzeń $\mathcal S$
jest przestrzenią Frécheta, a więc jest beczkowa.

Niech $(K_j)_{j\ge1}$ będzie wyczerpaniem $\Omega$
z Lematu~\ref{lem:exhaustion17}.
Dla ograniczonego $B\subseteq\mathcal E(\Omega)$,
$j\ge1$ i $m\in\N_0$ mamy
$M_{j,m}=\sup_{f\in B}p_{K_j,m}(f)<\infty$,
gdzie $p_{K_j,m}$ jest supremum modułów pochodnych
rzędu co najwyżej $m$ na $K_j$. Możemy założyć,
że $B$ i rozpatrywane $K_j$ są niepuste.
Ponieważ $K_j\subseteq\operatorname{int}K_{j+1}$,
istnieje $\delta_j>0$ takie, że domknięte
$\delta_j$-otoczenie $K_j$ jest zawarte w $K_{j+1}$.
Jeżeli $x,y\in K_j$ i $|x-y|<\delta_j$, odcinek
od $x$ do $y$ leży w $K_{j+1}$. Dla $|\alpha|\le m$
i $f\in B$ z twierdzenia Newtona--Leibniza wynika
\[
 |\partial^\alpha f(x)-\partial^\alpha f(y)|
 \le M_{j+1,m+1}\sum_{i=1}^d|x_i-y_i|.
\]
Rodzina pochodnych jest więc jednostajnie równociągła
i jednostajnie ograniczona na $K_j$, także gdy
$\Omega$ lub $K_j$ nie są wypukłe.
Dla dowolnego ciągu $(f_n)\subseteq B$ twierdzenie
Arzelà--Ascoliego i wybór podciągu przekątniowego
dla przeliczalnej rodziny par $(j,\alpha)$ dają
podciąg $(f_{n_k})$, którego każda pochodna zbiega
jednostajnie na każdym $K_j$. Zatem $(f_{n_k})$
jest ciągiem Cauchy'ego dla wszystkich $p_{K_j,m}$.
Zupełność i metryzowalność z Twierdzenia~
\ref{thm:frechet-cinf17} dają zbieżność w $\mathcal E$.
Dla ciągu $(g_n)$ w domknięciu $B$ wybieramy
$f_n\in B$ w odległości mniejszej niż $1/n$ od $g_n$
w metryce wyznaczającej topologię $\mathcal E$.
Zbieżny podciąg $(f_{n_k})$ daje zbieżny podciąg
$(g_{n_k})$ o tej samej granicy w domknięciu $B$.
Domknięcie jest więc zwarte z kryterium ciągowego
zwartości w przestrzeniach metrycznych.
Przestrzeń $\mathcal E$ jest beczkowa jako przestrzeń Frécheta.

Dla $\mathcal D(\Omega)$ względną zwartość zbiorów
ograniczonych wykazano już w Twierdzeniu~
\ref{thm:D-Montel18}. Pozostaje beczkowość.
Niech $U\subseteq\mathcal D(\Omega)$ będzie beczką,
czyli zbiorem domkniętym, absolutnie wypukłym i pochłaniającym.
Dla każdego etapu $\mathcal D_{K_j}$ przeciwobraz $U$
przez kanoniczne włożenie jest beczką w przestrzeni
Frécheta, a zatem otoczeniem zera w tym etapie.
Funkcjonał Minkowskiego $p_U$ jest seminormą, której
ograniczenia do wszystkich etapów są ciągłe.
Na mocy konstrukcji topologii granicy induktywnej
(Twierdzenie~\ref{thm:LF-universal18}) seminorma
$p_U$ jest ciągła na $\mathcal D(\Omega)$.
Ponieważ $\{f:p_U(f)<1\}\subseteq U$, zbiór $U$
jest otoczeniem zera. Dowodzi to beczkowości,
a tym samym własności Montela.
\end{proof}

\begin{wniosek}[Zbieżność na ograniczonych zbiorach testowych]
\label{cor:montel-dual-seq23}
Niech $E$ będzie przestrzenią Montela, a $(u_n)$ ciągiem w
$E'$ zbieżnym punktowo na $E$ do funkcjonału $u\in E'$.
Wtedy $u_n\to u$ w topologii silnej $\beta(E',E)$.
W szczególności dotyczy to każdej z przestrzeni
z Twierdzenia~\ref{thm:montel-test23}.
\end{wniosek}
\begin{proof}
Ponieważ $E$ jest beczkowa, punktowa ograniczoność
rodziny $\{u_n-u:n\ge1\}$, wynikająca ze zbieżności
punktowej, daje jej równociągłość na mocy twierdzenia
Banacha--Steinhausa. Niech $B\subseteq E$ będzie
niepustym zbiorem ograniczonym i $\varepsilon>0$.
Z równociągłości wybieramy otwarte otoczenie zera $V$
takie, że $|(u_n-u)(v)|<\varepsilon/3$ dla wszystkich
$v\in V$ i $n\ge1$. Zwartość $\overline B$ daje
skończone pokrycie zbiorami $x_i+V$,
$i=1,\ldots,r$, gdzie $x_i\in\overline B$.
Dla dostatecznie dużego $n$ wszystkie
$|(u_n-u)(x_i)|<\varepsilon/3$. Dla $x\in B$
wybieramy $i$ z $x-x_i\in V$; wówczas
$|(u_n-u)(x)|<2\varepsilon/3$.
Stąd $\sup_{x\in B}|(u_n-u)(x)|\le2\varepsilon/3
<\varepsilon$. Dla pustego $B$ warunek jest trywialny.
Jest to zbieżność silna. Odwrotna implikacja wynika
z ograniczoności singletonów.
\end{proof}

\begin{przyklad}[Dlaczego nie wystarczy przestrzeń Banacha]
\label{ex:banach-not-montel23}
Jeżeli $X$ jest nieskończenie wymiarową przestrzenią Banacha,
to jej domknięta kula jednostkowa nie jest zwarta w topologii
normowej, na mocy lematu Riesza. Dlatego taka przestrzeń
nie jest przestrzenią Montela. Własność Montela
jest tu skutkiem szczególnie mocnej topologii przestrzeni
funkcji gładkich, a nie samej zupełności.
\end{przyklad}

\subsection{Operatory nuklearne w przestrzeniach Banacha}

\begin{definicja}[Operator nuklearny]
\label{def:nuclear-banach23}
Niech $X,Y$ będą przestrzeniami Banacha. Operator liniowy
$T:X\to Y$ nazywamy \emph{nuklearnym}, jeżeli istnieją
$\ell_j\in X'$ i $y_j\in Y$, $j\ge1$, takie, że
\[
 Tx=\sum_{j=1}^\infty \ell_j(x)y_j\quad(x\in X),
 \qquad \sum_{j=1}^\infty\|\ell_j\|\,\|y_j\|<\infty.
\]
Normą nuklearną $\|T\|_{\mathcal N}$ nazywamy infimum
ostatniej sumy po wszystkich takich przedstawieniach.
Przedstawienie dopuszcza także skończenie wiele składników.
Dla ustalonego $x\in X$ szereg jest bezwzględnie
zbieżny w normie $Y$, ponieważ
$\sum_j\|\ell_j(x)y_j\|\le\|x\|\sum_j\|\ell_j\|\|y_j\|$,
a $Y$ jest zupełna.
\end{definicja}

\begin{twierdzenie}[Nuklearność implikuje zwartość]
\label{thm:nuclear-compact23}
Jeśli $X,Y$ są przestrzeniami Banacha i $T:X\to Y$
jest operatorem nuklearnym, to $T$ jest ograniczony i zwarty.
\end{twierdzenie}
\begin{proof}
Dla przedstawienia z Definicji~\ref{def:nuclear-banach23}
niech $T_Nx=\sum_{j=1}^N\ell_j(x)y_j$.
Dla każdego $x\in X$ mamy
$\|Tx\|\le\|x\|\sum_j\|\ell_j\|\|y_j\|$,
więc $T$ jest ograniczony. Każdy $T_N$ ma skończony
rząd i jest ograniczony. Ponadto
$\|T-T_N\|\le\sum_{j>N}\|\ell_j\|\|y_j\|\to0$.
Przeto $T$ jest ograniczony i jest granicą operatorów
skończonego rzędu w normie operatorowej. Granica
operatorów zwartych w normie operatorowej jest zwarta:
obraz kuli jednostkowej można dla dowolnego $\varepsilon>0$
pokryć skończoną liczbą kul o promieniu $\varepsilon$
za pomocą dostatecznie dobrego przybliżenia $T_N$.
Dokładniej, wybieramy $N$ z $\|T-T_N\|<\varepsilon/2$.
Obraz kuli jednostkowej przez $T_N$ jest ograniczony
w przestrzeni skończenie wymiarowej, więc ma skończone
pokrycie kulami o promieniu $\varepsilon/2$.
Te same środki z promieniem $\varepsilon$ pokrywają
obraz przez $T$. Zupełność $Y$ i całkowita ograniczoność
obrazu zapewniają zwartość jego domknięcia.
\end{proof}

\begin{stwierdzenie}[Norma nuklearna i składanie operatorów]
\label{prop:nuclear-ideal23}
Niech $X,Y$ będą przestrzeniami Banacha nad $\K$.
Operatory nuklearne $X\to Y$ tworzą przestrzeń liniową,
a $\|\cdot\|_{\mathcal N}$ jest na niej normą,
spełniającą $\|T\|\le\|T\|_{\mathcal N}$.
Jeżeli dodatkowo $X_0,Y_0$ są przestrzeniami Banacha
i $A:X_0\to X$, $B:Y\to Y_0$ są ograniczonymi
operatorami liniowymi, to $BTA$ jest nuklearny oraz
$\|BTA\|_{\mathcal N}\le\|B\|\|T\|_{\mathcal N}\|A\|$.
\end{stwierdzenie}
\begin{proof}
Przedstawienia nuklearne dwóch operatorów można połączyć
w jeden szereg dla ich sumy; suma norm składników jest
sumą odpowiednich oszacowań. Skalowanie jednego z czynników
każdego składnika przez skalar daje przedstawienie
dla wielokrotności operatora. Biorąc infima, otrzymujemy
nierówność trójkąta i jednorodność (dla niezerowego
skalaru stosujemy też skalowanie przez jego odwrotność).
Oszacowanie $\|T\|\le\sum_j\|\ell_j\|\|y_j\|$
z poprzedniego dowodu daje $\|T\|\le\|T\|_{\mathcal N}$,
więc norma nuklearna oddziela punkty.
Dla złożenia mamy
$BTAx=\sum_j(\ell_j\circ A)(x)By_j$ oraz
\[
 \sum_j\|\ell_j\circ A\|\|By_j\|
 \le\|B\|\|A\|\sum_j\|\ell_j\|\|y_j\|.
\]
Ciągłość $B$ uzasadnia zastosowanie go do sumy szeregu,
a wzięcie infimum daje żądane oszacowanie.
\end{proof}

\begin{twierdzenie}[Zupełność normy nuklearnej]
\label{thm:dt-nuclear-norm-complete}
Dla przestrzeni Banacha $X,Y$ przestrzeń
$\mathcal N(X,Y)$ operatorów nuklearnych, wyposażona
w normę z Definicji~\ref{def:nuclear-banach23},
jest przestrzenią Banacha.
\end{twierdzenie}
\begin{proof}
Na mocy Stwierdzenia~\ref{prop:nuclear-ideal23}
jest to przestrzeń unormowana. Najpierw niech
$(A_k)_{k\ge1}$ będzie ciągiem operatorów nuklearnych
takim, że $\sum_{k\ge1}\|A_k\|_{\mathcal N}<\infty$.
Dla każdego $k$ wybieramy przedstawienie
\[
 A_kx=\sum_{j\ge1}\ell_{k,j}(x)y_{k,j},\qquad
 \sum_{j\ge1}\|\ell_{k,j}\|\|y_{k,j}\|
 \le\|A_k\|_{\mathcal N}+2^{-k}.
\]
Wobec tego
$\sum_{k,j\ge1}\|\ell_{k,j}\|\|y_{k,j}\|<\infty$.
Dla każdego $x\in X$ szereg podwójny
$\sum_{k,j\ge1}\ell_{k,j}(x)y_{k,j}$ jest bezwzględnie
zbieżny w przestrzeni Banacha $Y$. Zatem dowolna
enumeracja par $(k,j)$ daje tę samą sumę, która
definiuje operator nuklearny $A$. Sumowanie najpierw
po $j$ daje $Ax=\sum_{k\ge1}A_kx$.
Dla ogona szeregu otrzymujemy przedstawienie nuklearne
przez ograniczenie do $k>m$, skąd
\[
 \left\|A-\sum_{k=1}^m A_k\right\|_{\mathcal N}
 \le\sum_{k>m}\bigl(\|A_k\|_{\mathcal N}+2^{-k}\bigr)
 \longrightarrow0.
\]
Każdy szereg bezwzględnie zbieżny w normie nuklearnej
jest więc zbieżny w tej normie.

Niech teraz $(T_n)$ będzie ciągiem Cauchy'ego
w normie nuklearnej. Wybieramy rosnący ciąg indeksów
$(n_k)$ taki, że
$\|T_{n_{k+1}}-T_{n_k}\|_{\mathcal N}\le2^{-k}$.
Jest to możliwe z warunku Cauchy'ego: wybieramy $n_k$
tak, aby dla wszystkich indeksów nie mniejszych
niż $n_k$ odległość była co najwyżej $2^{-k}$.
Stosując pierwszą część dowodu do różnic kolejnych
wyrazów podciągu, otrzymujemy zbieżność szeregu
teleskopowego w normie nuklearnej, a więc zbieżność
$(T_{n_k})$ do pewnego $T\in\mathcal N(X,Y)$.
Warunek Cauchy'ego i nierówność trójkąta dają
$T_n\to T$ w tej samej normie.
\end{proof}

\begin{przyklad}[Operator zwarty, który nie jest nuklearny]
\label{ex:compact-nonnuclear23}
Na przestrzeni Hilberta $\ell^2(\N;\K)$ z kanoniczną
bazą ortonormalną $(e_n)_{n\ge1}$ rozważmy operator
$D(x_n)_{n\ge1}=(x_n/n)_{n\ge1}$. Jest liniowy i
ograniczony, gdyż $\|Dx\|_2\le\|x\|_2$.
Iloczyn skalarny przyjmujemy liniowy w pierwszym argumencie.
Operator $D$ jest zwarty, gdyż operatory skończonego rzędu
$D_N=\sum_{n\le N}n^{-1}\langle\cdot,e_n\rangle e_n$
spełniają $\|D-D_N\|=1/(N+1)$:
oszacowanie z góry wynika z sumy kwadratów ogona,
a równość uzyskujemy dla $x=e_{N+1}$. Nie jest jednak
nuklearny. Istotnie, dla każdego przedstawienia nuklearnego
$D=\sum_j\ell_j(\cdot)y_j$ i każdego $N$ mamy
\[
\sum_{n=1}^N|\langle De_n,e_n\rangle|
\le\sum_j
\Big(\sum_{n=1}^N|\ell_j(e_n)|^2\Big)^{1/2}
\Big(\sum_{n=1}^N|\langle y_j,e_n\rangle|^2\Big)^{1/2}
\le\sum_j\|\ell_j\|\|y_j\|.
\]
Istotnie, w pierwszej nierówności stosujemy nierówność
trójkąta do szeregu dla każdego $n$, zamieniamy sumę
skończoną z szeregiem nieujemnym i stosujemy nierówność
Cauchy'ego--Schwarza do sumy po $n$. W drugiej nierówności
korzystamy z twierdzenia Riesza o reprezentacji funkcjonałów
na przestrzeni Hilberta oraz z nierówności Bessela:
$\sum_{n=1}^N|\ell_j(e_n)|^2\le\|\ell_j\|^2$ i
$\sum_{n=1}^N|\langle y_j,e_n\rangle|^2\le\|y_j\|^2$.
Lewa strona jest równa $\sum_{n=1}^N1/n$ i rośnie
nieograniczenie. Zatem zwartość jest własnością ściśle
słabszą niż nuklearność.
\end{przyklad}

\subsection{Nuklearność przestrzeni lokalnie wypukłych}

\begin{definicja}[Przestrzeń nuklearna]
\label{def:nuclear-lcs23}
Niech $E$ będzie przestrzenią lokalnie wypukłą Hausdorffa,
$\mathcal P$ rodziną wszystkich ciągłych seminorm na $E$.
Dla $p\in\mathcal P$ definiujemy przestrzeń Banacha $E_p$
jako uzupełnienie przestrzeni unormowanej $E/\ker p$
w normie $\|x+\ker p\|=p(x)$.
Przez $[x]_p$ oznaczamy obraz $x\in E$ w $E_p$.
Jeżeli $p\le q$, to $\ker q\subseteq\ker p$ i
odwzorowanie $[x]_q\mapsto[x]_p$ ma normę co najwyżej $1$,
więc jednoznacznie przedłuża się do ograniczonego operatora
$j_{q,p}:E_q\to E_p$ o gęstym obrazie.
Dla $p\le q\le r$ zachodzi $j_{r,p}=j_{q,p}j_{r,q}$:
obie strony zgadzają się na gęstym obrazie $E$ w $E_r$.
Przestrzeń $E$ nazywamy \emph{nuklearną}, gdy dla każdego
$p\in\mathcal P$ istnieje $q\in\mathcal P$, $q\ge p$,
takie, że $j_{q,p}$ jest operatorem nuklearnym.
\end{definicja}

\begin{lemat}[Kryterium dla rodziny seminorm wyznaczającej topologię]
\label{lem:nuclear-cofinal23}
Niech $E$ będzie lokalnie wypukłą przestrzenią Hausdorffa,
a $\mathcal P_0$ rodziną ciągłych seminorm wyznaczającą
jej topologię i skierowaną ku górze względem nierówności.
Jeżeli dla każdego $p_0\in\mathcal P_0$ istnieje
$q_0\in\mathcal P_0$ z $q_0\ge p_0$ i nuklearnym
$j_{q_0,p_0}$, to $E$ jest nuklearna.
\end{lemat}
\begin{proof}
Niech $p$ będzie dowolną ciągłą seminormą. Z ciągłości
i skierowania rodziny wynika $p\le C p_0$ dla pewnych
$C\ge1$ i $p_0\in\mathcal P_0$: odpowiednie skończone
przecięcie kul seminormowych zastępujemy jedną kulą
i stosujemy skalowanie. Wybierzmy $q_0$ jak w założeniu
i połóżmy $q=Cq_0$. Wtedy $q\ge p$.
Odwzorowanie $E_q\to E_{q_0}$ dane przez identyczność
na klasach jest ograniczone, podobnie jak
$E_{p_0}\to E_p$, którego norma nie przekracza $C$.
Ich złożenie z $j_{q_0,p_0}$ jest nuklearne na mocy
Stwierdzenia~\ref{prop:nuclear-ideal23} i zgadza się
z $j_{q,p}$ na gęstym obrazie $E$ w $E_q$.
\end{proof}

\begin{uwaga}[Ważne rozróżnienie]
\label{rem:nuclear-different23}
Nuklearność przestrzeni lokalnie wypukłej jest własnością
\emph{całej rodziny seminorm}, a nie twierdzeniem,
że identyczność na tej przestrzeni ma jedno nuklearne
przedstawienie. Jeśli przestrzeń Banacha $X$ jest nuklearna
w sensie Definicji~\ref{def:nuclear-lcs23}, to $X$
ma skończony wymiar: dla $p(x)=\|x\|$ nuklearność
daje ciągłą seminormę $q\ge p$. Z ciągłości $q$
w topologii normowej wynika $q\le Cp$, więc $q$
jest normą równoważną $p$ i oba uzupełnienia
identyfikujemy z $X$. Nuklearne $j_{q,p}$ po
złożeniu z ograniczoną identycznością $X_p\to X_q$
daje nuklearną identyczność na $X$, a ta jest zwarta
na mocy Twierdzenia~\ref{thm:nuclear-compact23}.
Lemat Riesza wyklucza nieskończony wymiar.
Odwrotnie, w skończenie wymiarowej przestrzeni Banacha
każde $j_{q,p}$ ma skończony rząd i jest nuklearne
(przedstawienie w bazie obrazu z ciągłymi współrzędnymi);
można wziąć $q=p$.
\end{uwaga}

\begin{stwierdzenie}[Nuklearność a całkowita ograniczoność]
\label{prop:dt-nuclear-precompact}
W nuklearnej przestrzeni lokalnie wypukłej Hausdorffa
$E$ każdy zbiór ograniczony $B$ jest całkowicie
ograniczony: dla każdego otoczenia zera $U$ istnieje
skończony zbiór $F\subseteq E$ taki, że $B\subseteq F+U$.
Jeżeli dodatkowo $E$ jest zupełna i beczkowa, to jest
przestrzenią Montela.
\end{stwierdzenie}
\begin{proof}
Dla niepustego $B$ i otoczenia zera $U$ wybieramy
ciągłą seminormę $p$ taką, że $\{x:p(x)<1\}\subseteq U$.
Jest to możliwe, ponieważ skończone przecięcie kul
seminormowych można zapisać jako kulę maksimum
odpowiednio przeskalowanych seminorm.
Z nuklearności wybieramy $q\ge p$ z nuklearnym
przejściem $j_{q,p}:E_q\to E_p$. Obraz $B$ w $E_q$
jest ograniczony, gdyż $\sup_{x\in B}q(x)<\infty$.
Na mocy Twierdzenia~\ref{thm:nuclear-compact23}
obraz $B$ w $E_p$ ma zwarte domknięcie.
Pokrywamy go skończenie wieloma kulami o promieniu
$1/3$ w $E_p$. Z każdej kuli przecinającej ten obraz
wybieramy jeden punkt postaci $[x_i]_p$, $x_i\in B$.
Dla każdego $x\in B$ istnieje więc $i$ taki, że
$p(x-x_i)<2/3<1$, czyli $x\in x_i+U$.
Dla pustego $B$ teza jest oczywista.

Domknięcie zbioru całkowicie ograniczonego pozostaje
całkowicie ograniczone: dla $U$ wybieramy otwarte
symetryczne otoczenie zera $V$ z $V+V\subseteq U$,
wtedy $\overline V\subseteq V+V$, a skończone
pokrycie $B$ przesunięciami $V$ daje pokrycie
$\overline B$ przesunięciami $U$.
Jeżeli $E$ jest zupełna, to $\overline B$ jest zupełne
jako zbiór domknięty w przestrzeni zupełnej.
Z kryterium zwartości w przestrzeniach jednostajnych
(zupełność i całkowita ograniczoność) wynika jego
zwartość. Wraz z beczkowością daje to własność Montela
z Definicji~\ref{def:montel23}. Sama nuklearność
nie stanowi tu założenia zupełności ani beczkowości.
\end{proof}

\begin{twierdzenie}[Nuklearność przestrzeni szybko malejących ciągów]
\label{thm:nuclear-s23}
Niech $d\ge1$ będzie liczbą całkowitą, $I=\N_0^d$
oraz $a_\nu\in\K$. Definiujemy
\[
 s(I)=\bigl\{a=(a_\nu)_{\nu\in I}:
 p_k(a)=\sup_{\nu\in I}(1+|\nu|)^k|a_\nu|<\infty
 \text{ dla każdego }k\in\N_0\bigr\},
 \qquad |\nu|=\nu_1+\cdots+\nu_d.
\]
Przestrzeń $s(I)$ z topologią seminorm $(p_k)$ jest
nuklearną przestrzenią Frécheta--Montela.
\end{twierdzenie}
\begin{proof}
Normy $p_k$ są rosnące i oddzielają punkty, więc
wyznaczają metryzowalną topologię lokalnie wypukłą.
Niech $(a^{(j)})$ będzie ciągiem Cauchy'ego w każdej
normie $p_k$. Granice współrzędnych
$a_\nu=\lim_j a_\nu^{(j)}$ istnieją dzięki zupełności $\K$.
Dla ustalonego $k$ i $\varepsilon>0$ istnieje $J$
takie, że dla $i,j\ge J$ i każdego $\nu$ zachodzi
$(1+|\nu|)^k|a_\nu^{(i)}-a_\nu^{(j)}|<\varepsilon$.
Przechodząc z $j$ do nieskończoności, otrzymujemy
$p_k(a^{(i)}-a)\le\varepsilon$ dla $i\ge J$.
W szczególności $p_k(a)<\infty$, więc $a\in s(I)$
i $a^{(i)}\to a$ w jego topologii. Jest to przestrzeń
Frécheta.

Dla $m\in\N_0$ oznaczmy $E_m=s(I)_{p_m}$.
Jest to przestrzeń
\[
 c_{0,m}(I)=\{a:(1+|\nu|)^m a_\nu\to0
 \text{ poza skończonymi podzbiorami }I\}
\]
z normą $p_m$. Mnożenie współrzędnych przez
$(1+|\nu|)^m$ jest izometrią tej przestrzeni na
zupełną przestrzeń $c_0(I;\K)$. Ponieważ
$(1+|\nu|)^m|a_\nu|\le p_{m+1}(a)/(1+|\nu|)$,
mamy $s(I)\subseteq c_{0,m}(I)$.
Ciągi o skończonym nośniku należą do $s(I)$
i są gęste w $c_{0,m}(I)$ przez obcinanie ogonów.
To uzasadnia identyfikację uzupełnienia.

Ustalmy $k$ i wybierzmy całkowite $m>k+d$.
Niech $e_\nu$ będzie ciągiem równym $1$ na indeksie $\nu$
i $0$ na pozostałych. Funkcjonał
$\mathrm{ev}_\nu(a)=a_\nu$ na $E_m$ ma normę
$(1+|\nu|)^{-m}$ (równość osiąga odpowiednio
przeskalowany $e_\nu$), a $p_k(e_\nu)=(1+|\nu|)^k$.
Zatem
\[
 \sum_{\nu\in I}\|\mathrm{ev}_\nu\|_{E_m'}
 \|e_\nu\|_{E_k}
 =\sum_{\nu\in I}(1+|\nu|)^{k-m}<\infty.
\]
Istotnie, liczba rozwiązań równania
$\nu_1+\cdots+\nu_d=r$ w $\N_0^d$ wynosi
$\binom{r+d-1}{d-1}$ (rozmieszczamy $d-1$
przegród między $r$ jednostkami).
Jest ograniczona przez $C_d(1+r)^{d-1}$.
Otrzymujemy szereg o wyrazach ograniczonych przez
$C_d(1+r)^{k-m+d-1}$, zbieżny, gdyż
$k-m+d-1<-1$. Szereg
$j_{p_m,p_k}a=\sum_{\nu\in I}a_\nu e_\nu$
jest zbieżny w normie $E_k$ i ma podane współrzędne.
Jest to przedstawienie nuklearne przejścia kanonicznego.
Lemat~\ref{lem:nuclear-cofinal23} daje nuklearność $s(I)$.

Niech $B\subseteq s(I)$ będzie niepustym zbiorem
ograniczonym i $M_k=\sup_{a\in B}p_k(a)$.
Dla każdego $R\ge0$ mamy
\[
 \sup_{a\in B,\,|\nu|>R}(1+|\nu|)^k|a_\nu|
 \le \frac{M_{k+1}}{1+R}.
\]
Dla dowolnego ciągu w $B$ wybieramy podciąg
przekątniowy zbieżny na każdej współrzędnej
(korzystając z ograniczoności w $p_0$).
Na skończonym zbiorze $|\nu|\le R$ jest on
ciągiem Cauchy'ego w każdej normie ważonej.
Na pozostałych indeksach różnica dwóch wyrazów
ma normę ważoną ograniczoną przez $2M_{k+1}/(1+R)$.
Wybierając najpierw $R$, a potem dostatecznie późne
wyrazy, otrzymujemy ciąg Cauchy'ego w każdej normie.
Z zupełności wynika jego zbieżność w $s(I)$.
Jak w dowodzie Twierdzenia~\ref{thm:montel-test23},
metryzowalność pozwala przejść do zwartości domknięcia
$B$. Beczkowość wynika z własności Frécheta.
\end{proof}

\begin{twierdzenie}[Nuklearność $C^\infty(\Omega)$]
\label{thm:nuclear-cinfty23}
Dla każdego niepustego zbioru otwartego $\Omega\subseteq\R^d$
przestrzeń $\mathcal E(\Omega)=C^\infty(\Omega)$
ze standardową topologią Frécheta jest nuklearna.
\end{twierdzenie}
\begin{proof}
Rozpatrzmy najpierw przypadek zespolony. Ustalmy
niepusty zwarty $K\Subset\Omega$ i $m\in\N_0$.
Niech $p_{K,m}(f)=\max_{|\alpha|\le m}\sup_{x\in K}
|\partial^\alpha f(x)|$. Z Lematu~\ref{lem:cutoffs17}
wybieramy $\chi\in\mathcal D(\Omega)$ równą $1$
w otoczeniu $K$, o nośniku $L=\supp\chi\Subset\Omega$.
W szczególności $K\subseteq L$. Funkcję $\chi f$
przedłużamy zerem na $\R^d$, co daje funkcję gładką,
ponieważ $L\Subset\Omega$.
Wybieramy $R>0$ z $L\subset(-R,R)^d$ i przedłużamy
tę funkcję okresowo z sześcianu $Q=[-R,R]^d$.
Przedłużenie jest gładkie, ponieważ funkcja znika
w otoczeniu brzegu sześcianu.
Przy znormalizowanej mierze $(2R)^{-d}\lambda_d$
na $Q$ funkcje $e_n(x)=e^{i\pi n\cdot x/R}$,
$n\in\mathbb Z^d$, tworzą bazę ortonormalną w $L^2$
torusa. Ortogonalność wynika z całkowania iloczynów
wykładników po kolejnych współrzędnych. Do zupełności
korzystamy z jednowymiarowego twierdzenia o gęstości
wielomianów trygonometrycznych, udowodnionego w dalszej
części skryptu (Twierdzenie~\ref{thm:dt-series-fejer}),
po liniowym przeskalowaniu przedziału.
Aby przejść do wymiaru $d$, stosujemy indukcję:
jeśli $g\in L^2(Q)$ jest prostopadła do wszystkich
$e_n$, to dla każdego $n'\in\mathbb Z^{d-1}$
funkcja ostatniej współrzędnej
\[
 g_{n'}(t)=(2R)^{-(d-1)}
 \int_{[-R,R]^{d-1}}g(x',t)e^{-i\pi n'\cdot x'/R}
 \,\mathrm d\lambda_{d-1}(x')
\]
należy do $L^2([-R,R])$ z nierówności
Cauchy'ego--Schwarza i twierdzenia Fubiniego.
Wszystkie jej współczynniki Fouriera są zerowe
z założenia i twierdzenia Fubiniego, zatem $g_{n'}=0$
prawie wszędzie. Ponieważ zbiór indeksów $n'$ jest
przeliczalny, dla prawie każdego $t$ jednocześnie
wszystkie współczynniki funkcji $g(\cdot,t)$ znikają.
Ta funkcja należy do $L^2([-R,R]^{d-1})$ dla prawie
każdego $t$, więc z założenia indukcyjnego jest zerowa
prawie wszędzie. Fubini daje $g=0$ w $L^2(Q)$,
co dowodzi zupełności układu ortonormalnego.
Niech
\[
 c_n(f)=(2R)^{-d}\int_Q(\chi f)(x)
 e^{-i\pi n\cdot x/R}\,\mathrm d\lambda_d(x).
\]
Ponieważ $(1-\Delta)e_n=(1+\pi^2|n|^2/R^2)e_n$,
dla każdego $r\in\N_0$ całkowanie przez części
bez wyrazów brzegowych daje
\[
 (1+\pi^2|n|^2/R^2)^r c_n(f)
 =(2R)^{-d}\int_Q (1-\Delta)^r(\chi f)(x)
 \overline{e_n(x)}\,\mathrm d\lambda_d(x).
\]
Z reguły Leibniza i zwartości $L$ otrzymujemy
\[
 |c_n(f)|\le C_r(1+|n|)^{-2r}p_{L,2r}(f).
\]
Stała zależy od $r,\chi,R,d$, lecz nie od $f,n$.
W szczególności $c_n$ znika na $\ker p_{L,2r}$
i przedłuża się do funkcjonału na uzupełnieniu
$E_{p_{L,2r}}$ z normą ograniczoną przez
$C_r(1+|n|)^{-2r}$.

Niech $v_n=[\chi e_n]_{p_{K,m}}$.
Ponieważ $\chi=1$ w otoczeniu $K$, klasa ta ma
na $K$ wszystkie pochodne do rzędu $m$ takie jak $e_n$;
jej norma nie przekracza $C_{R,m}(1+|n|)^m$.
Wybierzmy $r$ tak, by $2r>m+d$.
Szereg $\sum_n (1+|n|)^{m-2r}$ jest zbieżny:
liczba punktów $\mathbb Z^d$ z normą maksimum równą
$s\ge1$ wynosi $(2s+1)^d-(2s-1)^d
\le C_d(1+s)^{d-1}$, a normy euklidesowa i maksimum
są równoważne.

Rozwinięcie $\sum_n c_n(f)e_n$ jest jednostajnie
zbieżne na $Q$ wraz ze wszystkimi pochodnymi do rzędu
$m$, na mocy powyższych oszacowań.
Sumę szeregu bez pochodnych identyfikujemy z okresowym
przedłużeniem $\chi f$: sumy częściowe zbiegają
w $L^2(Q)$ do $\chi f$ z zupełności układu
trygonometrycznego, a jednostajna zbieżność daje
ten sam limit w $L^2(Q)$. Obie funkcje są ciągłe,
więc równość prawie wszędzie daje równość wszędzie.
Jednostajna zbieżność pochodnych pozwala różniczkować
szereg do rzędu $m$ (stosujemy twierdzenie
Newtona--Leibniza kolejno do każdej współrzędnej).
Ponieważ $\chi f=f$ w otoczeniu $K$ i $\chi e_n=e_n$
w tym samym otoczeniu, otrzymujemy w uzupełnieniu
\[
 [f]_{p_{K,m}}=\sum_{n\in\mathbb Z^d}c_n(f)v_n.
\]
Nie chodzi jedynie o wartości funkcji na $K$,
lecz także o pochodne do rzędu $m$.
Ponadto
\[
 \sum_n\|c_n\|_{(E_{p_{L,2r}})'}\|v_n\|_{E_{p_{K,m}}}
 \le C\sum_n(1+|n|)^{m-2r}<\infty.
\]
Ponieważ $p_{K,m}\le p_{L,2r}$, szereg definiuje
operator nuklearny $E_{p_{L,2r}}\to E_{p_{K,m}}$.
Zgadza się on z przejściem kanonicznym na gęstym
obrazie $E$, a zatem na całym uzupełnieniu.

W przypadku rzeczywistym stosujemy tę samą konstrukcję
do rzeczywistego $f$, ale bierzemy część rzeczywistą
szeregu. Składnik zastępujemy sumą
$\operatorname{Re}c_n(f)\operatorname{Re}v_n-
\operatorname{Im}c_n(f)\operatorname{Im}v_n$,
gdzie części rzeczywiste i urojone klas reprezentowane
są przez odpowiednie części funkcji $\chi e_n$
(możemy wybrać $\chi$ rzeczywistą).
Normy obu rzeczywistych funkcjonałów i obu wektorów
nie przekraczają odpowiadających oszacowań zespolonych,
więc szereg nuklearny ma sumę norm ograniczoną przez
dwukrotność poprzedniego oszacowania.

Standardowe seminormy $p_{K,m}$ tworzą rodzinę
skierowaną ku górze: dla dwóch zbiorów bierzemy
ich sumę, a dla rzędów maksimum.
Wyznaczają topologię $\mathcal E(\Omega)$.
Lemat~\ref{lem:nuclear-cofinal23} kończy dowód
w obu przypadkach.
\end{proof}

\begin{uwaga}[Hermite i przestrzeń Schwartza]
\label{rem:hermite-nuclear23}
Przestrzeń Schwartza $\mathcal S(\R^d)$ jest nuklearna:
współczynniki w bazie funkcji Hermite'a wyznaczają
izomorfizm liniowo-topologiczny
$\mathcal S(\R^d)\simeq s(\N_0^d)$.
Dokładne oszacowania i dowód tego izomorfizmu podano
w podsekcji~\ref{sec:hermite-supp23}.
Nuklearność $\mathcal D_K(\Omega)$ oraz
$\mathcal D(\Omega)$ zostanie poniżej wyprowadzona
z udowodnionych twierdzeń trwałości; nie jest to
formalna konsekwencja samej nuklearności $\mathcal S$.
\end{uwaga}

\subsection{Kryterium hilbertowskie i trwałość nuklearności}
\label{subsec:nuclear-permanence23}

W zastosowaniach do funkcji testowych potrzebujemy dwóch operacji:
przejścia do podprzestrzeni i do granicy induktywnej. Dowiedziemy
odpowiednich twierdzeń, wychodząc wyłącznie z
Definicji~\ref{def:nuclear-lcs23}. Ciągłą seminormę $h$ na przestrzeni lokalnie wypukłej $E$
nazywamy \emph{hilbertowską}, jeśli istnieje dodatnia
półokreślona hermitowska forma półtoraliniowa $b$ (nad $\R$:
symetryczna forma dwuliniowa), taka że $h(x)^2=b(x,x)$
dla każdego $x\in E$. Przyjmujemy liniowość w pierwszym
argumencie. Nierówność Cauchy'ego--Schwarza dla $b$
zapewnia $b(x,y)=0$ dla $x\in\ker h$ i dowolnego $y$,
więc $b$ wyznacza iloczyn skalarny na $E/\ker h$.
Przez $H_h=E_h$ oznaczamy jego uzupełnienie hilbertowskie.
Dla $h\le r$ operator $j_{r,h}:H_r\to H_h$ jest
przejściem kanonicznym z Definicji~\ref{def:nuclear-lcs23}.
Nie wymagamy, by $E$ była metryzowalna.

Przypomnijmy, że ograniczony operator liniowy $J:H\to G$
między przestrzeniami Hilberta jest operatorem
Hilberta--Schmidta, gdy $\sum_{i\in I}\|Ju_i\|^2<\infty$
dla bazy ortonormalnej $(u_i)_{i\in I}$ przestrzeni $H$.
Dla nieujemnych składników sumę po dowolnym zbiorze
indeksów definiujemy jako supremum sum skończonych.
Jeśli $(v_j)_{j\in J_0}$ jest bazą ortonormalną $G$,
to z tożsamości Parsevala i zamiany sum nieujemnych
otrzymujemy
\[
 \sum_{i\in I}\|Ju_i\|^2
 =\sum_{i\in I}\sum_{j\in J_0}|\langle Ju_i,v_j\rangle|^2
 =\sum_{j\in J_0}\|J^*v_j\|^2.
\]
Dowodzi to niezależności od wyboru obu baz oraz
$\|J^*\|_{\mathrm{HS}}=\|J\|_{\mathrm{HS}}$, gdzie
$\|J\|_{\mathrm{HS}}^2=\sum_i\|Ju_i\|^2$.

\begin{lemat}[Złożenie operatorów Hilberta--Schmidta]
\label{lem:hs-compose-nuclear23}
Niech $H_2,H_1,H_0$ będą przestrzeniami Hilberta nad $\K$, a
$A:H_2\to H_1$ oraz $B:H_1\to H_0$ operatorami
Hilberta--Schmidta. Wówczas $BA:H_2\to H_0$ jest
operatorem nuklearnym w sensie
Definicji~\ref{def:nuclear-banach23}, przy czym
$\|BA\|_{\mathcal N}\le\|A\|_{\mathrm{HS}}\|B\|_{\mathrm{HS}}$.
\end{lemat}
\begin{proof}
Niech $(e_j)_{j\in J}$ będzie dowolną bazą ortonormalną $H_1$.
Dla każdego $x\in H_2$ zachodzi
$BAx=\sum_{j\in J}\langle Ax,e_j\rangle Be_j$,
gdyż rozwinięcie $Ax$ w bazie jest zbieżne w $H_1$,
a $B$ jest ograniczony. Ponadto
\[
 \sum_{j\in J}\|A^*e_j\|\,\|Be_j\|
 \le\Bigl(\sum_{j\in J}\|A^*e_j\|^2\Bigr)^{1/2}
     \Bigl(\sum_{j\in J}\|Be_j\|^2\Bigr)^{1/2}
 =\|A\|_{\mathrm{HS}}\|B\|_{\mathrm{HS}}.
\]
Przy sumowaniu po dowolnym zbiorze indeksów rozumiemy supremum
sum po jego skończonych podzbiorach. Skończoność powyższych sum
kwadratów zapewnia, że zbiory indeksów z
$\|A^*e_j\|>0$ lub $\|Be_j\|>0$ są co najwyżej przeliczalne:
dla każdego $k\ge1$ tylko skończenie wiele składników
sumy kwadratów może przekraczać $1/k$.
Pozostałe składniki przedstawienia są zerowe,
nawet gdy $H_1$ nie jest ośrodkowa. Funkcjonały
$x\mapsto\langle Ax,e_j\rangle$ mają normy $\|A^*e_j\|$,
a wektory $Be_j$ normy $\|Be_j\|$. Otrzymujemy zatem
przedstawienie nuklearne $BA$ z wymaganym oszacowaniem.
\end{proof}

\begin{twierdzenie}[Hilbertowskie kryterium nuklearności]
\label{thm:hilbert-nuclear23}
Niech $E$ będzie przestrzenią lokalnie wypukłą Hausdorffa
nad $\K$. Następujące warunki są równoważne:
\begin{enumerate}[label=\textup{(\roman*)},leftmargin=*]
\item $E$ jest nuklearna w sensie
Definicji~\ref{def:nuclear-lcs23};
\item każda ciągła seminorma $p$ na $E$ jest ograniczona
z góry przez pewną ciągłą seminormę hilbertowską $h$,
a dla każdej ciągłej seminormy hilbertowskiej $h$ istnieje
ciągła seminorma hilbertowska $r\ge h$, dla której
kanoniczny operator $H_r\to H_h$ jest operatorem
Hilberta--Schmidta.
\end{enumerate}
\end{twierdzenie}
\begin{proof}
\emph{(i)$\Rightarrow$(ii).} Ustalmy dowolną ciągłą seminormę $p$.
Z nuklearności wybieramy $q\ge p$ i przedstawienie operatora
$j_{q,p}:E_q\to E_p$ w postaci
\[
 j_{q,p}z=\sum_{j\ge1}a_j\ell_j(z)y_j,\qquad
 A=\sum_{j\ge1}|a_j|<\infty,\quad
 \|\ell_j\|\le1,\quad\|y_j\|\le1.
\]
Takie przedstawienie uzyskujemy, usuwając składniki
zerowe i dzieląc każdy funkcjonał oraz każdy wektor
przez jego normę, a iloczyn norm przenosząc do
współczynnika $a_j$. Można więc nawet przyjąć $a_j\ge0$. Jeśli $A=0$, to $p=0$ i wystarcza
seminorma zerowa. W przeciwnym razie kładziemy
$h(x)^2=A\sum_j|a_j|\,|\ell_j([x]_q)|^2$.
Jest to seminorma hilbertowska zadana przez formę
\[
 b(x,y)=A\sum_j|a_j|\ell_j([x]_q)
 \overline{\ell_j([y]_q)}.
\]
W przypadku rzeczywistym pomijamy sprzężenie.
Szereg jest bezwzględnie zbieżny, ponieważ
$|\ell_j([x]_q)|\le q(x)$ i $\sum_j|a_j|=A$.
Forma jest dodatnia półokreślona, a $h(x)\le Aq(x)$
dowodzi ciągłości seminormy.
Z nierówności Cauchy'ego--Schwarza wynika
\[
 p(x)=\bigl\|j_{q,p}[x]_q\bigr\|
 \le\sum_j|a_j|\,|\ell_j([x]_q)|
 \le h(x).
\]

Weźmy teraz dowolną ciągłą seminormę hilbertowską $h$.
Jeżeli $h=0$, wybieramy $r=0$ i żądany operator jest zerowy. W przeciwnym razie
wybieramy $q\ge h$ i nuklearne przejście $j_{q,h}:E_q\to H_h$.
Niech $j_{q,h}z=\sum_ja_j\ell_j(z)y_j$ będzie jego
znormalizowanym przedstawieniem, z $A=\sum_j|a_j|>0$.
Definiujemy $r(x)^2=A\sum_j|a_j|\,|\ell_j([x]_q)|^2$.
Powyższy argument daje $r\ge h$ i ciągłość $r$.
Na $E/\ker r$ określamy
\[
 U([x]_r)=(\sqrt{A|a_j|}\,\ell_j([x]_q))_j\in\ell^2,
\qquad
 Ve_j=A^{-1/2}\sqrt{|a_j|}\,\operatorname{sgn}(a_j)y_j,
\]
gdzie $\operatorname{sgn}(a_j)=a_j/|a_j|$ dla $a_j\ne0$,
a $\operatorname{sgn}(0)=0$. Pierwsze odwzorowanie jest dobrze określone
na klasach, ponieważ dla $x\in\ker r$ wszystkie
ważone współrzędne są zerowe. Jest izometrią
i przedłuża się do $U:H_r\to\ell^2$.
Drugie odwzorowanie, początkowo określone na ciągach
o skończonym nośniku, spełnia
$\sum_j\|Ve_j\|^2\le A^{-1}\sum_j|a_j|=1$.
Dla takiego ciągu $z=(z_j)$ mamy
\[
 \Big\|\sum_j z_jVe_j\Big\|
 \le\|z\|_2\Big(\sum_j\|Ve_j\|^2\Big)^{1/2}.
\]
Zatem $V$ przedłuża się do ograniczonego operatora
$\ell^2\to H_h$, który jest operatorem
Hilberta--Schmidta na mocy powyższej sumy.
Na gęstym obrazie $E$ zachodzi
\[
 VU([x]_r)=\sum_j a_j\ell_j([x]_q)y_j=[x]_h.
\]
Wobec ciągłości $VU=j_{r,h}$.
Obraz izometrii $U$ jest domknięty, a obraz dowolnej
bazy ortonormalnej $H_r$ przez $U$ można uzupełnić
do bazy ortonormalnej $\ell^2$. Stąd
$\|VU\|_{\mathrm{HS}}^2\le\|V\|_{\mathrm{HS}}^2$,
co dowodzi żądanej własności.

\emph{(ii)$\Rightarrow$(i).} Dla dowolnej ciągłej seminormy
$p$ wybieramy $h\ge p$ hilbertowską, a następnie
hilbertowskie $r\ge h$ oraz $s\ge r$ tak, aby
$H_s\to H_r$ i $H_r\to H_h$ były operatorami
Hilberta--Schmidta. Ich złożenie $H_s\to H_h$ jest
nuklearne na mocy Lematu~\ref{lem:hs-compose-nuclear23}.
Nierówność $p\le h$ określa ograniczone odwzorowanie
$H_h\to E_p$. Na mocy Stwierdzenia~\ref{prop:nuclear-ideal23}
złożenie operatora nuklearnego z ograniczonym
jest nuklearne, zatem $j_{s,p}:E_s=H_s\to E_p$ jest
nuklearny (identyfikujemy złożenie na gęstym obrazie $E$). To jest warunek z Definicji~\ref{def:nuclear-lcs23}.
\end{proof}

\begin{twierdzenie}[Trwałość względem podprzestrzeni i ilorazów]
\label{thm:nuclear-subquot23}
Niech $E$ będzie nuklearną przestrzenią lokalnie wypukłą
Hausdorffa. Każda podprzestrzeń liniowa $F\subseteq E$
z topologią indukowaną jest nuklearna. Jeśli $M\subseteq E$
jest domkniętą podprzestrzenią liniową, to iloraz $E/M$
z topologią ilorazową także jest nuklearny.
\end{twierdzenie}
\begin{proof}
Wykorzystamy dwie własności operatorów Hilberta--Schmidta.
Jeżeli $J:H_1\to H_0$ jest operatorem Hilberta--Schmidta,
to jego ograniczenie do domkniętej podprzestrzeni $R\subseteq H_1$
z wartościami w $\overline{J(R)}$ też ma tę własność:
wystarczy dopełnić bazę ortonormalną $R$ do bazy $H_1$.
Jeżeli natomiast $N\subseteq H_1$ i $L\subseteq H_0$
są domkniętymi podprzestrzeniami oraz $J(N)\subseteq L$,
to operator na ilorazach $H_1/N\to H_0/L$ jest
Hilberta--Schmidta. Istotnie, po identyfikacjach ilorazów
z dopełnieniami ortogonalnymi jest to
$P_{L^\perp}J|_{N^\perp}$.
Dla bazy ortonormalnej $(u_i)$ podprzestrzeni $N^\perp$
mamy $\sum_i\|P_{L^\perp}Ju_i\|^2\le\sum_i\|Ju_i\|^2
\le\|J\|_{\mathrm{HS}}^2$, uzupełniając bazę do bazy $H_1$.
Nie potrzeba założenia $J(N)=L$.

Niech $F\subseteq E$ będzie dowolną podprzestrzenią.
Dla każdej ciągłej seminormy $p_F$ na $F$ istnieje
ciągła seminorma $p_E$ na $E$ taka, że
$p_F\le p_E|_F$. Aby to wykazać, wybieramy otwarte, absolutnie
wypukłe otoczenie zera $U$ w $E$ z
$U\cap F\subseteq\{x\in F:p_F(x)<1\}$ i bierzemy
jego ciągły funkcjonał Minkowskiego $p_E$.
Dla $x\in F$ i $t>p_E(x)$ mamy $x/t\in U\cap F$,
więc $p_F(x)<t$. Przechodząc do infimum po $t$,
otrzymujemy $p_F(x)\le p_E(x)$. Z kryterium hilbertowskiego
(Twierdzenie~\ref{thm:hilbert-nuclear23}) wybieramy
hilbertowskie $h\ge p_E$, $r\ge h$, $s\ge r$,
dla których $H_s\to H_r$ i $H_r\to H_h$ są
Hilberta--Schmidta. Dla $t\in\{h,r,s\}$ przestrzeń
$H_{t|_F}$ identyfikuje się izometrycznie z domknięciem
obrazu $F$ w $H_t$: odwzorowanie $[x]_{t|_F}\mapsto[x]_t$
jest izometrią na gęstej podprzestrzeni tego domknięcia.
Przejścia kanoniczne przenoszą te domknięcia do
odpowiednich domknięć na mocy ciągłości i zgodności na $F$. Wskazana własność ograniczeń zapewnia,
że $H_{s|_F}\to H_{r|_F}$ oraz
$H_{r|_F}\to H_{h|_F}$ są Hilberta--Schmidta.
Ich złożenie jest nuklearne. Ponieważ
$p_F\le h|_F$, istnieje ograniczone odwzorowanie
$H_{h|_F}\to F_{p_F}$; jego złożenie z poprzednim
przejściem dowodzi nuklearności $F$.

Niech teraz $M\subseteq E$ będzie domknięta i
$Q=E/M$ z odwzorowaniem ilorazowym $\pi:E\to Q$.
Dla ciągłej seminormy $p_Q$ na $Q$ zastosujmy kryterium
do $p_Q\circ\pi$ i wybierzmy hilbertowskie seminormy
$h\ge p_Q\circ\pi$, $r\ge h$, $s\ge r$,
których kolejne przejścia są Hilberta--Schmidta.
Dla $t\in\{h,r,s\}$ kładziemy
\[
 \bar t(\pi x)=\inf_{m\in M}t(x+m).
\]
Wartość nie zależy od reprezentanta klasy $\pi x$.
Niech $N_t$ oznacza domknięcie obrazu $M$ w $H_t$.
Wówczas
\[
 \bar t(\pi x)=\operatorname{dist}([x]_t,N_t)
 =\|P_{N_t^\perp}[x]_t\|.
\]
Ostatnia równość wynika z twierdzenia o rzucie
ortogonalnym. Jest to seminorma hilbertowska, zadana
przez iloczyn skalarny obrazów przez
$\pi x\mapsto P_{N_t^\perp}[x]_t$.
Jej ciągłość wynika z równości
$\{y:\bar t(y)<\varepsilon\}=
\pi(\{x:t(x)<\varepsilon\})$ i otwartości $\pi$.
Odwzorowanie $Q/\ker\bar t\to H_t/N_t$ dane przez
$[\pi x]_{\bar t}\mapsto [x]_t+N_t$ jest izometryczne
i ma gęsty obraz, ponieważ obraz $E$ w $H_t$ jest gęsty.
Przedłuża się więc do izometrycznego izomorfizmu
$H_{\bar t}\simeq H_t/N_t$.
Ponadto $\bar s\ge\bar r\ge\bar h\ge p_Q$,
gdyż $h(x+m)\ge p_Q(\pi(x+m))=p_Q(\pi x)$
dla każdego $m\in M$.
Operator $H_s\to H_r$ odwzorowuje $N_s$ w $N_r$,
podobnie $H_r\to H_h$ odwzorowuje $N_r$ w $N_h$.
Z własności operatorów na ilorazach otrzymujemy
przejścia Hilberta--Schmidta
$H_{\bar s}\to H_{\bar r}\to H_{\bar h}$.
Ich złożenie jest nuklearne, zaś
$H_{\bar h}\to Q_{p_Q}$ jest ograniczone. Zatem
$j_{\bar s,p_Q}:Q_{\bar s}\to Q_{p_Q}$ jest nuklearny,
co kończy dowód.
\end{proof}

\begin{lemat}[Nuklearność przeliczalnej sumy prostej]
\label{lem:nuclear-direct-sum23}
Niech $E_1,E_2,\ldots$ będą nuklearnymi przestrzeniami
lokalnie wypukłymi Hausdorffa. Wówczas ich lokalnie wypukła
suma prosta $E=\bigoplus_{n\ge1}E_n$, z najsilniejszą
topologią lokalnie wypukłą, dla której wszystkie kanoniczne
włożenia $E_n\to E$ są ciągłe, jest nuklearna.
\end{lemat}
\begin{proof}
Suma prosta jest Hausdorffa: dla niezerowego
$x=(x_n)$ wybieramy niezerową współrzędną $x_n$
i ciągłą seminormę na $E_n$ niezerową na $x_n$.
Jej złożenie z rzutem na $n$-tą współrzędną jest
ciągłe na $E$, ponieważ jego ograniczenia do wszystkich
składników są ciągłe.
Niech $p$ będzie ciągłą seminormą na $E$,
$p_n=p|_{E_n}$. Dla każdego $n$ z nuklearności $E_n$
wybieramy ciągłą seminormę $q_n^0\ge p_n$, dla której
$j_{q_n^0,p_n}$ jest nuklearny. Zastępując $q_n^0$
przez $q_n=c_nq_n^0$ dla dostatecznie dużego $c_n\ge1$,
możemy zapewnić
$\|j_{q_n,p_n}\|_{\mathcal N}<2^{-n}$:
przeskalowanie normy w przestrzeni źródłowej przez
$c_n$ dzieli normę każdego funkcjonału przez $c_n$,
więc dzieli także normę nuklearną przejścia przez $c_n$
(obie nierówności wynikają ze skalowania w obie strony).
Definiujemy na algebraicznej sumie prostej
$q((x_n))=\sum_nq_n(x_n)$; jest to suma skończona.
Seminorma $q$ jest ciągła, ponieważ jej ograniczenie
do każdego $E_n$ jest ciągłe, a topologia sumy prostej
jest końcową topologią lokalnie wypukłą. Z nierówności
trójkąta wynika $p\le\sum_n p_n\le q$.

Mamy $\ker q=\bigoplus_n\ker q_n$, a norma na ilorazie
$E/\ker q$ jest sumą norm współrzędnych.
Uzupełnienie $E/\ker q$ identyfikuje się izometrycznie
z $\bigoplus_n^{\ell^1}(E_n)_{q_n}$: jest to przestrzeń
Banacha, a obcięcie ogona i przybliżenie skończonej
liczby współrzędnych elementami $E_n/\ker q_n$
dowodzi gęstości obrazu $E/\ker q$.
Dla ustalonego $n$ kanoniczny operator
$(E_n)_{q_n}\to E_p$ jest złożeniem nuklearnego
$j_{q_n,p_n}$ i izometrycznego osadzenia
$(E_n)_{p_n}\to E_p$: na gęstym ilorazie
norma obrazu jest dokładnie $p_n(x_n)$.
Ze Stwierdzenia~\ref{prop:nuclear-ideal23} (własność
ideałowa) wynika, że norma nuklearna złożenia jest
mniejsza niż $2^{-n}$.
Wybieramy jego przedstawienie nuklearne, którego suma
norm składników jest mniejsza niż $2^{-n}$.
Każdy funkcjonał występujący w tym przedstawieniu
przedłużamy na sumę $\ell^1$ przez złożenie
z rzutem współrzędnym normy co najwyżej $1$ (rzut
na zerowy składnik ma normę $0$).
Indeksy par $(n,j)$ tworzą zbiór przeliczalny.
Po połączeniu wszystkich przedstawień w jeden szereg
suma iloczynów norm funkcjonałów i wektorów jest
ograniczona przez $\sum_{n\ge1}2^{-n}<\infty$.
Szereg wyznacza więc operator nuklearny bez potrzeby
odwoływania się do zupełności przestrzeni operatorów
nuklearnych. Ich suma jest dokładnie
kanonicznym operatorem $E_q\to E_p$:
najpierw sprawdzamy tę równość dla wektorów o skończonym
nośniku, a następnie korzystamy z ciągłości i gęstości.
To dowodzi nuklearności $E$.
\end{proof}

\begin{twierdzenie}[Trwałość względem przeliczalnych granic induktywnych]
\label{thm:nuclear-indlim23}
Niech $(E_n,u_n)_{n\ge1}$ będzie ciągiem nuklearnych
przestrzeni lokalnie wypukłych Hausdorffa i ciągłych
odwzorowań liniowych $u_n:E_n\to E_{n+1}$.
Jeśli lokalnie wypukła granica induktywna
$E=\varinjlim_n E_n$ jest Hausdorffa, to $E$ jest nuklearna.
Nie zakładamy metryzowalności $E$ ani ścisłości granicy.
\end{twierdzenie}
\begin{proof}
Niech $S=\bigoplus_{n\ge1}E_n$ będzie lokalnie wypukłą
sumą prostą, a $\iota_n:E_n\to E$ odwzorowaniami
kanonicznymi. Definiujemy
$\Pi:S\to E$ przez
$\Pi((x_n))=\sum_n\iota_n(x_n)$; suma jest skończona.
Z końcowej własności topologii sumy prostej wynika
ciągłość $\Pi$, a z definicji granicy induktywnej
wynika jej surjektywność.

Topologia ilorazowa na $S/\ker\Pi$ ma dokładnie tę samą
własność uniwersalną co lokalnie wypukła topologia
granicy induktywnej: odwzorowanie liniowe $T:E\to G$
do dowolnej przestrzeni lokalnie wypukłej $G$ jest ciągłe
wtedy i tylko wtedy, gdy wszystkie $T\circ\iota_n$ są ciągłe.
Zatem $E\simeq S/\ker\Pi$ jako przestrzenie lokalnie
wypukłe. Ponieważ $E$ jest Hausdorffa,
$\ker\Pi=\Pi^{-1}(\{0\})$ jest domknięte w $S$.
Suma prosta $S$ jest nuklearna na mocy
Lematu~\ref{lem:nuclear-direct-sum23}, a iloraz
przez domkniętą podprzestrzeń jest nuklearny na mocy
Twierdzenia~\ref{thm:nuclear-subquot23}. Wynika stąd
nuklearność $E$. Ścisłość i metryzowalność granicy
nie zostały wykorzystane.
\end{proof}

\begin{twierdzenie}[Nuklearność $\mathcal D_K$ i $\mathcal D$]
\label{thm:nuclear-D23}
Niech $d\ge1$, $\Omega\subseteq\R^d$ będzie niepustym
zbiorem otwartym, a $K\Subset\Omega$ zbiorem zwartym.
Wówczas $\mathcal D_K(\Omega)$ ze standardową topologią
Frécheta i $\mathcal D(\Omega)$ ze standardową topologią
LF są nuklearnymi przestrzeniami lokalnie wypukłymi.
\end{twierdzenie}
\begin{proof}
Dla $K=\varnothing$ przestrzeń $\mathcal D_K(\Omega)$
jest zerowa, a więc nuklearna. Dla niepustego $K$
rozważmy najpierw $\mathcal D_K(\Omega)$ jako
podprzestrzeń $\mathcal E(\Omega)=C^\infty(\Omega)$.
Ponieważ $K$ jest domknięty w $\Omega$,
warunek $\supp f\subseteq K$ jest równoważny
$f(x)=0$ dla wszystkich $x\in\Omega\setminus K$.
A zatem
\[
 \mathcal D_K(\Omega)=
 \bigcap_{x\in\Omega\setminus K}\ker(\mathrm{ev}_x),
 \qquad \mathrm{ev}_x(f)=f(x).
\]
Każda ewaluacja jest ciągła w topologii $C^\infty$,
więc $\mathcal D_K(\Omega)$ jest podprzestrzenią
domkniętą. Ponadto wszystkie pochodne jej elementów
znikają poza $K$, stąd dla dowolnego $L\Subset\Omega$
\[
 p_{L,m}(f)\le p_{K,m}(f),\qquad
 f\in\mathcal D_K(\Omega).
\]
Ponieważ $K$ samo jest zwarte, rodzina $p_{K,m}$
jest zarazem rodziną ograniczeń seminorm z $C^\infty(\Omega)$.
Topologia indukowana jest więc dokładnie standardową
topologią Frécheta na $\mathcal D_K(\Omega)$.
Twierdzenia~\ref{thm:nuclear-cinfty23} i
\ref{thm:nuclear-subquot23} dowodzą jej nuklearności.

Wybierzmy wyczerpanie $\Omega$ zbiorami zwartymi
$K_1\subseteq\operatorname{int}K_2\subseteq\cdots\Subset\Omega$
takie, że każdy zwarty podzbiór $\Omega$ zawiera się
w pewnym $K_n$. Definicja~\ref{def:D-LF18} identyfikuje standardową
topologię LF na $\mathcal D(\Omega)$ z topologią
lokalnie wypukłej granicy induktywnej
$\varinjlim_n\mathcal D_{K_n}(\Omega)$, a
Twierdzenie~\ref{thm:D-stage18} zapewnia warunek Hausdorffa.
Włożenia kolejnych etapów są ciągłe.
Jej nuklearność wynika teraz z
Twierdzenia~\ref{thm:nuclear-indlim23}.
\end{proof}

\begin{uwaga}[Dlaczego sam argument z $C^\infty$ nie wystarcza]
\label{rem:nuclear-D-summary23}
Nuklearność $\mathcal D_K$ jest konsekwencją trwałości
względem podprzestrzeni, ale nuklearność
$\mathcal D=\varinjlim\mathcal D_{K_n}$ wymaga osobnego
wyniku dotyczącego granic induktywnych. Powyższe dowody
ujawniają wszystkie użyte przejścia. Nie wynika z tego
równość topologii induktywnej tensorowej $i$ z topologią
rzutową $\pi$: dotyczą one innych własności uniwersalnych.
\end{uwaga}

\subsection{Funkcje Hermite'a i nuklearność przestrzeni Schwartza}
\label{sec:hermite-supp23}

\begin{lemat}[Rozwinięcie Hermite'a i seminormy]
\label{lem:hermite-schwartz23}
Niech $d\ge1$ będzie liczbą całkowitą. Dla $n\in\N_0$ definiujemy
\[
 P_n(x)=(-1)^n e^{x^2}\frac{\mathrm d^n}{\mathrm dx^n}e^{-x^2},
 \qquad h_n(x)=\frac{P_n(x)e^{-x^2/2}}{(2^n n!\sqrt\pi)^{1/2}}.
\]
Dla $\nu\in\N_0^d$ niech $h_\nu(x)=\prod_{j=1}^d h_{\nu_j}(x_j)$
oraz $|\nu|=\sum_j\nu_j$. Wówczas $(h_\nu)_\nu$
jest bazą ortonormalną $L^2(\R^d;\K)$, a odwzorowanie
\[
 H:\mathcal S(\R^d;\K)\to s(\N_0^d),\qquad
 H(f)=(\langle f,h_\nu\rangle_{L^2})_\nu
\]
jest izomorfizmem liniowo-topologicznym. Iloczyn skalarny
jest liniowy w pierwszym argumencie.
\end{lemat}
\begin{proof}
W jednym wymiarze z definicji $P_0=1$ i $P_{n+1}=2xP_n-P_n'$.
Indukcja daje stopień $n$ i współczynnik wiodący $2^n$.
Dla $n\ge m$, całkując $n$ razy przez części, otrzymujemy
\[
 \int_\R P_nP_m e^{-x^2}\,\mathrm d\lambda_1
 =\int_\R P_m^{(n)}e^{-x^2}\,\mathrm d\lambda_1.
\]
Wyrazy brzegowe znikają, gdyż pochodne Gaussa są
wielomianami pomnożonymi przez Gaussa. Wynik wynosi zero
dla $n>m$ oraz $2^n n!\sqrt\pi$ dla $n=m$.
Symetria daje ortogonalność również dla $n<m$.

Jeżeli $f\in L^2(\R;\C)$ jest ortogonalna do wszystkich
$h_n$, to $\int_\R f(x)x^n e^{-x^2/2}\,\mathrm d\lambda_1(x)=0$
dla każdego $n$, ponieważ $(P_0,\ldots,P_n)$ jest bazą
wielomianów stopnia co najwyżej $n$. Niech $g(x)=f(x)e^{-x^2/2}$.
Z nierówności Cauchy'ego--Schwarza mamy $g\in L^1\cap L^2$.
Dla każdego $R>0$ funkcja $e^{-x^2/2+R|x|}$ należy do $L^2$,
więc $|g(x)|e^{R|x|}$ jest całkowalna.
Rozwijając wykładniczą i stosując zbieżność zmajoryzowaną
dla $|z|\le R$, otrzymujemy
\[
 \int_\R g(x)e^{zx}\,\mathrm d\lambda_1(x)
 =\sum_{n=0}^\infty\frac{z^n}{n!}
 \int_\R x^ng(x)\,\mathrm d\lambda_1(x)=0
 \qquad(z\in\C).
\]
Dla $z=-i\xi$ klasyczna transformata całkowa $g$ znika.
Dla $\varphi\in\mathcal S$ z Fubiniego, z majorantą
$|g(x)\varphi(\xi)|$, mamy
\[
 \langle\mathcal F T_g,\varphi\rangle
 =(2\pi)^{-1/2}\int_\R\varphi(\xi)
 \int_\R g(x)e^{-ix\xi}\,\mathrm d\lambda_1(x)
 \,\mathrm d\lambda_1(\xi)=0.
\]
Twierdzenie~\ref{thm:fourier-auto-tempered22} i
Stwierdzenie~\ref{prop:schwartz-embeddings22} dają
$g=0$ prawie wszędzie, więc $f=0$. Przypadek rzeczywisty
wynika przez rozpatrzenie funkcji jako zespolonej.
W wymiarze $d$ ortonormalność iloczynów wynika z Fubiniego.
Zupełność wykazujemy indukcyjnie: jeśli $f$ jest
ortogonalna do wszystkich iloczynów, to dla każdego $n$
funkcja $g_n(x')=\int_\R f(x',t)h_n(t)\,\mathrm d\lambda_1(t)$
należy do $L^2(\R^{d-1})$ (Cauchy--Schwarz i Fubini)
i jest ortogonalna do wszystkich iloczynów w $d-1$
wymiarach. Zatem $g_n=0$ prawie wszędzie.
Przeliczalność rodziny $n$ pozwala wybrać wspólny zbiór
pełnej miary, na którym $f(x',\cdot)\in L^2(\R)$ i
wszystkie jej współczynniki znikają. Zupełność
jednowymiarowa i Fubini dają $f=0$.

Niech $L=-\Delta+|x|^2$ na $\mathcal S(\R^d)$.
Z rekurencji indukcyjnie dostajemy $P_n'=2nP_{n-1}$ dla $n\ge1$:
$P_{n+1}'=2P_n+2n(2xP_{n-1}-P_{n-1}')=2(n+1)P_n$,
z przypadkiem początkowym $P_1'=2P_0$.
Wprowadzamy $a_j=\partial_j+x_j$ i $a_j^*=-\partial_j+x_j$,
formalnie sprzężone na $\mathcal S$ przez całkowanie przez części.
Istotnie, w jednym wymiarze, oznaczając
$b_n=(2^n n!\sqrt\pi)^{-1/2}$, mamy
\[
 (\partial_x+x)h_n=b_n P_n'e^{-x^2/2},
 \qquad
 (-\partial_x+x)h_n=b_n(2xP_n-P_n')e^{-x^2/2}.
\]
Dla $n\ge1$ zachodzi $2n b_n/b_{n-1}=\sqrt{2n}$,
a dla $n\ge0$ mamy $b_n/b_{n+1}=\sqrt{2(n+1)}$.
Rekurencja, wzór na pochodną i te równości dają
\[
 a_jh_\nu=\sqrt{2\nu_j}\,h_{\nu-e_j},\qquad
 a_j^*h_\nu=\sqrt{2(\nu_j+1)}\,h_{\nu+e_j},
\]
gdzie pierwszy wyraz oznacza zero, gdy $\nu_j=0$.
Rozwinięcie iloczynu daje
$a_j^*a_j=-\partial_j^2+x_j^2-I$, gdyż
$\partial_j(x_j u)=u+x_j\partial_j u$.
Ponieważ $L=\sum_j a_j^*a_j+dI$, mamy
$Lh_\nu=(2|\nu|+d)h_\nu$.
Dla skończonej kombinacji $u=\sum_\nu c_\nu h_\nu$
i słowa $W$ długości $l$ z operatorów $a_j,a_j^*$
każdy $Wh_\nu$ jest zerem albo wielokrotnością
$h_{\nu+\gamma}$, gdzie $\gamma$ zależy wyłącznie od $W$.
Moduł współczynnika nie przekracza
$(2(|\nu|+l+1))^{l/2}$, ponieważ każdy krok zwiększa
stopień indeksu co najwyżej o jeden.
Niezerowe obrazy różnych indeksów są ortogonalne.
Dlatego
\[
 \|Wu\|_2^2\le C_l\sum_\nu(1+|\nu|)^l|c_\nu|^2.
\]
Z $x_j=(a_j+a_j^*)/2$ i $\partial_j=(a_j-a_j^*)/2$
oraz nierówności trójkąta wynika oszacowanie normy
każdego iloczynu $l$ operatorów mnożenia i różniczkowania
przez $C_l'\|(1+L)^r u\|_2$, gdy $2r\ge l$.

Dla całkowitego $s>d/2$ i $v\in\mathcal S$, ze wzoru
odwrotnego Fouriera (Twierdzenie~\ref{thm:fourier-schwartz17}),
nierówności Cauchy'ego--Schwarza i twierdzenia
Plancherela (Twierdzenie~\ref{thm:plancherel22}) wynika
\[
 \|v\|_\infty\le (2\pi)^{-d/2}\|\widehat v\|_1
 \le C_{d,s}\Big(\sum_{|\gamma|\le s}
 \|\partial^\gamma v\|_2^2\Big)^{1/2}.
\]
Stosujemy wagę $(1+|\xi|^2)^{-s/2}$, której kwadrat
jest całkowalny; rozwinięcie $(1+|\xi|^2)^s$ daje
skończoną sumę kwadratów monomów z dodatnimi współczynnikami.
Stosując to do $v=x^\alpha\partial^\beta u$ oraz
$(1+|x|)^N\le C_{N,d}\sum_{|\alpha|\le N}|x^\alpha|$,
co wynika z rozwinięcia $(1+\sum_j|x_j|)^N$,
otrzymujemy dla $2r\ge N+m+s$ poniższe oszacowanie.
Istotnie, reguła Leibniza rozpisuje każdą pochodną
$\partial^\gamma(x^\alpha\partial^\beta u)$
na skończoną sumę wyrazów, w których łączna liczba
mnożeń przez współrzędne i różniczkowań nie przekracza
$N+m+s$. Stosujemy do nich poprzednie oszacowanie
dla słów w operatorach $a_j,a_j^*$. Wybieramy
$r=\lceil(N+m+s)/2\rceil$, zatem
\[
 Q_{N,m}(u)\le C_{N,m,d,s}\|(1+L)^r u\|_2.
\]
Nierówność ta dotyczy na razie skończonych kombinacji.

Dla dowolnego $f\in\mathcal S$, rozwijając $(1+L)^r$
w skończoną sumę operatorów $x^\alpha\partial^\beta$
za pomocą reguły Leibniza, szacujemy każdy składnik
w $L^2$ przez stałą razy $Q_{|\alpha|+d+1,|\beta|}(f)$,
gdyż $(1+|x|)^{-d-1}\in L^2$.
Zatem $\|(1+L)^r f\|_2\le C_r Q_{N_r,m_r}(f)$.
Z formalnej symetrii $L$ na $\mathcal S$ mamy
$\langle(1+L)^r f,h_\nu\rangle
=(1+2|\nu|+d)^r\langle f,h_\nu\rangle$.
Parseval daje
\[
 \|(1+L)^r f\|_2^2
 =\sum_\nu(1+2|\nu|+d)^{2r}|H(f)_\nu|^2.
\]
Stąd $p_k(H(f))\le\|(1+L)^k f\|_2$, co dowodzi
ciągłości $H$ i przynależności współczynników do $s$.
Różnowartościowość wynika z zupełności bazy i ciągłości $f$.

Dla $c\in s(\N_0^d)$ i całkowitego $t$ z $2t>d$ mamy
\[
 \sum_\nu(1+2|\nu|+d)^{2r}|c_\nu|^2
 \le C_{r,d}p_{r+t}(c)^2
 \sum_\nu(1+|\nu|)^{-2t}<\infty
\]
na mocy zliczania wielowskaźników w dowodzie
Twierdzenia~\ref{thm:nuclear-s23}.
Sumy $u_M=\sum_{|\nu|\le M}c_\nu h_\nu$ są Cauchy'ego
w każdej normie $\|(1+L)^r\cdot\|_2$, a więc również
we wszystkich seminormach Schwartza. Zupełność
$\mathcal S$ daje granicę $u$. Ciągłość włożenia
$\mathcal S\to L^2$ (Stwierdzenie~\ref{prop:schwartz-operators17})
zapewnia $H(u)_\nu=c_\nu$ dla każdego $\nu$.
Otrzymujemy surjektywność $H$ i, przechodząc do granicy
w oszacowaniach sum częściowych,
$Q_{N,m}(H^{-1}c)\le C p_{r+t}(c)$.
Dowodzi to ciągłości odwrotności.
\end{proof}

\begin{wniosek}[Nuklearność przestrzeni Schwartza]
\label{cor:nuclear-schwartz23}
Przestrzeń Schwartza $\mathcal S(\R^d)$ jest nuklearną
przestrzenią Frécheta--Montela.
\end{wniosek}
\begin{proof}
Z Lematu~\ref{lem:hermite-schwartz23} mamy izomorfizm
liniowo-topologiczny $H:\mathcal S\to s(\N_0^d)$.
Dla ciągłej seminormy $p$ na $\mathcal S$ seminorma
$p\circ H^{-1}$ jest ciągła na $s(\N_0^d)$.
Z Twierdzenia~\ref{thm:nuclear-s23} wybieramy dominującą
seminormę $q'$ z nuklearnym przejściem między
uzupełnieniami. Po wzięciu $q=q'\circ H$ izometrie
indukowane przez $H$ na uzupełnieniach przenoszą
to przejście na $j_{q,p}$; własność ideałowa
(Stwierdzenie~\ref{prop:nuclear-ideal23}) daje jego
nuklearność. Zupełność, metryzowalność, beczkowość
i względna zwartość zbiorów ograniczonych są zachowywane
przez izomorfizmy liniowo-topologiczne, zatem również
własność Frécheta--Montela przechodzi na $\mathcal S$.
\end{proof}

\begin{wniosek}[Współczynniki Hermite'a dystrybucji temperowanych]
\label{cor:dt-hermite-tempered}
Niech $h_\nu$ będą funkcjami z Lematu~\ref{lem:hermite-schwartz23}.
Odwzorowanie
\[
 T\longmapsto (T(h_\nu))_{\nu\in\N_0^d}
\]
jest liniową bijekcją przestrzeni $\mathcal S'(\R^d)$
na przestrzeń ciągów $b=(b_\nu)$ o wzroście
wielomianowym, czyli takich, że dla pewnych
$C>0$, $k\in\N_0$ zachodzi
$|b_\nu|\le C(1+|\nu|)^k$ dla każdego $\nu$.
Odwrotna bijekcja jest dana wzorem
\[
 \langle T_b,\varphi\rangle
 =\sum_{\nu\in\N_0^d}b_\nu\langle\varphi,h_\nu\rangle_{L^2}
 \qquad(\varphi\in\mathcal S).
\]
Szereg jest bezwzględnie zbieżny.
\end{wniosek}
\begin{proof}
Niech $H$ będzie izomorfizmem z przywołanego lematu.
Ciągły liniowy funkcjonał $T\circ H^{-1}$ na
$s(\N_0^d)$ jest ograniczony przez $Cp_k$ dla
pewnych $C,k$, gdyż normy $p_k$ są rosnące i
wyznaczają topologię. Dla ciągu jednostkowego $e_\nu$
mamy $H^{-1}e_\nu=h_\nu$, a więc
$|T(h_\nu)|\le Cp_k(e_\nu)=C(1+|\nu|)^k$.

Odwrotnie, dla ciągu $b$ o takim wzroście
wybieramy całkowite $m>k+d$. Dla każdego $c\in s$
mamy
\[
 \sum_\nu |b_\nu c_\nu|
 \le Cp_m(c)\sum_\nu(1+|\nu|)^{k-m}<\infty,
\]
gdzie zbieżność ostatniej sumy wykazano w dowodzie
Twierdzenia~\ref{thm:nuclear-s23}.
Wzór z tezy definiuje zatem funkcjonał liniowy
i ciągły na $\mathcal S$ przez ciągłość $H$.
Ortonormalność daje $T_b(h_\nu)=b_\nu$.
Jednoznaczność wynika z gęstości skończonych
kombinacji funkcji Hermite'a w $\mathcal S$,
wykazanej w dowodzie Lematu~\ref{lem:hermite-schwartz23}:
ciągły funkcjonał jest wyznaczony przez wartości
na tej podprzestrzeni. Ponieważ dystrybucje są
liniowe, we wzorze na działanie $T_b$ nie sprzęgamy
współczynników $\langle\varphi,h_\nu\rangle_{L^2}$.
\end{proof}

% Materiał z AF 25; fragment do rozdziału o przestrzeniach Montela.
\subsection{Konsekwencje dla topologii przestrzeni sprzężonych}

\begin{uwaga}[Zupełność i lokalna wypukłość]
\label{rem:tabela-struktur25}
$\mathcal E(\Omega)$ i $\mathcal S(\R^d)$ są nuklearnymi
przestrzeniami Frécheta, a $\mathcal D(\Omega)$ jest nuklearną,
zupełną, niemetryzowalną przestrzenią LF
(Twierdzenia~\ref{thm:D-complete18}, \ref{thm:D-nonmetric18}).
Wszystkie trzy są przestrzeniami Montela
(Twierdzenie~\ref{thm:montel-test23}).
Nuklearność $\mathcal E$ i $\mathcal S$ wynika odpowiednio
z Twierdzenia~\ref{thm:nuclear-cinfty23}
i Wniosku~\ref{cor:nuclear-schwartz23}.
Trwałość nuklearności względem podprzestrzeni, ilorazów
i przeliczalnych granic induktywnych udowodniono w
Twierdzeniach~\ref{thm:nuclear-subquot23} i~\ref{thm:nuclear-indlim23};
nuklearność $\mathcal D_K$ i $\mathcal D$ wynika z
Twierdzenia~\ref{thm:nuclear-D23}. Nie należy mylić nuklearności
z własnością Montela: te dwie własności były w skrypcie
uzasadniane oddzielnie.
\end{uwaga}
\begin{twierdzenie}[Zbieżność ciągów w przestrzeniach sprzężonych]
\label{thm:weak-strong-sequences25}
Niech $X$ będzie beczkową przestrzenią lokalnie wypukłą
Hausdorffa nad $\K$, w której każdy zbiór ograniczony
ma zwarte domknięcie (w szczególności przestrzenią Montela).
Dla ciągu $(T_n)$ w $X'$ i $T\in X'$ następujące
warunki są równoważne: $T_n\to T$ w topologii
$\sigma(X',X)$ oraz $T_n\to T$ w topologii $\beta(X',X)$.
Stwierdzenie stosuje się w szczególności do
$X=\mathcal D(\Omega)$, $X=\mathcal E(\Omega)$ i
$X=\mathcal S(\R^d)$.
\end{twierdzenie}
\begin{proof}
Założenia oznaczają, że $X$ jest przestrzenią Montela.
Zbieżność słaba-* pociąga za sobą silną na mocy wniosku
\ref{cor:montel-dual-seq23}. Odwrotna implikacja wynika
z ograniczoności zbiorów jednoelementowych.
Zastosowania do $\mathcal D$, $\mathcal E$ i $\mathcal S$
wynikają z twierdzenia \ref{thm:montel-test23}.
\end{proof}


\begin{uwaga}[Zbieżność ciągów a równość topologii]
\label{rem:sequence-net25}
Twierdzenie~\ref{thm:weak-strong-sequences25} nie oznacza
równości topologii $\sigma$ i $\beta$.
W dowodzie używamy punktowej ograniczoności całej
rodziny wyrazów ciągu zbieżnego. Dla sieci zbieżnej
punktowo nie wynika stąd ograniczoność wszystkich
jej wyrazów.

Dla każdej z trzech rozważanych przestrzeni $X$
istnieje ograniczony zbiór $B_0$ zawierający nieskończenie
wiele liniowo niezależnych funkcji. Dla $\mathcal D$
jest to zbiór skonstruowany w dowodzie Twierdzenia~
\ref{thm:weak-strong-different19}; dla $\mathcal E$
i $\mathcal S$ korzystamy z jego obrazu przez ciągłe
włożenie z Twierdzenia~\ref{thm:zanurzenia25}.
Niech $I$ będzie zbiorem par $(F,n)$, gdzie
$F\subset X$ jest skończony, $n\in\N$, uporządkowanym
przez inkluzję pierwszego czynnika i zwykły porządek
drugiego. Dla każdej pary wybieramy
$u\in B_0\setminus\operatorname{lin}F$.
Skończenie wymiarowa podprzestrzeń $\operatorname{lin}F$
jest domknięta, więc twierdzenie Hahna--Banacha
zapewnia $T_{F,n}\in X'$ z
$T_{F,n}|_{\operatorname{lin}F}=0$ i $T_{F,n}(u)=n$.
Dla ustalonego $x\in X$ od indeksu $(\{x\},1)$
wszystkie wartości $T_{F,n}(x)$ są zerowe.
Sieć zbiega więc słabo-* do zera. Jednak
$\sup_{v\in B_0}|T_{F,n}(v)|\ge n$, więc
nie zbiega silnie do zera. Topologia silna
jest zatem ściśle silniejsza w każdym z trzech przypadków.
\end{uwaga}

\noindent\emph{Powiązania tematyczne w Bibliotece matematycznej:} \atlasref{PDE-07::def:pde7-test}{PDE-07: przestrzeń funkcji testowych}.

\setcounter{section}{17}
% Rozdział aktywny po restrukturyzacji DT.
% Wersja źródłowa: audyt II, jednostka 24

\section{Topologiczne iloczyny tensorowe przestrzeni lokalnie wypukłych}
\label{sec:tensor24}

\providecommand{\atlasref}[2]{#2}


\noindent\textbf{Wymagane wyniki.} Korzystamy z twierdzenia
Hahna--Banacha, definicji nuklearności przez przejścia seminormowe,
rozwinięć Fouriera i Hermite'a. Twierdzenia trwałości nuklearności
dowiedziono w rozdziale~\ref{sec:montel-nuklearne23}; tu korzystamy z nich
w teorii iloczynów tensorowych.
\medskip

W rachunku tensorowym przestrzeni skończenie wymiarowych pojęcie
$E\otimes F$ jest czysto algebraiczne. W analizie funkcjonalnej
potrzebujemy dodatkowo topologii: od niej zależy, które szeregi
tensorów są zbieżne, które formy dwuliniowe są ciągłe i czy
otrzymana przestrzeń ma interpretację jako przestrzeń funkcji.
W tym rozdziale $\K\in\{\R,\C\}$, a przestrzenie lokalnie wypukłe
są Hausdorffa, chyba że zaznaczono inaczej. Zapis $E'$ oznacza
przestrzeń sprzężoną topologicznie z $E$.

\subsection{Produkt tensorowy algebraiczny i własność uniwersalna}

\noindent\emph{Powiązanie tematyczne w Bibliotece matematycznej:} \atlasref{AA-18::aa:tw:tensor-universal}{AA-18: algebraiczna własność uniwersalna iloczynu tensorowego}.


\begin{definicja}[Produkt tensorowy algebraiczny]
\label{def:alg-tensor24}
Niech $E,F$ będą przestrzeniami liniowymi nad $\K$.
Ich \emph{produkt tensorowy algebraiczny} $E\otimes F$ jest
przestrzenią liniową wyposażoną w odwzorowanie dwuliniowe
$\otimes:E\times F\to E\otimes F$, $ (x,y)\mapsto x\otimes y$,
takie, że dla każdej przestrzeni liniowej $G$ i każdego
odwzorowania dwuliniowego $b:E\times F\to G$ istnieje dokładnie
jedno odwzorowanie liniowe $\widetilde b:E\otimes F\to G$
spełniające $\widetilde b(x\otimes y)=b(x,y)$ dla wszystkich
$x\in E$, $y\in F$.
\end{definicja}

\begin{stwierdzenie}[Istnienie i jednoznaczność]
\label{prop:alg-existence24}
Dla dowolnych przestrzeni liniowych $E,F$ produkt tensorowy
algebraiczny istnieje i jest jednoznaczny z dokładnością do
jednoznacznego izomorfizmu zachowującego tensory elementarne.
Każdy jego element jest skończoną sumą $\sum_{j=1}^N x_j\otimes y_j$.
\end{stwierdzenie}
\begin{proof}
Niech $V$ będzie wolną przestrzenią liniową o bazie
$[x,y]$, indeksowanej parami $(x,y)\in E\times F$.
Przez $N\subseteq V$ oznaczmy podprzestrzeń rozpiętą przez
relacje $[x+x',y]-[x,y]-[x',y]$,
$[x,y+y']-[x,y]-[x,y']$, $[ax,y]-a[x,y]$
oraz $[x,ay]-a[x,y]$ dla $a\in\K$.
Przyjmujemy $E\otimes F=V/N$ i $x\otimes y=[x,y]+N$.
Każde $b$ przedłuża się liniowo na $V$ i zeruje się na $N$,
więc przechodzi do jednoznacznego odwzorowania na ilorazie.
To daje własność uniwersalną. W dowolnej realizacji $T$ tensory elementarne rozpinają
całą przestrzeń: odwzorowanie ilorazowe z $T$ do ilorazu
przez ich rozpiętość i odwzorowanie zerowe mają to samo
złożenie z odwzorowaniem dwuliniowym $E\times F\to T$.
Jednoznaczność we własności uniwersalnej oznacza,
że iloraz jest zerowy. Jeśli $T$ i $T'$ są dwiema
realizacjami, własności uniwersalne dają odwzorowania liniowe
w obu kierunkach, których złożenia są identycznościami na
tensorach elementarnych, a więc na całych przestrzeniach.
\end{proof}

\begin{lemat}[Włożenia algebraicznych produktów tensorowych]
\label{lem:alg-inclusion24}
Niech $E,F$ będą przestrzeniami liniowymi nad $\K$,
a $E_0\subseteq E$, $F_0\subseteq F$ ich podprzestrzeniami.
Odwzorowanie kanoniczne $E_0\otimes F_0\to E\otimes F$
jest różnowartościowe. Jeśli $(e_i)_{i=1}^n$ jest bazą
skończenie wymiarowej przestrzeni $E_0$, to każdy element
$E_0\otimes F_0$ ma dokładnie jedno przedstawienie
$\sum_{i=1}^n e_i\otimes y_i$, gdzie $y_i\in F_0$.
\end{lemat}
\begin{proof}
Istnienie przedstawienia wynika z rozpinania przez tensory
elementarne i rozwinięcia pierwszego czynnika w bazie.
Dla funkcjonałów współrzędnych $e_i'$ odwzorowania
$e_i'\otimes I_{F_0}:E_0\otimes F_0\to F_0$,
określone dzięki własności uniwersalnej, odtwarzają
współczynniki $y_i$, co daje jednoznaczność.
W ogólnym przypadku wybieramy algebraiczne dopełnienia
podprzestrzeni $E_0,F_0$ (przez uzupełnienie ich baz).
Odpowiadające im rzuty $P:E\to E_0$ i $Q:F\to F_0$
spełniają $P|_{E_0}=I_{E_0}$, $Q|_{F_0}=I_{F_0}$.
Odwzorowanie $P\otimes Q$ jest lewą odwrotnością
włożenia tensorowego: złożenie jest identycznością na
tensorach elementarnych, więc na całej przestrzeni.
Nie twierdzimy, że rzuty te są ciągłe dla danych topologii.
\end{proof}

\begin{przyklad}[Tensory elementarne a funkcje]
\label{ex:separable-functions24}
Dla zbiorów $X,Y$ oznaczmy przez $\mathcal F(X)$
i $\mathcal F(Y)$ przestrzenie wszystkich funkcji o
wartościach w $\K$. Odwzorowanie
$\mathcal F(X)\otimes\mathcal F(Y)\to\mathcal F(X\times Y)$
określone przez $f\otimes g\mapsto ((x,y)\mapsto f(x)g(y))$
jest liniowe dzięki własności uniwersalnej i różnowartościowe.
Istotnie, zapisujemy tensor w postaci $\sum_{i=1}^n f_i\otimes g_i$
z liniowo niezależnymi $f_i$ (wybierając bazę rozpiętości
pierwszych czynników). Jeśli jego obraz znika, to dla
każdego $y\in Y$ mamy $\sum_i g_i(y)f_i=0$, zatem
wszystkie $g_i(y)=0$. To daje zerowy tensor.
Jeśli $X=Y=\mathbb N$, nie każda funkcja dwóch zmiennych
jest skończoną sumą funkcji rozdzielonych. Na przykład
funkcja $h(i,j)=\delta_{ij}$ nie ma takiego przedstawienia:
dla każdego $N$ jej ograniczenie do $\{1,\ldots,N\}^2$
ma macierz jednostkową rzędu $N$, podczas gdy suma $k$
funkcji rozdzielonych ma macierz rzędu co najwyżej $k$.
W przypadku jednego zbioru skończonego każda funkcja
jest natomiast skończoną sumą takich funkcji. Odpowiednie
topologie i uzupełnienia pozwalają rozważać granice
skończonych sum.
\end{przyklad}

\subsection{Topologia rzutowa}

\begin{definicja}[Seminorma rzutowa]
\label{def:pi-seminorm24}
Niech $E,F$ będą przestrzeniami lokalnie wypukłymi, a $p,q$
ciągłymi seminormami odpowiednio na $E$ i $F$. Dla
$u\in E\otimes F$ definiujemy
\[
 \pi_{p,q}(u)=\inf\Bigl\{\sum_{j=1}^N p(x_j)q(y_j):
 u=\sum_{j=1}^Nx_j\otimes y_j\Bigr\}.
\]
Topologię lokalnie wypukłą generowaną przez wszystkie
$\pi_{p,q}$ nazywamy \emph{topologią rzutową};
odpowiednią przestrzeń oznaczamy $E\otimes_\pi F$,
a jej uzupełnienie $E\widehat\otimes_\pi F$.
\end{definicja}

\begin{lemat}[Własności seminorm rzutowych]
\label{lem:pi-properties24}
Każda funkcja $\pi_{p,q}$ jest seminormą. Dla $x\in E$,
$y\in F$ zachodzi $\pi_{p,q}(x\otimes y)=p(x)q(y)$.
Odwzorowanie kanoniczne $E\times F\to E\otimes_\pi F$
jest ciągłe i dwuliniowe. Jeśli $E,F$ są Hausdorffa,
to również $E\otimes_\pi F$ jest Hausdorffa.
\end{lemat}
\begin{proof}
Jednorodność i nierówność trójkąta wynikają ze skalowania
oraz łączenia przedstawień tensorów. Dla tensora elementarnego
jednoskładnikowe przedstawienie daje nierówność $\le$.
Gdy $p(x)q(y)>0$, twierdzenie Hahna--Banacha daje
$\ell\in E'$ i $m\in F'$ z $|\ell(z)|\le p(z)$,
$|m(w)|\le q(w)$ oraz $\ell(x)=p(x)$, $m(y)=q(y)$
(po odpowiednim wyborze fazy w przypadku zespolonym).
Dla każdego przedstawienia $x\otimes y=\sum_jx_j\otimes y_j$
mamy $p(x)q(y)=|\sum_j\ell(x_j)m(y_j)|
\le\sum_jp(x_j)q(y_j)$. Gdy któryś czynnik zanika,
równość jest natychmiastowa. Ponieważ seminormy
$\pi_{p,q}(x\otimes y)$ są iloczynami $p(x)q(y)$,
ciągłość odwzorowania dwuliniowego w dowolnym punkcie
$(x_0,y_0)$ wynika z oszacowania
\[
\begin{aligned}
\pi_{p,q}(x\otimes y-x_0\otimes y_0)
&\le p(x-x_0)q(y-y_0)\\
&\quad+p(x-x_0)q(y_0)+p(x_0)q(y-y_0).
\end{aligned}
\] Dla $u\ne0$ istnieje para funkcjonałów
liniowych $\ell\in E'$, $m\in F'$ z
$(\ell\otimes m)(u)\ne0$: ograniczamy się do skończenie
wymiarowych przestrzeni rozpiętych przez współczynniki $u$
i przedłużamy funkcjonały współrzędnych twierdzeniem
Hahna--Banacha. Wybieramy $p=|\ell|$, $q=|m|$;
wtedy $\pi_{p,q}(u)\ge|(\ell\otimes m)(u)|>0$.
\end{proof}

\begin{twierdzenie}[Własność uniwersalna topologii rzutowej]
\label{thm:pi-universal24}
Niech $E,F,G$ będą przestrzeniami lokalnie wypukłymi
Hausdorffa. Przyporządkowanie $b\mapsto\widetilde b$ jest bijekcją
między ciągłymi odwzorowaniami dwuliniowymi
$b:E\times F\to G$ i ciągłymi odwzorowaniami liniowymi
$\widetilde b:E\otimes_\pi F\to G$. Jeśli $G$ jest
zupełna, każde takie $\widetilde b$ przedłuża się jednoznacznie
do $E\widehat\otimes_\pi F$.
\end{twierdzenie}
\begin{proof}
Niech $s$ będzie ciągłą seminormą na $G$ i niech $b$
będzie ciągłe. Istnieją zrównoważone otoczenia zera
$U\subset E$, $V\subset F$ takie, że $s(b(x,y))\le1$
na $U\times V$. Po zastąpieniu ich mniejszymi otwartymi
otoczeniami absolutnie wypukłymi i użyciu funkcjonałów
Minkowskiego otrzymujemy ciągłe seminormy $p,q$.
Dla $a>p(x)$ oraz $d>q(y)$, gdzie $a,d>0$,
mamy $x/a\in U$ i $y/d\in V$.
Dwuliniowość daje $s(b(x,y))\le ad$.
Przechodząc do infimum po $a,d$, otrzymujemy
$s(b(x,y))\le p(x)q(y)$, także gdy któraś
z seminorm znika. Stąd dla $u=\sum_jx_j\otimes y_j$ mamy
$s(\widetilde b(u))\le\sum_jp(x_j)q(y_j)$,
a po wzięciu infimum $s(\widetilde b(u))\le\pi_{p,q}(u)$.
To jest ciągłość $\widetilde b$. Odwrotnie, złożenie
ciągłego $\widetilde b$ z ciągłym odwzorowaniem
$(x,y)\mapsto x\otimes y$ jest ciągłe.
Jednoznaczność wynika z rozpinania produktu tensorowego
przez tensory elementarne. Przedłużenie do uzupełnienia
wynika z ciągłości jednostajnej odwzorowania liniowego
oraz zupełności $G$.
\end{proof}

\begin{wniosek}[Forma dwuliniowa i przestrzeń sprzężona]
\label{cor:pi-dual24}
Dla lokalnie wypukłych przestrzeni Hausdorffa $E,F$
przestrzeń liniowa $(E\otimes_\pi F)'$ jest w naturalnej
bijekcji z przestrzenią wszystkich łącznie ciągłych form
dwuliniowych $b:E\times F\to\K$.
\end{wniosek}
\begin{proof}
Stosujemy Twierdzenie~\ref{thm:pi-universal24} dla
$G=\K$. Przyporządkowanie jest liniowe, gdyż
sumowanie i mnożenie przez skalar zachowują
równości na tensorach elementarnych, które
jednoznacznie wyznaczają przedłużenie liniowe.
\end{proof}

\subsection{Topologia injektywna}

\begin{definicja}[Seminorma injektywna]
\label{def:eps-seminorm24}
Niech $E,F$ będą przestrzeniami lokalnie wypukłymi nad $\K$,
a $p,q$ ciągłymi seminormami odpowiednio na $E,F$.
Dla $u\in E\otimes F$ definiujemy
\[
 \varepsilon_{p,q}(u)=
 \sup\Bigl\{\big|\sum_{j=1}^N\ell(x_j)m(y_j)\big|:
 \substack{\ell\in E',\ m\in F',\\
             |\ell(x)|\le p(x)\ (x\in E),\\
             |m(y)|\le q(y)\ (y\in F)}\Bigr\},
 \quad u=\sum_{j=1}^N x_j\otimes y_j.
\]
Topologię wyznaczoną przez te seminormy nazywamy
\emph{topologią injektywną}; piszemy
$E\otimes_\varepsilon F$ i
$E\widehat\otimes_\varepsilon F$ dla jej uzupełnienia.
Występujące w definicji zbiory funkcjonałów są zbiorami
polarnymi domkniętych kul jednostkowych seminorm $p$ i $q$.
Dla przykładu warunek $|\ell(x)|\le1$ przy $p(x)\le1$
pociąga $|\ell(x)|\le p(x)$ przez skalowanie, także gdy
$p(x)=0$ (wtedy wszystkie wielokrotności $x$ leżą w kuli).
Dla ustalonej pary $(\ell,m)$ wartość sumy jest niezależna
od przedstawienia $u$: jest to wartość funkcjonału
$\ell\otimes m$, zadanego przez własność uniwersalną
produktu algebraicznego. Supremum jest skończone, ponieważ
dla dowolnego przedstawienia ogranicza je
$\sum_j p(x_j)q(y_j)$.
\end{definicja}

\begin{twierdzenie}[Porównanie topologii]
\label{thm:eps-pi24}
Dla dowolnych lokalnie wypukłych przestrzeni Hausdorffa
$E,F$ mamy $\varepsilon_{p,q}\le\pi_{p,q}$ dla każdej
pary ciągłych seminorm $p,q$. Obie są seminormami i
$\varepsilon_{p,q}(x\otimes y)=p(x)q(y)$.
Topologia injektywna jest Hausdorffa i nie jest silniejsza
od rzutowej. Kanoniczne odwzorowanie
$E\otimes_\pi F\to E\otimes_\varepsilon F$ jest ciągłe.
\end{twierdzenie}
\begin{proof}
Dla każdego przedstawienia $u=\sum_jx_j\otimes y_j$
i każdej dopuszczalnej pary $(\ell,m)$ zachodzi
$|\sum_j\ell(x_j)m(y_j)|\le\sum_jp(x_j)q(y_j)$.
Bierzemy supremum, a następnie infimum. Własności
seminormy są bezpośrednią konsekwencją definicji,
a równość dla tensora elementarnego wynika z osobnego
zastosowania twierdzenia Hahna--Banacha w obu czynnikach,
tak jak w Lemacie~\ref{lem:pi-properties24}; jeśli
$p(x)q(y)=0$, obie seminormy znikają.
Dla $u\ne0$ wybieramy skończenie wymiarowe podprzestrzenie
$E_0,F_0$ zawierające czynniki pewnego przedstawienia $u$.
Na mocy Lematu~\ref{lem:alg-inclusion24} tensor w
$E_0\otimes F_0$ jest niezerowy. W bazach obu
podprzestrzeni ma więc niezerowy współczynnik,
który wykrywają odpowiednie funkcjonały współrzędnych.
Skończenie wymiarowe podprzestrzenie przestrzeni
lokalnie wypukłych Hausdorffa mają standardową topologię,
więc funkcjonały te są ciągłe i przedłużają się przez
twierdzenie Hahna--Banacha do $\ell\in E'$ i $m\in F'$.
Dla $p=|\ell|$ i $q=|m|$ mamy
$\varepsilon_{p,q}(u)\ge|(\ell\otimes m)(u)|>0$.
To dowodzi oddzielania punktów.
Porównanie seminorm daje ciągłość identyczności z
topologii rzutowej do injektywnej.
\end{proof}

\begin{stwierdzenie}[Funktorialność i symetria iloczynów tensorowych]
\label{prop:dt-tensor-functoriality}
Niech $E,F,E_1,F_1$ będą przestrzeniami lokalnie
wypukłymi Hausdorffa i niech $A:E\to E_1$ oraz
$B:F\to F_1$ będą ciągłymi odwzorowaniami liniowymi.
Odwzorowanie algebraiczne $A\otimes B$, określone przez
$(A\otimes B)(x\otimes y)=Ax\otimes By$, jest ciągłe
zarówno dla topologii rzutowych, jak i injektywnych.
Dla ciągłych seminorm $p_1,q_1$ odpowiednio na $E_1,F_1$
zachodzi, dla $\tau\in\{\pi,\varepsilon\}$,
\[
 \tau_{p_1,q_1}((A\otimes B)u)
 \le\tau_{p_1\circ A,q_1\circ B}(u).
\]
Odwzorowania te przedłużają się jednoznacznie na
odpowiednie uzupełnienia. Ponadto zamiana czynników
$x\otimes y\mapsto y\otimes x$ jest izomorfizmem
liniowo-topologicznym dla obu topologii, również
po uzupełnieniu.
\end{stwierdzenie}
\begin{proof}
Istnienie i jednoznaczność $A\otimes B$ wynikają
z algebraicznej własności uniwersalnej zastosowanej
do odwzorowania dwuliniowego $(x,y)\mapsto Ax\otimes By$.
Oznaczmy $p=p_1\circ A$, $q=q_1\circ B$; są to
ciągłe seminormy. Każde przedstawienie
$u=\sum_j x_j\otimes y_j$ daje przedstawienie obrazu
o sumie iloczynów seminorm równej $\sum_jp(x_j)q(y_j)$.
Wzięcie infimum daje oszacowanie dla $\pi$.
Jeśli $\ell_1,m_1$ są funkcjonałami zdominowanymi
odpowiednio przez $p_1,q_1$, to $\ell_1\circ A$
i $m_1\circ B$ są zdominowane odpowiednio przez $p,q$.
Zatem
\[
 |(\ell_1\otimes m_1)((A\otimes B)u)|
 =|((\ell_1\circ A)\otimes(m_1\circ B))(u)|
 \le\varepsilon_{p,q}(u).
\]
Supremum po dopuszczalnych funkcjonałach daje
oszacowanie dla $\varepsilon$. Oba oszacowania
dowodzą ciągłości. Złożenia oraz identyczności
są zachowywane, ponieważ odpowiadające odwzorowania
zgadzają się na tensorach elementarnych, a te
rozpinają iloczyn algebraiczny.

Ciągłe odwzorowanie liniowe jest jednostajnie ciągłe.
Po złożeniu z kanonicznym włożeniem przestrzeni
docelowej do jej uzupełnienia przedłużamy je
jednoznacznie przez gęstość; uzupełnienie docelowe
jest zupełne i Hausdorffa.
Dla zamiany czynników definicje obu seminorm dają
$\tau_{q,p}(\sum_jy_j\otimes x_j)
=\tau_{p,q}(\sum_jx_j\otimes y_j)$, ponieważ zamiana
jest odwracalna i zamienia odpowiednio przedstawienia
lub pary funkcjonałów. To dowodzi ciągłości w obu
kierunkach. Przedłużenia zamiany i jej odwrotności
są wzajemnie odwrotne przez zgodność na gęstych
iloczynach algebraicznych.
\end{proof}

\begin{twierdzenie}[Dziedziczność normy injektywnej]
\label{thm:eps-injective24}
Niech $X_0\subseteq X$ i $Y_0\subseteq Y$ będą izometrycznie
włożonymi podprzestrzeniami przestrzeni unormowanych.
Dla każdego $u\in X_0\otimes Y_0$ jego norma injektywna
liczona w $X_0\otimes_\varepsilon Y_0$ jest równa normie
liczonej po utożsamieniu z elementem $X\otimes_\varepsilon Y$.
\end{twierdzenie}
\begin{proof}
Utożsamienie tensorów jest uprawnione przez
Lemat~\ref{lem:alg-inclusion24}.
Ograniczenie funkcjonałów z kul jednostkowych $X'$ i $Y'$
daje nierówność $\varepsilon_{X,Y}(u)\le
\varepsilon_{X_0,Y_0}(u)$, gdzie indeksy oznaczają
przestrzenie, w których liczymy normę. Odwrotnie, twierdzenie
Hahna--Banacha przedłuża dowolne $\ell\in X_0'$
i $m\in Y_0'$ z zachowaniem norm. Otrzymujemy więc
identyczny zbiór dopuszczalnych wartości na tensorze $u$.
\end{proof}

\begin{przyklad}[Przypadek skończenie wymiarowy]
\label{ex:finite-dim24}
Dla $E=\K^n$ z dowolną normą oraz przestrzeni Banacha $F$
obydwa uzupełnione produkty tensorowe są liniowo-topologicznie
izomorficzne z $F^n$. Rzeczywiście, przy wyborze bazy
$(e_i)$ i bazy sprzężonej $(e_i')$ każdy tensor ma dokładnie
jedno przedstawienie $u=\sum_{i=1}^ne_i\otimes y_i$.
Z twierdzenia Hahna--Banacha w $F$ mamy
\[
 \|y_i\|=\sup_{\|m\|\le1}|(e_i'\otimes m)(u)|
 \le\|e_i'\|\varepsilon(u).
\]
Z drugiej strony
\[
 \varepsilon(u)\le\pi(u)\le\sum_{i=1}^n\|e_i\|\|y_i\|
 \le\Big(\sum_{i=1}^n\|e_i\|\Big)\max_i\|y_i\|.
\]
Obie normy są zatem równoważne normie maksimum na $F^n$.
Ponieważ $F$ jest zupełna, oba produkty tensorowe są już
zupełne, a ich uzupełnienia nie wprowadzają nowych elementów.
\end{przyklad}

\begin{twierdzenie}[Rozbieżność norm w nieskończonym wymiarze]
\label{thm:hilbert-norms24}
Niech $H$ będzie nieskończenie wymiarową przestrzenią Hilberta
nad $\K$ i $(e_n)$ ortonormalnym ciągiem w $H$. Dla
$u_N=\sum_{n=1}^N e_n\otimes e_n\in H\otimes H$ zachodzi
\[ \varepsilon(u_N)=1,\qquad \pi(u_N)=N. \]
W szczególności na $H\otimes H$ topologie rzutowa i
injektywna są różne.
\end{twierdzenie}
\begin{proof}
Dla $\ell,m$ z kul jednostkowych $H'$ nierówność
Cauchy'ego--Schwarza i Bessela dają
$|\sum_{n=1}^N\ell(e_n)m(e_n)|\le1$. Wybór funkcjonałów
współrzędnych dla $e_1$ pokazuje równość. Z definicji
$\pi(u_N)\le N$. Aby uzyskać nierówność przeciwną,
rozważmy formę dwuliniową
$b(x,y)=\sum_{n=1}^N\langle x,e_n\rangle\langle y,e_n\rangle$.
Przy konwencji liniowości iloczynu skalarnego w pierwszym
argumencie jest ona dwuliniowa. Na mocy Bessela i
nierówności Cauchy'ego--Schwarza ma normę nie większą niż $1$,
a $\widetilde b(u_N)=N$. Dla każdej reprezentacji $u_N$ suma
$\sum_j\|x_j\|\|y_j\|$ jest więc nie mniejsza niż $N$.
Biorąc infimum, dostajemy $\pi(u_N)=N$.
Dla $v_N=u_N/N$ mamy $\varepsilon(v_N)=1/N\to0$,
lecz $\pi(v_N)=1$, co dowodzi różności topologii,
a nie tylko różności wartości norm.
\end{proof}

\subsection{Nuklearność i równość topologii tensorowych}

\begin{lemat}[Seminormy tensorowe i uzupełnienia ilorazów]
\label{lem:tensor-seminorm-completions24}
Niech $E,F$ będą przestrzeniami lokalnie wypukłymi
Hausdorffa, a $p,q$ ich ciągłymi seminormami.
Niech $E_p,F_q$ będą uzupełnieniami unormowanych
ilorazów (Definicja~\ref{def:nuclear-lcs23}).
Dla $u\in E\otimes F$ oznaczmy przez $u_{p,q}$
jego obraz w $E_p\otimes F_q$ pod działaniem
kanonicznych odwzorowań w obu czynnikach. Wówczas
\[
 \pi_{p,q}(u)=\|u_{p,q}\|_\pi,\qquad
 \varepsilon_{p,q}(u)=\|u_{p,q}\|_\varepsilon.
\]
Normy po prawej można liczyć również w odpowiednich
uzupełnionych produktach tensorowych.
\end{lemat}
\begin{proof}
Każde przedstawienie $u=\sum_jx_j\otimes y_j$
daje przedstawienie $u_{p,q}=\sum_j[x_j]_p\otimes[y_j]_q$
o tej samej sumie iloczynów seminorm. Stąd
$\|u_{p,q}\|_\pi\le\pi_{p,q}(u)$.
Jeżeli $\pi_{p,q}(u)=0$, otrzymujemy od razu równość.
W przeciwnym przypadku twierdzenie Hahna--Banacha
dla seminormy $\pi_{p,q}$ daje liniowy funkcjonał
$\ell$ na $E\otimes F$ taki, że
$|\ell(v)|\le\pi_{p,q}(v)$ dla każdego $v$
oraz $\ell(u)=\pi_{p,q}(u)$.
Istotnie, najpierw określamy funkcjonał na
jednowymiarowej podprzestrzeni rozpiętej przez $u$.

Forma $b(x,y)=\ell(x\otimes y)$ spełnia
$|b(x,y)|\le p(x)q(y)$ na mocy
Lematu~\ref{lem:pi-properties24}. Znika, gdy
$x\in\ker p$ lub $y\in\ker q$, więc przechodzi
do formy na ilorazach z normą co najwyżej $1$.
Przedłuża się jednoznacznie do formy dwuliniowej
$\bar b:E_p\times F_q\to\K$ o tym samym oszacowaniu.
Dokładniej, dla ciągów Cauchy'ego $(x_n),(y_n)$
w unormowanych ilorazach różnica wartości formy
jest ograniczona przez
\[
 \|x_n-x_m\|\,\|y_n\|+\|x_m\|\,\|y_n-y_m\|.
\]
Ograniczoność ciągów i to oszacowanie dają istnienie
granicy i jej niezależność od reprezentantów.
Dwuliniowość oraz oszacowanie normowe przechodzą
do granicy. Dla każdego przedstawienia
$u_{p,q}=\sum_j a_j\otimes b_j$ mamy zatem
\[
 \pi_{p,q}(u)=|\widetilde{\bar b}(u_{p,q})|
 \le\sum_j\|a_j\|\,\|b_j\|.
\]
Biorąc infimum, otrzymujemy przeciwną nierówność.
Argument nie zakłada dziedziczności normy rzutowej
dla dowolnych włożeń izometrycznych.

Dla normy injektywnej każdy funkcjonał $\ell$
z $|\ell|\le p$ znika na $\ker p$ i jednoznacznie
przedłuża się do funkcjonału na $E_p$ o normie
co najwyżej $1$. Odwrotnie, każdy taki funkcjonał
na $E_p$ daje przez złożenie z kanonicznym
odwzorowaniem funkcjonał na $E$ zdominowany przez $p$.
To samo dotyczy $q$ i $F_q$. Zbiory dopuszczalnych
wartości w obu definicjach supremum są więc identyczne.
Obie normy tensorowe są normami na iloczynie
przestrzeni unormowanych, a kanoniczne włożenia do
ich uzupełnień są izometriami.
\end{proof}

\begin{twierdzenie}[Nuklearność usuwa różnicę między $\pi$ i $\varepsilon$]
\label{thm:nuclear-pi-eps24}
Niech $E$ będzie nuklearną przestrzenią lokalnie wypukłą
Hausdorffa w sensie Definicji~\ref{def:nuclear-lcs23},
a $F$ dowolną przestrzenią lokalnie wypukłą Hausdorffa.
Wtedy topologie rzutowa i injektywna na
$E\otimes F$ są równe. Ich uzupełnienia są więc
kanonicznie izomorficzne:
$E\widehat\otimes_\pi F\simeq E\widehat\otimes_\varepsilon F$.
\end{twierdzenie}
\begin{proof}
Ustalmy seminormę rzutową $\pi_{p,r}$, gdzie $p$ jest
ciągłą seminormą na $E$, a $r$ na $F$. Z nuklearności
istnieje ciągła seminorma $q\ge p$, taka że kanoniczny
operator $j_{q,p}:E_q\to E_p$ między przestrzeniami
Banacha jest nuklearny. Wybierzmy jego przedstawienie
$j_{q,p}(z)=\sum_{k\ge1}\ell_k(z)a_k$ z
$\sum_k\|\ell_k\|\,\|a_k\|=C<\infty$.
Jeżeli $u=\sum_{j=1}^N x_j\otimes y_j$, to dla każdego
$j$ mamy $[x_j]_p=j_{q,p}([x_j]_q)$.
Wstawiając sumy częściowe przedstawienia nuklearnego
do skończonej sumy tensorów, otrzymujemy w
$E_p\widehat\otimes_\pi F_r$ granicę szeregu
\[
 \sum_{k\ge1}a_k\otimes
 \Bigl(\sum_{j=1}^N\ell_k([x_j]_q)[y_j]_r\Bigr).
\]
Istotnie, dla każdego $j$ sumy częściowe zbiegają
w $E_p$ do $[x_j]_p$, a norma rzutowa tensora
elementarnego jest iloczynem norm.
Z twierdzenia Hahna--Banacha w $F_r$ mamy
\[
 \left\|\sum_j\ell_k([x_j]_q)[y_j]_r\right\|_{F_r}
 \le\|\ell_k\|\,\varepsilon_{q,r}(u),
\]
ponieważ dla $\|m\|_{F_r'}\le1$ funkcjonał
$y\mapsto m([y]_r)$ jest zdominowany przez $r$,
zaś $x\mapsto\ell_k([x]_q)$ jest zdominowany przez
$\|\ell_k\|q$. Dla $\ell_k=0$ nierówność jest oczywista,
a w pozostałych przypadkach dzielimy przez $\|\ell_k\|$.
Z nierówności trójkąta dla normy rzutowej i zbieżności
szeregu otrzymujemy, stosując Lemat~
\ref{lem:tensor-seminorm-completions24},
\[
 \pi_{p,r}(u)=\|u_{p,r}\|_\pi
 \le\sum_k\|a_k\|\,\|\ell_k\|\,\varepsilon_{q,r}(u)
 =C\varepsilon_{q,r}(u).
\]
Otrzymaliśmy, że każda seminorma rzutowa jest
kontrolowana przez seminormę injektywną (z potencjalnie
większą seminormą w pierwszym czynniku). Odwrotna
ciągłość wynika z Twierdzenia~\ref{thm:eps-pi24}.
\end{proof}

\begin{uwaga}[Granica stosowalności]
\label{rem:nuclear-characterization24}
Twierdzenie~\ref{thm:nuclear-pi-eps24} wyjaśnia,
dlaczego dla przestrzeni nuklearnych można niekiedy
pisać $E\widehat\otimes F$ bez wskazywania, czy
użyto topologii $\pi$, czy $\varepsilon$.
Nie oznacza to równości wszystkich możliwych topologii
tensorowych: topologia induktywna $i$, która pojawi się
poniżej, może być silniejsza nawet dla przestrzeni nuklearnych.
Pełne twierdzenie odwrotne Grothendiecka (charakterystyka
nuklearności przez równość $\pi$ i $\varepsilon$ dla wszystkich
czynników) wymaga dodatkowej teorii i nie jest używane w dowodach.
\end{uwaga}

\begin{wniosek}[Przestrzenie Banacha nuklearne tylko w skończonym wymiarze]
\label{cor:banach-finite-nuclear24}
Przestrzeń Banacha $X$ jest nuklearna jako przestrzeń
lokalnie wypukła wtedy i tylko wtedy, gdy ma skończony wymiar.
\end{wniosek}
\begin{proof}
Obie implikacje wykazano w Uwadze~
\ref{rem:nuclear-different23}. Dla pełności przypomnijmy,
że w przypadku nieskończenie wymiarowym nuklearność
wymuszałaby zwartość identyczności na $X$ przez
równoważność normy z dominującą ciągłą seminormą.
W przypadku skończenie wymiarowym dla każdej ciągłej
seminormy $p$ można wybrać $q=p$: identyczność na
skończenie wymiarowym uzupełnieniu $X_p$ jest operatorem
nuklearnym (również gdy $X_p=\{0\}$).
\end{proof}


\begin{stwierdzenie}[Nuklearność produktów]
\label{prop:nuclear-products24}
Niech $(E_\alpha)_{\alpha\in A}$ będzie dowolną rodziną
nuklearnych przestrzeni lokalnie wypukłych Hausdorffa.
Wówczas przestrzeń $E=\prod_{\alpha\in A}E_\alpha$
z topologią produktu jest nuklearna.
\end{stwierdzenie}
\begin{proof}
Niech $p$ będzie ciągłą seminormą na $E$.
Z definicji topologii produktu istnieje skończony zbiór
$J\subseteq A$ taki, że $p(x)=0$, gdy $x_\alpha=0$
dla każdego $\alpha\in J$. Istotnie, w otoczeniu zera
$\{x:p(x)<1\}$ mieści się otoczenie cylindryczne
wyznaczone przez skończenie wiele współrzędnych.
Wektor zerujący się na nich można pomnożyć przez dowolny
skalar bez opuszczania cylindra, więc jego seminorma
musi być zerowa. Jeżeli $J=\varnothing$, to $p=0$ i wystarcza $q=0$.
W przeciwnym przypadku niech $\iota_\alpha:E_\alpha\to E$
będzie włożeniem na współrzędną $\alpha$, z zerami
na pozostałych współrzędnych. Jest ono ciągłe.
Dla $p_\alpha=p\circ\iota_\alpha$ mamy
\[
 p(x)=p\Bigl(\sum_{\alpha\in J}\iota_\alpha(x_\alpha)\Bigr)
 \le\sum_{\alpha\in J}p_\alpha(x_\alpha)=P(x),
\]
ponieważ różnica wektorów należy do $\ker p$.
W szczególności wszystkie $p_\alpha$ są ciągłe.
Dla każdego $\alpha\in J$ wybieramy $q_\alpha\ge p_\alpha$
takie, że $j_{q_\alpha,p_\alpha}$ jest nuklearny,
i przyjmujemy $q(x)=\sum_{\alpha\in J}q_\alpha(x_\alpha)$.
Dla $t=q$ oraz $t=P$ jądro seminormy składa się
z wektorów, których współrzędne z $J$ leżą w jądrach
odpowiednich seminorm składowych. Uzupełnienie $E_t$
identyfikujemy izometrycznie ze skończoną sumą
$\ell^1$ odpowiednich uzupełnień współrzędnych:
odwzorowanie ilorazu na tę sumę zachowuje normę
i ma gęsty obraz przez gęstość każdego składnika.
Przejście $E_q\to E_P$ jest więc sumą
\[
 j_{q,P}=\sum_{\alpha\in J}
 I_\alpha j_{q_\alpha,p_\alpha}R_\alpha,
\]
gdzie $R_\alpha$ jest rzutem współrzędnym, a $I_\alpha$
włożeniem do sumy $\ell^1$; ich normy nie przekraczają $1$.
Zgodność sprawdzamy na gęstym obrazie $E$.
Każdy składnik jest nuklearny na mocy
Stwierdzenia~\ref{prop:nuclear-ideal23}, zatem
skończona suma też jest nuklearna. Przez $p\le P$
operator $j_{P,p}$ jest ograniczony, a złożenie
$j_{q,p}=j_{P,p}j_{q,P}$ jest nuklearne.
Ponadto $q\ge p$, co kończy dowód.
\end{proof}

\begin{uwaga}[Trwałość względem podprzestrzeni i granic induktywnych]
\label{rem:nuclear-permanence24}
Trwałość nuklearności względem podprzestrzeni i ilorazów
została udowodniona w
Twierdzeniu~\ref{thm:nuclear-subquot23}, a względem
przeliczalnych granic induktywnych --- w
Twierdzeniu~\ref{thm:nuclear-indlim23}.
Dlatego nuklearność $\mathcal D_K(\Omega)$ i
$\mathcal D(\Omega)$ jest już wynikiem
Twierdzenia~\ref{thm:nuclear-D23}, a nie osobnym
nieudowodnionym założeniem. Wyniki te nie uprawniają
do utożsamienia topologii $i$ i $\pi$ na
$\mathcal D(\Omega)\otimes\mathcal D(\Lambda)$:
porównują one odmienne własności uniwersalne.
\end{uwaga}

\subsection{Przykłady produktów tensorowych przestrzeni funkcyjnych}

\begin{twierdzenie}[Produkt tensorowy $\ell^1$]
\label{thm:l1-tensor24}
Niech $F$ będzie przestrzenią Banacha. Wówczas istnieje
kanoniczny izometryczny izomorfizm
$\ell^1(\N)\widehat\otimes_\pi F\simeq\ell^1(\N;F)$,
gdzie prawa strona składa się z ciągów $(y_n)$ spełniających
$\sum_n\|y_n\|_F<\infty$, z normą tej sumy.
\end{twierdzenie}
\begin{proof}
Najpierw definiujemy liniowe odwzorowanie
$J:\ell^1\otimes F\to\ell^1(\N;F)$ wzorem
\[
 J\Bigl(\sum_{j=1}^M a^{(j)}\otimes y_j\Bigr)
 =\Bigl(\sum_{j=1}^M a_n^{(j)}y_j\Bigr)_{n\ge1}.
\]
Jest ono dobrze określone dzięki własności uniwersalnej
produktu algebraicznego: odwzorowanie
$(a,y)\mapsto(a_ny)_{n\ge1}$ jest dwuliniowe i jego
wartości należą do $\ell^1(\N;F)$, ponieważ
$\sum_n\|a_n y\|=\|a\|_1\|y\|$.
Przestrzeń $\ell^1(\N;F)$ jest zupełna: dla ciągu
Cauchy'ego granice współrzędnych istnieją w $F$,
a oszacowania norm różnic przechodzą do granicy
najpierw dla skończonych sum, następnie przez
supremum po długościach tych sum. Daje to
sumowalność norm współrzędnych granicy i zbieżność w $\ell^1$.
Dla dowolnego przedstawienia
$u=\sum_j a^{(j)}\otimes y_j$ zachodzi
$\|Ju\|_{\ell^1(F)}\le\sum_j\|a^{(j)}\|_1\|y_j\|$,
a więc $\|Ju\|_{\ell^1(F)}\le\pi(u)$.

Jeżeli $u=\sum_{n=1}^N e_n\otimes y_n$, to nierówność
przeciwna wynika z twierdzenia Hahna--Banacha.
Wybieramy $m_n\in F'$ z $\|m_n\|\le1$ i
$m_n(y_n)=\|y_n\|$ (dla $y_n=0$ bierzemy $m_n=0$).
Forma $b(a,y)=\sum_{n=1}^Na_nm_n(y)$ ma normę
nie większą niż $1$ na $\ell^1\times F$ oraz
$\widetilde b(u)=\sum_n\|y_n\|$. Zatem
$\pi(u)=\sum_n\|y_n\|=\|Ju\|_{\ell^1(F)}$
dla tensorów o skończonym nośniku współrzędnych.

Pozostaje uzasadnić, że wyliczenie normy obejmuje
\emph{dowolny} tensor algebraiczny. Jeśli
$u=\sum_{j=1}^M a^{(j)}\otimes y_j$, niech $P_N$
obcina ciąg po $N$ współrzędnych i połóżmy
$u_N=\sum_j(P_Na^{(j)})\otimes y_j$. Wtedy
\[
 \pi(u-u_N)\le\sum_{j=1}^M
 \|a^{(j)}-P_Na^{(j)}\|_1\,\|y_j\|\longrightarrow0.
\]
Ponadto $\|Ju_N-Ju\|_{\ell^1(F)}\le\pi(u_N-u)\to0$. Równość
norm dla $u_N$, wraz z ciągłością norm obu przestrzeni,
daje $\pi(u)=\|Ju\|_{\ell^1(F)}$.
Odwzorowanie $J$ jest więc izometrią o obrazie
zawierającym wszystkie ciągi o skończonym nośniku.
Przedłuża się do izometrii z zupełnego
$\ell^1\widehat\otimes_\pi F$ na $\ell^1(\N;F)$:
jej obraz jest domknięty i zarazem gęsty.
\end{proof}

\begin{twierdzenie}[Jądra gładkie i szybko malejące]
\label{thm:smooth-schwartz-tensor24}
Dla niepustych zbiorów otwartych $\Omega\subseteq\R^{d_1}$,
$\Lambda\subseteq\R^{d_2}$ oraz $d_1,d_2\ge1$ mamy
naturalne izomorfizmy liniowo-topologiczne
\[
 \mathcal E(\Omega)\widehat\otimes_\pi\mathcal E(\Lambda)
 \simeq \mathcal E(\Omega\times\Lambda),\qquad
 \mathcal S(\R^{d_1})\widehat\otimes_\pi\mathcal S(\R^{d_2})
 \simeq \mathcal S(\R^{d_1+d_2}).
\]
W obu wzorach można zastąpić $\pi$ przez $\varepsilon$.
\end{twierdzenie}
\begin{proof}
W obu przypadkach odwzorowanie kanoniczne na tensorach
elementarnych przypisuje $f\otimes g$ funkcję
$(x,y)\mapsto f(x)g(y)$.

Dla przestrzeni Schwartza stosujemy izomorfizmy z
przestrzeniami szybko malejących ciągów, udowodnione
w Lemacie~\ref{lem:hermite-schwartz23}.
Izomorfizmy w czynnikach indukują izomorfizm
uzupełnionych iloczynów rzutowych na mocy
Stwierdzenia~\ref{prop:dt-tensor-functoriality},
zastosowanego do obu izomorfizmów oraz ich odwrotności.
Z twierdzenia Fubiniego współczynniki Hermite'a
funkcji $(x,y)\mapsto f(x)g(y)$ są iloczynami
współczynników obu funkcji (całki są bezwzględnie
zbieżne dzięki nierówności Cauchy'ego--Schwarza).
Wystarczy więc zbadać odwzorowanie
$J:s(I)\otimes s(J_0)\to s(I\times J_0)$,
gdzie $I=\N_0^{d_1}$, $J_0=\N_0^{d_2}$,
a tensor elementarny $(a_i)\otimes(b_j)$
przechodzi na macierz $(a_i b_j)_{i,j}$.
Oznaczmy $w(i,j)=1+|i|+|j|$ oraz
$P_m(c)=\sup_{i,j}w(i,j)^m|c_{ij}|$.
Ponieważ $w(i,j)\le(1+|i|)(1+|j|)$,
dla każdego przedstawienia tensora $u$ otrzymujemy
$P_m(Ju)\le\sum_k p_m(a^{(k)})p_m(b^{(k)})$.
Biorąc infimum, mamy
$P_m(Ju)\le\pi_{p_m,p_m}(u)$.
Seminormy o równych indeksach są współkońcowe
wśród par seminorm obu przestrzeni, ponieważ
$p_k$ rosną z $k$. Zupełność $s(I\times J_0)$
(Twierdzenie~\ref{thm:nuclear-s23}) daje ciągłe
przedłużenie $\widehat J$ do uzupełnienia.

Dla $c\in s(I\times J_0)$ definiujemy odwzorowanie
odwrotne $R(c)$ jako sumę szeregu
$\sum_{i,j}c_{ij}e_i\otimes e_j$ w uzupełnieniu.
Dla $M>2m+d_1+d_2$ mamy
\[
 \sum_{i,j}|c_{ij}|(1+|i|)^m(1+|j|)^m
 \le P_M(c)\sum_{i,j}w(i,j)^{2m-M}<\infty.
\]
Ostatnia suma jest skończona na mocy zliczania
wielowskaźników w Twierdzeniu~\ref{thm:nuclear-s23}
dla wymiaru $d_1+d_2$. Sumy po skończonych zbiorach
indeksów tworzą sieć Cauchy'ego w każdej seminormie
rzutowej; z zupełności istnieje granica.
Powyższe oszacowanie dowodzi ciągłości $R$.
Współczynniki $\widehat J R(c)$ są równe $c_{ij}$,
ponieważ każda ewaluacja współrzędnej jest ciągła.
Dla tensora elementarnego $a\otimes b$ sumy po
prostokątnych zbiorach $|i|\le N$, $|j|\le N$
są tensorami obcięć $a^{[N]}\otimes b^{[N]}$.
Obcięcia zbiegają do $a,b$ we wszystkich normach
$p_m$, gdyż ogony są ograniczone przez
$p_{m+1}(a)/(1+N)$ i $p_{m+1}(b)/(1+N)$.
Z ciągłości kanonicznego odwzorowania dwuliniowego
wynika $R J(a\otimes b)=a\otimes b$.
Liniowość daje równość na iloczynie algebraicznym,
a ciągłość i gęstość dają $R\widehat J=I$ na
uzupełnieniu. Dowodzi to obu kierunków izomorfizmu.

Dla $\mathcal E$ oznaczmy przez $J_{\mathcal E}$
odwzorowanie kanoniczne na iloczynie algebraicznym.
Dla zwartych $K\Subset\Omega$, $L\Subset\Lambda$
i $m\in\N_0$ mamy
\[
 p_{K\times L,m}(J_{\mathcal E}u)
 \le\pi_{p_{K,m},p_{L,m}}(u),
\]
ponieważ pochodne tensora elementarnego w obu zmiennych
rozdzielają się, a następnie bierzemy infimum
po przedstawieniach. Każdy zwarty podzbiór iloczynu
mieści się w iloczynie swoich zwartych projekcji.
Stąd $J_{\mathcal E}$ jest ciągłe i przedłuża się
do ciągłego $\widehat J_{\mathcal E}$, ponieważ
$\mathcal E(\Omega\times\Lambda)$ jest zupełna.

Wybierzmy przeliczalne lokalnie skończone pokrycia
otwartymi sześcianami $V_i\Subset\Omega$,
$W_j\Subset\Lambda$, oraz gładkie rozkłady jedności
$(\eta_i)$, $(\theta_j)$ podporządkowane tym pokryciom.
Istnienie takich pokryć wynika z wyczerpania
$K_n\subseteq\operatorname{int}K_{n+1}$
(Lemat~\ref{lem:exhaustion17}):
dla warstwy $K_n\setminus\operatorname{int}K_{n-1}$
wybieramy skończenie wiele sześcianów o domknięciach
zawartych w $\operatorname{int}K_{n+1}\setminus K_{n-2}$
(przyjmujemy $K_0=K_{-1}=\varnothing$).
Pokrycie jest lokalnie skończone: dla punktu
w $\operatorname{int}K_N$ to otoczenie nie przecina
sześcianów z warstw $n\ge N+2$, a pozostałych
sześcianów jest skończenie wiele.
Stosujemy Twierdzenie~\ref{thm:dt-smooth-partition}.
Otrzymaną partycję, początkowo indeksowaną przez
rodzinę funkcji przypisanych elementom pokrycia,
grupujemy według tych elementów. Dla ustalonego $V_i$
tylko skończenie wiele niezerowych składników może
być przypisanych do $V_i$, gdyż ich nośniki tworzą
rodzinę lokalnie skończoną, a $\overline V_i$ jest
zwartym podzbiorem $\Omega$. Ich suma $\eta_i$ ma
więc zwarty nośnik zawarty w $V_i$; gdy nie ma
składników, przyjmujemy $\eta_i=0$. Tak samo
konstruujemy $\theta_j$ na $\Lambda$.
Wybieramy odcięcia rzeczywiste
$\chi_i\in\mathcal D(V_i)$, $\psi_j\in\mathcal D(W_j)$,
równe $1$ w otoczeniach tych nośników
(Lemat~\ref{lem:cutoffs17}; przy pustym nośniku
można wybrać odcięcie zerowe).
Istotna jest lokalna skończoność całych rodzin
$(\supp\chi_i)$ i $(\supp\psi_j)$, wynikająca
z lokalnej skończoności pokryć.

Na domknięciach sześcianów stosujemy rzeczywiste
ortonormalne układy trygonometryczne $(e_{ik})_k$,
$(h_{jl})_l$ względem znormalizowanych miar Lebesgue'a.
W każdej współrzędnej bierzemy $1$ oraz
$\sqrt2\cos(\pi n(t-a)/R)$, $\sqrt2\sin(\pi n(t-a)/R)$
dla $n\ge1$, gdzie $a$ jest środkiem, a $2R$ długością
boku. W wielu wymiarach bierzemy iloczyny tych funkcji.
Ortogonalność i normalizacja w jednym wymiarze
wynikają z całkowania iloczynów funkcji sinus
i cosinus po pełnym okresie; w wielu wymiarach
stosujemy twierdzenie Fubiniego. Jest to baza
ortonormalna także nad $\C$, ponieważ jej zespolona
rozpiętość zawiera wykładnicze Fouriera, których
zupełność wykazano w dowodzie
Twierdzenia~\ref{thm:nuclear-cinfty23} (po przesunięciu
środka sześcianu). W przypadku rzeczywistym zupełność
wynika przez rozpatrzenie rzeczywistej funkcji
ortogonalnej do tego układu jako funkcji zespolonej.
Przez $|k|$ i $|l|$ oznaczamy normy euklidesowe
wektorów częstości; każda częstość ma ograniczoną,
zależną tylko od wymiaru, krotność.
Funkcje $e_{ik},h_{jl}$ są funkcjami własnymi
$1-\Delta$ z wartościami własnymi odpowiednio
$1+\pi^2|k|^2/R_i^2$, $1+\pi^2|l|^2/S_j^2$.

Dla $\Phi\in\mathcal E(\Omega\times\Lambda)$
lokalizujemy $\Phi_{ij}=\eta_i\theta_j\Phi$,
przedłużamy zerem, a następnie okresowo z
$\overline V_i\times\overline W_j$.
Gładkość periodyzacji wynika ze znikania w otoczeniu
brzegu. Niech $c_{kl}^{ij}$ będą współczynnikami
w iloczynowej bazie trygonometrycznej: całkami
$\Phi_{ij}e_{ik}h_{jl}$ względem znormalizowanej
miary na iloczynie sześcianów. Kładziemy
$f_{ik}=\chi_i e_{ik}$, $g_{jl}=\psi_j h_{jl}$,
przedłużone zerem poza odpowiednie sześciany.
Dla $a,b,r,s\in\N_0$ całkowanie przez części z
$(1-\Delta_x)^a(1-\Delta_y)^b$ i reguła Leibniza dają
\[
 \begin{aligned}
 |c_{kl}^{ij}|&\le C_{ij,a,b}(1+|k|)^{-2a}
 (1+|l|)^{-2b}
 p_{\overline V_i\times\overline W_j,2a+2b}(\Phi),\\
 p_{K,r}(f_{ik})&\le C_{i,r}(1+|k|)^r,\qquad
 p_{L,s}(g_{jl})\le C_{j,s}(1+|l|)^s.
 \end{aligned}
\]
Stałe są niezależne od $\Phi,k,l$; pochodne ustalonych
odcięć są ograniczone. Dla $2a>r+d_1$, $2b>s+d_2$
szereg iloczynów tych trzech oszacowań jest zbieżny
przez zliczanie punktów kratowych.
Zatem
$T_{ij}(\Phi)=\sum_{k,l}c_{kl}^{ij}f_{ik}\otimes g_{jl}$
istnieje w uzupełnieniu rzutowym i jest liniowe
oraz ciągłe w $\Phi$.
Jednostajna zbieżność szeregu trygonometrycznego
wraz z pochodnymi dowolnego ustalonego rzędu oraz
zupełność układu w $L^2$ identyfikują jego sumę
z periodyzacją $\Phi_{ij}$ na iloczynie sześcianów.
Odcięcia eliminują pozostałe kopie okresowe, więc
$\widehat J_{\mathcal E}T_{ij}(\Phi)=\Phi_{ij}$
na całym $\Omega\times\Lambda$.

Definiujemy $T(\Phi)=\sum_{i,j}T_{ij}(\Phi)$.
W seminormie $\pi_{p_{K,r},p_{L,s}}$ niezerowy wkład
mogą dawać tylko pary z
$K\cap\supp\chi_i\ne\varnothing$ i
$L\cap\supp\psi_j\ne\varnothing$.
Jest ich skończenie wiele. Sieć sum skończonych
jest więc Cauchy'ego w każdej seminormie, a jej
granica istnieje z zupełności.
Dla każdej z tych seminorm $T(\Phi)$ szacujemy
skończoną sumą powyższych oszacowań, co dowodzi
ciągłości $T$. Ciągłość $\widehat J_{\mathcal E}$
oraz lokalnie skończony rozkład jedności dają
$\widehat J_{\mathcal E}T(\Phi)=\sum_{i,j}\Phi_{ij}=\Phi$.

Dla $\Phi(x,y)=f(x)g(y)$ współczynniki $c_{kl}^{ij}$
są iloczynami współczynników $\eta_i f$ i $\theta_j g$,
na mocy Fubiniego. Odpowiednie szeregi w poszczególnych grupach zmiennych
z odcięciami zbiegają w $\mathcal E$
do $\eta_i f$ i $\theta_j g$. Sumy po prostokątnych
zbiorach indeksów $(k,l)$ są tensorami ich sum
częściowych. Ciągłość odwzorowania dwuliniowego
daje $T_{ij}(fg)=(\eta_i f)\otimes(\theta_j g)$.
Sumując po prostokątnych skończonych zbiorach
indeksów $(i,j)$, otrzymujemy tensory sum lokalnych,
które zbiegają w $\mathcal E(\Omega)$ i
$\mathcal E(\Lambda)$ do $f$ i $g$.
Stąd $T J_{\mathcal E}(f\otimes g)=f\otimes g$.
Przez liniowość, ciągłość i gęstość algebraicznego
iloczynu w uzupełnieniu mamy $T\widehat J_{\mathcal E}=I$.
Oba izomorfizmy są zatem udowodnione nad $\R$ i $\C$.
Na mocy Twierdzenia~\ref{thm:nuclear-cinfty23}
i Wniosku~\ref{cor:nuclear-schwartz23} pierwsze
czynniki są nuklearne, więc
Twierdzenie~\ref{thm:nuclear-pi-eps24}
daje również izomorfizmy dla $\varepsilon$.
\end{proof}

\setcounter{section}{18}
% Rozdział aktywny po restrukturyzacji DT.
% Roboczy podział AF 23: etykiety historyczne pozostają niezmienione.

\section{Twierdzenie o jądrze Schwartza}
\label{sec:dt-schwartz-kernel}
\providecommand{\atlasref}[2]{#2}

Korzystamy z własności funkcji testowych, przestrzeni
dystrybucji i ciągłości form na etapach granicy induktywnej.
Najpierw udowodnimy rozkład funkcji testowych na szeregi
tensorów elementarnych, a następnie twierdzenie o jądrze
w wersji operatorowej. Dowód nie zakłada jeszcze
identyfikacji ukończonego iloczynu induktywnego funkcji
testowych z przestrzenią testową na iloczynie zbiorów otwartych.

\subsection{Rozkład tensorowy i reprezentacja operatorów}

\noindent\emph{Powiązania tematyczne w Bibliotece matematycznej:} \atlasref{LA-10::thm:lin10-istnienie}{LA-10: algebraiczna konstrukcja iloczynu tensorowego (przed uzupełnieniem topologicznym)}.


\begin{definicja}[Forma dwuliniowa zadana jądrem]
\label{def:kernel-pair23}
Niech $\Omega\subseteq\R^{d_1}$ i
$\Lambda\subseteq\R^{d_2}$ będą niepustymi zbiorami otwartymi.
Dla $K\in\mathcal D'(\Omega\times\Lambda)$ definiujemy
formę na parach funkcji testowych przez
$b_K(f,g)=\langle K,f\otimes g\rangle$,
dla $f\in\mathcal D(\Omega)$ i $g\in\mathcal D(\Lambda)$,
gdzie $(f\otimes g)(x,y)=f(x)g(y)$.
Parowanie jest dwuliniowe także nad $\C$; nie stosujemy
sprzężenia zespolonego.
\end{definicja}

\begin{lemat}[Rozkład testu na szeregi tensorowe]
\label{lem:tensor-density23}
Niech $d_1,d_2\ge1$, a $\Omega\subseteq\R^{d_1}$ i
$\Lambda\subseteq\R^{d_2}$ będą niepustymi zbiorami otwartymi.
Dla każdego $\Phi\in\mathcal D(\Omega\times\Lambda)$
istnieje przedstawienie $\Phi$ jako sumy skończenie
wielu szeregów postaci
$\sum_{n\in\mathbb Z^{d_2}}f_n\otimes g_n$,
zbieżnych w $\mathcal D(\Omega\times\Lambda)$,
gdzie wszystkie $f_n$ mają nośniki w jednym zwartym
$K\Subset\Omega$, wszystkie $g_n$ w jednym zwartym
$L\Subset\Lambda$ i dla każdego $m,r\in\N_0$
\[
 \sum_n p_{K,m}(f_n)\,p_{L,r}(g_n)<\infty.
\]
W szczególności algebraiczna przestrzeń tensorowa
$\mathcal D(\Omega)\otimes\mathcal D(\Lambda)$
jest gęsta w $\mathcal D(\Omega\times\Lambda)$.
\end{lemat}
\begin{proof}
Wybieramy lokalnie skończone gładkie rozkłady jedności
$(\eta_i)$ na $\Omega$ i $(\theta_j)$ na $\Lambda$
podporządkowane pokryciom stosunkowo zwartymi otwartymi
sześcianami. Istnienie takich partycji z nośnikami
zwartymi w poszczególnych sześcianach wynika z
Twierdzenia~\ref{thm:dt-smooth-partition}.
Dla każdego $i,j$ ustalamy zbiory zwarte
$K_i=\supp\eta_i$ i $L_j=\supp\theta_j$ oraz otwarty
sześcian $V_j$ z $L_j\subset V_j\Subset\Lambda$.
Wybieramy rzeczywistą $\chi_j\in\mathcal D(V_j)$,
równą $1$ w otoczeniu $L_j$ (Lemat~\ref{lem:cutoffs17}).
Funkcję $\Phi_{ij}=\eta_i\theta_j\Phi$ przedłużamy
zerem w zmiennej $y$ poza $V_j$ i okresowo z jego
domknięcia $Q_j$; przedłużenie jest gładkie,
gdyż nośnik w $y$ leży w $L_j$.
Niech $a_j$ będzie środkiem $Q_j$, a $2R_j$ jego
długością boku. Kładziemy
\[
 e_{jn}(y)=e^{i\pi n\cdot(y-a_j)/R_j},
 \quad
 f_{ijn}(x)=\frac{1}{(2R_j)^{d_2}}
 \int_{Q_j}\Phi_{ij}(x,y)\overline{e_{jn}(y)}
 \,\mathrm d\lambda_{d_2}(y),
 \quad g_{jn}=\chi_j e_{jn}.
\]
Funkcje $f_{ijn}$ są gładkie i mają nośniki w $K_i$:
pochodne względem $x$ obliczamy pod całką po zwartym
sześcianie, stosując lokalnie ograniczenie pochodnych
na iloczynie zwartego otoczenia punktu $x$ i $Q_j$.
Funkcję $g_{jn}$ przedłużamy przez zero poza $V_j$;
jest gładka, gdyż $\chi_j$ ma zwarty nośnik w $V_j$.
Niech $L_j^*=\supp\chi_j\cup L_j$.
Dla $m,r,t\in\N_0$, całkując przez części za pomocą
$(1-\Delta_y)^t$, otrzymujemy
\[
 \begin{aligned}
 p_{K_i,m}(f_{ijn})&\le C_{j,t}
 (1+|n|)^{-2t}p_{K_i\times L_j^*,m+2t}(\Phi_{ij}),\\
 p_{L_j^*,r}(g_{jn})&\le C_{j,r}'(1+|n|)^r .
 \end{aligned}
\]
Pierwsza nierówność wynika z tożsamości
$(1-\Delta_y)e_{jn}=(1+\pi^2|n|^2/R_j^2)e_{jn}$,
braku wyrazów brzegowych oraz ograniczenia pochodnych
funkcji podcałkowej; druga ze wzoru Leibniza.
Przy $2t>r+d_2$ szereg iloczynów seminorm jest zbieżny
(liczba punktów kratowych o normie maksimum $s$
jest rzędu co najwyżej $(1+s)^{d_2-1}$, jak w
dowodzie Twierdzenia~\ref{thm:nuclear-cinfty23}).
W szczególności dla każdego $q$ szereg
$\sum_n f_{ijn}\otimes g_{jn}$ jest bezwzględnie
zbieżny w seminormie $p_{K_i\times L_j^*,q}$,
stosując powyższe oszacowanie z $m=r=q$.
Zupełność etapu $\mathcal D_{K_i\times L_j^*}$
(Twierdzenie~\ref{thm:DK-frechet18}) daje granicę.

Dla ustalonego $x$ szereg $\sum_n f_{ijn}(x)e_{jn}(y)$
na $Q_j$ jest jednostajnie zbieżny, a jego granicę
identyfikujemy z $\Phi_{ij}(x,\cdot)$ dzięki
zupełności układu trygonometrycznego w $L^2(Q_j)$
(udowodnionej w dowodzie Twierdzenia~\ref{thm:nuclear-cinfty23}
dla dowolnego wymiaru, po przesunięciu sześcianu)
i ciągłości obu funkcji. Po pomnożeniu przez $\chi_j$
równość obowiązuje na całej $\Lambda$: na $Q_j$
mamy $\chi_j\Phi_{ij}=\Phi_{ij}$, a poza $Q_j$
obie strony znikają. Warunek $\supp\chi_j\subset V_j$
jest istotny: eliminuje inne kopie okresowe.
Stąd $\Phi_{ij}=\sum_n f_{ijn}\otimes g_{jn}$.
Tylko skończenie wiele par $(i,j)$ daje niezerowe
$\Phi_{ij}$, dzięki lokalnej skończoności i zwartości
projekcji nośnika $\Phi$. Sumujemy te rozwinięcia,
biorąc skończone sumy odpowiednich zwartych zbiorów
$K_i$ i $L_j^*$. Daje to wspólny zwarty nośnik
i żądaną gęstość.

Nad $\R$ w powyższych szeregach dla rzeczywistej $\Phi$
zastępujemy każdy składnik przez jego część rzeczywistą:
$\operatorname{Re}f_{ijn}\otimes\operatorname{Re}g_{jn}
-\operatorname{Im}f_{ijn}\otimes\operatorname{Im}g_{jn}$.
Otrzymujemy dwa szeregi rzeczywistych tensorów,
każdy indeksowany przez $\mathbb Z^{d_2}$; znak minus
włączamy do pierwszego czynnika drugiego szeregu.
Seminormy części rzeczywistej i urojonej nie przekraczają
seminorm funkcji zespolonej, więc suma pozostaje
zbieżna z tymi samymi oszacowaniami
(z dodatkowym czynnikiem $2$).
\end{proof}

\begin{lemat}[Ciągłość formy na ustalonych etapach]
\label{lem:separate-joint23}
Niech $E,F$ będą przestrzeniami Frécheta nad
$\K\in\{\R,\C\}$, a $b:E\times F\to\K$
formą dwuliniową ciągłą osobno w każdej zmiennej.
Wtedy $b$ jest łącznie ciągła.
\end{lemat}
\begin{proof}
Wybierzmy malejącą przeliczalną bazę
$(V_n)$ otoczeń zera w $F$, złożoną ze zbiorów
zrównoważonych. Niech
$A_n=\{x\in E: |b(x,y)|\le1\ \text{dla wszystkich }y\in V_n\}$.
Każdy $A_n$ jest domknięty, ponieważ $x\mapsto b(x,y)$
jest ciągły dla ustalonego $y$. Z ciągłości w drugiej
zmiennej wynika $E=\bigcup_n A_n$.
Z twierdzenia Baire'a pewien $A_n$ ma niepuste wnętrze;
ponieważ $A_n$ jest wypukły, zrównoważony i domknięty,
ma wówczas otoczenie zera we wnętrzu. Istotnie,
jeśli $x_0+U\subseteq A_n$ dla otwartego symetrycznego
otoczenia zera $U$, to również $-x_0+U\subseteq A_n$;
wypukłość daje $U\subseteq A_n$.
Stąd $|b(x,y)|\le1$ na $U\times V_n$.
Skalowanie daje ciągłość w $(0,0)$, a dla
dowolnego $(x_0,y_0)$ korzystamy z rozwinięcia
$b(x_0+x,y_0+y)-b(x_0,y_0)
=b(x_0,y)+b(x,y_0)+b(x,y)$ i ciągłości osobnej.
\end{proof}

\begin{uwaga}[Ciągłość osobna i łączna]
\label{rem:kernel-separate23}
Dla formy $b:\mathcal D(\Omega)\times\mathcal D(\Lambda)\to\K$
ciągłość osobna nie musi implikować łącznej ciągłości na
całym iloczynie przestrzeni LF. Implikacja ta jest jednak
prawdziwa po ograniczeniu formy do
$\mathcal D_K(\Omega)\times\mathcal D_L(\Lambda)$,
gdyż obie przestrzenie są przestrzeniami Frécheta.
To rozróżnienie jest istotne w twierdzeniu o jądrze;
właściwą globalną topologią tensorową jest topologia
induktywna, omówiona w następnym rozdziale.
\end{uwaga}

\begin{twierdzenie}[Twierdzenie o jądrze Schwartza]
\label{thm:schwartz-kernel23}
Niech $\Omega\subseteq\R^{d_1}$ i
$\Lambda\subseteq\R^{d_2}$ będą niepustymi zbiorami
otwartymi. Dla każdego ciągłego operatora liniowego
\[
 A:\mathcal D(\Omega)\longrightarrow
 \mathcal D'(\Lambda)_\beta
\]
istnieje dokładnie jedna dystrybucja
$K_A\in\mathcal D'(\Omega\times\Lambda)$ taka, że
\[
 \langle Af,g\rangle=\langle K_A,f\otimes g\rangle
 \qquad(f\in\mathcal D(\Omega),\ g\in\mathcal D(\Lambda)).
\]
Odwrotnie, każda dystrybucja
$K\in\mathcal D'(\Omega\times\Lambda)$ określa w ten
sposób ciągły operator liniowy
$A_K:\mathcal D(\Omega)\to\mathcal D'(\Lambda)_\beta$.
Indeks $\beta$ oznacza silną topologię przestrzeni sprzężonej,
czyli jednostajną zbieżność na ograniczonych zbiorach
funkcji testowych.
\end{twierdzenie}
\begin{proof}
Niech $b(f,g)=\langle Af,g\rangle$. Ciągłość $A$
i ciągłość ewaluacji na pojedynczej funkcji testowej
zapewniają osobną ciągłość $b$. Na przestrzeniach
$\mathcal D_{K_x}(\Omega)$ i $\mathcal D_{K_y}(\Lambda)$
forma $b$ jest wspólnie ciągła na mocy
Lematu~\ref{lem:separate-joint23}; zatem istnieją
$C,m,r$ takie, że
$|b(f,g)|\le C p_{K_x,m}(f)p_{K_y,r}(g)$.
Istotnie, ciągłość w $(0,0)$ daje iloczyn otoczeń
zera, na którym $|b|\le1$. Każde z tych otoczeń
zawiera kulę jednej seminormy etapowej o dodatnim
promieniu, ponieważ rodziny seminorm są rosnące.
Skalowanie obu argumentów, jak w dowodzie
Twierdzenia~\ref{thm:pi-universal24}, daje wskazane
oszacowanie, również gdy któraś seminorma znika.
Zbudujemy jeden liniowy funkcjonał na całej przestrzeni
testowej, a następnie sprawdzimy jego niezależność
od wyborów pomocniczych. Wybierzmy lokalnie skończone
rozkłady jedności $(\eta_i)$ na $\Omega$ oraz
$(\theta_j)$ na $\Lambda$, których elementy mają zwarte
nośniki w sześcianach stosunkowo zwartych. Produkty
$\eta_i(x)\theta_j(y)$ tworzą rozkład jedności na
$\Omega\times\Lambda$. Dla danej
$\Phi\in\mathcal D(\Omega\times\Lambda)$ tylko skończenie
wiele funkcji $\Phi_{ij}=\eta_i\theta_j\Phi$ jest niezerowych.
Dla każdej pary $(i,j)$ wybieramy z góry odcięcie
$\chi_j$ i układ Fouriera z dowodu
Lematu~\ref{lem:tensor-density23}, niezależnie od $\Phi$.
Otrzymujemy $\Phi_{ij}=\sum_n f_{ijn}\otimes g_{jn}$,
gdzie $f_{ijn}$ zależy liniowo od $\Phi$,
a $g_{jn}$ nie zależy od $\Phi$.
Nad $\R$ stosujemy opisane w lemacie dwa rzeczywiste
składniki; dalej zapisujemy wzory w przypadku zespolonym.
Definiujemy
\[
 \langle K_A,\Phi\rangle=\sum_{i,j}\sum_n b(f_{ijn},g_{jn}),
\]
gdzie suma po $i,j$ jest skończona, a wszystkie szeregi
wewnętrzne są bezwzględnie zbieżne na mocy oszacowań
seminormowych i ciągłości $b$ na odpowiednich etapach.
Konstrukcja jest liniowa, ponieważ wybory odcięć i baz
Fouriera nie zależą od $\Phi$, a współczynniki Fouriera
zależą od niej liniowo.

Aby sprawdzić ciągłość, ustalmy $M\Subset\Omega\times\Lambda$.
Lokalna skończoność rozkładów jedności daje skończoną
liczbę par $(i,j)$, które mogą wystąpić dla
$\Phi\in\mathcal D_M$. Dla każdej takiej pary istnieją $C_{ij},m_{ij},r_{ij}$
z oszacowania dwuliniowego na $\mathcal D_{K_i}
\times\mathcal D_{L_j^*}$. Wybieramy $t_{ij}$ z
$2t_{ij}>r_{ij}+d_2$. Z oszacowań w dowodzie
Lematu~\ref{lem:tensor-density23} i wzoru Leibniza
dla $\eta_i\theta_j$ otrzymujemy
\[
 |b(f_{ijn},g_{jn})|\le C'_{ij}
 (1+|n|)^{r_{ij}-2t_{ij}}
 p_{M,m_{ij}+2t_{ij}}(\Phi).
\]
Pochodne $\Phi$ znikają poza $M$, zatem po prawej
wystarcza seminorma na $M$. Sumujemy po $n$,
a następnie po skończonej liczbie par i bierzemy
$r_M=\max_{i,j}(m_{ij}+2t_{ij})$ (dla pustej
rodziny par funkcjonał jest zerowy). Otrzymujemy
stałą $C_M>0$ i $r_M\in\N_0$ takie, że
\[
 |\langle K_A,\Phi\rangle|\le C_M p_{M,r_M}(\Phi)
 \qquad (\Phi\in\mathcal D_M(\Omega\times\Lambda)).
\]
Jest to ciągłość na każdym etapie, więc własność
uniwersalna granicy induktywnej daje
$K_A\in\mathcal D'(\Omega\times\Lambda)$.

Dla $\Phi=f\otimes g$ współczynniki mają szczególną
postać $f_{ijn}=(\eta_i f)c_{jn}(\theta_j g)$,
gdzie $c_{jn}$ jest skalarnym współczynnikiem Fouriera
w $Q_j$. Z dowodu lematu
$\sum_n c_{jn}(\theta_j g)g_{jn}=\theta_j g$
w ustalonym etapie $\mathcal D_{L_j^*}$.
Liniowość i ciągłość $b$ względem drugiego argumentu
przy ustalonym $\eta_i f$ dają
\[
 \sum_n b(f_{ijn},g_{jn})=b(\eta_i f,\theta_j g).
\]
Nad $\R$ tę samą równość otrzymujemy, rozwijając
rzeczywistą część szeregu i korzystając z rzeczywistej
dwuliniowości. Sumując skończenie wiele składników
rozkładów jedności, dostajemy
$\langle K_A,f\otimes g\rangle=b(f,g)$.
Jeżeli zmienimy rozkłady jedności i odcięcia, otrzymamy
inny ciągły funkcjonał na całej przestrzeni $\mathcal D$,
ale o tych samych wartościach na tensorach elementarnych.
Ich liniowa powłoka jest gęsta na mocy
Lematu~\ref{lem:tensor-density23}; dwa funkcjonały muszą
więc być równe. To zarazem uzasadnia jednoznaczność,
bez zakładania gęstości tensorów o nośnikach w arbitralnym
zwartym zbiorze $M$.

W odwrotną stronę, dla $K\in\mathcal D'(\Omega\times\Lambda)$
forma $b_K(f,g)=\langle K,f\otimes g\rangle$ jest
łącznie ciągła na każdym
$\mathcal D_{K_x}(\Omega)\times\mathcal D_{K_y}(\Lambda)$,
ponieważ dla ustalonych niepustych zwartych $K_x,K_y$
z ciągłości $K$ na etapie $\mathcal D_{K_x\times K_y}$
otrzymujemy pewne $C>0$ i $q\in\N_0$ z
\[
 |\langle K,f\otimes g\rangle|
 \le C p_{K_x\times K_y,q}(f\otimes g)
 \le C p_{K_x,q}(f)p_{K_y,q}(g).
\]
Ostatnia nierówność wynika z rozdzielenia pochodnych
względem $x$ i $y$. Dla ustalonego $f$ wybieramy
$K_x$ zawierający jego nośnik. Forma $g\mapsto b_K(f,g)$
jest liniowa i ciągła na każdym etapie $\mathcal D_{K_y}$,
więc na mocy Twierdzenia~\ref{thm:LF-universal18}
wyznacza $A_Kf\in\mathcal D'(\Lambda)$. Dwuliniowość
zapewnia liniowość $A_K$.
Aby sprawdzić ciągłość w topologii silnej, niech
$B\subseteq\mathcal D(\Lambda)$ będzie ograniczony.
Z Twierdzenia~\ref{thm:D-bounded18} istnieje
$K_y\Subset\Lambda$, dla którego $B$ jest ograniczony
w $\mathcal D_{K_y}$. Na każdym etapie $\mathcal D_{K_x}$
mamy, dzięki wspólnej ciągłości formy,
\[
 \sup_{g\in B}|\langle A_Kf,g\rangle|
 \le C\,p_{K_x,m}(f)\sup_{g\in B}p_{K_y,r}(g).
\]
Ostatnie supremum po prawej jest skończone
z ograniczoności $B$ na etapie $\mathcal D_{K_y}$.
Lewa strona jest seminormą silnej topologii
(Definicja~\ref{def:strong19}) obliczoną na $A_Kf$.
Ponieważ $B$ jest dowolnym zbiorem ograniczonym,
oszacowanie dowodzi ciągłości $A_K$ na każdym etapie,
a więc, przez własność uniwersalną granicy induktywnej,
na całej $\mathcal D(\Omega)$. Jednoznaczność $K_A$
wynika ponownie z gęstości tensorów elementarnych.
\end{proof}

\begin{wniosek}[Słaba-* i silna ciągłość operatora o jądrze]
\label{cor:dt-kernel-weak-strong}
Niech $A:\mathcal D(\Omega)\to\mathcal D'(\Lambda)$
będzie odwzorowaniem liniowym. Jeżeli $A$ jest ciągłe
dla topologii $\sigma(\mathcal D'(\Lambda),\mathcal D(\Lambda))$
w przestrzeni docelowej, to jest również ciągłe
dla topologii silnej $\beta$. W szczególności twierdzenie
o jądrze obowiązuje przy założeniu ciągłości słabej-*.
\end{wniosek}
\begin{proof}
Dla ustalonego $g\in\mathcal D(\Lambda)$ funkcjonał
$f\mapsto\langle Af,g\rangle$ jest ciągły z definicji
topologii słabej-* (Definicja~\ref{def:weakstar19}).
Dla ustalonego $f$ ciągłość względem $g$ wynika
z tego, że $Af$ jest dystrybucją. Forma
$b(f,g)=\langle Af,g\rangle$ jest więc osobno ciągła.
Na dowolnej parze etapów $\mathcal D_{K_x}$,
$\mathcal D_{K_y}$ Lemat~\ref{lem:separate-joint23}
i oszacowanie z początku dowodu
Twierdzenia~\ref{thm:schwartz-kernel23} dają
\[
 |b(f,g)|\le C p_{K_x,m}(f)p_{K_y,r}(g).
\]
Dla niepustego ograniczonego $B\subseteq\mathcal D(\Lambda)$
wybieramy $K_y$ z Twierdzenia~\ref{thm:D-bounded18},
tak aby $B$ było ograniczone w $\mathcal D_{K_y}$.
Wówczas dla $f\in\mathcal D_{K_x}$
\[
 \sup_{g\in B}|\langle Af,g\rangle|
 \le C p_{K_x,m}(f)\sup_{g\in B}p_{K_y,r}(g).
\]
W ten sposób każda seminorma silnej topologii
złożona z $A$ jest ciągła na dowolnym etapie.
Twierdzenie~\ref{thm:LF-universal18} daje ciągłość
$A$ na całej $\mathcal D(\Omega)$.
Implikacja odwrotna wynika z ograniczoności
jednoelementowych zbiorów testowych. Nie oznacza
to równości topologii słabej-* i silnej na całej
$\mathcal D'(\Lambda)$.
\end{proof}

\begin{przyklad}[Jądro operatora o gładkim jądrze]
\label{ex:smooth-kernel23}
Jeżeli $k\in C^\infty(\Omega\times\Lambda)$, to dystrybucja
regularna $T_k$ jest jądrem operatora
$(Af)(y)=\int_\Omega k(x,y)f(x)\,\mathrm d\lambda_{d_1}(x)$.
Dla $g\in\mathcal D(\Lambda)$ mamy
\[
\langle Af,g\rangle
=\int_{\Lambda}\int_{\Omega}k(x,y)f(x)g(y)
\,\mathrm d\lambda_{d_1}(x)\,\mathrm d\lambda_{d_2}(y).
\]
Dla ustalonego $f$ różniczkowanie względem $y$
pod całką jest uzasadnione przez ograniczoność pochodnych
$k$ na iloczynie $\supp f$ i dowolnego zwartego
otoczenia punktu $y$; funkcja $Af$ jest gładka.
Twierdzenie Fubiniego stosujemy, ponieważ funkcja
podcałkowa jest ograniczona i znika poza
$\supp f\times\supp g$. Nie wymagamy, żeby $k$
miało zwarty nośnik. Identyfikując $Af$ z dystrybucją
regularną, z twierdzenia o jądrze otrzymujemy
ciągłość operatora $\mathcal D(\Omega)\to
\mathcal D'(\Lambda)_\beta$.
\end{przyklad}

\begin{przyklad}[Jądro osobliwe na przekątnej]
\label{ex:diagonal-kernel23}
Dla $\Omega=\Lambda\subseteq\R^d$ operator naturalnego
włożenia $\mathcal D(\Omega)\hookrightarrow\mathcal D'(\Omega)$
ma jądro $\delta_\Delta$ zadane przez
$\langle\delta_\Delta,\Phi\rangle
=\int_\Omega\Phi(x,x)\,\mathrm d\lambda_d(x)$.
Dla $\Phi\in\mathcal D_M(\Omega\times\Omega)$
całka jest ograniczona przez
$\lambda_d(\operatorname{pr}_1 M)p_{M,0}(\Phi)$,
co dowodzi ciągłości na każdym etapie.
Na tensorze $f\otimes g$ otrzymujemy $\int fg\,
\mathrm d\lambda_d$, czyli parowanie naturalnego
włożenia. Nośnik $\delta_\Delta$ jest zawarty
w przekątnej $\Delta=\{(x,x):x\in\Omega\}$;
każde otoczenie punktu $(x,x)$ zawiera iloczyn
małych kul o środku $x$, na którym można wybrać
nieujemny tensor $f\otimes f$ z $f\ne0$.
Wtedy $\langle\delta_\Delta,f\otimes f\rangle>0$,
więc nośnik jest równy $\Delta$.
Gdyby $\delta_\Delta=T_h$ dla
$h\in L^1_{\mathrm{loc}}(\Omega\times\Omega)$,
to ograniczenie do otwartego dopełnienia $\Delta$
byłoby zerowe. Różnowartościowość przyporządkowania
funkcji dystrybucji regularnej daje $h=0$ prawie
wszędzie poza $\Delta$. Przekątna ma miarę
$\lambda_{2d}$ zero: dla każdego $x\in\Omega$ jej
przekrój względem drugiej zmiennej jest singletonem
$\{x\}$ miary $\lambda_d$ zero, a zatem z twierdzenia
Tonellego dla $\chi_\Delta$ wynika
$\lambda_{2d}(\Delta)=0$. Stąd $h=0$ prawie wszędzie,
co przeczy $\delta_\Delta\ne0$.
\end{przyklad}

\begin{przyklad}[Jądra operatorów różniczkowych]
\label{ex:dt-kernel-differential}
Niech $\Omega\subseteq\R^d$ będzie niepustym zbiorem
otwartym i niech
\[
 P=\sum_{|\alpha|\le m}a_\alpha\partial^\alpha,
 \qquad a_\alpha\in C^\infty(\Omega).
\]
Operator $A_P:\mathcal D(\Omega)\to\mathcal D'(\Omega)_\beta$
zadany przez $A_P f=T_{Pf}$ ma jądro
\[
 K_P=\sum_{|\alpha|\le m}
 (-1)^{|\alpha|}a_\alpha(y)\partial_x^\alpha\delta_\Delta
 \quad\text{w }\mathcal D'(\Omega\times\Omega).
\]
Istotnie, pochodne i mnożenie przez funkcje gładkie
są określone na dystrybucjach (Definicja~\ref{def:dist-deriv20}
i Definicja~\ref{def:multiply20}). Działanie każdego
składnika na dowolnej funkcji testowej $\Phi$ jest
równe
\[
 \left\langle
 (-1)^{|\alpha|}a_\alpha(y)\partial_x^\alpha\delta_\Delta,
 \Phi\right\rangle
 =\int_\Omega a_\alpha(t)
       (\partial_x^\alpha\Phi)(t,t)\,\mathrm d\lambda_d(t).
\]
Wynika to z definicji pochodnej dystrybucyjnej:
znak pojawia się dwukrotnie, a czynnik $a_\alpha(y)$
nie zależy od $x$. Całka jest skończona, ponieważ
pochodna $\Phi$ ma zwarty nośnik.
Dla $\Phi=f\otimes g$ otrzymujemy więc
\[
 \langle K_P,f\otimes g\rangle
 =\int_\Omega (Pf)(t)g(t)\,\mathrm d\lambda_d(t).
\]
Twierdzenie~\ref{thm:schwartz-kernel23} daje ciągłość
$A_P$ oraz jednoznaczność wskazanego jądra.
Ponadto $\supp K_P\subseteq\Delta$: dla funkcji
testowej o nośniku poza przekątną wszystkie
jej pochodne również znikają na przekątnej.
W szczególności jądro operatora pochodnej
$\partial_j$ to $-\partial_{x_j}\delta_\Delta$.
Nie jest to to samo co pochodna wzdłuż przekątnej.
\end{przyklad}

\begin{uwaga}[Zakres twierdzenia o jądrze]
\label{rem:kernel-scope23}
Twierdzenie~\ref{thm:schwartz-kernel23} opisuje wszystkie
ciągłe operatory $\mathcal D(\Omega)\to\mathcal D'(\Lambda)_\beta$.
Nie oznacza ono, że jądro jest funkcją ani że operacja
składania dowolnych takich operatorów jest automatycznie
zdefiniowana. Twierdzenia trwałości nuklearności zostały udowodnione
w rozdziale~\ref{sec:montel-nuklearne23}; następny rozdział omawia
induktywny iloczyn tensorowy funkcji testowych. Pełne rozwiązywanie równań
różniczkowych pozostaje w osobnym skrypcie.
\end{uwaga}

\setcounter{section}{19}
% Rozdział aktywny po restrukturyzacji DT.
% Roboczy podział AF 24: etykiety historyczne pozostają niezmienione.

\section{Induktywne iloczyny tensorowe funkcji testowych}
\label{sec:dt-test-tensor}
\providecommand{\atlasref}[2]{#2}

Wykorzystujemy teorię topologicznych iloczynów tensorowych
z rozdziału~\ref{sec:tensor24}, rozkład funkcji testowych
z lematu~\ref{lem:tensor-density23} oraz twierdzenie
o jądrze Schwartza~\ref{thm:schwartz-kernel23}.
Przestrzenie rozpatrujemy nad wspólnym ciałem $\K$.

\subsection{Topologia induktywna i jej własność uniwersalna dla \texorpdfstring{$\mathcal D(\Omega)$}{D(Omega)}}

\begin{definicja}[Induktywna topologia tensorowa]
\label{def:i-topology24}
Niech $E,F$ będą przestrzeniami lokalnie wypukłymi.
\emph{Induktywną topologią tensorową} $i$ na $E\otimes F$
nazywamy najsilniejszą topologię lokalnie wypukłą,
dla której kanoniczne odwzorowanie
$E\times F\to E\otimes F$ jest osobno ciągłe.
Piszemy $E\otimes_i F$, a przez
$E\widehat\otimes_i F$ oznaczamy jego uzupełnienie.
\end{definicja}

\begin{stwierdzenie}[Rola topologii induktywnej]
\label{prop:i-vs-pi24}
Dla dowolnych przestrzeni lokalnie wypukłych Hausdorffa
$E,F$ topologia induktywna $i$ istnieje, jest Hausdorffa
i nie jest słabsza od rzutowej $\pi$.
Dla każdej przestrzeni lokalnie wypukłej Hausdorffa $G$
linearyzacja daje bijekcję między osobno ciągłymi
odwzorowaniami dwuliniowymi $E\times F\to G$
a ciągłymi odwzorowaniami liniowymi $E\otimes_i F\to G$.
Jeśli $G$ jest zupełna, te ostatnie przedłużają się
jednoznacznie do $E\widehat\otimes_i F$.
W szczególności ciągłe funkcjonały liniowe na
$E\otimes_i F$ odpowiadają dokładnie osobno ciągłym
formom dwuliniowym $E\times F\to\K$. W przypadku
funkcji testowych to właśnie topologia induktywna
odpowiada twierdzeniu o jądrze Schwartza.
\end{stwierdzenie}


\begin{proof}
Rozważmy rodzinę wszystkich seminorm $h$ na
$E\otimes F$, dla których dla ustalonych $x\in E$
i $y\in F$ seminormy $h(x\otimes\cdot)$ na $F$
i $h(\cdot\otimes y)$ na $E$ są ciągłe.
Topologia generowana przez tę rodzinę zapewnia
osobną ciągłość kanonicznego odwzorowania.
Każda seminorma ciągła w dowolnej dopuszczalnej
topologii należy do tej rodziny, więc otrzymana
topologia jest najsilniejsza. Zawiera topologię
$\pi$, gdyż jej odwzorowanie kanoniczne jest
łącznie ciągłe. Własność Hausdorffa wynika z
Lematu~\ref{lem:pi-properties24}.

Niech $b:E\times F\to G$ będzie osobno ciągłe.
Dla każdej ciągłej seminormy $s$ na $G$ funkcja
$h_s=s\circ\widetilde b$ jest seminormą należącą
do opisanej rodziny: jej ograniczenia do obu
czynników przy ustalonym drugim argumencie są
ciągłe. Stąd $\widetilde b:E\otimes_i F\to G$
jest ciągłe. Odwrotnie, złożenie ciągłego liniowego
odwzorowania z osobno ciągłym odwzorowaniem
kanonicznym jest osobno ciągłe. Jednoznaczność
wynika z własności algebraicznej. Ciągłość
jednostajna odwzorowania liniowego i zupełność
$G$ dają jednoznaczne przedłużenie na uzupełnienie.
\end{proof}

Dla osobno ciągłego odwzorowania dwuliniowego $b:E\times F\to G$ oznaczamy przez $\tau$ kanoniczne odwzorowanie dwuliniowe. Strzałka przerywana oznacza jedyne odwzorowanie wyznaczone przez własność uniwersalną.
\begin{center}
\begin{tikzpicture}[>=stealth]
 \node (a) at (0,2.0) {$E\times F$};
 \node (b) at (4.4,2.0) {$E\otimes_i F$};
 \node (c) at (2.2,0) {$G$};
 \draw[->] (a) -- node[above] {$\tau$} (b);
 \draw[->] (a) -- node[below left] {$b$} (c);
 \draw[->,dashed] (b) -- node[below right] {$\widetilde b$} (c);
\end{tikzpicture}
\end{center}

\begin{stwierdzenie}[Równość topologii induktywnej i rzutowej dla przestrzeni Frécheta]
\label{prop:dt-i-pi-frechet}
Jeżeli $E,F$ są przestrzeniami Frécheta, a $G$ jest
przestrzenią lokalnie wypukłą Hausdorffa, to każde
osobno ciągłe odwzorowanie dwuliniowe $b:E\times F\to G$
jest łącznie ciągłe. W szczególności na $E\otimes F$
topologie $i$ i $\pi$ są równe. Identyczność przedłuża
się do kanonicznego izomorfizmu
$E\widehat\otimes_i F\simeq E\widehat\otimes_\pi F$.
\end{stwierdzenie}
\begin{proof}
Ustalmy ciągłą seminormę $s$ na $G$. Niech
$(V_n)$ będzie malejącą przeliczalną bazą otwartych
absolutnie wypukłych otoczeń zera w $F$. Zbiory
\[
 A_n=\bigcap_{y\in V_n}\{x\in E:s(b(x,y))\le1\}
\]
są domknięte i absolutnie wypukłe: dla każdego $y$
odwzorowanie $x\mapsto b(x,y)$ jest liniowe i ciągłe.
Ciągłość w drugiej zmiennej daje $E=\bigcup_n A_n$.
Z twierdzenia Baire'a pewien $A_n$ ma niepuste wnętrze.
Argument z Lematu~\ref{lem:separate-joint23} pokazuje,
że $A_n$ zawiera otoczenie zera $U$ (korzystamy
z wypukłości i symetrii $A_n$).
Wobec tego $s(b(x,y))\le1$ na $U\times V_n$.
Wybieramy ciągłe seminormy $p,q$ odpowiednio na $E,F$,
których otwarte kule jednostkowe zawierają się
odpowiednio w $U,V_n$. Dla dodatnich $a>p(x)$
i $c>q(y)$ mamy $s(b(x,y))\le ac$; przechodząc
do infimum, otrzymujemy
\[
 s(b(x,y))\le p(x)q(y).
\]
Obejmuje to także przypadek znikania którejś seminormy.
Ponieważ $s$ była dowolna, $b$ jest ciągłe w $(0,0)$.
Rozwinięcie
$b(x_0+x,y_0+y)-b(x_0,y_0)
=b(x_0,y)+b(x,y_0)+b(x,y)$
i ciągłość osobna zapewniają ciągłość w każdym punkcie.

Stosujemy ten wynik do kanonicznego odwzorowania
$E\times F\to E\otimes_i F$, które jest osobno ciągłe
z definicji, a którego przeciwdziedzina jest Hausdorffa
na mocy Stwierdzenia~\ref{prop:i-vs-pi24}.
Własność uniwersalna topologii rzutowej
(Twierdzenie~\ref{thm:pi-universal24}) daje ciągłość
identyczności $E\otimes_\pi F\to E\otimes_i F$.
Odwrotna identyczność jest ciągła na mocy
Stwierdzenia~\ref{prop:i-vs-pi24}.
Izomorfizm przedłuża się wraz ze swoją odwrotnością
do uzupełnień, a ich złożenia są identycznościami
przez gęstość iloczynu algebraicznego.
\end{proof}

\begin{wniosek}[Trzy topologie dla nuklearnych przestrzeni Frécheta]
\label{cor:dt-three-tensor-frechet}
Jeśli $E,F$ są przestrzeniami Frécheta i przynajmniej
jedna z nich jest nuklearna, to na $E\otimes F$
mamy $\varepsilon=\pi=i$; ich uzupełnienia są
kanonicznie izomorficzne.
\end{wniosek}
\begin{proof}
Równość $\pi=i$ wynika z poprzedniego stwierdzenia.
Równość $\varepsilon=\pi$ wynika z
Twierdzenia~\ref{thm:nuclear-pi-eps24}, w razie potrzeby
po zamianie czynników na mocy
Stwierdzenia~\ref{prop:dt-tensor-functoriality}.
\end{proof}

\begin{twierdzenie}[Produkt tensorowy funkcji testowych]
\label{thm:D-i-product24}
Niech $d_1,d_2\ge1$ będą liczbami całkowitymi, a
$\Omega\subseteq\R^{d_1}$ i $\Lambda\subseteq\R^{d_2}$
niepustymi zbiorami otwartymi. Wówczas odwzorowanie
$f\otimes g\mapsto ((x,y)\mapsto f(x)g(y))$ przedłuża się
do kanonicznego izomorfizmu liniowo-topologicznego
\[
 \mathcal D(\Omega)\widehat\otimes_i\mathcal D(\Lambda)
 \simeq \mathcal D(\Omega\times\Lambda).
\]
\end{twierdzenie}
\begin{proof}
Najpierw dowodzimy następującego faktu. Jeżeli $G$ jest
zupełną przestrzenią lokalnie wypukłą Hausdorffa, to każde
osobno ciągłe odwzorowanie dwuliniowe o wartościach w $G$,
$b:\mathcal D(\Omega)\times\mathcal D(\Lambda)\to G$,
ma jednoznaczne ciągłe liniowe przedłużenie
$B:\mathcal D(\Omega\times\Lambda)\to G$, spełniające
$B(f\otimes g)=b(f,g)$.

Dla ustalonych zwartych $K\Subset\Omega$,
$L\Subset\Lambda$ ograniczenie $b$ do
$\mathcal D_K(\Omega)\times\mathcal D_L(\Lambda)$ jest
łącznie ciągłe na mocy Stwierdzenia~\ref{prop:dt-i-pi-frechet},
gdyż oba etapy są przestrzeniami Frécheta
(Twierdzenie~\ref{thm:DK-frechet18}).
W dowodzie tego stwierdzenia dla każdej ciągłej seminormy
$s$ na $G$ otrzymaliśmy oszacowanie iloczynem dwóch
ciągłych seminorm. Ponieważ na etapach topologie
wyznaczają rosnące rodziny seminorm $p_{K,m}$ i $p_{L,r}$,
otrzymujemy
$s(b(f,g))\le C p_{K,m}(f)p_{L,r}(g)$ dla pewnych
$m,r,C$, zależnych od $s,K,L$.

Ustalamy rozkłady jedności, sześciany i odcięcia
z dowodu Lematu~\ref{lem:tensor-density23}, niezależnie
od $\Phi$. Dla $\Phi\in\mathcal D(\Omega\times\Lambda)$
ten lemat daje skończenie wiele lokalnych szeregów
$\Phi=\sum_{i,j}\sum_n f_{ijn}\otimes g_{jn}$,
gdzie dla każdej pary $(i,j)$ wszystkie czynniki
mają nośniki w ustalonych zwartych $K_i,L_j^*$.
Nad $\R$ stosujemy rzeczywistą część rozwinięcia
w sposób opisany w tym lemacie, czyli dwa składniki
rzeczywiste dla każdego indeksu.
Definiujemy
\[
 B(\Phi)=\sum_{i,j}\sum_n b(f_{ijn},g_{jn}).
\]
Dla każdej ciągłej seminormy $s$ na $G$ szereg
wewnętrzny jest bezwzględnie sumowalny względem $s$,
na mocy oszacowania na etapach i lematu.
Sumy częściowe są więc ciągiem Cauchy'ego
w topologii $G$; zupełność zapewnia istnienie granicy.

Ustalmy zwarty $M\Subset\Omega\times\Lambda$.
Lokalna skończoność rozkładów jedności daje
skończony zbiór par $(i,j)$ mogących wystąpić
dla wszystkich $\Phi\in\mathcal D_M$.
Dla każdej pary i seminormy $s$ wybieramy
$C_{ij},m_{ij},r_{ij}$ z oszacowania dwuliniowego
na etapach $K_i,L_j^*$. Wybieramy całkowite
$t_{ij}$ z $2t_{ij}>r_{ij}+d_2$.
Oszacowania Fouriera z dowodu lematu i reguła
Leibniza dla ustalonych odcięć dają
\[
 s(b(f_{ijn},g_{jn}))\le C'_{ij}
 (1+|n|)^{r_{ij}-2t_{ij}}
 p_{M,m_{ij}+2t_{ij}}(\Phi).
\]
Pochodne $\Phi$ znikają poza $M$.
Sumując po $n$, a następnie po skończonej rodzinie
par, otrzymujemy $s(B(\Phi))\le C_s p_{M,r_s}(\Phi)$
z $r_s$ równym maksimum odpowiednich rzędów
(dla pustej rodziny $B$ jest zerowe na etapie).
Stałe nie zależą od $\Phi$. W przypadku rzeczywistym
oszacowania mają co najwyżej dodatkowy czynnik $2$.

Niezależność od przedstawienia trzeba uzasadnić bez
zakładania z góry, że obie sumy definiują funkcjonały
na całej przestrzeni. Dla każdego $\lambda\in G'$ forma
skalarna $(f,g)\mapsto\lambda(b(f,g))$ jest osobno ciągła.
Aby zastosować twierdzenie o jądrze Schwartza w jego
wersji operatorowej, zauważmy, że forma ta wyznacza
ciągłe odwzorowanie $\mathcal D(\Omega)\to
\mathcal D'(\Lambda)_\beta$: dla ograniczonego
$B\subset\mathcal D(\Lambda)$ istnieje etap
$\mathcal D_L(\Lambda)$ zawierający $B$ jako zbiór
ograniczony, a łączna ciągłość formy na
$\mathcal D_K(\Omega)\times\mathcal D_L(\Lambda)$
daje dla każdego $K\Subset\Omega$ oszacowanie
$\sup_{g\in B}|\lambda(b(f,g))|\le C_B p_{K,m}(f)$
na $\mathcal D_K(\Omega)$. Ciągłość na wszystkich
etapach jest równoważna ciągłości na $\mathcal D(\Omega)$.
Twierdzenie o jądrze Schwartza
(Twierdzenie~\ref{thm:schwartz-kernel23}) daje
jednoznaczną dystrybucję $K_\lambda$, dla której
$\lambda(b(f,g))=\langle K_\lambda,f\otimes g\rangle$.
Dla każdego rozkładu tensorowego
$\Phi=\sum_n f_n\otimes g_n$ zbieżnego w przestrzeni
testowej, dla którego $\sum_n b(f_n,g_n)$ jest
zbieżny w $G$, ciągłość $\lambda$ i $K_\lambda$ daje
$\lambda(\sum_n b(f_n,g_n))=\langle K_\lambda,\Phi\rangle$.
Prawa strona nie zależy od rozkładu. Ponieważ przestrzeń sprzężona $G'$ rozdziela punkty
$G$ na mocy twierdzenia Hahna--Banacha, suma nie zależy
od wyboru rozwinięcia. Równość
$\lambda(B(\Phi))=\langle K_\lambda,\Phi\rangle$
dla wszystkich $\lambda\in G'$ dowodzi liniowości $B$.
Na tensorze elementarnym dopuszczamy jednoskładnikowe
przedstawienie, więc $B(f\otimes g)=b(f,g)$.
Jednoznaczność ciągłego przedłużenia wynika z gęstości
algebraicznych tensorów z Lematu~\ref{lem:tensor-density23}.
Z własności uniwersalnej granicy induktywnej otrzymujemy
ciągłość $B$ na całej $\mathcal D(\Omega\times\Lambda)$.

Odwzorowanie $(f,g)\mapsto f\otimes g$ o wartościach
w $\mathcal D(\Omega\times\Lambda)$ jest osobno ciągłe:
dla ustalonego $g$ o nośniku w $L$ jego ograniczenie
do $\mathcal D_K(\Omega)$ ma wartości w
$\mathcal D_{K\times L}(\Omega\times\Lambda)$, a
$p_{K\times L,m}(f\otimes g)\le p_{K,m}(f)p_{L,m}(g)$.
Stosujemy własność uniwersalną granicy induktywnej
w pierwszym czynniku. Drugi czynnik traktujemy analogicznie.
Zupełność $\mathcal D(\Omega\times\Lambda)$ wynika z
Twierdzenia~\ref{thm:D-complete18}.
Stwierdzenie~\ref{prop:i-vs-pi24} daje więc ciągłe
odwzorowanie liniowe
$J:\mathcal D(\Omega)\widehat\otimes_i\mathcal D(\Lambda)
\to\mathcal D(\Omega\times\Lambda)$.
Do zupełnego $G=\mathcal D(\Omega)\widehat\otimes_i
\mathcal D(\Lambda)$ stosujemy wykazany fakt z
kanonicznym osobno ciągłym odwzorowaniem dwuliniowym.
Otrzymujemy ciągłe $B:\mathcal D(\Omega\times\Lambda)\to G$.
Oba złożenia $JB$ i $BJ$ są identycznościami
na tensorach elementarnych, a więc na ich liniowych
powłokach. Gęstość z Lematu~\ref{lem:tensor-density23}
oraz gęstość iloczynu algebraicznego w jego uzupełnieniu
dają $JB=I$ i $BJ=I$ na całych przestrzeniach.
To jest żądany izomorfizm.
\end{proof}

\begin{przyklad}[Forma osobno ciągła, która nie jest łącznie ciągła]
\label{ex:separate-not-joint24}
Na $\mathcal D(\R)\times\mathcal D(\R)$ definiujemy
\[
 b(f,g)=\sum_{m=1}^\infty (-1)^m f^{(m)}(0)g(m).
\]
Suma jest skończona dla ustalonej funkcji $g$, ponieważ
$\supp g$ jest zwarty. Dla ustalonej $g$ funkcjonał
$f\mapsto b(f,g)$ jest więc ciągłą kombinacją skończenie
wielu pochodnych w zerze. Dla ustalonej $f$ otrzymujemy
lokalnie skończoną sumę dystrybucji punktowych
$\sum_m(-1)^mf^{(m)}(0)\delta_m$: na każdym zwartym
etapie $\mathcal D_L$ jej moduł jest ograniczony przez
$p_{L,0}(g)\sum_{m\in L\cap\N}|f^{(m)}(0)|$.
Własność uniwersalna LF zapewnia ciągłość funkcjonału. Forma $b$ jest osobno ciągła.

Przypuśćmy, że jest także łącznie ciągła. Wtedy istnieją
ciągłe seminormy $P,Q$ na $\mathcal D(\R)$ takie, że
$|b(f,g)|\le P(f)Q(g)$ dla wszystkich $f,g$.
Wybierzmy przedział zwarty $K\Subset\R$ zawierający
otoczenie zera. Ciągłość $P$ na $\mathcal D_K(\R)$
daje pewne $r\in\N_0$ i $C>0$ takie, że
$P(f)\le Cp_{K,r}(f)$ dla $f\in\mathcal D_K(\R)$.
Dla $m>r$ wybieramy $g_m\in\mathcal D(\R)$ o nośniku
w $(m-1/3,m+1/3)$ i $g_m(m)=1$. Wtedy
$|f^{(m)}(0)|\le CQ(g_m)p_{K,r}(f)$ dla wszystkich
$f\in\mathcal D_K(\R)$. Niech $\chi\in\mathcal D_K(\R)$
będzie równa $1$ w otoczeniu zera i niech
$f_N(x)=\chi(x)\cos(Nx-m\pi/2)$, gdzie $N\ge1$.
Jest to funkcja rzeczywista, więc działa nad obydwoma
ciałami. Wówczas $f_N^{(m)}(0)=N^m$, ponieważ
$\chi=1$ w otoczeniu zera. Z reguły Leibniza
dla $k\le r$ mamy
\[
 \sup_K|f_N^{(k)}|
 \le\sum_{j=0}^k\binom{k}{j}\sup_K|\chi^{(k-j)}|N^j
 \le C_{\chi,r}N^r.
\]
Po podzieleniu nierówności przez $N^r$ otrzymujemy
$N^{m-r}\le C Q(g_m) C_{\chi,r}$ dla wszystkich $N$,
co jest niemożliwe.

Forma $b$ odpowiada dystrybucji
$\sum_{m\ge1}\partial_x^m\delta_0\otimes\delta_m$
na $\R^2$; suma jest lokalnie skończona ze względu
na drugą zmienną. Jest to jawny przykład, dla którego
własność uniwersalna topologii $i$ jest silniejsza
niż własność uniwersalna topologii $\pi$.
\end{przyklad}

\begin{wniosek}[Ścisłość porównania dla dowolnych otwartych zbiorów euklidesowych]
\label{cor:dt-strict-i-pi-open}
Niech $d_1,d_2\ge1$ i niech $\Omega\subseteq\R^{d_1}$
oraz $\Lambda\subseteq\R^{d_2}$ będą niepustymi
zbiorami otwartymi. Na iloczynie algebraicznym
$\mathcal D(\Omega)\otimes\mathcal D(\Lambda)$
zachodzi $\varepsilon=\pi<i$, gdzie ostatnia
nierówność oznacza ścisłe porównanie topologii.
\end{wniosek}
\begin{proof}
Równość $\varepsilon=\pi$ wynika z nuklearności
$\mathcal D(\Omega)$ (Twierdzenia~\ref{thm:nuclear-D23}
i~\ref{thm:nuclear-pi-eps24}).
Ustalmy $x_0\in\Omega$ i wyczerpanie zwarte
$(L_m)_{m\ge1}$ zbioru $\Lambda$ z
Lematu~\ref{lem:exhaustion17}. Indukcyjnie wybieramy
różne punkty $y_m\in\Lambda\setminus L_m$.
Na każdym kroku możemy ominąć także skończony
zbiór wcześniej wybranych punktów: niepusty
otwarty zbiór $\Lambda$ nie jest zwarty, więc
nie jest sumą $L_m$ i skończonego zbioru.
Każdy zwarty podzbiór $\Lambda$ jest zawarty
w pewnym $L_M$, a zatem zawiera co najwyżej
skończenie wiele punktów $y_m$. Ponieważ każdy
punkt $\Lambda$ ma zwarte otoczenie zawarte
w $\Lambda$, rodzina tych punktów jest lokalnie
skończona i każdy $y_m$ ma otoczenie niezawierające
żadnego innego punktu tej rodziny.

Definiujemy formę dwuliniową
\[
 b(f,g)=\sum_{m\ge1}(-1)^m
 (\partial_1^m f)(x_0)g(y_m).
\]
Dla ustalonego $g$ suma jest skończona,
a każda ewaluacja pochodnej w $x_0$ jest
ciągła na $\mathcal D(\Omega)$: na etapie
zawierającym $x_0$ ogranicza ją seminorma
rzędu $m$, a na pozostałych etapach jest zerowa.
Dla ustalonego $f$ i zwartego $L\Subset\Lambda$
mamy oszacowanie
\[
 |b(f,g)|\le p_{L,0}(g)
 \sum_{m:y_m\in L}|(\partial_1^m f)(x_0)|
 \qquad(g\in\mathcal D_L(\Lambda)).
\]
Z własności uniwersalnej granicy induktywnej
(Twierdzenie~\ref{thm:LF-universal18}) wynika
osobna ciągłość $b$.

Gdyby forma była łącznie ciągła, istniałyby
ciągłe seminormy $P,Q$ z $|b(f,g)|\le P(f)Q(g)$
na mocy oszacowania z dowodu
Twierdzenia~\ref{thm:pi-universal24}.
Wybieramy zwartą kulę $K\Subset\Omega$
o środku $x_0$ i dodatnim promieniu. Istnieją
$C>0$ i $r\in\N_0$ z $P(f)\le Cp_{K,r}(f)$
dla $f\in\mathcal D_K(\Omega)$.
Dla ustalonego $m>r$ wybieramy
$g_m\in\mathcal D(\Lambda)$ z $g_m(y_m)=1$
i nośnikiem w otoczeniu $y_m$ niezawierającym
innych punktów rodziny. Istnienie takiej funkcji,
jak również odcięcia $\chi\in\mathcal D_K(\Omega)$
równego $1$ w otoczeniu $x_0$, wynika z
Lematu~\ref{lem:cutoffs17}.
Kładziemy
\[
 f_N(x)=\chi(x)\cos\bigl(N(x_1-(x_0)_1)-m\pi/2\bigr),
 \qquad N\ge1.
\]
Mamy $(\partial_1^m f_N)(x_0)=N^m$.
Reguła Leibniza, jak w poprzednim przykładzie,
daje $p_{K,r}(f_N)\le C_{\chi,r}N^r$:
pochodne czynnika trygonometrycznego względem
pozostałych współrzędnych znikają, a pochodne
względem pierwszej współrzędnej rzędu $j\le r$
mają moduł co najwyżej $N^j$.
Zatem $N^m=|b(f_N,g_m)|\le
CC_{\chi,r}Q(g_m)N^r$, co jest sprzeczne
z $m>r$.

Forma $b$ jest osobno ciągła, ale nie jest łącznie
ciągła. Z własności uniwersalnych
(Stwierdzenie~\ref{prop:i-vs-pi24} i
Twierdzenie~\ref{thm:pi-universal24})
wynika, że odpowiadający jej funkcjonał jest
ciągły dla $i$, ale nie dla $\pi$. Dowodzi to
ścisłości porównania topologii.
\end{proof}

\begin{uwaga}[Nie wolno utożsamiać trzech topologii]
\label{rem:D-not-projective24}
Na $\mathcal D(\Omega)\otimes\mathcal D(\Lambda)$
zachodzi $\varepsilon=\pi\le i$, ponieważ
$\mathcal D(\Omega)$ jest nuklearna
(Twierdzenie~\ref{thm:nuclear-D23} i
Twierdzenie~\ref{thm:nuclear-pi-eps24}).
Druga nierówność jest ostra dla wszystkich rozważanych
niepustych zbiorów otwartych, na mocy
Wniosku~\ref{cor:dt-strict-i-pi-open}.
W szczególności
$\mathcal D(\Omega)\widehat\otimes_\pi\mathcal D(\Lambda)$
nie jest \emph{kanonicznie izomorficzna topologicznie}
z $\mathcal D(\Omega\times\Lambda)$ przez odwzorowanie
$f\otimes g\mapsto f(x)g(y)$: gdyby taki izomorfizm
istniał, forma z przykładu byłaby ograniczeniem
ciągłego funkcjonału na uzupełnieniu rzutowym,
a więc byłaby łącznie ciągła. Nuklearność nie wymusza
zatem równości $\pi=i$. Jest to powód, dla którego
w twierdzeniu o jądrze Schwartza posługiwaliśmy się
formami osobno ciągłymi, a nie wyłącznie łącznie ciągłymi.
\end{uwaga}

\subsection{Powiązanie z teorią jąder i dalsza geometria tensorowa}

\begin{wniosek}[Interpretacja jądra dystrybucyjnego]
\label{cor:kernel-i24}
Niech $d_1,d_2\ge1$ będą liczbami całkowitymi,
a $\Omega\subseteq\R^{d_1}$ i $\Lambda\subseteq\R^{d_2}$
niepustymi zbiorami otwartymi. Niech
$b:\mathcal D(\Omega)\times\mathcal D(\Lambda)\to\K$
będzie formą dwuliniową. Następujące warunki są równoważne:
\begin{enumerate}[label=\textup{(\roman*)},leftmargin=*]
\item $b$ jest osobno ciągła;
\item istnieje ciągły operator liniowy
$A:\mathcal D(\Omega)\to\mathcal D'(\Lambda)_\beta$
taki, że $b(f,g)=\langle Af,g\rangle$ dla wszystkich
$f\in\mathcal D(\Omega)$ i $g\in\mathcal D(\Lambda)$;
\item istnieje dystrybucja
$K\in\mathcal D'(\Omega\times\Lambda)$ taka, że
$b(f,g)=\langle K,f\otimes g\rangle$ dla wszystkich
$f\in\mathcal D(\Omega)$ i $g\in\mathcal D(\Lambda)$.
\end{enumerate}
Operator $A$ i dystrybucja $K$ są wyznaczone jednoznacznie.
Parowanie z dystrybucją jest liniowe w funkcji testowej,
również nad $\C$; $\beta$ oznacza silną topologię
przestrzeni sprzężonej.
\end{wniosek}
\begin{proof}
\emph{(i)$\Rightarrow$(iii).}
Przez własność uniwersalną iloczynu induktywnego
(Stwierdzenie~\ref{prop:i-vs-pi24}) forma $b$
wyznacza ciągły funkcjonał na
$\mathcal D(\Omega)\otimes_i\mathcal D(\Lambda)$.
Przedłużamy go jednoznacznie na uzupełnienie (ciało
$\K$ jest zupełne) i przenosimy na
$\mathcal D(\Omega\times\Lambda)$ przez izomorfizm
z Twierdzenia~\ref{thm:D-i-product24}.
Jednoznaczność dystrybucji wynika z gęstości
liniowej powłoki tensorów elementarnych.

\emph{(iii)$\Rightarrow$(ii).}
Stosujemy odwrotną część twierdzenia o jądrze
Schwartza (Twierdzenie~\ref{thm:schwartz-kernel23}):
$A_Kf$ jest dystrybucją zadaną wzorem
$\langle A_Kf,g\rangle=\langle K,f\otimes g\rangle$,
a $A_K$ jest ciągły dla silnej topologii.
Operator jest jednoznaczny, gdyż dystrybucje
są funkcjonałami na całej przestrzeni testowej.

\emph{(ii)$\Rightarrow$(i).}
Dla ustalonego $f$ ciągłość względem $g$ wynika
z tego, że $Af$ jest dystrybucją. Dla ustalonego
$g$ ewaluacja $T\mapsto\langle T,g\rangle$ jest
ciągła w silnej topologii, ponieważ $\{g\}$
jest zbiorem ograniczonym. Składając ją z $A$,
otrzymujemy ciągłość względem $f$.

Wszystkie trzy przyporządkowania są liniowe
i wzajemnie zgodne. Otrzymujemy w szczególności
izomorfizm algebraiczny
\[
 \bigl(\mathcal D(\Omega)\widehat\otimes_i
 \mathcal D(\Lambda)\bigr)'
 \simeq\mathcal D'(\Omega\times\Lambda).
\]
Nie przypisujemy tu przestrzeni form dwuliniowych
ani przestrzeni operatorów dodatkowej topologii:
sama bijekcja liniowa nie dowodzi zgodności
dowolnie wybranych topologii na tych przestrzeniach.
\end{proof}

\begin{przyklad}[Operator o jądrze na przekątnej]
\label{ex:diag-tensor24}
Niech $d\ge1$, a $\Omega\subseteq\R^d$ będzie
niepustym zbiorem otwartym. Dystrybucja
$\delta_\Delta$ na przekątnej $\Delta=\{(x,x):x\in\Omega\}$
z Przykładu~\ref{ex:diagonal-kernel23} zadaje formę
\[
 b(f,g)=\int_\Omega f(x)g(x)\,\mathrm d\lambda_d(x),
 \qquad f,g\in\mathcal D(\Omega).
\]
Odpowiadający jej operator to regularne włożenie
$f\mapsto T_f$ do $\mathcal D'(\Omega)_\beta$,
a nie operator o wartościach w funkcjach zadany
przez całkowanie względem zwykłego jądra funkcyjnego.
W szczególności $\supp\delta_\Delta=\Delta$,
a $\delta_\Delta$ nie jest dystrybucją regularną.

W tym konkretnym przykładzie forma jest nawet
łącznie ciągła: seminorma
$p(f)=\|f\|_{L^2(\Omega)}$ jest ciągła na
$\mathcal D(\Omega)$, ponieważ na każdym zwartym
etapie zachodzi
$p(f)\le\lambda_d(K)^{1/2}p_{K,0}(f)$.
Stosujemy własność uniwersalną granicy induktywnej.
Nierówność Cauchy'ego--Schwarza daje
$|b(f,g)|\le p(f)p(g)$ także nad $\C$
(bez wprowadzania sprzężenia do definicji formy).
Zatem $b$ wyznacza ciągły funkcjonał zarówno na
uzupełnionym iloczynie rzutowym, jak i induktywnym.
Nie wynika stąd równość tych topologii: przykład
\ref{ex:separate-not-joint24} rozróżnia ich
przestrzenie ciągłych funkcjonałów.
Osobliwość jądra nie oznacza braku łącznej
ciągłości odpowiadającej mu formy.
\end{przyklad}

\begin{uwaga}[Przyszłe rozszerzenia]
\label{rem:tensor-next24}
Topologiczne produkty tensorowe pojawią się ponownie
przy operatorach zwartych i nuklearnych, własności
aproksymacji oraz algebrach operatorowych.
W szczególności normy $\pi$ i $\varepsilon$ są dwiema
różnymi normami tensorowymi już na iloczynach dwóch
nieskończenie wymiarowych przestrzeni Hilberta.
Nie rozwijamy tu zastosowań do rozwiązywania równań
różniczkowych, które należą do odrębnego skryptu.
\end{uwaga}

% Materiał z AF 25; fragment do rozdziału o iloczynie induktywnym.
\subsection{Identyfikacje przestrzeni sprzężonych do iloczynów tensorowych}

\begin{stwierdzenie}[Zestawienie kanonicznych odwzorowań]
\label{prop:summary25}
Dla każdej liczby całkowitej $d\ge1$ i ustalonego
ciała $\K\in\{\R,\C\}$ zachodzi schemat
\[
 \begin{array}{ccccc}
 \mathcal D(\R^d)&\longrightarrow&\mathcal S(\R^d)
 &\longrightarrow&\mathcal E(\R^d)\\
 \mathcal D'(\R^d)&\longleftarrow&\mathcal S'(\R^d)
 &\longleftarrow&\mathcal E'(\R^d).
 \end{array}
\]
W pierwszym wierszu strzałki oznaczają ciągłe
różnowartościowe odwzorowania liniowe o gęstych obrazach,
w drugim zaś ich transpozycje, ciągłe i różnowartościowe
zarówno w topologiach słabych-*, jak i silnych.
Żadna ze strzałek nie jest izomorfizmem topologicznym
na obraz z topologią podprzestrzeni.

Dla dodatnich całkowitych $d_1,d_2$ oraz niepustych
zbiorów otwartych $\Omega\subseteq\R^{d_1}$,
$\Lambda\subseteq\R^{d_2}$ istnieją kanoniczne
izomorfizmy liniowo-topologiczne
\[
\begin{aligned}
 \mathcal E(\Omega)\widehat\otimes_\pi\mathcal E(\Lambda)
 &\simeq\mathcal E(\Omega\times\Lambda),\\
 \mathcal S(\R^{d_1})\widehat\otimes_\pi\mathcal S(\R^{d_2})
 &\simeq\mathcal S(\R^{d_1+d_2}),\\
 \mathcal D(\Omega)\widehat\otimes_i\mathcal D(\Lambda)
 &\simeq\mathcal D(\Omega\times\Lambda).
\end{aligned}
\]
Każdy izomorfizm przedłuża odwzorowanie
$f\otimes g\mapsto ((x,y)\mapsto f(x)g(y))$.
W pierwszych dwóch wzorach można zastąpić $\pi$
przez $\varepsilon$. W trzecim wzorze nie wolno
ogólnie zastąpić $i$ przez $\pi$, choć na algebraicznym
iloczynie tensorowym przestrzeni $\mathcal D$
zachodzi $\varepsilon=\pi$.
\end{stwierdzenie}
\begin{proof}
Pierwszy wiersz wynika z Twierdzenia~
\ref{thm:zanurzenia25}, drugi z Wniosku~
\ref{cor:dual-lancuch25}; różnicę topologii
na obrazach dualnych wykazano w Przykładzie~
\ref{ex:escaping-deltas25}. Pierwsze dwa izomorfizmy
oraz dopuszczalność $\varepsilon$ wynikają z
Twierdzenia~\ref{thm:smooth-schwartz-tensor24}.
Trzeci wynika z Twierdzenia~\ref{thm:D-i-product24};
równość $\varepsilon=\pi$ wynika z nuklearności
$\mathcal D$ (Twierdzenia~\ref{thm:nuclear-D23}
i~\ref{thm:nuclear-pi-eps24}). Nie można zastąpić
$i$ przez $\pi$ już dla $\Omega=\Lambda=\R$,
na mocy Uwagi~\ref{rem:D-not-projective24}.
\end{proof}

\begin{uwaga}[Odczytywanie schematu po stronie dualnej]
\label{rem:kernel-summary25}
Jeżeli $J:X\to Y$ jest izomorfizmem liniowo-topologicznym,
to jego transpozycja $J':Y'\to X'$ jest także
izomorfizmem dla obu topologii dualnych.
Istotnie, $J'$ jest ciągłe przez równość seminorm
silnych i słabych-* z dowodu Lematu~
\ref{lem:transpozycja25}, a jego odwrotnością jest
$(J^{-1})'$, ciągłe z tego samego powodu.
Dlatego transpozycje trzech izomorfizmów tensorowych
identyfikują przestrzenie dualne produktów
z odpowiednimi przestrzeniami dystrybucji na
iloczynie, z zachowaniem topologii słabej-* i silnej.

Konkretnie, w każdym przypadku funkcjonał $K$
na przestrzeni funkcji na iloczynie odpowiada
formie dwuliniowej
$b_K(f,g)=\langle K,f\otimes g\rangle$.
Dla $\mathcal E$ i $\mathcal S$ jest to bijekcja
liniowa między dystrybucjami z odpowiedniej
przestrzeni dualnej na iloczynie a \emph{łącznie}
ciągłymi formami dwuliniowymi na parze czynników.
Wynika to z własności uniwersalnej $\pi$
(Twierdzenie~\ref{thm:pi-universal24}) i jednoznacznego
przedłużenia funkcjonału na uzupełnienie.
Dla $\mathcal D$ analogiczna bijekcja dotyczy
\emph{osobno} ciągłych form, zgodnie z
Wnioskiem~\ref{cor:kernel-i24}.
W obu przypadkach jednoznaczność jądra wynika
z gęstości liniowej powłoki tensorów elementarnych
w odpowiednim uzupełnieniu.

Nie przypisujemy tutaj przestrzeni form dwuliniowych
dodatkowej topologii. Stwierdzenie o izomorfizmach
topologicznych dotyczy dualnych przestrzeni
uzupełnionych iloczynów tensorowych, nie dowolnie
wybranych topologii na przestrzeniach form.
\end{uwaga}

\begin{uwaga}[Zakres dalszych zastosowań]
\label{rem:future25}
Identyfikacje przestrzeni sprzężonych i twierdzenie o jądrze
opisują reprezentacje ciągłych operatorów na przestrzeniach
funkcji testowych. Zastosowania do konkretnych równań
różniczkowych rozwijamy w skrypcie RRC, a ogólną teorię
operatorów w analizie funkcjonalnej. W dalszej części
niniejszego skryptu przechodzimy do analizy harmonicznej.
\end{uwaga}

\noindent\emph{Powiązania tematyczne w Bibliotece matematycznej:} \atlasref{KAT-09::cat:ex:tensor-rep}{KAT-09: własność uniwersalna iloczynu tensorowego w języku funktorów}.

\setcounter{section}{20}
% Rozdział aktywny po restrukturyzacji DT.
\section{Rozwiązania fundamentalne operatorów różniczkowych}
\providecommand{\atlasref}[2]{#2}

\subsection{Definicja i związek ze splotem}

\begin{definicja}[Rozwiązanie fundamentalne]
\label{def:dt-fundamental-solution}
Niech $P(D)=\sum_{|\alpha|\le m}a_\alpha\partial^\alpha$
będzie operatorem różniczkowym o stałych współczynnikach
$a_\alpha\in\K$ na $\R^d$. Przez $D$ w tym zapisie
rozumiemy zwykłe pochodne cząstkowe, bez czynnika $i$.
Dystrybucję $E\in\mathcal D'(\R^d)$ nazywamy
\emph{rozwiązaniem fundamentalnym} operatora $P(D)$,
jeżeli
\[
 P(D)E=\delta_0.
\]
Równość rozumiemy w sensie pochodnych dystrybucyjnych
(Definicja~\ref{def:dist-deriv20}).
\end{definicja}

\begin{stwierdzenie}[Rozwiązanie równania z prawą stroną o zwartym nośniku]
\label{prop:dt-fundamental-convolution}
Niech $E$ będzie rozwiązaniem fundamentalnym $P(D)$.
Dla każdego $S\in\mathcal E'(\R^d)$ dystrybucja
$u=E*S$ jest rozwiązaniem równania $P(D)u=S$.
Jeżeli $S=T_f$ dla $f\in\mathcal D(\R^d)$, to $u$
jest reprezentowana przez funkcję gładką
\[
 u(x)=\langle E_y,f(x-y)\rangle,
\]
która spełnia równanie $P(D)u=f$ klasycznie.
Wszystkie dystrybucyjne rozwiązania dla ustalonego $S$
mają postać $E*S+v$, gdzie $P(D)v=0$.
W szczególności każde dwa rozwiązania fundamentalne
różnią się o rozwiązanie równania jednorodnego.
\end{stwierdzenie}
\begin{proof}
Splot jest określony, ponieważ $S$ ma zwarty nośnik
(Definicja~\ref{def:convolution21} i
Twierdzenie~\ref{thm:convolution-algebra21}).
Ze stałości współczynników i wzorów na pochodne
splotu (Twierdzenie~\ref{thm:convolution-derivative21})
otrzymujemy
\[
 P(D)(E*S)=(P(D)E)*S=\delta_0*S=S,
\]
gdzie ostatnia równość wynika z
Przykładu~\ref{ex:delta-convolution21}.
Dla $S=T_f$ Twierdzenie~\ref{thm:smoothing21}
daje wskazaną funkcję gładką. Zgodność pochodnych
klasycznych i dystrybucyjnych
(Stwierdzenie~\ref{prop:classical-weak20})
oraz różnowartościowość przyporządkowania funkcji
ciągłej odpowiadającej jej dystrybucji regularnej
dają równanie punktowe. Ostatnie twierdzenia
wynikają z liniowości $P(D)$: różnica dwóch
rozwiązań należy do jego jądra, a dodanie
elementu jądra nie zmienia prawej strony.
\end{proof}

\begin{uwaga}
Powyższe stwierdzenie nie zakłada istnienia
rozwiązania fundamentalnego dla każdego operatora.
Nie uzasadnia też splotu $E*S$ dla dowolnej
dystrybucji $S$ bez dodatkowych założeń
(por.~Uwaga~\ref{rem:convolution-limits21}).
Rozwiązanie fundamentalne nie jest automatycznie
jądrem funkcji Greena dla zagadnienia brzegowego:
nie narzucono tutaj żadnych warunków brzegowych.
\end{uwaga}

\subsection{Warunek istnienia rozwiązania o zwartym nośniku}

\begin{twierdzenie}[Rozwiązanie fundamentalne o zwartym nośniku]
\label{thm:dt-fundamental-compact}
Operator różniczkowy o stałych współczynnikach $P(D)$
na $\R^d$ ma rozwiązanie fundamentalne o zwartym
nośniku wtedy i tylko wtedy, gdy jest mnożeniem
przez niezerową stałą $a$. Wówczas jedynym
rozwiązaniem fundamentalnym jest $a^{-1}\delta_0$.
\end{twierdzenie}
\begin{proof}
Przypadek niezerowej stałej jest natychmiastowy,
a operator zerowy nie może mieć rozwiązania fundamentalnego.
Załóżmy zatem, że wielomian symbolu
$P(z)=\sum_\alpha a_\alpha z^\alpha$ ma stopień $m\ge1$.
Jego część jednorodna najwyższego stopnia $P_m$
nie jest wielomianem zerowym. Istnieje $v\in\R^d$
z $P_m(v)\ne0$: wielomian znikający na całym
$\R^d$ jest zerowy, co wynika przez indukcję
względem liczby zmiennych z jednowymiarowego
twierdzenia o liczbie pierwiastków wielomianu.
Wielomian $t\mapsto P(tv)$ ma stopień $m$, więc
z zasadniczego twierdzenia algebry istnieje
$t_0\in\C$ takie, że $P(t_0v)=0$.
Oznaczmy $z_0=t_0v$.

Jeżeli $E$ ma zwarty nośnik i $P(D)E=\delta_0$,
to równość zachodzi również w $\mathcal E'(\R^d)$
na mocy Twierdzenia~\ref{thm:compact-support-E25}.
Nad $\R$ przedłużamy działanie dystrybucji
zespolonoliniowo na funkcje zespolone.
Możemy więc zastosować obie strony do gładkiej
funkcji $h(x)=\exp(-z_0\cdot x)$, która nie musi
mieć zwartego nośnika. Z definicji pochodnej
i przedłużenia działania na funkcje gładkie
(Lemat~\ref{lem:compact-action21}) otrzymujemy
\[
 1=h(0)=\langle P(D)E,h\rangle
 =\langle E,P(-D)h\rangle
 =P(z_0)\langle E,h\rangle=0.
\]
Ostatnią równość poprzedza obliczenie
$\partial^\alpha h=(-z_0)^\alpha h$.
Sprzeczność dowodzi tezy.
\end{proof}

\subsection{Równanie symbolu w przestrzeni dystrybucji temperowanych}

\begin{stwierdzenie}[Charakterystyka przez transformację Fouriera]
\label{prop:dt-fundamental-symbol}
Przy normalizacji $c_d=(2\pi)^{-d/2}$ dystrybucja
$E\in\mathcal S'(\R^d)$ jest rozwiązaniem fundamentalnym
$P(D)$ wtedy i tylko wtedy, gdy
\[
 P(i\xi)\widehat E=c_d T_1.
\]
\end{stwierdzenie}
\begin{proof}
Wzory na transformację pochodnych dystrybucyjnych
i delty Diraca (Przykład~\ref{ex:fourier-delta22})
dają $\mathcal F(P(D)E)=P(i\xi)\widehat E$
oraz $\mathcal F\delta_0=c_d T_1$.
Implikacja odwrotna wynika z różnowartościowości
transformacji Fouriera na $\mathcal S'$
(Twierdzenie~\ref{thm:fourier-auto-tempered22}).
Równość w $\mathcal D'$ jest równoważna równości
w $\mathcal S'$ dla tych dystrybucji, przez
różnowartościowość ograniczenia do $\mathcal D$
(Stwierdzenie~\ref{prop:schwartz-embeddings22}).
\end{proof}

\begin{uwaga}
Równanie symbolu jest równaniem mnożenia dystrybucji
przez wielomian, a nie zwykłym dzieleniem funkcji.
Gdy symbol ma zera, funkcja $1/P(i\xi)$ może
nie być lokalnie całkowalna. Charakterystyka
nie stanowi sama w sobie dowodu istnienia
rozwiązania fundamentalnego dla każdego niezerowego
operatora.
\end{uwaga}

\subsection{Rozwiązanie fundamentalne operatora Laplace'a}
\begin{twierdzenie}[Rozwiązanie fundamentalne operatora $-\Delta$]
\label{thm:dt-laplace}
Niech $d\ge1$ i niech $\omega_d$ oznacza pole powierzchni
sfery jednostkowej $S^{d-1}\subset\R^d$ dla $d\ge2$.
Funkcja lokalnie całkowalna
\begin{equation}\label{eq:dt-laplace-fundamental}
 \Phi_d(x)=\begin{cases}
 -|x|/2,&d=1,\\
 (2\pi)^{-1}\log(1/|x|),&d=2,\\
 ((d-2)\omega_d)^{-1}|x|^{2-d},&d\ge3,
 \end{cases}
\end{equation}
spełnia $-\Delta T_{\Phi_d}=\delta_0$ na $\R^d$.
\end{twierdzenie}
\begin{proof}
Wartość $\Phi_d(0)$ w wymiarach $d\ge2$ ustalamy
dowolnie, na przykład równą $0$.
W wymiarze jeden $\Phi_1$ jest lokalnie całkowalna
i jej pierwsza pochodna dystrybucyjna jest funkcją
$-\tfrac12\operatorname{sgn}x$, ponieważ funkcja jest ciągła
i kawałkami liniowa. Skok tej pochodnej w zerze wynosi $-1$.
Z Twierdzenia~\ref{thm:jump20} wynika zatem
$\partial_x^2T_{\Phi_1}=-\delta_0$.

Dla $d\ge3$ całkowanie w zmiennych biegunowych daje
dla $r>0$
\[
 \int_{B_r(0)}|\Phi_d(x)|\,\mathrm d\lambda_d(x)
 =\frac1{d-2}\int_0^r s\,\mathrm d\lambda_1(s)
 =\frac{r^2}{2(d-2)}.
\]
Dla $d=2$ mamy $\omega_2=2\pi$ i dla $0<r<1$
\[
 \begin{aligned}
 \int_{B_r(0)}|\Phi_2(x)|\,\mathrm d\lambda_2(x)
 &=\int_0^r s\log(1/s)\,\mathrm d\lambda_1(s)\\
 &=\frac{r^2}{2}\log(1/r)+\frac{r^2}{4}.
 \end{aligned}
\]
Granica dolnego wyrazu pierwotnej jest zerowa,
gdyż $s^2\log(1/s)\to0$ dla $s\downarrow0$.
W obu przypadkach całka dąży do zera przy $r\downarrow0$.
Z gładkości poza zerem wynika więc lokalna
całkowalność na całej $\R^d$.

Niech $F(r)$ oznacza profil radialny $\Phi_d$ dla $r>0$.
Bezpośrednie różniczkowanie obu wzorów daje
\[
 F'(r)=-\frac1{\omega_dr^{d-1}},\qquad
 F''(r)=\frac{d-1}{\omega_dr^d}.
\]
Dla $x\ne0$, $r=|x|$,
$\partial_jr=x_j/r$ oraz
$\partial_j^2r=1/r-x_j^2/r^3$.
Z reguły łańcuchowej otrzymujemy
\[
 \partial_j^2\Phi_d(x)
 =F''(r)\frac{x_j^2}{r^2}
 +F'(r)\left(\frac1r-\frac{x_j^2}{r^3}\right).
\]
Sumując po $j$ i korzystając z $\sum_j x_j^2=r^2$,
dostajemy
\[
 \Delta\Phi_d(x)
 =F''(r)+\frac{d-1}{r}F'(r)=0.
\]
Dla $\varphi\in\mathcal D(\R^d)$ wybieramy $R>0$
tak, aby $\supp\varphi\Subset B_R(0)$.
Dla $0<\varepsilon<R$ stosujemy drugą tożsamość
Greena (wynikającą z twierdzenia o dywergencji
dla pola $\Phi_d\nabla\varphi-\varphi\nabla\Phi_d$) na pierścieniu
$A_{\varepsilon,R}=B_R(0)\setminus\overline{B_\varepsilon(0)}$
(do części rzeczywistej i urojonej testu osobno).
Ponieważ $\Delta\Phi_d=0$ na pierścieniu,
otrzymujemy
\[
 -\int_{A_{\varepsilon,R}}\Phi_d\Delta\varphi
 \,\mathrm d\lambda_d
 =\int_{\partial A_{\varepsilon,R}}
 (\varphi\partial_\nu\Phi_d-\Phi_d\partial_\nu\varphi)
 \,\mathrm d\mathcal H^{d-1}.
\]
Na zewnętrznej sferze znikają test i jego pochodne.
Na wewnętrznej sferze $\nu=-x/|x|$, zatem
$\partial_\nu\Phi_d=-F'(\varepsilon)
=1/(\omega_d\varepsilon^{d-1})$.
W konsekwencji
\[
 -\int_{A_{\varepsilon,R}}\Phi_d(x)\Delta\varphi(x)
 \,\mathrm d\lambda_d(x)
 =\int_{\partial B_\varepsilon}\varphi(x)
 \frac{1}{\omega_d\varepsilon^{d-1}}\,\mathrm d\mathcal H^{d-1}(x)
 -\int_{\partial B_\varepsilon}\Phi_d(x)
 \partial_\nu\varphi(x)\,\mathrm d\mathcal H^{d-1}(x).
\]
Miara powierzchniowa sfery o promieniu $\varepsilon$
wynosi $\omega_d\varepsilon^{d-1}$, przez jednorodność
miary Hausdorffa (por. dowód Lematu~\ref{lem:dt-fourier-sphere-kernel}). Po odjęciu
$\varphi(0)$ od pierwszej całki moduł różnicy
jest więc co najwyżej
$\sup_{|x|=\varepsilon}|\varphi(x)-\varphi(0)|$,
co dąży do zera z ciągłości testu.
Moduł drugiej całki nie przekracza
\[
 \|\nabla\varphi\|_\infty\,
 |F(\varepsilon)|\,\omega_d\varepsilon^{d-1}.
\]
Dla $d\ge3$ iloczyn bez normy gradientu wynosi
$\varepsilon/(d-2)$, dla $d=2$ i $\varepsilon<1$
wynosi $\varepsilon\log(1/\varepsilon)$.
Obie wielkości dążą do zera.
Ponieważ sfery mają miarę Lebesgue'a zero
i test znika poza $B_R(0)$, w całce po $\R^d$
pozostaje uwzględnić tylko wkład z $B_\varepsilon(0)$.
Jego moduł jest ograniczony przez
$\|\Delta\varphi\|_\infty
\int_{B_\varepsilon(0)}|\Phi_d(x)|\,\mathrm d\lambda_d(x)$,
co dąży do zera z wcześniejszych obliczeń.
Ze znaku $(-1)^2=1$ przy drugiej pochodnej
dystrybucyjnej otrzymujemy
\[
 \langle-\Delta T_{\Phi_d},\varphi\rangle
 =-\int_{\R^d}\Phi_d(x)\Delta\varphi(x)
 \,\mathrm d\lambda_d(x)=\varphi(0).
\]
Jest to żądana równość.
\end{proof}

\begin{uwaga}[Powiązanie z równaniami różniczkowymi cząstkowymi]
\noindent W skrypcie RRC rozwijamy zastosowania do
zagadnień brzegowych i ewolucyjnych:
\atlasref{PDE-03::sec:pde3-green}{funkcje Greena i tożsamości Greena},
\atlasref{PDE-07::thm:pde7-heat}{jądro cieplne jako rozwiązanie fundamentalne}
oraz
\atlasref{PDE-07::thm:pde7-weyl}{lemat Weyla}.
Tutaj koncentrujemy się na rachunku dystrybucji,
splotach, symbolach i ograniczeniach wynikających
z nośnika.
\end{uwaga}

\setcounter{part}{3}
\part{Szeregi Fouriera}
\setcounter{section}{23}

\section{Szereg Fouriera: współczynniki, ortogonalność i zbieżność}
\label{sec:dt-series-fourier}
\providecommand{\atlasref}[2]{#2}

Por. \atlasref{PDE-06::thm:pde6-fejer}{RRC-06: średnie Fejéra i szeregi Fouriera}.


Przez $\mathbb T=\R/(2\pi\mathbb Z)$ rozumiemy okrąg z
współrzędną $x\in[-\pi,\pi]$ i utożsamionymi końcami.
Funkcję z $L^1(-\pi,\pi)$ utożsamiamy z jej
$2\pi$-okresowym przedłużeniem, określonym z dokładnością
do zbioru miary zero; wartości na końcach przedziału
nie wpływają na całki. Dla takiej funkcji definiujemy
$\widehat f(k)=(2\pi)^{-1}\int_{-\pi}^{\pi}f(x)e^{-ikx}\,\mathrm d\lambda_1(x)$,
$k\in\mathbb Z$, oraz sumę częściową
$S_Nf(x)=\sum_{|k|\le N}\widehat f(k)e^{ikx}$.
W tym rozdziale $L^2(\mathbb T)$ oznacza zespoloną przestrzeń Hilberta
z iloczynem skalarnym
$\langle f,g\rangle=(2\pi)^{-1}\int_{-\pi}^{\pi}f(x)\overline{g(x)}\,\mathrm d\lambda_1(x)$.
Pojęcie całki Lebesgue'a i podstawowe własności $L^p$
traktujemy jako znane. Na okręgu używamy miary Haara
unormowanej do masy $1$, czyli $\mathrm d\lambda_1/(2\pi)$.
Na torusie $\mathbb T^d$ jej odpowiednikiem jest
$\mathrm d\lambda_d/(2\pi)^d$ na kostce fundamentalnej;
współczynniki Fouriera zawierają zatem czynnik $(2\pi)^{-d}$.
Norma $L^2$ względem nieznormalizowanej miary jest
$(2\pi)^{d/2}$ razy większa. Nie należy utożsamiać tej
normalizacji okresowej z unitarną normalizacją całki
Fouriera na $\R^d$ z definicji \ref{def:fourier-schwartz17}.
Ponieważ miara unormowana na
$\mathbb T$ ma masę $1$, nierówność Cauchy'ego--Schwarza
daje $(2\pi)^{-1}\|f\|_{L^1(-\pi,\pi)}\leq\|f\|_2$,
więc współczynniki Fouriera są określone dla każdego
$f\in L^2(\mathbb T)$.

\begin{lemat}[Ortogonalność i nierówność Bessela]
\label{lem:dt-series-bessel}
Niech $f\in L^2(\mathbb T)$ oraz $N\geq0$ będzie całkowite. Wówczas
$\langle e^{ikx},e^{i\ell x}\rangle=\delta_{k\ell}$,
$\|S_Nf\|_2^2=\sum_{|k|\le N}|\widehat f(k)|^2$ i
$\sum_{|k|\le N}|\widehat f(k)|^2\le\|f\|_2^2$.
\end{lemat}
\begin{proof}
Dla $k=\ell$ całka z jedynki względem miary
unormowanej wynosi $1$. Jeśli $k\ne\ell$, to
\[
 \frac1{2\pi}\int_{-\pi}^{\pi}e^{i(k-\ell)x}\,
 \mathrm d\lambda_1(x)
 =\frac{e^{i(k-\ell)\pi}-e^{-i(k-\ell)\pi}}
 {2\pi i(k-\ell)}=0,
\]
ponieważ $k-\ell$ jest całkowite. Otrzymujemy
ortonormalność. Rozwijając kwadrat normy skończonej
sumy $S_Nf$ i korzystając z tej ortonormalności,
dostajemy $\|S_Nf\|_2^2=\sum_{|k|\le N}|\widehat f(k)|^2$.
Ponadto dla $|\ell|\leq N$
\[
 \langle f-S_Nf,e^{i\ell x}\rangle
 =\widehat f(\ell)-\sum_{|k|\leq N}
 \widehat f(k)\delta_{k\ell}=0.
\]
Zatem $S_N$ jest rzutem ortogonalnym na przestrzeń
rozpiętą przez te wykładniki. W szczególności
$\langle f-S_Nf,S_Nf\rangle=0$, skąd
\[
 \|f\|_2^2=\|S_Nf\|_2^2+\|f-S_Nf\|_2^2.
\]
Nieujemność ostatniego składnika daje nierówność Bessela.
\end{proof}

Nie można jednak wywnioskować z ortogonalności, że $S_Nf\to f$.
Do tego potrzebujemy zupełności układu wykładniczego.


\setcounter{section}{25}
% Rozdział aktywny po restrukturyzacji DT.
\section{Średnie Fouriera, jądro Fejéra i tożsamość Parsevala}
\providecommand{\atlasref}[2]{#2}

\subsection{Dodatnie aproksymacje jedności na okręgu}

Niech $\mathrm dm(t)=(2\pi)^{-1}\mathrm d\lambda_1(t)$
oznacza unormowaną miarę Haara na $\mathbb T$.
Splot okresowy z funkcją $K\in L^1(\mathbb T)$
zapisujemy jako
\[
 (K*f)(x)=\int_{\mathbb T}K(t)f(x-t)\,\mathrm dm(t).
\]
Dla $f\in L^1$ całka istnieje prawie wszędzie:
z twierdzenia Tonellego i niezmienniczości miary
względem przesunięć całka podwójna z modułu
jest równa $\|K\|_1\|f\|_1$.
Jeżeli $K$ i $f$ są ciągłe, całka istnieje
w każdym punkcie i określa funkcję ciągłą
przez jednostajną ciągłość $f$.

\begin{definicja}[Dodatnia aproksymacja jedności]
\label{def:dt-periodic-approximate-identity}
Rodzinę $(K_\nu)$ nieujemnych funkcji z
$L^1(\mathbb T)$ nazywamy dodatnią aproksymacją
jedności, jeśli $\int_{\mathbb T} K_\nu\,\mathrm dm=1$
oraz dla każdego $0<\delta<\pi$ zachodzi
\[
 \int_{\delta\le |t|\le\pi}K_\nu(t)\,\mathrm dm(t)
 \longrightarrow0.
\]
Parametr $\nu$ może na przykład dążyć do
nieskończoności albo do brzegu przedziału parametrów.
\end{definicja}

\begin{twierdzenie}[Aproksymacja jednostajna i w normie $L^p$]
\label{thm:dt-periodic-approximate-identity}
Jeśli $(K_\nu)$ jest dodatnią aproksymacją jedności,
to $K_\nu*f\to f$ jednostajnie dla $f\in C(\mathbb T)$,
a dla $f\in L^p(\mathbb T)$, $1\le p<\infty$,
zbieżność zachodzi w normie $L^p$.
Ponadto $\|K_\nu*f\|_p\le\|f\|_p$
dla $1\le p\le\infty$.
\end{twierdzenie}
\begin{proof}
Dla skończonego $p$ nierówność Jensena względem
miary probabilistycznej $K_\nu\,\mathrm dm$ daje
\[
 |(K_\nu*f)(x)|^p
 \le\int_{\mathbb T}K_\nu(t)|f(x-t)|^p\,\mathrm dm(t).
\]
Całkując i stosując twierdzenie Tonellego oraz
niezmienniczość miary, otrzymujemy kontraktywność.
Dla $p=\infty$ wynika ona wprost z ograniczenia
modułu funkcji przez jej kres istotny.
Dla $f$ ciągłej, po odjęciu $f(x)$ pod całką,
jednostajna ciągłość daje dla dowolnego
$\varepsilon>0$ takie $\delta>0$, że
\[
 \|K_\nu*f-f\|_\infty
 \le\varepsilon+2\|f\|_\infty
 \int_{\delta\le|t|\le\pi}K_\nu(t)\,\mathrm dm(t).
\]
Drugi składnik dąży do zera.

Dla $1\le p<\infty$ funkcje ciągłe okresowe
są gęste w $L^p(\mathbb T)$: przedłużamy funkcję
z $(-\pi,\pi)$ przez zero na $\R$, najpierw obcinamy
ją do $[-\pi+\eta,\pi-\eta]$ (błąd dąży do zera
przez zbieżność zmajoryzowaną), a następnie
aproksymujemy przez funkcję testową z
Twierdzenia~\ref{thm:dt-test-density}.
Mnożąc tę funkcję przez gładkie odcięcie
o wartościach w $[0,1]$, równe $1$ na obciętym
przedziale i o nośniku w $(-\pi,\pi)$
(Lemat~\ref{lem:cutoffs17}), nie zwiększamy błędu
względem funkcji obciętej. Okresowe przedłużenie
wyniku jest ciągłe. Wybieramy więc $g\in C(\mathbb T)$
z $\|f-g\|_p<\varepsilon$.
Z kontraktywności i jednostajnej zbieżności dla $g$
wynika
\[
 \|K_\nu*f-f\|_p
 \le 2\varepsilon+\|K_\nu*g-g\|_\infty.
\]
Dowolność $\varepsilon$ kończy dowód.
\end{proof}

\subsection{Jądro Fejéra i średnie Cesàra}

\begin{lemat}[Jądro Fejéra]
\label{lem:dt-series-fejer}
Dla całkowitego $N\geq0$ niech
\begin{equation}\label{eq:dt-series-fejer}
 F_N(t)=\frac1{N+1}\left|\sum_{j=0}^Ne^{ijt}\right|^2
 =\frac1{N+1}\left(\frac{\sin((N+1)t/2)}{\sin(t/2)}\right)^2
\end{equation}
(w punktach $t\in2\pi\mathbb Z$ przyjmujemy
$F_N(t)=N+1$, zgodnie z pierwszą postacią wzoru). Wówczas $F_N\ge0$,
$(2\pi)^{-1}\int_{-\pi}^{\pi}F_N(t)\,\mathrm d\lambda_1(t)=1$ oraz
dla każdego $\delta\in(0,\pi)$
$\int_{\delta\le|t|\le\pi}F_N(t)\,\mathrm d\lambda_1(t)\to0$.
\end{lemat}
\begin{proof}
Nieujemność jest bezpośrednia z definicji.
Dla $t\notin2\pi\mathbb Z$ z sumy geometrycznej
$\sum_{j=0}^N e^{ijt}=(1-e^{i(N+1)t})/(1-e^{it})$
i tożsamości $|1-e^{i\alpha}|=2|\sin(\alpha/2)|$
otrzymujemy drugą postać~\eqref{eq:dt-series-fejer}.
Rozwijając kwadrat modułu, mamy
\[
 F_N(t)=\frac1{N+1}\sum_{j,\ell=0}^N e^{i(j-\ell)t}
 =\sum_{|k|\leq N}\left(1-\frac{|k|}{N+1}\right)e^{ikt}.
\]
Dla ustalonego $|k|\leq N$ istnieje dokładnie
$N+1-|k|$ par $(j,\ell)$ z różnicą $j-\ell=k$:
dla $k\geq0$ wybieramy $\ell=0,\ldots,N-k$,
a dla $k<0$ zamieniamy role indeksów.
Wyraz stały jest równy $1$, a pozostałe
wykładniki mają całkę zero na pełnym okresie,
skąd normalizacja. Gdy $\delta\leq|t|\leq\pi$,
mamy $|\sin(t/2)|\geq\sin(\delta/2)>0$.
Dlatego
\[
 \int_{\delta\leq|t|\leq\pi}F_N(t)\,\mathrm d\lambda_1(t)
 \leq\frac{2\pi}{(N+1)\sin^2(\delta/2)}
 \longrightarrow0.
\]
\end{proof}

\begin{twierdzenie}[Fejér i zupełność układu trygonometrycznego]
\label{thm:dt-series-fejer}
Niech $f\in C(\mathbb T)$. Średnie Cesàra
$\sigma_N f=(N+1)^{-1}\sum_{j=0}^N S_jf$
zbieżne są jednostajnie do $f$ na $\mathbb T$.
W konsekwencji wielomiany trygonometryczne są gęste
w $C(\mathbb T)$ w normie supremum i w $L^2(\mathbb T)$
w normie $L^2$.
\end{twierdzenie}
\begin{proof}
W średniej z sum częściowych współczynnik
przy $\widehat f(k)e^{ikx}$ wynosi
$(N+1-|k|)/(N+1)$ dla $|k|\leq N$, ponieważ
wyraz ten występuje dla $j=|k|,\ldots,N$.
Z drugiej strony podstawienie $y=x-t$ daje
\[
 \frac1{2\pi}\int_{-\pi}^{\pi}e^{ikt}f(x-t)\,
 \mathrm d\lambda_1(t)=e^{ikx}\widehat f(k).
\]
Po podstawieniu całkujemy po odcinku
$[x-\pi,x+\pi]$, który możemy zastąpić
$[-\pi,\pi]$, gdyż $f(y)e^{-iky}$ jest
$2\pi$-okresowa i całka po pełnym okresie
nie zależy od jego początku (przez podział
odcinka i przesunięcie o wielokrotność okresu).
Rozwinięcie skończone z lematu~\ref{lem:dt-series-fejer}
daje więc
\[
 \sigma_Nf(x)=\frac1{2\pi}\int_{-\pi}^{\pi}
 F_N(t)f(x-t)\,\mathrm d\lambda_1(t).
\]
Z Lematu~\ref{lem:dt-series-fejer} wynika, że
$(F_N)$ jest dodatnią aproksymacją jedności.
Stąd Twierdzenie~\ref{thm:dt-periodic-approximate-identity}
daje jednostajną zbieżność dla funkcji ciągłych
oraz zbieżność w normie $L^2$ dla funkcji z $L^2$:
tożsamość splotowa obowiązuje także dla $f\in L^1$,
ponieważ użyto jedynie skończonych sum i całek.
Każda $\sigma_Nf$ jest wielomianem trygonometrycznym,
co dowodzi obu twierdzeń o gęstości.
\end{proof}

\begin{wniosek}[Średnie Fejéra w przestrzeniach $L^p$]
\label{cor:dt-fejer-lp}
Dla każdego $1\le p<\infty$ i $f\in L^p(\mathbb T)$
średnie $\sigma_Nf$ zbiegają do $f$ w normie $L^p$.
Wielomiany trygonometryczne są zatem gęste
w każdej z tych przestrzeni.
\end{wniosek}
\begin{proof}
Lemat~\ref{lem:dt-series-fejer} pokazuje, że
$(F_N)$ jest dodatnią aproksymacją jedności.
Tożsamość $\sigma_N f=F_N*f$ z dowodu
Twierdzenia~\ref{thm:dt-series-fejer} obowiązuje
również dla $f\in L^1$, ponieważ wymaga jedynie
zamiany skończonej sumy z całką.
Stosujemy Twierdzenie~\ref{thm:dt-periodic-approximate-identity}.
Każda średnia $\sigma_Nf$ jest wielomianem trygonometrycznym.
\end{proof}

\subsection{Zupełność układu i tożsamość Parsevala}

\begin{twierdzenie}[Zbieżność w $L^2$ i tożsamość Parsevala]
\label{thm:dt-series-parseval}
Dla dowolnego $f\in L^2(\mathbb T)$ zachodzą
$S_Nf\to f$ w normie $L^2$ oraz
\begin{equation}\label{eq:dt-series-parseval}
 \frac1{2\pi}\int_{-\pi}^{\pi}|f(x)|^2\,\mathrm d\lambda_1(x)
 =\sum_{k\in\mathbb Z}|\widehat f(k)|^2.
\end{equation}
\end{twierdzenie}
\begin{proof}
Dla $\varepsilon>0$ wybieramy wielomian trygonometryczny $p$
z $\|f-p\|_2<\varepsilon$. Dla $N$ nie mniejszych niż stopień
$p$ mamy $S_Np=p$. Ponieważ $f-S_Nf$ jest
prostopadła do przestrzeni, do której należy
$S_Nf-p$, zachodzi
\[
 \|f-p\|_2^2=\|f-S_Nf\|_2^2+\|S_Nf-p\|_2^2.
\]
Zatem $\|f-S_Nf\|_2\leq\|f-p\|_2<\varepsilon$,
co dowodzi zbieżności.
W konsekwencji $\|S_Nf\|_2\to\|f\|_2$,
a z tożsamości $\|S_Nf\|_2^2=
\sum_{|k|\leq N}|\widehat f(k)|^2$
w lemacie~\ref{lem:dt-series-bessel} otrzymujemy
\eqref{eq:dt-series-parseval}, ponieważ sumy
nieujemnych wyrazów rosną do pełnej sumy po $\mathbb Z$.
\end{proof}

\begin{uwaga}[Jakiego rodzaju zbieżność została udowodniona?]
Sama zbieżność w normie $L^2$ nie pociąga
za sobą zbieżności jednostajnej ani punktowej
całego ciągu. Nie rozstrzygamy tutaj odrębnego
zagadnienia zbieżności prawie wszędzie sum
częściowych Fouriera. Dla funkcji ciągłych jednostajnie zbieżne są
\emph{średnie Fejéra}, niekoniecznie zwykłe sumy $S_N$.
Dla gładkich funkcji zachodzi dodatkowe szybkie zanikanie
współczynników, co wykażemy w następnym rozdziale.
\end{uwaga}

\subsection{Jądro Poissona i sumowanie Abela}

\begin{twierdzenie}[Jądro Poissona na okręgu]
\label{thm:dt-periodic-poisson}
Dla $0\le r<1$ kładziemy
\[
 P_r(t)=\frac{1-r^2}{1-2r\cos t+r^2}.
\]
Wówczas
\[
 P_r(t)=\sum_{k\in\mathbb Z}r^{|k|}e^{ikt}
 =1+2\sum_{k=1}^\infty r^k\cos(kt),
\]
przy czym wyraz dla $k=0$ wynosi $1$ także
dla $r=0$. Szereg jest bezwzględnie i jednostajnie
zbieżny dla każdego ustalonego $r<1$.
Rodzina $(P_r)$ jest dodatnią aproksymacją jedności
przy $r\uparrow1$.
\end{twierdzenie}
\begin{proof}
Z sumy szeregu geometrycznego i jego sprzężenia
zespolonego otrzymujemy
\[
 1+\frac{re^{it}}{1-re^{it}}
   +\frac{re^{-it}}{1-re^{-it}}
 =\frac{1-r^2}{|1-re^{it}|^2}.
\]
Jednostajna zbieżność wynika z majoranty
$\sum_{k\ge1}2r^k$. Nieujemność wynika
z ilorazu, a całkowanie wyraz po wyrazie
i ortogonalność dają $\int_{\mathbb T}P_r\,\mathrm dm=1$.
Dla $\delta\le|t|\le\pi$ i $r\ge1/2$ mamy
\[
 |1-re^{it}|^2=(1-r)^2+2r(1-\cos t)
 \ge1-\cos\delta>0.
\]
Stąd całka poza otoczeniem zera jest ograniczona
przez $(1-r^2)/(1-\cos\delta)$ i dąży do zera.
\end{proof}

\begin{wniosek}[Sumowanie Abela szeregów Fouriera]
\label{cor:dt-poisson-abel}
Dla $f\in L^1(\mathbb T)$ i $0\le r<1$ zachodzi
\[
 (P_r*f)(x)=\sum_{k\in\mathbb Z}r^{|k|}
 \widehat f(k)e^{ikx}.
\]
Jeżeli $f\in C(\mathbb T)$, lewa strona zbiega
jednostajnie do $f$ przy $r\uparrow1$.
Jeżeli $f\in L^p$, $1\le p<\infty$,
zbieżność zachodzi w normie $L^p$.
\end{wniosek}
\begin{proof}
Jednostajna zbieżność szeregu jądra i całkowalność
$f$ pozwalają zamienić szereg z całką splotową.
Podstawienie $y=x-t$ daje współczynniki Fouriera.
Ponadto $|\widehat f(k)|\le\|f\|_{L^1(\mathbb T)}$,
więc wynikowy szereg jest bezwzględnie i jednostajnie
zbieżny względem $x$ dla ustalonego $r<1$.
Oba rodzaje zbieżności wynikają z
Twierdzenia~\ref{thm:dt-periodic-poisson}
i Twierdzenia~\ref{thm:dt-periodic-approximate-identity}.
\end{proof}

Jądra Fejéra i Poissona mają masę $1$ i są nieujemne,
w odróżnieniu od jądra Dirichleta. Dzięki temu
średnie Cesàra i sumy Abela odzyskują każdą
funkcję ciągłą jednostajnie, mimo że zwykłe
sumy częściowe nie mają takiej własności
(Wniosek~\ref{cor:dt-dirichlet-unbounded}).

\begin{center}
\begin{tikzpicture}[x=1.6cm,y=0.37cm,>=stealth]
 \draw[->,gray] (-3.35,0)--(3.43,0) node[right] {$t$};
 \draw[->,gray] (0,-2.5)--(0,10) node[above] {$K(t)$};
 \foreach \x/\lab in {-3.14159/{-\pi},-1.5708/{-\pi/2},1.5708/{\pi/2},3.14159/{\pi}}
   {\draw[gray] (\x,0.12)--(\x,-0.12) node[below=2pt,fill=white,inner sep=1pt] {$\lab$};}
 \foreach \y in {3,5,9}
   {\draw[gray] (-0.04,\y)--(0.04,\y) node[left=3pt,fill=white,inner sep=1pt] {\small $\y$};}
 \draw[blue!75!black,thick,domain=-3.14159:3.14159,samples=321,variable=\x]
 plot ({\x},{1+2*cos(deg(\x))+2*cos(deg(2*\x))+2*cos(deg(3*\x))+2*cos(deg(4*\x))});
 \draw[red!75!black,thick,dashed,domain=-3.14159:3.14159,samples=321,variable=\x]
 plot ({\x},{1+1.6*cos(deg(\x))+1.2*cos(deg(2*\x))+0.8*cos(deg(3*\x))+0.4*cos(deg(4*\x))});
 \draw[teal!75!black,thick,densely dotted,domain=-3.14159:3.14159,samples=321,variable=\x]
 plot ({\x},{0.75/(1.25-cos(deg(\x)))});
 \draw[blue!75!black,thick] (1.3,9)--(1.82,9) node[right,black] {\small $D_4$};
 \draw[red!75!black,thick,dashed] (1.3,7.5)--(1.82,7.5) node[right,black] {\small $F_4$};
 \draw[teal!75!black,thick,densely dotted] (1.3,6)--(1.82,6) node[right,black] {\small $P_{1/2}$};
\end{tikzpicture}

\small Porównanie trzech jąder na jednym okresie.
Wszystkie mają unormowaną całkę równą $1$,
ale tylko jądro Dirichleta przyjmuje wartości ujemne.
Wykresy korzystają z rozwinięć trygonometrycznych
i wzoru wymiernego na jądro Poissona.
\end{center}

\subsection{Iloczynowe aproksymacje jedności na torusie}

Na $\mathbb T^d$ piszemy
$\mathrm dm_d(t)=(2\pi)^{-d}\mathrm d\lambda_d(t)$
na kostce fundamentalnej $Q=[-\pi,\pi]^d$.
Splot definiujemy analogicznie jak w jednym wymiarze.
Otoczenia zera możemy opisywać przez
$\max_j|t_j|<\delta$, gdzie $0<\delta<\pi$.

\begin{twierdzenie}[Aproksymacje jedności na torusie]
\label{thm:dt-torus-approximate-identity}
Niech $K_\nu\ge0$ na $\mathbb T^d$,
$\int K_\nu\,\mathrm dm_d=1$ oraz
\[
 \int_{\max_j|t_j|\ge\delta} K_\nu(t)\,\mathrm dm_d(t)
 \longrightarrow0
 \qquad(0<\delta<\pi).
\]
Wówczas $K_\nu*f\to f$ jednostajnie dla
$f\in C(\mathbb T^d)$ i w normie $L^p$ dla
$f\in L^p(\mathbb T^d)$, $1\le p<\infty$.
Operator splotu jest kontrakcją na $L^p$
dla $1\le p\le\infty$.
\end{twierdzenie}
\begin{proof}
Dowód Twierdzenia~\ref{thm:dt-periodic-approximate-identity}
przenosi się na torus: nierówność Jensena,
twierdzenie Tonellego i niezmienniczość $m_d$
względem przesunięć dają
$\|K_\nu*f\|_p\le\|f\|_p$.
Dla $f$ ciągłej wybieramy $\delta$ tak,
aby $\max_j|t_j|<\delta$ implikowało
$|f(x-t)-f(x)|<\varepsilon$ dla wszystkich $x$.
Podział całki na to otoczenie i jego dopełnienie
daje oszacowanie
\[
 \|K_\nu*f-f\|_\infty
 \le\varepsilon+2\|f\|_\infty
 \int_{\max_j|t_j|\ge\delta}K_\nu(t)\,\mathrm dm_d(t).
\]
W konsekwencji zachodzi jednostajna zbieżność.

Uzasadnimy gęstość $C(\mathbb T^d)$ w $L^p$
dla skończonego $p$. Przedłużamy $f|_Q$ zerem
na $\R^d$ i obcinamy ją do kostki
$Q_\eta=[-\pi+\eta,\pi-\eta]^d$.
Błąd w normie $L^p$ dąży do zera przez
zbieżność zmajoryzowaną, gdyż brzeg $Q$
ma miarę Lebesgue'a zero.
Aproksymujemy obciętą funkcję przez
$g\in\mathcal D(\R^d)$ (Twierdzenie~\ref{thm:dt-test-density}).
Wybieramy $\chi\in\mathcal D((-\pi,\pi)^d)$
o wartościach w $[0,1]$ i równą $1$ na $Q_\eta$
(Lemat~\ref{lem:cutoffs17}).
Mnożenie $g$ przez $\chi$ nie zwiększa błędu
względem obciętej funkcji, ponieważ ta ostatnia
nie zmienia się przy mnożeniu przez $\chi$.
Przedłużenie okresowe $\chi g$ jest gładkie,
gdyż funkcja znika w otoczeniu brzegu $Q$.
W ten sposób otrzymujemy żądaną gęstość.
Kontraktywność i aproksymacja przez funkcje ciągłe
kończą dowód tak jak w jednym wymiarze.
\end{proof}

\begin{wniosek}[Jądra Fejéra i Poissona na torusie]
\label{cor:dt-torus-fejer-poisson}
Rodziny
\[
 F_N^{[d]}(t)=\prod_{j=1}^dF_N(t_j),\qquad
 P_r^{[d]}(t)=\prod_{j=1}^dP_r(t_j)
\]
są dodatnimi aproksymacjami jedności na $\mathbb T^d$
odpowiednio przy $N\to\infty$ i $r\uparrow1$.
Dla $f\in L^1(\mathbb T^d)$ zachodzą wzory
\[
 \begin{aligned}
 (F_N^{[d]}*f)(x)
 &=\sum_{\max_j|k_j|\le N}
 \widehat f(k)e^{ik\cdot x}
 \prod_{j=1}^d\left(1-\frac{|k_j|}{N+1}\right),\\
 (P_r^{[d]}*f)(x)
 &=\sum_{k\in\mathbb Z^d}r^{\sum_j|k_j|}
 \widehat f(k)e^{ik\cdot x}.
 \end{aligned}
\]
W szczególności wielomiany trygonometryczne są
gęste w $C(\mathbb T^d)$ oraz w $L^p(\mathbb T^d)$
dla $1\le p<\infty$.
\end{wniosek}
\begin{proof}
Niech $H$ oznacza jedno z jąder jednowymiarowych.
Twierdzenie Fubiniego daje
$\int_{\mathbb T^d}\prod_j H(t_j)\,\mathrm dm_d=1$.
Zbiór $\{\max_j|t_j|\ge\delta\}$ jest sumą
zbiorów $\{|t_j|\ge\delta\}$, więc z nieujemności
oraz normalizacji $H$ dostajemy
\[
 \int_{\max_j|t_j|\ge\delta}\prod_j H(t_j)\,\mathrm dm_d
 \le d\int_{|t|\ge\delta}H(t)\,\mathrm dm(t).
\]
Prawa strona dąży do zera na mocy
Lematu~\ref{lem:dt-series-fejer} albo
Twierdzenia~\ref{thm:dt-periodic-poisson}.
Wstawiając rozwinięcia jednowymiarowych jąder
i mnożąc je, otrzymujemy podane szeregi splotowe.
W przypadku Poissona uzasadnienie zamiany szeregu
z całką wynika z bezwzględnej i jednostajnej zbieżności:
\[
 \sum_{k\in\mathbb Z^d} r^{\sum_j|k_j|}
 =\left(1+2\sum_{n\ge1}r^n\right)^d
 =\left(\frac{1+r}{1-r}\right)^d<\infty.
\]
Dla $r=0$ współczynnik przy $k=0$ wynosi $1$.
Twierdzenie~\ref{thm:dt-torus-approximate-identity}
daje zbieżność splotów, a sploty z jądrem Fejéra
są wielomianami trygonometrycznymi.
\end{proof}

\subsection{Twierdzenie Riesza--Fischera i unitarność}

\begin{twierdzenie}[Parseval i Riesz--Fischer na torusie]
\label{thm:dt-torus-riesz-fischer}
Układ $(e^{ik\cdot x})_{k\in\mathbb Z^d}$
jest bazą ortonormalną przestrzeni Hilberta
$L^2(\mathbb T^d,m_d)$.
Odwzorowanie
\[
 \mathcal F_{\mathbb T^d}: L^2(\mathbb T^d,m_d)
 \longrightarrow\ell^2(\mathbb Z^d),\qquad
 f\longmapsto(\widehat f(k))_{k\in\mathbb Z^d}
\]
jest izomorfizmem unitarnym. W szczególności
dla $f,g\in L^2(\mathbb T^d)$ mamy
\[
 \langle f,g\rangle=
 \sum_{k\in\mathbb Z^d}\widehat f(k)
 \overline{\widehat g(k)}.
\]
Każdy ciąg $(c_k)\in\ell^2(\mathbb Z^d)$
jest ciągiem współczynników Fouriera dokładnie
jednego elementu $L^2(\mathbb T^d)$, a odwrotne
odwzorowanie jest dane szeregiem
\[
 f=\sum_{k\in\mathbb Z^d}c_ke^{ik\cdot x}
\]
zbieżnym bezwarunkowo w normie $L^2$.
\end{twierdzenie}
\begin{proof}
Ortonormalność wynika ze
Stwierdzenia~\ref{prop:dt-series-torus},
a zupełność z gęstości wielomianów trygonometrycznych
we Wniosku~\ref{cor:dt-torus-fejer-poisson}.
Niech $A_N=\{-N,\ldots,N\}^d$.
Rzut $S_{A_N}f$ jest najlepszym przybliżeniem
$f$ w rozpiętości odpowiednich harmonicznych:
dla wielomianu $p$ o indeksach w $A_N$
rozkład $f-p=(f-S_{A_N}f)+(S_{A_N}f-p)$
jest ortogonalny. Gęstość wielomianów daje
$S_{A_N}f\to f$ w $L^2$. Przechodząc do granicy
w tożsamości norm ze
Stwierdzenia~\ref{prop:dt-series-torus},
otrzymujemy
\[
 \|f\|_2^2=\sum_{k\in\mathbb Z^d}|\widehat f(k)|^2.
\]
To daje izometryczność i różnowartościowość.
Tożsamość iloczynów skalarnych wynika przez
przejście do granicy w
$\langle S_{A_N}f,S_{A_N}g\rangle$;
szereg iloczynów współczynników jest bezwzględnie
zbieżny z nierówności Cauchy'ego--Schwarza.

Dla dowolnego $(c_k)\in\ell^2$ kładziemy
$p_N=\sum_{k\in A_N}c_ke^{ik\cdot x}$.
Z ortonormalności dla $M>N$ mamy
\[
 \|p_M-p_N\|_2^2=\sum_{k\in A_M\setminus A_N}|c_k|^2
 \longrightarrow0
\]
jednostajnie względem $M>N$ przy $N\to\infty$,
ponieważ ogon szeregu kwadratów dąży do zera.
Z zupełności przestrzeni Hilberta istnieje
$f\in L^2$ takie, że $p_N\to f$ w normie.
Dla ustalonego $k$ funkcjonał
$f\mapsto\langle f,e^{ik\cdot x}\rangle$ jest ciągły
z nierówności Cauchy'ego--Schwarza, więc
$\widehat f(k)=c_k$. Dowodzi to surjektywności.
Jeżeli $A$ jest dowolnym skończonym zbiorem
zawierającym $A_N$, to udowodniona tożsamość norm
daje
\[
 \left\|f-\sum_{k\in A}c_ke^{ik\cdot x}\right\|_2^2
 =\sum_{k\notin A}|c_k|^2
 \le\sum_{k\notin A_N}|c_k|^2.
\]
Stąd wynika bezwarunkowa zbieżność (granica
po skończonych zbiorach skierowanych przez inkluzję),
a w szczególności niezależność od enumeracji indeksów.
\end{proof}

\subsection{Jednoznaczność i zastosowania tożsamości Parsevala}

\begin{wniosek}[Jednoznaczność współczynników Fouriera w $L^1$]
\label{cor:dt-fourier-l1-uniqueness}
Jeżeli $f,g\in L^1(\mathbb T^d)$ i
$\widehat f(k)=\widehat g(k)$ dla wszystkich
$k\in\mathbb Z^d$, to $f=g$ prawie wszędzie.
Jeżeli obie funkcje są ciągłe, równość zachodzi
w każdym punkcie.
\end{wniosek}
\begin{proof}
Dla $h=f-g$ wszystkie współczynniki znikają.
Wzór z Wniosku~\ref{cor:dt-torus-fejer-poisson}
daje $F_N^{[d]}*h=0$ dla każdego $N$.
Z Twierdzenia~\ref{thm:dt-torus-approximate-identity}
wynika $F_N^{[d]}*h\to h$ w $L^1$, a więc
$\|h\|_1=0$. Dla funkcji ciągłych niezerowa
wartość w jednym punkcie pociągałaby
niezerowość na otwartym zbiorze dodatniej miary.
\end{proof}

\begin{wniosek}[Odtwarzanie funkcji z bezwzględnie sumowalnych współczynników]
\label{cor:dt-fourier-l1-absolute}
Niech $f\in L^1(\mathbb T^d)$ i niech
$\sum_{k\in\mathbb Z^d}|\widehat f(k)|<\infty$.
Wówczas szereg
\[
 g(x)=\sum_{k\in\mathbb Z^d}\widehat f(k)e^{ik\cdot x}
\]
zbiega bezwzględnie i jednostajnie, a jego suma
jest ciągłym reprezentantem $f$.
Jeżeli $f$ jest ciągła, to $g=f$ punktowo.
\end{wniosek}
\begin{proof}
Kryterium Weierstrassa daje jednostajną
zbieżność, niezależną od enumeracji indeksów.
Całkowanie wyraz po wyrazie jest uzasadnione
jednostajną zbieżnością i skończonością miary.
Ortonormalność ze Stwierdzenia~\ref{prop:dt-series-torus}
daje $\widehat g(k)=\widehat f(k)$.
Stosujemy Wniosek~\ref{cor:dt-fourier-l1-uniqueness}.
\end{proof}

\begin{przyklad}[Sumy odwrotności kwadratów i czwartych potęg]
\label{ex:dt-parseval-zeta}
Dla okresowego przedłużenia $f(x)=x$ z
Przykładu~\ref{ex:dt-series-sawtooth}
mamy $\widehat f(0)=0$ oraz
$|\widehat f(n)|=1/|n|$ dla $n\ne0$,
na podstawie Stwierdzenia~\ref{prop:dt-series-real}.
Tożsamość Parsevala daje
\[
 \frac{\pi^2}{3}
 =\frac1{2\pi}\int_{-\pi}^{\pi}x^2\,\mathrm d\lambda_1(x)
 =2\sum_{n=1}^{\infty}\frac1{n^2}.
\]
Stąd $\sum_{n=1}^{\infty}n^{-2}=\pi^2/6$.

Dla $g(x)=x^2$ z Przykładu~\ref{ex:dt-series-quadratic}
mamy $\widehat g(0)=\pi^2/3$ oraz
$\widehat g(n)=2(-1)^n/n^2$ dla $n\ne0$.
Ze Wniosku~\ref{cor:dt-fourier-l1-absolute}
wynika tożsamość punktowa
\[
 x^2=\frac{\pi^2}{3}+4\sum_{n=1}^\infty
 \frac{(-1)^n}{n^2}\cos(nx),\qquad -\pi\le x\le\pi.
\]
Tożsamość Parsevala zastosowana do $g$ daje
\[
 \frac{\pi^4}{5}
 =\frac1{2\pi}\int_{-\pi}^{\pi}x^4\,\mathrm d\lambda_1(x)
 =\frac{\pi^4}{9}+8\sum_{n=1}^{\infty}\frac1{n^4}.
\]
Po odjęciu wyrazu stałego otrzymujemy
$\sum_{n=1}^{\infty}n^{-4}=\pi^4/90$.
\end{przyklad}

\subsection{Półgrupa Poissona}

\begin{twierdzenie}[Własność półgrupowa jąder Poissona]
\label{thm:dt-poisson-semigroup}
Niech $0\le r,s<1$. Na $\mathbb T^d$ zachodzi
$P_r^{[d]}*P_s^{[d]}=P_{rs}^{[d]}$.
Dla $t>0$ oznaczmy przez $Q_t$ operator splotu
z $P_{e^{-t}}^{[d]}$ i połóżmy $Q_0=I$.
Wówczas $(Q_t)_{t\ge0}$ jest półgrupą kontrakcji,
silnie ciągłą na $L^p(\mathbb T^d)$ dla
$1\le p<\infty$ oraz na $C(\mathbb T^d)$
z normą supremum.
\end{twierdzenie}
\begin{proof}
Dla $f,g\in L^1(\mathbb T^d)$ twierdzenie
Fubiniego i podstawienie $y=x-t$ dają
\[
 \widehat{f*g}(k)=\widehat f(k)\widehat g(k).
\]
Zastosowanie twierdzenia Fubiniego jest
uzasadnione całkowalnością modułu:
$\int\!\int |f(t)g(x-t)|\,\mathrm dm_d(x)\,\mathrm dm_d(t)
=\|f\|_1\|g\|_1$.
Z rozwinięcia jąder Poissona we
Wniosku~\ref{cor:dt-torus-fejer-poisson} wynika
$\widehat P_r^{[d]}(k)=r^{\sum_j|k_j|}$,
z konwencją $r^0=1$ także dla $r=0$.
Stąd współczynniki splotu są współczynnikami
$P_{rs}^{[d]}$. Jednoznaczność w $L^1$ i ciągłość
obu stron dają równość punktową.
Łączność splotu wynika z twierdzenia Fubiniego
dla trzech funkcji całkowalnych, a zatem
$Q_tQ_u=Q_{t+u}$; przypadek zerowego parametru
wynika z definicji $Q_0$.

Kontraktywność i silna ciągłość w zerze
wynikają z Twierdzenia~\ref{thm:dt-torus-approximate-identity}.
Dla $t\ge u\ge0$ i dla każdej rozważanej normy mamy
\[
 \|Q_tf-Q_uf\|
 =\|Q_u(Q_{t-u}f-f)\|
 \le\|Q_{t-u}f-f\|.
\]
To daje ciągłość przy każdym $t\ge0$.
\end{proof}

\begin{uwaga}
Dla $d>1$ jest to półgrupa odpowiadająca mnożnikowi
$\exp(-t\sum_j|k_j|)$, a nie mnożnikowi
$\exp(-t|k|)$ z normą euklidesową.
Obie konstrukcje zgadzają się w jednym wymiarze,
ale w wyższych wymiarach wyznaczają różne operatory.
\end{uwaga}

\begin{twierdzenie}[Generator półgrupy iloczynowej Poissona na $L^2$]
\label{thm:dt-poisson-generator-l2}
Niech $H=L^2(\mathbb T^d,m_d)$ i
$\rho(k)=\sum_{j=1}^d|k_j|$.
Generator $A$ półgrupy $(Q_t)$, określony
na zbiorze tych $f\in H$, dla których istnieje
granica w normie $H$
\[
 Af=\lim_{t\downarrow0}\frac{Q_tf-f}{t},
\]
ma dziedzinę
\[
 \operatorname{Dom}(A)=
 \left\{ f\in H:
 \sum_{k\in\mathbb Z^d}\rho(k)^2|\widehat f(k)|^2
 <\infty\right\},
 \qquad
 \widehat{Af}(k)=-\rho(k)\widehat f(k).
\]
Operator $A$ jest gęsto określony, domknięty,
samosprzężony i niedodatni, tj.
$\langle Af,f\rangle\le0$.
Dla $d=1$ jest to realizacja operatora
$-(-\partial_x^2)^{1/2}$ przez mnożniki Fouriera.
Dla $d>1$ odpowiada ona sumie operatorów
$-\sum_{j=1}^d(-\partial_{x_j}^2)^{1/2}$,
gdzie każdy pierwiastek oznacza mnożnik $|k_j|$
na właściwej dziedzinie w $L^2$.
\end{twierdzenie}
\begin{proof}
Na mocy Wniosku~\ref{cor:dt-torus-fejer-poisson}
i Twierdzenia~\ref{thm:dt-torus-riesz-fischer}
półgrupa na $H$ odpowiada mnożnikom
$e^{-t\rho(k)}$ na $\ell^2(\mathbb Z^d)$.
Ciągłość funkcjonału wyznaczającego współczynnik
Fouriera pokazuje, że jeżeli ilorazy różnicowe
zbiegają do $g\in H$, to dla każdego $k$
\[
 \widehat g(k)=
 \lim_{t\downarrow0}\frac{e^{-t\rho(k)}-1}{t}
 \widehat f(k)=-\rho(k)\widehat f(k).
\]
Z tożsamości Parsevala wynika konieczność
warunku sumowalności.

Odwrotnie, dla $f$ spełniającej ten warunek
twierdzenie Riesza--Fischera daje $g\in H$
o współczynnikach $-\rho(k)\widehat f(k)$.
Dla $u\ge0$ zachodzi $0\le1-e^{-u}\le u$
(z całkowania $e^{-v}\le1$ po $[0,u]$).
Wobec tego
\[
 \left|
 \left(\frac{e^{-t\rho(k)}-1}{t}+\rho(k)\right)
 \widehat f(k)\right|^2
 \le \rho(k)^2|\widehat f(k)|^2.
\]
Lewy składnik dąży do zera dla każdego $k$,
a prawa strona jest sumowalna.
Z twierdzenia o zbieżności zmajoryzowanej
dla miary liczącej i z Parsevala wynika
zbieżność ilorazów różnicowych do $g$ w $H$.

Dziedzina zawiera wielomiany trygonometryczne,
a więc jest gęsta. Unitarne odwzorowanie Fouriera
przenosi $A$ na mnożenie przez rzeczywisty ciąg
$-\rho$ na maksymalnej dziedzinie w $\ell^2$.
Domkniętość widać bezpośrednio: jeśli $c_n\to c$
oraz $-\rho c_n\to b$ w $\ell^2$, to zbieżność
każdej współrzędnej daje $b_k=-\rho(k)c_k$.
W szczególności $c$ należy do dziedziny.
Symetryczność wynika z rzeczywistości $\rho$.
Jeżeli $b$ należy do dziedziny operatora
sprzężonego i $h=A^*b$ w reprezentacji ciągowej,
to z definicji sprzężenia, testując na ciągu
jednostkowym w pozycji $k$, otrzymujemy
$h_k=-\rho(k)b_k$. Ponieważ $h\in\ell^2$,
wektor $b$ należy do dziedziny mnożenia.
Symetryczność daje inkluzję przeciwną,
a więc operator jest samosprzężony.
Ponadto z Parsevala
\[
 \langle Af,f\rangle
 =-\sum_{k\in\mathbb Z^d}\rho(k)|\widehat f(k)|^2
 \le0.
\]
Szereg jest bezwzględnie zbieżny z nierówności
Cauchy'ego--Schwarza.
Dla $d=1$ symbol $-\partial_x^2$ na wykładnikach
wynosi $k^2$, a jego nieujemny pierwiastek
w reprezentacji Fouriera ma symbol $|k|$.
Analogicznie dla każdej współrzędnej.
Dziedzina sumy jest przecięciem dziedzin tych
pierwiastków: $|k_j|\le\rho(k)$ oraz
$\rho(k)^2\le d\sum_j|k_j|^2$.
\end{proof}

\begin{wniosek}[Regularność półgrupy Poissona]
\label{cor:dt-poisson-smoothing}
Dla $f\in L^2(\mathbb T^d)$ i $t>0$
funkcja $Q_tf$ ma gładkiego reprezentanta.
Dla każdego $m\in\N_0$ istnieje stała $C_m$
taka, że
\[
 \|A^mQ_tf\|_2\le C_m t^{-m}\|f\|_2.
\]
Ponadto $u(t)=Q_tf$ jest różniczkowalna w normie
$L^2$ dla $t>0$, spełnia $u'(t)=Au(t)$
i $u(t)\to f$ w $L^2$ dla $t\downarrow0$.
\end{wniosek}
\begin{proof}
Dla ustalonego $t>0$ i wielowskaźnika $\alpha$
szereg
\[
 \sum_k(ik)^\alpha e^{-t\rho(k)}\widehat f(k)e^{ik\cdot x}
\]
jest jednostajnie i bezwzględnie zbieżny:
$|\widehat f(k)|\le\|f\|_2$,
$|k^\alpha|\le\rho(k)^{|\alpha|}$ dla $k\ne0$,
a $\rho(k)^{|\alpha|}e^{-t\rho(k)/2}$ jest ograniczone.
Sumowalność $\sum_k e^{-t\rho(k)/2}$ wynika
z iloczynu szeregów geometrycznych, jak w dowodzie
Wniosku~\ref{cor:dt-torus-fejer-poisson}.
Kryterium Weierstrassa i twierdzenie o różniczkowaniu
szeregu funkcji, stosowane kolejno do każdej zmiennej,
dają gładkość reprezentanta.
Iteracja opisu generatora pokazuje, że $Q_tf$
należy do dziedziny wszystkich jego potęg.
Z Parsevala oraz
$\sup_{v\ge0}v^me^{-v}<\infty$
(dla $m=0$ przyjmujemy $C_0=1$)
otrzymujemy żądane oszacowanie.
Różniczkowanie mnożnika $e^{-t\rho(k)}$
względem $t$ daje $-\rho(k)e^{-t\rho(k)}$.
Dla $|h|<t/2$ odpowiednie ilorazy różnicowe
są ograniczone przez $\rho(k)e^{-t\rho(k)/2}$,
z twierdzenia o wartości średniej.
Ten mnożnik jest ograniczony niezależnie od $k$,
więc zbieżność zmajoryzowana dla miary liczącej
i Parseval uzasadniają różniczkowanie w normie $L^2$.
Warunek początkowy wynika z silnej ciągłości półgrupy.
\end{proof}

\setcounter{section}{26}
% Kandydat restrukturyzacji; nie jest aktywnym rozdziałem Atlasu.
\section{Zbieżność jednostajna szeregów Fouriera}
\providecommand{\atlasref}[2]{#2}

\begin{lemat}[Zanikanie współczynników gładkiej funkcji]
\label{lem:dt-series-decay}
Jeśli $f\in C^m(\mathbb T)$, $m\geq1$ całkowite, to dla $k\ne0$
$|\widehat f(k)|\le\|f^{(m)}\|_{L^1(-\pi,\pi)}/(2\pi|k|^m)$.
Jeśli $f\in C^\infty(\mathbb T)$, to szereg Fouriera wraz
ze wszystkimi szeregami otrzymanymi przez różniczkowanie
zbiega jednostajnie.
\end{lemat}
\begin{proof}
Dla $k\ne0$ całkowanie przez części daje
\[
 \widehat f(k)=\frac1{ik}\widehat{f'}(k),
\]
ponieważ $f(\pi)e^{-ik\pi}=f(-\pi)e^{ik\pi}$.
Powtarzając ten rachunek $m$ razy (pochodne
do rzędu $m-1$ mają zgodne wartości na końcach
przedziału), dostajemy
$\widehat f(k)=(ik)^{-m}\widehat{f^{(m)}}(k)$,
co po oszacowaniu modułu całki daje żądaną nierówność.

Dla $f\in C^\infty$ i ustalonego całkowitego
$r\geq0$ wybieramy $m\geq r+2$. Wówczas
\[
 \sum_{k\ne0}|k|^r|\widehat f(k)|
 \leq\frac{\|f^{(m)}\|_{L^1(-\pi,\pi)}}{2\pi}
 \sum_{k\ne0}|k|^{r-m}<\infty.
\]
Kryterium Weierstrassa daje bezwzględną
i jednostajną zbieżność szeregu
$\sum_{k\in\mathbb Z}(ik)^r\widehat f(k)e^{ikx}$,
przy czym dla $r=0$ wyraz $k=0$ oznacza
$\widehat f(0)$, a dla $r>0$ wynosi zero.
Oznaczmy jego ciągłą sumę przez $G_r$.
Z całkowania przez części także dla $k=0$
(jeśli $r>0$, całka z $f^{(r)}$ po okresie
jest zerowa) współczynniki tego szeregu
są dokładnie współczynnikami Fouriera $f^{(r)}$.
Twierdzenie~\ref{thm:dt-series-parseval} daje więc
zbieżność tych samych sum częściowych do $f^{(r)}$
w normie $L^2$. Z jednostajnej zbieżności
mamy zarazem zbieżność w $L^2$ do $G_r$.
Jednoznaczność granicy w $L^2$ daje
$G_r=f^{(r)}$ prawie wszędzie, a ponieważ
obie funkcje są ciągłe, równość zachodzi
w każdym punkcie: niezerowa różnica
w jednym punkcie byłaby niezerowa na
pewnym otwartym łuku dodatniej miary.
W szczególności $G_0=f$ i wszystkie
różniczkowania wyraz po wyrazie prowadzą
do właściwych pochodnych.
\end{proof}

\setcounter{part}{6}
\part{Całki osobliwe i klasyczna analiza harmoniczna}
\setcounter{section}{40}
% Kandydat restrukturyzacji; nie jest aktywnym rozdziałem Atlasu.

\section{Operator maksymalny Hardy'ego--Littlewooda}
\label{sec:dt-maximal}

Rozdział opiera się na
\atlasref{ME-24::thm:HL-weak-24}{ME-24: nierówności słabego typu dla funkcji maksymalnej}.
Używamy również lematu o pięciokrotnym powiększeniu
\atlasref{ME-24::thm:5r-24}{ME-24: lemat 5r}.
Dla $d\ge1$ niech $v_d=\lambda_d(B(0,1))$.

\begin{definicja}[Funkcja maksymalna Hardy'ego--Littlewooda]
\label{def:dt-HL-max}
Dla $f\in L^1_{\mathrm{loc}}(\mathbb R^d)$ definiujemy
\emph{centrowaną funkcję maksymalną}
\[
 Mf(x):=\sup_{r>0}\frac1{\lambda_d(B(x,r))}
 \int_{B(x,r)}|f(y)|\,\mathrm d\lambda_d(y)
 \in[0,+\infty].
\]
Definicja nie zależy od reprezentanta klasy $f$ w
$L^1_{\mathrm{loc}}$, ponieważ całki funkcji równych
prawie wszędzie są równe na każdej kuli.
Funkcja $Mf$ jest mierzalna jako supremum przeliczalnej
rodziny ciągłych średnich dla dodatnich promieni wymiernych. Dokładniej, gdy $(x_n,r_n)\to(x,r)$ z $r>0$,
wszystkie kule $B(x_n,r_n)$ od pewnego indeksu
mieszczą się w jednej ograniczonej kuli $D$.
Dla $y\notin\partial B(x,r)$ zachodzi
$\chi_{B(x_n,r_n)}(y)\to\chi_{B(x,r)}(y)$.
Sfera ma miarę zero, a funkcja $|f|\chi_D$
jest całkowalna. Twierdzenie Lebesgue'a o zbieżności
zmajoryzowanej daje ciągłość licznika; mianownik
$v_dr_n^d\to v_dr^d>0$. W szczególności można
ograniczyć supremum do dodatnich promieni wymiernych.
\end{definicja}

\begin{twierdzenie}[Nierówność słabego typu $(1,1)$]
\label{thm:dt-HL-weak}
Jeśli $f\in L^1(\mathbb R^d)$, to dla każdego $t>0$
\[
 \lambda_d\bigl(\{x\in\mathbb R^d:Mf(x)>t\}\bigr)
 \le\frac{5^d}{t}\int_{\mathbb R^d}|f(x)|\,
 \mathrm d\lambda_d(x).
\]
\end{twierdzenie}
\begin{proof}
Jeśli $\{Mf>t\}=\varnothing$, teza jest oczywista.
W przeciwnym razie ustalmy niepusty zwarty zbiór
$K\subseteq\{Mf>t\}$ (dla pustego $K$ żądane
oszacowanie jest natychmiastowe).
Dla każdego $x\in K$ wybieramy promień $r_x>0$ taki,
że $\int_{B(x,r_x)}|f|\,\mathrm d\lambda_d>
 t\lambda_d(B(x,r_x))$. Promienie takich kul są
ograniczone z góry: ponieważ ich objętość jest mniejsza
od $\|f\|_1/t$, zachodzi $r_x^d<\|f\|_1/(tv_d)$.
Z przytoczonego lematu $5r$ wybieramy rozłączne kule $(B_j)$, których
pięciokrotne powiększenia pokrywają $K$. Wówczas
\[
 \lambda_d(K)\le5^d\sum_j\lambda_d(B_j)
 <\frac{5^d}{t}\sum_j\int_{B_j}|f|\,\mathrm d\lambda_d
 \le\frac{5^d}{t}\|f\|_1.
\]
Zbiór $\{Mf>t\}$ jest otwarty: dla dowolnego $x$
warunek średniej większej od $t$ zachowuje się przy małej
zmianie środka kuli dzięki ciągłości całki z $|f|$.
Jego miara jest supremum miar zwartych podzbiorów,
na mocy regularności miary Lebesgue'a. Biorąc supremum
po $K$, kończymy dowód.
\end{proof}

\begin{wniosek}[Kontrola zbioru, na którym średnie są duże]
\label{cor:dt-HL-level}
Dla $f\in L^1(\mathbb R^d)$ oraz $a>0$ funkcja $Mf$
jest skończona prawie wszędzie, a
$\lambda_d(\{Mf>a\})\le5^d\|f\|_1/a$.
\end{wniosek}
\begin{proof}
Oszacowanie wynika z twierdzenia~\ref{thm:dt-HL-weak}.
Przechodząc do granicy $a\to\infty$, dostajemy
$\lambda_d(\{Mf=+\infty\})=0$.
\end{proof}


\subsection{Ograniczoność w przestrzeniach Lebesgue'a}

\begin{twierdzenie}[Silne oszacowanie operatora maksymalnego]
\label{thm:dt-HL-strong}
Niech $1<p<\infty$ i $f\in L^p(\R^d)$.
Wówczas $Mf\in L^p(\R^d)$ oraz
\[
 \|Mf\|_p^p\le \frac{2^p5^d p}{p-1}\|f\|_p^p.
\]
Dla $p=\infty$ zachodzi $\|Mf\|_\infty\le\|f\|_\infty$.
\end{twierdzenie}
\begin{proof}
Dla $1<p<\infty$ nierówność Höldera na każdej
ograniczonej kuli zapewnia $f\in L^1_{\mathrm{loc}}$.
Zatem $Mf$ jest określona i mierzalna na mocy
definicji \ref{def:dt-HL-max}. Dla $t>0$ rozkładamy
$f=f_1+f_2$, gdzie
$f_1=f\chi_{\{|f|>t/2\}}$ i $f_2=f-f_1$.
Z nierówności $|f_1|\le(2/t)^{p-1}|f|^p$
wynika $f_1\in L^1$. Ponadto $|f_2|\le t/2$
prawie wszędzie, więc średnia modułu $f_2$
na każdej kuli nie przekracza $t/2$.
Nierówność trójkąta pod całką i supremum po promieniach
dają $Mf\le Mf_1+Mf_2\le Mf_1+t/2$.
Stąd $\{Mf>t\}\subseteq\{Mf_1>t/2\}$.
Z twierdzenia \ref{thm:dt-HL-weak} otrzymujemy
\[
 \lambda_d(\{Mf>t\})
 \le \frac{2\,5^d}{t}
 \int_{\{|f|>t/2\}}|f(x)|\,\mathrm d\lambda_d(x).
\]
Dla każdej liczby $a\in[0,+\infty]$ zachodzi
$a^p=\int_0^\infty p t^{p-1}\chi_{\{t<a\}}\,\mathrm dt$;
dla $a=+\infty$ rozumiemy obie strony jako $+\infty$.
Podstawiając $a=Mf(x)$ i stosując twierdzenie Tonellego
do nieujemnej funkcji mierzalnej, otrzymujemy
\[
 \begin{aligned}
 \int_{\R^d}(Mf)^p\,\mathrm d\lambda_d
 &=p\int_0^\infty t^{p-1}
            \lambda_d(\{Mf>t\})\,\mathrm dt\\
 &\le 2\,5^d p\int_{\R^d}|f(x)|
       \left(\int_0^{2|f(x)|}t^{p-2}\,\mathrm dt\right)
       \mathrm d\lambda_d(x)\\
 &=\frac{2^p5^d p}{p-1}\int_{\R^d}|f|^p\,\mathrm d\lambda_d.
 \end{aligned}
\]
W drugiej linii ponownie stosujemy Tonellego;
całka wewnętrzna jest skończona dla skończonego
$|f(x)|$, ponieważ $p-2>-1$. Ostatnia strona
jest skończona, co daje tezę i skończoność $Mf$
prawie wszędzie.
Dla $p=\infty$ średnia na każdej kuli jest
ograniczona przez $\|f\|_\infty$, a supremum
zachowuje nierówność.
\end{proof}

\begin{uwaga}
Oszacowania maksymalne i twierdzenie Lebesgue'a
o różniczkowaniu całek
(\atlasref{ME-25::thm:lebesgue-diff-25}{ME-25: różniczkowanie całek})
są podstawą dalszego rozkładu Calderóna--Zygmunda.
Przypadek $p=1$ wymaga słabego oszacowania;
nie wynika z powyższej nierówności silnego typu.
\end{uwaga}

\setcounter{section}{41}
% Kandydat restrukturyzacji; nie jest aktywnym rozdziałem Atlasu.
\section{Interpolacja rzeczywista i twierdzenie Marcinkiewicza}
\providecommand{\atlasref}[2]{#2}

\begin{lemat}[Interpolacja Marcinkiewicza pomiędzy słabym $L^1$ i silnym $L^2$]
\label{lem:dt-cz-interp}
Niech $T$ będzie ograniczonym operatorem liniowym
na $L^2(\R^d)$, spełniającym oszacowanie słabego typu
$(1,1)$ na $L^1\cap L^2$ ze stałą $C_1$.
Dla każdego $1<p<2$ jego ograniczenie do
$L^1\cap L^2\cap L^p$ spełnia
$\|Tf\|_p\le C_p\|f\|_p$, gdzie $C_p$ zależy również
od stałych obu oszacowań. Ma ono jednoznaczne
ograniczone przedłużenie na $L^p$.
\end{lemat}
\begin{proof}
Dla $f\in L^1\cap L^2\cap L^p$ i $t>0$ kładziemy
$f_t^{>}=f\chi_{\{|f|>t\}}$ i
$f_t^{\le}=f\chi_{\{|f|\le t\}}$.
Obie funkcje należą do $L^1\cap L^2$.
Z liniowości i nierówności trójkąta wynika
$\{|Tf|>2t\}\subset\{|Tf_t^{>}|>t\}
\cup\{|Tf_t^{\le}|>t\}$ z dokładnością do zbioru
miary zero. Oznaczając $\|T\|_{L^2\to L^2}=C_2$,
otrzymujemy
\[
 \lambda_d(\{|Tf|>2t\})
 \le \frac{C_1}{t}\int_{\{|f|>t\}}|f|\,\mathrm d\lambda_d
 +\frac{C_2^2}{t^2}\int_{\{|f|\le t\}}|f|^2\,\mathrm d\lambda_d .
\]
Tożsamość
$|Tf(x)|^p=\int_0^{|Tf(x)|}p s^{p-1}\,\mathrm d\lambda_1(s)$
i twierdzenie Tonellego dają tożsamość warstwową
\[
 \|Tf\|_p^p=p\int_0^\infty s^{p-1}
 \lambda_d(\{|Tf|>s\})\,\mathrm d\lambda_1(s)
\]
również wtedy, gdy początkowo obie strony są nieskończone.
Po podstawieniu $s=2t$ i użyciu Tonellego otrzymujemy
\[
 \begin{aligned}
 \|Tf\|_p^p
 &\le p2^p C_1\int_{\{|f|>0\}} |f(x)|
       \int_0^{|f(x)|}t^{p-2}\,\mathrm d\lambda_1(t)
       \,\mathrm d\lambda_d(x)\\
 &\quad+p2^p C_2^2\int_{\{|f|>0\}} |f(x)|^2
       \int_{|f(x)|}^\infty t^{p-3}\,\mathrm d\lambda_1(t)
       \,\mathrm d\lambda_d(x)\\
 &=p2^p\left(\frac{C_1}{p-1}+\frac{C_2^2}{2-p}\right)
   \|f\|_p^p .
 \end{aligned}
\]
Punkty $f=0$ pomijamy, gdyż pierwotne całki po $t$
są w nich zerowe. Warunki $p>1$ i $p<2$ odpowiadają
zbieżności obu całek wewnętrznych.
Gęstość $C_c^\infty$ w $L^p$
(twierdzenie~\ref{thm:dt-test-density}) i zupełność $L^p$
dają jednoznaczne przedłużenie.
Jest ono zgodne z wyjściowym operatorem $T$ na
$L^2\cap L^p$: dla $f$ z tej przestrzeni funkcje
$f_n=f\chi_{B_n(0)}\chi_{\{|f|\le n\}}$
należą do $L^1\cap L^2\cap L^p$ i zbiegają do $f$
w obu normach przez zbieżność zmajoryzowaną.
Ich obrazy zbiegają w $L^2$ do $Tf$ oraz w $L^p$
do wartości przedłużenia. Z obu zbieżności wybieramy
wspólny podciąg zbieżny prawie wszędzie, więc granice
są równe prawie wszędzie.
\end{proof}

\setcounter{section}{42}

\section{Wartości główne, jądra całek osobliwych i mnożniki Fouriera}
\label{sec:dt-singular}
\providecommand{\atlasref}[2]{#2}
Por. \atlasref{PDE-07::thm:pde7-laplace}{RRC-07: rozwiązanie fundamentalne laplasjanu}
oraz \atlasref{PDE-16::lem:pde16-distr-kernel}{RRC-16: jądra osobliwe}.

\subsection{Mnożniki Fouriera w przestrzeni $L^2$}
\label{sec:dt-multipliers-l2}

W tym rozdziale zachowujemy unitarną normalizację
transformacji Fouriera z definicji \ref{def:fourier-schwartz17}.
Materiał przenoszony z RRC-08 wymaga w tej konwencji
usunięcia czynników $(2\pi)^{-d/2}$ i $(2\pi)^{-d}$ w
tożsamościach norm oraz iloczynów skalarnych.
Por. \atlasref{PDE-08::thm:pde8-multiplier}{RRC-08: mnożniki Fouriera}.


\begin{twierdzenie}[Mnożnik w $L^2$]
\label{thm:dt-multiplier-l2}
Niech $m:\R^d\to\C$ będzie funkcją mierzalną i
istotnie ograniczoną. Dla $f\in L^2(\R^d)$ istnieje
dokładnie jedna klasa $T_mf\in L^2(\R^d)$ spełniająca
$\widehat{T_mf}=m\widehat f$ w sensie $L^2$.
Operator $T_m$ jest liniowy i ograniczony oraz
\[
 \|T_m\|_{L^2\to L^2}=\|m\|_{L^\infty}.
\]
Odwrotnie, jeśli $m$ jest mierzalna i skończona
prawie wszędzie, a ograniczony operator liniowy
$A:L^2\to L^2$ spełnia $\widehat{Af}=m\widehat f$
dla każdego $f\in L^2$ takiego, że
$m\widehat f\in L^2$, to $m\in L^\infty$
oraz $A=T_m$.
\end{twierdzenie}
\begin{proof}
Niech $M=\|m\|_{L^\infty}$. Z definicji istotnego
supremum dla każdego $k\ge1$ zbiór
$\{|m|>M+1/k\}$ ma miarę zero.
Ich przeliczalna suma jest zbiorem $\{|m|>M\}$,
więc $|m|\le M$ prawie wszędzie.
Dla $h\in L^2$ funkcja $mh$ jest więc kwadratowalna
i spełnia $\|mh\|_2\le M\|h\|_2$.
Mnożenie nie zależy od wyboru reprezentantów
klas funkcji $m$ i $h$.
Z twierdzenia~\ref{thm:plancherel22}
transformata Fouriera na $L^2$ jest bijekcją.
Możemy zatem określić
\[
 T_mf=\mathcal F^{-1}(m\widehat f).
\]
Jest to jedyna klasa spełniająca warunek
z tezy, a liniowość wynika z liniowości
wszystkich trzech operatorów w tym złożeniu.
Stosując dwukrotnie twierdzenie Plancherela,
otrzymujemy
\[
 \|T_mf\|_2=\|m\widehat f\|_2
 \le M\|\widehat f\|_2
 =M\|f\|_2.
\]
Stąd $\|T_m\|\le M$.

Gdy $M=0$, funkcja $m$ jest zerowa prawie
wszędzie i $T_m=0$, co daje równość norm.
Załóżmy $M>0$ i ustalmy $0<\varepsilon<M$.
Z definicji istotnego supremum zbiór
$E_\varepsilon=\{\xi:|m(\xi)|>M-\varepsilon\}$
ma dodatnią miarę: gdyby miał miarę zero,
$M-\varepsilon$ byłoby istotnym ograniczeniem
z góry dla $|m|$, wbrew definicji $M$.
Ponieważ $\R^d=\bigcup_{n\ge1}B_n(0)$,
dla pewnego $n$ zbiór
$E=E_\varepsilon\cap B_n(0)$ ma dodatnią
i skończoną miarę. Niech
\[
 h=\frac{\chi_E}{\lambda_d(E)^{1/2}},
 \qquad f=\mathcal F^{-1}h.
\]
Wówczas $\|h\|_2=1$ i
$\|f\|_2=1$. Ponadto
\[
 \|mh\|_2^2
 =\frac1{\lambda_d(E)}\int_E|m(\xi)|^2
 \,\mathrm d\lambda_d(\xi)
 \ge(M-\varepsilon)^2.
\]
W konsekwencji
$\|T_mf\|_2/\|f\|_2=\|mh\|_2\ge M-\varepsilon$.
Przechodząc z $\varepsilon$ do zera,
otrzymujemy $\|T_m\|\ge M$.

Dla implikacji odwrotnej połóżmy $C=\|A\|$.
Jeśli zbiór $\{|m|>C+\varepsilon\}$ dla pewnego
$\varepsilon>0$ ma dodatnią miarę, to z
prawie wszędzie skończoności $m$ wynika, że
z dokładnością do zbioru miary zero jest on
przeliczalną sumą poniższych zbiorów po $n\ge1$.
Dodatnia miara sumy oznacza, że dla pewnego $n$ zbiór
\[
 E=\{\xi:C+\varepsilon<|m(\xi)|\le n\}
       \cap B_n(0)
\]
ma dodatnią miarę. Ma też miarę skończoną.
Wybierając $f=\mathcal F^{-1}\chi_E$,
mamy $m\widehat f\in L^2$ dzięki ograniczeniu
$|m|\le n$ na $E$. Z założenia i Plancherela
wynika wtedy
\[
 \frac{\|Af\|_2}{\|f\|_2}
 =\frac{\|m\chi_E\|_2}{\|\chi_E\|_2}
 \ge C+\varepsilon>C,
\]
co przeczy definicji normy operatorowej.
Zatem dla każdego $\varepsilon>0$ zbiór
$\{|m|>C+\varepsilon\}$ ma miarę zero.
Biorąc przeliczalną sumę dla
$\varepsilon=1/k$, $k\ge1$, dostajemy
$|m|\le C$ prawie wszędzie. Tym samym
$m\in L^\infty$, a warunek na $A$ obowiązuje
już dla wszystkich $f\in L^2$.
Jednoznaczność daje $A=T_m$.
\end{proof}

\begin{wniosek}[Rachunek ograniczonych mnożników]
Dla $m,n\in L^\infty(\R^d)$ i $a,b\in\C$ zachodzi
\[
 T_{am+bn}=aT_m+bT_n,\qquad
 T_mT_n=T_{mn},\qquad T_1=I,\qquad
 T_m^*=T_{\overline m}.
\]
Gwiazdka oznacza operator sprzężony w zespolonej
przestrzeni Hilberta $L^2$.
\end{wniosek}
\begin{proof}
Dla $f\in L^2$ z definicji mnożników mamy
\[
 \begin{aligned}
 \mathcal F(aT_mf+bT_nf)&=(am+bn)\widehat f,\\
 \mathcal F(T_mT_nf)&=m\mathcal F(T_nf)=mn\widehat f,\\
 \mathcal F(T_1f)&=\widehat f.
 \end{aligned}
\]
Wszystkie symbole po prawej stronie są istotnie
ograniczone. Równości operatorów wynikają
z iniektywności transformaty na $L^2$.
Dla $f,g\in L^2$ tożsamość
 z twierdzenia~\ref{thm:plancherel22} daje
\[
 \begin{aligned}
 \langle T_m f,g\rangle_2
 &=\int_{\R^d}m(\xi)\widehat f(\xi)
       \overline{\widehat g(\xi)}\,\mathrm d\lambda_d(\xi)\\
 &=\int_{\R^d}\widehat f(\xi)
       \overline{\overline{m(\xi)}\widehat g(\xi)}
       \,\mathrm d\lambda_d(\xi)\\
 &=\langle f,T_{\overline m}g\rangle_2.
 \end{aligned}
\]
Całki są bezwzględnie zbieżne przez nierówność
Cauchy'ego--Schwarza i istotną ograniczoność $m$.
To dowodzi równości operatorów sprzężonych.
\end{proof}

\begin{przyklad}[Transformata Riesza w $L^2$]
Dla $1\le j\le d$ połóżmy
$m_j(\xi)=-i\xi_j/|\xi|$ dla $\xi\ne0$ i $m_j(0)=0$.
Twierdzenie~\ref{thm:dt-multiplier-l2} definiuje
ograniczony operator $R_j=T_{m_j}:L^2\to L^2$,
zwany transformatą Riesza.
Mamy $|m_j(\xi)|\le1$. Dla każdego
$0<\varepsilon<1$ zbiór
$\{\xi\ne0:|m_j(\xi)|>1-\varepsilon\}$
zawiera otwartą kulę o dodatnim promieniu
wokół $e_j$: funkcja $|\xi_j|/|\xi|$ jest
ciągła w otoczeniu $e_j$ i przyjmuje tam
wartość $1$ w punkcie $e_j$.
Zatem $\|m_j\|_\infty=1$ i w konsekwencji
$\|R_j\|_{L^2\to L^2}=1$, także dla $d=1$.

Ponieważ $\overline{m_j}=-m_j$ oraz dla
$\xi\ne0$ zachodzą
\[
 \sum_{j=1}^d |m_j(\xi)|^2=1,\qquad
 \sum_{j=1}^d m_j(\xi)^2=-1,
\]
z rachunku mnożników otrzymujemy
\[
 R_j^*=-R_j,\qquad
 \sum_{j=1}^d R_j^2=-I,\qquad
 \sum_{j=1}^d\|R_j f\|_2^2=\|f\|_2^2.
\]
Ostatnia tożsamość wynika również bezpośrednio
z Plancherela:
\[
 \sum_{j=1}^d\|R_j f\|_2^2
 =\int_{\R^d}
 \left(\sum_{j=1}^d|m_j(\xi)|^2\right)
 |\widehat f(\xi)|^2\,\mathrm d\lambda_d(\xi)
 =\|f\|_2^2.
\]
Wartość mnożnika w pojedynczym punkcie $\xi=0$
nie zmienia operatora na $L^2$, ponieważ
$\lambda_d(\{0\})=0$. Nie oznacza to jednak,
że można przez ten sam symbol mnożyć dowolną
dystrybucję temperowaną: wzdłuż $\xi=te_j$
symbol przyjmuje wartość $-i$ dla $t>0$
oraz $i$ dla $t<0$. Nie ma więc ciągłego
przedłużenia do zera i nie spełnia warunku
mnożenia dystrybucji z poprzedniej sekcji.
Na przykład dla testu $\varphi$ z $\varphi(0)\ne0$
iloczyn $m_j\varphi$ nie jest ciągły w zerze,
a zatem nie należy do $\mathcal S$.
\end{przyklad}

\begin{uwaga}[Zakres twierdzenia o mnożnikach]
Kryterium istotnej ograniczoności symbolu
dotyczy normy $L^2$ i korzysta z twierdzenia
Plancherela. Dla $p\ne2$ powyższy dowód
nie daje twierdzenia o ograniczoności
na $L^p$. W szczególności nie należy
wyciągać z niego wniosku o ograniczoności
transformat Riesza w normach $L^1$ lub
$L^\infty$. Badanie mnożników na $L^p$
dla $1<p<\infty$ wymaga dalszych narzędzi
analizy harmonicznej; do tej teorii wrócimy po omówieniu operatorów maksymalnych i interpolacji.
\end{uwaga}

\subsection{Jądra drugich pochodnych potencjału Newtona}
Niech $\omega_{d-1}=\mathcal H^{d-1}(S^{d-1})$ oznacza pole
jednostkowej sfery, a $\Gamma$ rozwiązanie fundamentalne $-\Delta$:
$-\Delta\Gamma=\delta_0$. Z twierdzenia~\ref{thm:dt-laplace} dla $d\ge3$ można wziąć
$\Gamma(x)=\bigl((d-2)\omega_{d-1}\bigr)^{-1}|x|^{2-d}$,
a dla $d=2$ funkcję $\Gamma(x)=-(2\pi)^{-1}\log|x|$.
Dla $x\ne0$ w każdym z obu przypadków
$\partial_i\Gamma(x)=-x_i/(\omega_{d-1}|x|^d)$.
Różniczkując względem $x_j$, otrzymujemy
\begin{equation}\label{eq:dt-cz-kernel}
 K_{ij}(x):=\partial_{ij}\Gamma(x)
 =\frac{d x_ix_j-\delta_{ij}|x|^2}
 {\omega_{d-1}|x|^{d+2}},\qquad 1\le i,j\le d,
\end{equation}
gdzie $\delta_{ij}$ oznacza deltę Kroneckera.
Jądro jest gładkie poza zerem i spełnia
$K_{ij}(tx)=t^{-d}K_{ij}(x)$ dla $t>0$.
Różniczkowanie tej tożsamości daje jednorodność gradientu
stopnia $-d-1$. Z ograniczoności obu funkcji na zwartej
sferze wynika $|K_{ij}(x)|\le C_d|x|^{-d}$ oraz
$|\nabla K_{ij}(x)|\le C_d|x|^{-d-1}$.
Dla $i\ne j$ odbicie względem $\theta_i=0$ zmienia znak
$\theta_i\theta_j$ i zachowuje miarę powierzchniową.
Dla $i=j$ całki z $\theta_i^2$ są równe przez symetrię
permutacji współrzędnych, a ich suma wynosi $\omega_{d-1}$.
Stąd
\[
 \int_{S^{d-1}}\theta_i\theta_j\,\mathrm d\mathcal H^{d-1}
 =\frac{\omega_{d-1}}d\delta_{ij},\qquad
 \int_{S^{d-1}}K_{ij}(\theta)\,\mathrm d\mathcal H^{d-1}(\theta)=0.
\]
Każde jądro jest niezerowe na pewnym otwartym podzbiorze
sfery (dla $i=j$ weźmy $\theta=e_i$, dla $i\ne j$
$\theta=(e_i+e_j)/\sqrt2$). Współrzędne biegunowe dają
\[
 \int_{B_1(0)}|K_{ij}(z)|\,\mathrm d\lambda_d(z)
 =\int_{S^{d-1}}|K_{ij}(\theta)|\,\mathrm d\mathcal H^{d-1}(\theta)
 \int_0^1 r^{-1}\,\mathrm d\lambda_1(r)=+\infty .
\]
Jeżeli $f$ jest ciągła i $f(x)\ne0$, to $|f(x-z)|$
jest oddzielone od zera w otoczeniu zera. Zwykła całka
Lebesgue'a z $K_{ij}(z)f(x-z)$ nie jest więc bezwzględnie
zbieżna, nawet dla takich $f\in C_c^\infty$.

\subsection{Główna wartość całki i dystrybucja związana z jądrem}

\begin{definicja}[Główna wartość całki]
\label{def:dt-principal-value}
Niech $f\in C_c^1(\R^d)$, a $K$ będzie jądrem klasy $C^1$ poza
zerem, jednorodnym stopnia $-d$ i o zerowej średniej sferycznej.
Definiujemy
\[
 T_Kf(x)=\operatorname{p.v.}\!\int_{\R^d}K(x-y)f(y)\,\mathrm d\lambda_d(y)
 :=\lim_{\varepsilon\downarrow0}
 \int_{|x-y|>\varepsilon}K(x-y)f(y)\,\mathrm d\lambda_d(y).
\]
Dla $0<\varepsilon<1$ podstawiamy $z=x-y$.
Jednorodność i zerowa średnia dają zerową całkę z $K$
po każdym pierścieniu. Zatem po odjęciu $f(x)$
w pierścieniu $\varepsilon<|z|<1$ dostajemy w granicy
\[
 \begin{aligned}
 T_Kf(x)&=\int_{|z|<1}K(z)(f(x-z)-f(x))\,\mathrm d\lambda_d(z)\\
 &\quad+\int_{|z|\ge1}K(z)f(x-z)\,\mathrm d\lambda_d(z).
 \end{aligned}
\]
Obie całki są bezwzględnie zbieżne: w pierwszej
$|f(x-z)-f(x)|\le\|\nabla f\|_\infty |z|$,
więc całka radialna jest ograniczona przez stałą razy
$\int_0^1\mathrm d\lambda_1(r)$; druga ma moduł
nie większy niż $C\|f\|_1$. Uzasadnia to istnienie
granicy dla każdego $x$ oraz
$|T_Kf(x)|\le C(\|\nabla f\|_\infty+\|f\|_1)$.
\end{definicja}

Dystrybucję $\operatorname{p.v.}K$ określamy przez
\[
 \begin{aligned}
 \langle\operatorname{p.v.}K,\varphi\rangle
 &=\int_{|z|<1}K(z)(\varphi(z)-\varphi(0))\,\mathrm d\lambda_d(z)\\
 &\quad+\int_{|z|\ge1}K(z)\varphi(z)\,\mathrm d\lambda_d(z).
 \end{aligned}
\]
Dla testów o nośniku w ustalonym zwartym $L$ moduł
działania jest ograniczony przez
$C\|\nabla\varphi\|_\infty+C\lambda_d(L)\|\varphi\|_\infty$.
Liniowość i oszacowanie przez półnormy na każdym zwartym
nośniku dowodzą ciągłości dystrybucyjnej na mocy
twierdzenia \ref{thm:distribution-estimate19}.
Jest to także dystrybucja temperowana. Dla
$\varphi\in\mathcal S$ pierwsza całka jest ograniczona
przez stałą razy $\|\nabla\varphi\|_\infty$,
a druga przez
$C Q_{d+1,0}(\varphi)\int_1^\infty r^{-1}(1+r)^{-d-1}\,\mathrm d\lambda_1(r)$,
z dodatkowym stałym czynnikiem powierzchni sfery.
To daje ciągłość na $\mathcal S$.

\begin{lemat}[Hessian rozwiązania fundamentalnego jako dystrybucja]
\label{lem:dt-cz-distr-kernel}
Dla każdego $1\le i,j\le d$ zachodzi w $\mathcal D'(\R^d)$
\begin{equation}\label{eq:dt-cz-distr-kernel}
 \partial_{ij}\Gamma
 =\operatorname{p.v.}K_{ij}-\frac{\delta_{ij}}d\delta_0.
\end{equation}
W szczególności dla $f\in C_c^\infty(\R^d)$ i $u=\Gamma*f$,
\begin{equation}\label{eq:dt-cz-repr}
 \partial_{ij}u=T_{ij}f-\frac{\delta_{ij}}d f,
 \qquad T_{ij}=T_{K_{ij}}.
\end{equation}
\end{lemat}
\begin{proof}
Niech $\varphi\in C_c^\infty(\R^d)$, $R>1$ tak duże,
aby nośnik $\varphi$ leżał w $B_R(0)$, oraz
$D_\varepsilon=\{x:\varepsilon<|x|<R\}$.
Z lokalnej całkowalności $\Gamma$ działanie
$\partial_{ij}\Gamma$ jest granicą całki
$\int_{D_\varepsilon}\Gamma\partial_{ij}\varphi\,\mathrm d\lambda_d$.
Całkując przez części najpierw względem $x_i$, potem
$x_j$, otrzymujemy
\[
 \begin{aligned}
 \int_{D_\varepsilon}\Gamma\partial_{ij}\varphi\,\mathrm d\lambda_d
 &=\int_{D_\varepsilon}K_{ij}\varphi\,\mathrm d\lambda_d\\
 &\quad+\int_{\partial D_\varepsilon}
 \Gamma n_i\partial_j\varphi\,\mathrm d\mathcal H^{d-1}\\
 &\quad-\int_{\partial D_\varepsilon}
 (\partial_i\Gamma)n_j\varphi\,\mathrm d\mathcal H^{d-1}.
 \end{aligned}
\]
Na zewnętrznej sferze oba wyrazy brzegowe znikają.
Na wewnętrznej sferze $x=\varepsilon\theta$ mamy
$n=-\theta$, $\partial_i\Gamma(x)=
-\theta_i/(\omega_{d-1}\varepsilon^{d-1})$,
a miara powierzchniowa skaluje się przez
$\varepsilon^{d-1}$. Ostatni wyraz wynosi więc
\[
 -\frac1{\omega_{d-1}}\int_{S^{d-1}}\theta_i\theta_j
 \varphi(\varepsilon\theta)\,\mathrm d\mathcal H^{d-1}(\theta)
 \longrightarrow-\frac{\delta_{ij}}d\varphi(0).
\]
Przedostatni wyraz jest ograniczony co do modułu przez
$C\|\partial_j\varphi\|_\infty\varepsilon|\log\varepsilon|$
dla $d=2$ i przez $C_d\|\partial_j\varphi\|_\infty\varepsilon$
dla $d\ge3$, więc znika w granicy.
Pierwsza całka dąży do działania $\operatorname{p.v.}K_{ij}$.
Daje to~\eqref{eq:dt-cz-distr-kernel}.
Z twierdzeń~\ref{thm:smoothing21} i
\ref{thm:convolution-derivative21} splot
$u=\Gamma*f$ jest gładki, $-\Delta u=\delta_0*f=f$ oraz
$\partial_{ij}u=(\partial_{ij}\Gamma)*f$.
Działanie $\operatorname{p.v.}K_{ij}$ na teście
$z\mapsto f(x-z)$ jest dokładnie $T_{ij}f(x)$,
a $\delta_0*f=f$. Otrzymujemy~\eqref{eq:dt-cz-repr}.
\end{proof}

\subsection{Mnożniki związane z jądrem Newtona}
\begin{lemat}[Oszacowanie w $L^2$]
\label{lem:dt-cz-l2}
Dla $f\in C_c^\infty(\R^d)$ zachodzi
$\|T_{ij}f\|_2\le C_d\|f\|_2$. Operator $T_{ij}$
przedłuża się jednoznacznie do ograniczonego operatora na
$L^2(\R^d)$.
\end{lemat}
\begin{proof}
Najpierw wykażemy przynależność do $L^2$, bez jednostajnego
oszacowania normy przez $\|f\|_2$. Z definicji głównej
wartości $T_{ij}f$ jest ograniczona na całej przestrzeni.
Niech $\operatorname{supp}f\subset B_R(0)$, $R\ge1$.
Dla $|x|>2R$ mamy $|x-y|\ge|x|/2$ na nośniku $f$,
więc całka nie jest osobliwa i
$|T_{ij}f(x)|\le C_d|x|^{-d}\|f\|_1$.
Całka kwadratu tego ograniczenia poza $B_{2R}$ jest
skończona, gdyż $\int_{2R}^\infty r^{-d-1}
\,\mathrm d\lambda_1(r)<\infty$.
Na $B_{2R}$ wystarcza ograniczoność. Stąd $T_{ij}f\in L^2$,
a z~\eqref{eq:dt-cz-repr} także $H_{ij}:=\partial_{ij}u\in L^2$.

Z $-\Delta u=f$ i przemienności pochodnych mamy
$-\Delta H_{ij}=\partial_{ij}f$ dystrybucyjnie.
Obie strony są dystrybucjami temperowanymi; równość
w $\mathcal S'$ wynika z gęstości $C_c^\infty$ w $\mathcal S$
(twierdzenie~\ref{thm:density-schwartz17}). Twierdzenie~\ref{thm:fourier-auto-tempered22} daje
\[
 |\xi|^2\widehat H_{ij}=-\xi_i\xi_j\widehat f
 \qquad\text{w }\mathcal S'(\R^d).
\]
Obie strony są funkcjami lokalnie całkowalnymi, ponieważ
$\widehat H_{ij},\widehat f\in L^2$ z twierdzenia
Plancherela. Z jednoznaczności dystrybucji regularnych
(stwierdzenie~\ref{prop:regular-injective19}) mamy równość
prawie wszędzie. Dzieląc dla $\xi\ne0$, otrzymujemy
$\widehat H_{ij}(\xi)=-(\xi_i\xi_j/|\xi|^2)\widehat f(\xi)$
prawie wszędzie. Nie ma dodatkowego składnika skupionego
w zerze, bo $\widehat H_{ij}$ jest funkcją $L^2$.
Z~\eqref{eq:dt-cz-repr} wynika
$\widehat{T_{ij}f}=m_{ij}\widehat f$, gdzie
$m_{ij}(\xi)=\delta_{ij}/d-\xi_i\xi_j/|\xi|^2$
poza zerem, a wartość w zerze jest nieistotna.
Mamy $|m_{ij}|\le1+1/d\le2$. Z twierdzenia
~\ref{thm:plancherel22} przy przyjętej konwencji
\[
 \|T_{ij}f\|_2=\|m_{ij}\widehat f\|_2
 \le 2\|\widehat f\|_2=2\|f\|_2.
\]
Argument działa również dla $d=2$, bez wymagania, aby
sam potencjał $u$ należał do $L^2$. Liniowość, gęstość
$C_c^\infty$ w $L^2$ i zupełność $L^2$ dają jednoznaczne
ciągłe przedłużenie (gęstość w $L^2$ jest wykorzystana w dowodzie
twierdzenia~\ref{thm:plancherel22}).
\end{proof}

\setcounter{section}{45}

\section{Rozkład Calderóna--Zygmunda i interpolacja}
\label{sec:dt-calderon-zygmund}
\providecommand{\atlasref}[2]{#2}
Por. \atlasref{PDE-16::lem:pde16-cz-decomp}{RRC-16: rozkład Calderóna--Zygmunda}
oraz \atlasref{ME-25::cor:cubes-differentiation-25}{ME-25: różniczkowanie po kostkach}.
\subsection{Rozkład dyadyczny}
\begin{lemat}[Rozkład Calderóna--Zygmunda]
\label{lem:dt-cz-cz-decomp}
Niech $f\in L^1(\R^d)\cap L^2(\R^d)$ i $s>0$. Istnieje
co najwyżej przeliczalna rodzina parami rozłącznych półotwartych
kostek dyadycznych $Q_k$ oraz rozkład $f=g+\sum_k b_k$
zbieżny w $L^1$ i $L^2$, taki że
\begin{enumerate}[label=(\roman*)]
\item $|g|\le 2^ds$ prawie wszędzie i $\|g\|_1\le\|f\|_1$;
\item $b_k=0$ prawie wszędzie poza $Q_k$,
 $\int_{Q_k} b_k\,\mathrm d\lambda_d=0$,
 $\sum_k\|b_k\|_1\le2\|f\|_1$;
\item $\sum_k|Q_k|\le s^{-1}\|f\|_1$ oraz
 $\|g\|_2^2\le 2^ds\|f\|_1$.
\end{enumerate}
\end{lemat}
\begin{proof}
Ustalamy siatkę kostek $2^m(k+[0,1)^d)$,
gdzie $m\in\mathbb Z$, $k\in\mathbb Z^d$.
Dwie kostki z tej siatki są rozłączne lub jedna
zawiera się w drugiej. Każda ma jedynego rodzica
o dwukrotnie większym boku. Dla $w\in L^1$
i $s>0$ nazywamy kostkę złą, jeśli średnia $|w|$
na niej przekracza $s$. Każda zła kostka zawiera
się w maksymalnej złej kostce: w łańcuchu rodziców
objętości rosną do nieskończoności, a średnie
nie przekraczają $\|w\|_1/|Q|$, więc tylko skończenie
wiele z nich może być złych. Maksymalne złe kostki
są parami rozłączne. Oznaczając przez $M_{\mathrm{dy}}w(x)$
supremum średnich $|w|$ po kostkach dyadycznych
zawierających $x$, dostajemy
\[
 \lambda_d(\{M_{\mathrm{dy}}w>s\})
 =\sum_{Q\ \mathrm{maks.}}|Q|
 \le s^{-1}\sum_{Q\ \mathrm{maks.}}\int_Q|w|\,\mathrm d\lambda_d
 \le s^{-1}\|w\|_1.
\]
Supremum jest mierzalne, gdyż siatka jest przeliczalna.

Uzasadnijmy także różniczkowanie całek po tych kostkach.
Niech $Q_n(x)$ oznacza jedyną kostkę o boku $2^{-n}$
zawierającą $x$ i $A_n w(x)$ średnią $w$ na $Q_n(x)$.
Dla $v\in C_c^\infty(\R^d)$ mamy $A_nv(x)\to v(x)$
w każdym punkcie przez ciągłość i średnicę kostki dążącą
do zera. Dla dowolnego $w\in L^1$ wybieramy
mierzalnego reprezentanta; zachodzi prawie wszędzie
\[
 \limsup_n |A_nw(x)-w(x)|
 \le M_{\mathrm{dy}}(w-v)(x)+|w(x)-v(x)|.
\]
Miara zbioru, gdzie lewa strona przekracza $\tau>0$,
jest więc ograniczona przez $4\|w-v\|_1/\tau$,
na mocy powyższej nierówności i nierówności Czebyszewa.
Z gęstości $C_c^\infty$ w $L^1$
(twierdzenie~\ref{thm:dt-test-density}) miara ta wynosi zero.
Biorąc $\tau=1/l$ dla $l\in\N$, otrzymujemy
$A_nw\to w$ prawie wszędzie.

Wracamy do $w=f$ i oznaczamy maksymalne złe kostki
przez $Q_k$. Ich rodzice nie są źli, więc
\[
 s<\frac1{|Q_k|}\int_{Q_k}|f|\,\mathrm d\lambda_d
 \le 2^ds.
\]
Poza sumą tych kostek wszystkie średnie $|f|$ są
nie większe niż $s$; zastosowane do $|f|$ różniczkowanie
daje $|f|\le s$ prawie wszędzie. Kładziemy
$g=f$ poza kostkami oraz $g=f_{Q_k}$ na $Q_k$, gdzie
$f_{Q_k}=|Q_k|^{-1}\int_{Q_k}f\,\mathrm d\lambda_d$,
a $b_k=(f-f_{Q_k})\chi_{Q_k}$.
Z definicji średniej wynika zerowa całka $b_k$.
Ponadto $|f_{Q_k}|\le 2^ds$ i
\[
 \|g\|_1
 \le\int_{\R^d\setminus\bigcup_k Q_k}|f|\,\mathrm d\lambda_d
    +\sum_k\int_{Q_k}|f|\,\mathrm d\lambda_d
 =\|f\|_1,\qquad
 \|b_k\|_1\le2\int_{Q_k}|f|\,\mathrm d\lambda_d.
\]
Otrzymujemy $\|g\|_2^2\le2^ds\|f\|_1$ oraz
oszacowanie sumy norm $b_k$. Oszacowanie objętości
wynika z nierówności dla $M_{\mathrm{dy}}f$.
Każdy $b_k$ należy do $L^2$, a punktowo prawie wszędzie
$\sum_k b_k=f-g$, przy czym w każdym punkcie co najwyżej
jeden składnik jest niezerowy.
Dla $q=1,2$ z rozłączności i twierdzenia Tonellego
\[
 \sum_k\|b_k\|_q^q=\|f-g\|_q^q<\infty .
\]
Potęga $q$ normy ogona jest sumą odpowiednich potęg
norm, więc szeregi zbiegają w obu przestrzeniach.
Wybierając reprezentanta $b_k$ równego zero poza $Q_k$,
mamy jego nośnik topologiczny zawarty w $\overline{Q_k}$,
niekoniecznie w samej półotwartej kostce.
\end{proof}

\subsection{Oszacowanie słabego typu}
\begin{twierdzenie}[Słaby typ $(1,1)$ operatora Calderóna--Zygmunda]
\label{thm:dt-cz-weak11}
Niech $d\ge1$ i $K\in C^1(\R^d\setminus\{0\};\C)$
spełnia $|K(x)|\le A|x|^{-d}$ i
$|\nabla K(x)|\le A|x|^{-d-1}$, gdzie $A>0$.
Niech $T:L^2(\R^d)\to L^2(\R^d)$ będzie ograniczonym
operatorem liniowym, którego działanie dla każdego
$f\in C_c^\infty(\R^d)$ poza $\supp f$ jest dane przez
zwykłą całkę $\int K(x-y)f(y)\,\mathrm d\lambda_d(y)$
prawie wszędzie.
Istnieje $C>0$, zależna tylko od $d,A,\|T\|_{2\to2}$,
taka, że dla każdego $f\in L^1\cap L^2$ i $t>0$ zachodzi
\begin{equation}\label{eq:dt-cz-weak11}
 \lambda_d\bigl(\{x:|Tf(x)|>t\}\bigr)
 \le \frac{C}{t}\,\|f\|_{L^1(\R^d)}.
\end{equation}
\end{twierdzenie}
\begin{proof}
Stosujemy lemat~\ref{lem:dt-cz-cz-decomp} z $s=t$.
W dowodzie litera $C$ może oznaczać różne dodatnie
stałe zależne wyłącznie od podanych parametrów.
Niech $Q_k^*$ oznacza domkniętą kostkę o tym samym środku
$c_k$ co $Q_k$ i boku $3\sqrt d\,\ell_k$,
gdzie $\ell_k$ jest bokiem $Q_k$. Dla
$E=\bigcup_k Q_k^*$ mamy
$|E|\le(3\sqrt d)^d\sum_k|Q_k|\le C\|f\|_1/t$.
Dla części dobrej nierówność Czebyszewa i
ograniczoność $T$ na $L^2$ dają
\[
 \lambda_d(\{|Tg|>t/2\})
 \le 4t^{-2}\|Tg\|_2^2
 \le Ct^{-1}\|f\|_1.
\]

Uzasadnijmy reprezentację przedłużonego operatora $T$
dla $b_k$, które nie muszą być gładkie.
Wygładzamy $b_k$ standardowymi jądrami
$\rho_\varepsilon$ o nośnikach w $B_\varepsilon(0)$.
Funkcje $b_k*\rho_\varepsilon\in C_c^\infty$
zbiegają do $b_k$ w $L^1$ i $L^2$; wynika to
z twierdzenia~\ref{thm:dt-approx-lp} na ustalonym
zwartym otoczeniu $\overline{Q_k}$, gdyż wszystkie
nośniki są w nim zawarte dla małych $\varepsilon$.
Z ciągłości $T$ na $L^2$ i wyboru podciągu
$T(b_k*\rho_\varepsilon)\to Tb_k$ prawie wszędzie.
Dla ustalonego $x\notin Q_k^*$ odległość od
$\overline{Q_k}$ jest dodatnia; dla małych
$\varepsilon$ jądro $K(x-y)$ jest ograniczone
na wspólnym zwartym otoczeniu nośników.
Zbieżność w $L^1$ daje tam granicę zwykłych całek.
Stąd dla prawie każdego $x\notin Q_k^*$, dzięki zerowej
całce $b_k$,
\[
 Tb_k(x)=\int_{Q_k}[K(x-y)-K(x-c_k)]
 b_k(y)\,\mathrm d\lambda_d(y).
\]
Dla $y\in Q_k$ mamy $|y-c_k|\le\sqrt d\,\ell_k/2$,
a dla $x\notin Q_k^*$ zachodzi
$|x-c_k|\ge3\sqrt d\,\ell_k/2$.
Każdy punkt odcinka od $x-c_k$ do $x-y$ ma więc
moduł co najmniej $2|x-c_k|/3$. Stosując podstawowe
twierdzenie rachunku całkowego wzdłuż odcinka oraz
oszacowanie gradientu jądra, dostajemy
\[
 |Tb_k(x)|\le C\int_{Q_k}
 \frac{|y-c_k|\,|b_k(y)|}{|x-c_k|^{d+1}}
 \,\mathrm d\lambda_d(y)
 \le C\ell_k |x-c_k|^{-d-1}\|b_k\|_1.
\]
Całkując i powiększając zbiór całkowania do
$\{|x-c_k|\ge3\sqrt d\,\ell_k/2\}$, otrzymujemy
\[
 \int_{\R^d\setminus Q_k^*}|Tb_k|\,\mathrm d\lambda_d
 \le C\ell_k\|b_k\|_1
 \int_{3\sqrt d\,\ell_k/2}^\infty r^{-2}\,\mathrm d\lambda_1(r)
 \le C\|b_k\|_1.
\]
Dla skończonych sum nierówność trójkąta daje
$\int_{\R^d\setminus E}|T(\sum_{k\le N}b_k)|
\,\mathrm d\lambda_d\le C\sum_k\|b_k\|_1$.
Ponieważ $\sum_{k\le N}b_k\to f-g$ w $L^2$,
obrazy zbiegają w $L^2$ do $T(f-g)$.
Wybieramy podciąg zbieżny prawie wszędzie
i stosujemy lemat Fatou. Otrzymujemy
\[
 \int_{\R^d\setminus E}|T(f-g)|\,\mathrm d\lambda_d
 \le C\sum_k\|b_k\|_1\le C\|f\|_1.
\]
Z liniowości $Tf=Tg+T(f-g)$ prawie wszędzie.
Zbiór $\{|Tf|>t\}$ zawiera się z dokładnością
do zbioru miary zero w sumie $E$,
$\{|Tg|>t/2\}$ i
$\{x\notin E:|T(f-g)(x)|>t/2\}$.
Ostatnią miarę szacujemy przez $2/t$ razy
poprzednią całkę. To kończy dowód.
\end{proof}

\subsection{Interpolacja i oszacowania w przestrzeniach Lebesgue'a}
\begin{twierdzenie}[Ograniczoność w $L^p$ dla $1<p<\infty$]
\label{thm:dt-cz-general-lp}
Przy hipotezach twierdzenia \ref{thm:dt-cz-weak11}
operator $T$ ma dla każdego $1<p<\infty$ jednoznaczne
ograniczone przedłużenie z $L^2\cap L^p$ na $L^p$,
zgodne z działaniem na $L^2$ na wspólnej dziedzinie.
\end{twierdzenie}
\begin{proof}
Dla $1<p<2$ wynik daje lemat
\ref{lem:dt-cz-interp} zastosowany do oszacowania
\ref{thm:dt-cz-weak11}; dla $p=2$ jest to założenie.
Dla $2<p<\infty$ oznaczmy $q=p/(p-1)$.
Operator sprzężony $T^*$ na $L^2$ jest ograniczony
i spełnia hipotezy twierdzenia \ref{thm:dt-cz-weak11}
z jądrem $K^*(z)=\overline{K(-z)}$.
Istotnie, dla testów $f,h$ o rozłącznych nośnikach
z definicji sprzężenia oraz reprezentacji $Th$ mamy
\[
 \langle T^*f,h\rangle_2
 =\int_{\R^d}\int_{\R^d}f(y)
 \overline{K(y-x)}\,\overline{h(x)}
 \,\mathrm d\lambda_d(y)\,\mathrm d\lambda_d(x).
\]
Twierdzenie Fubiniego ma zastosowanie, gdyż nośniki
są zwarte i oddalone od siebie o dodatnią odległość.
Dla ustalonego $f$ równość zachodzi dla każdego testu
$h$ o nośniku w otwartym dopełnieniu $\supp f$.
Na tym zbiorze funkcja
$x\mapsto\int K^*(x-y)f(y)\,\mathrm d\lambda_d(y)$
jest lokalnie całkowalna, ponieważ na każdym jego
zwartym podzbiorze jądro jest jednostajnie ograniczone
dla $y\in\supp f$. Z jednoznaczności dystrybucji
regularnych (stwierdzenie \ref{prop:regular-injective19})
wynika żądana reprezentacja prawie wszędzie.
Oszacowania $K^*$ i jego gradientu wynikają
bezpośrednio z odpowiednich oszacowań $K$.
Wobec tego $T^*$ jest ograniczony w $L^q$
na mocy pierwszej części dowodu.

Dla $f\in L^2\cap L^p$ oraz $h\in L^2\cap L^q$
otrzymujemy
\[
 |\langle Tf,h\rangle_2|
 =|\langle f,T^*h\rangle_2|
 \le\|f\|_p\|T^*h\|_q
 \le C_q\|f\|_p\|h\|_q.
\]
Aby nie zakładać z góry przynależności $Tf$ do $L^p$,
połóżmy $v=Tf\in L^2$ i
$E_n=B_n(0)\cap\{|v|\le n\}$. Funkcja
$h_n=\chi_{E_n}|v|^{p-2}v$ należy do $L^2\cap L^q$,
gdyż jest ograniczona i ma nośnik skończonej miary.
Dla $a_n=\int_{E_n}|v|^p\,\mathrm d\lambda_d$ mamy
$\langle v,h_n\rangle_2=a_n$ oraz $\|h_n\|_q=a_n^{1/q}$.
Oszacowanie daje $a_n^{1/p}\le C_q\|f\|_p$
(gdy $a_n=0$, nierówność jest automatyczna).
Z twierdzenia o zbieżności monotonicznej dostajemy
$\|Tf\|_p\le C_q\|f\|_p$.
Gęstość $L^2\cap L^p$ w $L^p$ wynika z obcinania
do kul i ograniczonych wartości. Zupełność $L^p$
daje jednoznaczne przedłużenie; jest ono zgodne
na $L^2\cap L^p$ z konstrukcji.
\end{proof}

\begin{uwaga}
Powyższe twierdzenie dotyczy przedłużenia operatora,
który już jest ograniczony w $L^2$ i ma reprezentację
jądrową poza nośnikiem. Nie wynika z niego automatycznie
istnienie granicy głównych wartości prawie wszędzie
dla dowolnej funkcji $L^p$. Do tego potrzebna jest
osobna kontrola obciętych operatorów.
\end{uwaga}

\subsection{Zastosowanie do potencjału Newtona i przypadki końcowe}
\begin{twierdzenie}[Oszacowanie Calderóna--Zygmunda dla laplasjanu]
\label{thm:dt-cz-cz}
Niech $d\ge2$ i $1<p<\infty$. Dla
$f\in C_c^\infty(\R^d)$ oraz $u=\Gamma*f$ zachodzi
\begin{equation}\label{eq:dt-cz-cz}
 \|D^2u\|_{L^p(\R^d)}\le C_{d,p}\|f\|_{L^p(\R^d)}.
\end{equation}
Oszacowanie przedłuża operator $f\mapsto D^2(\Gamma*f)$
do $L^p$; nie oznacza ono, że samo $u$ musi należeć
do $L^p(\R^d)$.
\end{twierdzenie}
\begin{proof}
Jądra $K_{ij}$ i operatory $T_{ij}$ spełniają hipotezy
twierdzenia~\ref{thm:dt-cz-weak11}: oszacowania jądra
wykazaliśmy przy wzorze \eqref{eq:dt-cz-kernel}, reprezentacja
poza nośnikiem wynika z definicji głównej wartości,
a ograniczoność w $L^2$ z lematu~\ref{lem:dt-cz-l2}.
Z ogólnego twierdzenia~\ref{thm:dt-cz-general-lp}
wynika więc, dla każdego $1<p<\infty$,
$\|T_{ij}f\|_p\le C_{d,p}\|f\|_p$.
Stałą możemy wybrać wspólną dla wszystkich $i,j$,
ponieważ dla ustalonego $d$ jest tylko skończenie wiele par.
Ze wzoru~\eqref{eq:dt-cz-repr} mamy
$\|\partial_{ij}u\|_p\le(C_{d,p}+1/d)\|f\|_p$.
Norma Frobeniusa hesjanu jest punktowo nie większa
niż suma modułów jego $d^2$ składowych, więc nierówność
trójkąta w $L^p$ daje~\eqref{eq:dt-cz-cz}.
Gęstość testów (twierdzenie~\ref{thm:dt-test-density})
i zupełność przestrzeni pól macierzowych $L^p$
zapewniają jednoznaczne ograniczone przedłużenie.
\end{proof}

\begin{uwaga}[Dlaczego przypadki $1$ i $\infty$ są osobne?]
\label{rem:dt-cz-endpoints}
Silne oszacowanie $L^1\to L^1$ jest fałszywe nawet
dla gładkich danych. Istotnie, wybierzmy
$f\in C_c^\infty(B_1(0))$, $f\ge0$, $\int f\,\mathrm d\lambda_d=1$.
Dla $|x|>2$ mamy
\[
 |T_{12}f(x)-K_{12}(x)|
 \le C_d|x|^{-d-1}\int |y|f(y)\,\mathrm d\lambda_d(y)
 \le C_d|x|^{-d-1}.
\]
Wynika to z oszacowania gradientu jądra na odcinku
od $x$ do $x-y$, którego punkty mają moduł co najmniej
$|x|/2$. Na pewnym otwartym fragmencie sfery wokół
$(e_1+e_2)/\sqrt2$ mamy $K_{12}(\theta)\ge c_d>0$.
Dzięki jednorodności, dla odpowiednio dużego $r$ i takich
$\theta$, $T_{12}f(r\theta)\ge(c_d/2)r^{-d}$.
Całka po tym stożku jest rozbieżna, ponieważ
$\int_R^\infty r^{-1}\,\mathrm d\lambda_1(r)=\infty$.
Ponieważ $\partial_{12}(\Gamma*f)=T_{12}f$,
nie zachodzi też silne oszacowanie Hessianu w $L^1$.
Oszacowanie~\eqref{eq:dt-cz-weak11} jest słabszym
oszacowaniem końcowym.

Dla końca $p=\infty$ w wymiarze $d=2$ weźmy
$f=\chi_{(0,1)^2}$ i $u=\Gamma*f$.
Poza domkniętym kwadratem wolno różniczkować pod całką
(lokalnie odległość od nośnika jest dodatnia).
Dla $x=(-t,-t)$, $0<t<1/8$, otrzymujemy
\[
 \partial_{12}u(x)=\frac1\pi
 \int_{(0,1)^2}
 \frac{(t+y_1)(t+y_2)}
 {((t+y_1)^2+(t+y_2)^2)^2}\,\mathrm d\lambda_2(y).
\]
Całkowana funkcja jest nieujemna. Na podsektorze
$\pi/6<\theta<\pi/3$, $2t<r<1/2$, gdzie
$y=(r\cos\theta,r\sin\theta)$, obie współrzędne $y$
są co najmniej $r/2$, a $|x-y|\le r+\sqrt2t\le2r$.
Stąd całkowana funkcja jest nie mniejsza niż $c r^{-2}$,
z $c>0$ niezależnym od $t$. Zatem
$\partial_{12}u(-t,-t)\ge c'\log(1/(4t))\to+\infty$.
Pochodna jest ciągła w otwartym dopełnieniu kwadratu,
więc każda dostatecznie duża wartość daje także
przekroczenie zadanego progu na otwartym zbiorze
dodatniej miary. Jest to brak ograniczoności istotnej,
a nie jedynie rozbieżność na ciągu punktów.
Przykład można przenieść na gładkie dane: wygładzenia
$f_\varepsilon=f*\rho_\varepsilon$ mają normę
$L^\infty$ nie większą niż $1$, a w każdym ustalonym
punkcie poza domkniętym kwadratem
$\partial_{12}(\Gamma*f_\varepsilon)$ zbiega
do powyższej całki dzięki zbieżności w $L^1$
i temu, że dla dostatecznie małych $\varepsilon$
wszystkie nośniki leżą w ustalonym zwartym otoczeniu
kwadratu, oddalonym od rozważanego punktu.
Jednostajne oszacowanie $L^\infty$ dla gładkich danych
prowadziłoby zatem do sprzeczności.
\end{uwaga}

\setcounter{section}{49}

\section{Potencjały Bessela i przestrzenie \texorpdfstring{$H^s$}{Hs}}
\label{sec:dt-bessel}
\providecommand{\atlasref}[2]{#2}

Ujednolicamy normalizację przestrzeni Bessela z rozdziałem
o dystrybucjach temperowanych. Por.
\atlasref{PDE-08::thm:pde8-Hs}{RRC-08: potencjały Bessela}.


Dla $s\in\R$ połóżmy $b_s(\xi)=(1+|\xi|^2)^{s/2}$.
Najpierw uzasadnimy, że $b_s$ jest dopuszczalnym
mnożnikiem dystrybucji temperowanych.
Każda jego pochodna jest skończoną sumą wyrazów
postaci
\[
 P(\xi)(1+|\xi|^2)^{s/2-k},
 \qquad k\in\N_0,
\]
gdzie $P$ jest wielomianem. Dla pochodnej rzędu
zero mamy $P=1$, $k=0$. Różniczkowanie takiego
wyrazu względem $\xi_j$ daje
\[
 (\partial_j P)b_{s-2k}
 +(s-2k)\xi_j P b_{s-2k-2},
\]
co uzasadnia indukcję.
Dla dowolnego rzeczywistego $a$ mamy
$(1+|\xi|^2)^{a/2}\le(1+|\xi|)^{\max\{a,0\}}$:
dla $a\le0$ lewa strona nie przekracza $1$,
a dla $a>0$ korzystamy z
$\sqrt{1+|\xi|^2}\le1+|\xi|$.
Każda pochodna $b_s$ ma zatem wzrost co najwyżej
wielomianowy (można zwiększyć wykładnik do
liczby całkowitej). Na mocy twierdzenia~\ref{thm:tempered-operators22}
o mnożnikach dystrybucji temperowanych
mnożenie przez $b_s$ działa na $\mathcal S$
i na $\mathcal S'$.

\begin{definicja}[Przestrzeń Sobolewa typu Bessela]
\label{def:dt-Hs}
Dla dowolnego $s\in\R$ definiujemy
\[
 H^s(\R^d)=\bigl\{u\in\mathcal S'(\R^d):
 b_s\widehat u=T_h\text{ dla pewnego }
 h\in L^2(\R^d)\bigr\}.
\]
Z liniowości transformaty i mnożenia przez $b_s$
oraz liniowości odwzorowania $h\mapsto T_h$
wynika, że $H^s$ jest zespoloną podprzestrzenią
wektorową $\mathcal S'$.
Funkcja $h$ jest określona jednoznacznie jako
klasa $L^2$ dzięki iniektywności odwzorowania
funkcji lokalnie całkowalnych w dystrybucje.
Definiujemy
$\|u\|_{H^s}=\|h\|_2$.
Mnożąc równość $b_s\widehat u=T_h$ przez
$b_{-s}$, otrzymujemy dla $\varphi\in\mathcal S$
\[
 \langle\widehat u,\varphi\rangle
 =\langle T_h,b_{-s}\varphi\rangle
 =\int_{\R^d}b_{-s}(\xi)h(\xi)\varphi(\xi)
 \,\mathrm d\lambda_d(\xi).
\]
Całka jest bezwzględnie zbieżna z nierówności
Cauchy'ego--Schwarza, ponieważ $h\in L^2$
i $b_{-s}\varphi\in\mathcal S\subset L^2$.
Zatem $\widehat u=T_{b_{-s}h}$ i możemy zapisać
\[
 \|u\|_{H^s}^2=\int_{\R^d}
 (1+|\xi|^2)^s|\widehat u(\xi)|^2\,
 \mathrm d\lambda_d(\xi),
\]
gdzie pod całką $\widehat u$ oznacza funkcję
$b_{-s}h$, lokalnie należącą do $L^2$.
Iloczyn jest lokalnie kwadratowalny, ponieważ
$b_{-s}$ jest ograniczony na każdym zbiorze zwartym.
Operator
\[
 J^s u=(I-\Delta)^{s/2}u
 :=\mathcal F^{-1}(b_s\widehat u)
\]
określamy na $\mathcal S'$ przez mnożnik Fouriera,
a nie przez formalne potęgowanie operatora
różniczkowego.
Dla ujemnego $s$ element $H^s$ może być
dystrybucją nieregularną; wymóg $L^2$ dotyczy
ważonej transformaty, a nie samego $u$.
\end{definicja}

Z równości $b_sb_t=b_{s+t}$ wynika
$J^sJ^t=J^{s+t}$ i $J^0=I$, a więc
$(J^s)^{-1}=J^{-s}$ na $\mathcal S'$.
Istotnie, mnożenie dystrybucji przez dwa takie
symbole odpowiada mnożeniu testu przez ich iloczyn.
Dla $s=2$ mamy $b_2=1+|\xi|^2$, toteż
$J^2=I-\Delta$ dzięki twierdzeniu
~\ref{thm:fourier-auto-tempered22}.

Dla $m\in\N_0$ zapis $W^{m,2}(\R^d)$ oznacza
przestrzeń funkcji $u\in L^2(\R^d)$, których
wszystkie pochodne słabe $\partial^\alpha u$,
$|\alpha|\le m$, należą do $L^2$.
Przyjmujemy normę
\[
 \|u\|_{W^{m,2}}^2=\sum_{|\alpha|\le m}
 \|\partial^\alpha u\|_2^2,
\]
włączając $\partial^0u=u$.

\begin{uwaga}[Szeregi Fouriera i transformata na $\R^d$]
W rozdziale \ref{sec:dt-series-fourier} na okręgu używamy
znormalizowanej miary Haara. Tutaj stosujemy miarę
Lebesgue'a i unitarną transformatę Fouriera z definicji
\ref{def:fourier-schwartz17}; w normie $H^s$ nie występuje
dodatkowy czynnik normalizujący.
\end{uwaga}

\begin{twierdzenie}[Podstawowe własności $H^s$]
\label{thm:dt-Hs}
Dla dowolnych $s,t\in\R$ zachodzą:
\begin{enumerate}
\item $H^s(\R^d)$ jest przestrzenią Hilberta, a
 $J^s:H^s\to L^2$ jest izometrycznym izomorfizmem
 na całą przestrzeń $L^2$.
\item Jeśli $s\ge t$, to $H^s\hookrightarrow H^t$
 ciągle i $\|u\|_{H^t}\le\|u\|_{H^s}$.
\item Dla $\alpha\in\N_0^d$, $\partial^\alpha:H^s\to
 H^{s-|\alpha|}$ jest ograniczony, z normą
 operatorową co najwyżej $1$.
\item Dla $m\in\N_0$ zachodzi $H^m(\R^d)=W^{m,2}(\R^d)$
 ze wzajemnie równoważnymi normami, po utożsamieniu
 funkcji z dystrybucjami regularnymi.
\item $\mathcal S(\R^d)$ jest gęsta w $H^s(\R^d)$.
 Ponadto $J^t:H^s\to H^{s-t}$ jest izometrycznym
 izomorfizmem na całą przestrzeń $H^{s-t}$.
\end{enumerate}
\end{twierdzenie}
\begin{proof}
Niech $u\in H^s$ i $b_s\widehat u=T_h$,
gdzie $h\in L^2$. Niech $v=\mathcal F^{-1}h$
w sensie $L^2$. Zgodność transformat na $L^2$
i $\mathcal S'$ z twierdzenia~\ref{thm:plancherel22}
daje $\mathcal FT_v=T_h$, a odwracalność
na $\mathcal S'$ daje $J^su=T_v$.
Wobec tego $J^s$ przyjmuje wartości w $L^2$
i z twierdzenia~\ref{thm:plancherel22}
otrzymujemy
\[
 \|J^s u\|_2=\|b_s\widehat u\|_2
 =\|u\|_{H^s}.
\]
Zapis normy pośrodku oznacza normę reprezentanta
$h$ z definicji. Dla dowolnego $v\in L^2$
traktujemy $v$ jako dystrybucję temperowaną
i określamy $u=J^{-s}v$. Wówczas
$b_s\widehat u=\widehat v\in L^2$, a $J^su=v$.
Iniektywność $J^s$ wynika z istnienia $J^{-s}$.
Iloczyn skalarny definiujemy przez
\[
 \langle u,v\rangle_{H^s}
 =\langle J^s u,J^s v\rangle_2.
\]
Liniowość $J^s$ i własności iloczynu skalarnego
na $L^2$ dają liniowość względem pierwszego
argumentu, antyliniowość względem drugiego
i symetrię hermitowską.
Ponadto $\langle u,u\rangle_{H^s}=\|J^su\|_2^2$
jest nieujemne, a równość z zerem daje $J^su=0$
i stąd $u=0$ dzięki iniektywności.
Iloczyn skalarny wyznacza podaną normę.
Jeśli $(u_n)$ jest ciągiem Cauchy'ego w $H^s$,
to $(J^su_n)$ jest ciągiem Cauchy'ego w $L^2$
i ma tam granicę $v$. Dla $u=J^{-s}v$ mamy
$\|u_n-u\|_{H^s}=\|J^s u_n-v\|_2\to0$.
Stąd zupełność i pierwsza teza.

Jeśli $s\ge t$, to $0<b_{t-s}\le1$ i
$b_t\widehat u=b_{t-s}(b_s\widehat u)$.
Prawa strona jest reprezentowana przez funkcję
$L^2$, a jej norma jest nie większa niż norma
$b_s\widehat u$. Daje to drugą tezę.

Dla dowolnego wielowskaźnika $\alpha$ mamy
na mocy twierdzenia~\ref{thm:fourier-auto-tempered22}
\[
 b_{s-|\alpha|}\mathcal F(\partial^\alpha u)
 =b_{-|\alpha|}(i\xi)^\alpha b_s\widehat u.
\]
Moduł mnożnika po prawej stronie nie przekracza
$1$, ponieważ
$|\xi^\alpha|\le(1+|\xi|^2)^{|\alpha|/2}$.
Zatem $\partial^\alpha u\in H^{s-|\alpha|}$
i $\|\partial^\alpha u\|_{H^{s-|\alpha|}}
\le\|u\|_{H^s}$.

Niech teraz $m\in\N_0$. Z rozwinięcia
wielomianowego otrzymujemy
\[
 (1+|\xi|^2)^m
 =\sum_{|\alpha|\le m}c_\alpha|\xi^\alpha|^2,
 \qquad c_\alpha=\frac{m!}{(m-|\alpha|)!\alpha!}.
\]
Każde $c_\alpha$ jest dodatnią liczbą całkowitą,
a z podstawienia $\xi=(1,\ldots,1)$ wynika
$\sum_{|\alpha|\le m}c_\alpha=(d+1)^m$.
W szczególności
\[
 \sum_{|\alpha|\le m}|\xi^\alpha|^2
 \le (1+|\xi|^2)^m
 \le (d+1)^m\sum_{|\alpha|\le m}|\xi^\alpha|^2.
\]
Jeśli $u\in H^m$, to z drugiej i trzeciej tezy
wynika, że $\partial^\alpha u\in H^{m-|\alpha|}$
dla $|\alpha|\le m$ (włącznie z $\alpha=0$).
Ponieważ $m-|\alpha|\ge0$, drugą tezę
stosujemy z wykładnikami $m-|\alpha|$ i $0$.
Zatem wszystkie te pochodne należą do $H^0=L^2$.
Utożsamienie $H^0=L^2$ wynika z pierwszej tezy
dla $s=0$, ponieważ $J^0=I$.
Po ograniczeniu do testów o zwartym nośniku
pochodne dystrybucyjne reprezentowane przez
funkcje $L^2$ są pochodnymi słabymi w sensie
rozdziału~\ref{sec:pochodne20}.
A zatem $u\in W^{m,2}$.

Odwrotnie, dla $u\in W^{m,2}$ każda słaba
pochodna $v_\alpha=\partial^\alpha u$ należy
do $L^2$. Równość dystrybucyjna
$\partial^\alpha T_u=T_{v_\alpha}$ zachodzi
na $\mathcal D$, więc zachodzi też na $\mathcal S$:
obie strony są temperowane, a ograniczenie
$\mathcal S'\to\mathcal D'$ jest iniektywne.
Po transformacji dostajemy
$(i\xi)^\alpha\widehat u=\widehat v_\alpha$
jako dystrybucje temperowane.
Lewą stronę reprezentuje funkcja lokalnie
całkowalna, ponieważ $\widehat u\in L^2$
i wielomian jest ograniczony na zwartych zbiorach.
Z iniektywności dystrybucji regularnych
równość obowiązuje prawie wszędzie.
Stąd z Plancherela i skończonej sumy
\[
 \int_{\R^d}
 \left(\sum_{|\alpha|\le m}|\xi^\alpha|^2\right)
 |\widehat u(\xi)|^2\,\mathrm d\lambda_d(\xi)
 =\sum_{|\alpha|\le m}\|v_\alpha\|_2^2<\infty.
\]
Porównanie wielomianów daje $u\in H^m$.
Ten sam rachunek obowiązuje dla $u\in H^m$,
ponieważ wykazaliśmy już przynależność
do $W^{m,2}$. W obu kierunkach otrzymujemy
\[
 \|u\|_{W^{m,2}}\le\|u\|_{H^m}
 \le(d+1)^{m/2}\|u\|_{W^{m,2}}.
\]

Na koniec $J^s$ i $J^{-s}$ odwzorowują
$\mathcal S$ w siebie, ponieważ transformata
i odwrotność zachowują $\mathcal S$,
a symbole $b_s,b_{-s}$ są dopuszczalnymi
mnożnikami. Każda funkcja Schwartza należy
do $H^s$, gdyż $b_s\widehat u\in\mathcal S\subset L^2$.
Dla $u\in H^s$ wybieramy $v_n\in\mathcal S$
zbiegające do $J^su$ w $L^2$ na mocy
gęstości $\mathcal S$ w $L^2$ wykorzystanej w dowodzie
twierdzenia~\ref{thm:plancherel22}. Wtedy
$u_n=J^{-s}v_n\in\mathcal S$ i
$\|u_n-u\|_{H^s}=\|v_n-J^su\|_2\to0$.
Dla dowolnego $t$ mamy ponadto
$b_{s-t}\mathcal F(J^tu)=b_s\widehat u$,
więc $\|J^tu\|_{H^{s-t}}=\|u\|_{H^s}$.
Dla $v\in H^{s-t}$ stosujemy tę samą tożsamość
z wykładnikami $s-t$ i $-t$, otrzymując
$J^{-t}v\in H^s$.
Ponieważ $J^tJ^{-t}v=v$ i $J^{-t}J^tu=u$,
operator $J^{-t}$ jest odwrotnością na całym
$H^{s-t}$.
\end{proof}

\begin{wniosek}[Kiedy delta należy do ujemnej przestrzeni Sobolewa?]
\label{cor:dt-delta-Hs}
Dla $s\in\R$ zachodzi
\[
 \delta_0\in H^s(\R^d)\quad\Longleftrightarrow\quad
 s<-\frac d2.
\]
\end{wniosek}
\begin{proof}
Z przykładu~\ref{ex:fourier-delta22} mamy
$\widehat{\delta_0}=(2\pi)^{-d/2}$ jako dystrybucję regularną.
Wobec tego $b_s\widehat{\delta_0}$ jest
reprezentowana przez funkcję $(2\pi)^{-d/2}b_s$.
Stały niezerowy czynnik nie wpływa na warunek
kwadratowalności.
Z iniektywności dystrybucji regularnych
warunek przynależności do $H^s$ jest równoważny
\[
 \int_{\R^d}(1+|\xi|^2)^s\,
 \mathrm d\lambda_d(\xi)<\infty.
\]
Na domkniętej kuli jednostkowej funkcja podcałkowa
jest ciągła i ograniczona. Dla $r\ge1$ zachodzi
\[
 \min\{1,2^s\}r^{2s}\le(1+r^2)^s
 \le\max\{1,2^s\}r^{2s}.
\]
Dla $d\ge2$ całkowanie radialne poza kulą
sprowadza zatem warunek do skończoności
$\int_1^\infty r^{2s+d-1}\,\mathrm d\lambda_1(r)$;
dla $d=1$ całkujemy osobno na obu półprostych
i otrzymujemy ten sam warunek z $d=1$.
Całka potęgowa jest skończona dokładnie wtedy,
gdy $2s+d-1<-1$, czyli $s<-d/2$.
Dla $s=-d/2$ jest to całka z $r^{-1}$,
rozbieżna logarytmicznie. W szczególności
$H^s$ dla $s<-d/2$ zawiera dystrybucję
niereprezentowaną przez funkcję lokalnie całkowalną:
nieregularność $\delta_0$ wykazano w przykładzie
~\ref{ex:delta19}.
\end{proof}


\subsection{Zasada nieoznaczoności}
\label{sec:dt-uncertainty}
\providecommand{\atlasref}[2]{#2}

Por. \atlasref{PDE-08::thm:pde8-uncertainty}{RRC-08: nierówność Heisenberga}.


\begin{twierdzenie}[Nierówność Heisenberga w jednej współrzędnej]
\label{thm:dt-uncertainty}
Jeśli $0\ne f\in\mathcal S(\R^d)$ i $1\le j\le d$, to
\[
 \|x_jf\|_{L^2}\,\|\xi_j\widehat f\|_{L^2}
 \ge\frac12\,\|f\|_{L^2}\,\|\widehat f\|_{L^2}.
\]
Stała $1/2$ jest optymalna. Równość zachodzi
dokładnie wtedy, gdy dla pewnego $\lambda>0$
funkcja $f$ ma postać
\[
 f(x)=e^{-\lambda x_j^2/2}\,g(x'),
\]
gdzie dla $d>1$ przez $x'$ oznaczamy pozostałe
$d-1$ współrzędnych, a $g\in\mathcal S(\R^{d-1})$
jest niezerowa. Dla $d=1$ czynnik $g(x')$
zastępujemy niezerową stałą zespoloną.
\end{twierdzenie}
\begin{proof}
Najpierw zauważmy, że $|f|^2=f\overline f$,
a także $x_j|f|^2$ należą do $\mathcal S$.
Sprzężenie zespolone nie zmienia półnorm.
Dla $u,v\in\mathcal S$ wzór Leibniza daje
\[
 Q_{N,m}(uv)\le 2^m Q_{N,m}(u)Q_{0,m}(v)
 \qquad(N,m\in\N_0),
\]
gdyż każda pochodna iloczynu jest sumą składników
$(\partial^\gamma u)(\partial^{\alpha-\gamma}v)$
ze współczynnikami dwumianowymi o sumie $2^{|\alpha|}\le2^m$.
Iloczyn dwóch funkcji Schwartza należy zatem do
$\mathcal S$, a mnożenie przez współrzędną zachowuje
$\mathcal S$ na mocy stwierdzenia~\ref{prop:schwartz-operators17}.
W szczególności $\partial_j(x_j|f|^2)\in L^1$.
Dla ustalonych pozostałych współrzędnych
podstawowe twierdzenie rachunku całkowego daje
\[
 \int_{-R}^R\partial_j(x_j|f(x)|^2)
 \,\mathrm d\lambda_1(x_j)
 =\left[x_j|f(x)|^2\right]_{x_j=-R}^{x_j=R}.
\]
Prawa strona dąży do zera dla $R\to\infty$,
ponieważ dla $|x_j|=R$ mamy
\[
 |x_j||f(x)|^2
 \le R\,q_{1,0}(f)^2(1+R)^{-2}.
\]
Dla $d>1$ dzięki całkowalności pochodnej
na $\R^d$ twierdzenie Fubiniego zapewnia, że
dla prawie każdego $x'$ pochodna ta jest
bezwzględnie całkowalna na prostej zmiennej $x_j$.
Na takim przekroju całka po $[-R,R]$ zbiega
do całki po $\R$ z twierdzenia o zbieżności
zmajoryzowanej, z majorantą
$|\partial_j(x_j|f|^2)|$ na przekroju.
Z oszacowania składnika brzegowego wynika,
że ta całka po $\R$ jest zerowa. Fubini
pozwala teraz całkować tę równość względem $x'$.
Dla $d=1$ przechodzimy do granicy bezpośrednio,
ponieważ pochodna jest w $L^1(\R)$.
Otrzymujemy
\[
 \int_{\R^d}\partial_j(x_j|f|^2)
 \,\mathrm d\lambda_d(x)=0.
\]
Z reguły Leibniza
\[
 \partial_j(x_j|f|^2)
 =|f|^2+2\operatorname{Re}(x_j\overline f\,\partial_j f).
\]
Oba składniki są całkowalne: dla drugiego
wynika to z nierówności Cauchy'ego--Schwarza
dla $x_jf,\partial_j f\in L^2$.
W konsekwencji
\[
 \begin{aligned}
 \|f\|_2^2
 &=-2\operatorname{Re}
 \int_{\R^d}x_j\overline f(x)\,\partial_j f(x)
 \,\mathrm d\lambda_d(x)\\
 &\le 2\left|
 \int_{\R^d}x_j\overline f(x)\,\partial_j f(x)
 \,\mathrm d\lambda_d(x)\right|\\
 &\le 2\|x_j f\|_2\|\partial_j f\|_2.
 \end{aligned}
\]
Z twierdzeń~\ref{thm:fourier-schwartz17}
i \ref{thm:plancherel22} mamy
\[
 \|\partial_j f\|_2=\|\xi_j\widehat f\|_2,
 \qquad
 \|f\|_2=\|\widehat f\|_2.
\]
Po podstawieniu otrzymujemy
$\|f\|_2^2\le2\|x_jf\|_2\|\xi_j\widehat f\|_2$.
Dzieląc przez $2$ i wykorzystując zachowanie normy
przez unitarną transformatę, otrzymujemy tezę.

Ustalmy teraz warunek równości. Ponieważ
$\|f\|_2>0$, powyższe oszacowanie pokazuje,
że zarówno $x_jf$, jak i $\partial_j f$
mają dodatnie normy $L^2$.
Równość w nierówności Heisenberga jest
równoważna równości w obu nierównościach
w powyższym łańcuchu.
Równość w nierówności Cauchy'ego--Schwarza
oznacza liniową zależność obu wektorów w $L^2$:
$\partial_j f=cx_jf$ dla pewnego $c\in\C$.
Istotnie, dla $p=\partial_j f$, $q=x_jf$
oraz $c=\langle p,q\rangle_2/\|q\|_2^2$
mamy
\[
 \|p-cq\|_2^2
 =\|p\|_2^2-\frac{|\langle p,q\rangle_2|^2}{\|q\|_2^2}.
\]
Równość w nierówności Cauchy'ego--Schwarza
daje zatem $p=cq$ w $L^2$, a zależność ta
jest również wystarczająca dla równości.
Wtedy
\[
 \int_{\R^d}x_j\overline f\,\partial_j f
 \,\mathrm d\lambda_d(x)=c\|x_jf\|_2^2.
\]
Równość w pierwszej nierówności wymaga,
by liczba po lewej stronie była rzeczywista
i ujemna (jej część rzeczywista wynosi
$-\|f\|_2^2/2$). Stąd $c=-\lambda$
dla pewnego $\lambda>0$.
Ponieważ $\partial_j f$ i $x_j f$ są ciągłe,
ich równość prawie wszędzie daje równość
w każdym punkcie: niezerowa różnica w punkcie
byłaby niezerowa na pewnej otwartej kuli
o dodatniej mierze.
Dla ustalonych $x'$ otrzymujemy więc
\[
 \partial_j\bigl(e^{\lambda x_j^2/2}f(x)\bigr)=0,
\]
czyli $f(x)=e^{-\lambda x_j^2/2}g(x')$,
gdzie $g(x')=f(x)|_{x_j=0}$.
Dla $d>1$ funkcja $g$ jest gładka jako
ograniczenie funkcji gładkiej do hiperpłaszczyzny.
Jeśli $\alpha',\beta'\in\N_0^{d-1}$,
to przez $\alpha,\beta\in\N_0^d$ oznaczmy
wielowskaźniki powstałe przez wstawienie
zera na $j$-tej pozycji.
Wtedy
\[
 \sup_{x'\in\R^{d-1}}
 |(x')^{\alpha'}\partial^{\beta'}g(x')|
 \le p_{\alpha,\beta}(f)<\infty.
\]
Zatem $g\in\mathcal S(\R^{d-1})$ i $g\ne0$,
bo w przeciwnym razie wzór na $f$ dawałby
$f=0$. Dla $d=1$ tym samym argumentem
otrzymujemy niezerową stałą $f(0)$.

Odwrotnie, dla dowolnej $\lambda>0$ i takiej
niezerowej $g$ funkcja
$f(x)=e^{-\lambda x_j^2/2}g(x')$ należy
do $\mathcal S$. Dla $d>1$, oznaczając
$G(t)=e^{-\lambda t^2/2}$, mamy
\[
 p_{\alpha,\beta}(f)=
 p_{\alpha_j,\beta_j}(G)\,
 p_{\alpha',\beta'}(g)<\infty,
\]
gdzie $\alpha',\beta'$ powstają przez
usunięcie $j$-tej współrzędnej.
Wynika to z rozdzielenia pochodnych i wag
oraz supremum niezależnych nieujemnych czynników.
Dla $d=1$ wystarcza przynależność Gaussa do $\mathcal S$.
Niezerowość $f$ wynika z niezerowości $g$.
Spełnia $\partial_j f=-\lambda x_j f$,
więc obie nierówności w dowodzie stają się
równościami. Takie funkcje istnieją w każdym
wymiarze, na przykład $f(x)=e^{-|x|^2/2}$.
Zatem stałej $1/2$ nie można zwiększyć.
\end{proof}

\begin{wniosek}[Nierówność z przesuniętymi środkami]
Dla $f\in\mathcal S(\R^d)$ i dowolnych $a,b\in\R$
mamy
\[
 \|(x_j-a)f\|_2\,\|(\xi_j-b)\widehat f\|_2
 \ge\tfrac12\|f\|_2\|\widehat f\|_2.
\]
\end{wniosek}
\begin{proof}
Dla $e_j$ oznaczającego $j$-ty wektor
bazy standardowej połóżmy
$h(x)=e^{-ibx_j}f(x+ae_j)$.
Przesunięcie zachowuje $\mathcal S$ przez
nierówność $1+|x|\le(1+|a|)(1+|x+ae_j|)$,
a mnożenie przez $e^{-ibx_j}$ zachowuje
$\mathcal S$, gdyż wszystkie pochodne
mnożnika są ograniczone.
Ponieważ $h\in\mathcal S\subset L^1$,
podstawienie $y=x+ae_j$ w całce Fouriera daje
\[
 \begin{aligned}
 \widehat h(\xi)
 &=(2\pi)^{-d/2}\int_{\R^d}e^{-ix\cdot(\xi+be_j)}
 f(x+ae_j)\,\mathrm d\lambda_d(x)\\
 &=e^{ia(\xi_j+b)}(2\pi)^{-d/2}
 \int_{\R^d}e^{-iy\cdot(\xi+be_j)}f(y)
 \,\mathrm d\lambda_d(y)\\
 &=e^{ia(\xi_j+b)}\widehat f(\xi+be_j).
 \end{aligned}
\]
Czynniki wykładnicze mają moduł $1$,
a miara Lebesgue'a jest niezmiennicza
względem przesunięć. Stąd
\[
 \begin{aligned}
 \|h\|_2&=\|f\|_2,&
 \|x_jh\|_2&=\|(x_j-a)f\|_2,\\
 \|\widehat h\|_2&=\|\widehat f\|_2,&
 \|\xi_j\widehat h\|_2&=\|(\xi_j-b)\widehat f\|_2.
 \end{aligned}
\]
Dla $f\ne0$ stosujemy twierdzenie
~\ref{thm:dt-uncertainty} do $h$;
dla $f=0$ nierówność jest oczywista.
\end{proof}

Dla $f\ne0$ można w szczególności wybrać
\[
 a=\frac{\int_{\R^d}x_j|f(x)|^2\,
 \mathrm d\lambda_d(x)}{\|f\|_2^2},
 \qquad
 b=\frac{\int_{\R^d}\xi_j|\widehat f(\xi)|^2\,
 \mathrm d\lambda_d(\xi)}{\|\widehat f\|_2^2}.
\]
Całki są bezwzględnie zbieżne, ponieważ
\[
 \begin{aligned}
 \int_{\R^d}|x_j||f(x)|^2\,\mathrm d\lambda_d(x)
 &\le\|x_jf\|_2\|f\|_2,\\
 \int_{\R^d}|\xi_j||\widehat f(\xi)|^2\,
 \mathrm d\lambda_d(\xi)
 &\le\|\xi_j\widehat f\|_2\|\widehat f\|_2.
 \end{aligned}
\]
Oba mianowniki są dodatnie z niezerowości $f$
i twierdzenia Plancherela. Liczby $a,b$ są
rzeczywiste, a obie niżej wymienione gęstości
mają całkę $1$ z definicji normy $L^2$.
Drugie momenty są skończone, ponieważ
$x_jf,\xi_j\widehat f\in L^2$.
Jeśli $\sigma_x^2$ i $\sigma_\xi^2$
oznaczają wariancje $j$-tej współrzędnej
względem miar probabilistycznych o gęstościach
$|f|^2/\|f\|_2^2$ oraz
$|\widehat f|^2/\|\widehat f\|_2^2$, to
\[
 \sigma_x=\frac{\|(x_j-a)f\|_2}{\|f\|_2},
 \qquad
 \sigma_\xi=\frac{\|(\xi_j-b)\widehat f\|_2}
 {\|\widehat f\|_2}.
\]
Z poprzedniego wniosku dostajemy
$\sigma_x\sigma_\xi\ge1/2$.
Stała odpowiada konwencji z definicji~\ref{def:fourier-schwartz17},
w której wykładnik transformaty to $-ix\cdot\xi$.

\setcounter{part}{7}
\part{Zaawansowana analiza harmoniczna}
\setcounter{section}{56}
% Kandydat restrukturyzacji; nie jest aktywnym rozdziałem Atlasu.

\section{Splot na grupach lokalnie zwartych}
\label{sec:dt-haar-convolution}
\providecommand{\atlasref}[2]{#2}

Konstrukcję i jednoznaczność miary Haara pozostawiamy
w teorii miary. W tym rozdziale korzystamy z wyników
ME-23 dla lokalnie zwartej grupy Hausdorffa $G$
z lewostronną miarą Haara $\mu$. Nie zakładamy
$\sigma$-skończoności $G$.
Przyjmujemy konwencję funkcji modularnej
$\mu(Ea)=\Delta(a)\mu(E)$, a więc
\[
 \int_G h(x^{-1})\,\mathrm d\mu(x)
 =\int_G h(x)\Delta(x)^{-1}\,\mathrm d\mu(x)
\]
dla funkcji nieujemnych lub w przypadku bezwzględnej
zbieżności odpowiednich całek. Są to wyniki
\atlasref{ME-23::thm:haar-modular-23}{ME-23: modularność}
oraz \atlasref{ME-23::cor:haar-integral-inversion-23}{ME-23: inwersja}.
Por. także \atlasref{ME-23::it:def:haar-convolution}{ME-23: splot na grupie}.

\subsection{Splot funkcji ciągłych o zwartym nośniku}
\begin{definicja}[Splot na grupie]
\label{def:dt-haar-convolution}
Niech $G$ będzie lokalnie zwartą grupą Hausdorffa,
$\mu$ jej lewostronną miarą Haara, a
$\mathbb K=\mathbb R$ lub $\mathbb C$.
Dla $f,g\in C_c(G;\mathbb K)$ definiujemy
\[
 (f*g)(x)=\int_G f(y)g(y^{-1}x)\,\mathrm d\mu(y).
\]
Dla ustalonego $x$ funkcja podcałkowa jest ciągła
i znika poza zwartym nośnikiem $f$, więc całka
istnieje i jest skończona. Na grupie addytywnej
$\mathbb R^d$, po wyborze normalizacji $\mu=\lambda_d$,
otrzymujemy zwykły splot
$(f*g)(x)=\int_{\mathbb R^d}f(y)g(x-y)
\,\mathrm d\lambda_d(y)$.
Dla innej miary Haara $\mu=c\lambda_d$, $c>0$,
splot jest $c$ razy większy od splotu względem
$\lambda_d$. Wynika to z jednoznaczności miary Haara
(twierdzenie \atlasref{ME-23::thm:haar-unique-23}{ME-23: jednoznaczność miary Haara}) i liniowości
całki względem miary.
\end{definicja}

\begin{twierdzenie}[Podstawowe własności splotu]
\label{thm:dt-haar-convolution}
Dla $f,g,h\in C_c(G;\mathbb K)$ splot jest
dwuliniowy, $f*g\in C_c(G;\mathbb K)$ oraz
$\operatorname{supp}(f*g)\subseteq
(\operatorname{supp}f)(\operatorname{supp}g)$.
Ponadto
\[
 \int_G(f*g)\,\mathrm d\mu
 =\left(\int_G f\,\mathrm d\mu\right)
  \left(\int_G g\,\mathrm d\mu\right),
 \qquad
 \|f*g\|_1\le\|f\|_1\|g\|_1,
\]
\[
 \|f*g\|_\infty\le\|f\|_1\|g\|_\infty,\qquad
 (f*g)*h=f*(g*h).
\]
\end{twierdzenie}
\begin{proof}
Dwuliniowość wynika z liniowości całki.
Oznaczmy $K=\operatorname{supp}f$,
$L=\operatorname{supp}g$ i $M=\operatorname{supp}h$.
Jeżeli $f(y)g(y^{-1}x)\ne0$, to $y\in K$,
$y^{-1}x\in L$, a więc $x\in KL$. Zbiór $KL$
jest zwarty jako obraz $K\times L$ przez mnożenie,
a zatem domknięty. Stąd nośnik splotu zawiera się
w $KL$. Dla $a\to e$ mamy
\[
 |(f*g)(xa)-(f*g)(x)|
 \le\|f\|_1\sup_{z\in G}|g(za)-g(z)|\longrightarrow0
\]
na mocy lematu \atlasref{ME-23::lem:haar-uniform-23}{ME-23: jednostajna ciągłość translacji},
który w przypadku zespolonym stosujemy do części
rzeczywistej i urojonej. Ponieważ każde $x'$
ma postać $xa$ z $a=x^{-1}x'$, otrzymujemy
ciągłość splotu w każdym $x$.
Nierówność z normą supremum wynika bezpośrednio
z oszacowania $|g(y^{-1}x)|\le\|g\|_\infty$.

Jądro $(y,x)\mapsto f(y)g(y^{-1}x)$ jest ciągłe,
o nośniku zawartym w zwartym $K\times KL$.
Z lematu \atlasref{ME-23::lem:haar-fubini-compact-23}{ME-23: Fubini dla zwartych nośników}
i lewostronnej niezmienniczości mamy
\[
 \int_G(f*g)(x)\,\mathrm d\mu(x)
 =\int_G f(y)\left(\int_G g(y^{-1}x)
                 \,\mathrm d\mu(x)\right)\mathrm d\mu(y)
 =\int_G f\,\mathrm d\mu\int_G g\,\mathrm d\mu.
\]
Punktowa nierówność $|f*g|\le |f|*|g|$ oraz
powyższy wzór dla funkcji nieujemnych dają
nierówność norm $L^1$.

Dla ustalonego $x$ jądro
$(y,z)\mapsto f(y)g(y^{-1}z)h(z^{-1}x)$
jest ciągłe i znika poza $K\times KL$.
Możemy zatem zamienić kolejność całkowania
tym samym lematem Fubiniego. Po podstawieniu
$z=yu$ w całce względem $z$ dostajemy
\begin{align*}
 ((f*g)*h)(x)
 &=\int_G\int_G f(y)g(y^{-1}z)h(z^{-1}x)
       \,\mathrm d\mu(y)\,\mathrm d\mu(z)\\
 &=\int_G f(y)\int_G g(u)h(u^{-1}y^{-1}x)
       \,\mathrm d\mu(u)\,\mathrm d\mu(y)\\
 &= (f*(g*h))(x).
\end{align*}
Nie zakładamy $\sigma$-skończoności całej grupy.
\end{proof}

\begin{wniosek}[Zmiana postaci splotu]
\label{cor:dt-haar-convolution-change}
W przyjętej konwencji $\mu(Ea)=\Delta(a)\mu(E)$
dla $f,g\in C_c(G;\mathbb K)$ zachodzi
\[
 (f*g)(x)=\int_G f(xz^{-1})g(z)\Delta(z)^{-1}
                  \,\mathrm d\mu(z).
\]
Jeżeli grupa jest przemienna, to $f*g=g*f$.
\end{wniosek}
\begin{proof}
W definicji podstawiamy $y=xu$ (lewa translacja),
co daje
\[
 (f*g)(x)=\int_G f(xu)g(u^{-1})\,\mathrm d\mu(u).
\]
Następnie stosujemy wzór inwersji przytoczony na początku rozdziału
do funkcji $z\mapsto f(xz^{-1})g(z)$. Ma ona zwarty nośnik
zawarty w nośniku $g$, więc obie całki są skończone.
W grupie przemiennej $\Delta=1$ oraz
$xz^{-1}=z^{-1}x$, co daje przemienność splotu.
\end{proof}

\begin{przyklad}[Splot nieprzemienny]
Na skończonej nieprzemiennej grupie $G$ z miarą
liczącą wybierzmy $a,b$ takie, że $ab\ne ba$.
Wówczas $\chi_{\{a\}}*\chi_{\{b\}}=\chi_{\{ab\}}$,
ponieważ jedyny niezerowy składnik sumy odpowiada
$y=a$ i warunkowi $a^{-1}x=b$.
W odwrotnej kolejności otrzymujemy
$\chi_{\{ba\}}$, a zatem inny splot.
\end{przyklad}

\subsection{Algebra splotu i przedłużenia do przestrzeni Lebesgue'a}

\begin{twierdzenie}[Algebra splotu $L^1$]
\label{thm:dt-haar-l1-convolution}
Splot na $C_c(G;\mathbb K)$ ma dokładnie jedno
ograniczone dwuliniowe rozszerzenie do
$L^1(G,\mu;\mathbb K)\times L^1(G,\mu;\mathbb K)$
o wartościach w $L^1(G,\mu;\mathbb K)$.
Rozszerzenie jest łączne i spełnia
$\|f*g\|_1\le\|f\|_1\|g\|_1$.
Zatem $L^1(G,\mu;\mathbb K)$ ze splotem jest
algebrą Banacha, także gdy $\mu$ nie jest
$\sigma$-skończona.
\end{twierdzenie}
\begin{proof}
Miara Haara jest miarą Radona, więc z twierdzenia
\atlasref{ME-20::thm:cc-dense-lp-radon-20}{ME-20: gęstość funkcji ciągłych o zwartym nośniku} dla $p=1$ funkcje
z $C_c(G;\mathbb K)$ są gęste w $L^1(G,\mu;\mathbb K)$.
Twierdzenie to nie wymaga $\sigma$-skończoności. Ich klasy w $L^1$
jednoznacznie wyznaczają funkcje ciągłe:
jeżeli ciągła funkcja $f$ ma normę $L^1$ zero,
to dla każdego $\varepsilon>0$ otwarty zbiór
$\{|f|>\varepsilon\}$ ma miarę zero, a z wniosku
\atlasref{ME-23::cor:haar-open-positive-23}{ME-23: dodatnia miara otwartych zbiorów} musi być pusty.

Wybieramy ciągi $f_n,g_n\in C_c(G;\mathbb K)$
zbieżne w $L^1$ odpowiednio do $f,g$.
Z twierdzenia \ref{thm:dt-haar-convolution}
mamy
\[
 \|f_n*g_n-f_m*g_m\|_1
 \le\|f_n-f_m\|_1\|g_n\|_1+
    \|f_m\|_1\|g_n-g_m\|_1.
\]
Ciągi norm są ograniczone, więc prawa strona
dąży do zera przy $m,n\to\infty$. Z zupełności
$L^1$ (twierdzenie \atlasref{ME-14::thm:riesz-fischer-14}{ME-14: zupełność przestrzeni $L^1$})
istnieje granica, którą oznaczamy $f*g$.
Jeżeli $\widetilde f_n\to f$ i $\widetilde g_n\to g$
w $L^1$ dla innych ciągów z $C_c(G;\mathbb K)$, to
\[
 \|f_n*g_n-\widetilde f_n*\widetilde g_n\|_1
 \le \|f_n-\widetilde f_n\|_1\|g_n\|_1
   +\|\widetilde f_n\|_1\|g_n-\widetilde g_n\|_1
 \longrightarrow0.
\]
To dowodzi niezależności granicy od wyboru ciągów. Przechodząc do granicy
w nierówności norm, otrzymujemy
$\|f*g\|_1\le\|f\|_1\|g\|_1$.
Dwuliniowość wynika z dwuliniowości splotu
na $C_c(G)$ i liniowości granicy.
Jeśli także $h_n\in C_c(G;\mathbb K)$ i $h_n\to h$ w $L^1$,
to z otrzymanej nierówności i dwuliniowości mamy
\begin{align*}
 \|(f_n*g_n)*h_n-(f*g)*h\|_1
 &\le\|f_n*g_n-f*g\|_1\|h_n\|_1
   +\|f*g\|_1\|h_n-h\|_1,\\
 \|f_n*(g_n*h_n)-f*(g*h)\|_1
 &\le\|f_n-f\|_1\|g_n*h_n\|_1
   +\|f\|_1\|g_n*h_n-g*h\|_1.
\end{align*}
Obie prawe strony dążą do zera, ponieważ ciągi norm
$\|h_n\|_1$ oraz $\|g_n*h_n\|_1\le\|g_n\|_1\|h_n\|_1$
są ograniczone.
Równość wyrazów obu ciągów daje łączność.
Każde ograniczone dwuliniowe rozszerzenie
musi zachowywać te granice, co dowodzi
jednoznaczności. Zupełność algebry wynika
z zupełności $L^1$.

W tej ogólności definiujemy splot klas $L^1$
przez rozszerzenie ciągłe, nie przez nieuzasadnione
użycie globalnego twierdzenia Fubiniego.
\end{proof}


\begin{twierdzenie}[Nierówność Younga dla pary $L^1$ i $L^p$]
\label{thm:dt-haar-young}
Niech $1\le p<\infty$. Splot ma jednoznaczne
ograniczone dwuliniowe przedłużenie
$L^1(G,\mu)\times L^p(G,\mu)\to L^p(G,\mu)$
i spełnia
\[
 \|f*g\|_p\le \|f\|_1\|g\|_p.
\]
Dla $p=1$ jest to przedłużenie z twierdzenia
\ref{thm:dt-haar-l1-convolution}.
\end{twierdzenie}
\begin{proof}
Najpierw niech $f,g\in C_c(G;\C)$.
Dla $f=0$ teza jest oczywista. W przeciwnym razie
$\|f\|_1>0$, ponieważ niezerowa ciągła funkcja
jest oddzielona od zera na niepustym otwartym zbiorze.
Dla ustalonego $x$ nierówność Jensena względem
probabilistycznej miary $|f(y)|\,\mathrm d\mu(y)/\|f\|_1$
daje
\[
 |(f*g)(x)|^p
 \le \|f\|_1^{p-1}
 \int_G |f(y)|\,|g(y^{-1}x)|^p\,\mathrm d\mu(y).
\]
Funkcja pod całką podwójną jest ciągła i ma nośnik
w zwartym zbiorze
$\supp f\times(\supp f)(\supp g)$.
Możemy więc zastosować twierdzenie Fubiniego
dla zwartych nośników z ME-23.
Lewostronna niezmienniczość $\mu$ daje
\[
 \begin{aligned}
 \int_G |f*g|^p\,\mathrm d\mu
 &\le \|f\|_1^{p-1}\int_G|f(y)|
       \int_G |g(y^{-1}x)|^p\,\mathrm d\mu(x)
       \,\mathrm d\mu(y)\\
 &=\|f\|_1^p\|g\|_p^p.
 \end{aligned}
\]
Po wyciągnięciu pierwiastka stopnia $p$ otrzymujemy
żądaną nierówność.

Z gęstości $C_c$ w obu przestrzeniach Lebesgue'a
(por. \atlasref{ME-20::thm:cc-dense-lp-radon-20}{ME-20: gęstość})
wybieramy $f_n\to f$ w $L^1$ i $g_n\to g$ w $L^p$,
gdzie $f_n,g_n\in C_c$. Mamy
\[
 \|f_n*g_n-f_m*g_m\|_p
 \le\|f_n-f_m\|_1\|g_n\|_p+
 \|f_m\|_1\|g_n-g_m\|_p.
\]
Oba ciągi norm są ograniczone, a różnice w odpowiednich
normach dążą do zera. Z zupełności $L^p$ wynika istnienie
granicy. Dla innych ciągów przybliżających
$\widetilde f_n,\widetilde g_n$ analogiczne oszacowanie
\[
 \|f_n*g_n-\widetilde f_n*\widetilde g_n\|_p
 \le\|f_n-\widetilde f_n\|_1\|g_n\|_p+
 \|\widetilde f_n\|_1\|g_n-\widetilde g_n\|_p
\]
daje niezależność granicy. Liniowość w każdym
argumencie i nierówność norm przechodzą do granicy.
Każde ograniczone rozszerzenie musi przyjmować te
same granice, co dowodzi jednoznaczności.
Dla $p=1$ oba rozszerzenia zgadzają się na gęstej
podprzestrzeni $C_c\times C_c$, więc są równe.
Dla funkcji rzeczywistych dowód jest identyczny.
\end{proof}

\begin{uwaga}
W grupie nieunimodularnej nie wolno zamieniać ról
$f$ i $g$ w nierówności Younga bez dodatkowego
czynnika modularnego. Ponadto dla ogólnej grupy
lokalnie zwartej powyższe przedłużenia definiujemy
w normach przestrzeni Lebesgue'a, a nie przez
nieuzasadnione użycie globalnego twierdzenia Fubiniego.
\end{uwaga}

\end{document}
